\RequirePackage{fix-cm}
\documentclass[11pt,twoside]{book}

% Format Grand In-Quarto (240 x 340 mm, plus grand que le A4 210 x 297 mm)
% Marges réduites et harmonieuses : bloc de texte de 21.2 cm x 31.0 cm
\usepackage[
  paperwidth=240mm,
  paperheight=340mm,
  twoside,
  inner=15mm,
  outer=13mm,
  top=15mm,
  bottom=15mm,
  headheight=22pt,
  headsep=4.5mm,
  footskip=6.5mm
]{geometry}

\usepackage{graphicx}
\usepackage{fontspec}
\usepackage{calc}
\usepackage{fancyhdr}
\usepackage{hyperref}
\usepackage{amsmath}
\usepackage{amssymb}
\usepackage{cancel}
\usepackage{quoting}
\usepackage{ulem}
\usepackage{xcolor}
\usepackage{xspace}
\usepackage{adjustbox}
\usepackage{array}
\usepackage{longtable}
\usepackage{booktabs}
\usepackage{multicol}

% Palette typographique d'édition humaniste (encre sépia, rubrique rouge, filet or ancien)
\definecolor{inkDark}{HTML}{18120E}
\definecolor{rubricRed}{HTML}{6B1D1D}
\definecolor{sepiaNote}{HTML}{523A28}
\definecolor{ruleGold}{HTML}{8C7355}
\definecolor{bgBanner}{HTML}{F7F2E7}
\definecolor{bgScholion}{HTML}{FCF9F2}
\definecolor{framePlate}{HTML}{D6C7B2}

\pagecolor{white}
\color{inkDark}

% Polices : Libertinus Serif pour le français et l'apparat ; GFS Didot pour le grec ancien
\setmainfont{Libertinus Serif}[
  Ligatures=TeX,
  Numbers=Lining
]
\setsansfont{Libertinus Sans}[Ligatures=TeX]
\setmonofont{Libertinus Mono}[Scale=MatchLowercase, BoldFont={Libertinus Mono}]

\newfontfamily{\policeFr}{Libertinus Serif}[
  Ligatures=TeX,
  Numbers=Lining
]
\newfontfamily{\policeGr}{GFSDidot.otf}[
  Path=/usr/share/fonts/opentype/didot/,
  BoldFont=GFSDidotBold.otf,
  ItalicFont=GFSDidotItalic.otf,
  BoldItalicFont=GFSDidotBoldItalic.otf,
  Scale=MatchLowercase,
  SmallCapsFont=GFSDidot.otf,
  BoldFeatures={SmallCapsFont=GFSDidotBold.otf}
]

% Bascule automatique Latin <-> Grec via ucharclasses
\usepackage[Latin,Greek]{ucharclasses}
\setTransitionsForLatin{\policeFr}{}
\setTransitionsForGreek{\policeGr}{\policeFr}
% Empêcher la bascule vers Libertinus Serif entre GreekExtended (1F00-1FFF) et GreekAndCoptic (0370-03FF)
\XeTeXinterchartoks\GreekExtendedClass\GreekAndCopticClass={}
\XeTeXinterchartoks\GreekAndCopticClass\GreekExtendedClass={}
\XeTeXcharclass"0300=\GreekAndCopticClass
\XeTeXcharclass"0301=\GreekAndCopticClass
\XeTeXcharclass"0304=\GreekAndCopticClass
\XeTeXcharclass"0306=\GreekAndCopticClass
\XeTeXcharclass"0308=\GreekAndCopticClass
\XeTeXcharclass"0313=\GreekAndCopticClass
\XeTeXcharclass"0314=\GreekAndCopticClass
\XeTeXcharclass"0342=\GreekAndCopticClass
\XeTeXcharclass"0345=\GreekAndCopticClass
\XeTeXcharclass"1FBD=\GreekAndCopticClass
\XeTeXcharclass"1FBF=\GreekAndCopticClass
\XeTeXcharclass"1FFE=\GreekAndCopticClass

% Symboles paléographiques et critiques absents de GFS Didot -> Libertinus Serif
\usepackage{newunicodechar}
\newcommand{\fbchar}[1]{{\XeTeXinterchartokenstate=0\policeFr #1}}
\newunicodechar{Ϲ}{\fbchar{Ϲ}}
\newunicodechar{ϲ}{\fbchar{ϲ}}
\newunicodechar{ϗ}{\fbchar{ϗ}}
\newunicodechar{Ϛ}{\fbchar{Ϛ}}
\newunicodechar{ϛ}{\fbchar{ϛ}}
\newunicodechar{⊢}{\fbchar{⊢}}
\newunicodechar{⊣}{\fbchar{⊣}}
\newunicodechar{∞}{\fbchar{∞}}
\newunicodechar{※}{\fbchar{※}}
\newunicodechar{☉}{\fbchar{☉}}
\newunicodechar{⟦}{\fbchar{⟦}}
\newunicodechar{⟧}{\fbchar{⟧}}

\usepackage{polyglossia}
\setdefaultlanguage{french}
\usepackage[maxlevel=3]{csquotes}

% Garantir la transition vers GFS Didot après la ponctuation française de polyglossia ( ( [ { : ; « etc. )
\AtBeginDocument{%
  \count255=1
  \loop\ifnum\count255<255
    \ifnum\count255=\GreekAndCopticClass\else
      \ifnum\count255=\GreekExtendedClass\else
        \XeTeXinterchartoks \count255 \GreekAndCopticClass = {\policeGr}%
        \XeTeXinterchartoks \count255 \GreekExtendedClass = {\policeGr}%
        \XeTeXinterchartoks \GreekAndCopticClass \count255 = {\policeFr}%
        \XeTeXinterchartoks \GreekExtendedClass \count255 = {\policeFr}%
      \fi
    \fi
    \advance\count255 by 1
  \repeat
  \XeTeXinterchartoks 0 \GreekAndCopticClass = {\policeGr}%
  \XeTeXinterchartoks 0 \GreekExtendedClass = {\policeGr}%
  \XeTeXinterchartoks 4095 \GreekAndCopticClass = {\policeGr}%
  \XeTeXinterchartoks 4095 \GreekExtendedClass = {\policeGr}%
  \DeclareRobustCommand{\textlbrackdbl}{\fbchar{⟦}}%
  \DeclareRobustCommand{\textrbrackdbl}{\fbchar{⟧}}%
}

\hypersetup{
  colorlinks=true,
  linkcolor=rubricRed,
  urlcolor=rubricRed,
  citecolor=sepiaNote,
  pdftitle={Atlas et Edition des Scholies Marginales, Gloses Interlineaires et Collation Diplomatique du Codex Ravennas 429 - Les Thesmophories (ff. 162v-177r)},
  pdfauthor={Dossier Philologique et Paleographique - Edition Louis Hauseux (Bibi)}
}

% Tolérance typographique pour garantir zéro Overfull \hbox sur 21.2 cm
\tolerance=2000
\emergencystretch=3em
\hbadness=10000
\hfuzz=0.5pt

% En-têtes et pieds de page courants
\pagestyle{fancy}
\fancyhf{}
\newcommand{\currHeadLE}{ΑΡΙΣΤΟΦΑΝΟΥΣ ΘΕΣΜΟΦΟΡΙΑΖΟΥΣΑΙ · SCHOLIES ({\policeGr Σ}\expo{R}) \& GLOSES DU CODEX RAVENNAS 429}
\newcommand{\currHeadRE}{Atlas \& Édition Folio par Folio (ff. 162v–177r)}
\newcommand{\currHeadLO}{ATLAS DES SCHOLIES ({\policeGr Σ}\expo{R}), GLOSES INTERLINÉAIRES \& COLLATION DE R}
\newcommand{\currHeadRO}{Édition Louis Hauseux (Bibi)}
\fancyhead[LE]{\footnotesize\color{sepiaNote}\currHeadLE}
\fancyhead[RE]{\footnotesize\color{sepiaNote}\currHeadRE}
\fancyhead[LO]{\footnotesize\color{sepiaNote}\currHeadLO}
\fancyhead[RO]{\footnotesize\color{sepiaNote}\currHeadRO}
\fancyfoot[LE,RO]{\small\textbf{\color{rubricRed}\thepage}}
\fancyfoot[RE,LO]{\scriptsize\color{sepiaNote}\textit{Aristophane · Les Thesmophories — Atlas des Scholies ({\policeGr Σ}\expo{R}), Gloses \& Collation du Codex Ravennas 429 (ff. 162v–177r)}}
\renewcommand{\headrulewidth}{0.4pt}
\renewcommand{\footrulewidth}{0.3pt}
\renewcommand{\headrule}{\hbox to\headwidth{\color{ruleGold}\leaders\hrule height \headrulewidth\hfill}}
\renewcommand{\footrule}{\hbox to\headwidth{\color{ruleGold}\leaders\hrule height \footrulewidth\hfill}}

\setlength{\columnsep}{12pt}
\setlength{\multicolsep}{2pt}
\setlength{\parindent}{0pt}

% Macros philologiques, paléographiques et métriques
\newcommand{\Petit}{\fontsize{9pt}{11pt}\selectfont}
\newcommand{\Taille}[1]{\fontsize{#1pt}{12pt}\selectfont}
\newcommand{\TaillePt}[1]{\fontsize{#1}{15pt}\selectfont}
\newcommand{\Automatique}[1]{#1}
\newcommand{\renverse}[1]{\rotatebox[origin=c]{180}{#1}}
\newcommand{\italique}[1]{\textit{#1}}
\newcommand{\citComm}[1]{\begin{quoting}[vskip=1.2pt,leftmargin=0.9em,rightmargin=0.3em]#1\end{quoting}}
\newcommand{\espace}[1]{\hspace*{#1}}
\newcommand{\cesure}{ [...] }
\newcommand{\nieme}[1]{#1\textsuperscript{e}}
\providecommand{\ieme}{}\renewcommand{\ieme}{\textsuperscript{e}\xspace}
\providecommand{\ier}{}\renewcommand{\ier}{\textsuperscript{er}\xspace}
\DeclareRobustCommand{\textlbrackdbl}{\fbchar{⟦}}
\DeclareRobustCommand{\textrbrackdbl}{\fbchar{⟧}}
\newcommand{\lettreBreve}[1]{\leavevmode\vbox{\offinterlineskip\ialign{\hfil##\hfil\cr\fontsize{5.5pt}{6pt}\selectfont{\policeFr ˘}\cr\noalign{\kern0.05ex}#1\cr}}}
\newcommand{\widearc}[1]{\leavevmode\vbox{\offinterlineskip\ialign{\hfil##\hfil\cr\resizebox{\widthof{#1}}{0.35ex}{\rotatebox[origin=c]{180}{\policeFr ˘}}\cr\noalign{\kern0.1ex}#1\cr}}}
\newcommand{\synizese}[1]{\widearc{#1}}
\newcommand{\auDessus}[2]{\leavevmode\vbox{\offinterlineskip\ialign{\hfil##\hfil\cr\fontsize{5.5pt}{6pt}\selectfont #2\cr\noalign{\kern0.1ex}#1\cr}}}
\newcommand{\inRas}[2]{\if\relax\detokenize{#2}\relax\sout{#1}\else\auDessus{\sout{#1}}{#2}\fi}
\newcommand{\expo}[1]{\textsuperscript{#1}}
\newcommand{\alinea}{\hspace{1.2em}}

% Sigles des manuscrits, papyrus et lexicographes
\newcommand{\R}{\textbf{R}}
\newcommand{\Run}{R\expo{1}}
\newcommand{\Rdeux}{R\expo{2}}
\newcommand{\Rtrois}{R\expo{3}}
\newcommand{\Rpc}{R\expo{pc}}
\newcommand{\Rac}{R\expo{ac}}
\newcommand{\Mon}{M}
\newcommand{\PsiUn}{PSI-1194}
\newcommand{\PsiDeux}{POxy-3839}
\newcommand{\PsiTrois}{POxy-3840}
\newcommand{\PsiQuatre}{POxy-4935}
\newcommand{\PsiCinq}{POxy-1176}
\newcommand{\Fr}{Ar. \textit{fragm.} (Kassel-Austin) }
\newcommand{\Syn}{Συναγωγή }
\newcommand{\SynB}{Συναγωγή\expo{B} }
\newcommand{\ClemAlex}{Clemens Alexandrinus }
\newcommand{\Eustathe}{Eustathius }
\newcommand{\AelDion}{Aelius Dionysius }
\newcommand{\EtymMagnum}{\textit{Etymologicum Magnum} }
\newcommand{\EtymGen}{\textit{Etymologicum Genuinum} }
\newcommand{\EtymSym}{\textit{Etymologicum Symeonis} }
\newcommand{\Apostol}{Michael Apostolius }
\newcommand{\Macaire}{Macarius }
\newcommand{\Phryn}{Phrynichus \textit{Praep. soph.} }
\newcommand{\Stobee}{Joannes Stobaeus }
\newcommand{\EpimAlphab}{\textit{Epimerismi alphabetici in Homerum} (ed. Cramer) }
\newcommand{\GregCor}{Gregorius Corinthius }
\newcommand{\Antiatt}{\textit{Antiatticista} (ed. Bekker) }
\newcommand{\Harp}{Harpocration }
\newcommand{\Nauck}{N\expo{2}}
\newcommand{\LSJ}{LSJ}
\newcommand{\cf}{\textit{cf.}\xspace}
\newcommand{\chap}[1]{\^{#1}}

\newcommand{\francais}[1]{{\policeFr #1}}
\newcommand{\grec}[1]{{\policeGr #1}}
\newcommand{\latin}[1]{\textit{#1}}
\newcommand{\allemand}[1]{\textit{#1}}
\newcommand{\anglais}[1]{\textit{#1}}
\newcommand{\italien}[1]{\textit{#1}}
\newcommand{\espagnol}[1]{\textit{#1}}
\newcommand{\didascalie}[1]{{\policeFr\itshape\color{sepiaNote}#1}}

\newlength{\larg}
\newcommand{\decalage}[1]{\settowidth{\larg}{#1}\hspace*{\larg}}
\newcommand{\decalageN}[1]{\settowidth{\larg}{#1}\hspace*{-\larg}}

\newcommand{\lecon}[1]{\textbf{#1}}
\newcommand{\douteux}[1]{\leavevmode\vtop{\offinterlineskip\ialign{\hfil##\hfil\cr#1\cr\noalign{\kern0.15ex}\fontsize{6pt}{6pt}\selectfont .\cr}}}
\newcommand{\difference}[1]{\dotuline{#1}}
\newcommand{\espApp}{\ {\color{ruleGold}{\policeFr ∥}}\ }
\newcommand{\espTest}{\ {\color{ruleGold}{\policeFr ∥}}\ }
\newcommand{\ligneApp}[1]{\ \ {\color{rubricRed}\textbf{#1}}\ }
\newcommand{\marqueur}[1]{\textsuperscript{\color{rubricRed}\textbf{#1}}}
\newcommand{\cat}{\texorpdfstring{\textsubscript{\textasciicircum}}{cat}}

% Bandeau de titre en haut de la page de collation diplomatique de R
\newcommand{\bandeauRav}[4]{%
  \noindent\fcolorbox{ruleGold}{bgBanner}{%
    \parbox{\dimexpr\textwidth-2\fboxsep-2\fboxrule\relax}{%
      \vspace{1.5pt}%
      \hbox to \linewidth{\adjustbox{max width=0.56\linewidth}{{\small\scshape\bfseries\color{rubricRed}Codex Ravennas 429 (\R) — #1 (#2) · Vers #3}}\hfill\adjustbox{max width=0.42\linewidth}{{\footnotesize\itshape\color{sepiaNote}#4}}}%
      \vspace{1.5pt}%
    }%
  }\par\vspace{3pt}%
}

% Bandeau de titre de la section Scholies & Gloses de chaque folio
\newcommand{\bandeauScholies}[4]{%
  \noindent\fcolorbox{ruleGold}{bgBanner}{%
    \parbox{\dimexpr\textwidth-2\fboxsep-2\fboxrule\relax}{%
      \vspace{2pt}%
      \hbox to \linewidth{\adjustbox{max width=0.60\linewidth}{{\small\scshape\bfseries\color{rubricRed}Scholies Marginales ({\policeGr Σ}\expo{R}) \& Gloses Interlinéaires — #1 (#2) · #3}}\hfill\adjustbox{max width=0.38\linewidth}{{\footnotesize\itshape\color{sepiaNote}#4}}}%
      \vspace{2pt}%
    }%
  }\par\vspace{4pt}%
}

% Ligne de transcription diplomatique corrigée de R + apparat critique en ligne (largeur 21.2 cm)
\newcommand{\ligneRav}[4]{%
  \noindent\hbox to \textwidth{%
    \hbox to 0.70cm{\hfill\footnotesize\color{sepiaNote}\adjustbox{max width=0.56cm}{#1}\hspace{0.12cm}}%
    \hbox to 1.45cm{\hfill\footnotesize\color{rubricRed}\adjustbox{max width=1.25cm}{#2}\hspace{0.18cm}}%
    \hbox to 10.35cm{\adjustbox{max width=10.25cm}{\policeGr #3}\hfill}%
    \hspace{0.20cm}%
    \parbox[t]{8.50cm}{\raggedright\fontsize{7.8pt}{9.1pt}\selectfont #4}%
  }\par\vspace{0.45pt}%
}

% Bloc structuré pour chaque scholie ou glose (Σ 001 à Σ 386)
\newcommand{\entreeScholie}[7]{%
  \par\vspace{3.5pt}\noindent\fcolorbox{framePlate}{bgScholion}{%
    \parbox{\dimexpr\textwidth-2\fboxsep-2\fboxrule\relax}{%
      \vspace{2pt}%
      {\fontsize{9.5pt}{11.2pt}\selectfont
      \noindent{\bfseries\color{rubricRed}{\policeGr Σ}~#1 — v.~#2}\hspace{0.6em}{\color{ruleGold}\textbullet}\hspace{0.6em}{\bfseries\color{sepiaNote}Lemme :} {\policeGr #3}\hspace{0.6em}{\color{ruleGold}\textbullet}\hspace{0.6em}{\itshape\color{sepiaNote}#4}\par\vspace{2pt}%
      {\color{ruleGold}\hrule height 0.3pt}\vspace{2.5pt}%
      \noindent\textbf{\color{rubricRed}Texte grec ({\policeGr Σ}\expo{R}) :} {\policeGr #5}\par\vspace{1.8pt}%
      \noindent\textbf{\color{sepiaNote}Tachygraphie de \Rdeux{} :} {\policeGr #6}\par\vspace{1.8pt}%
      \noindent\textbf{\color{inkDark}Traduction française :} #7\par}%
      \vspace{2pt}%
    }%
  }\par\nopagebreak\vspace{1.5pt}%
}

\newcommand{\commScholie}[1]{%
  \noindent\hspace*{0.3cm}\parbox{\dimexpr\textwidth-0.6cm\relax}{%
    \fontsize{8.8pt}{10.5pt}\selectfont\color{sepiaNote}\textbf{\color{rubricRed}Commentaire philologique (\textbf{Bibi}) :} #1%
  }\par\vspace{3.5pt}%
}

\begin{document}

% ===========================================================================
% PAGE DE TITRE MONUMENTALE (FORMAT GRAND IN-QUARTO 240 x 340 MM)
% ===========================================================================
\thispagestyle{empty}
\vspace*{0.8cm}
\begin{center}
  {\color{ruleGold}\rule{\textwidth}{1.4pt}}\par\vspace{2.5pt}
  {\color{ruleGold}\rule{\textwidth}{0.5pt}}\par\vspace{0.7cm}

  {\fontsize{26pt}{32pt}\selectfont\bfseries\color{rubricRed}\policeGr ΣΧΟΛΙΑ ΕΙΣ ΤΑΣ ΘΕΣΜΟΦΟΡΙΑΖΟΥΣΑΣ}\par\vspace{0.35cm}
  {\fontsize{21pt}{26pt}\selectfont\scshape\bfseries ATLAS \& ÉDITION INTÉGRALE DES SCHOLIES MARGINALES ({\policeGr Σ}\expo{R}),\\GLOSES INTERLINÉAIRES \& COLLATION DIPLOMATIQUE\\DU CODEX RAVENNAS 429 (FF. 162v–177r)}\par\vspace{0.35cm}
  {\large\itshape\color{sepiaNote}Édition critique, diplomatique, tachygraphique, traduite et commentée des 386 scholies et gloses de \textbf{R\expo{2}}\\et des 1\,238 lignes poétiques de \textbf{R\expo{1}} pour \textbf{Les Thesmophories} d'Aristophane (vv. 1–1231)}\par\vspace{0.22cm}
  {\normalsize\bfseries\color{rubricRed}Établie en accompagnement de l'édition critique et commentée de Louis Hauseux (sigle d'apparat : \textbf{Bibi})\\Format Grand In-Quarto (240~×~340~mm) — Libertinus Serif \& GFS Didot}\par\vspace{0.6cm}

  {\color{ruleGold}\rule{0.50\textwidth}{0.6pt}}\par\vspace{0.6cm}
\end{center}

\noindent\fcolorbox{ruleGold}{bgBanner}{%
  \parbox{\dimexpr\textwidth-2\fboxsep-2\fboxrule\relax}{%
    \vspace{5pt}
    \begin{center}
      {\large\scshape\bfseries\color{rubricRed}Architecture et contenu de l'Atlas scholiographique et diplomatique (ff. 162v–177r)}
    \end{center}
    \vspace{3pt}
    {\fontsize{9.6pt}{11.8pt}\selectfont
    \begin{multicols}{2}
      \noindent\textbf{\color{rubricRed}Première Partie — Monographie sur le corpus {\policeGr Σ}\expo{R} :}
      \begin{itemize}
        \item \textbf{1. Statut unique du \textit{Codex Ravennas 429} (\R{}) :} \textit{Codex unicus} tant pour le texte poétique (\Run{}) que pour le corpus scholiographique (\Rdeux{}), intégralement omis par l'apographe \Mon{} (\textit{Monacensis gr. 492}).
        \item \textbf{2. Topographie des 30 folios (\texttt{ff. 162v–177r}) :} Alternance des marges extérieures (gauche aux versos, droite aux rectos), marges supérieures et inférieures, gloses interlinéaires (\textit{supra lineam}), \textit{parepigraphai} scéniques et gloses intruses.
        \item \textbf{3. Tachygraphie de la Main~B (\Rdeux{}) :} Décodage systématique des compendia casuels, suspensions, finales en exposant, signes critiques (χ, \textit{σημεῖον}, \textit{διχῶς}, \textit{γράφεται}) et \textit{notae personarum}.
        \item \textbf{4. Stratigraphie alexandrine et lexiques anciens :} D'Aristarque, Callimaque, Callistrate, Ératosthène et Didyme à Juba~II, Séleucos et Symmaque ; confrontation avec Hésychius, Photius, la \textit{Souda} et Zonaras dans l'édition de \textbf{Louis Hauseux (\textbf{Bibi})}.
      \end{itemize}
      \columnbreak
      \noindent\textbf{\color{rubricRed}Deuxième Partie — Atlas \& Édition Folio par Folio (01 à 30) :}
      \begin{itemize}
        \item \textbf{Collation diplomatique intégrale ligne par ligne :} Les 1\,238 lignes des 30 folios de \R{} (\texttt{f. 162v} à \texttt{f. 177r}, vv. 1–1231) avec \textit{notae personarum}, \textit{paragraphoi}, \textit{dikola}, grattages (\textit{in rasura}), \textit{Loci paralleli} et apparat critique en regard.
        \item \textbf{Corpus exhaustif des 386 scholies et gloses (\texttt{{\policeGr Σ} 001} à \texttt{{\policeGr Σ} 386}) :}
        \begin{enumerate}
          \item Référence séquentielle, vers et lemme grec dans \R{} et \textbf{Bibi} ;
          \item Position topographique exacte sur le folio ;
          \item Texte grec intégral de la scholie ou glose en \textit{GFS Didot} ;
          \item Tracé diplomatique et abréviations tachygraphiques de \Rdeux{} ;
          \item Traduction française intégrale en \textit{Libertinus Serif} ;
          \item Commentaire philologique détaillé reliant la scholie aux choix critiques de \textbf{Louis Hauseux (\textbf{Bibi})}.
        \end{enumerate}
      \end{itemize}
    \end{multicols}
    }
    \vspace{4pt}
  }%
}

\vfill
\begin{center}
  {\color{ruleGold}\rule{\textwidth}{0.5pt}}\par\vspace{2pt}
  {\color{ruleGold}\rule{\textwidth}{1.4pt}}\par\vspace{0.2cm}
  {\footnotesize\scshape\color{sepiaNote}Biblioteca Classense de Ravenne (Ms. 429) · Fac-similé phototypique de 1904 (BnF, FOL-FAC SIM-182 (9)) · Édition critique Louis Hauseux (Bibi)}
\end{center}
\newpage

% ===========================================================================
% PREMIÈRE PARTIE : ÉTUDE CODICOLOGIQUE, PALÉOGRAPHIQUE, TACHYGRAPHIQUE ET CRITIQUE
% ===========================================================================
\renewcommand{\currHeadLE}{PREMIÈRE PARTIE · ÉTUDE DU CORPUS DES SCHOLIES ({\policeGr Σ}\expo{R}) ET GLOSES}
\renewcommand{\currHeadRE}{Topographie, Tachygraphie de R\expo{2}, Stratigraphie alexandrine \& Édition Bibi}
\renewcommand{\currHeadLO}{PREMIÈRE PARTIE · ÉTUDE DU CORPUS DES SCHOLIES ({\policeGr Σ}\expo{R}) ET GLOSES}
\renewcommand{\currHeadRO}{Topographie, Tachygraphie de R\expo{2}, Stratigraphie alexandrine \& Édition Bibi}

\bandeauScholies{Première Partie}{Étude méthodologique et critique}{ff. 162v–177r}{Codicologie, topographie marginale, tachygraphie de R\expo{2} et critique textuelle (Bibi)}

\vspace{2pt}
{\fontsize{10pt}{12.2pt}\selectfont
\noindent{\Large\bfseries\color{rubricRed}1. Statut unique du corpus {\policeGr Σ}\expo{R} et topographie des 30 folios (\texttt{ff. 162v–177r})}\par\vspace{4pt}

\noindent Pour \textit{Les Thesmophories} d'Aristophane (\texttt{ff. 162v–177r}, pagination moderne au crayon \texttt{164v–179r}, vv.~1–1231), le \textbf{\textit{Codex Ravennas 429}} (\R{}, Constantinople, vers 950–970) est l'\textbf{unique manuscrit médiéval} qui nous ait transmis les scholies anciennes ({\policeGr Σ}\expo{R}) et les gloses interlinéaires. En effet, le \textit{Monacensis gr. 492} (\Mon{}, XV\ieme{} siècle), copie indirecte du \textit{Ravennas}, a délibérément laissé de côté l'\textit{Hypothesis} d'Aristophane de Byzance, la liste des personnages, les 386 scholies et gloses de la Main~B (\Rdeux{}), ainsi que les vers rétablis en marge par \Rdeux{} (le \textbf{v.~33} au \texttt{f. 162v} et les \textbf{vv.~1102–1103} dans la marge inférieure du \texttt{f. 175v}).

\par\vspace{4pt}
\noindent L'organisation matérielle du commentaire de \Rdeux{} obéit rigoureusement à l'alternance des pages \textbf{versos} (pages de gauche, où la grande marge extérieure est \textbf{à gauche}) et des pages \textbf{rectos} (pages de droite, où la grande marge extérieure est \textbf{à droite}). Le tableau récapitulatif ci-dessous dresse l'inventaire complet des \textbf{30 folios}, des \textbf{1\,238 lignes poétiques} et des \textbf{386 unités scholiographiques et glosématiques} (\texttt{{\policeGr Σ} 001} à \texttt{{\policeGr Σ} 386}) :

\par\vspace{4pt}
{\fontsize{8.1pt}{9.5pt}\selectfont
\noindent\begin{adjustbox}{max totalheight=23.5cm, max width=\textwidth, keepaspectratio}\setlength{\tabcolsep}{4pt}\begin{tabular}{@{}c l l c c c c l p{12.2cm}@{}}
\toprule
\textbf{N°} & \textbf{Folio anc.} & \textbf{Classense} & \textbf{R°/V°} & \textbf{Vers} & \textbf{Lignes} & \textbf{{\policeGr Σ}\expo{R}} & \textbf{Numéros} & \textbf{Pièces maîtresses du commentaire {\policeGr Σ}\expo{R} et faits paléographiques} \\
\midrule
\textbf{01} & \textbf{f. 162v} & 164v & V° (g.) & Hyp., 1–42 & 42 & \textbf{21} & 001–021 & \textit{Hypothesis} ({\policeGr Μνησίλοχος}, {\policeGr Χοιρίνη}), {\policeGr Σ} 1 (hirondelle), {\policeGr Σ} 3 ({\policeGr οἷόν τε}), {\policeGr Σ} 18 ({\policeGr λείπει ὡς}), {\policeGr Σ} 21 (\textit{Ajax le Locrien} de Sophocle, Didyme), {\policeGr Σ} 31 (Aristarque et Didyme), {\policeGr Σ} 32 (416 av. J.-C.), {\policeGr Σ} 39 ({\policeGr λαός}) \\
\textbf{02} & \textbf{f. 163r} & 165r & R° (dr.) & 43–79 & 37 & \textbf{17} & 022–038 & {\policeGr Σ} 50 ({\policeGr λείπει ἕτοιμός ἐστιν}), {\policeGr Σ} 52–57 (les 7 métaphores artisanales d'Agathon), {\policeGr Σ} 68 ({\policeGr δυσκολοκάμπτους}, \textit{Nuées} 971), {\policeGr Σ} 71 et {\policeGr Σ} 74 ({\policeGr τὸ σημεῖον}, sens de {\policeGr κηδεστής}) \\
\textbf{03} & \textbf{f. 163v} & 165v & V° (g.) & 80–122 & 44 & \textbf{13} & 039–051 & {\policeGr Σ} 80 (grande scholie de Callimaque sur le 3\ieme{} jour / jour du milieu des Thesmophories), {\policeGr Σ} 94 ({\policeGr πυραμοῦς}), {\policeGr Σ} 96 (\textit{ekkyklêma}), {\policeGr Σ} 98 (Cyrène), {\policeGr Σ} 100 (\textit{parepigraphê} {\policeGr μινυρισμός}), {\policeGr Σ} 101 (monodie dialoguée) \\
\textbf{04} & \textbf{f. 164r} & 166r & R° (dr.) & 123–164 & 44 & \textbf{22} & 052–073 & {\policeGr Σ} 125 ({\policeGr δοκίμῳ}), {\policeGr Σ} 129 (\textit{parepigraphê} {\policeGr ὀλολύζει}), {\policeGr Σ} 135 (la tétralogie de la \textit{Lycurgie} d'Eschyle), {\policeGr Σ} 137 (\textit{Denys} d'Euboule), {\policeGr Σ} 162 (polémique contre Didyme sur Alcée / Achaïos) \\
\textbf{05} & \textbf{f. 164v} & 166v & V° (g.) & 165–207 & 43 & \textbf{8} & 074–081 & {\policeGr Σ} 168 (Philoclès \textit{Korydos}, fr. de Téléclide), {\policeGr Σ} 179 (glose {\policeGr μεγάλῃ καὶ θαυμαστῇ} prouvant {\policeGr καινῇ} contre {\policeGr κοινῇ}), {\policeGr Σ} 190 (mort d'Euripide 6 ans plus tard), {\policeGr Σ} 194 (\textit{Alceste} 691) \\
\textbf{06} & \textbf{f. 165r} & 167r & R° (dr.) & 208–251 & 44 & \textbf{13} & 082–094 & {\policeGr Σ} 215 (\textit{Idéens} de Cratinos), {\policeGr Σ} 220 ({\policeGr ξυροδόκης} sur {\policeGr δουροδόκη}), {\policeGr Σ} 224 (sanctuaire des Érinyes), {\policeGr Σ} 231 ({\policeGr μύζειν}), {\policeGr Σ} 236 ({\policeGr δελφάκια}), {\policeGr Σ} 242 ({\policeGr οἰκίαν}), {\policeGr Σ} 246 ({\policeGr τράμις}) \\
\textbf{07} & \textbf{f. 165v} & 167v & V° (g.) & 252–294 & 44 & \textbf{17} & 095–111 & {\policeGr Σ} 254 ({\policeGr ποσθίου}), {\policeGr Σ} 258 (\textit{Charlatans} d'Aristomène), {\policeGr Σ} 261 ({\policeGr ἔγκυκλον} vs {\policeGr κροκωτός}), {\policeGr Σ} 272 (\textit{Mélanippe}), {\policeGr Σ} 273 ({\policeGr Ὑποκράτους}), {\policeGr Σ} 277 (\textit{parepigraphê}), {\policeGr Σ} 289 ({\policeGr Χοιρίλ'}), {\policeGr Σ} 291 ({\policeGr θυλάκιον}) \\
\textbf{08} & \textbf{f. 166r} & 168r & R° (dr.) & 295–336 & 42 & \textbf{10} & 112–121 & {\policeGr Σ} 295 ({\policeGr γυνὴ μιμουμένη κήρυκα}, grattage de {\policeGr Κῆρυξ}), {\policeGr Σ} 297–300 (gloses intruses {\policeGr τῇ Δήμητρῃ...} et prologue de Kalligéneia dans les \textit{Thesmophories II}), {\policeGr Σ} 336 ({\policeGr φιλίαν σπένδει}) \\
\textbf{09} & \textbf{f. 166v} & 168v & V° (g.) & 337–379 & 42 & \textbf{14} & 122–135 & {\policeGr Σ} 339 (imprécation civique contre les Pisistratides et Hippias), {\policeGr Σ} 346 (variante {\policeGr ἑταίρᾳ} avec iota adoptée par \textbf{Bibi}), {\policeGr Σ} 365 (les Mèdes), {\policeGr Σ} 373 ({\policeGr Τιμόκλεια} épistate) \\
\textbf{10} & \textbf{f. 167r} & 169r & R° (dr.) & 380–422 & 43 & \textbf{18} & 136–153 & {\policeGr Σ} 380 (couronne de l'orateur), {\policeGr Σ} 389 ({\policeGr ἐπισμῇ}, Diodore et \textit{Cléoboulines} de Cratinos), {\policeGr Σ} 393 (Symmaque et la vraie leçon {\policeGr οἰνοπίπας}), {\policeGr Σ} 395 ({\policeGr ἴκρια}), {\policeGr Σ} 401 (Sappho fr. 125), {\policeGr Σ} 404 (\textit{Sthénébée}) \\
\textbf{11} & \textbf{f. 167v} & 169v & V° (g.) & 423–465 & 49 & \textbf{12} & 154–165 & {\policeGr Σ} 423 (clés laconiennes à trois dents, \textit{Misoumenos} de Ménandre), {\policeGr Σ} 427 ({\policeGr θριπήδεστα}), {\policeGr Σ} 441 (Xénoclès fils de Carcinos), {\policeGr Σ} 458 ({\policeGr συνθηματιαίους} = {\policeGr ἐκδοτικούς}) \\
\textbf{12} & \textbf{f. 168r} & 170r & R° (dr.) & 466–508 & 43 & \textbf{14} & 166–179 & {\policeGr Σ} 469 (discours de Mnésiloque), {\policeGr Σ} 481 ({\policeGr ἔκνυεν}), {\policeGr Σ} 486 ({\policeGr σφάκος} = sauge), {\policeGr Σ} 487 (eau sur le gond), {\policeGr Σ} 489 (Apollon Aguieus), {\policeGr Σ} 491 ({\policeGr σταβλῖται / βουρδωνάριοι}), {\policeGr Σ} 500 (glose), {\policeGr Σ} 506 (miel des nouveau-nés, Ménandre) \\
\textbf{13} & \textbf{f. 168v} & 170v & V° (g.) & 509–550 & 44 & \textbf{9} & 180–188 & {\policeGr Σ} 509 ({\policeGr τὸ ἦτρον τῆς χύτρας}), {\policeGr Σ} 516 ({\policeGr κύτταρος} selon Théophraste et Ératosthène), {\policeGr Σ} 519 (\textit{Télèphe}), {\policeGr Σ} 529 (\textit{Scolie} de Praxilla), {\policeGr Σ} 531 et 547 (notes métriques d'Héliodore), {\policeGr Σ} 533 (serment par Agraulos) \\
\textbf{14} & \textbf{f. 169r} & 171r & R° (dr.) & 551–589 & 39 & \textbf{15} & 189–203 & {\policeGr Σ} 557 ({\policeGr σιφωνίζομεν}), {\policeGr Σ} 560–564 (Clytemnestre et faits divers attiques), {\policeGr Σ} 567 ({\policeGr ποκάδας}, traité de Lycophron), {\policeGr Σ} 572 ({\policeGr ὁμοῦ} chez Ménandre), {\policeGr Σ} 575 (Clisthène), {\policeGr Σ} 585 (topographie de l'\textit{Anodos}) \\
\textbf{15} & \textbf{f. 169v} & 171v & V° (g.) & 590–634 & 45 & \textbf{12} & 204–215 & {\policeGr Σ} 600 (les {\policeGr κρύπται} à Sparte et à Thasos), {\policeGr Σ} 603 (nom de {\policeGr Μνησίλοχος}), {\policeGr Σ} 616 (le cresson chez les Perses, Xénophon), {\policeGr Σ} 620 (dème de Cothôcides), {\policeGr Σ} 624 (les tentes {\policeGr σκηναί}), {\policeGr Σ} 633 ({\policeGr σκάφιον} et {\policeGr ἀμίς}) \\
\textbf{16} & \textbf{f. 170r} & 172r & R° (dr.) & 635–678 & 45 & \textbf{8} & 216–223 & {\policeGr Σ} 641 ({\policeGr στερίφη}), {\policeGr Σ} 646–648 (le lazzi du phallus halé à travers l'Isthme de Corinthe, {\policeGr διισθμονίσαι}), {\policeGr Σ} 656 ({\policeGr ἐγκύκλων}), {\policeGr Σ} 658 (la Pnyx et les tentes) \\
\textbf{17} & \textbf{f. 170v} & 172v & V° (g.) & 679–720 & 40 & \textbf{10} & 224–233 & {\policeGr Σ} 682 (fragment d'Alexandros), {\policeGr Σ} 688 (parodie du \textit{Télèphe} : l'outre de vin déguisée en bébé), {\policeGr Σ} 697 (accentuation propérispomène attique {\policeGr τροπαῖον}), {\policeGr Σ} 701 ({\policeGr παροιμία}), {\policeGr Σ} 710 (Symmaque sur {\policeGr ἥκειν γ' ὅθεν ἥκεις}) \\
\textbf{18} & \textbf{f. 171r} & 173r & R° (dr.) & 721–764 & 43 & \textbf{15} & 234–248 & {\policeGr Σ} 728 (quantité longue de l'α de {\policeGr Μανία}), {\policeGr Σ} 730 et 734 ({\policeGr Κρητικόν} et {\policeGr Περσικαί} de l'outre), {\policeGr Σ} 735 ({\policeGr ὦ μέθυσοι} mal rattaché par la Souda), {\policeGr Σ} 742 (grossesse de 10 mois), {\policeGr Σ} 750 (calcul de l'âge), {\policeGr Σ} 754 ({\policeGr σφαγεῖον}), {\policeGr Σ} 760 ({\policeGr Μίκα}) \\
\textbf{19} & \textbf{f. 171v} & 173v & V° (g.) & 765–799 & 35 & \textbf{8} & 249–256 & {\policeGr Σ} 771 (l'intrigue d'Oïax et des rames dans le \textit{Palamède} d'Euripide), {\policeGr Σ} 773 (les tablettes votives {\policeGr πινάκια}), {\policeGr Σ} 779 ({\policeGr σμίλης ὁλκούς}), {\policeGr Σ} 781 ({\policeGr τουτὶ τὸ ῥῶ μοχθηρόν}), {\policeGr Σ} 785 (début de la Parabase) \\
\textbf{20} & \textbf{f. 172r} & 174r & R° (dr.) & 800–837 & 37 & \textbf{11} & 257–267 & {\policeGr Σ} 804 (défaite navale de Charminos en 412/411, Nausimaché et Salabakchô), {\policeGr Σ} 805 (Cléophon), {\policeGr Σ} 808 (les 2 fragments de l'\textit{Hyperbolos} de Platon le Comique sur l'\textit{épilachèsis}), {\policeGr Σ} 825–829 ({\policeGr κανών} et {\policeGr σκιάδειον}), {\policeGr Σ} 834 ({\policeGr Στηνίοισι} et {\policeGr Σκίροις}) \\
\textbf{21} & \textbf{f. 172v} & 174v & V° (g.) & 838–875 & 38 & \textbf{15} & 268–282 & {\policeGr Σ} 838 (coupe {\policeGr σκάφιον}), {\policeGr Σ} 840 (le \textit{Prôtatikos} de l'Homérique et Xénophantos père de Lamachos), {\policeGr Σ} 841 (mort de Lamachos en Sicile 4 ans plus tôt), {\policeGr Σ} 846 (Sophron fr. 64), {\policeGr Σ} 850–857 (\textit{Hélène} d'Euripide), {\policeGr Σ} 868 ({\policeGr (Γυνή)}), {\policeGr Σ} 870 (\textit{Pélée} de Sophocle) \\
\textbf{22} & \textbf{f. 173r} & 175r & R° (dr.) & 876–919 & 43 & \textbf{11} & 283–293 & {\policeGr Σ} 876 (le stratège athénien Protéas mort depuis dix ans), {\policeGr Σ} 897 (Idothée / Théonoé), {\policeGr Σ} 898 (dème de Gargettos et tombeau d'Eurysthée), {\policeGr Σ} 910 ({\policeGr ἐκ τῶν ἰφύων} vs {\policeGr λαιφίων}), {\policeGr Σ} 917 (le signe critique χ et Callistrate) \\
\textbf{23} & \textbf{f. 173v} & 175v & V° (g.) & 920–958 & 40 & \textbf{11} & 294–304 & {\policeGr Σ} 922 ({\policeGr ᾐγυπτιάζετ'}), {\policeGr Σ} 928 (proverbe {\policeGr αὕτη μὲν ἡ μήρινθος}), {\policeGr Σ} 935 ({\policeGr ἱστιορράφος}), {\policeGr Σ} 940 ({\policeGr πρὸς τῇ σανίδι}, fr. de Cratinos), {\policeGr Σ} 941 (\textit{Autolycos II} d'Eupolis), {\policeGr Σ} 947 ({\policeGr παίσωμεν} attique), {\policeGr Σ} 949 (le peintre Pauson) \\
\textbf{24} & \textbf{f. 174r} & 176r & R° (dr.) & 959–1000 & 43 & \textbf{6} & 305–310 & {\policeGr Σ} 968 (passage de la ronde à l'ode), {\policeGr Σ} 969 (Apollon {\policeGr Εὐλύρας}), {\policeGr Σ} 973 (culte nuptial de Zeus Téléios et d'Héra Téléia, {\policeGr προτέλεια}), {\policeGr Σ} 977 (Hermès Nomios), {\policeGr Σ} 986 ({\policeGr τόρνευε τορῶς}) \\
\textbf{25} & \textbf{f. 174v} & 176v & V° (g.) & 1001–1043 & 40 & \textbf{18} & 311–328 & {\policeGr Σ} 1001–1007 (phonétique barbare de l'Archer scythe), {\policeGr Σ} 1012 (\textit{Andromède} jouée avec l'\textit{Hélène} en 412), {\policeGr Σ} 1015, 1018 et 1022 (fragments originaux de l'\textit{Andromède} d'Euripide), {\policeGr Σ} 1031 ({\policeGr κημός}), {\policeGr Σ} 1033 (Glaukétès), {\policeGr Σ} 1040 (variante {\policeGr φῶτα ἀντομένα}) \\
\textbf{26} & \textbf{f. 175r} & 177r & R° (dr.) & 1044–1081 & 40 & \textbf{15} & 329–343 & {\policeGr Σ} 1050–1051 ({\policeGr διχῶς} : {\policeGr τὸν βάρβαρον} vs {\policeGr τὸν ἄθλιον}), {\policeGr Σ} 1054 (glose déplacée {\policeGr σκοτεινήν} = {\policeGr ψεφαυγοῦς} v. 1041), {\policeGr Σ} 1059 (l'\textit{Adonis} de Ptolémée IV Philopator et Agathocle), {\policeGr Σ} 1060 ({\policeGr πέρυσιν} = 412), {\policeGr Σ} 1065–1072 (prologue de l'\textit{Andromède}) \\
\textbf{27} & \textbf{f. 175v} & 177v & V° (g.) & 1082–1127 & 46 & \textbf{18} & 344–361 & {\policeGr Σ} 1083 ({\policeGr λαλῖς} sans ε), {\policeGr Σ} 1085 ({\policeGr σί κακόν} = {\policeGr σοὶ κακόν}), {\policeGr Σ} 1086 ({\policeGr πόθεν ἡ φωνή}), {\policeGr Σ} 1090 (attribution au Chœur), {\policeGr Σ} 1094 (restitution de {\policeGr <οὐ καιρήσεις>}), {\policeGr Σ} 1098 et 1105 (\textit{Andromède}), {\policeGr Σ} 1103 (Gorgos le greffier, vv. 1102–1103 en marge inférieure) \\
\textbf{28} & \textbf{f. 176r} & 178r & R° (dr.) & 1128–1168 & 36 & \textbf{3} & 362–364 & {\policeGr Σ} 1130 (parodie de \textit{Médée} 298), {\policeGr Σ} 1151 (glose intra-textuelle {\policeGr ἵνα} = {\policeGr ὅπου}), {\policeGr Σ} 1168 (clause du traité de paix d'Euripide avec les femmes) \\
\textbf{29} & \textbf{f. 176v} & 178v & V° (g.) & 1169–1210 & 43 & \textbf{15} & 365–379 & {\policeGr Σ} 1172 (Euripide en vieille maquerelle et la danseuse Élaphion), {\policeGr Σ} 1174 ({\policeGr κἀνακάλπασον}), {\policeGr Σ} 1175 (\textit{Histoire théâtrale} de Juba II et \textit{Prôtatikos} de Séleucos sur l'{\policeGr ὄκλασμα}), {\policeGr Σ} 1187 (\textit{parepigraphê}), {\policeGr Σ} 1190 ({\policeGr οὐκὶ πιλῆσι}), {\policeGr Σ} 1197 ({\policeGr συβήνη}) \\
\textbf{30} & \textbf{f. 177r} & 179r & R° (dr.) & 1211–1231 & 24 & \textbf{7} & 380–386 & {\policeGr Σ} 1214 ({\policeGr διέβαλλέ με} = {\policeGr ἐξηπάτησε}, \textit{Oiseaux} 1648), {\policeGr Σ} 1218 (ruse du Chœur et {\policeGr πηκτίς}), {\policeGr Σ} 1224 (jeu de scène final), {\policeGr Σ} 1227 (glose {\policeGr πέπαικται} restituant {\policeGr πέπαισται} contre {\policeGr πέπυσται}), {\policeGr Σ} 1230 (Déméter et Perséphone) \\
\bottomrule
\end{tabular}\end{adjustbox}
}

\newpage
\noindent{\Large\bfseries\color{rubricRed}2. Paléographie et système tachygraphique de la Main~B (\Rdeux{}, \textit{scholiorum manus})}\par\vspace{4pt}

\noindent Tandis que le scribe principal \Run{} (Main~A) copie la colonne poétique centrale en minuscule bouletée calligraphique en écrivant tous les mots en toutes lettres (avec l'iota muet systématiquement adscrit), le réviseur et scholiaste contemporain \Rdeux{} (Main~B, \textit{scholiorum manus}) utilise une \textbf{semi-onciale de glose} (petites capitales détachées mêlées d'éléments cursifs) fortement condensée par un \textbf{système d'abréviations tachygraphiques} qui lui permet de loger dans les marges l'héritage des commentaires alexandrins :

\par\vspace{4pt}
{\fontsize{8.8pt}{10.8pt}\selectfont
\noindent\begin{adjustbox}{max width=\textwidth}\begin{tabular}{@{}p{3.2cm} p{3.6cm} p{3.4cm} p{9.6cm}@{}}
\toprule
\textbf{Tracé tachygraphique} & \textbf{Position graphique} & \textbf{Développement} & \textbf{Emploi et exemples caractéristiques dans {\policeGr Σ}\expo{R} (\texttt{ff. 162v–177r})} \\
\midrule
\textbf{Trait oblique \texttt{\textbackslash}} & Suscrit sur la consonne finale & \textbf{{\policeGr -ων}} (gén. plur. / part.) & Constant : {\policeGr τ\expo{ῶν}, γυναικ\expo{ῶν}, θεσμοφοριαζουσ\expo{ῶν}, οἰκοδομούντ\expo{ων}, Ἀθηναί\expo{ων}}. \\
\textbf{Godet / circonflexe} & En exposant après la consonne & \textbf{{\policeGr -ου}} (gén. sing.) & Constant : {\policeGr τ\expo{οῦ}, αὐτ\expo{οῦ}, Εὐριπίδ\expo{ου}, Ἀγάθων\expo{ος}, Σοφοκλέ\expo{ους}, Παλαμήδ\expo{ους}}. \\
\textbf{Anneau / spirale \texttt{@}} & En exposant en fin de mot & \textbf{{\policeGr -ος}} (nom./gén.) & Constant : {\policeGr Μνησίλοχ\expo{ος}, χορ\expo{ὸς}, χειμῶν\expo{ος}, τραγῳδοποι\expo{ός}, Ἀριστοφάν\expo{ους}}. \\
\textbf{Sigma lunaire haut} & En exposant après la consonne & \textbf{{\policeGr -ης}} ou \textbf{{\policeGr -ας}} & Constant : {\policeGr τ\expo{ῆς}, τ\expo{ὰς}, γυναικ\expo{ὸς} / γυναῖκ\expo{ας}, τραγῳδί\expo{ας}, Ἀνδρομέδ\expo{ας}, Ἑλέν\expo{ης}}. \\
\textbf{Double oblique \texttt{//}} & En exposant en fin de mot & \textbf{{\policeGr -αις}} ou \textbf{{\policeGr -οις}} (dat. pl.) & Constant : {\policeGr τ\expo{αῖς}, τ\expo{οῖς}, Ἀττικ\expo{οῖς}, τραγικ\expo{οῖς}, κωμικ\expo{οῖς}, Θεσμοφορί\expo{οις}}. \\
\textbf{Voyelle/nasale haute} & Au-dessus de la consonne & \textbf{{\policeGr -ον, -ην, -αν, -εν, -ιν}} & Suspensions courantes : {\policeGr τ(ὸν), τ(ὴν), εἶπ(εν), φησ(ίν), ἐστ(ίν), πάλ(ιν)}. \\
{\policeGr ἀν\auDessus{τ}{ι} τ\expo{οῦ}} & En tête de glose explicative & \textbf{{\policeGr ἀντὶ τοῦ}} (« au sens de ») & 38 occurrences dans nos 30 folios (\texttt{{\policeGr Σ} 007, 008, 013, 032, 040, 055, 062, 065, 066, 086...}). \\
{\policeGr λείπ\expo{ει}} & En tête de glose syntaxique & \textbf{{\policeGr λείπει}} (« il manque ») & 9 occurrences (\texttt{{\policeGr Σ} 009, 010, 014, 019, 025, 170, 332, 341}) signalant une ellipse. \\
{\policeGr γρ/} et {\policeGr διχ\expo{ῶς}} & En marge ou dans la scholie & \textbf{{\policeGr γράφεται}} / \textbf{{\policeGr διχῶς}} & Variantes antiques : {\policeGr γρ. σύ τ' αὐτός} (v.~141), {\policeGr γρ. ἑταίρα} (v.~346), {\policeGr γρ. φῶτα ἀντομένα} (v.~1040), {\policeGr διχῶς} (vv.~1051, 1176). \\
{\policeGr τ(ὸ) σημεῖ\expo{ον}} / {\policeGr τ(ὸ) χ} & En tête de scholie alexandrine & \textbf{{\policeGr τὸ σημεῖον}} / \textbf{{\policeGr τὸ χ}} & Renvoi aux signes critiques d'Aristarque et Callistrate : \texttt{{\policeGr Σ} 038} (v.~74), \texttt{{\policeGr Σ} 293} (v.~917), \texttt{{\policeGr Σ} 367} (v.~1174). \\
\textbf{Signe \textbf{{\policeGr ς}} et {\policeFr ÷}} & Dans le corps / en fin de note & \textbf{{\policeGr καί}} et \textbf{clôture} & Le signe en S plongeant note {\policeGr καί} ; le lemnisque ({\policeFr ÷}) clôt chaque scholie. \\
\bottomrule
\end{tabular}\end{adjustbox}
}

\par\vspace{8pt}
\noindent{\Large\bfseries\color{rubricRed}3. Stratigraphie alexandrine du commentaire {\policeGr Σ}\expo{R} et transmission lexicographique}\par\vspace{4pt}

\noindent L'analyse interne des 386 entrées du \textit{Ravennas} révèle quatre strates historiques superposées qui confèrent à ce corpus une valeur documentaire inestimable :
\begin{enumerate}
  \item \textbf{La strate alexandrine ancienne (III\ieme{}–II\ieme{} s. av. J.-C.) :}
  \begin{itemize}
    \item \textbf{Callimaque de Cyrène} (\texttt{{\policeGr Σ} 039}, v.~80) : cité nommément ({\policeGr κατὰ Καλλιμάχου}) à propos du problème célèbre du « troisième jour / jour du milieu » des Thesmophories et du comput incluant le proagôn d'Halimous (10 Pyanepsion) ;
    \item \textbf{Ératosthène de Cyrène} et \textbf{Théophraste} (\texttt{{\policeGr Σ} 181}, v.~516) : cités côte à côte pour l'explication botanique et entomologique du mot {\policeGr κύτταρος} ;
    \item \textbf{Lycophron de Chalcis} (\texttt{{\policeGr Σ} 197}, v.~567) : cité pour son traité \textit{Sur la Comédie} à propos de {\policeGr ποκάδας} ;
    \item \textbf{Aristophane de Byzance} (\texttt{{\policeGr Σ} 001}, \textit{Hypothesis} ; \texttt{{\policeGr Σ} 071}, v.~162) : auteur de l'\textit{Hypothesis} liminaire (qui conserve les noms de \textbf{{\policeGr Μνησίλοχος}} et de \textbf{{\policeGr Χοιρίνη / Χοιρίλη}}) et de l'émendation décisive {\policeGr Ἀλκαῖος} (contre {\policeGr Ἀχαιός} des plus anciens exemplaires) au v.~162 ;
    \item \textbf{Aristarque de Samothrace} (\texttt{{\policeGr Σ} 015}, v.~31) et son élève \textbf{Callistrate} (\texttt{{\policeGr Σ} 293}, v.~917, commentant le signe critique {\policeGr χ}).
  \end{itemize}
  \item \textbf{La synthèse de Didyme Chalcentère (fin du I\ier{} s. av. J.-C.) et Diodore de Tarse :} Didyme est cité nommément aux vv.~31 (\texttt{{\policeGr Σ} 015}) et 162 (\texttt{{\policeGr Σ} 071}), et sa main se reconnaît dans la grande scholie du v.~21 (\texttt{{\policeGr Σ} 011}, sur le vers de l'\textit{Ajax le Locrien} de Sophocle attribué par erreur à Euripide par Platon et Antisthène, comme l'ont montré Pierre Boudreaux et Louis Hauseux), ainsi que dans les notices chronologiques exactes (vv.~32, 190, 804, 841, 850, 876, 1012, 1059, 1060). Diodore de Tarse est cité au v.~389 (\texttt{{\policeGr Σ} 138}, sur {\policeGr ἐπισμῇ}).
  \item \textbf{L'édition commentée de Symmaque (fin du I\ier{} s. ap. J.-C.), intégrant Juba~II et Séleucos :} Symmaque est cité nommément aux vv.~393 (\texttt{{\policeGr Σ} 139}, où le rédacteur le critique pour rétablir l'étymologie exacte de {\policeGr οἰνοπίπας}) et 710 (\texttt{{\policeGr Σ} 231}). C'est lui qui cite l'\textit{Histoire théâtrale} du roi \textbf{Juba~II de Maurétanie} et le \textit{Prôtatikos} de \textbf{Séleucos d'Alexandrie (« l'Homérique »)} aux vv.~840 (\texttt{{\policeGr Σ} 269}) et 1175 (\texttt{{\policeGr Σ} 368}, sur la danse perse {\policeGr ὄκλασμα}).
  \item \textbf{Les gloses tardives et le recoupement avec Hésychius, Photius, la \textit{Souda} et Zonaras :} Une seule scholie (\texttt{{\policeGr Σ} 173}, v.~491) porte la marque d'une addition protobyzantine par les latinismes {\policeGr σταβλίτων ἢ βουρδωναρίων}. Surtout, comme le démontre constamment \textbf{Louis Hauseux (\textbf{Bibi})}, la confrontation entre {\policeGr Σ}\expo{R}, l'unicodex d'\textbf{Hésychius}, le \textit{Lexique} de \textbf{Photius}, la \textbf{\textit{Souda}} et \textbf{Zonaras} constitue l'instrument décisif de l'\textit{emendatio} du texte des \textit{Thesmophories}.
\end{enumerate}

\vfill\newpage
\noindent{\Large\bfseries\color{rubricRed}4. Tableau synoptique de l'apport de {\policeGr Σ}\expo{R} à l'édition critique de Louis Hauseux (\textbf{Bibi})}\par\vspace{4pt}

{\fontsize{8.3pt}{9.9pt}\selectfont
\noindent\begin{adjustbox}{max width=\textwidth}\begin{tabular}{@{}c l p{3.9cm} p{4.0cm} p{9.1cm}@{}}
\toprule
\textbf{Vers} & \textbf{Folio} & \textbf{Leçon de \Run{}} & \textbf{Témoignage de {\policeGr Σ}\expo{R} / Lexiques} & \textbf{Texte établi et démonstration de Louis Hauseux (\textbf{Bibi})} \\
\midrule
\textbf{ante 1} & \texttt{f. 162v} & \textit{(omis dans \Mon{})} & {\policeGr Μνησίλοχος ... Χοιρίνη / Χοιρίλη} & \textbf{Dossier Mnésiloque} : authenticité historique du nom du beau-père d'Euripide et de sa fille Choirilè/Choirinè (\texttt{Introduction.tex}, ll.~325–357). \\
\textbf{v. 3} & \texttt{f. 162v} & {\policeGr οἷόν τι} & {\policeGr οἷόν τε : δυνατὸν} ({\policeGr Σ}\expo{R}, Souda) & \textbf{{\policeGr οἷόν τε}} : {\policeGr Σ}\expo{R} sauve la particule authentique contre l'iotacisme de \Run{}. \\
\textbf{v. 18} & \texttt{f. 162v} & {\policeGr ἀκοὴν δὲ χοάνης} & {\policeGr λείπει ὡς} ({\policeGr Σ}\expo{R}) & \textbf{{\policeGr ἀκοὴν δὲ χοάνης}} : la glose {\policeGr λείπει ὡς} confirme la leçon de \R{} défendue par Willems et \textbf{Bibi}. \\
\textbf{v. 21} & \texttt{f. 162v} & {\policeGr οἷόν τέ που} & {\policeGr οἷόν γε πού 'στιν} ({\policeGr Σ}\expo{R}) & \textbf{{\policeGr οἷόν γέ που}} : le lemme de Didyme dans {\policeGr Σ}\expo{R} conserve la vraie particule {\policeGr γε}. \\
\textbf{v. 32} & \texttt{f. 162v} & {\policeGr [μν\expo{Η}] οὐχ ἑόρακα} & {\policeGr ἐπειδὴ οὐ πάλαι ἤρξατο...} & \textbf{{\policeGr (Εὐρ.) ... οὐχ ἑόρακας}} : la scholie prouve qu'Euripide continue (d'où le grattage du sigle {\policeGr μν\expo{Η}}). \\
\textbf{v. 39} & \texttt{f. 162v} & {\policeGr λέως} & {\policeGr πᾶς ἔστω λαός} ({\policeGr Σ}\expo{R}, Souda GIVM) & \textbf{{\policeGr λαός}} : {\policeGr Σ}\expo{R} garantit expressément le vocalisme tragique {\policeGr λαός}. \\
\textbf{vv. 71–74} & \texttt{f. 163r} & {\policeGr (Μν.) 71 / (Εὐρ.) 72–74} & {\policeGr στενάζει φοβούμενος... τὸ σημεῖον} & \textbf{{\policeGr (Μνη.) 71–74}} : {\policeGr Σ}\expo{R} 71 et 74 justifient l'attribution de tout le groupe à Mnésiloque. \\
\textbf{v. 80} & \texttt{f. 163v} & {\policeGr ἐπεὶ τρίτη 'στὶ} & Callimaque dans {\policeGr Σ}\expo{R} 80, Hésychius & Corruption pré-alexandrine déjà commentée par Callimaque ; analysée par \textbf{Bibi}. \\
\textbf{v. 125} & \texttt{f. 164r} & {\policeGr δοκίμῳ} & {\policeGr δοκίμῳ} ({\policeGr Σ}\expo{1}, Souda) : {\policeGr δόκιμον} {\policeGr Σ}\expo{2} & \textbf{{\policeGr δοκίμῳ}} : maintenu par \textbf{Bibi} avec \R{}, {\policeGr Σ}\expo{1} et la Souda α 4016. \\
\textbf{v. 162} & \texttt{f. 164r} & {\policeGr κἀλκαῖος ... ἐχύμησαν} & {\policeGr Ἀλκαῖος} Ar. Byz. ({\policeGr Σ}\expo{R}) : {\policeGr Ἀχαιός} ant. & \textbf{{\policeGr κἀλκαῖος ... ἐκεχήνεσαν}} : \textbf{Bibi} garde {\policeGr κἀλκαῖος} avec {\policeGr Σ}\expo{R} et corrige {\policeGr ἐχύμησαν} en {\policeGr ἐκεχήνεσαν}. \\
\textbf{v. 179} & \texttt{f. 164v} & {\policeGr κοινῇ ξυμφορᾷ} & {\policeGr μεγάλῃ καὶ θαυμαστῇ} ({\policeGr Σ}\expo{R}) & \textbf{{\policeGr καινῇ ξυμφορᾷ}} : la glose de {\policeGr Σ}\expo{R} traduit {\policeGr καινῇ} (tandis que la Souda κ 2505 accole les deux !). \\
\textbf{v. 220} & \texttt{f. 165r} & {\policeGr ξυροδίκης} & {\policeGr παρὰ τὸ δουροδόκης} ({\policeGr Σ}\expo{R}, Pollux) & \textbf{{\policeGr ξυροδόκης}} : l'étymologie donnée par {\policeGr Σ}\expo{R} corrige la faute d'iotacisme de \Run{}. \\
\textbf{v. 289} & \texttt{f. 165v} & {\policeGr χοῖρον} & {\policeGr χοιρίον ... ὡς βοίδιον} ({\policeGr Σ}\expo{R}) & \textbf{{\policeGr Χοιρίλ'}} : \textbf{Bibi} restitue le nom propre de la fille de Mnésiloque (\cf{} {\policeGr Σ} \textit{ante} 1). \\
\textbf{v. 291} & \texttt{f. 165v} & {\policeGr καὶ πρὸς θάληκον} & {\policeGr τὸν παιδαρίσκον...} ({\policeGr Σ}\expo{R}) vs Hésych. & \textbf{{\policeGr καὶ πρὸς θυλάκιον}} : \textbf{Bibi} corrige par Hésychius θ 63 ({\policeGr θάλικον· σακίον}) contre la bévue de {\policeGr Σ}\expo{R}. \\
\textbf{vv. 295–300} & \texttt{f. 166r} & {\policeGr [Κῆρυξ] ... [τῇ Δήμητρῃ...]} & {\policeGr γυνὴ μιμουμένη κήρυκα} ({\policeGr Σ}\expo{R} 295, 299) & \textbf{{\policeGr (Κορ.-Ἱέρεια)}} et expulsion des deux gloses intruses des vv.~297 et 300 (absentes de {\policeGr Σ}\expo{R} 299). \\
\textbf{v. 346} & \texttt{f. 166v} & {\policeGr ἑταίρα} (sans ι) & {\policeGr ἐὰν δὲ σὺν τῷ ι, ἑταίρᾳ} ({\policeGr Σ}\expo{R}) & \textbf{{\policeGr ἑταίρᾳ}} : adopté par Fritzsche et \textbf{Bibi} d'après la variante attestée par {\policeGr Σ}\expo{R}. \\
\textbf{v. 393} & \texttt{f. 167r} & {\policeGr οἰνοπότιδας} & {\policeGr οἰνοπίπης} ({\policeGr Σ}\expo{R}, Souda {\policeGr οι} 128) & \textbf{{\policeGr τὰς οἰνοπίπας}} : la scholie de Symmaque restitue le composé authentique contre la banalisation de \Run{}. \\
\textbf{v. 735} & \texttt{f. 171r} & {\policeGr ὦ θερμόταται ... ὦ ποτίσταται} & {\policeGr ὦ μέθυσοι} ({\policeGr Σ}\expo{R} sur {\policeGr ποτίσταται}) & \textbf{Bibi} montre comment la Souda (θ 251) a rattaché à tort la glose {\policeGr μέθυσοι} à {\policeGr θερμόταται}. \\
\textbf{v. 834} & \texttt{f. 172r} & {\policeGr Τηνίοισι} & {\policeGr τὰ μὲν Στήνια...} ({\policeGr Σ}\expo{R}, Souda) & \textbf{{\policeGr Στηνίοισι}} : {\policeGr Σ}\expo{R} conserve le {\policeGr Σ} initial tombé par haplographie dans \Run{}. \\
\textbf{v. 868} & \texttt{f. 172v} & {\policeGr (Μνη.)} \textit{continuat} & {\policeGr οὐκ ἐσπάραξάν σε} ({\policeGr Σ}\expo{R}) & \textbf{{\policeGr (Γυνή) τῶν κοράκων πονηρίᾳ}} : la 2\ieme{} pers. de {\policeGr Σ}\expo{R} prouve le changement d'interlocuteur. \\
\textbf{v. 947} & \texttt{f. 173v} & {\policeGr πέσωμεν} & {\policeGr τὸ παίξωμεν οἱ Ἀττικοὶ παίσωμεν} ({\policeGr Σ}\expo{R}) & \textbf{{\policeGr παίσωμεν}} : la remarque d'atticiste de {\policeGr Σ}\expo{R} restitue la forme authentique contre l'iotacisme de \Run{}. \\
\textbf{vv. 1041/54} & \texttt{f. 175r} & {\policeGr φεύγουσαν ... αἰόλαν} & {\policeGr σκοτεινήν} ({\policeGr Σ}\expo{R} ad 1054, Souda) & \textbf{{\policeGr ψεφαυγοῦς}} (v.~1041) : \textbf{Bibi} montre que la glose {\policeGr σκοτεινήν} égarée au v.~1054 glosait {\policeGr ψεφαυγοῦς}. \\
\textbf{vv. 1083–85} & \texttt{f. 175v} & {\policeGr λαλεῖς} \Rac{} / {\policeGr σί κακόν} & {\policeGr χωρὶς τοῦ ε} / {\policeGr ἀντὶ τοῦ σοὶ κακόν} & \textbf{{\policeGr λαλῖς}} et \textbf{{\policeGr σί κακόν}} : {\policeGr Σ}\expo{R} défend le sabir authentique du Scythe contre les éditeurs modernes. \\
\textbf{vv. 1090–94} & \texttt{f. 175v} & \textit{om.} {\policeGr οὐ καιρήσεις} & {\policeGr τοῦτο ὁ χορός} / {\policeGr ἀντὶ τοῦ οὐ χαιρήσεις} & Attribution au \textbf{Chœur} (v.~1090) et restitution de \textbf{{\policeGr <οὐ καιρήσεις>}} (v.~1094) grâce à {\policeGr Σ}\expo{R}. \\
\textbf{v. 1174} & \texttt{f. 176v} & {\policeGr κἀνακόλπασον} & {\policeGr ἁβρῶς} [{\policeGr ἄκρως} Zon.] {\policeGr βαδίζειν} & \textbf{{\policeGr κἀνακάλπασον}} : restitué grâce au croisement de {\policeGr Σ}\expo{R}, Photius α 1514, Souda κ 623 et Zonaras. \\
\textbf{v. 1190} & \texttt{f. 176v} & {\policeGr τί οὐκ ἐπιλήσει} & {\policeGr οὐχὶ φιλήσει} ({\policeGr Σ}\expo{R}) & \textbf{{\policeGr οὐκὶ πιλῆσι πρῶτά με;}} : la glose de {\policeGr Σ}\expo{R} restaure le barbarisme scythe défiguré par \Run{}. \\
\textbf{v. 1227} & \texttt{f. 177r} & {\policeGr πέπυσται} & {\policeGr πέπαικται} ({\policeGr Σ}\expo{R}) & \textbf{{\policeGr πέπαισται}} : la glose de {\policeGr Σ}\expo{R} restitue le parfait attique de {\policeGr παίζω} contre {\policeGr πέπυσται} de \Run{}. \\
\bottomrule
\end{tabular}\end{adjustbox}
}
}
\newpage
\clearpage
\renewcommand{\currHeadLE}{ΑΡΙΣΤΟΦΑΝΟΥΣ ΘΕΣΜΟΦΟΡΙΑΖΟΥΣΑΙ · CODEX RAVENNAS 429 (\R) · FOLIO 01/30}
\renewcommand{\currHeadRE}{f. 162v (Classense 164v) · Hypothesis + vv. 1–42 · Collation diplomatique}
\renewcommand{\currHeadLO}{ΑΡΙΣΤΟΦΑΝΟΥΣ ΘΕΣΜΟΦΟΡΙΑΖΟΥΣΑΙ · CODEX RAVENNAS 429 (\R) · FOLIO 01/30}
\renewcommand{\currHeadRO}{f. 162v (Classense 164v) · Hypothesis + vv. 1–42 · Collation diplomatique}
\noindent\begin{adjustbox}{max totalheight=30.4cm, max width=\textwidth, keepaspectratio}\begin{minipage}{\textwidth}
\bandeauRav{f. 162v}{Classense 164v}{1–42}{Folio 01/30 · Verso (page de gauche) · 21 scholies \& gloses ({\policeGr Σ} 001–021)}
{\fontsize{8.8pt}{10.2pt}\selectfont
\ligneRav{1}{μ\auDessus{ν}{Η}}{ὦ Ζεῦ χελιδὼν \difference{ἆρά ποτε} φανήσεται;}{\textbf{1} ἄρα πότε R : \lecon{ἆρά ποτε} Biset}
\ligneRav{2}{}{ἀπολεῖ μ᾽ ἀλοῶν \difference{ἅνθρωπος} ἐξ ἑωθινοῦ·}{\textbf{2} ἄνθρωπος R : \lecon{ἅνθρωπος} Boissonade}
\ligneRav{3}{}{οἷόν \difference{τε}, πρὶν τὸν σπλῆνα κομιδῇ μ᾽ ἐκβαλεῖν \hfill\textbf{[β \Rpc]}}{\textbf{3} τι R : \lecon{τε} Σ (δυνατόν), Suid. οι 150 \espApp vv. 3–4 inv. \Rac{} : litt. β', α' corr. \Rpc{}}
\ligneRav{4}{}{παρὰ σοῦ πυθέσθαι ποῖ μ᾽ ἄγεις \difference{ωὖριπίδη}; \hfill\textbf{[α \Rpc]}}{\textbf{4} εὐριπίδη R : \lecon{ωὖριπίδη} Dindorf}
\ligneRav{5}{\auDessus{εὐ}{ρ}}{ἀλλ᾽ οὐκ ἀκούειν δεῖ σε πάνθ᾽ \difference{ὅσ᾽} αὐτίκα}{\textbf{5} ὅσ R : \lecon{ὅσ᾽} \Mon}
\ligneRav{}{—}{ὄψει παρεστώς. \auDessus{:}{\tiny μν} πῶς λέγεις; αὖθις φράσον·}{}
\ligneRav{}{—}{οὐ δεῖ μ᾽ ἀκούειν; \auDessus{:}{\tiny εὐ} οὐχ ἅ γ᾽ ἂν μέλλῃς ὁρᾶν·}{}
\ligneRav{}{μ\auDessus{ν}{Η}}{οὐδ᾽ ἆρ᾽ ὁρᾶν δεῖ μ᾽; \auDessus{:}{\tiny εὐ} οὐχ ἅ γ᾽ ἂν ἀκούειν δέῃ·}{}
\ligneRav{}{μ\auDessus{ν}{Η}}{πῶς μοι παραινεῖς; δεξιῶς μέντοι λέγεις·}{}
\ligneRav{10}{}{οὐ φῂς σὺ χρῆναί μ᾽ οὔτ᾽ ἀκούειν οὔθ᾽ ὁρᾶν·}{\textbf{10} \lecon{ὁρᾶν.} R : ὁρᾶν; Gelenius}
\ligneRav{}{\auDessus{εὐ}{ρ}}{χωρὶς γὰρ αὐτοῖν ἑκατέρου 'στὶν ἡ φύσις·}{}
\ligneRav{}{μ\auDessus{ν}{Η}}{τοῦ μήτ᾽ ἀκούειν μήθ᾽ ὁρᾶν; \auDessus{:}{\tiny εὐ} εὖ ἴσθ᾽ ὅτι·}{}
\ligneRav{}{μ\auDessus{ν}{Η}}{πῶς χωρίς; \auDessus{:}{\tiny εὐ} οὕτω ταῦτα διεκρίθη τότε·}{}
\ligneRav{}{}{αἰθὴρ γὰρ ὅτε τὰ πρῶτα διεχωρίζετο}{}
\ligneRav{15}{}{καὶ ζᾦ᾽ ἐν αὑτῷ ξυνετέκνου κινούμενα,}{}
\ligneRav{16}{}{ᾧ μὲν βλέπειν χρὴ \difference{πρῶτ᾽} ἐμηχανήσατο}{\textbf{16} πρῶτα R : \lecon{πρῶτ᾽} Grynaeus}
\ligneRav{}{}{ὀφθαλμὸν ἀντίμιμον ἡλίου τροχῷ,}{}
\ligneRav{18}{}{ἀκοὴν δὲ χοάνης ὦτα διετετρήνατο·}{\textbf{18} \lecon{ἀκοὴν δὲ χοάνης} R, Willems : ἀκοὴς δὲ χοάνην Biset : ἀκοῇ δὲ χοάνην Boissonade}
\ligneRav{}{μ\auDessus{ν}{Η}}{διὰ τὴν χοάνην οὖν μήτ᾽ ἀκούω μήθ᾽ ὁρῶ;}{}
\ligneRav{20}{}{νὴ τὸν Δί᾽ ἥδομαί γε τουτὶ προσμαθών·}{}
\ligneRav{21}{}{οἷόν \difference{γέ} που ᾽στιν αἱ σοφαὶ ξυνουσίαι·}{\textbf{21} τε R : \lecon{γε} Σ}
\ligneRav{}{\auDessus{εὐ}{ρ}}{πόλλ᾽ ἂν μάθοις τοιαῦτα παρ᾽ ἐμοῦ. \auDessus{:}{\tiny μν} πῶς ἂν οὖν}{}
\ligneRav{(23)}{}{\textlbrackdbl\sout{πρὸς τοῖς ἀγαθοῖς τούτοισιν ἐξεύροιμ᾽ ὅπως}\textrbrackdbl}{\textbf{v. 23} πρὸς τοῖς ἀγαθοῖς τούτοισιν ἐξεύροιμ᾽ ὅπως R : \lecon{del.} Willems}
\ligneRav{24}{}{ἔτι \difference{προσμάθοιμι} χωλὸς εἶναι τὼ σκέλει·}{\textbf{24} προσμάθοι· μὴ R : \lecon{προσμάθοιμι} Ellebodius}
\ligneRav{25}{\auDessus{εὐ}{ρ}}{βάδιζε δευρὶ καὶ πρόσεχε τὸν νοῦν. \auDessus{:}{\tiny μν} ἰδού·}{}
\ligneRav{}{\auDessus{εὐ}{ρ}}{ὁρᾷς τὸ θύριον τοῦτο; \auDessus{:}{\tiny μν} νὴ τὸν Ἡρακλέα}{}
\ligneRav{27}{—}{οἶμαί γε \auDessus{:}{\tiny εὐ} σίγα \difference{νυν} \auDessus{:}{\tiny μν} σιωπῶ τὸ θύριον· \hfill\textbf{[α \Rpc]}}{\textbf{27} νῦν R : \lecon{νυν} Boissonade \espApp θύριον R : \lecon{θύριον;} Dindorf}
\ligneRav{28}{\auDessus{εὐ}{ρ}}{ἄκου᾽ \auDessus{:}{\tiny μν} \difference{ἀκούω} καὶ σιωπῶ τὸ θύριον· \hfill\textbf{[β \Rpc]}}{\textbf{28} ἀκούσω R : \lecon{ἀκούω} Brunck \espApp vers. 28 post 30 transp. \Rac : \lecon{corr.} \Rpc}
\ligneRav{29}{\auDessus{εὐ}{ρ}}{ἐνταῦθ᾽ Ἀγάθων ὁ κλεινὸς οἰκῶν τυγχάνει \hfill\textbf{[γ \Rpc]}}{}
\ligneRav{30}{—}{ὁ τραγῳδοποιός \auDessus{:}{\tiny μν} ποῖος οὗτος \difference{Ἁγάθων}· \hfill\textbf{[δ \Rpc]}}{\textbf{30-1} Ἀγάθων (bis) R : \lecon{Ἁγάθων} et \lecon{ἀγαθῶν} Bibi}
\ligneRav{31}{}{ἔστιν τις \difference{ἀγαθῶν}; μῶν ὁ μέλας ὁ καρτερός·}{\textbf{31} ἔστιν τις ἀγαθῶν \lecon{(Μνη.) continuat} R : (Εὐρ.) \Mon \espApp Ἀγάθων R : \lecon{ἀγαθῶν;} Bibi}
\ligneRav{32}{\auDessus{εὐ}{ρ}}{οὔκ, ἀλλ᾽ ἕτερός τις· \textlbrackdbl\auDessus{\sout{μν}}{Η}\textrbrackdbl\ οὐχ \difference{ἑόρακας} πώποτε;}{\textbf{32} (Μνη.) οὐχ ... πώποτε \Rac : \lecon{(Εὐρ.) continuat} \Rpc, Σ \espApp ἑόρακα R : \lecon{ἑόρακας} Bentley \ligneApp{v. 33} omm. \Rac, \Mon : in marg. \Rpc}
\ligneRav{33}{μ\auDessus{ν}{Η}}{\textlbrackdbl{\policeFr in marg. \Rpc{} : }μῶν ὁ δασυπώγων; : οὐχ ἑόρακας πώποτε·\textrbrackdbl}{\textbf{33} omm. \Rac, \Mon : in marg. suppl. \Rpc (\lecon{μῶν ὁ δασυπώγων} \Rpc : οὔκ, ἀλλ᾽ ἕτερός τις Rutherford)}
\ligneRav{34}{μ\auDessus{ν}{Η}}{μὰ τὸν Δί᾽ οὔτοι γ᾽ ὥστε \difference{κἀμέ γ᾽} (\inRas{καὶ}{κἀ}μέ γ᾽) εἰδέναι·}{\textbf{34} καὶ μεγ᾽ \Rac \ (ὶ eras. \Rpc) : καὶ μεγ᾽ \Mon : \lecon{κἀμέ γ᾽} edd. omn.}
\ligneRav{35}{\auDessus{εὐ}{ρ}}{καὶ μὴν \difference{†}βε\inRas{α}{β}ίνηκας\difference{†} σύ γ᾽, ἀλλ᾽ οὐκ οἶσθ᾽ ἴσως·}{\textbf{35} βε\inRas{α}{β}ίνηκας \Rpc{} (ex βεαίνηκας \Rac{} ut vid.) : \lecon{†βεβίνηκας†} Bibi}
\ligneRav{}{}{ἀλλ᾽ ἐκποδὼν πτήξωμεν, ὡς ἐξέρχεται}{}
\ligneRav{}{}{θεράπων τις αὐτοῦ πῦρ ἔχων καὶ μυρρίνας}{}
\ligneRav{}{}{προθυσόμενος ἔοικε τῆς ποιήσεως·}{}
\ligneRav{39}{θερ/}{\hspace{1.1em}εὔφημος πᾶς ἔστω \difference{λαός},}{\textbf{39} λέως R, Suid. ε 3792 codd. AF : \lecon{λαός} Σ, Suid. codd. GIVM}
\ligneRav{40}{}{\hspace{1.1em}στόμα συγκλῄσας· ἐπιδημεῖ γὰρ}{}
\ligneRav{}{}{\hspace{1.1em}θίασος Μουσῶν ἔνδον μελάθρων}{}
\ligneRav{}{}{\hspace{1.1em}τῶν δεσποσύνων μελοποιῶν·}{}
}
\vspace{3pt}{\color{ruleGold}\hrule height 0.6pt}\vspace{3pt}
{\fontsize{8.0pt}{9.4pt}\selectfont
\noindent{\bfseries\scshape\color{rubricRed}Loci paralleli et testimonia (f. 162v, vv. 1–42) :} \textbf{1} \Fr 617 πυθοῦ χελιδὼν πηνίκ᾽ ἄττα φαίνεται \espTest \textbf{2} (ἀλοῶν) Suidas α 1330 ; Photius α 1029 = \SynB α 984 \espTest \textbf{18-9} (χοάνη = χώνη) Moeris χ 1 ~ Hesychius χ 576 χοάνη \espTest \textbf{24} \Macaire VI 90 ; \Apostol II 94 \espTest \textbf{33} (δασυπώγων) Pollux II 88 \espTest \textbf{36} (ἀλλ᾽ -- πτήξωμεν) Suidas π 3024 \espTest \textbf{39-41} Suidas ε 3791 (λεὼς codd. AF : λαὸς codd. GIVM, συγκλείσας, ἐνδημεῖ pro ἐπιδημεῖ ) \espTest \textbf{40-2} (ἐπιδημεῖ -- μελοποιῶν) Suidas θ 378\par\vspace{2.5pt}
{\color{ruleGold}\hrule height 0.35pt}\vspace{2.5pt}
\noindent{\bfseries\scshape\color{rubricRed}Suite (en commentaire) de l'apparat critique et paléographique :} \textbf{Apport du fac-similé phototypique de 1904 (BnF, \textit{FOL-FAC SIM-182 (9)}, éd. J. Van Leeuwen) :} Au sommet du \textit{f. 162v}, les trois lignes de l'\textit{Hypothesis} copiées par \Rdeux{} au ras de la tranche supérieure, aujourd'hui fortement voilées sur le cliché numérique de Ravenne par l'assombrissement du bord du parchemin, sont intégralement nettes sur la planche phototypique de 1904 (restaurée en tête de la planche ci-contre), confirmant sans ambiguïté les lectures \textbf{Χοιρίνη}, \textbf{μή(τ)ηρ δὲ Κλειτώ} et \textbf{προλογίζει Μνησίλοχ(ος) ὁ κηδεστὴς Εὐριπίδου}. --- En haut de la page, \Rdeux{} (Main~B) a copié en semi-onciale de glose une notice liminaire de trois lignes (ὁ χορὸς ἐκ θεσμοφοριαζουσῶν · καὶ τοῦτο τὸ δρᾶμα τῶν κατ' Εὐριπίδου πεποιημένων · Εὐριπίδου γυνὴ μὲν Χοιρίνη, / μή(τ)ηρ δὲ Κλειτώ :{\policeFr ∼} προλογίζει Μνησίλοχ(ος) ὁ κηδεστὴς Εὐριπίδου {\policeFr ÷}, sans liste séparée des \italique{Dramatis personae}), source unique des noms de \textbf{Μνησίλοχος} et de \textbf{Χοιρίνη / Χοιρίλη} défendus par Louis Hauseux (« Le dossier Mnésiloque », \textit{Introduction.tex}, ll. 325–357). --- Au-dessous, le titre en semi-onciale d'apparat précédé d'un \italique{asteriscus} : *~·~* ἀριστοφάνους θεσμοφοριάζουσαι —. --- \textbf{Inversion des vv. 3–4 et transposition des vv. 27–30 corrigées par lettres-chiffres :} Aux vv.~3–4, \Rac{} avait copié le v.~4 avant le v.~3 (ordre rétabli par β', α' de \Rpc{}) ; aux vv.~27–30, \Rac{} avait copié le v. 28 après le v. 30 ; \Rpc{} a rétabli l'ordre exact en inscrivant à l'extrémité droite des quatre lignes les lettres numérales \textbf{α, γ, δ, β}. --- \textbf{Sigle tachygraphique gratté et scholie au v. 32 :} Comme le relève Louis Hauseux (\textit{vv 1-641.tex}, ll. 777–787), le manuscrit comportait à l'intérieur du vers, devant οὐχ ἑόρακα πώποτε, le sigle tachygraphique μ\auDessus{ν}{Η} (\auDessus{μν}{Η} = Μνησίλοχος), qui a été ensuite \textbf{gratté (\italique{ut videtur})} ; l'erreur de ce sigle est en outre prouvée par la scholie {\policeGr Σ}\expo{R} du v. 32 (ἐπειδὴ οὐ πάλαι ἤρξατο διδάσκειν, ἀλλὰ πρὸ τριῶν ἐτῶν πρὸ τούτου), qui ne se comprend qu'avec la 2\ieme{} personne ἑόρακας prononcée par Euripide. --- \textbf{Omission de vers réparée en marge (v. 33) :} Le \textbf{v. 33} (μῶν ὁ δασυπώγων; --- οὐχ ἑόρακας πώποτε) avait été omis par saut du même au même dans la colonne par \Rac{} (et reste \textbf{omis dans \Mon{}}) ; il est rajouté sur deux lignes dans la marge droite par \Rpc{}. --- \textbf{Grattage au v. 34 et correction \latin{in rasura} au v. 35 :} Au v.~34, \Rac{} avait écrit καὶ μεγ' (copié tel quel par \Mon{}) ; \Rpc{} a gratté l'iota (ὶ eras.) pour amorcer la crase κἀμέ γ'. Au \textbf{v.~35}, le second β de \textbf{βεβίνηκας} est corrigé \latin{in rasura} sur un α gratté (\textbf{βεαίνηκας} \Rac{}), ce qui, joint au silence complet de {\policeGr Σ}\expo{R}, motive les \latin{cruces} \lecon{†βεβίνηκας†} de \lecon{Bibi}. --- \textbf{v. 1 :} ἄρα πότε R : \lecon{ἆρά ποτε} Biset \espApp χελιδὼν R : \lecon{Χελιδὼν} \lecon{Bibi} ap. Clements (dans l'\italique{Index nominum} et le commentaire : référence au devin \italique{Khelidôn}, πυθοῦ Χελιδόνος). --- \textbf{v. 3 :} τι R : \lecon{τε} {\policeGr Σ} (δυνατόν), Souda οι 150. --- \textbf{v. 10 :} \lecon{ὁρᾶν.} R (ponctuation haute maintenue par L. Hauseux) : ὁρᾶν; Gelenius. --- \textbf{v. 18 :} \lecon{ἀκοὴν δὲ χοάνης} R, défendu par Alphonse Willems et adopté par L. Hauseux : ἀκοὴς δὲ χοάνην Biset : ἀκοῇ δὲ χοάνην Boissonade. --- \textbf{v. 21 :} τε R : \lecon{γε} {\policeGr Σ}. --- \textbf{v. 23 :} πρὸς τοῖς ἀγαθοῖς τούτοισιν ἐξεύροιμ' ὅπως R : \textbf{vers athétisé (\italique{del.}) par Alphonse Willems} et L. Hauseux comme interpolation issue d'une glose syntaxique. --- \textbf{v. 24 :} προσμάθοι· μὴ R (médivision avec \italique{dikolon} fautif) : \lecon{προσμάθοιμι} Ellebodius. --- \textbf{v. 28 :} ἀκούσω R : \lecon{ἀκούω} Brunck \espApp vers. 28 post 30 transp. \Rac{} : \lecon{corr.} \Rpc{} (litt. α, γ, δ, β in mg.). --- \textbf{vv. 30–31 :} Ἀγάθων (bis) R : \lecon{Ἁγάθων} (v. 30, avec crase ὁ Ἀγάθων) et \lecon{ἀγαθῶν} (v. 31, génitif pluriel en calembour) \lecon{Bibi} \espApp Au \textbf{v. 31}, ἔστιν τις ἀγαθῶν \lecon{(Μνη.) continuat} R (maintenu par \textbf{Bibi}) : (Εὐρ.) \Mon{}. --- \textbf{v. 32 :} (Μνη.) (μ\auDessus{ν}{Η}) οὐχ ... πώποτε \Rac{} (sigle gratté \italique{ut vid.}) : \lecon{(Εὐρ.) continuat} \Rpc{}, {\policeGr Σ} \espApp ἑόρακα R : \lecon{ἑόρακας} Bentley. --- \textbf{v. 35 :} βε\inRas{α}{β}ίνηκας \Rpc{} (ex βεαίνηκας \Rac{} ut vid.) : \lecon{†βεβίνηκας†} \lecon{Bibi}. --- \textbf{v. 39 :} λέως R, Souda ε 3792 codd. AF : \lecon{λαός} {\policeGr Σ}, Souda codd. GIVM.\par}
\end{minipage}\end{adjustbox}
\newpage
\renewcommand{\currHeadLE}{SCHOLIES MARGINALES ({\policeGr Σ}\expo{R}) \& GLOSES INTERLINÉAIRES · FOLIO 01/30}
\renewcommand{\currHeadRE}{f. 162v (Classense 164v) · Hypothesis + vv. 1–42 · {\policeGr Σ} 001 à {\policeGr Σ} 021}
\renewcommand{\currHeadLO}{SCHOLIES MARGINALES ({\policeGr Σ}\expo{R}) \& GLOSES INTERLINÉAIRES · FOLIO 01/30}
\renewcommand{\currHeadRO}{f. 162v (Classense 164v) · Hypothesis + vv. 1–42 · {\policeGr Σ} 001 à {\policeGr Σ} 021}
\bandeauScholies{Folio 01/30 — f. 162v}{Classense 164v}{Hypothesis + vv. 1–42}{Verso (page de gauche) · 21 entrées ({\policeGr Σ} 001 à {\policeGr Σ} 021)}
{\fontsize{8.8pt}{10.5pt}\selectfont\noindent\textbf{\color{rubricRed}Disposition topographique des marges sur le f. 162v :} Marge extérieure à gauche ; marge intérieure à droite. — \textit{Fac-similé de référence :} \texttt{folios\_individuels\_ff162v-177r/01\_f162v\_classense164v\_vv0001-0042.jpg}.\par\vspace{3pt}}
\entreeScholie{001}{ante 1}{Ὑπόθεσις \& Dramatis personae (f. 162v, bandeau supérieur)}{Marge supérieure (3 lignes en semi-onciale de R\expo{2} au ras de la tranche supérieure, restaurées d'après le fac-similé phototypique BnF de 1904)}{Προλογίζει Μνησίλοχος κηδεστὴς Εὐριπίδου. ὁ χορὸς ἐκ θεσμοφοριαζουσῶν. καὶ τοῦτο τὸ δρᾶμα τῶν κατ’ Εὐριπίδου πεποιημένων. ἀπὸ δὲ τῶν θεσμοφορίων τὰς θεσμοφοριαζούσας γυναῖκας ἐπέγραψεν, ἀφ’ ὧν καὶ ὁ χορός. Εὐριπίδου γυνὴ μὲν Χοιρίλη, μήτηρ δὲ Κλειτώ.}{ὁ χορ\expo{ὸς} ἐκ θεσμοφοριαζουσ\expo{ῶν} · \textbf{ς} τοῦτ(ο) τ(ὸ) δρᾶμ(α) τ\expo{ῶν} κατ' Εὐριπίδ\expo{ου} πεποιημέν\expo{ων} · ἀπ(ὸ) δ(ὲ) τ\expo{ῶν} θεσμοφορί\expo{ων} τ\expo{ὰς} θεσμοφοριαζούσ\expo{ας} γυναῖκ\expo{ας} ἐπέγραψ(εν), ἀφ' ὧν \textbf{ς} ὁ χορός · Εὐριπίδ\expo{ου} γυνὴ μ(ὲν) Χοιρίνη, μή(τ)ηρ δ(ὲ) Κλειτώ :{\policeFr ∼} προλογίζ(ει) Μνησίλοχ\expo{ος} ὁ κηδεστ\expo{ὴς} Εὐριπίδ\expo{ου} {\policeFr ÷}}{« Le prologue est dit par Mnésiloque, parent par alliance (beau-père) d'Euripide. Le Chœur est formé de femmes célébrant les Thesmophories. Cette pièce fait elle aussi partie de celles qui ont été composées contre Euripide ; c'est d'après la fête des Thesmophories qu'il l'a intitulée Les Thesmophoriazuses, d'où provient aussi le Chœur. La femme d'Euripide s'appelait Choirilè (Choirinè sur R), et sa mère Clitô. »}
\commScholie{Notice liminaire capitale d'Aristophane de Byzance, copiée par R\expo{2} en tête du f. 162v (omise par M). C'est l'unique témoin antique conservant à la fois le nom de Μνησίλοχος et celui de Χοιρίνη / Χοιρίλη, dont Louis Hauseux (« Le dossier Mnésiloque », Introduction.tex, ll. 325-357, et v. 289 Χοιρίλ') démontre l'authenticité historique grâce aux règles de l'onomastique familiale attique (Démosthène, C. Macartatos, 43.74).}
\entreeScholie{002}{1}{ὦ Ζεῦ Χελιδὼν ἆρά ποτε φανήσεται;}{Marge gauche (extérieure, haut de colonne, en face des vv. 1–6)}{ἐπεὶ εἰώθασιν ἀπὸ χειμῶνος εὔχεσθαι ἔαρ, τῷ δὲ ἔαρι χελιδόνες φαίνονται, οὗτος δὲ ὥσπερ ἐχειμάσθη περιαγόμενος ὑπὸ Εὐριπίδου ἀλύοντος. τοῦτο ἔφη ἐν ἤθει οἷον πότε ἀπαλλαγήσομαι τοῦ κακοῦ τούτου, ὥσπερ οἱ ἐκ χειμῶνος ἐπιθυμοῦντες ἔαρ ἀφικέσθαι.}{ὦ ζεῦ χελιδών · ἐπεὶ εἰώθασ(ιν) ἀπ(ὸ) χειμῶν\expo{ος} εὔχεσθ(αι) ἔαρ, τῷ δ(ὲ) ἔαρι χελιδόνες φαίνοντ(αι), οὗτ\expo{ος} δ(ὲ) ὥσπ(ερ) ἐχειμάσθη περιαγόμεν\expo{ος} ὑπ(ὸ) Εὐριπίδου ἀλύοντ\expo{ος} · τοῦτ(ο) ἔφη ἐν ἤθ(ει) οἷ(ον) πότε ἀπαλλαγήσομ(αι) τ\expo{οῦ} κακ\expo{οῦ} τούτου, ὥσπ(ερ) οἱ ἐκ χειμῶν\expo{ος} ἐπιθυμοῦντες ἔαρ ἀφικέσθ(αι) {\policeFr ÷}}{« Puisqu'on a coutume, au sortir de l'hiver, de souhaiter la venue du printemps, et qu'au printemps paraissent les hirondelles, tandis que lui-même a été comme battu par la tempête à force d'être traîné partout par Euripide désemparé. Il dit cela avec caractère (ἐν ἤθει), au sens de : “Quand donc serai-je délivré de ce malheur ?”, tout comme ceux qui, après l'hiver, aspirent au retour du printemps. »}
\commScholie{Grande scholie exégétique d'ouverture. Le scholiaste interprète χελιδών au sens proverbial de l'hirondelle printanière ; Louis Hauseux (Bibi, v. 1 et Index nominum), tout en conservant ce premier niveau, relève avec Clements le double sens comique visant le devin Χελιδών (« interroge Chélidôn pour savoir quand paraît l'hirondelle », Ar. fr. 617 K.-A.).}
\entreeScholie{003}{2}{ἀπολεῖ μ᾽ ἀλοῶν ἅνθρωπος ἐξ ἑωθινοῦ·}{Marge gauche (extérieure, sous Σ 1)}{ἀλοῶν: Ἔξωθεν ἐν κύκλῳ περιάγων ὡς οἱ ἐν ταῖς ἅλωσι. καὶ τὸ τύπτειν ἀλοᾶν λέγουσιν ἀπὸ τῶν κοπτόντων τοὺς στάχυας.}{ἀλοῶν : ἔξωθ(εν) ἐν κύκλ(ῳ) περιάγ\expo{ων} ὡς οἱ ἐν τ\expo{αῖς} ἅλωσι · \textbf{ς} τ(ὸ) τύπτ(ειν) ἀλοᾶν λέγουσ(ιν) ἀπ(ὸ) τ\expo{ῶν} κοπτόντ\expo{ων} τ\expo{οὺς} στάχυ\expo{ας} {\policeFr ÷}}{« ἀλοῶν (“me foulant sur l'aire / me rouant de coups”) : me faisant tourner en rond par l'extérieur comme font les bouviers sur les aires à battre (ἅλωσι). Et l'on appelle aussi ἀλοᾶν le fait de frapper, par métaphore de ceux qui battent les épis. »}
\commScholie{Repris mot pour mot par la Souda α 1330, Photius α 1029 et la Συναγωγή B α 984. Le scholiaste distingue finement les deux sens d'ἀλοᾶν : 1° faire tourner en rond comme les bœufs sur l'aire ; 2° rouer de coups, briser de fatigue.}
\entreeScholie{004}{3}{οἷόν τε, πρὶν τὸν σπλῆνα κομιδῇ μ᾽ ἐκβαλεῖν}{Interlinéaire (supra lineam au-dessus de οἷόν τι de R\expo{1}) et marge intérieure droite}{οἷόν τε: δυνατὸν, ἐρωτηματικῶς.}{οἷόν τ(ε) : δυνατ(όν), ἐρωτηματικ\expo{ῶς} {\policeFr ÷}}{« οἷόν τε : est-il possible (δυνατόν), pris interrogativement. »}
\commScholie{Exemple majeur où {\policeGr Σ}\expo{R} conserve la bonne leçon (οἷόν τε, confirmée par la glose δυνατὸν et par la Souda οι 150) contre la faute d'iotacisme/mélecture οἷόν τι de R\expo{1}. Louis Hauseux (Bibi) édite οἷόν τε d'après {\policeGr Σ}\expo{R}.}
\entreeScholie{005}{5}{ἀλλ᾽ οὐκ ἀκούειν δεῖ σε πάνθ᾽ ὅσ᾽ αὐτίκα / ὄψει παρεστώς.}{Marge gauche (extérieure, en face du v. 5)}{ὁ μὲν τραγικώτερον καὶ ὑψηλότερον φράζει, ὁ δὲ ταπεινότερον ἢ δεῖ ἀκούει.}{ὁ μ(ὲν) τραγικώτερ\expo{ον} \textbf{ς} ὑψηλότερ\expo{ον} φράζ(ει), ὁ δ(ὲ) ταπεινότερ\expo{ον} ἢ δεῖ ἀκούει {\policeFr ÷}}{« L'un (Euripide) s'exprime sur un ton trop tragique et trop élevé, tandis que l'autre (Mnésiloque) comprend ses paroles de façon plus terre à terre qu'il ne faudrait. »}
\commScholie{Cité et approuvé par Louis Hauseux (vv 1-641.tex, l. 521) : le scholiaste définit exactement le ressort comique du dialogue philosophique initial entre le jargon sophistique d'Euripide et le bon sens trivial de Mnésiloque.}
\entreeScholie{006}{11}{χωρὶς γὰρ αὐτοῖν ἑκατέρου 'στὶν ἡ φύσις·}{Marge gauche (extérieure) et glose interlinéaire sur αὐτοῖν}{ταῦτα τῶν φυσικῶν λόγων. χωρὶς, τοῦ ἀκούειν καὶ ὁρᾶν.}{ταῦτ(α) τ\expo{ῶν} φυσικ\expo{ῶν} λόγ\expo{ων} {\policeFr ÷} // χωρὶς, τ\expo{οῦ} ἀκού(ειν) \textbf{ς} ὁρᾶν {\policeFr ÷}}{« Ceci relève des doctrines des physiciens (philosophes de la nature). — Séparément : à savoir le fait d'entendre et le fait de voir. »}
\commScholie{Allusion aux cosmogonies présocratiques (Anaxagore, Diogène d'Apollonie) parodiées par Euripide (notamment dans Mélanippe la Sage).}
\entreeScholie{007}{12}{τοῦ μήτ᾽ ἀκούειν μήθ᾽ ὁρᾶν;}{Marge gauche (extérieure, en face du v. 12)}{ἀντὶ τοῦ εἰπεῖν τοῦ ἀκούειν καὶ ὁρᾶν εἶπε τοῦ μήτ’ ἀκούειν μήθ’ ὁρᾶν.}{ἀν\auDessus{τ}{ι} τ\expo{οῦ} εἰπ(εῖν) τ\expo{οῦ} ἀκού(ειν) \textbf{ς} ὁρᾶν εἶπ(ε) τ\expo{οῦ} μήτ' ἀκού(ειν) μήθ' ὁρᾶν {\policeFr ÷}}{« Au lieu de dire “du fait d'entendre et de voir”, il a dit par plaisanterie “du fait de ne ni entendre ni voir”. »}
\commScholie{Glose explicative soulignant le paralogisme comique de Mnésiloque, qui transforme la distinction épistémologique d'Euripide en double négation.}
\entreeScholie{008}{15}{καὶ ζῷ᾽ ἐν αὑτῷ ξυνετέκνου κινούμενα,}{Marge gauche (extérieure) et gloses interlinéaires sur κινούμενα et ξυνετέκνου}{ἐπεὶ ὁ αἰθὴρ δοκεῖ τρέφειν ἡμᾶς. κινούμενα: οὐδὲν γὰρ δίχα κινήσεως γίνεται. ξυνετέκνου δὲ ἀντὶ τοῦ ἐτεκνοποίει.}{ἐπεὶ ὁ αἰθ(ὴρ) δοκ(εῖ) τρέφ(ειν) ἡμᾶς {\policeFr ÷} // κινούμ(ενα) : οὐδ(ὲν) γ(ὰρ) δίχα κινήσ(εως) γίνετ(αι) {\policeFr ÷} // ξυνετέκν\expo{ου} δ(ὲ) ἀν\auDessus{τ}{ι} τ\expo{οῦ} ἐτεκνοποίει {\policeFr ÷}}{« Parce que l'Éther passe pour nous nourrir. — κινούμενα (“doués de mouvement”) : car rien ne vient à l'être sans mouvement. — Et ξυνετέκνου (“il enfantait en lui”) est employé au sens de ἐτεκνοποίει (“il procréait”). »}
\commScholie{Trois notes fusionnées par Dübner : une scholie cosmologique sur l'Éther nourricier (doctrine d'Anaxagore et d'Euripide, fr. 839 et 941 K.) et deux gloses interlinéaires.}
\entreeScholie{009}{17}{ὀφθαλμὸν ἀντίμιμον ἡλίου τροχῷ,}{Interlinéaire (supra lineam) et marge intérieure droite}{ἀντίμιμον: ὅμοιον. λείπει τὸ περιφερές.}{ἀντίμιμ\expo{ον} : ὅμοι\expo{ον} · λείπ(ει) τ(ὸ) περιφερές {\policeFr ÷}}{« ἀντίμιμον : semblable, imitant en contrepartie. Il manque (par ellipse) “quant à sa forme circulaire” (τὸ περιφερές). »}
\commScholie{Glose reprise par la Souda α 2678. Le néologisme euripidéen ἀντίμιμον est glosé par l'adjectif usuel ὅμοιον.}
\entreeScholie{010}{18}{ἀκοὴν δὲ χοάνης ὦτα διετετρήνατο·}{Interlinéaire (supra lineam au-dessus de χοάνης)}{λείπει ὡς.}{λείπ(ει) ὡς {\policeFr ÷}}{« Il manque (par ellipse) la conjonction ὡς (“comme un entonnoir”). »}
\commScholie{Comme le souligne Alphonse Willems (approuvé par Louis Hauseux, vv 1-641.tex, l. 636, et Annexes.tex, l. 583), cette courte glose syntaxique (λείπει ὡς) prouve que le scholiaste lisait bien le texte authentique de R : ἀκοὴν δὲ χοάνης ὦτα διετετρήνατο (maintenu par Bibi contre les conjectures de Biset et Boissonade).}
\entreeScholie{011}{21}{οἷόν γέ που ᾽στιν αἱ σοφαὶ ξυνουσίαι·}{Marge gauche (extérieure, grand bloc de 15 lignes en semi-onciale de R\expo{2})}{διὰ τούτου ὑπονοεῖ Εὐριπίδου εἶναι ἐκεῖνο τὸ «σοφοὶ τύραννοι τῶν σοφῶν ξυνουσίᾳ.» ἔστι δὲ Σοφοκλέους ἐξ Αἴαντος Λοκροῦ. ἢ οὖν ἐπίτηδες, ἵνα καὶ τοὺς ἄλλους ἐξαπατήσῃ, ἢ συνέμπτωσις Σοφοκλεῖ καὶ Εὐριπίδῃ ἐγένετο. οἷόν γε πού ’στιν αἱ σοφαί: καὶ διὰ τούτου φαίνεται ὑπονοῶν Εὐριπίδου εἶναι τὸ «σοφοὶ τύραννοι τῶν σοφῶν συνουσίᾳ.» ἔστι δὲ Σοφοκλέους ἐξ Αἴαντος Λοκροῦ. ἐνταῦθα μέντοι ὑπονοεῖ μόνον, ἐν δὲ τοῖς Ἥρωσιν ἄντικρυς ἀποφαίνεται. καὶ Ἀντισθένης καὶ Πλάτων [Rep. 8, p. 568, A; Theag. p. 125, D] Εὐριπίδου αὐτὸ εἶναι ἡγοῦνται, οὐκ ἔχω εἰπεῖν ὅ τι παθόντες. ἔοικε δὲ ἤτοι πεπλανημένος ἢ συνεξαπατῆσαι τοὺς ἄλλους, ὥσπερ ὑπονοοῦσί τινες συμπτώσεις τῷ τε Σοφοκλεῖ καὶ τῷ Εὐριπίδῃ, ὥσπερ καὶ ἐπὶ ἄλλων τινῶν. τὸ μέντοι δρᾶμα ἐν ᾧ Εὐριπίδης ταῦτα εἶπεν οὐ σώζεται.}{διὰ τούτου ὑπονο(εῖ) Εὐριπίδου εἶν(αι) ἐκεῖνο τ(ὸ) «σοφοὶ τύραννοι τ\expo{ῶν} σοφ\expo{ῶν} ξυνουσίᾳ · » ἔστι δ(ὲ) Σοφοκλέους ἐξ Αἴαντ\expo{ος} Λοκρ\expo{οῦ} · ἢ οὖν ἐπίτηδες, ἵνα \textbf{ς} τ\expo{οὺς} ἄλλ\expo{ους} ἐξαπατήσῃ, ἢ συνέμπτωσις Σοφοκλ(εῖ) \textbf{ς} Εὐριπίδῃ ἐγένετο · οἷόν γε πού ’στ(ιν) αἱ σοφαί: \textbf{ς} διὰ τούτου φαίνετ(αι) ὑπονο\expo{ῶν} Εὐριπίδου εἶν(αι) τ(ὸ) «σοφοὶ τύραννοι τ\expo{ῶν} σοφ\expo{ῶν} συνουσίᾳ · » ἔστι δ(ὲ) Σοφοκλέους ἐξ Αἴαντ\expo{ος} Λοκρ\expo{οῦ} · ἐνταῦθα μέντοι ὑπονο(εῖ) μόνον, ἐν δ(ὲ) τ\expo{οῖς} Ἥρωσ(ιν) ἄντικρυς ἀποφαίνετ(αι) · \textbf{ς} Ἀντισθένης \textbf{ς} Πλάτων Εὐριπίδου αὐτὸ εἶν(αι) ἡγοῦντ(αι), οὐκ ἔχω εἰπ(εῖν) ὅ τι παθόντες · ἔοικε δ(ὲ) ἤτοι πεπλανημένος ἢ συνεξαπατῆσ(αι) τ\expo{οὺς} ἄλλ\expo{ους}, ὥσπ(ερ) ὑπονοοῦσί τινες συμπτώσεις τῷ τε Σοφοκλ(εῖ) \textbf{ς} τῷ Εὐριπίδῃ, ὥσπ(ερ) \textbf{ς} ἐπὶ ἄλλ\expo{ων} τιν\expo{ῶν} · τ(ὸ) μέντοι δρᾶμα ἐν ᾧ Εὐριπίδης ταῦτ(α) εἶπ(εν) οὐ σώζετ(αι) {\policeFr ÷}}{« Par ce vers, il laisse entendre qu'est d'Euripide le fameux vers : “Les tyrans deviennent sages par la fréquentation des sages.” Or ce vers est de Sophocle, dans son Ajax le Locrien (fr. 14 Radt). Ou bien donc Aristophane le fait exprès afin d'induire aussi les autres en erreur, ou bien il y a eu rencontre fortuite (συνέμπτωσις) entre Sophocle et Euripide. — Autre rédaction : Par là aussi il paraît croire que “Les tyrans sont sages par la fréquentation des sages” est d'Euripide, alors qu'il est de Sophocle dans l'Ajax le Locrien. Ici toutefois il le suggère seulement, tandis que dans les Héros (fr. 323 K.-A.) il l'affirme ouvertement. Antisthène et Platon (République VIII, 568a ; Théagès, 125d) croient eux aussi qu'il est d'Euripide, je ne saurais dire par quelle méprise. Il semble soit s'être trompé, soit avoir entraîné les autres dans son erreur, à moins qu'il n'y ait eu coïncidence entre Sophocle et Euripide, comme on le soupçonne pour d'autres passages ; cependant le drame dans lequel Euripide aurait dit cela n'est pas conservé. »}
\commScholie{L'une des scholies les plus érudites de la pièce, attribuée à Didyme Chalcentère par Pierre Boudreaux et Louis Hauseux (vv 1-641.tex, ll. 661-668) en raison de son ton polémique à la première personne (οὐκ ἔχω εἰπεῖν ὅ τι παθόντες). De plus, le lemme de {\policeGr Σ}\expo{R} conserve la bonne particule οἷόν γέ που (adoptée par Bibi) contre οἷόν τέ de R\expo{1}.}
\entreeScholie{012}{24}{ἔτι προσμάθοιμι χωλὸς εἶναι τὼ σκέλει·}{Marge intérieure droite (en face du v. 24)}{ἵνα μὴ περιπατήσῃ μετ’ αὐτοῦ.}{ἵνα μὴ περιπατήσ(ῃ) μετ' αὐτ\expo{οῦ} {\policeFr ÷}}{« Afin de ne plus avoir à marcher en sa compagnie (avec Euripide). »}
\commScholie{Comme le note Louis Hauseux (vv 1-641.tex, ll. 686-692, et Annexes.tex, l. 646) à la suite d'Alphonse Willems, l'explication du scholiaste est lumineuse : Mnésiloque souhaite apprendre à boiter des deux jambes pour être dispensé de suivre Euripide, tout en raillant ce dernier comme « faiseur de boiteux » (χωλοποιός). Le v. 23 de R n'est qu'une glose syntaxique versifiée interpolée dans l'archétype et athétisée par Willems et Bibi.}
\entreeScholie{013}{25}{βάδιζε δευρὶ καὶ πρόσεχε τὸν νοῦν.}{Interlinéaire (supra lineam au-dessus de πρόσεχε)}{πρόσεχε ἀντὶ τοῦ μνῆσαι.}{πρόσεχ(ε) ἀν\auDessus{τ}{ι} τ\expo{οῦ} μνῆσαι {\policeFr ÷}}{« πρόσεχε (“prête attention”) : au sens de “souviens-toi” (μνῆσαι). »}
\commScholie{Glose interlinéaire élémentaire de R\expo{2} expliquant l'impératif idiomatique πρόσεχε τὸν νοῦν.}
\entreeScholie{014}{27}{σίγα νυν. — σιωπῶ τὸ θύριον·}{Marge gauche (extérieure, en face du v. 27)}{λείπει ἡ διά, οἷον σιωπῶ διὰ τὸ θύριον. ἔπαιξε δέ.}{λείπ(ει) ἡ διά, οἷον σιωπῶ διὰ τ(ὸ) θύρι\expo{ον} · ἔπαιξ(ε) δέ {\policeFr ÷}}{« Il manque la préposition διά (“je garde le silence à cause de la petite porte”) ; mais il a fait un jeu de mots comique (en construisant σιωπῶ transitivement avec τὸ θύριον : “je tais la petite porte”). »}
\commScholie{Le scholiaste saisit fort bien la construction plaisante de Mnésiloque, qui applique littéralement et transitivement l'ordre σίγα à l'objet du vers précédent (τὸ θύριον), d'où la réplique suivante : ἀκούω καὶ σιωπῶ τὸ θύριον (v. 28, rétabli à sa place par les lettres-chiffres α', γ', δ', β' de R\expo{2}).}
\entreeScholie{015}{31}{ἔστιν τις ἀγαθῶν; μῶν ὁ μέλας ὁ καρτερός;}{Marge gauche (extérieure, en face du v. 31)}{οἱ περὶ Ἀρίσταρχον καὶ Δίδυμόν φασιν εἶναι τοιοῦτον Ἀγάθωνα. ἐγὼ δὲ οὐχ ἡγοῦμαι, ἀλλ’ ἐπειδὴ βούλεται κωμῳδεῖν τὸν Ἀγάθωνα ὡς μήτε καρτερὸν μήτε μέλανα, τοῦτο εἶπεν. ἢ ἄσημόν τινα.}{οἱ π(ερὶ) Ἀρίσταρχ\expo{ον} \textbf{ς} Δίδυμ\expo{όν} φασ(ιν) εἶν(αι) τοιοῦτ\expo{ον} Ἀγάθωνα · ἐγὼ δ(ὲ) οὐχ ἡγοῦμ(αι), ἀλλ' ἐπειδὴ βούλετ(αι) κωμῳδ(εῖν) τ(ὸν) Ἀγάθωνα ὡς μήτ(ε) καρτερ(ὸν) μήτ(ε) μέλανα, τοῦτ(ο) εἶπ(εν) · ἢ ἄσημόν τινα {\policeFr ÷}}{« Aristarque et Didyme affirment qu'il existait réellement un autre Agathon présentant ce physique (brun et vigoureux). Pour ma part, je ne le crois pas : c'est parce qu'il veut railler Agathon comme n'étant ni vigoureux ni hâlé qu'il a dit cela ; ou alors il désigne un homonyme obscur. »}
\commScholie{Scholie alexandrine majeure opposant l'interprétation prosopographique d'Aristarque et de Didyme à l'interprétation antiphrastique du commentateur (Symmaque). Louis Hauseux (Bibi, vv. 30-31) résout brillamment la difficulté en lisant au v. 30 ποῖος οὗτος Ἁγάθων; (avec crase ὁ Ἀγάθων) et au v. 31 ἔστιν τις ἀγαθῶν; (« Y a-t-il quelqu'un parmi les gens de bien ? »), calembour sur Ἀγάθων / ἀγαθῶν.}
\entreeScholie{016}{32}{οὔκ, ἀλλ᾽ ἕτερός τις· οὐχ ἑόρακας πώποτε;}{Marge gauche (extérieure, en face du v. 32)}{ἐπειδὴ οὐ πάλαι ἤρξατο διδάσκειν, ἀλλὰ τρισὶν πρὸ τούτων ἔτεσιν.}{ἐπειδὴ οὐ πάλαι ἤρξατ(ο) διδάσκ(ειν), ἀλλὰ τρισὶν πρὸ τούτ\expo{ων} ἔτεσ(ιν) {\policeFr ÷}}{« Parce qu'il n'avait pas commencé à faire représenter des tragédies depuis longtemps, mais seulement trois ans avant cette pièce (aux Lénéennes de 416 av. J.-C.). »}
\commScholie{Témoignage chronologique capital (rapproché par L. Hauseux, Introduction.tex, l. 175, du Banquet de Platon, 173a, qui fixe la première victoire d'Agathon à 416). En outre, comme le démontre L. Hauseux (vv 1-641.tex, ll. 777-787), cette scholie prouve qu'οὐχ ἑόρακας πώποτε est dit par Euripide (d'où le grattage sur R du sigle fautif μν[Η] devant οὐχ).}
\entreeScholie{017}{33}{μῶν ὁ δασυπώγων; — οὐχ ἑόρακας πώποτε;}{Marge intérieure droite (sous le v. 33 rétabli en marge par R\expo{2})}{διαβάλλει τὸν Ἀγάθωνα πάλιν ὡς λεῖον καὶ ἀποτίλλοντα αὑτοῦ τὰς τρίχας τοῦ πώγωνος.}{διαβάλλ(ει) τ(ὸν) Ἀγάθωνα πάλ(ιν) ὡς λεῖ\expo{ον} \textbf{ς} ἀποτίλλοντ(α) αὑτ\expo{οῦ} τ\expo{ὰς} τρίχ\expo{ας} τ\expo{οῦ} πώγων\expo{ος} {\policeFr ÷}}{« Il raille de nouveau Agathon comme étant glabre et s'épilant lui-même les poils de la barbe. »}
\commScholie{Le v. 33, omis dans la colonne par R\expo{1} (et totalement absent de M) par saut du même au même sur οὐχ ἑόρακας πώποτε, a été rétabli sur deux lignes dans la marge droite par R\expo{2}, accompagné de cette scholie qui atteste l'ancienneté du vers.}
\entreeScholie{018}{36}{ἀλλ᾽ ἐκποδὼν πτήξωμεν, ὡς ἐξέρχεται}{Interlinéaire (supra lineam sur πτήξωμεν) et marge gauche}{κρυβῶμεν. κυρίως δὲ πτήσσειν ἐστὶ τὸ φεύγειν καὶ φοβεῖσθαι.}{κρυβῶμ(εν) · κυρί\expo{ως} δ(ὲ) πτήσσ(ειν) ἐστὶ τ(ὸ) φεύγ(ειν) \textbf{ς} φοβεῖσθ(αι) {\policeFr ÷}}{« κρυβῶμεν (“cachons-nous”). Au sens propre, πτήσσειν signifie fuir et se blottir de peur. »}
\commScholie{Repris intégralement par la Souda π 3024 (s.v. Πτήξωμεν).}
\entreeScholie{019}{38}{προθυσόμενος ἔοικε τῆς ποιήσεως·}{Marge gauche (extérieure) et glose interlinéaire sur ἔοικε}{ὅτε γὰρ δρᾶμα ποιῆσαι ἤθελον, πρότερον θυσίας ἐποίουν. λείπει τὸ ὡς. ὡς ἔοικε.}{ὅτε γ(ὰρ) δρᾶμα ποιῆσ(αι) ἤθελ\expo{ον}, πρότερ\expo{ον} θυσί\expo{ας} ἐποί\expo{ουν} {\policeFr ÷} // λείπ(ει) τ(ὸ) ὡς · ὡς ἔοικ(ε) {\policeFr ÷}}{« Car lorsqu'ils voulaient composer un drame, ils offraient au préalable un sacrifice. — Il manque ὡς : au sens de ὡς ἔοικε (“à ce qu'il semble”). »}
\commScholie{Explique le génitif du composé προ-θύομαι et la construction paratactique d'ἔοικε en incise.}
\entreeScholie{020}{39}{εὔφημος πᾶς ἔστω λαός,}{Marge gauche (extérieure, bas du f. 162v)}{παρὰ τοῖς τραγικοῖς οὕτω λέγεται, πᾶς ἔστω λαός.}{παρὰ τ\expo{οῖς} τραγικ\expo{οῖς} οὕτω λέγετ(αι), πᾶς ἔστω λαός {\policeFr ÷}}{« C'est ainsi que l'on s'exprime chez les poètes tragiques : πᾶς ἔστω λαός (avec la forme non attique λαός au lieu de λεώς). »}
\commScholie{Crucial pour l'établissement du texte de Bibi : alors que R\expo{1} porte la forme attique banalisée λέως (suivie par les mss AF de la Souda ε 3791), {\policeGr Σ}\expo{R} atteste expressément que le Serviteur emploie le vocalisme tragique λαός (confirmé par les mss GIVM de la Souda et adopté par Louis Hauseux).}
\entreeScholie{021}{41}{θίασος Μουσῶν ἔνδον μελάθρων}{Interlinéaire (supra lineam au-dessus de θίασος, bas du f. 162v)}{ὁ ἱερὸς χορός.}{ὁ ἱερ\expo{ὸς} χορός {\policeFr ÷}}{« La troupe sacrée (le chœur sacré des Muses). »}
\commScholie{Glose interlinéaire reprise par la Souda θ 378 (s.v. Θίασος).}
\clearpage
\renewcommand{\currHeadLE}{ΑΡΙΣΤΟΦΑΝΟΥΣ ΘΕΣΜΟΦΟΡΙΑΖΟΥΣΑΙ · CODEX RAVENNAS 429 (\R) · FOLIO 02/30}
\renewcommand{\currHeadRE}{f. 163r (Classense 165r) · vv. 43–79 · Collation diplomatique}
\renewcommand{\currHeadLO}{ΑΡΙΣΤΟΦΑΝΟΥΣ ΘΕΣΜΟΦΟΡΙΑΖΟΥΣΑΙ · CODEX RAVENNAS 429 (\R) · FOLIO 02/30}
\renewcommand{\currHeadRO}{f. 163r (Classense 165r) · vv. 43–79 · Collation diplomatique}
\noindent\begin{adjustbox}{max totalheight=30.4cm, max width=\textwidth, keepaspectratio}\begin{minipage}{\textwidth}
\bandeauRav{f. 163r}{Classense 165r}{43–79}{Folio 02/30 · Recto (page de droite) · 17 scholies \& gloses ({\policeGr Σ} 022–038)}
{\fontsize{9.3pt}{10.8pt}\selectfont
\ligneRav{43}{θερ/}{\hspace{1.1em}\auDessus{ἐχέτω}{\tiny ἐπεχέτω} δὲ πνοὰς νήνεμος αἰθήρ,}{\textbf{43} \lecon{ἐχέτω} R (avec la glose interne ἐπεχέτω)}
\ligneRav{44}{}{\hspace{1.1em}κῦμα δὲ πόντου μὴ κελαδείτω}{\textbf{44} \lecon{δὲ} R : τε Austin, cf. v.46, 47}
\ligneRav{45}{—}{\hspace{1.1em}γλαυκόν — : βομβάξ· : σίγα· τί λέγει\inRas{ς}{};}{\textbf{45} λέγει\inRas{ς}{} \Rpc}
\ligneRav{}{θερ/}{\hspace{1.1em}πτηνῶν τε γένη κατακοιμάσθω,}{}
\ligneRav{}{}{\hspace{1.1em}θηρῶν τ᾽ ἀγρίων πόδες ὑλοδρόμων}{}
\ligneRav{}{—}{\hspace{1.1em}μὴ λυέσθων. : βομβαλοβομβάξ·}{}
\ligneRav{}{θερ/}{\hspace{1.1em}μέλλει γὰρ ὁ καλλιεπὴς Ἀγάθων}{}
\ligneRav{50}{—}{\hspace{1.1em}\difference{πρανοῦν προτέρους} — : μῶν βινεῖσθαι;}{\textbf{50} πράμος R : πρόμος Grynaeus : \lecon{πρανοῦν} Bibi ap. Hesych. π 3188 \espApp ἡμέτερος R : \lecon{προτέρους} Bibi, cf. \italique{Nubes} v.562}
\ligneRav{}{θερ/}{\hspace{1.1em}τίς ὁ φωνήσας; : νήνεμος αἰθήρ·}{}
\ligneRav{}{θερ/}{\hspace{1.1em}δρυόχους, τιθέναι δράματος ἀρχάς·}{}
\ligneRav{53}{}{\hspace{1.1em}κάμπτει δὲ νέας \difference{ἁψῖδας} ἐπῶν,}{\textbf{53} ἀσπίδας R : \lecon{ἁψῖδας} Suid. δ 1547, λ 181}
\ligneRav{}{}{\hspace{1.1em}τὰ δὲ τορνεύει, τὰ δὲ κολλομελεῖ,}{}
\ligneRav{55}{}{\hspace{1.1em}καὶ γνωμοτυπεῖ κἀντονομάζει}{}
\ligneRav{56}{}{\hspace{1.1em}καὶ κηροχυτεῖ καὶ \difference{γογγύλλει}}{\textbf{56} γογγύλίζει R, Suid. γ 361 : \lecon{γογγύλλει} Porson : γογγυλίει Hermann (ad \italique{Nub.} 992)}
\ligneRav{}{—}{\hspace{1.1em}καὶ χοανεύει— : καὶ λαικάζει·}{}
\ligneRav{58}{θερ/}{\hspace{1.1em}τίς \difference{ἀγροιώτας} πελάθει θριγκοῖς;}{\textbf{58} ἀγριώτας R : \lecon{ἀγροιώτας} anonym. Parisinus}
\ligneRav{}{μ\auDessus{ν}{Η}}{\hspace{1.1em}ὃς ἕτοιμος σοῦ τοῦ τε ποιητοῦ}{}
\ligneRav{60}{}{\hspace{1.1em}τοῦ καλλιεποῦς \expo{κατὰ} τοῦ θριγκοῦ}{\textbf{60} κατὰ om. \Rac : add. \Rpc}
\ligneRav{61}{}{\hspace{1.1em}\difference{συγγογγύλας} καὶ συστρέψας}{\textbf{61} γογγυλίσας R : \lecon{συγγογγύλας} Enger}
\ligneRav{}{}{\hspace{1.1em}τουτὶ τὸ πέος χοανεῦσαι·}{}
\ligneRav{}{θερ/}{ἦ που νέος γ᾽ ὢν ἦσθ᾽ ὑβριστὴς ὦ γέρον·}{}
\ligneRav{}{\auDessus{εὐ}{ρ}}{ὦ δαιμόνιε τοῦτον μὲν ἔα χαίρειν, σὺ δὲ}{}
\ligneRav{65}{}{Ἀγάθωνά μοι δεῦρ᾽ ἐκκάλεσον πάσῃ τέχνῃ·}{}
\ligneRav{}{θερ/}{μηδὲν ἱκέτευ᾽· αὐτὸς γὰρ ἔξεισιν τάχα·}{}
\ligneRav{}{}{καὶ γὰρ μελοποιεῖν ἄρχεται· χειμῶνος οὖν}{}
\ligneRav{}{}{ὄντος κατακάμπτειν τὰς στροφὰς οὐ ῥᾴδιον,}{}
\ligneRav{69}{}{ἢν μὴ προίῃ θύρασι πρὸς τὸν ἥλιον·}{\textbf{69} \lecon{θύρασι} R : θύραζε Zanetti}
\ligneRav{70}{μ\auDessus{ν}{Η}}{τί οὖν ἐγὼ δρῶ; : περίμεν᾽, ὡς ἐξέρχεται·}{}
\ligneRav{71}{μ\auDessus{ν}{Η}}{ὦ Ζεῦ τί δρᾶσαι διανοεῖ με τήμερον;}{\textbf{71-4} (Μνη.) (71) et (Εὐρ.) (72-4) R : \lecon{(Μνη.) (71-4)} Bibi ap. Σ : alii alia}
\ligneRav{}{}{νὴ τοὺς θεοὺς ἐγὼ πυθέσθαι βούλομαι}{}
\ligneRav{}{}{τί τὸ πρᾶγμα τουτί. τί στένεις; τί δυσφορεῖς;}{}
\ligneRav{74}{}{οὐ χρῆν σε κρύπτειν ὄντα κηδεστὴν ἐμόν·}{\textbf{74} \lecon{χρῆν} R, Suid. κ 1487 codd. FV : χρή Suid. codd. AGM \espApp \lecon{ἐμόν} R : ἐμέ Valckenaer ap. Dobree (p. 237)}
\ligneRav{75}{\auDessus{εὐ}{ρ}}{ἔστιν κακόν μοι μέγα τι προπεφυραμένον·}{}
\ligneRav{}{μ\auDessus{ν}{Η}}{ποῖόν τι; : τῇδε θἠμέρᾳ κριθήσεται}{}
\ligneRav{}{}{εἴτ᾽ ἔστ᾽ ἔτι ζῶν εἴτ᾽ ἀπόλωλ᾽ Εὐριπίδης·}{}
\ligneRav{}{μ\auDessus{ν}{Η}}{καὶ πῶς; ἐπεὶ νῦν γ᾽ οὔτε τὰ δικαστήρια}{}
\ligneRav{}{}{μέλλει δικάζειν οὔτε βουλῆς ἐσθ᾽ ἕδρα,}{}
}
\vspace{3pt}{\color{ruleGold}\hrule height 0.6pt}\vspace{3pt}
{\fontsize{8.0pt}{9.4pt}\selectfont
\noindent{\bfseries\scshape\color{rubricRed}Loci paralleli et testimonia (f. 163r, vv. 43–79) :} \textbf{43-5} (ἐχέτω -- γλαυκόν) Suidas ε 3792 ; Zonaras 941 ἐχέτω \espTest \textbf{44} (κελαδείτω) Suidas κ 1278 \espTest \textbf{45} (βομβάξ) Hesychius β 787 βομβάξ ; Suidas ε 2807 ; \Eustathe 1718.9 \espTest \textbf{48} (βομβαλοβομβάξ) Suidas ε 2807 ; \Eustathe 790.13, 1718.13 \espTest \textbf{49} (καλλιεπὴς) Philostrate \italique{Vit. soph.} I 9.1 [= II 12.8 Kayser] \espTest \textbf{50} (πραῖνοι) Hesychius π 3188 †πραῖνοι· πρηνίζειν, καταστρέφειν \espTest \textbf{52-7} Suidas δ 1547 (sed τορεύει pro τορνεύει, χιανεύει pro χοανεύει) \espTest \textbf{52} Zonaras 569 δρύοχοι \espTest \textbf{53} (κάμπτει) Hesychius κ 617 κάμπτειν \espTest \textbf{54} (τορνεύει) Hesychius τ 1168 τορνεύει \espTest \textbf{56} (γογγύλλει) Hesychius γ 764 γογγύλλειν \espTest Suidas γ 361 = Zonaras 449 γογγυλίζει \espTest \textbf{57} (χοανεύει) \Eustathe 1153.40 \espTest \textbf{58} (θριγκοῖς) Hesychius θ 750 θριγκοῖς· φραγμοῖς \espTest \textbf{62} (χοανεῦσαι = χωνεῦσαι) Moeris χ 21 ; Hesychius χ 575 χοανεῦσαι \espTest \textbf{74} Suidas κ 1487, cf. Moeris κ 38\par\vspace{2.5pt}
{\color{ruleGold}\hrule height 0.35pt}\vspace{2.5pt}
\noindent{\bfseries\scshape\color{rubricRed}Suite (en commentaire) de l'apparat critique et paléographique :} \textbf{Collation avec le fac-similé phototypique de 1904 (BnF, \textit{FOL-FAC SIM-182 (9)}) :} La planche phototypique de 1904 du \textit{f. 163r} confirme dans la marge intérieure gauche la parfaite lisibilité des sigles tachygraphiques de personnages ({\policeGr μν}\expo{{\policeGr η}}, {\policeGr εὐ}\expo{{\policeGr ρ}}, θερ/) et des \textit{paragraphoi} des vv. 43–79, ainsi que le tracé de la scholie marginale droite au v. 68 ({\policeGr Σ}\expo{R}). --- \textbf{v. 43 :} Glose interlinéaire/interne dans R : ἐπεχέτω glosant ἐχέτω. --- \textbf{v. 45 :} Grattage du sigma final dans λέγει\inRas{ς}{} par \Rpc{}. --- \textbf{v. 50 :} Lieu célèbre de corruption paléographique (πράμος ἡμέτερος R), magistralement résolu par Louis Hauseux (\textbf{Bibi}). --- \textbf{v. 60 :} La préposition κατὰ, omise par \Rac{}, a été rajoutée par \Rpc{} au-dessus de la ligne. --- \textbf{v. 43 :} \lecon{ἐχέτω} R (avec la glose interne ἐπεχέτω). --- \textbf{v. 44 :} \lecon{δὲ} R (maintenu par \textbf{Bibi}) : τε Austin (cf. vv. 46, 47). --- \textbf{v. 50 :} πράμος R (πρόμος Grynaeus) : \lecon{πρανοῦν} \lecon{Bibi} ap. Hésychius π 3188 (et πραῖνοι dans l'\italique{Index}) \espApp ἡμέτερος R : \lecon{προτέρους} \lecon{Bibi}, cf. \italique{Nuées} v. 562. \textbf{Explication paléographique (Commentaire ad v. 50) :} Le texte primitif πρανοῦν προτέρους (ou πρανοι ν) a subi : 1° une médivision avec confusion minuscule ν → μ et ι → ς donnant πράμος ; 2° la syllabe προ- de προτέρους, écrite au moyen du signe tachygraphique en forme de Ʒ inversé (semblable à un ε, cf. Gardthausen II p. 340), a été rattachée au ν précédent et lue <η>μ- + -ετερος {\policeFr →} ἡμέτερος ! --- \textbf{v. 53 :} ἀσπίδας R (« boucliers », faute par attraction d'un mot militaire banal) : \lecon{ἁψῖδας} Souda δ 1547, λ 181 (« cintres, voûtes d'hymnes »). --- \textbf{v. 60 :} κατὰ om. \Rac{} : \textbf{add.} \Rpc{}. --- \textbf{v. 69 :} \lecon{θύρασι} R (maintenu par \textbf{Bibi}) : θύραζε Zanetti. --- \textbf{vv. 71–74 :} (Μνη.) (71) et (Εὐρ.) (72–74) R : \lecon{(Μνη.) (71–74)} \lecon{Bibi} d'après {\policeGr Σ}\expo{R} (toute la tirade appartient à Mnésiloque).\par}
\end{minipage}\end{adjustbox}
\newpage
\renewcommand{\currHeadLE}{SCHOLIES MARGINALES ({\policeGr Σ}\expo{R}) \& GLOSES INTERLINÉAIRES · FOLIO 02/30}
\renewcommand{\currHeadRE}{f. 163r (Classense 165r) · vv. 43–79 · {\policeGr Σ} 022 à {\policeGr Σ} 038}
\renewcommand{\currHeadLO}{SCHOLIES MARGINALES ({\policeGr Σ}\expo{R}) \& GLOSES INTERLINÉAIRES · FOLIO 02/30}
\renewcommand{\currHeadRO}{f. 163r (Classense 165r) · vv. 43–79 · {\policeGr Σ} 022 à {\policeGr Σ} 038}
\bandeauScholies{Folio 02/30 — f. 163r}{Classense 165r}{vv. 43–79}{Recto (page de droite) · 17 entrées ({\policeGr Σ} 022 à {\policeGr Σ} 038)}
{\fontsize{8.8pt}{10.5pt}\selectfont\noindent\textbf{\color{rubricRed}Disposition topographique des marges sur le f. 163r :} Marge extérieure à droite ; marge intérieure à gauche. — \textit{Fac-similé de référence :} \texttt{folios\_individuels\_ff162v-177r/02\_f163r\_classense165r\_vv0043-0079.jpg}.\par\vspace{3pt}}
\entreeScholie{022}{43}{ἐχέτω δὲ πνοὰς νήνεμος αἰθήρ,}{Interlinéaire (supra lineam au-dessus d'ἐχέτω, haut du f. 163r)}{ἐπεχέτω.}{ἐπεχέτω {\policeFr ÷}}{« Qu'il retienne (qu'il suspende ses souffles). »}
\commScholie{Précise la valeur du simple poétique ἐχέτω employé pour le composé ἐπεχέτω dans l'invocation paratragique du Serviteur.}
\entreeScholie{023}{45}{βομβάξ.}{Marge intérieure gauche (en face du v. 45, f. 163r)}{ἐπίρρημα ἐπὶ θαυμασμοῦ λαμβανόμενον.}{ἐπίρρημα ἐπὶ θαυμασμ\expo{οῦ} λαμβανόμεν\expo{ον} {\policeFr ÷}}{« Adverbe (interjection) employé pour marquer l'étonnement (ironique : “Boum ! Peste !”). »}
\commScholie{Repris par la Souda β 369 et Hésychius β 790. Souligne la rupture bouffonne introduite par Mnésiloque au milieu des dimètres anapestiques solennels du Serviteur.}
\entreeScholie{024}{48}{βομβαλοβομβάξ.}{Marge intérieure gauche (en face du v. 48, f. 163r)}{ἀναδιπλασιάζει τὴν λέξιν.}{ἀναδιπλασιάζ(ει) τ(ὴν) λέξ(ιν) {\policeFr ÷}}{« Il redouble comiquement le mot (par réduplication burlesque : “Boum-bala-boum !”). »}
\commScholie{Repris par la Souda β 368. Le scholiaste note le procédé comique d'amplification syllabique.}
\entreeScholie{025}{50}{μέλλει γὰρ ὁ καλλιεπὴς Ἀγάθων / πρανοῦν προτέρους [πράμος ἡμέτερος R] — μῶν βινεῖσθαι;}{Marge intérieure gauche (f. 163r, en face du v. 50, Planche 22)}{λείπει ἕτοιμός ἐστιν.}{λείπ(ει) ἕτοιμός ἐστ(ιν) {\policeFr ÷}}{« Il manque (par ellipse) “il est prêt à” (ou dépend de μέλλει interrompu). »}
\commScholie{Glose marginale de R\expo{2} (décrite dans Manuel\_Dechiffrement, Planche 22, Exercice 3) expliquant l'interruption syntaxique de la phrase du Serviteur (μέλλει...) par la question obscène de Mnésiloque (μῶν βινεῖσθαι;).}
\entreeScholie{026}{52}{δρυόχους τιθέναι δράματος ἀρχάς·}{Marge droite (extérieure, f. 163r, en face du v. 52)}{προοίμια. κυρίως δὲ δρύοχοί εἰσιν οἱ ἐντιθέμενοι πάτταλοι ναυπηγουμένης νεώς.}{προοίμια · κυρί\expo{ως} δ(ὲ) δρύοχοί εἰσ(ιν) οἱ ἐντιθέμεν(οι) πάτταλ(οι) ναυπηγουμέν\expo{ης} νεώς {\policeFr ÷}}{« Les préludes (les fondements). Au sens propre, les δρύοχοι sont les étais (ou quilles de chêne) que l'on pose lors de la construction d'un navire. »}
\commScholie{Repris par la Souda δ 1552 et commenté par Louis Hauseux (vv 1-641.tex, ll. 1030-1042) à propos de la série des sept métaphores artisanales appliquées à la composition poétique d'Agathon.}
\entreeScholie{027}{53}{κάμπτει δὲ νέας ἁψῖδας ἐπῶν,}{Marge droite (extérieure, f. 163r, sous Σ 52)}{ὡς ἀπὸ τῶν οἰκοδομούντων.}{ὡς ἀπ(ὸ) τ\expo{ῶν} οἰκοδομούντ\expo{ων} · [ἀν\auDessus{τ}{ι} τ\expo{οῦ} ἐκ ξύλ\expo{ου} κάμπτ(ει)] {\policeFr ÷}}{« Métaphore tirée des bâtisseurs (ou des charrons qui cintrent les jantes de roues en bois). »}
\commScholie{Le scholiaste explicite le passage de la charpenterie navale (v. 52) à l'architecture et au charronnage (ἁψῖδας, cintres ou jantes).}
\entreeScholie{028}{54}{τὰ δὲ τορνεύει, τὰ δὲ κολλομελεῖ,}{Marge droite (extérieure) et glose interlinéaire sur κολλομελεῖ}{ὡς ἐπὶ τεκτόνων. μέλη κολλᾷ.}{ὡς ἐπὶ τεκτόν\expo{ων} {\policeFr ÷} // μέλη κολλᾷ {\policeFr ÷}}{« Métaphore prise des menuisiers-tourneurs. — Il colle ensemble des chants (ou ajuste des membres mélodiques). »}
\commScholie{Explique le composé comique κολλο-μελεῖ, forgé par Aristophane sur κολλάω (coller, assembler en ébénisterie) et μέλος (membre / chant lyrique).}
\entreeScholie{029}{55}{καὶ γνωμοτυπεῖ κἀντονομάζει}{Interlinéaire (supra lineam au-dessus de κἀντονομάζει)}{ὄνομα ἀντὶ ὀνόματος τίθησιν.}{ὄνομα ἀν\auDessus{τ}{ι} ὀνόματ\expo{ος} τίθησ(ιν) {\policeFr ÷}}{« Il met un mot à la place d'un autre (procédé rhétorique de l'antonomase et de la métaphore précieuse). »}
\commScholie{Repris par la Souda α 2760. Désigne la recherche gorgianique des tropes et substitutions lexicales propre au style d'Agathon.}
\entreeScholie{030}{56}{καὶ κηροχυτεῖ καὶ γογγύλλει [γογγυλίζει R]}{Marge droite (extérieure) et glose interlinéaire sur γογγυλίζει}{ὡς ἐπιπηγνύντος καὶ συντιθέντος. μεταστρέφει.}{ὡς ἐπιπηγνύντ\expo{ος} \textbf{ς} συντιθέντ\expo{ος} {\policeFr ÷} // μεταστρέφ(ει) {\policeFr ÷}}{« Comme quelqu'un qui modèle de la cire fondue et l'assemble. — Il arrondit et tourne en tous sens (μεταστρέφει). »}
\commScholie{Rapproché par Louis Hauseux (vv 1-641.tex, l. 1130 n. 13) d'Hésychius γ 764 (γογγύλλειν· συστρέφειν) et de la Souda γ 364. Sur R\expo{1}, la leçon amétrique γογγυλίζει a été corrigée par Biset (adopté par Bibi) en γογγύλλει.}
\entreeScholie{031}{57}{καὶ χοανεύει —}{Interlinéaire (supra lineam au-dessus de χοανεύει)}{χωνεύει.}{χωνεύει {\policeFr ÷}}{« Il coule au creuset (forme contracte courante χωνεύει pour la forme ouverte poétique χοανεύει). »}
\commScholie{Achève la métaphore de la fonte à la cire perdue (κηροχυτεῖ ... χοανεύει) étudiée en détail par L. Hauseux.}
\entreeScholie{032}{58}{τίς ἀγροιώτας πελάθει θριγκοῖς;}{Interlinéaire (supra lineam au-dessus de θριγκοῖς)}{τοῖς φραγμοῖς τῶν οἴκων.}{τ\expo{οῖς} φραγμ\expo{οῖς} τ\expo{ῶν} οἴκ\expo{ων} {\policeFr ÷}}{« Aux clôtures (au chaperon ou corniche) de la demeure. »}
\commScholie{Glose le terme tragique θριγκοῖς (corniche supérieure du mur d'enceinte).}
\entreeScholie{033}{60}{τοῦ καλλιεποῦς κατὰ [add. R\expo{2}] τοῦ θριγκοῦ}{Marge droite (extérieure, f. 163r, en face du v. 60)}{ἀντὶ τοῦ ἐπὶ τοῦ θριγκοῦ.}{ἀν\auDessus{τ}{ι} τ\expo{οῦ} ἐπὶ τ\expo{οῦ} θριγκ\expo{οῦ} {\policeFr ÷}}{« Au sens de : sur la corniche (contre la clôture). »}
\commScholie{Dans R, le réviseur R\expo{2} a ajouté κατὰ au-dessus de la ligne devant τοῦ θριγκοῦ (tandis que M porte uniquement τοῦ θριγκοῦ) ; cette glose explique la construction locative de θριγκοῦ.}
\entreeScholie{034}{61}{συγγογγύλας [γογγυλίσας R] καὶ συστρέψας}{Interlinéaire (supra lineam au-dessus de συγκύψας)}{συγκάμψας.}{συγκάμψ\expo{ας} {\policeFr ÷}}{« L'ayant fait se courber en deux (l'ayant plié). »}
\commScholie{Glose interlinéaire de R\expo{2} expliquant le participe métaphorique συγγογγύλας (R γογγυλίσας) / συστρέψας dans la menace obscène de Mnésiloque : « ayant courbé, replié en boule ».}
\entreeScholie{035}{62}{τουτὶ τὸ πέος χοανεῦσαι·}{Marge droite (extérieure, f. 163r, en face du v. 62)}{ὡς εἰς χόανον τὸν πρωκτὸν ἐμβαλεῖν.}{ὡς εἰς χόαν\expo{ον} τ(ὸν) πρωκτ(ὸν) ἐμβαλεῖν {\policeFr ÷}}{« Pour ainsi dire : enfourner dans son derrière comme dans un creuset de fondeur (χόανον). »}
\commScholie{Explique le retournement bouffon par Mnésiloque des métaphores métallurgiques du Serviteur (χοανεύει v. 57, γογγυλίσας v. 61).}
\entreeScholie{036}{68}{καὶ γὰρ μελοποιεῖν ἄρχεται· χειμῶνος οὖν / ὄντος κατακάμπτειν τὰς στροφὰς οὐ ῥᾴδιον, / ἢν μὴ προίῃ θύρασι [θύραζε Zanetti] πρὸς τὸν ἥλιον·}{Marge droite (extérieure, f. 163r, bloc de 6 lignes en face des vv. 67–69)}{ὡς ἀπὸ τῶν λυόντων τὸν κηρὸν ἐν τῷ ἡλίῳ. ἐπειδὴ δύσκαμπτοι αἱ στροφαί. καὶ Ἀριστοφάνης [Nub. 971] τὰς δυσκολοκάμπτους αὐτάς φησιν. εἴρηται δὲ καὶ διὰ τὴν φωνὴν, ἐπεὶ λέγεταί τι κάμπτειν καὶ ἐν ταῖς ᾠδαῖς.}{ὡς ἀπ(ὸ) τ\expo{ῶν} λυόντ\expo{ων} τ(ὸν) κηρ(ὸν) ἐν τῷ ἡλίῳ · ἐπειδὴ δύσκαμπτοι αἱ στροφαί · \textbf{ς} Ἀριστοφάνης τ\expo{ὰς} δυσκολοκάμπτ\expo{ους} αὐτάς φησ(ιν) · εἴρητ(αι) δ(ὲ) \textbf{ς} διὰ τ(ὴν) φων(ὴν), ἐπεὶ λέγεταί τι κάμπτ(ειν) \textbf{ς} ἐν τ\expo{αῖς} ᾠδ\expo{αῖς} {\policeFr ÷}}{« Métaphore prise de ceux qui ramollissent la cire au soleil, parce que les “strophes” (les tournants / modulations musicales) sont difficiles à assouplir. Aristophane lui-même, dans les Nuées (v. 971), les appelle des modulations “difficiles à infléchir” (δυσκολοκάμπτους). Et cela est dit aussi à propos de la voix, puisque l'on parle d'“infléchir” (κάμπτειν) dans les chants. »}
\commScholie{Excellente scholie ancienne (reprise par la Souda κ 577) qui éclaire le double sens technique de στροφαί et de κατακάμπτειν (assouplir la cire ou le bois au soleil / exécuter les vocalises sinueuses du nouveau dithyrambe), en citant le parallèle exact des Nuées v. 971.}
\entreeScholie{037}{71}{ὦ Ζεῦ τί δρᾶσαι διανοεῖ με τήμερον;}{Marge intérieure gauche (f. 163r, en face du v. 71)}{στενάζει φοβούμενος τὰς γυναῖκας.}{στενάζ(ει) φοβούμεν\expo{ος} τ\expo{ὰς} γυναῖκ\expo{ας} {\policeFr ÷}}{« Il gémit par crainte des femmes. »}
\commScholie{Observation capitale exploitée par Louis Hauseux (Bibi, vv. 70-74, ll. 1169-1185) : alors que R attribue le v. 71 à Mnésiloque et les vv. 72-74 à Euripide (tandis que les éditeurs modernes donnent tout à Euripide), {\policeGr Σ}\expo{R} 71 et {\policeGr Σ}\expo{R} 74 montrent que dans l'édition alexandrine commentée par le scholiaste, Mnésiloque prononçait les vv. 71-74 ou faisait un jeu de scène muet (τὸ σημεῖον).}
\entreeScholie{038}{74}{ὄντα κηδεστὴν ἐμόν.}{Marge droite (extérieure, bas du f. 163r)}{τὸ σημεῖον· οὐ μόνον κηδεστὴς ὁ ἐκδιδοὺς, ἀλλὰ καὶ ὁ λαμβάνων.}{τ(ὸ) σημεῖ\expo{ον} · οὐ μόν\expo{ον} κηδεστ\expo{ὴς} ὁ ἐκδιδ\expo{οὺς}, ἀλλὰ \textbf{ς} ὁ λαμβάν\expo{ων} {\policeFr ÷}}{« Le signe critique (τὸ σημεῖον, la diplè ou le chi) est apposé ici : on appelle κηδεστής (parent par alliance) non seulement celui qui donne sa fille en mariage (le beau-père), mais aussi celui qui la reçoit (le gendre). »}
\commScholie{Scholie alexandrine (Aristarque / Didyme) signalant la présence d'un signe critique marginal (τὸ σημεῖον) et précisant la réciprocité du terme κηδεστής en attique (reprise par la Souda κ 1485, Photius κ 651 et Hésychius κ 2441 ; discutée par L. Hauseux, Introduction.tex, l. 327, et vv 1-641.tex, ll. 1169-1184).}
\clearpage
\renewcommand{\currHeadLE}{ΑΡΙΣΤΟΦΑΝΟΥΣ ΘΕΣΜΟΦΟΡΙΑΖΟΥΣΑΙ · CODEX RAVENNAS 429 (\R) · FOLIO 03/30}
\renewcommand{\currHeadRE}{f. 163v (Classense 165v) · vv. 80–122 · Collation diplomatique}
\renewcommand{\currHeadLO}{ΑΡΙΣΤΟΦΑΝΟΥΣ ΘΕΣΜΟΦΟΡΙΑΖΟΥΣΑΙ · CODEX RAVENNAS 429 (\R) · FOLIO 03/30}
\renewcommand{\currHeadRO}{f. 163v (Classense 165v) · vv. 80–122 · Collation diplomatique}
\noindent\begin{adjustbox}{max totalheight=30.4cm, max width=\textwidth, keepaspectratio}\begin{minipage}{\textwidth}
\bandeauRav{f. 163v}{Classense 165v}{80–122}{Folio 03/30 · Verso (page de gauche) · 13 scholies \& gloses ({\policeGr Σ} 039–051)}
{\fontsize{8.8pt}{10.2pt}\selectfont
\ligneRav{80}{μ\auDessus{ν}{Η}}{\difference{ἐπείπερ ἐστὶ} Θεσμοφορίων ἡ μέση·}{\textbf{80} ἐπεὶ τρίτη 'στὶ R : \lecon{ἐπείπερ ἐστὶ} Nauck (\italique{Mél.} (1866), II pp.468-71)}
\ligneRav{}{\auDessus{εὐ}{ρ}}{τοῦτ᾽ αὐτὸ γὰρ τοι κἀπολεῖν με προσδοκῶ·}{}
\ligneRav{}{}{αἱ γὰρ γυναῖκες ἐπιβεβουλεύκασί μοι}{}
\ligneRav{83}{}{κἀν Θεσμοφόροιν μέλλουσι περί \difference{μου} τήμερον}{\textbf{83} ἐμοῦ R : \lecon{μου} Biset}
\ligneRav{}{—}{ἐκκλησιάζειν ἐπ᾽ ὀλέθρῳ. : τιὴ τί δή;}{}
\ligneRav{85}{\auDessus{εὐ}{ρ}}{ὁτιὴ τραγῳδῶ καὶ κακῶς αὐτὰς λέγω·}{}
\ligneRav{86}{μ\auDessus{ν}{Η}}{νὴ τὸν Ποσειδῶ καὶ \difference{δικαίαν} ἂν πάθοις·}{\textbf{86} δί᾽ και᾽ R : δίκαιά γ᾽ Grynaeus : \lecon{δικαίαν} Wilamowitz}
\ligneRav{}{}{ἀτὰρ τίν᾽ ἐκ ταύτης σὺ μηχανὴν ἔχεις;}{}
\ligneRav{}{\auDessus{εὐ}{ρ}}{Ἀγάθωνα πεῖσαι τὸν τραγῳδοδιδάσκαλον}{}
\ligneRav{}{—}{ἐς Θεσμοφόροιν ἐλθεῖν. : τί δράσοντ᾽; εἰπέ μοι·}{}
\ligneRav{90}{\auDessus{εὐ}{ρ}}{ἐκκλησιάσοντ᾽ ἐν ταῖς \difference{γυναιξὶ} \difference{χἂν} δέῃ}{\textbf{90} γυναιξίν R : \lecon{γυναιξὶ} \Mon \espApp κἂν R : \lecon{χἂν} Markland (1763) ad Eur. \italique{Suppl.} 364}
\ligneRav{91}{—}{λέξονθ᾽ ὑπὲρ ἐμοῦ. : πότερα φανερὸν ἢ λάθρᾳ;}{\textbf{91} \lecon{φανερὸν} R : φανερῶς Meineke}
\ligneRav{}{\auDessus{εὐ}{ρ}}{λάθρᾳ, στολὴν γυναικὸς ἠμφιεσμένον·}{}
\ligneRav{}{μ\auDessus{ν}{Η}}{τὸ πρᾶγμα κομψὸν καὶ σφόδρ᾽ ἐκ τοῦ σοῦ τρόπου·}{}
\ligneRav{}{}{τοῦ γὰρ τεχνάζειν ἡμέτερος ὁ πυραμοῦς·}{}
\ligneRav{95}{\auDessus{εὐ}{ρ}}{σίγα. : τί δ᾽ ἔστιν; : Ἁγάθων ἐξέρχεται·}{\textbf{95} \lecon{τί δ᾽} R : τί Kaibel}
\ligneRav{96}{μ\auDessus{ν}{Η}}{καὶ \difference{ποῖός} ἐστιν; : οὗτος οὑκκυκλούμενος·}{\textbf{96} ποῖο ἐστὶν οὗτος — οὑκκυκλούμενος R : ποῖος ἐστὶν οὗτος — οὑκκυκλούμενος \Mon : ποῦ ᾽σθ᾽; — ὅπου ᾽στίν; οὗτος οὑκκυκλούμενος Meineke : \lecon{ποῖος ἐστὶν; — οὗτος οὑκκυκλούμενος} Tyrwhitt (p. 65)}
\ligneRav{97}{μ\auDessus{ν}{Η}}{ἀλλ᾽ \difference{ἦ} τυφλὸς μέν εἰμ᾽· ἐγὼ γὰρ οὐχ ὁρῶ}{\textbf{97} ἢ R : \lecon{ἦ} Biset}
\ligneRav{}{}{ἄνδρ᾽ οὐδέν᾽ ἐνθάδ᾽ ὄντα, Κυρήνην δ᾽ ὁρῶ·}{}
\ligneRav{99}{\auDessus{εὐ}{ρ}}{σίγα· μελῳδεῖν \difference{νῦν} παρασκευάζεται·}{\textbf{99} ἂν R : \lecon{νῦν} Rogers : alii alia}
\ligneRav{}{}{\hspace{2.2em}\textbf{[παρεπιγραφή R\expo{1} :]} \lecon{μινυρισμός}}{\textbf{post 99} \lecon{μινυρισμός} (παρεπιγραφή) R, {\policeGr Σ}\expo{R}}
\ligneRav{100}{μ\auDessus{ν}{Η}}{\difference{Μύρμηκος} ἀτραπούς, ἢ τί διαμινύρεται;}{\textbf{100} μύρμ- R : \lecon{Μύρμ-} Bibi ap Hesych. μ 1904}
\ligneRav{101}{ἀγάθ/}{\hspace{1.1em}ἱερὰν \difference{Χθονίαιν} δε-}{\textbf{101} χθονίαις R : \lecon{χθονίαιν} Meineke}
\ligneRav{}{}{\hspace{1.1em}ξάμεναι λαμπάδα κοῦραι}{}
\ligneRav{103}{}{\hspace{1.1em}ξὺν ἐλευθέρᾳ πατρί-}{\textbf{103} \lecon{πατρίδι} R : πραπίδι Wecklein}
\ligneRav{}{}{\hspace{1.1em} δι χορεύσασθε βοάν·}{}
\ligneRav{}{}{\hspace{1.1em}τίνι δαιμόνων ὁ κῶμος;}{}
\ligneRav{105}{}{\hspace{1.1em}λέγε \difference{νιν} \difference{εὐπείστως} δὲ τοὐμὸν}{\textbf{105} νυν R : \lecon{νιν} Dindorf \espApp εὐπίστως R, Suid. σ 178, 183 : \lecon{εὐπείστως} Reiske}
\ligneRav{106}{}{\hspace{1.1em}δαίμονας \difference{ἔχει} σεβίσαι·}{\textbf{106} ἔχεις R : \lecon{ἔχει} Suid.}
\ligneRav{107}{}{\hspace{1.1em}ἄγε νῦν \difference{ὄλβιζε} Μοῦσα}{\textbf{107} ὅπλιζε R : \lecon{ὄλβιζε} Bentley \espApp \lecon{Μοῦσα} R : μούσᾳ Bergk}
\ligneRav{}{}{\hspace{1.1em}χρυσέων ῥύτορα τόξων}{}
\ligneRav{}{}{\hspace{1.1em}Φοῖβον, ὃς ἱδρύσατο χώρας}{}
\ligneRav{110}{}{\hspace{1.1em}γύαλα Σιμουντίδι γᾷ·}{}
\ligneRav{}{}{\hspace{1.1em}χαῖρε καλλίσταις ἀοιδαῖς}{}
\ligneRav{}{}{\hspace{1.1em}Φοῖβ᾽ ἐν εὐμούσοισι τιμαῖς}{}
\ligneRav{}{}{\hspace{1.1em}γέρας ἱερὸν προφέρων·}{}
\ligneRav{}{}{\hspace{1.1em}τάν τ᾽ ἐν ὄρεσι δρυογόνοισι}{}
\ligneRav{115}{}{\hspace{1.1em}κόραν \difference{ἀείσατ᾽} Ἄρτεμιν ἀγροτέραν·}{\textbf{115} ἀεισαντ᾽ R : \lecon{ἀείσατ᾽} Zanetti (\italique{laudate} Divus)}
\ligneRav{117}{}{\hspace{1.1em}ἕπομαι κλῄζουσα \difference{σεμνὰν}}{\textbf{117} σεμνὸν R : \lecon{σεμνὰν} Dindorf}
\ligneRav{}{}{\hspace{1.1em}γόνον ὀλβίζουσα Λατοῦς}{}
\ligneRav{}{}{\hspace{1.1em}Ἄρτεμιν ἀπειρολεχῆ·}{}
\ligneRav{}{}{\hspace{1.1em}Λατώ τε κρούματά τ' Ἀσιάδος ποδὶ}{}
\ligneRav{120}{}{\hspace{1.1em}παράρυθμ' εὔρυθμα Φρυγίων}{}
\ligneRav{122}{}{\hspace{1.1em}\difference{διὰ νεύματα} Χαρίτων·}{\textbf{122} διανεύματα R : δινεύματα Bentley : \lecon{διὰ νεύματα} Fritzsche}
}
\vspace{3pt}{\color{ruleGold}\hrule height 0.6pt}\vspace{3pt}
{\fontsize{8.0pt}{9.4pt}\selectfont
\noindent{\bfseries\scshape\color{rubricRed}Loci paralleli et testimonia (f. 163v, vv. 80–122) :} \textbf{80} Hesychius τ 1440 τρίτη Θεσμοφορίων \espTest \textbf{93-4} Suidas τ 439 \espTest \textbf{93} Suidas κ 2025 \espTest \textbf{94} (ἡμέτερος ὁ πυραμοῦς) Suidas π 3192, cf. Tzetzes \italique{Chil.} XIII 281-288 \espTest \textbf{96} (οὑκκυκλούμενος) Zonaras 678 ἐκκυκλεῖ \espTest \textbf{100} Hesychius μ 1904 Μύρμηκος ἀτραπούς ; Photius μ 606 Μύρμηκος ἀτραπός ; Suidas μ 1445 \espTest \textbf{104-6} (εὐπίστως -- σεβίσαι) Suidas σ 183 \espTest \textbf{115-6} (Ἄρτεμιν ἀγροτέραν) \Eustathe 361.37 \espTest \textbf{120} \EtymMagnum 153.21 Ἀσιάδος κρούματα· τῆς κιθάρας. οὔτως Αριστοφάνης εἶπε, παρῳδῶν τὸ ἐξ Ἐρεχθέως Εὐριπίδου ; \SynB α 2226 = Photius α 2957 = Suidas 4147 = \EtymGen (AB) α 1271 ed. Lasserre-Livadaras\par\vspace{2.5pt}
{\color{ruleGold}\hrule height 0.35pt}\vspace{2.5pt}
\noindent{\bfseries\scshape\color{rubricRed}Suite (en commentaire) de l'apparat critique et paléographique :} \textbf{Première \italique{parepigraphè} (entre vv. 99 et 100) :} \Run{} copie dans la colonne le mot scénique \lecon{μινυρισμός} (« fredonnement d'Agathon »), accompagné de la précieuse scholie {\policeGr Σ}\expo{R} : μεταξὺ τῶν δυοῖν ἀξιοῦσί τινες γράφειν 'μινυρισμός', ὡς πολλὰ τοιαῦτα παρεπιγράφεται. --- \textbf{Erreurs d'attribution de \Rdeux{} dans la monodie d'Agathon (vv. 101–129) :} Prenant le dialogue lyrique fictif d'Agathon (qui chante successivement le rôle de la coryphée troyenne et celui du chœur) pour un vrai dialogue dramatique, \Rdeux{} inscrit à tort en marge des \italique{notae personarum} faisant alterner αγαθ/ et χορ/. --- \textbf{v. 80 :} ἐπεὶ τρίτη 'στὶ R : \lecon{ἐπείπερ ἐστὶ} Nauck (\italique{Mélanges gréco-romains}, 1866, II, pp. 468–471), adopté et expliqué paléographiquement par L. Hauseux : dans l'archétype en onciale, l'abréviation \textbf{Π' = περ} dans ΕΠΕΙΠ'ΕΣΤΙ a été lue \textbf{Γ' = τρίτη} (ΕΠΕΙΓ'ΕΣΤΙ), sous l'influence d'une glose savante sur le troisième jour des Thesmophories (la \italique{Nesteia}). --- \textbf{post v. 99 :} \italique{parepigraphè} \lecon{μινυρισμός} R ({\policeGr Σ}\expo{R} μεταξὺ τῶν δυοῖν ἀξιοῦσί τινες γράφειν 'μινυρισμός', ὡς πολλὰ τοιαῦτα παρεπιγράφεται). --- \textbf{v. 100 :} μύρμ- R : \lecon{Μύρμ-} (Μύρμηκος ἀτραπούς) \lecon{Bibi} d'après Hésychius μ 1904. --- \textbf{v. 107 :} ὅπλιζε R (métathèse de consonnes) : \lecon{ὄλβιζε} Bentley \espApp \lecon{Μοῦσα} R (μούσᾳ Bergk).\par}
\end{minipage}\end{adjustbox}
\newpage
\renewcommand{\currHeadLE}{SCHOLIES MARGINALES ({\policeGr Σ}\expo{R}) \& GLOSES INTERLINÉAIRES · FOLIO 03/30}
\renewcommand{\currHeadRE}{f. 163v (Classense 165v) · vv. 80–122 · {\policeGr Σ} 039 à {\policeGr Σ} 051}
\renewcommand{\currHeadLO}{SCHOLIES MARGINALES ({\policeGr Σ}\expo{R}) \& GLOSES INTERLINÉAIRES · FOLIO 03/30}
\renewcommand{\currHeadRO}{f. 163v (Classense 165v) · vv. 80–122 · {\policeGr Σ} 039 à {\policeGr Σ} 051}
\bandeauScholies{Folio 03/30 — f. 163v}{Classense 165v}{vv. 80–122}{Verso (page de gauche) · 13 entrées ({\policeGr Σ} 039 à {\policeGr Σ} 051)}
{\fontsize{8.8pt}{10.5pt}\selectfont\noindent\textbf{\color{rubricRed}Disposition topographique des marges sur le f. 163v :} Marge extérieure à gauche ; marge intérieure à droite. — \textit{Fac-similé de référence :} \texttt{folios\_individuels\_ff162v-177r/03\_f163v\_classense165v\_vv0080-0122.jpg}.\par\vspace{3pt}}
\entreeScholie{039}{80}{ἐπείπερ ἐστὶ [ἐπεὶ τρίτη ᾽στὶ R] Θεσμοφορίων ἡ μέση·}{Marge supérieure et marge extérieure gauche (f. 163v, immense scholie de 18 lignes encadrant le haut de la colonne)}{τοῦτο τῶν ζητουμένων ἐστὶ, πῶς καὶ τρίτην καὶ μέσην εἶπεν. ια' γὰρ ἄνοδος, εἶτα ιβ' νηστεία, εἶτα τρισκαιδεκάτη καλλιγένεια, ὥστε μέση μὲν εἶναι δύναται, τρίτη μὲν οὒ, ἀλλὰ δευτέρα. ἀλλ’ οὐδὲ ψυχρεύεσθαί τις δύναται λέγων ὅτι τρίτη ἡ τρισκαιδεκαταία, ὡς ἑκταία ἑκατηβόλος σελάνα, ἑκκαιδεκαταία· ἡ γὰρ μέση οὐκ ἔστιν ιγ', ἀλλὰ δωδεκάτη. ἡ λύσις οὖν ἥδε. δεκάτῃ ἐν Ἁλιμοῦντι Θεσμοφόρια ἄγεται, ὥστε τρίτην μὲν ἀπὸ δεκάτης ιβ' εἶναι, μέσην δὲ μὴ συναριθμουμένης τῆς δεκάτης. τοῦτο δὲ αἰνιγματῶδες κατὰ Καλλιμάχου ἄν τις φαίη, ἔνθα μὲν τρίτη λέγεται, συναριθμεῖν τὴν ι', ἔνθα δὲ μέση, μηκέτι συναριθμεῖν. καὶ ὅπου γε λιμώττουσιν, ἀστεϊζόμενοι τὴν μέσην τῶν Θεσμοφορίων ἄγειν φασὶν, ἐπεὶ αὕτη ἡ νηστεία. ἑνδεκάτη Πυανεψιῶνος ἄνοδος. δωδεκάτη νηστεία, ἐν ᾗ σχολὴν ἄγουσαι ὑπόκεινται αἱ γυναῖκες ἐκκλησιάζουσαι περὶ Εὐριπίδου. τρισκαιδεκάτη καλλιγένεια. γ' οὖν συναριθμουμένης ι', μέσην δὲ τῆς ἀνόδου καὶ τῆς καλλιγενείας.}{τοῦτ(ο) τ\expo{ῶν} ζητουμένων ἐστ(ὶ), πῶς \textbf{ς} τρίτην \textbf{ς} μέσην εἶπ(εν) · ια' γ(ὰρ) ἄνοδ\expo{ος}, εἶτα ιβ' νηστεία, εἶτα τρισκαιδεκάτη καλλιγένεια, ὥστε μέση μ(ὲν) εἶν(αι) δύνατ(αι), τρίτη μ(ὲν) οὒ, ἀλλὰ δευτέρα · ἀλλ’ οὐδὲ ψυχρεύεσθαί τις δύνατ(αι) λέγων ὅτι τρίτη ἡ τρισκαιδεκαταία, ὡς ἑκταία ἑκατηβόλος σελάνα, ἑκκαιδεκαταία· ἡ γ(ὰρ) μέση οὐκ ἔστ(ιν) ιγ', ἀλλὰ δωδεκάτη · ἡ λύσις οὖν ἥδε · δεκάτῃ ἐν Ἁλιμοῦντι Θεσμοφόρια ἄγετ(αι), ὥστε τρίτην μ(ὲν) ἀπ(ὸ) δεκάτης ιβ' εἶν(αι), μέσην δ(ὲ) μὴ συναριθμουμένης τ\expo{ῆς} δεκάτης · τοῦτ(ο) δ(ὲ) αἰνιγματῶδες κατὰ Καλλιμάχου ἄν τις φαίη, ἔνθα μ(ὲν) τρίτη λέγετ(αι), συναριθμ(εῖν) τ(ὴν) ι', ἔνθα δ(ὲ) μέση, μηκέτι συναριθμ(εῖν) · \textbf{ς} ὅπ\expo{ου} γε λιμώττουσ(ιν), ἀστεϊζόμενοι τ(ὴν) μέσην τ\expo{ῶν} Θεσμοφορίων ἄγ(ειν) φασὶν, ἐπεὶ αὕτη ἡ νηστεία · ἑνδεκάτη Πυανεψιῶν\expo{ος} ἄνοδ\expo{ος} · δωδεκάτη νηστεία, ἐν ᾗ σχολ(ὴν) ἄγουσ(αι) ὑπόκειντ(αι) αἱ γυναῖκες ἐκκλησιάζουσ(αι) π(ερὶ) Εὐριπίδου · τρισκαιδεκάτη καλλιγένεια · γ' οὖν συναριθμουμένης ι', μέσην δ(ὲ) τ\expo{ῆς} ἀνόδου \textbf{ς} τ\expo{ῆς} καλλιγενείας {\policeFr ÷}}{« Ceci fait partie des problèmes classiques (τῶν ζητουμένων) : comment a-t-il pu appeler le même jour à la fois “le troisième” et “celui du milieu” ? En effet, le 11 (de Pyanepsion) est l'Anodos (la Montée), puis le 12 la Nêsteia (le Jeûne), puis le 13 la Kalligéneia (la Belle-Naissance), de sorte que ce jour peut bien être “celui du milieu”, mais non “le troisième”, puisque c'est le deuxième. Et l'on ne saurait non plus soutenir froidement que τρίτη signifie “le treizième jour”, comme ἑκταία (“du sixième jour”) chez Pindare pour “du seizième jour” : car le jour du milieu n'est pas le 13, mais le 12. Voici donc la solution (ἡ λύσις) : le 10, on célèbre les Thesmophories à Halimous, de sorte que le 12 est bien le troisième jour en comptant depuis le 10, tout en étant le jour du milieu à Athènes si l'on ne compte pas le 10. Mais l'on pourrait dire contre Callimaque qu'il serait énigmatique de compter le 10 quand on l'appelle “le troisième”, et de ne plus le compter quand on l'appelle “celui du milieu” ! Aussi, quand des gens jeûnent, disent-ils par plaisanterie urbaine qu'ils célèbrent “le jour du milieu des Thesmophories”, parce que c'est le jour du Jeûne. Le 11 de Pyanepsion : Anodos ; le 12 : Nêsteia, jour où les femmes, étant de loisir, sont supposées tenir l'assemblée au sujet d'Euripide ; le 13 : Kalligéneia. C'est donc le 3e jour en comptant le 10, et le jour du milieu entre l'Anodos et la Kalligéneia. »}
\commScholie{La plus célèbre et la plus longue scholie des Thesmophories, citant nommément Callimaque (fr. 461 Pfeiffer) et reprise par la Souda τ 1030 et Hésychius τ 1440. Comme le démontre Louis Hauseux (Bibi, v. 80, vv 1-641.tex, ll. 1226-1238), toute cette casuistique alexandrine sur le dème d'Halimous provient d'une corruption très ancienne du texte : au lieu de τρίτη 'στὶ, Aristophane avait écrit ἐπὶ τρίτῃ (« en vue du troisième jour, voici le jour du milieu des Thesmophories », ou selon la conjecture de Bibi en apparat).}
\entreeScholie{040}{83}{κἀν Θεσμοφόροιν μέλλουσι περί μου τήμερον}{Marge gauche (extérieure, f. 163v, en face du v. 83)}{κἀν θεσμοφόροιν: Ἀττικὴ σύνταξις καὶ δυϊκὴ, ἀντὶ τοῦ ἐν τῷ τῶν θεσμοφόρων ἱερῷ, Δήμητρος καὶ Περσεφόνης.}{κἀν θεσμοφόροιν : Ἀττικὴ σύνταξ\expo{ις} \textbf{ς} δυϊκὴ, ἀν\auDessus{τ}{ι} τ\expo{οῦ} ἐν τῷ τ\expo{ῶν} θεσμοφόρ\expo{ων} ἱερῷ, Δήμητρ\expo{ος} \textbf{ς} Περσεφόν\expo{ης} {\policeFr ÷}}{« κἀν Θεσμοφόροιν : construction attique au duel (avec ellipse de ἱερῷ), au sens de “dans le sanctuaire des Deux Déesses Thesmophores”, Déméter et Perséphone. »}
\commScholie{Cité par Louis Hauseux (vv 1-641.tex, l. 1239) et repris par la Souda θ 271 : tour elliptique attique régulier (ἐν + génitif duel du nom divin = dans le sanctuaire des Deux Déesses).}
\entreeScholie{041}{93}{τὸ πρᾶγμα κομψὸν καὶ σφόδρ᾽ ἐκ τοῦ σοῦ τρόπου·}{Marge gauche (extérieure, f. 163v, en face du v. 93)}{ἀεὶ γὰρ τὸν Εὐριπίδην ὡς πανοῦργον διαβάλλει.}{ἀεὶ γ(ὰρ) τ(ὸν) Εὐριπίδ(ην) ὡς πανοῦργ\expo{ον} διαβάλλ(ει) {\policeFr ÷}}{« Car il raille toujours Euripide comme un homme fertile en ruses (πανοῦργος). »}
\commScholie{Justifie l'expression ἐκ τοῦ σοῦ τρόπου (« tout à fait dans ta manière »).}
\entreeScholie{042}{94}{τοῦ γὰρ τεχνάζειν ἡμέτερος ὁ πυραμοῦς.}{Marge gauche (extérieure, f. 163v, bloc de 5 lignes en face du v. 94)}{εἶδος πλακοῦντος ὁ πυραμοῦς. οὗτος δὲ ἐτίθετο τοῖς διαπαννυχίζουσιν ἔπαθλον, καὶ ὁ ἐγρηγορὼς ἕως πρωΐας ὡς νικῶν αὐτὸν ἐλάμβανε. νῦν οὖν ταύτῃ τῇ μεταφορᾷ ἐχρήσατο ὡς νικῶντος αὐτοῦ τῇ πανουργίᾳ πάντας, ἤτοι τὰς γυναῖκας.}{εἶδ\expo{ος} πλακοῦντ\expo{ος} ὁ πυραμ\expo{οῦς} · οὗτ\expo{ος} δ(ὲ) ἐτίθετο τ\expo{οῖς} διαπαννυχίζουσ(ιν) ἔπαθλ(ον), \textbf{ς} ὁ ἐγρηγορὼς ἕως πρωΐας ὡς νικ\expo{ῶν} αὐτ(ὸν) ἐλάμβανε · νῦν οὖν ταύτῃ τῇ μεταφορᾷ ἐχρήσατο ὡς νικῶντ\expo{ος} αὐτ\expo{οῦ} τῇ πανουργίᾳ πάντ\expo{ας}, ἤτοι τ\expo{ὰς} γυναῖκ\expo{ας} {\policeFr ÷}}{« Le πυραμοῦς est une espèce de gâteau (de froment grillé et de miel). On le proposait comme prix dans les veillées nocturnes (παννυχίδες), et celui qui restait éveillé jusqu'à l'aube le recevait en tant que vainqueur. Ici donc, il a usé de cette métaphore comme si Euripide l'emportait en ruse sur tout le monde, ou plutôt sur les femmes. »}
\commScholie{Repris par la Souda π 3172 et commenté par Louis Hauseux (vv 1-641.tex, l. 1318) avec le parallèle des Cavaliers v. 277.}
\entreeScholie{043}{96}{καὶ ποῖός ἐστιν; — οὗτος οὑκκυκλούμενος·}{Interlinéaire / marge intérieure droite (f. 163v, v. 96)}{ἐπὶ ἐκκυκλήματος γὰρ φαίνεται.}{ἐπὶ ἐκκυκλήματ\expo{ος} γ(ὰρ) φαίνετ(αι) {\policeFr ÷}}{« Car il (Agathon) apparaît amené au dehors sur la machine roulante (l'ekkyklêma). »}
\commScholie{Note scénique précieuse confirmant le texte de R ἀλλ' ἐκκυκλούμενος φαίνεται (défendu ou discuté à propos de la sortie d'Agathon sur l'ekkyklêma).}
\entreeScholie{044}{98}{Κυρήνην ὁρῶ.}{Marge gauche (extérieure, f. 163v, en face du v. 98)}{Κυρήνην: Τὴν ἑταίραν, περὶ ἧς πολλάκις φησί. τοιοῦτον δὲ τὸν Ἀγάθωνα καὶ οἱ ἄλλοι κωμῳδοῦσιν.}{Κυρήν(ην) : τ(ὴν) ἑταίρ\expo{αν}, π(ερὶ) ἧς πολλάκ(ις) φησί · τοιοῦτ\expo{ον} δ(ὲ) τ(ὸν) Ἀγάθωνα \textbf{ς} οἱ ἄλλ(οι) κωμῳδοῦσ(ιν) {\policeFr ÷}}{« Κυρήνην : la célèbre courtisane (Cyrène), dont Aristophane parle souvent (Grenouilles v. 1328). Et les autres poètes comiques raillent eux aussi Agathon sous ce rapport. »}
\commScholie{Rapproché par Louis Hauseux (vv 1-641.tex, ll. 1346-1365) de la scholie aux Grenouilles v. 1328 sur la courtisane Cyrène « aux douze postures » (δωδεκαμήχανος).}
\entreeScholie{045}{100}{Μύρμηκος ἀτραπούς, ἢ τί διαμινύρεται;}{Marge gauche (extérieure, f. 163v, en face du v. 100, Planche 19)}{ὡς λεπτὰ καὶ ἀγκύλα ἀνακρουομένου μέλη τοῦ Ἀγάθωνος· τοιαῦται γὰρ αἱ τῶν μυρμήκων ὁδοί. μεταξὺ δὲ τῶν δυοῖν ἀξιοῦσί τινες γράφειν μινυρισμὸς, ὡς πολλὰ τοιαῦτα παρεπιγράφεται.}{ὡς λεπτὰ \textbf{ς} ἀγκύλα ἀνακρουομέν\expo{ου} μέλ(η) τ\expo{οῦ} Ἀγάθων\expo{ος} · τοιαῦτ(αι) γ(ὰρ) αἱ τ\expo{ῶν} μυρμήκ\expo{ων} ὁδοί · μεταξὺ δ(ὲ) τ\expo{ῶν} δυοῖν ἀξιοῦσί τιν\expo{ες} γράφ(ειν) μινυρισμὸς, ὡς πολλὰ τοιαῦτ(α) παρεπιγράφεται {\policeFr ÷}}{« Parce qu'Agathon prélude par des mélodies menues et tortueuses : tels sont en effet les sentiers des fourmis. Et entre les deux vers (vv. 99 et 100), certains critiques veulent que l'on écrive μινυρισμός (“fredonnement”), comme bien d'autres indications scéniques de ce genre sont notées en parépigraphè (παρεπιγράφεται). »}
\commScholie{Scholie capitale pour l'histoire du texte dramatique antique (commentée par L. Hauseux, vv 1-641.tex, ll. 1367-1385, et illustrée sur la Planche 19 du Manuel) : elle atteste l'existence de la première parepigraphê scénique (μινυρισμός) entre les vv. 99 et 100.}
\entreeScholie{046}{101}{ἱερὰν Χθονίαιν [Χθονίαις R] δε- / ξάμεναι λαμπάδα κοῦραι / ξὺν ἐλευθέρᾳ πατρί- / δι χορεύσασθε βοάν·}{Marge gauche (extérieure, f. 163v, bloc de 7 lignes en face des vv. 101–104)}{ὁ Ἀγάθων ὑποκριτικὰ μέλη τέως ποιεῖ. ἀμφότερα δὲ αὐτὸς ὑποκρίνεται. μονῳδεῖ ὁ Ἀγάθων ὡς πρὸς χορὸν, οὐχ ὡς ἐπὶ σκηνῆς, ἀλλ’ ὡς ποιήματα συντιθείς. διὸ καὶ χορικὰ λέγει μέλη αὐτὸς πρὸς αὑτὸν, ὡς χορικὰ δέ. χθονίαις δὲ τῇ Δήμητρι καὶ τῇ Κόρῃ.}{ὁ Ἀγάθων ὑποκριτικὰ μέλη τέως ποι(εῖ) · ἀμφότερα δ(ὲ) αὐτὸς ὑποκρίνετ(αι) · μονῳδ(εῖ) ὁ Ἀγάθων ὡς πρ(ὸς) χορ(ὸν), οὐχ ὡς ἐπὶ σκην\expo{ῆς}, ἀλλ’ ὡς ποιήματα συντιθείς · διὸ \textbf{ς} χορικὰ λέγει μέλη αὐτὸς πρ(ὸς) αὑτ(ὸν), ὡς χορικὰ δέ · χθ(ον)ίαις δ(ὲ) τῇ Δήμητρι \textbf{ς} τῇ Κόρῃ {\policeFr ÷}}{« Agathon compose pour le moment des chants dialogués d'acteur (ὑποκριτικὰ μέλη), et il interprète lui-même les deux rôles (le coryphée et le chœur de jeunes filles troyennes). — Autre rédaction : Agathon chante une monodie comme s'il s'adressait à un chœur, non pas comme s'il était en représentation sur scène, mais en poète en train de composer ses chants ; c'est pourquoi il prononce des chants choraux en se répondant à lui-même, sur le mode choral. — Quant à Χθονίαις (“aux Déesses Chthoniennes”), il désigne Déméter et Coré. »}
\commScholie{Double scholie dramaturgique essentielle : elle explique pourquoi le manuscrit R alterne les rubriques d'Agathon dans cette monodie dialoguée où Agathon joue successivement la prêtresse-coryphée et le chœur des jeunes filles.}
\entreeScholie{047}{105}{τίνι δαιμόνων ὁ κῶμος; / λέγε νιν [νυν R]· εὐπείστως [εὐπίστως R, Σ] δὲ τοὐμὸν / δαίμονας ἔχει σεβίσαι·}{Marge gauche (extérieure, f. 163v, en face du v. 105)}{εὐπίστως ἔχω πρὸς τὸ τοὺς θεοὺς ὑμνῆσαι καὶ εὐλογῆσαι.}{εὐπίστ\expo{ως} ἔχω πρ(ὸς) τ(ὸ) τ\expo{οὺς} θε\expo{οὺς} ὑμνῆσ(αι) \textbf{ς} εὐλογῆσ(αι) {\policeFr ÷}}{« Je suis tout disposé (plein de foi et de docilité) à chanter et célébrer les dieux. »}
\commScholie{Paraphrase de la périphrase lyrique εὐπίστως δὲ τοὐμὸν δαίμονας ἔχει σεβίσαι.}
\entreeScholie{048}{109}{Φοῖβον, ὃς ἱδρύσατο χώρας / γύαλα Σιμουντίδι γᾷ·}{Interlinéaire (supra lineam au-dessus de v. 109, f. 163v)}{ὃς τὴν Ἴλιον ἐτείχισεν.}{ὃς τ(ὴν) Ἴλι\expo{ον} ἐτείχισ(εν) {\policeFr ÷}}{« Celui (Apollon) qui a bâti les remparts d'Ilion. »}
\commScholie{Explique pourquoi le chœur supposé de jeunes filles troyennes invoque Apollon comme fondateur des murailles de Troie (Iliade VII, 452 ; XXI, 446).}
\entreeScholie{049}{117}{γόνον ὀλβίζουσα Λατοῦς / Ἄρτεμιν ἀπειρολεχῆ·}{Interlinéaire (supra lineam au-dessus du v. 117, f. 163v)}{τὴν παρθένον.}{τ(ὴν) παρθέν\expo{ον} {\policeFr ÷}}{« La Vierge (Artémis). »}
\commScholie{Glose interlinéaire de R\expo{2} expliquant l'épithète composée ἀπειρολεχῆ (« qui ignore la couche nuptiale », c'est-à-dire vierge, τὴν παρθένον) appliquée à Artémis.}
\entreeScholie{050}{121}{κρούματά τ᾽ Ἀσιάδος}{Marge inférieure et marge gauche (bas du f. 163v, en face du v. 121)}{τῇ Φρυγίᾳ ἁρμονίᾳ ἡρμοσμένα. καὶ τὰ κρούματα τοῦ Ἀπόλλωνος κιθάρας σέβομαι, ἐπειδὴ εἶπε κάτω «κιθάραν ματέρα ὕμνων.» Ἀσιάδα δὲ τὴν κιθάραν λέγει.}{τῇ Φρυγίᾳ ἁρμονίᾳ ἡρμοσμένα · \textbf{ς} τὰ κρούματ(α) τ\expo{οῦ} Ἀπόλλων\expo{ος} κιθάρ\expo{ας} σέβομ(αι), ἐπειδὴ εἶπ(ε) κάτω «κιθάρ(αν) ματέρα ὕμν\expo{ων}» · Ἀσιάδα δ(ὲ) τ(ὴν) κιθάρ(αν) λέγ(ει) {\policeFr ÷}}{« Accordés selon l'harmonie phrygienne. Et je vénère les accords frappés sur la cithare d'Apollon, puisqu'il dit plus bas (v. 124) “la cithare mère des hymnes”. Il appelle Ἀσιάς (“l'Asiatique”) la cithare à trois cordes (ou luth lydien). »}
\commScholie{Repris par la Souda α 4149 (s.v. Ἀσιάδος κρούματα) et Hésychius α 7650 : l'Ἀσιάς désigne la grande cithare d'Asie Mineure.}
\entreeScholie{051}{122}{παράρυθμ᾽ εὔρυθμα Φρυγίων / διὰ νεύματα [δινεύματα Σ] Χαρίτων·}{Interlinéaire (supra lineam sur δινεύματα, dernière ligne du f. 163v)}{ὀρχήματα.}{ὀρχήματα {\policeFr ÷}}{« Les danses tournoyantes (les évolutions chorégraphiques). »}
\commScholie{Glose le terme poétique δινεύματα par le mot usuel ὀρχήματα.}
\clearpage
\renewcommand{\currHeadLE}{ΑΡΙΣΤΟΦΑΝΟΥΣ ΘΕΣΜΟΦΟΡΙΑΖΟΥΣΑΙ · CODEX RAVENNAS 429 (\R) · FOLIO 04/30}
\renewcommand{\currHeadRE}{f. 164r (Classense 166r) · vv. 123–164 · Collation diplomatique}
\renewcommand{\currHeadLO}{ΑΡΙΣΤΟΦΑΝΟΥΣ ΘΕΣΜΟΦΟΡΙΑΖΟΥΣΑΙ · CODEX RAVENNAS 429 (\R) · FOLIO 04/30}
\renewcommand{\currHeadRO}{f. 164r (Classense 166r) · vv. 123–164 · Collation diplomatique}
\noindent\begin{adjustbox}{max totalheight=30.4cm, max width=\textwidth, keepaspectratio}\begin{minipage}{\textwidth}
\bandeauRav{f. 164r}{Classense 166r}{123–164}{Folio 04/30 · Recto (page de droite) · 22 scholies \& gloses ({\policeGr Σ} 052–073)}
{\fontsize{8.8pt}{10.2pt}\selectfont
\ligneRav{123}{ἀγάθ/}{\hspace{1.1em}σέβομαι Λατώ τ' ἄνασσαν}{}
\ligneRav{}{}{\hspace{1.1em}κίθαρίν τε ματέρ᾽ ὕμνων}{}
\ligneRav{125}{}{\hspace{1.1em}ἄρσενι βοᾷ δοκίμῳ·}{\textbf{125} \lecon{δοκιμῳ} R, Σ\expo{1}, Suid. α 4016 : δόκιμον Σ\expo{2} : δοκίμων Schöne}
\ligneRav{126}{}{\hspace{1.1em}τᾷ \difference{φάος} ἔσσυτο δαιμονίοις}{\textbf{126} φῶς R : \lecon{φάος} Burges}
\ligneRav{127}{}{\hspace{1.1em}ὄμμασιν \difference{ἁμετέρας} τε δι᾽ αἰφνι-}{\textbf{127} ἡμετέρας R : ὑμετέρας Nietzsche : \lecon{ἁμετέρας} Meineke}
\ligneRav{}{}{\hspace{1.1em}δίου ὀπός. ὧν χάριν}{}
\ligneRav{128}{}{\hspace{1.1em}ἄνακτ᾽ ἄγαλλε \difference{Φοῖβον}·}{\textbf{128} \lecon{ἄγαλλε} R : ἀγάλλετε Austin (1987) \espApp Φοῖβον τιμᾷ R : \lecon{Φοῖβον} Dindorf}
\ligneRav{}{}{\hspace{1.1em}χαῖρ᾽ ὄλβιε παῖ Λατοῦς·}{}
\ligneRav{}{}{\hspace{2.2em}\textbf{[παρεπιγραφή R\expo{1} :]} \lecon{ὀλολύζει ὁ γέρων}}{\textbf{post 129} \lecon{ὀλολύζει ὁ γέρων} (παρεπιγραφή) R}
\ligneRav{130}{μ\auDessus{ν}{Η}}{ὡς ἡδὺ τὸ μέλος ὦ \difference{πότνιαι} Γενετυλλίδες}{\textbf{130} πότνια R : \lecon{πότνιαι} Suid. γ 141}
\ligneRav{}{}{καὶ θηλυδριῶδες καὶ κατεγλωττισμένον}{}
\ligneRav{}{}{καὶ μανδαλωτόν, ὥστ᾽ ἐμοῦ γ᾽ ἀκροωμένου}{}
\ligneRav{}{}{ὑπὸ τὴν ἕδραν αὐτὴν ὑπῆλθε γάργαλος·}{}
\ligneRav{134}{}{καί σ᾽ ὦ \difference{νεανίσκ᾽ εἴ τις} εἶ, κατ᾽ Αἰσχύλον}{\textbf{134} νεανίσχ᾽ εἴ τις R : \lecon{νεανίσκ᾽ εἴ τις} \Mon : νεανίσχ᾽ ἥτις Gelenius}
\ligneRav{135}{}{ἐκ τῆς \difference{Λυκουργείας} ἐρέσθαι βούλομαι·}{\textbf{135} Λυκουργίας R : \lecon{Λυκουργείας} Dindorf}
\ligneRav{}{}{ποδαπὸς ὁ γύννις; τίς πάτρα; τίς ἡ στολή;}{}
\ligneRav{}{}{τίς ἡ τάραξις τοῦ βίου; τί βάρβιτος}{}
\ligneRav{}{}{λαλεῖ κροκωτῷ; τί δὲ λύρα κεκρυφάλῳ;}{}
\ligneRav{139}{}{τί λήκυθος καὶ στρόφιον; ὡς οὐ ξύμφορον·}{\textbf{139} \lecon{οὐ ξύμ-} R, Suid. β 190 V : ἀξυμ- Suid. M : ὀξύμ- Suid. A,G,I : ξύμ- Suid. F}
\ligneRav{140}{}{τίς δαὶ \difference{κατρόπτου} καὶ ξίφους κοινωνία;}{\textbf{140} κατόπτρου R, Suid. β 110 : \lecon{κατρόπτου} Dover ad \italique{Ran.} 47}
\ligneRav{141}{}{\difference{σύ τ᾽} αὐτὸς ὦ παῖ; πότερον ὡς ἀνὴρ τρέφει;}{\textbf{141} \lecon{σύ τ᾽} \PsiUn, Σ : τίς δ᾽ R}
\ligneRav{142}{}{καὶ ποῦ \difference{πέος}; ποῦ χλαῖνα; ποῦ Λακωνικαί;}{\textbf{142} \lecon{πέος} \PsiUn, R : τὸ πέος Suid. λ 63, π 987}
\ligneRav{}{}{ἀλλ᾽ ὡς γυνὴ δῆτ᾽ εἶτα ποῦ τὰ τιτθία;}{}
\ligneRav{}{}{τί φῄς; τί σιγᾷς; ἀλλὰ δῆτ᾽ ἐκ τοῦ μέλους}{}
\ligneRav{145}{}{ζητῶ σ᾽, ἐπειδή γ᾽ αὐτὸς οὐ βούλει φράσαι;}{}
\ligneRav{}{ἀγάθ/}{ὦ πρέσβυ πρέσβυ, τοῦ φθόνου μὲν τὸν ψόγον}{}
\ligneRav{}{}{ἤκουσα, τὴν δ᾽ ἄλγησιν οὐ παρεσχόμην}{}
\ligneRav{148}{}{ἐγὼ δὲ τὴν ἐσθῆθ᾽ ἅμα \difference{τῇ γνώμῃ} φορῶ·}{\textbf{148} γνώμη R : \lecon{τῇ γνώμῃ} Meineke}
\ligneRav{}{}{χρὴ γὰρ ποιητὴν ἄνδρα πρὸς τὰ δράματα}{}
\ligneRav{150}{}{ἃ δεῖ ποιεῖν πρὸς ταῦτα τοὺς τρόπους ἔχειν·}{}
\ligneRav{}{}{αὐτίκα γυναικεῖ᾽ ἢν ποιῇ τις δράματα,}{}
\ligneRav{}{}{μετουσίαν δεῖ τῶν τρόπων τὸ σῶμ᾽ ἔχειν·}{}
\ligneRav{}{μ\auDessus{ν}{Η}}{οὐκοῦν κελητίζεις, ὅταν Φαίδραν ποιῇς;}{}
\ligneRav{154}{ἀγάθ/}{ἀνδρεῖα δ᾽ \difference{ἄν} ποιῇ τις, ἐν τῷ σώματι}{\textbf{154} \lecon{ἄν} \PsiUn, R : ἢν Dindorf}
\ligneRav{155}{}{ἔνεσθ᾽ ὑπάρχον τοῦθ᾽. ἃ δ᾽ οὐ κεκτήμεθα,}{}
\ligneRav{}{}{μίμησις ἤδη ταῦτα συνθηρεύεται·}{}
\ligneRav{}{μ\auDessus{ν}{Η}}{ὅταν σατύρους τοίνυν ποιῇς, καλεῖν ἐμέ,}{}
\ligneRav{}{}{ἵνα συμποιῶ σοὔπισθεν ἐστυκὼς ἐγώ·}{}
\ligneRav{}{ἀγάθ/}{ἄλλως τ᾽ ἄμουσόν ἐστι ποιητὴν ἰδεῖν}{}
\ligneRav{160}{}{ἀγρεῖον ὄντα καὶ δασύν σκέψαι δ᾽ ὅτι}{}
\ligneRav{}{}{Ἴβυκος ἐκεῖνος κἀνακρέων ὁ Τήιος}{}
\ligneRav{162}{}{κἀλκαῖος, οἳ περὶ ἁρμονίαν \difference{ἐκεχήνεσαν},}{\textbf{162} \lecon{κἀλκαῖος} R, Suid. : (κ)ἀχαιός Σ \espApp ἐχύμησαν R : ἐχύμισαν Suid. : \lecon{ἐκεχήνεσαν} Bibi}
\ligneRav{163}{}{ἐμιτροφόρουν τε καὶ \difference{διέκειντ'} Ἰωνικῶς,}{\textbf{163} διεκίνων R : διεκίνουν Suid. ε 989, ι 495 : διεκλῶντ᾽ Toup : \lecon{διέκειντ'} Bibi}
\ligneRav{}{}{καὶ Φρύνιχος, τοῦτον γὰρ οὖν ἀκήκοας,}{}
}
\vspace{3pt}{\color{ruleGold}\hrule height 0.6pt}\vspace{3pt}
{\fontsize{8.0pt}{9.4pt}\selectfont
\noindent{\bfseries\scshape\color{rubricRed}Loci paralleli et testimonia (f. 164r, vv. 123–164) :} \textbf{125} (ἄρσενι βοᾷ δοκιμῳ) Suidas α 4016 = Scholium \espTest \textbf{128} (ἄγαλλε) Photius α 164 = \SynB α 145 \espTest \textbf{130-2} (ἡδὺ -- μανδαλωτόν) Suidas γ 141 \espTest \textbf{130} (Γενετυλλίδες) Hesychius γ 343 Γενετυλλίς \espTest \textbf{131-2} (θηλυδριῶδες -- μανδαλωτόν) Suidas μ 134 \espTest \textbf{131} (θηλυδριῶδες καὶ κατεγλωττισμένον) Suidas κ 912 \espTest Hesychius θ 491 θηλυδριῶδες \espTest (κατεγλωττισμένον) Pollux II.109 \espTest \textbf{132} Hesychius μ 224 μανδαλωτόν {\policeFr ≈} Photius μ 86, cf. Photius γ 115 \espTest \textbf{132-3} (ὥστ᾽ -- γάργαλος) Suidas γ 83 (sed τοῦ μέλους post ἀκροωμένου addito) \espTest \textbf{133} (γάργαλος = γαργαλισμός) Moeris γ 23, cf. \Phryn 56.9 \espTest \textbf{134} Hesychius η 948 ἥτις· ὅστις (?) \espTest \textbf{136} (ποδαπὸς -- πάτρα) Suidas γ 504 \espTest (γύννις) Hesychius γ 1015 γύννις, cf. \Phryn 60.19 \espTest \textbf{136-40} (τίς ἡ στολή -- κοινωνία) Suidas β 110 (δὲ pro δαὶ) \espTest \textbf{137} (βάρβιτος) Hesychius β 211 βάρβιτος, cf. Photius β 62 \espTest \textbf{141-2} (πότερον -- πέος) Suidas π 987 \espTest \textbf{142} (ποῦ τὸ πέος -- Λακωνικαί) Suidas λ 63, π 987 (τὸ πέος pro πέος) \espTest Moeris χ 19 χλαῖνα \espTest (Λακωνικαί) Hesychius λ 225 Λακωνικαί ; Photius λ 57 Λακωνικαί ; scholium \italique{Vespae} 1158-60 \espTest \textbf{142} (ποῦ τὸ πέος -- Λακωνικαί) Suidas λ 63 \espTest \textbf{146-7} Suidas α 1097 (μὲν omisso et δὲ pro δ') \espTest \textbf{149-50} \Apostol XVIII 41h \espTest \textbf{159-61} Suidas α 1633 (κ' Ἀνακρ.) \espTest \textbf{161-2} Suidas ε 4049 (Ἀνακρέων, Ἀλκαῖος, οἵπερ ἁρμονίαν ἐχύμισαν) \espTest \textbf{162-3} (κ' Ἀλκαῖος -- Ἰωνικῶς) Suidas ε 989 (sed οἵπερ ἁρμονίαν ἐχύμισαν) {\policeFr ≈} Zonaras 707 ἐμιτροφόρουν \espTest \textbf{162} Zonaras 942 οἵπερ ἁρμονίαν ἐχύμισαν \espTest (ἐχύμισαν) Hesychius ε 7681 ἐχύμισαν· ἤρτυσαν \espTest \textbf{162-3} (κ' Ἀλκαῖος -- Ἰωνικῶς) Suidas ε 989 {\policeFr ≈} Zonaras 707 ἐμιτροφόρουν \espTest (ἐχύμισαν) Hesychius ε 7681 ἐχύμισαν \espTest \textbf{163} Suidas ι 495 ἐμιτροφόρουν τε καὶ διεκίνουν Ἰωνικῶς\par\vspace{2.5pt}
{\color{ruleGold}\hrule height 0.35pt}\vspace{2.5pt}
\noindent{\bfseries\scshape\color{rubricRed}Suite (en commentaire) de l'apparat critique et paléographique :} \textbf{Témoin papyrologique :} \textbf{\PsiUn{} (\italique{PSI} XI 1194 + XIV p. XV)} couvre les \textbf{vv. 133–156}. --- \textbf{Deuxième \italique{parepigraphè} (entre vv. 129 et 130) :} \Run{} copie sur une ligne propre \lecon{ὀλολύζει ὁ γέρων}. --- \textbf{Confrontation avec \PsiUn{} (v. 141) :} Au v. 141, R porte la leçon fautive τίς δ' αὐτὸς ὦ παῖ· πότερον ὡς ἀνὴρ τρέφει;, tandis que le papyrus \PsiUn{} et la scholie {\policeGr Σ}\expo{R} conservent l'authentique \lecon{σύ τ'}. --- \textbf{Double émendation paléographique de \textbf{Bibi} aux vv. 162–163 :} ἐχύμησαν ... διεκίνων R {\policeFr →} \lecon{ἐκεχήνεσαν ... διέκειντ'} \textbf{Bibi}. --- \textbf{v. 125 :} \lecon{δοκίμῳ} R, {\policeGr Σ}\expo{1}, Souda α 4016 : δόκιμον {\policeGr Σ}\expo{2} : δοκίμων Schöne. --- \textbf{v. 127 :} ἡμετέρας R (ὑμετέρας Nietzsche) : \lecon{ἁμετέρας} Meineke (dorisme lyrique). --- \textbf{v. 128 :} \lecon{ἄγαλλε} R (ἀγάλλετε Austin 1987) \espApp Φοῖβον τιμᾷ R (τιμᾷ est une glose intruse) : \lecon{Φοῖβον} Dindorf. --- \textbf{vv. 134–135 :} νεανίσχ' εἴ τις R (aspiration fautive par anticipation de l'esprit rude) : \lecon{νεανίσκ' εἴ τις} \Mon{} (νεανίσχ' ἥτις Gelenius) \espApp \textbf{v. 135 :} Λυκουργίας R (iotacisme) : \lecon{Λυκουργείας} Dindorf. --- \textbf{v. 141 :} \lecon{σύ τ'} {\policeGr Ψ}\expo{1}, {\policeGr Σ} : τίς δ' R. --- \textbf{v. 162 :} \lecon{κἀλκαῖος} R, Souda ((κ)ἀχαιός {\policeGr Σ}) \espApp ἐχύμησαν R (ἐχύμισαν Souda) : \lecon{ἐκεχήνεσαν} \lecon{Bibi}. \textbf{Explication paléographique (Commentaire ad v. 162) :} Dans ἐκεχήνεσαν (plus-que-parfait de χάσκω, « ils restaient bouche bée », cf. \italique{Ach.} 10), le saut visuel εκεχ- {\policeFr →} εχ- (haplographie), joint à l'iotacisme η → υ et à la confusion minuscule constante ν → μ, a produit le monstre ἐχύμησαν de R ! --- \textbf{v. 163 :} διεκίνων R (διεκίνουν Souda ε 989, ι 495 ; διεκλῶντ' Toup) : \lecon{διέκειντ'} \lecon{Bibi} (ou δίκανον ap. Hésychius δ 1805).\par}
\end{minipage}\end{adjustbox}
\newpage
\renewcommand{\currHeadLE}{SCHOLIES MARGINALES ({\policeGr Σ}\expo{R}) \& GLOSES INTERLINÉAIRES · FOLIO 04/30}
\renewcommand{\currHeadRE}{f. 164r (Classense 166r) · vv. 123–164 · {\policeGr Σ} 052 à {\policeGr Σ} 073}
\renewcommand{\currHeadLO}{SCHOLIES MARGINALES ({\policeGr Σ}\expo{R}) \& GLOSES INTERLINÉAIRES · FOLIO 04/30}
\renewcommand{\currHeadRO}{f. 164r (Classense 166r) · vv. 123–164 · {\policeGr Σ} 052 à {\policeGr Σ} 073}
\bandeauScholies{Folio 04/30 — f. 164r}{Classense 166r}{vv. 123–164}{Recto (page de droite) · 22 entrées ({\policeGr Σ} 052 à {\policeGr Σ} 073)}
{\fontsize{8.8pt}{10.5pt}\selectfont\noindent\textbf{\color{rubricRed}Disposition topographique des marges sur le f. 164r :} Marge extérieure à droite ; marge intérieure à gauche. — \textit{Fac-similé de référence :} \texttt{folios\_individuels\_ff162v-177r/04\_f164r\_classense166r\_vv0123-0164.jpg}.\par\vspace{3pt}}
\entreeScholie{052}{123}{σέβομαι Λατώ τ᾽ ἄνασσαν}{Interlinéaire (supra lineam en tête du f. 164r, v. 123)}{τὸ σέβομαι κοινόν.}{τ(ὸ) σέβομ(αι) κοινόν {\policeFr ÷}}{« Le verbe σέβομαι (“je vénère”) est en facteur commun (pour Létô et pour la cithare). »}
\commScholie{Précise la syntaxe de la strophe finale du chant d'Agathon.}
\entreeScholie{053}{125}{ἄρσενι βοᾷ δοκίμῳ,}{Marge supérieure et marge droite (extérieure, f. 164r, en face du v. 125)}{εὐτόνῳ. προεῖπε γὰρ ὅτι ἐκ παρθένων ἐστὶν ὁ χορός. ἢ ἄρσενι βοᾷ δοκίμῳ, ἐπεὶ Ἀπόλλων ἐστὶν ὁ κιθαρίζων. τῆς οὖν κιθάρας ἄρσενι βοᾷ δοκίμου οὔσης φῶς ἔσσυτο. ἐπάγει δὲ σημεῖον τοῦ παρὰ Ἀπόλλωνος λελέχθαι. τῇ θαυμαστῇ φωνῇ καὶ δεδοκιμασμένῃ.}{εὐτόνῳ · προεῖπε γ(ὰρ) ὅτι ἐκ παρθένων ἐστ(ὶν) ὁ χορός · ἢ ἄρσενι βοᾷ δοκίμῳ, ἐπεὶ Ἀπόλλ\expo{ων} ἐστ(ὶν) ὁ κιθαρίζων · τ\expo{ῆς} οὖν κιθάρας ἄρσενι βοᾷ δοκίμου οὔσ\expo{ης} φῶς ἔσσυτο · ἐπάγει δ(ὲ) σημεῖ(ον) τ\expo{οῦ} παρὰ Ἀπόλλων\expo{ος} λελέχθ(αι) · τῇ θαυμαστῇ φωνῇ \textbf{ς} δεδοκιμασμένῃ {\policeFr ÷}}{« Vigoureuse, bien tendue (εὐτόνῳ), car il a dit plus haut que le chœur est composé de jeunes filles (ou par contraste avec la voix féminine). Ou bien : “par une clameur mâle éprouvée” (ἄρσενι βοᾷ δοκίμῳ), parce que c'est Apollon qui joue de la cithare ; la cithare étant donc éprouvée par une voix mâle, une lumière s'est élancée (φῶς ἔσσυτο). Il ajoute ensuite la preuve que cela est dit d'Apollon. — Autre glose : Par la voix admirable et éprouvée (δοκίμῳ). »}
\commScholie{Comme le note Louis Hauseux (Bibi, v. 125, vv 1-641.tex, ll. 1521-1522), la première rédaction de cette scholie (Σ\expo{1}) ainsi que la Souda α 4016 confirment le datif δοκίμῳ de R (retenu par Bibi) contre la variante δόκιμον de Σ\expo{2}.}
\entreeScholie{054}{126}{τᾷ φάος ἔσσυτο δαιμονίοις / ὄμμασιν ἁμετέρας [ἡμετέρας R] τε δι᾽ αἰφνιδίου ὀπός.}{Marge droite (extérieure, f. 164r, bloc de 6 lignes sous Σ 125)}{τῆς κιθάρας φῶς ἔσσυτο τοῖς σοφοῖς· ἀπὸ γὰρ τῆς κιθάρας ἐδιδάχθησαν οἱ παλαιοὶ ᾄδοντες καὶ γράφοντες ποιήματα. ἢ τῇ Λητοῖ λαμπρότατα ὄμματα λέγει. πάλιν δὲ φῶς ἐγένετο αὐτοῖς ἀπὸ τῆς ἡμετέρας φωνῆς· ἄνευ γὰρ κιθάρας ἢ ἄνευ ἀνδρὸς τοῦ διδάσκοντος οὐ μαθήσεταί τις.}{τ\expo{ῆς} κιθάρας φῶς ἔσσυτο τ\expo{οῖς} σοφ\expo{οῖς}· ἀπ(ὸ) γ(ὰρ) τ\expo{ῆς} κιθάρας ἐδιδάχθησ(αν) οἱ παλαιοὶ ᾄδοντες \textbf{ς} γράφοντες ποιήματα · ἢ τῇ Λητοῖ λαμπρότατα ὄμματα λέγει · πάλιν δ(ὲ) φῶς ἐγένετο αὐτ\expo{οῖς} ἀπ(ὸ) τ\expo{ῆς} ἡμετέρας φων\expo{ῆς}· ἄνευ γ(ὰρ) κιθάρας ἢ ἄνευ ἀνδρὸς τ\expo{οῦ} διδάσκοντ\expo{ος} οὐ μαθήσεταί τις {\policeFr ÷}}{« De la cithare une lumière a jailli pour les sages : car c'est à partir de la cithare que les anciens ont appris à chanter et à composer des poèmes. Ou bien il parle des yeux très brillants de Létô (ou des dieux). Et en retour, une lumière leur est venue de notre propre voix : car sans cithare ou sans maître qui enseigne, nul ne saurait apprendre. »}
\commScholie{Glose exégétique essayant de démêler la syntaxe volontairement alambiquée et dithyrambique du chant d'Agathon (sur R φῶς, corrigé métriquement en φάος par Burges et Bibi).}
\entreeScholie{055}{127}{ἁμετέρας [ἡμετέρας R] τε δι᾽ αἰφνιδίου ὀπός.}{Interlinéaire (supra lineam au-dessus d'αἰφνιδίου, f. 164r)}{ἀντὶ τοῦ ἐνθουσιαστικῆς.}{ἀν\auDessus{τ}{ι} τ\expo{οῦ} ἐνθουσιαστικ\expo{ῆς} {\policeFr ÷}}{« (αἰφνιδίου, “soudaine”) : au sens d'inspirée par un transport divin (ἐνθουσιαστικῆς). »}
\commScholie{Glose fine saisissant la valeur religieuse de l'adjectif αἰφνίδιος appliqué au jaillissement soudain du chant rituel.}
\entreeScholie{056}{129}{[παρεπιγραφή :] ὀλολύζει ὁ γέρων}{Intra-textuel (dans la colonne de R\expo{1} après le v. 129) et Marge droite (extérieure, Σ 129 de R\expo{2}, Planche 19)}{ὀλολύζει. παρεπιγραφὴ τοῦτο· εἰπόντος γὰρ τοῦ Ἀγάθωνος ὁ κηδεστὴς τοῦ Εὐριπίδου ὀλολύζει εἰς τὴν θηλύτητα αὐτοῦ καὶ μόνον οὐχὶ γυναικὸς εὐχήν.}{ὀλολύζ(ει) · παρεπιγραφὴ τοῦτ(ο) · εἰπόντ\expo{ος} γ(ὰρ) τ\expo{οῦ} Ἀγάθων\expo{ος} ὁ κηδεστ\expo{ὴς} τ\expo{οῦ} Εὐριπίδ\expo{ου} ὀλολύζ(ει) εἰς τ(ὴν) θηλύτητ(α) αὐτ\expo{οῦ} \textbf{ς} μόν\expo{ον} οὐχὶ γυναικ\expo{ὸς} εὐχήν {\policeFr ÷}}{« ὀλολύζει (“il pousse le cri rituel des femmes”) : ceci est une didascalie scénique (παρεπιγραφὴ τοῦτο). En effet, quand Agathon a fini de chanter, le parent par alliance d'Euripide pousse un ὀλολυγμός pour se moquer de son caractère efféminé et de son chant qui n'est presque rien d'autre qu'une prière de femme. »}
\commScholie{Deuxième des quatre parepigraphai scéniques antiques conservées par le Ravennas (étudiée dans le Manuel, Planche 19). R\expo{1} a copié dans la colonne ὀλολύζει γέρων, et R\expo{2} confirme expressément en marge : παρεπιγραφὴ τοῦτο.}
\entreeScholie{057}{130}{ὡς ἡδὺ τὸ μέλος ὦ πότνιαι Γενετυλλίδες}{Marge droite (extérieure, f. 164r, en face du v. 130)}{δαίμων ἡ Γενετυλλὶς περὶ τὴν Ἀφροδίτην. φασὶ δὲ παρὰ τὴν γέννησιν πεποιῆσθαι τὸ ὄνομα. τινὲς δὲ περὶ τὴν Ἄρτεμίν φασιν αὐτὰς τῶν τοκετῶν ἐφόρους, καὶ πάλιν παρὰ τὴν γέννησιν τὸ ὄνομα πεποιῆσθαι.}{δαίμων ἡ Γενετυλλὶς π(ερὶ) τ(ὴν) Ἀφροδίτην · φασὶ δ(ὲ) παρὰ τ(ὴν) γέννησ(ιν) πεποιῆσθ(αι) τ(ὸ) ὄνομα · τινὲς δ(ὲ) π(ερὶ) τ(ὴν) Ἄρτεμίν φασ(ιν) αὐτὰς τ\expo{ῶν} τοκετ\expo{ῶν} ἐφόρους, \textbf{ς} πάλιν παρὰ τ(ὴν) γέννησ(ιν) τ(ὸ) ὄνομα πεποιῆσθ(αι) {\policeFr ÷}}{« Génetyllis est une divinité du cortège d'Aphrodite ; on dit que son nom a été formé d'après la génération (γέννησις). Certains disent que les Génetyllides appartiennent au cortège d'Artémis et président aux accouchements, leur nom étant là encore tiré de la naissance. »}
\commScholie{Repris par la Souda γ 142 (cf. Nuées v. 52 et Lysistrata v. 2). Mnésiloque invoque ironiquement ces déesses spécifiquement féminines et lascives.}
\entreeScholie{058}{131}{καὶ θηλυδριῶδες καὶ κατεγλωττισμένον}{Marge droite (extérieure, f. 164r, sous Σ 130)}{ποικίλον καὶ ἡδύ. ἔστι δὲ εἶδος φιλήματος πολλαῖς γλώτταις μεμιγμένον. γλῶττα δέ ἐστι καὶ ἡ λέξις· κυρίως δὲ κατεγλωττισμένον ἐστὶ τὸ ἔγγλωττον φίλημα.}{ποικίλον \textbf{ς} ἡδύ · ἔστι δ(ὲ) εἶδ\expo{ος} φιλήματ\expo{ος} πολλ\expo{αῖς} γλώττ\expo{αις} μεμιγμένον · γλῶττα δέ ἐστ(ι) \textbf{ς} ἡ λέξις· κυρίως δ(ὲ) κατεγλωττισμένον ἐστ(ὶ) τ(ὸ) ἔγγλωττ(ον) φίλημα {\policeFr ÷}}{« Varié et voluptueux. C'est une sorte de baiser mêlé de nombreux jeux de langue ; mais γλῶττα signifie aussi le mot rare (la parole), et au sens propre κατεγλωττισμένον désigne le baiser avec la langue (τὸ ἔγγλωττον φίλημα, “baiser colombien”). »}
\commScholie{Repris par la Souda κ 915. Le scholiaste explique admirablement le double sens : musical/verbal (brodé de fioritures vocales) et érotique (baiser lascif avec la langue), qui prépare le v. 132 (ὥστ' ἐγὼγ' ἀκροώμενος / ὑπὸ τὴν ἕδραν αὐτὴν ὑπῆλθον γαργαλός).}
\entreeScholie{059}{135}{ἐκ τῆς Λυκουργείας ἐρέσθαι βούλομαι· / ποδαπὸς ὁ γύννις;}{Marge droite (extérieure, f. 164r, en face des vv. 135–136)}{τὴν τετραλογίαν λέγει Λυκουργίαν, Ἠδωνοὺς, Βασσαρίδας, Νεανίσκους, Λυκοῦργον τὸν σατυρικόν. λέγει δὲ ἐν τοῖς Ἠδωνοῖς πρὸς τὸν συλληφθέντα Διόνυσον «ποδαπὸς ὁ γύννις;»}{τ(ὴν) τετραλογί(αν) λέγ(ει) Λυκουργί(αν), Ἠδων\expo{οὺς}, Βασσαρίδ\expo{ας}, Νεανίσκ\expo{ους}, Λυκοῦργ\expo{ον} τ(ὸν) σατυρικόν · λέγ(ει) δ(ὲ) ἐν τ\expo{οῖς} Ἠδων\expo{οῖς} πρ(ὸς) τ(ὸν) συλληφθέντ(α) Διόνυσ\expo{ον} «ποδαπ\expo{ὸς} ὁ γύννις;» {\policeFr ÷}}{« Il désigne la tétralogie d'Eschyle intitulée la Lycurgie, comprenant : les Édoniens, les Bassarides, les Jeunes Gens, et le drame satyrique Lycurgue. Dans les Édoniens (fr. 61 Radt), Lycurgue dit à Dionysos qui vient d'être capturé : “De quel pays sort cet homme-femme (ποδαπὸς ὁ γύννις) ?” »}
\commScholie{Joyau d'érudition alexandrine (Aristarque / Didyme) cité par Louis Hauseux (vv 1-641.tex, l. 1566) : c'est grâce à cette unique scholie du Ravennas que nous connaissons la composition exacte des quatre pièces de la Lycurgie d'Eschyle !}
\entreeScholie{060}{137}{τίς ἡ τάραξις τοῦ βίου; τί βάρβιτος / λαλεῖ κροκωτῷ; τί δὲ λύρα κεκρυφάλῳ;}{Marge droite (extérieure, f. 164r, en face du v. 137) et glose interlinéaire sur βάρβιτος}{ἐντεῦθεν τὴν ἀρχὴν Εὔβουλος ἐποιήσατο τοῦ Διονυσίου, τὰ ἀνόμοια τῶν ἐν τῇ Διονυσίου οἰκίᾳ καταλέγων· ἐπὶ πλέον μέντοι. ἡ κιθάρα.}{ἐντεῦθ(εν) τ(ὴν) ἀρχ(ὴν) Εὔβουλ\expo{ος} ἐποιήσατ(ο) τ\expo{οῦ} Διονυσί\expo{ου}, τὰ ἀνόμοια τ\expo{ῶν} ἐν τῇ Διονυσί\expo{ου} οἰκίᾳ καταλέγ\expo{ων} · ἐπὶ πλέον μέντοι {\policeFr ÷} // ἡ κιθάρα {\policeFr ÷}}{« C'est de ce passage qu'Euboule (poète de la Comédie Moyenne, fr. 26 K.-A.) a tiré le début de son Denys (Διονύσιος), en énumérant les objets disparates qui se trouvaient dans la maison de Denys le Tyran, mais en développant davantage. — βάρβιτος : la cithare (lyre grave). »}
\commScholie{Autre notice d'histoire littéraire alexandrine de premier ordre, signalant l'imitation de cette tirade d'Aristophane par le poète comique Euboule dans sa pièce sur Denys de Syracuse.}
\entreeScholie{061}{139}{τί λήκυθος καὶ στρόφιον; ὡς οὐ ξύμφορον.}{Interlinéaire (supra lineam au-dessus de στρόφιον, f. 164r)}{ζωνάριον.}{ζωνάρι\expo{ον} {\policeFr ÷}}{« Une petite ceinture de poitrine (bandeau de gorge des femmes, soutien-gorge). »}
\commScholie{Glose opposant la fiole d'huile du gymnase masculin (λήκυθος) au bandeau de gorge féminin (στρόφιον).}
\entreeScholie{062}{141}{τίς δαὶ σύ; πότερον ὡς ἀνὴρ τρέφει;}{Marge intérieure gauche (variante γρ.) et interlinéaire sur τρέφει (f. 164r)}{γρ. σύ τ’ αὐτός. ἀντὶ τοῦ τρέφῃ· τὸ γὰρ τρέφει Ἀττικόν.}{γρ(άφεται) σύ τ' αὐτός {\policeFr ÷} // ἀν\auDessus{τ}{ι} τ\expo{οῦ} τρέφῃ · τ(ὸ) γ(ὰρ) τρέφει Ἀττικόν {\policeFr ÷}}{« On écrit aussi (variante) : σύ τ' αὐτός (“et toi-même”). — τρέφει : au sens de τρέφῃ (2e pers. sing. moyen-passif “es-tu élevé”), car la désinence -ει pour -ῃ est attique. »}
\commScholie{Double note grammaticale et critique de R\expo{2} : elle atteste une variante antique (σύ τ' αὐτός pour τίς δαὶ σύ) et justifie l'orthographe attique de 2e personne en -ει (τρέφει).}
\entreeScholie{063}{142}{καὶ ποῦ πέος; ποῦ χλαῖνα; ποῦ Λακωνικαί;}{Interlinéaire (supra lineam au-dessus de Λακωνικαί, f. 164r)}{ἀνδρεῖα ὑποδήματα.}{ἀνδρεῖα ὑποδήματα {\policeFr ÷}}{« Des chaussures d'homme (les chaussures rouges de Laconie). »}
\commScholie{Repris par la Souda λ 63 et π 987 (cf. L. Hauseux, vv 1-641.tex, l. 1632).}
\entreeScholie{064}{144}{τί φῄς; τί σιγᾷς; ἀλλὰ δῆτ᾽ ἐκ τοῦ μέλους / ζητῶ σ᾽, ἐπειδή γ᾽ αὐτὸς οὐ βούλει φράσαι;}{Marge droite (extérieure, f. 164r, en face du v. 144)}{στοχάζομαί γε ἐκ τῆς μελοποιίας οἷος ἄν τις εἴης.}{στοχάζομαί γε ἐκ τ\expo{ῆς} μελοποιί\expo{ας} οἷ\expo{ος} ἄν τις εἴης {\policeFr ÷}}{« Je conjecture d'après la nature de ta composition musicale quel genre d'être tu peux bien être. »}
\commScholie{Paraphrase libre expliquant le lien établi par Mnésiloque entre le costume efféminé et la mollesse de la mélodie.}
\entreeScholie{065}{147}{ἤκουσα, τὴν δ᾽ ἄλγησιν οὐ παρεσχόμην·}{Interlinéaire (supra lineam, f. 164r, v. 147)}{ἀντὶ τοῦ οὐκ ἐλυπήθην.}{ἀν\auDessus{τ}{ι} τ\expo{οῦ} οὐκ ἐλυπήθ(ην) {\policeFr ÷}}{« Au sens de : je n'ai pas éprouvé de chagrin (j'ai écarté la peine de mon esprit). »}
\commScholie{Explique le sens de la réplique sentencieuse d'Agathon face aux sarcasmes de Mnésiloque.}
\entreeScholie{066}{152}{μετουσίαν δεῖ τῶν τρόπων τὸ σῶμ᾽ ἔχειν.}{Interlinéaire (supra lineam, f. 164r, v. 152)}{ἀντὶ τοῦ τὴν σκέψιν.}{ἀν\auDessus{τ}{ι} τ\expo{οῦ} τ(ὴν) σκέψ(ιν) {\policeFr ÷}}{« Au sens de : la méditation (l'invention poétique). »}
\commScholie{Glose rapportant l'accord nécessaire du corps à l'effort de conception dramatique (vv. 149-152).}
\entreeScholie{067}{154}{ἤν τις ποιῇ γυναικεῖα δράματα (v. 151) / ἀνδρεῖα δ᾽ ἂν [δ᾽ ἢν R] ποιῇ τις, ἐν τῷ σώματι (v. 154)}{Marge droite (extérieure, f. 164r, en face du v. 154)}{γυναικεῖα δράματα λέγεται ἐν οἷς ὁ χορὸς ἐκ γυναικῶν ἐστὶν, ἀνδρεῖα δὲ ἐν οἷς ἐξ ἀνδρῶν.}{γυναικεῖα δράματα λέγετ(αι) ἐν οἷς ὁ χορ(ὸς) ἐκ γυναικ\expo{ῶν} ἐστὶν, ἀνδρεῖα δ(ὲ) ἐν οἷς ἐξ ἀνδρ\expo{ῶν} {\policeFr ÷}}{« On appelle “drames féminins” (γυναικεῖα δράματα) ceux dans lesquels le chœur est composé de femmes, et “drames masculins” (ἀνδρεῖα) ceux dans lesquels il est composé d'hommes. »}
\commScholie{Définition technique issue de la critique dramatique péripatéticienne et alexandrine.}
\entreeScholie{068}{158}{ὅταν σατύρους τοίνυν ποιῇς, καλεῖν ἐμέ, / ἵνα συμποιῶ σοὔπισθεν ἐστυκὼς ἐγώ·}{Marge droite (extérieure, f. 164r, en face du v. 158)}{ἐπεὶ ὡς ἐπίπαν ἐστυκότας τοὺς σατύρους εἰσάγουσιν.}{ἐπεὶ ὡς ἐπίπ(αν) ἐστυκότ\expo{ας} τ\expo{οὺς} σατύρ\expo{ους} εἰσάγουσ(ιν) {\policeFr ÷}}{« Parce que la plupart du temps on met en scène les satyres le phallus en érection (ἐστυκότας). »}
\commScholie{Explique l'allusion au costume théâtral des choreutes de drame satyrique (munis du phallus de cuir dressé) et la grivoiserie de Mnésiloque proposant de collaborer “par derrière”.}
\entreeScholie{069}{160}{ἄλλως τ᾽ ἄμουσόν ἐστι ποιητὴν ἰδεῖν / ἀγρεῖον ὄντα καὶ δασύν·}{Interlinéaire (supra lineam au-dessus du v. 160, f. 164r)}{ἀγροῖκον καὶ πώγωνα μέγαν ἔχοντα.}{ἀγροῖκ\expo{ον} \textbf{ς} πώγωνα μέγ(αν) ἔχοντ(α) {\policeFr ÷}}{« D'aspect rustre et portant une grosse barbe hirsute. »}
\commScholie{Glose explicative sur l'aspect physique incompatible, selon Agathon, avec la délicatesse lyrique.}
\entreeScholie{070}{161}{Ἴβυκος ἐκεῖνος κἀνακρέων ὁ Τήιος / κἀλκαῖος,}{Marge droite (extérieure, f. 164r, en face du v. 161)}{Ἴβυκος ὁ Ῥηγῖνος καὶ Ἀνακρέων ὁ ἀπὸ Τέω (αὕτη δὲ πόλις τῆς Ἰωνίας), καὶ Ἀλκαῖος ὁ Λέσβιος μέλη ἔγραψαν πρὸς λύραν.}{Ἴβυκ\expo{ος} ὁ Ῥηγῖν\expo{ος} \textbf{ς} Ἀνακρέ(ων) ὁ ἀπ(ὸ) Τέω (αὕτη δ(ὲ) πόλ(ις) τ\expo{ῆς} Ἰωνί\expo{ας}), \textbf{ς} Ἀλκαῖ\expo{ος} ὁ Λέσβι\expo{ος} μέλ(η) ἔγραψ(αν) πρ(ὸς) λύρ(αν) {\policeFr ÷}}{« Ibycos de Rhégion, Anacréon de Téos (cité d'Ionie) et Alcée de Lesbos ont composé des chants lyriques accompagnés à la lyre. »}
\commScholie{Notice biographique sur la triade des grands lyriques archaïques évoqués par Agathon (commentée par L. Hauseux, vv 1-641.tex, l. 1685).}
\entreeScholie{071}{162}{κἀλκαῖος, οἳ περὶ ἁρμονίαν ἐκεχήνεσαν [ἐχύμισαν R],}{Marge inférieure et marge extérieure droite (f. 164r, grande scholie alexandrine de 16 lignes au bas du recto)}{ἔγχυμον ἐποίησαν. ἐν ἐνίοις δὲ Ἀχαιὸς γέγραπται, καὶ τὰ παλαιότερα ἀντίγραφα οὕτως εἶχεν. Ἀριστοφάνης δέ ἐστιν ὁ μεταγράψας Ἀλκαῖος· περὶ γὰρ παλαιῶν ἐστὶν ὁ λόγος, ὁ δὲ Ἀχαιὸς νεώτερος. τὸ δὲ λεγόμενον ὑπὸ Διδύμου πρὸς Ἀριστοφάνην, ὅτι οὐ δύναται Ἀλκαίου μνημονεύειν (οὐ γὰρ ἐπεπόλαζε, φησὶ, τὰ Ἀλκαίου διὰ τὴν διάλεκτον), λελήρηται ἄντικρυς. καὶ ἐν τῷ πρὸ τούτου δράματι, τοῖς Ὄρνισι, παρῴδηται τὸ «ὄρνιθες τίνες οἵδ’ ὠκεανῶ γᾶς ἀπὺ περράτων» οὕτως «ὄρνιθες τίνες οἵδ’ οὐδὲν ἔχοντες» καὶ ἐν Σφηξὶν «ὤνηρ οὗτος ὁ μαιόμενος τὸ μέγα κράτος.» ἀλλαχοῦ δὲ ὁ Δίδυμός φησιν, ἡ μὲν γραφὴ δύναται μένειν, οὐκ ἂν δὲ τούτου τοῦ μελοποιοῦ μέμνηται, πάλιν τὸ αὐτὸ λέγω ὅτι οὐκ ἐπεπόλαζε τὰ μέλη, ἀλλ’ Ἀλκαίου τοῦ κιθαρῳδοῦ, οὗ καὶ Εὔπολις ἐν Χρυσῷ Γένει μέμνηται «ὦ ’λκαῖε Σικελιῶτα Πελοποννήσιε.» τί δὲ ἐνταῦθα κιθαρῳδοῦ, περὶ ποιητοῦ ὄντος τοῦ λόγου;}{ἔγχυμ(ον) ἐποίησ(αν) · ἐν ἐνίοις δ(ὲ) Ἀχαιὸς γέγραπτ(αι), \textbf{ς} τὰ παλ(αι)ότερα ἀντίγραφα οὕτ\expo{ως} εἶχ(εν) · Ἀριστοφάνης δέ ἐστ(ιν) ὁ μεταγράψας Ἀλκαῖ\expo{ος}· π(ερὶ) γ(ὰρ) παλαι\expo{ῶν} ἐστ(ὶν) ὁ λόγος, ὁ δ(ὲ) Ἀχαιὸς νεώτερ\expo{ος} · τ(ὸ) δ(ὲ) λεγόμεν(ον) ὑπ(ὸ) Διδύμου πρ(ὸς) Ἀριστοφάνην, ὅτι οὐ δύνατ(αι) Ἀλκαίου μνημονεύειν (οὐ γ(ὰρ) ἐπεπόλαζε, φησ(ὶ), τὰ Ἀλκαίου διὰ τ(ὴν) διάλεκτ(ον)), λελήρητ(αι) ἄντικρυς · \textbf{ς} ἐν τῷ πρὸ τούτου δράματι, τ\expo{οῖς} Ὄρνισι, παρῴδητ(αι) τ(ὸ) «ὄρνιθες τίνες οἵδ’ ὠκεανῶ γᾶς ἀπὺ περράτων» οὕτ\expo{ως} «ὄρνιθες τίνες οἵδ’ οὐδὲν ἔχοντες» \textbf{ς} ἐν Σφηξὶν «ὤνηρ οὗτ\expo{ος} ὁ μαιόμεν\expo{ος} τ(ὸ) μέγα κράτος · » ἀλλαχ\expo{οῦ} δ(ὲ) ὁ Δίδυμός φησ(ιν), ἡ μ(ὲν) γραφὴ δύνατ(αι) μένε(ιν), οὐκ ἂν δ(ὲ) τούτου τ\expo{οῦ} μελοποι\expo{οῦ} μέμνητ(αι), πάλιν τ(ὸ) αὐτὸ λέγω ὅτι οὐκ ἐπεπόλαζε τὰ μέλη, ἀλλ’ Ἀλκαίου τ\expo{οῦ} κιθαρῳδ\expo{οῦ}, οὗ \textbf{ς} Εὔπολις ἐν Χρυσῷ Γένει μέμνητ(αι) «ὦ ’λκαῖε Σικελιῶτα Πελοποννήσιε · » τί δ(ὲ) ἐνταῦθα κιθαρῳδ\expo{οῦ}, π(ερὶ) ποιητ\expo{οῦ} ὄντ\expo{ος} τ\expo{οῦ} λόγου; {\policeFr ÷}}{« ἐχύμισαν : ils ont donné de la saveur (du suc, ἔγχυμον ἐποίησαν) à l'harmonie. — Dans certains manuscrits, on lit Ἀχαιός (Achaïos), et c'est ce qu'avaient les plus anciens exemplaires ; mais c'est Aristophane (de Byzance) qui a corrigé en Ἀλκαῖος (Alcée), car il est question ici des poètes anciens, tandis qu'Achaïos est plus récent. Quant à l'objection faite par Didyme contre Aristophane de Byzance, selon laquelle le poète ne pourrait mentionner Alcée parce que les chants d'Alcée n'étaient pas répandus à Athènes à cause de leur dialecte éolien, c'est pur radotage (λελήρηται ἄντικρυς) ! En effet, dans la pièce représentée avant celle-ci, les Oiseaux (v. 1410), est parodié le vers d'Alcée (fr. 345 L.-P.) “Quels sont ces oiseaux venus des confins de l'Océan et de la terre ?” sous la forme “Quels sont ces oiseaux qui n'ont rien ?”, et dans les Guêpes (v. 1234), “Cet homme qui convoite le grand pouvoir” (Alcée, fr. 141 L.-P.). Ailleurs, Didyme dit que la leçon Ἀλκαῖος peut rester, mais qu'il ne s'agirait pas du poète lyrique — je répète que c'est faux de dire que ses chants n'étaient pas connus —, mais d'Alcée le citharède, que mentionne aussi Eupolis dans l'Âge d'or (fr. 303 K.-A.) : “Ô Alcée de Sicile, Péloponnésien !” Mais que viendrait faire ici un citharède, alors qu'il est question de poètes ? »}
\commScholie{Chef-d'œuvre de la critique alexandrine (Symmaque réfutant Didyme pour défendre l'émendation d'Aristophane de Byzance κἀλκαῖος, que conserve R\expo{1}). Louis Hauseux (Bibi, v. 162, ll. 1685-1748) adopte κἀλκαῖος avec Aristophane de Byzance, R et la Souda, mais montre que la glose ἔγχυμον ἐποίησαν sur le barbarisme amétrique ἐχύμησαν (R) / ἐχύμισαν (Souda) est une fantaisie de glossateur, et restitue brillamment le plus-que-parfait attique ἐκεχήνεσαν (« qui bayaient après l'harmonie »).}
\entreeScholie{072}{163}{ἐμιτροφόρουν τε καὶ διέκειντ᾽ [διεκίνων R] Ἰωνικῶς,}{Interlinéaire (supra lineam au-dessus d'Ἰωνικῶς, bas du f. 164r)}{ἁβρῶς.}{ἁβρῶς {\policeFr ÷}}{« Avec mollesse (avec une délicatesse voluptueuse, à la mode ionienne). »}
\commScholie{Comme le souligne Louis Hauseux (vv 1-641.tex, ll. 1692-1706), cette glose (jointe à la Souda ε 989 et ι 495) confirme l'adverbe Ἰωνικῶς et le sens lascif du mouvement corporel (διεκίνων R : διεκλῶντ' Toup : διεκίνουν Souda : conjecture de Bibi discutée en commentaire).}
\entreeScholie{073}{164}{καὶ Φρύνιχος — τοῦτον γὰρ οὖν ἀκήκοας —}{Marge intérieure gauche (dernière ligne du f. 164r)}{Φρύνιχος ὁ τραγικὸς, ἐπεὶ Ἀθηναῖος ἦν.}{Φρύνιχ\expo{ος} ὁ τραγικ\expo{ὸς}, ἐπεὶ Ἀθηναῖ\expo{ος} ἦν {\policeFr ÷}}{« Il s'agit de Phrynichos le poète tragique : “celui-là du moins, tu en as entendu parler”, parce qu'il était Athénien (à la différence d'Ibycos, d'Anacréon et d'Alcée). »}
\commScholie{Cité et approuvé par Louis Hauseux (vv 1-641.tex, l. 1766) : le scholiaste explique très finement l'incise d'Agathon à Mnésiloque, qui pouvait ignorer les lyriques ioniens et éoliens mais connaissait forcément le vieux tragique athénien Phrynichos.}
\clearpage
\renewcommand{\currHeadLE}{ΑΡΙΣΤΟΦΑΝΟΥΣ ΘΕΣΜΟΦΟΡΙΑΖΟΥΣΑΙ · CODEX RAVENNAS 429 (\R) · FOLIO 05/30}
\renewcommand{\currHeadRE}{f. 164v (Classense 166v) · vv. 165–207 · Collation diplomatique}
\renewcommand{\currHeadLO}{ΑΡΙΣΤΟΦΑΝΟΥΣ ΘΕΣΜΟΦΟΡΙΑΖΟΥΣΑΙ · CODEX RAVENNAS 429 (\R) · FOLIO 05/30}
\renewcommand{\currHeadRO}{f. 164v (Classense 166v) · vv. 165–207 · Collation diplomatique}
\noindent\begin{adjustbox}{max totalheight=30.4cm, max width=\textwidth, keepaspectratio}\begin{minipage}{\textwidth}
\bandeauRav{f. 164v}{Classense 166v}{165–207}{Folio 05/30 · Verso (page de gauche) · 8 scholies \& gloses ({\policeGr Σ} 074–081)}
{\fontsize{8.8pt}{10.2pt}\selectfont
\ligneRav{165}{ἀγάθ/}{αὐτός τε καλὸς ἦν καὶ καλῶς ἠμπέσχετο}{\textbf{165} \lecon{ἠμπέσχετο} R : ἠμπισχετο Elmsley ad Eur. \italique{Med.} 1128 (\italique{induebatur} Divus)}
\ligneRav{}{}{διὰ τοῦτ᾽ ἄρ᾽ αὐτοῦ καὶ κάλ᾽ ἦν τὰ δράματα·}{}
\ligneRav{}{}{ὅμοια γὰρ ποιεῖν ἀνάγκη τῇ φύσει·}{}
\ligneRav{168}{μ\auDessus{ν}{Η}}{ταῦτ᾽ ἄρ᾽ ὁ \difference{Φιλοκλέης} αἰσχρὸς ὢν αἰσχρῶς ποιεῖ,}{\textbf{168} Φιλοκλῆς R : \lecon{Φιλοκλέης} Bentley}
\ligneRav{169}{}{ὁ δ᾽ αὖ \difference{Ξενοκλέης} ὢν κακὸς κακῶς ποιεῖ,}{\textbf{169} \lecon{ὁ δ᾽ αὖ} R : ὁ δὲ Brunck \espApp Ξενοκλῆς R : \lecon{Ξενοκλέης} Bentley}
\ligneRav{170}{}{ὁ δ᾽ αὖ Θέογνις ψυχρὸς ὢν ψυχρῶς ποιεῖ·}{}
\ligneRav{}{ἀγάθ/}{ἅπασ᾽ ἀνάγκη ταῦτα γάρ τοι γνοὺς ἐγὼ}{}
\ligneRav{}{—}{ἐμαυτὸν ἐθεράπευσα. : πῶς πρὸς τῶν θεῶν;}{}
\ligneRav{173}{\auDessus{εὐ}{ρ}}{παῦσαι βαΰζων· καὶ γὰρ ἐγὼ τοιοῦτος ἦν}{\textbf{173} \lecon{ἦν} R, Suid. β 192 : ἦ Dindorf (1825)}
\ligneRav{}{}{ὢν τηλικοῦτος, ἡνίκ᾽ ἠρχόμην ποιεῖν·}{}
\ligneRav{175}{μ\auDessus{ν}{Η}}{μὰ τὸν Δί᾽ οὐ ζηλῶ σε τῆς παιδεύσεως·}{}
\ligneRav{176}{\auDessus{εὐ}{ρ}}{ἀλλ᾽ ὧνπερ οὕνεκ᾽ ἦλθον, ἔα μ᾽ εἰπεῖν. : λέγε·}{\textbf{176} \lecon{(Ἀγάθ.)} λέγε R, Fritzsche : (Μνη.) Küster}
\ligneRav{}{\auDessus{εὐ}{ρ}}{Ἀγάθων, σοφοῦ πρὸς ἀνδρός, ὅστις ἐν βραχεῖ}{}
\ligneRav{178}{}{πολλοὺς καλῶς \difference{οἷός} τε συντέμνειν λόγους·}{\textbf{178} οἷόν R : \lecon{οἷός} Suid. σ 792, οι 154, Stob. III,35.3 \espApp \lecon{συντέμνειν} R, Suid. οι 154, Stob. : συντεμεῖν Suid. σ 792}
\ligneRav{179}{}{ἐγὼ δὲ \difference{καινῇ} ξυμφορᾷ πεπληγμένος}{\textbf{179} κοινῇ R, Souda κ 2556, σ 792 (confusion minuscule α → ο) : \lecon{καινῇ} Biset d'après {\policeGr Σ}\expo{R}}
\ligneRav{180}{—}{ἱκέτης ἀφῖγμαι πρὸς σέ. : τοῦ χρείαν ἔχων;}{}
\ligneRav{}{\auDessus{εὐ}{ρ}}{μέλλουσί μ᾽ αἱ γυναῖκες ἀπολεῖν τήμερον}{}
\ligneRav{}{}{τοῖς Θεσμοφορίοις, ὅτι κακῶς αὐτὰς λέγω·}{}
\ligneRav{}{ἀγάθ/}{τίς οὖν παρ᾽ ἡμῶν ἐστιν ὠφέλειά σοι;}{}
\ligneRav{184}{\auDessus{εὐ}{ρ}}{\difference{ἡ πᾶσ᾽·} ἐὰν γὰρ ἐγκαθεζόμενος λάθρᾳ}{\textbf{184} ἡ πᾶσα R : \lecon{ἡ πᾶσ᾽·} Küster : ἅπασ' Scaliger}
\ligneRav{185}{}{ἐν ταῖς γυναιξίν, ὡς δοκῶν εἶναι γυνή,}{}
\ligneRav{}{}{ὑπεραποκρίνῃ μου, σαφῶς σώσεις ἐμέ·}{}
\ligneRav{}{}{μόνος γὰρ ἂν λέξειας ἀξίως ἐμοῦ·}{}
\ligneRav{}{ἀγάθ/}{ἔπειτα πῶς οὐκ αὐτὸς ἀπολογεῖ παρών;}{}
\ligneRav{}{\auDessus{εὐ}{ρ}}{ἐγὼ φράσω σοι. πρῶτα μὲν γιγνώσκομαι}{}
\ligneRav{190}{}{ἔπειτα πολιός εἰμι καὶ πώγων᾽ ἔχω,}{}
\ligneRav{}{}{σὺ δ᾽ εὐπρόσωπος λευκὸς ἐξυρημένος}{}
\ligneRav{}{}{γυναικόφωνος ἁπαλὸς εὐπρεπὴς ἰδεῖν·}{}
\ligneRav{}{ἀγάθ/}{Εὐριπίδη— : τί ἔστιν; : ἐποίησάς ποτε,}{}
\ligneRav{}{}{`χαίρεις ὁρῶν φῶς, πατέρα δ᾽ οὐ χαίρειν δοκεῖς;'}{}
\ligneRav{195}{\auDessus{εὐ}{ρ}}{ἔγωγε. : μή νυν ἐλπίσῃς τὸ σὸν κακὸν}{}
\ligneRav{196}{}{ἡμᾶς ὑφέξειν. καὶ \difference{γὰρ ἂν} μαινοίμεθ᾽ ἂν·}{\textbf{196} γὰρ R : \lecon{γὰρ ἂν} Suid. υ 709}
\ligneRav{}{}{ἀλλ᾽ αὐτὸς ὅ γε σόν ἐστιν οἰκείως φέρε·}{}
\ligneRav{}{}{τὰς συμφορὰς γὰρ οὐχὶ τοῖς τεχνάσμασιν}{}
\ligneRav{}{}{φέρειν δίκαιον ἀλλὰ τοῖς παθήμασιν·}{}
\ligneRav{200}{μ\auDessus{ν}{Η}}{καὶ μὴν σύ γ᾽ ὦ κατάπυγον εὐρύπρωκτος εἶ}{}
\ligneRav{}{}{οὐ τοῖς λόγοισιν ἀλλὰ τοῖς παθήμασιν·}{}
\ligneRav{}{\auDessus{εὐ}{ρ}}{τί δ᾽ ἔστιν ὅτι δέδοικας ἐλθεῖν αὐτόσε;}{}
\ligneRav{}{ἀγάθ/}{κάκιον ἀπολοίμην ἂν ἢ σύ. : πῶς; : ὅπως;}{}
\ligneRav{204}{}{δοκῶν γυναικῶν ἔργα νυκτερείσια}{\textbf{204} \lecon{νυκτερείσια} R : νυκτερήσια Bothe}
\ligneRav{205}{}{κλέπτειν ὑφαρπάζειν τε θήλειαν Κύπριν·}{}
\ligneRav{206}{μ\auDessus{ν}{Η}}{ἰδού γε κλέπτειν· νὴ Δία βινεῖσθαι μὲν οὖν·}{\textbf{206} \lecon{νὴ Δία} R : μà Δία Fritzsche}
\ligneRav{}{}{ἀτὰρ ἡ πρόφασίς γε νὴ Δί᾽ εἰκότως ἔχει·}{}
}
\vspace{3pt}{\color{ruleGold}\hrule height 0.6pt}\vspace{3pt}
{\fontsize{8.0pt}{9.4pt}\selectfont
\noindent{\bfseries\scshape\color{rubricRed}Loci paralleli et testimonia (f. 164v, vv. 165–207) :} \textbf{168} (ὁ Φιλοκλέης αἰσχρὸς = κορυδός) scholium \italique{Aves} 1295a.α \espTest \textbf{173} (βαύ̈ζων -- ἦν) Suidas β 192 {\policeFr ≈} Zonaras 379 (sed γὰρ καὶ pro καὶ γὰρ) \espTest \textbf{175} (οὐ -- παιδεύσεως) Suidas ζ 63 \espTest \textbf{176} Scholium v. 500 (sed ἀλλ' εἰπεῖν βούλομαι pro ἔα μ' εἰπεῖν) \espTest \textbf{177-80} (σοφοῦ -- πρὸς σέ) Suidas σ 792 \espTest \textbf{177-8} (σοφοῦ -- λόγους) Suidas οι 154 \espTest \textbf{179} Suidas κ 2556 \espTest \textbf{192} (γυναικόφωνος) Pollux II 111 \espTest \textbf{193} (ἐποίησάς ποτε) \Eustathe 4.9 \espTest \textbf{195-7} (μή νυν -- φέρε) Suidas υ 709 (νῦν pro νυν et ἀλλ' omisso) \espTest \textbf{198-201} (τὰς συμφορὰς -- παθήμασιν) Suidas τ 440 (-σι pro -σιν) %(τεχνάσμασι et παθήμασι pro -σιν)\par\vspace{2.5pt}
{\color{ruleGold}\hrule height 0.35pt}\vspace{2.5pt}
\noindent{\bfseries\scshape\color{rubricRed}Suite (en commentaire) de l'apparat critique et paléographique :} \textbf{Contexte :} Fin de la tirade d'Agathon (vv. 165–172), requête d'Euripide priant Agathon d'aller plaider sa cause aux Thesmophories et refus prudent d'Agathon (vv. 173–207). --- Page très régulière de 43 lignes. On relève les formes contractes non métriques de noms propres tragiques (Φιλοκλῆς, Ξενοκλῆς R aux vv. 168–169 pour -κλέης), la confusion minuscule α/ο dans κοινῇ (v. 179 R et Souda) pour καινῇ (Biset d'après {\policeGr Σ}\expo{R}), et la graphie iotacisante νυκτερείσια (v. 204 R, maintenue par L. Hauseux). --- \textbf{v. 179 :} κοινῇ R, Souda κ 2556, σ 792 (confusion minuscule α → ο) : \lecon{καινῇ} Biset d'après {\policeGr Σ}\expo{R}.\par}
\end{minipage}\end{adjustbox}
\newpage
\renewcommand{\currHeadLE}{SCHOLIES MARGINALES ({\policeGr Σ}\expo{R}) \& GLOSES INTERLINÉAIRES · FOLIO 05/30}
\renewcommand{\currHeadRE}{f. 164v (Classense 166v) · vv. 165–207 · {\policeGr Σ} 074 à {\policeGr Σ} 081}
\renewcommand{\currHeadLO}{SCHOLIES MARGINALES ({\policeGr Σ}\expo{R}) \& GLOSES INTERLINÉAIRES · FOLIO 05/30}
\renewcommand{\currHeadRO}{f. 164v (Classense 166v) · vv. 165–207 · {\policeGr Σ} 074 à {\policeGr Σ} 081}
\bandeauScholies{Folio 05/30 — f. 164v}{Classense 166v}{vv. 165–207}{Verso (page de gauche) · 8 entrées ({\policeGr Σ} 074 à {\policeGr Σ} 081)}
{\fontsize{8.8pt}{10.5pt}\selectfont\noindent\textbf{\color{rubricRed}Disposition topographique des marges sur le f. 164v :} Marge extérieure à gauche ; marge intérieure à droite. — \textit{Fac-similé de référence :} \texttt{folios\_individuels\_ff162v-177r/05\_f164v\_classense166v\_vv0165-0207.jpg}.\par\vspace{3pt}}
\entreeScholie{074}{168}{ταῦτ᾽ ἄρ᾽ ὁ Φιλοκλέης αἰσχρὸς ὢν αἰσχρῶς ποιεῖ,}{Marge gauche (extérieure, f. 164v, bloc de 7 lignes en haut du verso)}{ἐκ τούτου [...] καὶ ἐν τοῖς Ὄρνισιν [1295] οὐκ ἄλλοθεν κορυδὸν αὐτὸν ὠνόμασεν ἀλλ’ ἢ ὅτι μικρὸς καὶ αἰσχρός. καὶ Τηλεκλείδης Ἡσιόδοις «ἀλλ’ ἡ τάλαινα Φιλοκλέα [...] οὖν, εἰ δ’ ἐστὶν Αἰσχύλου φρόνημ’ ἔχων.» τραγῳδίας ποιητὴς αἰσχρὸς τὸ πρόσωπον καὶ σιμὸς καὶ μικρὸς τὸ σῶμα.}{ἐκ τούτου \textbf{ς} ἐν τ\expo{οῖς} Ὄρνισ(ιν) οὐκ ἄλλοθ(εν) κορυδ(ὸν) αὐτ(ὸν) ὠνόμασ(εν) ἀλλ’ ἢ ὅτι μικρὸς \textbf{ς} αἰσχρός · \textbf{ς} Τηλεκλείδης Ἡσιόδοις «ἀλλ’ ἡ τάλαινα Φιλοκλέα [...] οὖν, εἰ δ’ ἐστ(ὶν) Αἰσχύλου φρόνημ’ ἔχ\expo{ων} · » τραγῳδίας ποιητὴς αἰσχρὸς τ(ὸ) πρόσωπ(ον) \textbf{ς} σιμὸς \textbf{ς} μικρὸς τ(ὸ) σῶμα {\policeFr ÷}}{« C'est pour cela que dans les Oiseaux (v. 1295) aussi, il ne l'a appelé “Alouette huppée” (κορυδός) pour aucune autre raison sinon qu'il était petit et laid. Et Téléclide dit dans ses Hésiodes (fr. 15 K.-A.) : “Mais la malheureuse [...] Philoclès donc, s'il a l'orgueil d'Eschyle.” — Autre scholie : Poète tragique (neveu d'Eschyle), laid de visage, camus et petit de corps. »}
\commScholie{Double scholie érudite (reprise par la Souda φ 379 et discutée par L. Hauseux, vv 1-641.tex, ll. 1782-1800) conservant un précieux fragment de Téléclide et expliquant le sobriquet physique de Κόρυδος donné à Philoclès.}
\entreeScholie{075}{173}{ἐμαυτὸν ἐθεράπευσα. — πῶς πρὸς τῶν θεῶν; / παῦσαι βαΰζων· καὶ γὰρ ἐγὼ τοιοῦτος ἦν}{Interlinéaire (supra lineam, f. 164v, v. 173)}{κακὸς ἤμην.}{κακ(ὸς) ἤμην {\policeFr ÷}}{« J'étais atteint de ce défaut (avant de m'en guérir). »}
\commScholie{Explique l'ellipse dans la réplique de Mnésiloque (οὐδ' ἂν ἐγὼ γὰρ ἦν, corr. par le sens : “je me suis soigné quand j'ai connu Agathon”).}
\entreeScholie{076}{177}{σοφοῦ πρὸς ἀνδρός, ὅστις ἐν βραχεῖ / πολλοὺς καλῶς οἷός τε συντέμνειν λόγους.}{Marge gauche (extérieure, f. 164v, en face du v. 177)}{σοφοὺς μᾶλλον ἔλεγον καὶ σοφιστὰς τοὺς περὶ μουσικὴν καὶ ποιητικήν.}{σοφ\expo{οὺς} μᾶλλ\expo{ον} ἔλεγ\expo{ον} \textbf{ς} σοφιστ\expo{ὰς} τ\expo{οὺς} π(ερὶ) μουσικ(ὴν) \textbf{ς} ποιητικήν {\policeFr ÷}}{« Les anciens appelaient plutôt “sages” (σοφοί) et “sophistes” (σοφισταί) ceux qui s'occupaient de musique et de poésie. »}
\commScholie{Observation lexicologique classique (remontant à Aristarque) sur l'emploi archaïque et classique de σοφός = poète.}
\entreeScholie{077}{179}{ἐγὼ δὲ καινῇ ξυμφορᾷ πεπληγμένος}{Interlinéaire (supra lineam au-dessus de κοινῇ de R\expo{1}, f. 164v)}{μεγάλῃ καὶ θαυμαστῇ.}{μεγάλῃ \textbf{ς} θαυμαστῇ {\policeFr ÷}}{« Par un malheur grand et extraordinaire (inouï, étrange : καινῇ). »}
\commScholie{Cas exemplaire analysé en détail par Louis Hauseux (Bibi, v. 179, vv 1-641.tex, ll. 1869-1878) : alors que R\expo{1} a commis la faute d'onciale/minuscule κοινῇ (« commun »), la glose de R\expo{2} μεγάλῃ καὶ θαυμαστῇ (« grand et extraordinaire ») traduit indubitablement la vraie leçon καινῇ (« nouveau, inouï », restituée par Biset et adoptée par Bibi). Le compilateur de la Souda (κ 2505), disposant à la fois du texte déjà altéré en κοινῇ et de la glose de R\expo{2}, a accolé mécaniquement les deux : ἀνθρωπίνῃ (« humain, commun » = κοινῇ), μεγάλῃ, θαυμαστῇ (= καινῇ) !}
\entreeScholie{078}{190}{ἔπειτα πολιός εἰμι καὶ πώγων᾽ ἔχω, / σὺ δ᾽ εὐπρόσωπος λευκὸς ἐξυρημένος}{Marge gauche (extérieure, f. 164v, en face du v. 190)}{γέρων γὰρ τότε Εὐριπίδης ἦν· ἕκτῳ γοῦν ἔτει ὕστερον τελευτᾷ.}{γέρ(ων) γ(ὰρ) τότε Εὐριπίδ(ης) ἦν · ἕκτῳ γοῦν ἔτ(ει) ὕστερ\expo{ον} τελευτᾷ {\policeFr ÷}}{« Car Euripide était alors un vieillard : de fait, il meurt la sixième année après cette pièce (en 406/405 av. J.-C., en comptant inclusivement depuis 411). »}
\commScholie{Indice chronologique fondamental (commenté par L. Hauseux, Introduction.tex, l. 175) : Euripide étant mort au début de 406 av. J.-C. (sous l'archontat de Callias), le compte inclusif ἕκτῳ ἔτει ὕστερον confirme la date de 411 av. J.-C. pour les Thesmophories.}
\entreeScholie{079}{194}{χαίρεις ὁρῶν φῶς, πατέρα δ᾽ οὐ χαίρειν δοκεῖς;}{Marge gauche (extérieure, f. 164v, en face du v. 194)}{ἐξ Ἀλκήστιδος Εὐριπίδου [694].}{ἐξ Ἀλκήστιδ\expo{ος} Εὐριπίδ\expo{ου} {\policeFr ÷}}{« Ce vers est tiré de l'Alceste d'Euripide (v. 691). »}
\commScholie{Identification exacte de la citation tragique (Alceste v. 691, réplique de Phérès à son fils Admète qui lui demandait de mourir à sa place), qu'Agathon retourne spirituellement contre Euripide.}
\entreeScholie{080}{199}{τὰς συμφορὰς γὰρ οὐχὶ τοῖς τεχνάσμασιν / φέρειν δίκαιον, ἀλλὰ τοῖς παθήμασιν·}{Interlinéaire (supra lineam au-dessus de φέρειν, f. 164v)}{ἐκφεύγειν.}{ἐκφεύγ(ειν) {\policeFr ÷}}{« Échapper à (esquiver les malheurs par des ruses). »}
\commScholie{Précise le sens zeugmatique de φέρειν associé à τοῖς τεχνάσμασιν.}
\entreeScholie{081}{204}{δοκῶν γυναικῶν ἔργα νυκτερείσια / κλέπτειν ὑφαρπάζειν τε θήλειαν Κύπριν·}{Interlinéaire (supra lineam, f. 164v, v. 203)}{ἐὰν τοῦτο ποιήσωμεν.}{ἐὰν τοῦτ(ο) ποιήσωμ(εν) {\policeFr ÷}}{« Si nous faisons cela (si je me glisse secrètement parmi les femmes). »}
\commScholie{Glose explicative sur la crainte exprimée par Agathon d'être pris pour un galant nocturne.}
\clearpage
\renewcommand{\currHeadLE}{ΑΡΙΣΤΟΦΑΝΟΥΣ ΘΕΣΜΟΦΟΡΙΑΖΟΥΣΑΙ · CODEX RAVENNAS 429 (\R) · FOLIO 06/30}
\renewcommand{\currHeadRE}{f. 165r (Classense 167r) · vv. 208–251 · Collation diplomatique}
\renewcommand{\currHeadLO}{ΑΡΙΣΤΟΦΑΝΟΥΣ ΘΕΣΜΟΦΟΡΙΑΖΟΥΣΑΙ · CODEX RAVENNAS 429 (\R) · FOLIO 06/30}
\renewcommand{\currHeadRO}{f. 165r (Classense 167r) · vv. 208–251 · Collation diplomatique}
\noindent\begin{adjustbox}{max totalheight=30.4cm, max width=\textwidth, keepaspectratio}\begin{minipage}{\textwidth}
\bandeauRav{f. 165r}{Classense 167r}{208–251}{Folio 06/30 · Recto (page de droite) · 13 scholies \& gloses ({\policeGr Σ} 082–094)}
{\fontsize{8.8pt}{10.2pt}\selectfont
\ligneRav{208}{\auDessus{εὐ}{ρ}}{τί \difference{οὖν}; ποιήσεις \difference{ταῦτα}; : μὴ δόκει γε σύ·}{\textbf{208} οὖν et ταῦτα R : \lecon{οὖν;} et \lecon{ταῦτα;} Ellebodius}
\ligneRav{209}{\auDessus{εὐ}{ρ}}{ὦ τρισκακοδαίμον, ὡς ἀπόλωλ᾽ Εὐριπίδης·}{\textbf{209} \lecon{τρισκακόδαιμον} R : τρισκακοδαίμων \Mon \espApp \lecon{(Εὐρ.)} R : (Μνη.) Elmsley ad \italique{Ach}. 475}
\ligneRav{210}{μ\auDessus{ν}{Η}}{ὦ φίλτατ᾽ ὦ κηδεστὰ μὴ σαυτὸν προδῷς·}{}
\ligneRav{}{\auDessus{εὐ}{ρ}}{πῶς οὖν ποιήσω δῆτα; : τοῦτον μὲν μακρὰ}{}
\ligneRav{}{}{κλάειν κέλευ᾽, ἐμοὶ δ᾽ ὅ τι βούλει χρῶ λαβών·}{}
\ligneRav{}{\auDessus{εὐ}{ρ}}{ἄγε νυν ἐπειδὴ σαυτὸν ἐπιδίδως ἐμοί,}{}
\ligneRav{}{—}{ἀπόδυθι τουτὶ θοἰμάτιον. : καὶ δὴ χαμαί·}{}
\ligneRav{215}{—}{ἀτὰρ τί μέλλεις δρᾶν \difference{μ᾽}; : ἀποξυρεῖν ταδί,}{\textbf{215} με R : \lecon{μ᾽} edd. omnes}
\ligneRav{216}{—}{τὰ κάτω δ᾽ ἀφεύειν. : ἀλλὰ πρᾶττ᾽, \difference{εἴ σοι} δοκεῖ·}{\textbf{216} \lecon{ἀφεύειν} R : ἀφαύειν Suid. α 4565 \espApp είσω R : \lecon{εἴ σοι} Bonini}
\ligneRav{217}{}{ἢ μὴ \difference{διδόναι γ' ἑαυτὸν} ὤφελόν ποτε·}{\textbf{217} διδόναι γ' αὐτὸν R : διδόναι γ' ἑμαυτὸν Scaliger : \lecon{διδόναι γ' ἑαυτὸν} Bibi auctore Willems}
\ligneRav{}{\auDessus{εὐ}{ρ}}{Ἀγάθων σὺ μέντοι ξυροφορεῖς ἑκάστοτε,}{}
\ligneRav{}{—}{χρῆσόν τι νῦν ἡμῖν ξυρόν. : αὐτὸς λάμβανε}{}
\ligneRav{220}{—}{ἐντεῦθεν ἐκ τῆς \difference{ξυροδόκης}. : γενναῖος εἶ·}{\textbf{220} ξυροδίκης R : \lecon{ξυροδόκης} Σ, Pollux II.32, X.140}
\ligneRav{}{}{κάθιξε· φύσα τὴν γνάθον τὴν δεξιάν·}{}
\ligneRav{222}{μ\auDessus{ν}{Η}}{\difference{οἴμοι}. : τί κέκραγας; ἐμβαλῶ σοι πάτταλον,}{\textbf{222} ὤιμοι R : \lecon{οἴμοι} Dindorf}
\ligneRav{223}{—}{ἢν μὴ σιωπᾷς. : \difference{ἀτταταῖ ἰατταταῖ}·}{\textbf{223} ἀττατα· ατταται R : \lecon{ἀτταταῖ ἰατταταῖ} Biset}
\ligneRav{}{\auDessus{εὐ}{ρ}}{οὗτος σὺ ποῖ θεῖς; : εἰς τὸ τῶν σεμνῶν θεῶν·}{}
\ligneRav{225}{}{οὐ γὰρ μὰ τὴν Δήμητρά γ᾽ ἐνταυθοῖ μενῶ}{\textbf{225} \lecon{Δήμητρά γ᾽} R : Δήμητρ᾽ ἔτ᾽ Porson, cf. ad \italique{Eq}. 698 : Δημήτερ᾽ Fritzsche}
\ligneRav{226}{—}{τεμνόμενος. : \difference{οὔκουν} καταγέλαστος δῆτ᾽ ἔσει}{\textbf{226} οὐκ οὖν R : \lecon{οὔκουν} Brunck}
\ligneRav{}{}{τὴν ἡμίκραιραν τὴν ἑτέραν ψιλὴν ἔχων;}{}
\ligneRav{}{μ\auDessus{ν}{Η}}{ὀλίγον μέλει μοι. : μηδαμῶς πρὸς τῶν θεῶν}{}
\ligneRav{}{—}{προδῷς με· χώρει δεῦρο. : κακοδαίμων ἐγώ·}{}
\ligneRav{230}{\auDessus{εὐ}{ρ}}{ἔχ᾽ ἀτρέμα σαυτὸν κἀνάκυπτε· ποῖ στρέφει;}{}
\ligneRav{}{μ\auDessus{ν}{Η}}{μυμῦ. : τί μύζεις; πάντα πεποίηται καλῶς·}{}
\ligneRav{232}{μ\auDessus{ν}{Η}}{οἴμοι κακοδαίμων, ψιλὸς αὖ στρατεύσομαι·}{\textbf{232} \lecon{αὖ} R, Suid. ψ 105 : οὖν von Velsen}
\ligneRav{}{\auDessus{εὐ}{ρ}}{μὴ φροντίσῃ· ὡς εὐπρεπὴς φανεῖ πάνυ·}{}
\ligneRav{234}{—}{βούλει \difference{θεάσασθ᾽ αὑτόν}; : εἰ δοκεῖ, φέρε·}{\textbf{234} θεάσασθαι σεαυτόν R : \lecon{θεάσασθ᾽ αὑτόν} Fritzsche : θεᾶσθαι σαυτόν Porson}
\ligneRav{235}{\auDessus{εὐ}{ρ}}{ὁρᾷς σεαυτόν; : οὐ μὰ Δί᾽ ἀλλὰ \difference{Κλεισθένη}·}{\textbf{235} Κλεισθένην R, Suid. κ 1756 : \lecon{Κλεισθένη} Brunck}
\ligneRav{}{\auDessus{εὐ}{ρ}}{ἀνίστασ᾽, ἵν᾽ ἀφεύσω σε, κἀγκύψας ἔχε·}{}
\ligneRav{}{μ\auDessus{ν}{Η}}{οἴμοι κακοδαίμων δελφάκιον γενήσομαι·}{}
\ligneRav{}{\auDessus{εὐ}{ρ}}{ἐνεγκάτω τις ἔνδοθεν δᾷδ᾽ ἢ λύχνον·}{}
\ligneRav{}{}{ἐπίκυπτε· τὴν κέρκον φυλάττου νυν ἄκραν·}{}
\ligneRav{240}{μ\auDessus{ν}{Η}}{ἐμοὶ μελήσει νὴ Δία, πλήν γ᾽ ὅτι κάομαι·}{}
\ligneRav{}{}{οἴμοι τάλας. ὕδωρ ὕδωρ ὦ γείτονες}{}
\ligneRav{242}{}{πρὶν ἀντιλαβέσθαι πρωκτὸν \difference{[ἓ]τερον} τῆς φλογός·}{\textbf{242} \lecon{[ἓ]τερον} \PsiUn \ (restituit Medaglia) : ἓτερον om. R}
\ligneRav{}{\auDessus{εὐ}{ρ}}{θάρρει. : τί θαρρῶ καταπεπυρπολημένος;}{}
\ligneRav{}{\auDessus{εὐ}{ρ}}{ἀλλ᾽ οὐκ ἔτ᾽ οἰδὲν πρᾶγμά σοι· τὰ πλεῖστα γὰρ}{}
\ligneRav{245}{—}{ἀποπεπόνηκας. : φεῦ ἰοὺ τῆς ἀσβόλου·}{\textbf{245} \lecon{φεῦ} R, Suid. α 4132 : φῦ Dindorf, Wilamowitz \italique{Lysistr.} v.294}
\ligneRav{246}{}{αἰθὸς \difference{γεγένημαι} πάντα τὰ περὶ τὴν τράμιν·}{\textbf{246} γεγέννημαι R : \lecon{γεγένημαι} Suid. τ 903}
\ligneRav{247}{\auDessus{εὐ}{ρ}}{μὴ φροντίσῃς· ἕτερος γὰρ αὐτὰ σπογγιεῖ·}{\textbf{247} \lecon{σπογγιεῖ} R : σφογγιεῖ Elmsley ad \italique{Ach.} 463}
\ligneRav{248}{μ\auDessus{ν}{Η}}{\difference{οἰμώξετ᾽} ἆρ᾽ \difference{εἴ τις} τὸν ἐμὸν πρωκτὸν πλυνεῖ·}{\textbf{248} οἰμώζετ᾽ R : \lecon{οἰμώξετ᾽} Zanetti \espApp εἰς R : \lecon{εἴ τις} Brunck}
\ligneRav{}{\auDessus{εὐ}{ρ}}{Ἀγάθων, ἐπειδὴ σαυτὸν ἐπιδοῦναι φθονεῖς,}{}
\ligneRav{250}{}{ἀλλ᾽ ἱμάτιον γοῦν χρῆσον ἡμῖν \difference{τουτῳὶ}}{\textbf{250} τουτωί R : \lecon{τουτῳὶ} edd. plerique}
\ligneRav{}{}{καὶ στρόφιον· οὐ γὰρ ταῦτά γ᾽ ὡς οὐκ ἔστ᾽ ἐρεῖς·}{}
}
\vspace{3pt}{\color{ruleGold}\hrule height 0.6pt}\vspace{3pt}
{\fontsize{8.0pt}{9.4pt}\selectfont
\noindent{\bfseries\scshape\color{rubricRed}Loci paralleli et testimonia (f. 165r, vv. 208–251) :} \textbf{215-6} (ἀποξυρεῖν -- ἀφεύειν) Suidas α 4565 (ἀφαύειν pro ἀφεύειν), cf. Scholium ταῦτα δὲ ἔλαβεν ἐκ τῶν Ἰδαίων Κρατίνου (fr. 90 Kassel-Austin) ; \ClemAlex, \italique{Στρώματα} VI 26.4 Ἀριστοφάνης δὲ ὁ κωμικὸς ἐν ταῖς πρώταις Θεσμοφοριαζούσαις τὰ ἐκ τῶν Κρατίνου Ἐμπιπραμένων μετήνεγκεν ἔπη. \espTest \textbf{219-20} (χρῆσόν -- ξυροδόκης) Pollux X 140 (νῦν omisso) \espTest \textbf{220} (ξυροδόκης) Pollux II 32 ; \Eustathe 130.43, 1399.38 ; cf. Scholium \espTest \textbf{223} (ἀτταταῖ ἰατταταῖ) Suidas ε 2807 \espTest \textbf{227} Suidas η 334 \espTest (ἡμίκραιραν) Hesychius η 496 ἡμίκραιρα ; \EtymGen (AB) 150 ed. Miller ; \Eustathe 710.50, 1127.35 ; \AelDion η 9 \espTest \textbf{231} \EtymGen (AB) 218 ed. Miller \espTest (μῦ μῦ) Suidas ε 2807 ; cf. \Eustathe 723.9 \espTest (τί -- καλῶς) Suidas μ 1383 ; \Apostol XVI 60c \espTest (μύζεις) \Eustathe 440.23 \espTest \textbf{232} Suidas ψ 105 \espTest \textbf{235} Suidas κ 1756 Κλεισθένην ὁρῶ \espTest \textbf{237} Suidas δ 205 (sed γέγονα pro γενήσομαι) \espTest \textbf{244-5} (τὰ πλεῖστα -- ἀσβόλου) Photius α 2946 = \SynB α 2213 ; \AelDion α 186 \espTest \textbf{245} (φεῦ -- ἀσβόλου) Suidas α 4132 \espTest (ἄσβολος = ἀσβόλη) Moeris α 17 {\policeFr ≈} Hesychius α 7634 ; cf. Phrynicus \italique{Praep. soph} 28.1 (ἄσβολος fem.) \espTest \textbf{246} Suidas τ 903 (sed τὸν pro τὴν) \espTest (αἰθὸς γεγένημαι) \Eustathe 83.6, 992.60, 1314.15 \espTest (αἰθὸς) \Eustathe 1385.59 \espTest (τράμιν) Hesychius τ 1243 ; Photius τ 415 \espTest \textbf{248} (πλυνεῖ) Moeris π 79\par\vspace{2.5pt}
{\color{ruleGold}\hrule height 0.35pt}\vspace{2.5pt}
\noindent{\bfseries\scshape\color{rubricRed}Suite (en commentaire) de l'apparat critique et paléographique :} \textbf{Témoin papyrologique :} \textbf{\PsiUn{} (\italique{PSI} XI 1194)} couvre les \textbf{vv. 237–245} et comble au \textbf{v. 242} l'omission de [ἓ]τερον dans R. --- \textbf{v. 242 :} R avait omis un crétique avant πρὶν ἀντιλαβέσθαι ; le papyrus \PsiUn{} (lu par Medaglia) restitue \lecon{[ἓ]τερον}, tombé par saut visuel dans l'ancêtre de R. --- \textbf{v. 250 :} Exemple typique de \textbf{iota muet adscrit} dans \Run{} : τουτωί. --- \textbf{v. 217 :} διδόναι γ' αὐτὸν R (διδόναι γ' ἐμαυτὸν Scaliger) : \lecon{διδόναι γ' ἑαυτὸν} \lecon{Bibi} \italique{auctore Willems}. --- \textbf{v. 220 :} ξυροδίκης R (iotacisme/attraction de δίκη) : \lecon{ξυροδόκης} {\policeGr Σ}, Pollux II.32, X.140 (« étui à rasoir »). --- \textbf{v. 234 :} θεάσασθαι σεαυτόν R (glose pronominale non contracte) : \lecon{θεάσασθ' αὑτόν} Fritzsche (θεᾶσθαι σαυτόν Porson). --- \textbf{v. 242 :} \lecon{[ἓ]τερον} \PsiUn{} (restituit Medaglia) : ἕτερον om. R. --- \textbf{v. 248 :} οἰμώζετ' R : \lecon{οἰμώξετ'} Zanetti \espApp εἰς R (haplographie de ΤΙ dans ΕΙΤΙΣ) : \lecon{εἴ τις} Brunck.\par}
\end{minipage}\end{adjustbox}
\newpage
\renewcommand{\currHeadLE}{SCHOLIES MARGINALES ({\policeGr Σ}\expo{R}) \& GLOSES INTERLINÉAIRES · FOLIO 06/30}
\renewcommand{\currHeadRE}{f. 165r (Classense 167r) · vv. 208–251 · {\policeGr Σ} 082 à {\policeGr Σ} 094}
\renewcommand{\currHeadLO}{SCHOLIES MARGINALES ({\policeGr Σ}\expo{R}) \& GLOSES INTERLINÉAIRES · FOLIO 06/30}
\renewcommand{\currHeadRO}{f. 165r (Classense 167r) · vv. 208–251 · {\policeGr Σ} 082 à {\policeGr Σ} 094}
\bandeauScholies{Folio 06/30 — f. 165r}{Classense 167r}{vv. 208–251}{Recto (page de droite) · 13 entrées ({\policeGr Σ} 082 à {\policeGr Σ} 094)}
{\fontsize{8.8pt}{10.5pt}\selectfont\noindent\textbf{\color{rubricRed}Disposition topographique des marges sur le f. 165r :} Marge extérieure à droite ; marge intérieure à gauche. — \textit{Fac-similé de référence :} \texttt{folios\_individuels\_ff162v-177r/06\_f165r\_classense167r\_vv0208-0251.jpg}.\par\vspace{3pt}}
\entreeScholie{082}{214}{ἀπόδυθι τουτὶ θοἰμάτιον. — καὶ δὴ χαμαί·}{Interlinéaire (supra lineam, f. 165r, v. 214)}{ἀληθῶς.}{ἀληθ\expo{ῶς} {\policeFr ÷}}{« Vraiment (tout de bon, sérieusement). »}
\commScholie{Glose interlinéaire de R\expo{2} expliquant la particule intensive et d'acquiescement immédiat καὶ δὴ (« voici qui est fait en vérité, effectivement », ἀληθῶς) lorsque Mnésiloque laisse tomber son manteau à terre.}
\entreeScholie{083}{215}{ἀποξυρεῖν ταδί, / τὰ κάτω δ᾽ ἀφεύειν.}{Marge droite (extérieure, f. 165r, en face du v. 215)}{τὰ γένεια. ταῦτα δὲ ἔλαβεν ἐκ τῶν Ἰδαίων Κρατίνου.}{τὰ γένεια · ταῦτ(α) δ(ὲ) ἔλαβ(εν) ἐκ τ\expo{ῶν} Ἰδαί\expo{ων} Κρατίν\expo{ου} {\policeFr ÷}}{« τὰ γένεια (“le menton, la barbe” pour ταδί). Et il a emprunté cette scène d'épilation aux Idéens de Cratinos (fr. 90 K.-A.). »}
\commScholie{Précieuse scholie d'histoire littéraire (reprise par la Souda α 4565 et citée par L. Hauseux, vv 1-641.tex, l. 2063) qui nous apprend que la scène de rasage et de flambage de Mnésiloque s'inspirait des Idéens (Ἰδαῖοι) de Cratinos.}
\entreeScholie{084}{220}{ἐντεῦθεν ἐκ τῆς ξυροδόκης.}{Marge droite (extérieure, f. 165r, en face du v. 220, légèrement mutilée en fin de ligne)}{παρὰ τὸ δουροδόκης, ὡς πολλὰ αὐ[τῶν].}{παρὰ τ(ὸ) δουροδόκ(ης) ὡς πολλὰ αὐ[τῶν] {\policeFr ÷}}{« Forgé par parodie sur le modèle homérique δουροδόκη (“étui à lances”, Odyssée I, 128), comme beaucoup d'autres composés de ce genre. »}
\commScholie{Crucial pour le texte (L. Hauseux, vv 1-641.tex, ll. 2095-2096) : alors que R\expo{1} porte la coquille d'iotacisme ξυροδίκης, {\policeGr Σ}\expo{R} (ainsi que Pollux II, 32 et X, 140) garantit la forme authentique ξυροδόκης (« étui à rasoirs », sur δουροδόκη).}
\entreeScholie{085}{221}{γενναῖος εἶ· / κάθιζε [κάθιξε R]· φύσα τὴν γνάθον τὴν δεξιάν·}{Marge droite (extérieure, f. 165r, en face du v. 221)}{ἵνα ἀποτεταμένης αὐτῆς καλ[ῶς ξυρήσῃ].}{ἵνα ἀποτεταμέν\expo{ης} αὐτ\expo{ῆς} καλ[ῶς ξυρήσῃ] {\policeFr ÷}}{« (“Gonfle la joue droite”) : afin que, la joue étant bien tendue, il puisse la raser proprement. »}
\commScholie{Note pratique expliquant le geste des barbiers antiques qui faisaient gonfler la joue du client pour tendre la peau sous la lame.}
\entreeScholie{086}{224}{οὗτος σὺ ποῖ θεῖς; — εἰς τὸ τῶν σεμνῶν θεῶν·}{Marge droite (extérieure, f. 165r, en face du v. 224)}{ἀντὶ τοῦ εἰς τῶν Ἐρινύων. ἄσυλον γὰρ εἶχον οἱ καταλαμβάνοντες τὰ ἱερὰ τῶν Ἐρινύων.}{ἀν\auDessus{τ}{ι} τ\expo{οῦ} εἰς τ\expo{ῶν} Ἐρινύ\expo{ων} · ἄσυλ\expo{ον} γ(ὰρ) εἶχ\expo{ον} οἱ καταλαμβάνοντ\expo{ες} τὰ ἱερὰ τ\expo{ῶν} Ἐρινύ\expo{ων} {\policeFr ÷}}{« Au sens de : “au sanctuaire des Érinyes” (les Déesses Vénérables, au pied de l'Aréopage). Car ceux qui se réfugiaient dans le sanctuaire des Érinyes jouissaient du droit d'asile inviolable. »}
\commScholie{Cité par Louis Hauseux (vv 1-641.tex, l. 2131) et repris par la Souda σ 219 (cf. Cavaliers v. 1312).}
\entreeScholie{087}{227}{τὴν ἡμίκραιραν τὴν ἑτέραν ψιλὴν ἔχων;}{Marge droite (extérieure, f. 165r, en face des vv. 226–227)}{τὴν μίαν κομῶν γνάθον, τὴν ἑτέραν ἐψιλωμένος.}{τ(ὴν) μίαν κομ\expo{ῶν} γνάθ\expo{ον}, τ(ὴν) ἑτέρ(αν) ἐψιλωμέν\expo{ος} {\policeFr ÷}}{« Ayant encore une joue barbue et l'autre complètement rasée. »}
\commScholie{Glose expliquant τὰ ἡμίεκτα / ἡμίκραιρα (v. 226) : à moitié tondu d'un seul côté de la tête.}
\entreeScholie{088}{231}{μυμῦ. — τί μύζεις; πάντα πεποίηται καλῶς·}{Marge droite (extérieure, f. 165r, en face du v. 231)}{ὅτι τὸ μύζειν παρῆκται ἀπὸ τοῦ μῦ πολλοῖς ἄλλοις ὁμοίως. μύζειν δέ ἐστι τὸ τοῖς μυκτῆρσι ποιὸν ἦχον ἀποτελεῖν.}{ὅτι τ(ὸ) μύζ(ειν) παρῆκτ(αι) ἀπ(ὸ) τ\expo{οῦ} μῦ πολλ\expo{οῖς} ἄλλ\expo{οις} ὁμοί\expo{ως} · μύζ(ειν) δέ ἐστ(ι) τ(ὸ) τ\expo{οῖς} μυκτῆρσι ποι(ὸν) ἦχ\expo{ον} ἀποτελεῖν {\policeFr ÷}}{« Notez que le verbe μύζειν (“geindre à bouche fermée”) est dérivé de l'onomatopée μῦ, comme beaucoup d'autres verbes semblables (cf. γρῦ -> γρύζειν, κόκκυ -> κοκκύζειν). Et μύζειν consiste à produire un certain son sourd par les narines. »}
\commScholie{Observation grammaticale et phonétique d'origine alexandrine (reprise par la Souda μ 1380 et l'Etymologicum Magnum).}
\entreeScholie{089}{232}{οἴμοι κακοδαίμων, ψιλὸς αὖ στρατεύσομαι·}{Marge droite (extérieure, f. 165r, bloc de 5 lignes en face du v. 232)}{ψιλὴ τάξις ἐστὶν ἀγεννεστάτη, ἥτις βάλλεται ἔμπροσθεν τοῦ στρατοπέδου γυμνὴ οὖσα ὅπλων, καὶ ἐὰν ταραχὴ γένηται, ἐκείνη πρώτη ἀναιρεῖται. τοῦτο δὲ λέγει, ἐπεὶ ἐψίλωται.}{ψιλὴ τάξις ἐστ(ὶν) ἀγεννεστάτη, ἥτις βάλλετ(αι) ἔμπροσθ(εν) τ\expo{οῦ} στρατοπέδου γυμνὴ οὖσα ὅπλ\expo{ων}, \textbf{ς} ἐὰν ταραχὴ γένητ(αι), ἐκείνη πρώτη ἀναιρεῖτ(αι) · τοῦτ(ο) δ(ὲ) λέγει, ἐπεὶ ἐψίλωτ(αι) {\policeFr ÷}}{« Le corps des troupes légères (ψιλὴ τάξις, vélites sans armure lourde) est le moins estimé : on le place en avant de l'armée sans armes défensives, et s'il survient une mêlée, il est le premier massacré. Il dit cela par calembour, parce qu'il vient d'être rasé à nu (ἐψίλωται). »}
\commScholie{Explique le calembour militaire sur ψιλός : 1° rasé, glabre ; 2° fantassin léger sans hoplon.}
\entreeScholie{090}{235}{ὁρᾷς σεαυτόν; — οὐ μὰ Δί᾽ ἀλλὰ Κλεισθένη·}{Interlinéaire / marge intérieure gauche (f. 165r, v. 235)}{διαβάλλει τοῦτον ὡς ξυρούμενον.}{διαβάλλ(ει) τοῦτ\expo{ον} ὡς ξυρούμεν\expo{ον} {\policeFr ÷}}{« Il raille ce personnage (Clisthène) comme se rasant les joues (comme un efféminé). »}
\commScholie{Prépare l'entrée en scène effective de Clisthène au v. 574 (cf. Σ 575).}
\entreeScholie{091}{236}{οἴμοι κακοδαίμων, δελφάκιον γενήσομαι·}{Marge droite (extérieure, f. 165r, en face du v. 236)}{μετὰ γὰρ τὸ τυθῆναι τὰ δελφάκια φλογίζονται, ἵνα ψιλωθῶσιν.}{μετὰ γ(ὰρ) τ(ὸ) τυθῆν(αι) τὰ δελφάκια φλογίζοντ(αι), ἵνα ψιλωθῶσ(ιν) {\policeFr ÷}}{« Parce qu'après avoir été égorgés en sacrifice, les cochons de lait (δελφάκια) sont passés à la flamme (grillés au flambeau) afin qu'on leur brûle les soies. »}
\commScholie{Repris par la Souda δ 209 : explique pourquoi Mnésiloque, voyant Euripide approcher une torche allumée pour lui flamber les poils du derrière, s'écrie qu'on va le traiter en cochon de lait.}
\entreeScholie{092}{242}{πρὶν ἀντιλαβέσθαι πρωκτὸν [ἕ]τερον τῆς φλογός·}{Interlinéaire (supra lineam au-dessus de πρωκτὸν, f. 165r)}{ὡς εἰ ἔλεγεν οἰκίαν.}{ὡς εἰ ἔλεγ(εν) οἰκί(αν) {\policeFr ÷}}{« Comme s'il avait dit “la maison” (avant que le feu ne prenne à la maison -> au derrière, παρ' ὑπόνοιαν). »}
\commScholie{Cité par Louis Hauseux (vv 1-641.tex, l. 2267 n. 1) : l'appel au secours avec des seaux d'eau est celui qu'on lance lors d'un incendie domestique, mais οἰκίαν est remplacé παρὰ προσδοκίαν par πρωκτόν.}
\entreeScholie{093}{245}{αἰθὸς γεγένημαι πάντα τὰ περὶ τὴν τράμιν.}{Marge droite (extérieure, f. 165r, en face du v. 245)}{ἄσβολος καλεῖται ἡ αἰθάλη τοῦ πυρός.}{ἄσβολ\expo{ος} καλεῖτ(αι) ἡ αἰθάλ(η) τ\expo{οῦ} πυρός {\policeFr ÷}}{« On appelle suie (ἄσβολος, cf. ἐσβόλη v. 244 / ἀσβόλη) la fumée noire laissée par le feu. »}
\commScholie{Glose lexicographique se rapportant à la fin du v. 244 / v. 245.}
\entreeScholie{094}{246}{τὰ περὶ τὴν τράμιν.}{Interlinéaire (supra lineam au-dessus de τράμιν, f. 165r)}{τράμις ὁ πρωκτός.}{τράμις ὁ πρωκτός {\policeFr ÷}}{« τράμις : le périnée, le derrière (la raie des fesses). »}
\commScholie{Repris par la Souda τ 899, Photius τ 433 et Hésychius τ 1247 (τράμις· ὁ ὄρρος).}
\clearpage
\renewcommand{\currHeadLE}{ΑΡΙΣΤΟΦΑΝΟΥΣ ΘΕΣΜΟΦΟΡΙΑΖΟΥΣΑΙ · CODEX RAVENNAS 429 (\R) · FOLIO 07/30}
\renewcommand{\currHeadRE}{f. 165v (Classense 167v) · vv. 252–294 · Collation diplomatique}
\renewcommand{\currHeadLO}{ΑΡΙΣΤΟΦΑΝΟΥΣ ΘΕΣΜΟΦΟΡΙΑΖΟΥΣΑΙ · CODEX RAVENNAS 429 (\R) · FOLIO 07/30}
\renewcommand{\currHeadRO}{f. 165v (Classense 167v) · vv. 252–294 · Collation diplomatique}
\noindent\begin{adjustbox}{max totalheight=30.4cm, max width=\textwidth, keepaspectratio}\begin{minipage}{\textwidth}
\bandeauRav{f. 165v}{Classense 167v}{252–294}{Folio 07/30 · Verso (page de gauche) · 17 scholies \& gloses ({\policeGr Σ} 095–111)}
{\fontsize{8.8pt}{10.2pt}\selectfont
\ligneRav{252}{ἀγάθ/}{λαμβάνετε καὶ χρῆσθ᾽· οὐ φθονῶ. : τί οὖν λάβω;}{\textbf{252} (Εὐρ.) R : \lecon{(Μνη.)} Bentley}
\ligneRav{253}{\auDessus{εὐ}{ρ}}{ὅ τι; τὸν κροκωτὸν πρῶτον ἐνδύου λαβών·}{\textbf{253} (Ἀγάθ.) R : \lecon{(Εὐρ.)} Dobree}
\ligneRav{254}{μ\auDessus{ν}{Η}}{νὴ τὴν Ἀφροδίτην ἡδύ γ᾽ ὄζει ποσθίου·}{\textbf{254} (Εὐρ.) R}
\ligneRav{255}{ἀγάθ/}{σύζωσον ἀνύσας. : αἶρε νῦν στρόφιον. : ἰδού·}{\textbf{255} (Ἀγάθ.) et dicolon ante αἶρε praefixit R \espApp (Εὐρ.) ἰδού — (Μνη.) ἴθι ... σκέλη R}
\ligneRav{256}{μ\auDessus{ν}{Η}}{ἴθι νυν κατάστειλόν με τὰ περὶ τὼ \difference{σκέλει}·}{\textbf{256} σκέλη R : \lecon{σκέλει} van Leeuwen}
\ligneRav{}{\auDessus{εὐ}{ρ}}{κεκρυφάλου δεῖ καὶ μίτρας. : ἡδὶ μὲν οὖν}{}
\ligneRav{}{}{κεφαλὴ περίθετος, ἣν ἐγὼ νύκτωρ φορῶ·}{}
\ligneRav{}{\auDessus{εὐ}{ρ}}{νὴ τὸν Δί᾽ ἀλλὰ κἀπιτηδεία πάνυ·}{}
\ligneRav{260}{μ\auDessus{ν}{Η}}{\difference{ἆρ᾽ ἁρμόσει} μοι; : νὴ Δί᾽ ἀλλ᾽ ἄριστ᾽ ἔχει·}{\textbf{260} ἠραρμοσηι R : \lecon{ἆρ᾽ ἁρμόσει} Küster \espApp (Ἀγάθ.) R : \lecon{(Εὐρ.)} Ellebodius}
\ligneRav{261}{—}{φέρ᾽ ἔγκυκλον· : \difference{τουτὶ} \difference{λάβ᾽} ἀπὸ τῆς κλινίδος·}{\textbf{261} τοῦτὶ R : \lecon{τουτὶ} Suid. ε 135, κ 1809 \espApp λάμβαν᾽ R, Suid. : \lecon{λάβ᾽} Daubuz, Küster}
\ligneRav{}{\auDessus{εὐ}{ρ}}{ὑποδημάτων δεῖ. : τἀμὰ ταυτὶ λάμβανε·}{}
\ligneRav{263}{μ\auDessus{ν}{Η}}{ἆρ᾽ ἁρμόσει \difference{μοι;} : \auDessus{ἀγά}{θ} χαλαρὰ γ᾽ ο\auDessus{ῦ}{ν} \difference{χαίροις} φορῶν·}{\textbf{263} μοι R : \lecon{μοι;} Küster \espApp \lecon{(Ἀγάθ.)} χαλαρὰ ... φορῶν R : (Μνη.) Brunck, edd. post. \espApp γ᾽ ο\auDessus{ῦ}{ν} \Rpc \espApp χαίρεις R, Suid. AGM : \lecon{χαίροις} Suid. FV}
\ligneRav{264}{μ\auDessus{ν}{Η}}{σὺ τοῦτο γίγνωσκ᾽· \difference{:} ἀλλ᾽ ἔχεις γὰρ ὧν δέει,}{\textbf{264} nulla nota mut. pers. R : \lecon{(Ἀγάθ.)} ἀλλ᾽ ... Bibi}
\ligneRav{265}{}{εἴσω τις ὡς τάχιστά μ᾽ ἐσκυκλησάτω·}{}
\ligneRav{266}{\auDessus{εὐ}{ρ}}{\difference{ἁνὴρ} μὲν ἡμῖν οὑτοσὶ καὶ δὴ γυνὴ}{\textbf{266} ἀνὴρ R : \lecon{ἁνὴρ} Fritzsche}
\ligneRav{267}{}{\difference{τό γ᾽} εἶδος· ἢν λαλῇς δ᾽, ὅπως τῷ φθέγματι}{\textbf{267} τόδ᾽ R : \lecon{τό γ᾽} Küster}
\ligneRav{}{—}{γυναικιεῖς εὖ καὶ πιθανῶς. : πειράσομαι·}{}
\ligneRav{}{\auDessus{εὐ}{ρ}}{βάδιζε τοίνυν. : μὰ τὸν Ἀπόλλω οὔκ, ἤν γε μὴ}{}
\ligneRav{270}{—}{ὀμόσῃς ἐμοί— : τί χρῆμα; : συσσώσειν ἐμὲ}{}
\ligneRav{}{}{πάσαις τέχναις, ἤν μοί τι περιπίπτῃ κακόν·}{}
\ligneRav{}{\auDessus{εὐ}{ρ}}{`ὄμνυμι τοίνυν αἰθέρ᾽ οἴκησιν Διός.'}{}
\ligneRav{273}{μ\auDessus{ν}{Η}}{τί μᾶλλον ἢ τὴν Ὑποκράτους ξυνοικίαν;}{\textbf{273} \lecon{Ὑποκράτους} R : Ἱπποκράτους Σ}
\ligneRav{}{\auDessus{εὐ}{ρ}}{ὄμνυμι τοίνυν πάντας ἄρδην τοὺς θεούς·}{}
\ligneRav{275}{μ\auDessus{ν}{Η}}{μέμνησο τοίνυν ταῦθ᾽, ὅτι ἡ φρὴν ὤμοσεν,}{}
\ligneRav{}{}{ἡ γλῶττα δ᾽ οὐκ ὀμώμοκ᾽· οὐδ᾽ ὥρκωσ᾽ ἐγώ·}{}
\ligneRav{}{}{\hspace{2.2em}\textbf{[παρεπιγραφή R\expo{1}, \PsiUn{} :]} \lecon{ὀλολύζουσι τε ἱερὸν ὠθεῖται}}{\textbf{post 276} \lecon{ὀλολύζουσι τε ἱερὸν ὠθεῖται} R, ολολυ[ \PsiUn, {\policeGr Σ}\expo{R}}
\ligneRav{277}{\auDessus{εὐ}{ρ}}{ἔκσπευδε ταχέως· ὡς τὸ τῆς ἐκκλησίας}{\textbf{277} \lecon{ἔκσπευδε} R : εἄ· σπεῦδε Maas}
\ligneRav{278}{}{σημεῖον ἐν τῷ Θεσμοφορίῳ φαίνεται·}{\textbf{278} \lecon{Θεσμοφορίῳ} R, cf. \italique{IG} II\expo{2} 1177.24 ; 2498.12 : Θεσμοφόρῳ Suid. σ 325 : Θεσμοφορείῳ Scaliger}
\ligneRav{}{—}{ἐγὼ δ᾽ ἄπειμι. : δεῦρό νυν ὦ Θρᾷτθ᾽ ἕπου·}{}
\ligneRav{280}{}{ὦ Θρᾷττα θέασαι, καομένων τῶν λαμπάδων}{}
\ligneRav{}{}{ὅσον τὸ χρῆμ᾽ ἀνέρχεθ᾽ ὑπὸ τῆς λιγνύος·}{}
\ligneRav{282}{}{ἀλλ᾽, ὦ \difference{περικαλλεῖ} \difference{Θεσμοφόρω}, δέξασθέ με}{\textbf{282} Θεσμοφόρῳ R : \lecon{Θεσμοφόρω} Photius α 1386 \espApp περὶ καλλῆ R : \lecon{περικαλλεῖ} Kaibel}
\ligneRav{283}{}{ἀγαθῇ τύχῃ \difference{καὶ δεῦρο καὶ πάλιν} οἴκαδε·}{\textbf{283} καὶ δεῦρο πάλιν R : \lecon{καὶ δεῦρο καὶ πάλιν} Bentley}
\ligneRav{}{}{ὦ Θρᾷττα τὴν κίστην κάθελε, κᾆτ᾽ ἔξελε}{}
\ligneRav{285}{}{\difference{τὰ πόπαν᾽}, ὅπως λαβοῦσα θύσω ταῖν θεαῖν·}{\textbf{285} τὸ πόπανον R : \lecon{τὰ πόπαν᾽} \PsiUn, Porson}
\ligneRav{}{}{δέσποινα πολυτίμητε Δήμητερ φίλη}{}
\ligneRav{287}{}{καὶ \difference{Φερρέφαττα}, πολλὰ πολλάκις μέ σοι}{\textbf{287} \lecon{Φερρέφαττα} \PsiUn, Thiersch : Φερσέφαττα R}
\ligneRav{}{}{θύειν ἔχουσαν, εἰ δὲ μή, ἀλλὰ νῦν λαθεῖν}{}
\ligneRav{289}{}{καὶ τὴν θυγατέρα \difference{Χοιρίλ'} ἀνδρός μοι τυχεῖν}{\textbf{289} χοῖρον R : χοιρίον Σ : \lecon{Χοιρίλ'} Bibi}
\ligneRav{290}{}{πλουτοῦντος, ἄλλως τ᾽ ἠλιθίου κἀβελτέρου,}{}
\ligneRav{291}{}{καὶ \difference{πρὸς θυλάκιον} νοῦν ἔχειν μοι καὶ φρένας·}{\textbf{291} πρὸς θάληκον R : Ποσθαλίσκον Dindorf : \lecon{πρὸς θυλάκιον} Bibi}
\ligneRav{}{}{ποῦ ποῦ καθίζωμ᾽ ἐν καλῷ, τῶν ῥητόρων}{}
\ligneRav{}{}{ἵν᾽ ἐξακούω; σὺ δ᾽ ἄπιθ᾽ ὦ Θρᾷττ᾽ ἐκποδών·}{}
\ligneRav{}{}{δούλοις γὰρ οὐκ ἔξεστ᾽ ἀκούειν τῶν λόγων·}{}
}
\vspace{3pt}{\color{ruleGold}\hrule height 0.6pt}\vspace{3pt}
{\fontsize{8.0pt}{9.4pt}\selectfont
\noindent{\bfseries\scshape\color{rubricRed}Loci paralleli et testimonia (f. 165v, vv. 252–294) :} \textbf{254} Suidas π 2114 \espTest (ποσθίου) Photius π 1113 πόσθιον· αἰδοῖον· ἢ (sic) τὴν Ἀφροδίτην ; Hesychius π 3096 \espTest \textbf{257-8} (ἡδὶ -- φορῶ) Suidas π 1165 (sed ἥδιον pro ἡδὶ) \espTest \textbf{258} Photius π 690 περίθετος κεφαλή ; \Eustathe 1280.63 \espTest (κεφαλὴ περίθετος) Pollux II 35 \espTest \textbf{261} Suidas ε 135 (sed τουτὶ pro τι) \espTest (ἔγκυκλον -- κλινίδος) Suidas κ 1809 \espTest \textbf{263} Suidas χ 7 χαλαρὰ γοῦν χαίρεις φορῶν ὑποδήματα (sed χαίροις codd. FV) \espTest \textbf{274} (ἄρδην) Hesychius α 7097 \espTest \textbf{275-6} cf. \Eustathe 443.37 \espTest \textbf{277-8} (ὡς -- φαίνεται) Suidas σ 325 (sed Θεσμοφόρῳ pro Θεσμοφορίῳ) \espTest \textbf{282} Photius α 1386 ; \AelDion α 113 \espTest \textbf{285} (πόπαν᾽) \Eustathe 437.4 \espTest \textbf{289} (χοῖρον) Euripides, testimonia 4, 18 Kannicht Χοιρίλη \espTest \textbf{291} (θάληκον) Hesychius θ 50 θάλλικα ; Pollux X 152 θυλάκιον [...] ἐπὶ τοῦ αὐτοῦ δὲ καὶ σακίον, ὅταν φῇ ἐν Θεσμοφοριαζούσαις \espTest \textbf{292} (τῶν ῥητόρων) Photius ρ 106 ῥήτορα· καὶ τὴν γυναῖκα. Ἀριστοφάνης\par\vspace{2.5pt}
{\color{ruleGold}\hrule height 0.35pt}\vspace{2.5pt}
\noindent{\bfseries\scshape\color{rubricRed}Suite (en commentaire) de l'apparat critique et paléographique :} \textbf{Témoin papyrologique :} \textbf{\PsiUn{} (\italique{PSI} XI 1194)} couvre les \textbf{vv. 273–288}. --- \textbf{Troisième \italique{parepigraphè} (post v. 276) :} R écrit dans la colonne \lecon{ὀλολύζουσι τε ἱερὸν ὠθεῖται}, avec une \textbf{\italique{koronis}} dans la marge gauche et la scholie {\policeGr Σ}\expo{R} : παρεπιγραφή. ἐκκυκλεῖται ἐπὶ τὸ ἔξω τὸ θεσμοφόριον. Le papyrus \PsiUn{} confirme ολολυ[. --- \textbf{Abréviation supralinéaire au v. 263 :} γ' ο\auDessus{ῦ}{ν} \Rpc{}. --- \textbf{Deux conjectures majeures de \textbf{Bibi} (vv. 289 et 291) :} Au v. 289, \lecon{Χοιρίλ'} (nom de la fille de Mnésiloque et épouse d'Euripide, corrompu en χοῖρον dans R et χοιρίον dans {\policeGr Σ}\expo{R}) ; au v. 291, \lecon{πρὸς θυλάκιον} (ou πρὸς θάλικα ap. Hésychius θ 50) pour πρὸς θάληκον de R. --- \textbf{v. 260 :} ἠραρμοσηι R (iotacisme et médivision) : \lecon{ἆρ' ἁρμόσει} Küster \espApp (Ἀγάθ.) R : \lecon{(Εὐρ.)} Ellebodius. --- \textbf{v. 263 :} μοι R : \lecon{μοι;} Küster \espApp \textbf{(Ἀγάθ.) χαλαρὰ ... φορῶν} R (maintenu par \textbf{Bibi} contre (Μνη.) Brunck) \espApp \lecon{γ' ο\auDessus{ῦ}{ν}} \Rpc{} \espApp χαίρεις R, Souda AGM : \lecon{χαίροις} Souda FV. --- \textbf{v. 264 :} nulla nota mut. pers. R : \lecon{(Ἀγάθ.)} ἀλλ' ... \lecon{Bibi}. --- \textbf{v. 273 :} \lecon{Ὑποκράτους} R (maintenu par \textbf{Bibi} : calembour comique sur ὑπό, contre la normalisation Ἱπποκράτους de {\policeGr Σ}\expo{R}). --- \textbf{v. 289 :} χοῖρον R (χοιρίον {\policeGr Σ}) : \lecon{Χοιρίλ'} (ou Χοιρίν') \lecon{Bibi} (fille de Mnésiloque et épouse d'Euripide). --- \textbf{v. 291 :} πρὸς θάληκον R (Ποσθαλίσκον Dindorf) : \lecon{πρὸς θυλάκιον} (ou πρὸς θάλικα ap. Hésychius θ 50) \lecon{Bibi} (par métathèse).\par}
\end{minipage}\end{adjustbox}
\newpage
\renewcommand{\currHeadLE}{SCHOLIES MARGINALES ({\policeGr Σ}\expo{R}) \& GLOSES INTERLINÉAIRES · FOLIO 07/30}
\renewcommand{\currHeadRE}{f. 165v (Classense 167v) · vv. 252–294 · {\policeGr Σ} 095 à {\policeGr Σ} 111}
\renewcommand{\currHeadLO}{SCHOLIES MARGINALES ({\policeGr Σ}\expo{R}) \& GLOSES INTERLINÉAIRES · FOLIO 07/30}
\renewcommand{\currHeadRO}{f. 165v (Classense 167v) · vv. 252–294 · {\policeGr Σ} 095 à {\policeGr Σ} 111}
\bandeauScholies{Folio 07/30 — f. 165v}{Classense 167v}{vv. 252–294}{Verso (page de gauche) · 17 entrées ({\policeGr Σ} 095 à {\policeGr Σ} 111)}
{\fontsize{8.8pt}{10.5pt}\selectfont\noindent\textbf{\color{rubricRed}Disposition topographique des marges sur le f. 165v :} Marge extérieure à gauche ; marge intérieure à droite. — \textit{Fac-similé de référence :} \texttt{folios\_individuels\_ff162v-177r/07\_f165v\_classense167v\_vv0252-0294.jpg}.\par\vspace{3pt}}
\entreeScholie{095}{254}{νὴ τὴν Ἀφροδίτην ἡδύ γ᾽ ὄζει ποσθίου·}{Marge gauche (extérieure, f. 165v, en haut du verso, en face du v. 254)}{δέον εἰπεῖν μύρου εἶπε ποσθίου. πόσθιον δέ ἐστι τὸ αἰδοῖον τοῦ ἀνδρός. διαβάλλει δὲ αὐτὸν ὡς μετ’ αὐτοῦ ἑταιρίζοντα.}{δέον εἰπ(εῖν) μύρ\expo{ου} εἶπ(ε) ποσθί\expo{ου} · πόσθι\expo{ον} δέ ἐστ(ι) τ(ὸ) αἰδοῖ\expo{ον} τ\expo{οῦ} ἀνδρός · διαβάλλ(ει) δ(ὲ) αὐτ(ὸν) ὡς μετ' αὐτ\expo{οῦ} ἑταιρίζοντ(α) {\policeFr ÷}}{« Alors qu'on attendait “il sent bon le parfum” (μύρου), il a dit (παρ' ὑπόνοιαν) “il sent la petite verge” (ποσθίου). Le πόσθιον est le membre viril ; il raille ainsi Agathon comme se prostituant aux hommes. »}
\commScholie{Repris par la Souda π 2116. Explique l'effet de surprise comique sur le parfum de la robe jaune safran (κροκωτός) prêtée par Agathon.}
\entreeScholie{096}{258}{κεφαλὴ περίθετος, ἣν ἐγὼ νύκτωρ φορῶ.}{Marge gauche (extérieure, f. 165v, bloc de 5 lignes en face du v. 258)}{ὅτι τὴν περίθετον οὕτως. καὶ τοῦτο δὲ εἰς μαλακίαν, ἵνα λανθάνῃ νυκτὸς καὶ ὡς γυνὴ πάσχῃ. λέγουσι δὲ περίθετον πρόσωπον, ὡς Ἀριστομένης ἐν Γόησιν «παντευχίαν δὲ τοῦ θεοῦ ταύτην λαβεῖν καὶ περίθετον πρόσωπον, ὃ λαβὼν ἕσταθι.»}{ὅτι τ(ὴν) περίθετ(ον) οὕτ\expo{ως} · \textbf{ς} τοῦτ(ο) δ(ὲ) εἰς μαλακίαν, ἵνα λανθάνῃ νυκτὸς \textbf{ς} ὡς γυνὴ πάσχῃ · λέγουσι δ(ὲ) περίθετ(ον) πρόσωπ(ον), ὡς Ἀριστομένης ἐν Γόησ(ιν) «παντευχίαν δ(ὲ) τ\expo{οῦ} θε\expo{οῦ} ταύτην λαβ(εῖν) \textbf{ς} περίθετ(ον) πρόσωπ(ον), ὃ λαβὼν ἕσταθι · » {\policeFr ÷}}{« Notez qu'il appelle ainsi la perruque (κεφαλὴ περίθετος, “coiffure postiche”). Ceci encore vise sa mollesse efféminée, afin qu'il ne soit pas reconnu la nuit et subisse l'amour comme une femme. On dit de même περίθετον πρόσωπον (“masque postiche”), comme Aristomène dans ses Charlatans (Γόητες, fr. 5 K.-A.) : “Prends cette panoplie du dieu et ce masque postiche, et une fois que tu l'auras pris, tiens-toi debout.” »}
\commScholie{Scholie alexandrine précieuse (citée par L. Hauseux, vv 1-641.tex, l. 2420) qui nous a conservé un fragment de deux trimètres du poète de l'Ancienne Comédie Aristomène (contemporain d'Aristophane).}
\entreeScholie{097}{259}{νὴ τὸν Δί᾽ ἀλλὰ κἀπιτηδεία πάνυ·}{Interlinéaire (supra lineam, f. 165v, v. 259)}{πάντα ἔχουσα τὰ τοιαῦτα ἐν ἑαυτῇ.}{πάντα ἔχουσα τὰ τοιαῦτ(α) ἐν ἑαυτῇ {\policeFr ÷}}{« Possédant en lui-même (ou dans sa garde-robe) tous les accessoires de ce genre. »}
\commScholie{Glose interlinéaire de R\expo{2} expliquant l'adjectif κἀπιτηδεία (καὶ ἐπιτηδεία, « tout à fait appropriée ») appliqué par Mnésiloque à la coiffe nocturne d'Agathon, parce qu'elle réunit déjà en un seul accessoire (ἐν ἑαυτῇ) la perruque, la résille (κεκρύφαλος) et la mitre.}
\entreeScholie{098}{261}{φέρ᾽ ἔγκυκλον· / λάβ᾽ ἀπὸ τῆς κλινίδος.}{Marge gauche (extérieure, f. 165v, en face du v. 261)}{τῆς κλίνης. στρῶμα δίδωσιν αὐτῷ. βάθριόν τι ἡ κλινὶς εἰστεταμένον. δῆλον δὲ ὅτι τὸ ἔγκυκλον ἱμάτιον, ὁ δὲ κροκωτὸς ἔνδυμα.}{τ\expo{ῆς} κλίν\expo{ης} · στρῶμα δίδωσ(ιν) αὐτῷ · βάθριόν τι ἡ κλινὶς εἰστεταμέν\expo{ον} · δῆλ\expo{ον} δ(ὲ) ὅτι τ(ὸ) ἔγκυκλ\expo{ον} ἱμάτι\expo{ον}, ὁ δ(ὲ) κροκωτ\expo{ὸς} ἔνδυμα {\policeFr ÷}}{« κλινίδος : de la couchette (du lit de repos sur l'ekkyklêma) ; il lui tend un manteau de dessus. La κλινίς est une sorte de banquette tendue de sangles. Et il est clair d'après ce passage que l'ἔγκυκλον est un manteau de dessus (drapé autour du corps), tandis que le κροκωτός est une tunique de dessous (ἔνδυμα). »}
\commScholie{Observation vestimentaire très exacte (discutée par L. Hauseux, vv 642-fin-bis.tex, l. 558, et reprise par la Souda ε 126) distinguant la tunique safran (κροκωτός) du manteau circulaire féminin (ἔγκυκλον).}
\entreeScholie{099}{263}{ἆρ᾽ ἁρμόσει μοι; — χαλαρὰ γ᾽ οὖν χαίροις [χαίρεις R] φορῶν·}{Marge gauche (extérieure, f. 165v, en face du v. 263)}{χαῦνα. διαβάλλει δὲ πάλιν τὸν Ἀγάθωνα ὡς χαῦνον.}{χαῦνα · διαβάλλ(ει) δ(ὲ) πάλ(ιν) τ(ὸν) Ἀγάθωνα ὡς χαῦν\expo{ον} {\policeFr ÷}}{« χαλαρά (“lâches, larges”) : molles, béantes (χαῦνα). Il raille de nouveau Agathon comme étant relâché (χαῦνος, aux mœurs dissolues). »}
\commScholie{Louis Hauseux (vv 1-641.tex, ll. 2461-2468) critique avec esprit le contresens des éditeurs modernes qui ont suivi aveuglément l'allusion obscène supposée par le scholiaste, alors que Mnésiloque plaisante sur la pointure des chaussures féminines d'Agathon.}
\entreeScholie{100}{268}{ἀνὴρ μὲν ἡμῖν οὑτοσὶ καὶ δὴ γυνὴ / τό γ᾽ εἶδος· ἢν λαλῇς δ᾽, ὅπως τῷ φθέγματι / γυναικιεῖς εὖ καὶ πιθανῶς.}{Interlinéaire (supra lineam au-dessus de γυναικιεῖς, f. 165v)}{ὡς γυνὴ λαλήσεις.}{ὡς γυνὴ λαλήσ(εις) {\policeFr ÷}}{« Tu parleras en imitant la voix d'une femme. »}
\commScholie{Glose interlinéaire de R\expo{2} expliquant le futur attique contracte γυναικιεῖς (de γυναικίζω, « imiter la voix et les manières d'une femme ») par ὡς γυνὴ λαλήσεις (« tu parleras comme une femme »).}
\entreeScholie{101}{272}{ὄμνυμι τοίνυν αἰθέρ᾽ οἴκησιν Διός.}{Marge gauche (extérieure, f. 165v, en face du v. 272)}{ἐκ τῆς Μελανίππης Εὐριπίδου.}{ἐκ τ\expo{ῆς} Μελανίππ\expo{ης} Εὐριπίδ\expo{ου} {\policeFr ÷}}{« Ce vers est tiré de la Mélanippe (la Sage) d'Euripide (fr. 487 Kannicht). »}
\commScholie{Cité par Louis Hauseux (vv 1-641.tex, l. 2537) : la même parodie du célèbre serment cosmologique de Mélanippe revient dans les Grenouilles (vv. 100 et 311).}
\entreeScholie{102}{273}{τί μᾶλλον ἢ τὴν Ὑποκράτους ξυνοικίαν;}{Marge gauche (extérieure, f. 165v, en face du v. 273)}{Ἱπποκράτης ἐγένετο Ἀθηναῖος, ὃς εἶχεν υἱοὺς [...] διεβάλλετο δὲ σὺν αὐτοῖς ὡς ἀπαίδευτος.}{Ἱπποκράτ(ης) ἐγένετ(ο) Ἀθηναῖ\expo{ος}, ὃς εἶχ(εν) υἱ\expo{οὺς} [τρεῖς · Τελέσιππ\expo{ον}, Δημοφῶντ(α), Περικλέα] · διεβάλλετ(ο) δ(ὲ) σὺν αὐτ\expo{οῖς} ὡς ἀπαίδευτ\expo{ος} {\policeFr ÷}}{« Hippocrate était un Athénien (neveu de Périclès et stratège) qui avait trois fils [Télésippos, Démophon, Périclès], et on le raillait avec eux pour sa grossièreté et sa bêtise (litt. comme des “cochons”, ὑο-κράτους). »}
\commScholie{Analyse décisive de Louis Hauseux (Bibi, v. 273, vv 1-641.tex, ll. 2544-2565) : R\expo{1} a conservé le calembour authentique d'Aristophane Ὑποκράτους (forgé sur ὗς, « porc », + Ἱπποκράτης, les fils d'Hippocrate étant surnommés les « pourceaux », cf. Nuées v. 1001 et Photius σ 845), tandis que le scholiaste normalise en Ἱπποκράτους.}
\entreeScholie{103}{274}{ὄμνυμι τοίνυν πάντας ἄρδην τοὺς θεούς.}{Marge gauche (extérieure, f. 165v, en face du v. 274)}{ὅτι κατὰ δύναμιν τῆς λέξεως, ἀντὶ τοῦ ἀθρόους πάντας. τὸ γὰρ ἄρδην ὅμοιόν ἐστι τῷ φοράδην ἐνέγκαι· ἡ γὰρ ἄρσις τοῦτο δηλοῖ. εἰ μὴ νῦν τοῦτο βούλεται, πάντας ὁμοῦ αἴρω.}{ὅτι κατὰ δύναμ(ιν) τ\expo{ῆς} λέξ(εως), ἀν\auDessus{τ}{ι} τ\expo{οῦ} ἀθρό\expo{ους} πάντ\expo{ας} · τ(ὸ) γ(ὰρ) ἄρδ(ην) ὅμοιόν ἐστ(ι) τῷ φοράδ(ην) ἐνέγκ(αι) · ἡ γ(ὰρ) ἄρσ(ις) τοῦτ(ο) δηλοῖ · εἰ μὴ νῦν τοῦτ(ο) βούλετ(αι), πάντ\expo{ας} ὁμοῦ αἴρω {\policeFr ÷}}{« Notez que le mot ἄρδην est employé selon sa force étymologique, au sens de “tous ensemble d'un seul bloc” (ἀθρόους πάντας) : car ἄρδην est analogue à φοράδην ἐνέγκαι (“emporter d'une traite”), ce qu'indique le nom verbal ἄρσις (“action de soulever”). À moins qu'il ne veuille dire ici : “je les enlève tous d'un coup”. »}
\commScholie{Fine remarque grammaticale et étymologique sur l'adverbe ἄρδην (dérivé de αἴρω).}
\entreeScholie{104}{276}{μέμνησο τοίνυν ταῦθ᾽, ὅτι ἡ φρὴν ὤμοσεν, / ἡ γλῶττα δ᾽ οὐκ ὀμώμοκ᾽·}{Marge gauche (extérieure, f. 165v, en face des vv. 275–276)}{ἐπειδὴ ἐν Ἱππολύτῳ [608] τοῦτο εἶπεν Εὐριπίδης, νῦν εὐλαβούμενος λέγει ὅτι ἡ φρήν σου ὀμώμοκε καὶ οὐχ ἡ γλῶττα.}{ἐπειδὴ ἐν Ἱππολύτ(ῳ) τοῦτ(ο) εἶπ(εν) Εὐριπίδ(ης), νῦν εὐλαβούμεν\expo{ος} λέγ(ει) ὅτι ἡ φρήν σ\expo{ου} ὀμώμοκ(ε) \textbf{ς} οὐχ ἡ γλῶττα {\policeFr ÷}}{« Puisque dans l'Hippolyte (v. 612) Euripide avait fait dire : “Ma langue a juré, mais mon esprit n'a pas juré”, ici Mnésiloque, prenant ses précautions, inverse la formule et dit : “C'est ton esprit qui a juré, et non ta langue !” »}
\commScholie{Le scholiaste souligne le retournement comique du vers le plus scandaleux de l'Hippolyte d'Euripide (v. 612, parodié aussi dans les Grenouilles vv. 101-102 et 1471).}
\entreeScholie{105}{277}{[παρεπιγραφή R\expo{1}, \PsiUn{} :] ὀλολύζουσί τε· ἱερὸν ὠθεῖται}{Intra-textuel (dans la colonne de R\expo{1} après le v. 276) et Marge gauche (Σ 277 de R\expo{2}, f. 165v, Planche 19)}{παρεπιγραφή. ἐκκυκλεῖται ἐπὶ τὸ ἔξω τὸ θεσμοφόριον.}{παρεπιγραφή · ἐκκυκλεῖτ(αι) ἐπὶ τ(ὸ) ἔξω τ(ὸ) θεσμοφόρι\expo{ον} {\policeFr ÷}}{« Didascalie scénique intercalée (παρεπιγραφή) : le sanctuaire du Thesmophorion est poussé au dehors sur l'ekkyklêma (ἐκκυκλεῖται ἐπὶ τὸ ἔξω τὸ Θεσμοφόριον). »}
\commScholie{Troisième parepigraphê scénique antique du Ravennas (citée par L. Hauseux, vv 1-641.tex, l. 2582, et étudiée sur la Planche 19 du Manuel) : son antiquité alexandrine est prouvée par le papyrus du IIe s. ap. J.-C. Ψ\expo{2} (POxy 3839) qui conserve la trace marginale ολολυ[ !}
\entreeScholie{106}{278}{ἔκσπευδε ταχέως· ὡς τὸ τῆς ἐκκλησίας / σημεῖον ἐν τῷ Θεσμοφορίῳ φαίνεται·}{Marge gauche (extérieure, f. 165v, en face du v. 278)}{ὅτε ἔμελλε γίνεσθαι ἐκκλησία, σημεῖον ἐτίθετο. οὕτως οὖν καὶ τῶν γυναικῶν μελλουσῶν ἐκκλησιάζειν σημεῖον τίθεται.}{ὅτε ἔμελλ(ε) γίνεσθ(αι) ἐκκλησία, σημεῖ\expo{ον} ἐτίθετ(ο) · οὕτ\expo{ως} οὖν \textbf{ς} τ\expo{ῶν} γυναικ\expo{ῶν} μελλουσ\expo{ῶν} ἐκκλησιάζ(ειν) σημεῖ\expo{ον} τίθετ(αι) {\policeFr ÷}}{« Lorsqu'une séance de l'Assemblée du peuple (ἐκκλησία) allait se tenir sur la Pnyx, on hissait un signal (un fanion). Ainsi donc, comme les femmes vont elles aussi tenir assemblée, on hisse le signal au Thesmophorion. »}
\commScholie{Note institutionnelle exacte sur le fanion (σημεῖον) hissé au début des séances de l'Ecclésia athénienne (cf. Acharniens v. 171, Guêpes v. 690).}
\entreeScholie{107}{279}{ἐγὼ δ᾽ ἄπειμι. — δεῦρό νυν ὦ Θρᾷτθ᾽ ἕπου·}{Interlinéaire / marge intérieure droite (f. 165v, v. 279)}{μιμεῖται τὴν φωνὴν γυναικός.}{μιμεῖτ(αι) τ(ὴν) φων(ὴν) γυναικός {\policeFr ÷}}{« Il contrefait la voix d'une femme (en interpellant une servante imaginaire, Thratta). »}
\commScholie{Note de jeu scénique indiquant qu'à partir du v. 279, Mnésiloque parle d'une voix de fausset et s'adresse à une esclave thrace fictive.}
\entreeScholie{108}{282}{ὦ Θρᾷττα θέασαι, καομένων τῶν λαμπάδων / ὅσον τὸ χρῆμ᾽ ἀνέρχεθ᾽ ὑπὸ τῆς λιγνύος· / ἀλλ᾽, ὦ περικαλλεῖ Θεσμοφόρω [τὼ θεώ Σ], δέξασθέ με}{Interlinéaire (supra lineam au-dessus de περικαλλεῖ, f. 165v)}{δυϊκόν. τὸ ἐντελὲς περικαλλέε.}{δυϊκόν · τ(ὸ) ἐντελ(ὲς) περικαλλέε {\policeFr ÷}}{« Forme de duel (contracte) ; la forme pleine non contractée est περικαλλέε. »}
\commScholie{Remarque morphologique sur le duel attique contracte en -εῖ des thèmes sigmatiques (type τὼ γένεε > γένει).}
\entreeScholie{109}{289}{καὶ τὴν θυγατέρα Χοιρίλ᾽ ἀνδρός μοι τυχεῖν}{Marge gauche (extérieure, f. 165v, en face du v. 289, Planche 23)}{ὡς τῆς θυγατρὸς αὐτοῦ κατ’ ἐπίκλησιν οὕτω καλουμένης. οἷον βοίδιον ἢ χρυσίον ἢ μύρτιον. ἴσως δὲ τὸ αἰδοῖον τὸ γυναικεῖον αἰνίττεται.}{ὡς τ\expo{ῆς} θυγατρ\expo{ὸς} αὐτ\expo{οῦ} κατ' ἐπίκλησ(ιν) οὕτω καλουμέν\expo{ης} · οἷον βοίδι\expo{ον} ἢ χρυσί\expo{ον} ἢ μύρτι\expo{ον} · ἴσ\expo{ως} δ(ὲ) τ(ὸ) αἰδοῖ\expo{ον} τ(ὸ) γυναικεῖ\expo{ον} αἰνίττετ(αι) [κύστ\expo{ον} ἴσ\expo{ως} αἰνίττετ(αι)] {\policeFr ÷}}{« Comme si sa fille était appelée ainsi par petit nom affectueux (ὑποκόρισμα, à savoir Χοιρίον, “Petit Cochon”), à la manière de Boïdion (“Petite Génisse”), Chrysion (“Petit Or”) ou Myrtion (“Petite Myrte”). Mais peut-être fait-il aussi allusion à mots couverts au sexe féminin (χοῖρος / κύσθος). »}
\commScholie{Pierre angulaire du « dossier Mnésiloque » de Louis Hauseux (Bibi, v. 289, Introduction.tex, ll. 325-357, et vv 1-641.tex, ll. 2620-2639) : alors que R\expo{1} porte χοῖρον et que la scholie glose le diminutif χοιρίον, Louis Hauseux démontre que sous χοῖρον / χοιρίον se cache avec double entente le nom même de la fille de Mnésiloque et épouse d'Euripide attestée par l'Hypothesis (Σ ante 1) : Χοιρίλ' (ou Χοιρίν').}
\entreeScholie{110}{290}{πλουτοῦντος, ἄλλως δ᾽ ἠλιθίου κἀβελτέρου,}{Marge gauche (extérieure, f. 165v, en face du v. 290)}{ἀναισθήτου καὶ μωροῦ, ἵνα κρατῇ αὐτοῦ.}{ἀναισθήτ\expo{ου} \textbf{ς} μωρ\expo{οῦ}, ἵνα κρατῇ αὐτ\expo{οῦ} {\policeFr ÷}}{« D'un mari stupide et sot, afin qu'elle puisse le mener par le bout du nez (le dominer). »}
\commScholie{Glose psychologique expliquant le vœu intéressé de la mère de famille.}
\entreeScholie{111}{291}{καὶ πρὸς θυλάκιον νοῦν ἔχειν μοι καὶ φρένας.}{Marge gauche (extérieure, bas du f. 165v, en face du v. 291)}{τὸν παιδαρίσκον. ἴσως δὲ παρὰ τὴν πόσθην αὐτὸ συνέθηκεν. ἀντὶ τοῦ κατωφερῆ τὸν νοῦν μου ποίησον.}{τ(ὸν) παιδαρίσκ\expo{ον} · ἴσ\expo{ως} δ(ὲ) παρὰ τ(ὴν) πόσθ(ην) αὐτ(ὸ) συνέθηκ(εν) · ἀν\auDessus{τ}{ι} τ\expo{οῦ} κατωφερῆ τ(ὸν) νοῦν μ\expo{ου} ποίησ\expo{ον} {\policeFr ÷}}{« Le petit garçon (que le scholiaste croit lire dans ποσθαλίσκον / θάληκον). Peut-être le poète a-t-il forgé ce mot sur πόσθη (la verge enfantine). — Autre glose : Au sens de “rends mon esprit enclin aux plaisirs d'en bas”. »}
\commScholie{Longuement discuté par Louis Hauseux (Bibi, v. 291, vv 1-641.tex, ll. 2640-2705) : contre la conjecture usuelle ποσθαλίσκον tirée de cette scholie, L. Hauseux s'appuie sur la leçon diplomatique de R\expo{1} (καὶ πρὸς θάληκον) et sur Hésychius θ 63 (θάλικον· σακίον, ἐν οἷσπερ τἀργύριον ταμιεύεται) pour restituer par résolution de métathèse καὶ πρὸς θυλάκιον (« et qu'il n'ait d'esprit et de tête que pour son escarcelle ! »).}
\clearpage
\renewcommand{\currHeadLE}{ΑΡΙΣΤΟΦΑΝΟΥΣ ΘΕΣΜΟΦΟΡΙΑΖΟΥΣΑΙ · CODEX RAVENNAS 429 (\R) · FOLIO 08/30}
\renewcommand{\currHeadRE}{f. 166r (Classense 168r) · vv. 295–336 · Collation diplomatique}
\renewcommand{\currHeadLO}{ΑΡΙΣΤΟΦΑΝΟΥΣ ΘΕΣΜΟΦΟΡΙΑΖΟΥΣΑΙ · CODEX RAVENNAS 429 (\R) · FOLIO 08/30}
\renewcommand{\currHeadRO}{f. 166r (Classense 168r) · vv. 295–336 · Collation diplomatique}
\noindent\begin{adjustbox}{max totalheight=30.4cm, max width=\textwidth, keepaspectratio}\begin{minipage}{\textwidth}
\bandeauRav{f. 166r}{Classense 168r}{295–336}{Folio 08/30 · Recto (page de droite) · 10 scholies \& gloses ({\policeGr Σ} 112–121)}
{\fontsize{8.4pt}{9.7pt}\selectfont
\ligneRav{295}{\inRas{Κῆρυξ}{}}{\hspace{1.1em}\difference{εὐφημία} ἔστω,}{\textbf{295} εὐφημί᾽ R : \lecon{εὐφημία} Suid. ε 3797}
\ligneRav{}{}{\hspace{1.1em}εὐφημία ἔστω·}{}
\ligneRav{297}{}{\hspace{1.1em}εὔχεσθε ταῖν Θεσμοφόροιν \textlbrackdbl\sout{τῇ Δήμητρῃ καὶ τῇ Κόρῃ}\textrbrackdbl}{\textbf{297} Θεσμοφόροιν τῇ Δήμητρῃ καὶ τῇ Κόρῃ R : \lecon{Θ.} Blaydes}
\ligneRav{}{}{\hspace{1.1em}καὶ τῷ Πλούτῳ καὶ τῇ Καλλιγενείᾳ}{}
\ligneRav{300}{}{\hspace{1.1em}καὶ τῇ Κουροτρόφῳ \textlbrackdbl\sout{τῇ Γῇ}\textrbrackdbl}{\textbf{300} τῇ Κουροτρόφῳ τῇ Γῇ R : \lecon{τῇ Κουροτρόφῳ} Dobree ex {\policeGr Σ}\expo{R}}
\ligneRav{301}{}{\hspace{1.1em}καὶ τῷ Ἑρμῇ καὶ \difference{ταῖς Χάρισιν}}{\textbf{301} Χάρισιν R : \lecon{ταῖς Χ.} Meineke}
\ligneRav{}{}{\hspace{1.1em}ἐκκλησίαν τήνδε καὶ σύνοδον τὴν νῦν}{}
\ligneRav{}{}{\hspace{1.1em}κάλλιστα καὶ ἄριστα ποῆσαι,}{}
\ligneRav{304}{}{\hspace{1.1em}πολυωφελῶς μὲν \difference{τῇ πόλει} τῇ Ἀθηναίων,}{\textbf{304} πόλει R : \lecon{τῇ π.} Meineke}
\ligneRav{305}{}{\hspace{1.1em}τυχηρῶς δ᾽ ἡμῖν αὐταῖς·}{\textbf{305} \lecon{ἡμῖν} R : ὑμῖν Fritzsche}
\ligneRav{306}{}{\hspace{1.1em}\expo{καὶ}τὴν δρῶσαν καὶ ἀγορεύουσαν}{\textbf{306} \expo{καὶ}τὴν ἀγορεύουσαν \Rdeux}
\ligneRav{307}{}{\hspace{1.1em}τὰ βέλτιστα περὶ τὸν δῆμον \difference{τὸν Ἀθηναίων}}{\textbf{307} τῶν Ἀθηναίων R : \lecon{τὸν Ἀθηναίων} Grynaeus}
\ligneRav{}{}{\hspace{1.1em}καὶ τὸν τῶν γυναικῶν}{}
\ligneRav{}{}{\hspace{1.1em}ταύτην νικᾶν·}{}
\ligneRav{310}{}{\hspace{1.1em}ταῦτ᾽ εὔχεσθε καὶ ὑμῖν αὐταῖς τἀγαθά·}{\textbf{310} \lecon{ὑμῖν} R : ἡμῖν \Mon}
\ligneRav{}{}{\hspace{1.1em}ἰὴ παιὼν ἰὴ παιὼν ἰὴ παιών. χαίρωμεν·}{}
\ligneRav{312}{χορ/}{\hspace{1.1em}\difference{δεχόμεθα} καὶ θεῶν γένος}{\textbf{312} δεχόμεσθα R : \lecon{-μεθα} Dindorf}
\ligneRav{313}{}{\hspace{1.1em}\difference{λιτόμεθα} ταῖσδ᾽ ἐπ᾽ εὐχαῖς}{\textbf{313} λιτόμεσθα R : \lecon{λιτόμεθα} Bourdin}
\ligneRav{}{}{\hspace{1.1em}φανέντας ἐπιχαρῆναι·}{}
\ligneRav{315}{}{\hspace{1.1em}Ζεῦ μεγαλώνυμε χρυσολύρα τε}{}
\ligneRav{}{}{\hspace{1.1em}Δῆλον ὃς ἔχεις ἱεράν,}{}
\ligneRav{317}{}{\hspace{1.1em}καὶ σὺ \difference{παγκρατὲς} \difference{κόρα}}{\textbf{317} παμκρατὲς R : \lecon{παγκρατὲς} ed. pr. \espApp κόρᾳ R : \lecon{κόρα} \Mon}
\ligneRav{}{}{\hspace{1.1em}γλαυκῶπι χρυσόλογχε πόλιν}{}
\ligneRav{}{}{\hspace{1.1em}οἰκοῦσα περιμάχητον, ἐλθὲ δεῦρο·}{}
\ligneRav{320}{}{\hspace{1.1em}καὶ πολυώνυμε \difference{θηροφόνη} (θηροφ\inRas{ω}{ό}νε παῖ—)}{\textbf{320} θηροφ\inRas{ω}{ό}νε παῖ-- \Rpc : \lecon{θηροφόνη} Hermann}
\ligneRav{}{}{\hspace{1.1em}Λατοῦς χρυσώπιδος ἔρνος,}{}
\ligneRav{}{}{\hspace{1.1em}σύ τε πόντιε σεμνὲ Πόσειδον}{}
\ligneRav{}{}{\hspace{1.1em}ἁλιμέδον}{}
\ligneRav{324}{}{\hspace{1.1em}προλιπὼν μυχὸν \difference{ἰχθυόεντα},}{\textbf{324} ἰχθυόεντ᾽ R, Suid. οι 181 : \lecon{ἰχθυόεντα} Wilamowitz}
\ligneRav{}{}{\hspace{1.1em}οἰστροδόνητον}{}
\ligneRav{325}{}{\hspace{1.1em}Νηρέος \difference{εἰναλίου} τε κόραι}{\textbf{325} ἐναλίου R : ἐνάλιοί Fritzsche : \lecon{εἰναλίου} Bentley}
\ligneRav{326}{}{\hspace{1.1em}Νύμφαι τ᾽ ὀρείπλαγκτοι·}{\textbf{326} \lecon{ὀρείπλαγκτοι} R : ὀρίπλαγκτοι Austin}
\ligneRav{}{}{\hspace{1.1em}χρυσέα τε φόρμιγξ}{}
\ligneRav{}{}{\hspace{1.1em}ἰαχήσειεν ἐπ᾽ εὐχαῖς}{}
\ligneRav{329}{}{\hspace{1.1em}ἡμετέραις· τελέως δ᾽ ἐκκλησιάσαιμεν \difference{Ἀθηνῶν}}{\textbf{329} Ἀθηναίων R : \lecon{Ἀθηνῶν} Reisig}
\ligneRav{330}{}{\hspace{1.1em}εὐγενεῖς γυναῖκες·}{}
\ligneRav{331}{κηρ/}{εὔχεσθε τοῖς θεοῖσι τοῖς Ὀλυμπίοις}{\textbf{331} (Κηρ.) R : \lecon{(Κορ.)} van Leeuwen}
\ligneRav{}{}{καὶ ταῖς Ὀλυμπίαισι καὶ τοῖς Πυθίοις}{}
\ligneRav{}{}{καὶ ταῖσι Πυθίαισι καὶ τοῖς Δηλίοις}{}
\ligneRav{334}{}{καὶ ταῖσι Δηλίαισι \difference{τοῖς τ᾽} ἄλλοις θεοῖς,}{\textbf{334} καὶ τοῖς R : \lecon{τοῖς τ᾽} Grynaeus}
\ligneRav{335}{}{εἴ τις ἐπιβουλεύει τι τῷ δήμῳ κακὸν}{}
\ligneRav{}{}{τῷ τῶν γυναικῶν ἢ 'πικηρυκεύεται}{}
}
\vspace{3pt}{\color{ruleGold}\hrule height 0.6pt}\vspace{3pt}
{\fontsize{8.0pt}{9.4pt}\selectfont
\noindent{\bfseries\scshape\color{rubricRed}Loci paralleli et testimonia (f. 166r, vv. 295–336) :} \textbf{299} (Καλλιγενείᾳ = Δημήτερι) Hesychius κ 472 Καλλιγένειαν ; cf. Photius κ 118 (= fr. Aristophanes 331 Kassel-Austin) \espTest \textbf{300} (Κουροτρόφῳ = Δημήτηρ) Hesychius κ 3865 κουροτρόφος ; scholium vet. in Hesiodus \italique{Opera et dies} 227-9 \espTest \textbf{322-4} (σύ τε -- οἰστροδόνητον) Suidas οι 181 (sed σεμνὲ et ἁλιμέδον omissis) \espTest \textbf{335-7} Satyrus \italique{Vit. Euripid.} (= \italique{The Oxyrhynchus Papyri} IX 1176) fr. 39 col. XII 8-16 (Εὐριπιδίῳ pro Εὐριπίδῃ (?)) \espTest \textbf{336} ('πικηρυκεύεται) Photius ε 1580 ἐπικηρυκεύεται = \Syn ε 682 = Suidas ε 2375 = \EtymMagnum 360.21 = \EtymGen (B) ; cf. Hesychius ε 4853\par\vspace{2.5pt}
{\color{ruleGold}\hrule height 0.35pt}\vspace{2.5pt}
\noindent{\bfseries\scshape\color{rubricRed}Suite (en commentaire) de l'apparat critique et paléographique :} \textbf{Témoin papyrologique :} \textbf{\PsiCinq{} (\italique{P.Oxy.} IX 1176, Satyrus)} cite les \textbf{vv. 335–337}. --- \textbf{Grattage de la rubrique d'interlocuteur au v. 295 :} \Rdeux{} a gratté Κῆρυξ (\inRas{Κῆρυξ}{} \Rdeux{} : \auDessus{κη}{ρυ} \Mon{}, tandis que {\policeGr Σ}\expo{R} explique γυνὴ μιμουμένη κήρυκα ; L. Hauseux édite Κορ.-Ἱέρεια). --- \textbf{Intrusions de gloses religieuses (vv. 297 et 300) :} Dans la prière en prose de la Héraut, R a incorporé dans le texte les deux gloses interlinéaires τῇ Δήμητρῃ καὶ τῇ Κόρῃ (glosant τοῖν Θεσμοφόροιν) et τῇ Γῇ (glosant τῇ Κουροτρόφῳ). --- \textbf{Ajout supralinéaire de \Rdeux{} (v. 306) :} \expo{καὶ}τὴν ἀγορεύουσαν. --- \textbf{Correction \italique{in rasura} au v. 320 :} θηροφ\inRas{ω}{ό}νε παῖ-- \Rpc{} (ω de \Rac{} gratté et corrigé en ό). --- \textbf{v. 295 :} \inRas{Κῆρυξ}{} \Rdeux{} : \auDessus{κη}{ρυ} \Mon{} : γυνὴ μιμουμένη κήρυκα {\policeGr Σ} : Κρίτυλλα Sommerstein : \lecon{(Κορ.-Ἱέρεια)} \lecon{Bibi} \espApp εὐφημί' R : \lecon{εὐφημία} Souda ε 3797. --- \textbf{v. 297 :} Θεσμοφόροιν τῇ Δήμητρῃ καὶ τῇ Κόρῃ R : \lecon{Θεσμοφόροιν} Blaydes (suppression de la glose). --- \textbf{v. 300 :} τῇ Κουροτρόφῳ τῇ Γῇ R : \lecon{τῇ Κουροτρόφῳ} Dobree ex {\policeGr Σ}\expo{R}.\par}
\end{minipage}\end{adjustbox}
\newpage
\renewcommand{\currHeadLE}{SCHOLIES MARGINALES ({\policeGr Σ}\expo{R}) \& GLOSES INTERLINÉAIRES · FOLIO 08/30}
\renewcommand{\currHeadRE}{f. 166r (Classense 168r) · vv. 295–336 · {\policeGr Σ} 112 à {\policeGr Σ} 121}
\renewcommand{\currHeadLO}{SCHOLIES MARGINALES ({\policeGr Σ}\expo{R}) \& GLOSES INTERLINÉAIRES · FOLIO 08/30}
\renewcommand{\currHeadRO}{f. 166r (Classense 168r) · vv. 295–336 · {\policeGr Σ} 112 à {\policeGr Σ} 121}
\bandeauScholies{Folio 08/30 — f. 166r}{Classense 168r}{vv. 295–336}{Recto (page de droite) · 10 entrées ({\policeGr Σ} 112 à {\policeGr Σ} 121)}
{\fontsize{8.8pt}{10.5pt}\selectfont\noindent\textbf{\color{rubricRed}Disposition topographique des marges sur le f. 166r :} Marge extérieure à droite ; marge intérieure à gauche. — \textit{Fac-similé de référence :} \texttt{folios\_individuels\_ff162v-177r/08\_f166r\_classense168r\_vv0295-0336.jpg}.\par\vspace{3pt}}
\entreeScholie{112}{295}{εὐφημία ἔστω, / εὐφημία ἔστω· / εὔχεσθε ταῖν Θεσμοφόροιν [τοῖν Θεσμοφόροιν Σ]}{Marge supérieure et marge droite (extérieure, f. 166r, en tête de la prière en prose)}{γυνὴ μιμουμένη κήρυκα. πεζῇ εὐφημίᾳ χρῶνται οἱ κωμικοὶ, ἐπειδὰν εὐχὴν ἢ ψήφισμα εἰσάγωσιν.}{γυνὴ μιμουμέν(η) κήρυκα · πεζῇ εὐφημίᾳ χρῶντ(αι) οἱ κωμικ(οὶ), ἐπειδὰν εὐχ(ὴν) ἢ ψήφισμ(α) εἰσάγωσ(ιν) {\policeFr ÷}}{« C'est une femme qui imite le héraut (γυνὴ μιμουμένη κήρυκα). Et les poètes comiques emploient une formule de silence sacré en prose (πεζῇ εὐφημίᾳ) chaque fois qu'ils introduisent une prière officielle ou un décret. »}
\commScholie{Scholie essentielle (citée par L. Hauseux, vv 1-641.tex, ll. 2716-2720) : dans R, le scribe avait d'abord inscrit la rubrique Κῆρυξ devant εὐφημί' ἔστω, puis R\expo{2} a gratté Κῆρυξ (laissé en rasura, mais encore copié κη[ρυ] par M) en notant en marge qu'il ne s'agit pas d'un héraut masculin mais d'une femme faisant office de héraut (édité Κορ.-Ἱέρεια par Bibi).}
\entreeScholie{113}{297, 300}{{\policeFr [gloses intruses dans R\expo{1} :]} εὔχεσθε ταῖν Θεσμοφόροιν [τῇ Δήμητρῃ καὶ τῇ Κόρῃ] ... καὶ τῇ Κουροτρόφῳ [τῇ Γῇ]}{Intra-textuel (glose interlinéaire incorporée dans la prose de R\expo{1} aux vv. 297 et 300, f. 166r)}{[τῇ Δήμητρῃ καὶ τῇ Κόρῃ] (v. 297) · [τῇ Γῇ] (v. 300) (gloses intruses incorporées dans la prose de R\expo{1}).}{[τῇ Δήμητρῃ \textbf{ς} τῇ Κόρῃ] (v. 297) · [τῇ Γῇ] (v. 300)}{« [à Déméter et à Coré] ... [à la Terre]. »}
\commScholie{Comme l'ont reconnu Dobree et Blaydes (suivis par Louis Hauseux, Bibi, vv. 297-298, ll. 2717-2725), les mots τῇ Δήμητρῃ καὶ τῇ Κόρῃ (après τοῖν Θεσμοφόροιν) et τῇ Γῇ (après τῇ Κουροτρόφῳ) sont deux gloses explicatives interlinéaires qui se sont glissées dans le texte en prose de R\expo{1} (comme le prouve {\policeGr Σ}\expo{R} 299 qui hésite encore : εἴτε τῇ Γῇ εἴτε τῇ Ἑστίᾳ).}
\entreeScholie{114}{298}{καὶ τῷ Πλούτῳ [Πλούτωνι Σ] καὶ τῇ Καλλιγενείᾳ}{Marge droite (extérieure, f. 166r, deux scholies en face du v. 298)}{ὡς ἀνδρὶ τῆς Περσεφόνης. δαίμων περὶ τὴν Δήμητραν, ἣν προλογίζουσαν ἐν ταῖς ἑτέραις Θεσμοφοριαζούσαις ἐποίησεν.}{ὡς ἀνδρὶ τ\expo{ῆς} Περσεφόν\expo{ης} {\policeFr ÷} // δαίμ(ων) π(ερὶ) τ(ὴν) Δήμητρ(αν), ἣν προλογίζουσ(αν) ἐν τ\expo{αῖς} ἑτέρ\expo{αις} Θεσμοφοριαζούσ\expo{αις} ἐποίησ(εν) {\policeFr ÷}}{« τῷ Πλούτωνι : en tant qu'époux de Perséphone. — τῇ Καλλιγενείᾳ : divinité du cortège de Déméter, dont Aristophane a fait la récitante du prologue dans ses secondes Thesmophories (Θεσμοφοριάζουσαι β', fr. 331 K.-A.). »}
\commScholie{Information littéraire inestimable (commentée par L. Hauseux, vv 1-641.tex, l. 2672) : c'est par cette scholie du Ravennas que nous savons que la déesse Kalligéneia prononçait le prologue de la seconde pièce d'Aristophane intitulée Les Thesmophories.}
\entreeScholie{115}{299}{καὶ τῇ Κουροτρόφῳ}{Marge droite (extérieure, f. 166r, en face du v. 299)}{εἴτε τῇ γῇ εἴτε τῇ ἑστίᾳ, ὁμοίως πρὸ τοῦ Διὸς θύουσιν αὐτῇ.}{εἴτε τῇ γῇ εἴτε τῇ ἑστίᾳ, ὁμοί\expo{ως} πρὸ τ\expo{οῦ} Δι\expo{ὸς} θύουσ(ιν) αὐτῇ {\policeFr ÷}}{« Qu'il s'agisse de la Terre (Gê) ou d'Hestia, on lui offre pareillement un sacrifice préliminaire avant Zeus. »}
\commScholie{Repris par la Souda κ 2193. Prouve que dans l'exemplaire commenté par le scholiaste alexandrin, la glose τῇ Γῇ (présente dans R\expo{1}) ne figurait pas encore dans le texte d'Aristophane !}
\entreeScholie{116}{310}{ταῦτ᾽ εὔχεσθε, καὶ ὑμῖν αὐταῖς τἀγαθά.}{Interlinéaire (supra lineam, f. 166r, v. 310)}{ἀπὸ κοινοῦ εὔχεσθε.}{ἀπ(ὸ) κοιν\expo{οῦ} εὔχεσθ(ε) {\policeFr ÷}}{« Le verbe εὔχεσθε (“demandez dans vos prières”) est sous-entendu en facteur commun. »}
\commScholie{Glose syntaxique de R\expo{2} signalant que l'impératif εὔχεσθε est construit en facteur commun (ἀπὸ κοινοῦ) avec les deux compléments coordonnés ταῦτα et καὶ ὑμῖν αὐταῖς τἀγαθά dans la formule liturgique finale de la Héraut.}
\entreeScholie{117}{312}{δεχόμεθα [δεχόμεσθα R] καὶ θεῶν γένος / λιτόμεθα [λιτόμεσθα R] ταῖσδ᾽ ἐπ᾽ εὐχαῖς}{Marge intérieure gauche (f. 166r, en face du v. 312, entrée du Chœur)}{ὁ χορὸς ἀπὸ γυναικῶν θεσμοφοριαζουσῶν.}{ὁ χορ(ὸς) ἀπ(ὸ) γυναικ\expo{ῶν} θεσμοφοριαζουσ\expo{ῶν} {\policeFr ÷}}{« Ici chante le Chœur formé des femmes célébrant les Thesmophories. »}
\commScholie{Rubrique scénique marquant le premier chant choral (vv. 312-330).}
\entreeScholie{118}{324}{σύ τε πόντιε σεμνὲ Πόσειδον / ἁλιμέδον / προλιπὼν μυχὸν ἰχθυόεντα / οἰστροδόνητον}{Marge droite (extérieure, f. 166r, en face du v. 324)}{τὸν ὑπὸ τῶν ἀνέμων κινούμενον. πᾶσαν δὲ κίνησιν καλοῦσιν οἶστρον.}{τ(ὸν) ὑπ(ὸ) τ\expo{ῶν} ἀνέμ\expo{ων} κινούμεν\expo{ον} · πᾶσ(αν) δ(ὲ) κίνησ(ιν) καλοῦσ(ιν) οἶστρ\expo{ον} {\policeFr ÷}}{« Agité par les vents : car les poètes appellent οἶστρος (“aiguillon, fureur”) toute agitation impétueuse (ici le tumulte marin). »}
\commScholie{Explique la variante/conjecture sur μυχὸν ἰχθυόεντ' οἰστροδόνητον (vv. 324-325).}
\entreeScholie{119}{327}{χρυσέα τε φόρμιγξ / ἰαχήσειεν ἐπ᾽ εὐχαῖς}{Interlinéaire (supra lineam, f. 166r, v. 327)}{ἡ τοῦ Ἀπόλλωνος.}{ἡ τ\expo{οῦ} Ἀπόλλων\expo{ος} {\policeFr ÷}}{« Celle (la lyre dorée) d'Apollon. »}
\commScholie{Glose interlinéaire de R\expo{2} identifiant la « lyre d'or » (χρυσέα τε φόρμιγξ) comme celle d'Apollon (ἡ τοῦ Ἀπόλλωνος), reprenant l'invocation initiale au dieu citharède du v. 315.}
\entreeScholie{120}{329}{τελέως δ᾽ ἐκκλησιάσαιμεν Ἀθηνῶν [Ἀθηναίων R] / εὐγενεῖς γυναῖκες·}{Marge droite (extérieure) et interlinéaire sur τελέως (f. 166r, v. 329)}{ἐπὶ τὸ τελεστικὰ τὰ πράγματα γενέσθαι. εἰς τέλος ἄγοιμεν τὰ δόξαντα.}{ἐπὶ τ(ὸ) τελεστικὰ τὰ πράγματ(α) γενέσθ(αι) {\policeFr ÷} // εἰς τέλ\expo{ος} ἄγοιμ(εν) τὰ δόξαντ(α) {\policeFr ÷}}{« Afin que nos délibérations aboutissent à un résultat accompli. — Puissions-nous mener à bonne fin nos résolutions. »}
\commScholie{Double glose expliquant l'adverbe τελέως dans le vœu final du Chœur.}
\entreeScholie{121}{336}{εἴ τις ἐπιβουλεύει τι τῷ δήμῳ κακὸν / τῷ τῶν γυναικῶν, ἢ ᾽πικηρυκεύεται}{Interlinéaire (supra lineam au-dessus d'ἐπικηρυκεύεται, dernière ligne du f. 166r)}{φιλίαν σπένδει.}{φιλί(αν) σπένδ(ει) {\policeFr ÷}}{« Négocie secrètement la paix (par envoi de héraut) avec l'ennemi. »}
\commScholie{Cité par Satyrus dans sa Vie d'Euripide (Ψ\expo{5} = POxy 1176, fr. 39 col. XIII) et commenté par L. Hauseux (vv 1-641.tex, l. 2864 n. 5).}
\clearpage
\renewcommand{\currHeadLE}{ΑΡΙΣΤΟΦΑΝΟΥΣ ΘΕΣΜΟΦΟΡΙΑΖΟΥΣΑΙ · CODEX RAVENNAS 429 (\R) · FOLIO 09/30}
\renewcommand{\currHeadRE}{f. 166v (Classense 168v) · vv. 337–379 · Collation diplomatique}
\renewcommand{\currHeadLO}{ΑΡΙΣΤΟΦΑΝΟΥΣ ΘΕΣΜΟΦΟΡΙΑΖΟΥΣΑΙ · CODEX RAVENNAS 429 (\R) · FOLIO 09/30}
\renewcommand{\currHeadRO}{f. 166v (Classense 168v) · vv. 337–379 · Collation diplomatique}
\noindent\begin{adjustbox}{max totalheight=30.4cm, max width=\textwidth, keepaspectratio}\begin{minipage}{\textwidth}
\bandeauRav{f. 166v}{Classense 168v}{337–379}{Folio 09/30 · Verso (page de gauche) · 14 scholies \& gloses ({\policeGr Σ} 122–135)}
{\fontsize{8.8pt}{10.2pt}\selectfont
\ligneRav{337}{χορ/}{†\difference{Εὐριπίδῃ Μήδοις τ᾽}† ἐπὶ βλάβῃ τινὶ}{\textbf{337} ευριπιδιω\douteux{ι}[. .]\douteux{υ}τ \PsiCinq : Εὐριπίδῃ Μήδοις R : \lecon{Εὐριπίδῃ Μήδοις τ᾽} Biset}
\ligneRav{}{}{τῇ τῶν γυναικῶν, ἢ τυραννεῖν ἐπινοεῖ}{}
\ligneRav{}{}{ἢ τὸν τύραννον συγκατάγειν, ἢ παιδίον}{}
\ligneRav{340}{}{ὑποβαλλομένης \difference{κατεῖπέ τις}, ἢ δούλη τινὸς}{\textbf{340} κατεῖπεν τις R : \lecon{κατεῖπέ τις} Brunck}
\ligneRav{341}{}{προαγωγὸς οὖσ᾽ ἐνετρύλλισεν τῷ δεσπότῃ}{\textbf{341} \lecon{ἐνετρύλλισεν} R, Suid. ε 1271 : ἐνετρύλισεν Dindorf}
\ligneRav{}{}{ἢ πεμπομένη τις ἀγγελίας ψευδεῖς φέρει,}{}
\ligneRav{}{}{ἢ μοιχὸς εἴ τις ἐξαπατᾷ ψευδῆ λέγων}{}
\ligneRav{344}{}{καὶ μὴ δίδωσιν \difference{ἃν} ὑπόσχηταί ποτε,}{\textbf{344} ἅ ἂν R : \lecon{ἃν} Faber}
\ligneRav{345}{}{ἢ δῶρά τις δίδωσι μοιχῷ γραῦς γυνή,}{}
\ligneRav{346}{}{ἢ καὶ δέχεται προδιδοῦσ᾽ \difference{ἑταίρᾳ} τὸν φίλον,}{\textbf{346} ἑταίρα R : \lecon{ἑταίρᾳ} Fritzsche ap Σ}
\ligneRav{347}{}{κεἴ τις κάπηλος ἢ καπηλὶς τοῦ χοὸς}{\textbf{347} \lecon{χοὸς} R, Suid. λ 394 : χοῶς Brunck}
\ligneRav{}{}{ἢ τῶν κοτυλῶν τὸ νόμισμα διαλυμαίνεται,}{}
\ligneRav{349}{}{κακῶς ἀπολέσθαι τοῦτον αὐτὸν \difference{κᾠκίαν}}{\textbf{349} κ' οἰκίαν R : \lecon{κᾠκίαν} Brunck}
\ligneRav{350}{}{ἀρᾶσθε, ταῖς δ᾽ ἄλλαισιν ὑμῖν τοὺς θεοὺς}{\textbf{350} \lecon{ὑμῖν} R : ἡμῖν \Mon}
\ligneRav{351}{}{εὔχεσθε πάσαις πολλὰ δοῦναι κἀγαθά·}{\textbf{351} \lecon{κἀγαθά} R : τἀγαθά Zanetti}
\ligneRav{352}{χορ/}{\hspace{1.1em}\difference{ξυνευχόμεσθα} τέλεα μὲν}{\textbf{352} ξυνευχόμεθα R : \lecon{ξυνευχόμεσθα} Daubruz}
\ligneRav{353}{}{\hspace{1.1em}πόλει τέλεα τε δήμῳ}{\textbf{353} \lecon{τε} R : δὲ \Mon}
\ligneRav{354}{}{\hspace{1.1em}τάδ᾽ εὔγματ᾽ \difference{ἀποτελεῖσθαι},}{\textbf{354} γενέσθαι R : ἐκγενέσθαι Dindorf : \lecon{ἀποτελεῖσθαι} Willems}
\ligneRav{355}{}{\hspace{1.1em}τὰ δ᾽ ἄρισθ᾽ ὅσαις προσήκει}{}
\ligneRav{}{}{\hspace{1.1em}νικᾶν λεγούσαις· ὁπόσαι δ᾽}{}
\ligneRav{357}{}{\hspace{1.1em}\difference{ἐξαπατῶσιν} παραβαίνουσί τε τοὺς}{\textbf{357} ἐξαπατῶσι R : \lecon{ἐξαπατῶσιν} Bentley}
\ligneRav{}{}{\hspace{1.1em}ὅρκους τοὺς νενομισμένους}{}
\ligneRav{360}{}{\hspace{1.1em}κερδῶν \difference{οὕνεκ᾽} ἐπὶ βλάβῃ,}{\textbf{360} ἕνεκ᾽ R : \lecon{οὕνεκ᾽} Bentley}
\ligneRav{}{}{\hspace{1.1em}ἢ ψηφίσματα καὶ νόμον}{}
\ligneRav{}{}{\hspace{1.1em}ζητοῦσ᾽ ἀντιμεθιστάναι,}{}
\ligneRav{}{}{\hspace{1.1em}τἀπόρρητά τε τοῖσιν ἐχ-}{}
\ligneRav{364}{}{\hspace{1.1em} θροῖς τοῖς ἡμετέροις \difference{λέγουσ᾽},}{\textbf{364} λέγουσιν R, Suid. α 3501 : \lecon{λέγουσ᾽} Bentley}
\ligneRav{365}{}{\hspace{1.1em}ἢ Μήδους ἐπάγουσι τῆς}{}
\ligneRav{}{}{\hspace{1.1em}χώρας οὕνεκ᾽ ἐπὶ βλάβῃ,}{}
\ligneRav{367}{}{\hspace{1.1em}\difference{ἀσεβοῦσ᾽ ἀδικοῦσί} τε τὴν πόλιν·}{\textbf{367} ἀσεβοῦσιν ἀδικοῦσίν R : \lecon{ἀσεβοῦσ᾽ ἀδικοῦσί} Bothe}
\ligneRav{}{}{\hspace{1.1em}ἀλλ᾽ ὦ παγκρατὲς}{}
\ligneRav{}{}{\hspace{1.1em}Ζεῦ ταῦτα κυρώσειας, ὥσθ᾽}{}
\ligneRav{370}{}{\hspace{1.1em}ἡμῖν θεοὺς παραστατεῖν}{}
\ligneRav{}{}{\hspace{1.1em}καίπερ γυναιξὶν οὔσαις·}{}
\ligneRav{372}{χορ/}{ἄκουε \difference{πᾶσ᾽}. ἔδοξε τῇ βουλῇ τάδε}{\textbf{372} πᾶς R : \lecon{πᾶσ᾽} Voss}
\ligneRav{373}{}{τῇ τῶν γυναικῶν· Τιμόκλει᾽ ἐπεστάτει,}{\textbf{373} Ἀρ[. . . .]\douteux{λ}(vel \douteux{a}, \douteux{δ}, \douteux{μ})ειa \PsiCinq, Ἀρ[χίκλ]ει᾽ suppl. Hunt : \lecon{Τιμόκλει᾽} R, Σ, Suid. ε 2616}
\ligneRav{374}{}{Λύσιλλ᾽ ἐγραμμάτευεν, εἶπε Ζωστράτη·}{}
\ligneRav{375}{}{ἐκκλησίαν ποιεῖν ἕωθεν τῇ μέσῃ}{}
\ligneRav{376}{}{τῶν Θεσμοφορίων, ᾗ μάλισθ᾽ ἡμῖν \difference{σχολή},}{\textbf{376} σχολῇ R : \lecon{σχολή} Zanetti}
\ligneRav{}{}{καὶ χρηματίζειν πρῶτα περὶ Εὐριπίδου,}{}
\ligneRav{}{}{ὅ τι χρὴ παθεῖν ἐκεῖνον· ἀδικεῖν γὰρ δοκεῖ}{}
\ligneRav{}{}{ἡμῖν ἁπάσαις. τίς ἀγορεύειν βούλεται;}{}
}
\vspace{3pt}{\color{ruleGold}\hrule height 0.6pt}\vspace{3pt}
{\fontsize{8.0pt}{9.4pt}\selectfont
\noindent{\bfseries\scshape\color{rubricRed}Loci paralleli et testimonia (f. 166v, vv. 337–379) :} \textbf{341} Suidas ε 1271 \espTest (προαγωγὸς) Hesychius π 3308 προαγωγός \espTest (ἐνετρύλλισεν) Zonaras 745 οὔ σ' ἐνετρύλλησεν \espTest \textbf{347-8} Suidas χ 394 (sed εἴ pro κεἴ) \espTest \textbf{363-4} (τἀπόρρητά -- λέγουσ᾽) Suidas α 3501 ἀπόρρητά τε τοῖσιν ἐχθροῖσι τοῖς ἡμετέροις λέγουσιν \espTest \textbf{369-71} (ὥσθ᾽ -- οὔσαις) Suidas π 444 \espTest \textbf{373-4} Satyrus \italique{Vit. Euripid.} (= \italique{The Oxyrhynchus Papyri} IX 1176) fr. 39 col. XII 1 (vid. ad 335-7 ; sed Ἀρ[χίκλ]ει᾽ pro Τιμόκλει᾽, suppl. Hunt) \espTest (Τιμόκλει᾽ -- ἐγραμμάτευεν) Suidas ε 2616 (sed ἐπιστάτει pro ἐπεστάτει et ἐγραμμάτευε pro ἐγραμμάτευεν) \espTest \textbf{377-8} (καὶ χρηματίζειν -- παθεῖν) Suidas χ 481 ; Lexicon Vindobonense 190.10\par\vspace{2.5pt}
{\color{ruleGold}\hrule height 0.35pt}\vspace{2.5pt}
\noindent{\bfseries\scshape\color{rubricRed}Suite (en commentaire) de l'apparat critique et paléographique :} \textbf{Témoin papyrologique :} \textbf{\PsiCinq{} (\italique{P.Oxy.} IX 1176, Satyrus)} cite les \textbf{vv. 335–337} et \textbf{vv. 373–375}. --- \textbf{v. 346 :} ἑταίρα R : \lecon{ἑταίρᾳ} Fritzsche ap. {\policeGr Σ}\expo{R}. --- \textbf{v. 354 :} γενέσθαι R (ἐκγενέσθαι Dindorf) : \lecon{ἀποτελεῖσθαι} Alphonse Willems (adopté par L. Hauseux : γενέσθαι est une glose banale substituée au verbe rituel ἀποτελεῖσθαι). --- \textbf{v. 374 :} Ἀρ[....]λεια \PsiCinq{} (Ἀρ[χίκλ]ει' suppl. Hunt) : \lecon{Τιμόκλει'} R, {\policeGr Σ}\expo{R}, Souda ε 2616 (maintenu par L. Hauseux).\par}
\end{minipage}\end{adjustbox}
\newpage
\renewcommand{\currHeadLE}{SCHOLIES MARGINALES ({\policeGr Σ}\expo{R}) \& GLOSES INTERLINÉAIRES · FOLIO 09/30}
\renewcommand{\currHeadRE}{f. 166v (Classense 168v) · vv. 337–379 · {\policeGr Σ} 122 à {\policeGr Σ} 135}
\renewcommand{\currHeadLO}{SCHOLIES MARGINALES ({\policeGr Σ}\expo{R}) \& GLOSES INTERLINÉAIRES · FOLIO 09/30}
\renewcommand{\currHeadRO}{f. 166v (Classense 168v) · vv. 337–379 · {\policeGr Σ} 122 à {\policeGr Σ} 135}
\bandeauScholies{Folio 09/30 — f. 166v}{Classense 168v}{vv. 337–379}{Verso (page de gauche) · 14 entrées ({\policeGr Σ} 122 à {\policeGr Σ} 135)}
{\fontsize{8.8pt}{10.5pt}\selectfont\noindent\textbf{\color{rubricRed}Disposition topographique des marges sur le f. 166v :} Marge extérieure à gauche ; marge intérieure à droite. — \textit{Fac-similé de référence :} \texttt{folios\_individuels\_ff162v-177r/09\_f166v\_classense168v\_vv0337-0379.jpg}.\par\vspace{3pt}}
\entreeScholie{122}{339}{ἢ τυραννεῖν ἐπινοεῖ / ἢ τὸν τύραννον συγκατάγειν,}{Marge gauche (extérieure, f. 166v, en haut du verso, en face des vv. 338–339)}{ἐκ τῆς κατάρας τῆς ἐπὶ τοῖς Πεισιστρατίδαις γενομένης ταῦτα παρέμιξεν. ἔοικέ τινα παραμιγνύναι τῶν ἐφ’ Ἱππίᾳ ψηφισθέντων.}{ἐκ τ\expo{ῆς} κατάρ\expo{ας} τ\expo{ῆς} ἐπὶ τ\expo{οῖς} Πεισιστρατίδ\expo{αις} γενομέν\expo{ης} ταῦτ(α) παρέμιξ(εν) · ἔοικέ τινα παραμιγνύν(αι) τ\expo{ῶν} ἐφ' Ἱππίᾳ ψηφισθέντ\expo{ων} {\policeFr ÷}}{« Il a mêlé ici des formules tirées de l'imprécation officielle prononcée contre les Pisistratides ; il semble y incorporer certaines clauses des décrets votés contre le tyran Hippias. »}
\commScholie{Remarquable scholie historique (remontant à Cratéros ou Aristarque) qui identifie la source officielle de la parodie d'Aristophane : la grande imprécation civique prononcée par le héraut à l'ouverture de chaque séance de la Boulè et de l'Ecclésia contre quiconque tenterait de rétablir la tyrannie des Pisistratides ou de négocier avec le Grand Roi.}
\entreeScholie{123}{340}{ἢ παιδίον / ὑποβαλλομένης κατεῖπέ τις [κατεῖπεν τις R],}{Interlinéaire (supra lineam au-dessus de κατεῖπεν, f. 166v, v. 341)}{κατεῖπε πρὸς τὸν δεσπότην.}{κατεῖπ(ε) πρ(ὸς) τ(ὸν) δεσπότ(ην) {\policeFr ÷}}{« A dénoncé au maître (au mari) une femme qui supposait un enfant. »}
\commScholie{Précise le complément implicite du verbe κατεῖπεν.}
\entreeScholie{124}{345}{ἢ μοιχὸς εἴ τις ἐξαπατᾷ ψευδῆ λέγων / καὶ μὴ δίδωσιν ἃν ὑπόσχηταί ποτε, / ἢ δῶρά τις δίδωσι μοιχῷ γραῦς γυνή,}{Interlinéaire (supra lineam, f. 166v, v. 345)}{ἵνα αὐτῇ συγγένηται.}{ἵνα αὐτῇ συγγένητ(αι) {\policeFr ÷}}{« (Lorsque l'amant a fait de fausses promesses) afin d'obtenir ses faveurs (de coucher avec elle). »}
\commScholie{Glose interlinéaire de R\expo{2} précisant le mobile de la vieille femme qui offre des présents à un amant (ἢ δῶρά τις δίδωσι μοιχῷ γραῦς γυνή, v. 345) ou de l'amant qui trompe par de fausses promesses (v. 343) : « afin qu'il s'unisse à elle » (ἵνα αὐτῇ συγγένηται).}
\entreeScholie{125}{346}{ἢ καὶ δέχεται προδιδοῦσ᾽ ἑταίρᾳ τὸν φίλον,}{Marge gauche (extérieure, f. 166v, bloc de 6 lignes en face du v. 346)}{ἢ καὶ δέχεται δῶρα παρ’ ἑταίρας ἐπὶ τῷ προδιδόναι καὶ συγχωρῆσαι αὐτῇ τὸν φίλον. γράφεται ἔν τισιν ἑταίρα καὶ ἔστιν ἕτερος νοῦς. Ἄλλως. ἐὰν μὲν ἄνευ τοῦ ι γρ. ἑταίρα, οὕτως· ἢ δέχεται δῶρα ὑπὲρ τοῦ ἑταίρα οὖσα προδιδόναι τὸν φίλον. ἐὰν δὲ σὺν τῷ ι, προδιδοῦσάν τινα ἑταίραν τὸν ἑαυτῆς φίλον.}{ἢ \textbf{ς} δέχετ(αι) δῶρα παρ’ ἑταίρας ἐπὶ τῷ προδιδόναι \textbf{ς} συγχωρῆσ(αι) αὐτῇ τ(ὸν) φίλον · γράφετ(αι) ἔν τισ(ιν) ἑταίρα \textbf{ς} ἔστ(ιν) ἕτερ\expo{ος} νοῦς · ἄλλ\expo{ως} · ἐὰν μ(ὲν) ἄνευ τ\expo{οῦ} ι γρ · ἑταίρα, οὕτ\expo{ως}· ἢ δέχετ(αι) δῶρα ὑπὲρ τ\expo{οῦ} ἑταίρα οὖσα προδιδόναι τ(ὸν) φίλον · ἐὰν δ(ὲ) σὺν τῷ ι, προδιδοῦσάν τινα ἑταίραν τ(ὸν) ἑαυτ\expo{ῆς} φίλον {\policeFr ÷}}{« Ou bien accepte des présents de la part d'une courtisane pour trahir son amant et le lui céder. Dans certains exemplaires, on écrit ἑταίρα (au nominatif sans iota), et le sens est différent. — Autre rédaction : Si l'on écrit ἑταίρα sans iota adscrit (au nominatif), le sens est : “ou si une femme, étant courtisane, reçoit des présents pour trahir son amant” ; mais si l'on écrit ἑταίρᾳ avec l'iota (au datif, ou ἑταίραν à l'accusatif), le sens est : “trahissant son propre amant au profit d'une amie”. »}
\commScholie{Scholie philologique décisive exploitée par Fritzsche et Louis Hauseux (Bibi, v. 346, vv 1-641.tex, ll. 2889-2898) : alors que R\expo{1} porte le nominatif sans iota ἑταίρα, la première explication du scholiaste (ἐὰν δὲ σὺν τῷ ι) atteste et justifie le datif ἑταίρᾳ (« trahissant son amant au profit d'une amie »), adopté par Bibi.}
\entreeScholie{126}{347}{κεἴ τις κάπηλος ἢ καπηλὶς τοῦ χοὸς / ἢ τῶν κοτυλῶν τὸ νόμισμα διαλυμαίνεται,}{Marge gauche (extérieure, f. 166v, en face du v. 347)}{τοῦ διξέστου. χοὺς γὰρ δύο ξέσται, χοεὺς δὲ ἕξ.}{τ\expo{οῦ} διξέστ\expo{ου} · χοὺς γ(ὰρ) δύο ξέστ(αι), χοεὺς δ(ὲ) ἕξ {\policeFr ÷}}{« De la mesure de deux setiers (contenance métrologique tardive) ; le χοῦς attique contient douze cotyles (soit six setiers). »}
\commScholie{Note métrologique byzantine sur le conge (χοῦς, gén. attique χοός maintenu par R, la Souda λ 394 et Bibi contre χοῶς de Brunck).}
\entreeScholie{127}{350}{ἀρᾶσθε}{Interlinéaire (supra lineam au-dessus d'ἀρᾶσθε, f. 166v, v. 350)}{καταρᾶσθε.}{καταρᾶσθ(ε) {\policeFr ÷}}{« Prononcez des imprécations (maudissez-les). »}
\commScholie{Précise le sens imprécatoire de ἀρᾶσθε avec κακῶς ἀπολέσθαι.}
\entreeScholie{128}{355}{τὰ δ᾽ ἄρισθ᾽ ὅσαις προσήκει / νικᾶν λεγούσαις·}{Marge gauche (extérieure, f. 166v, en face des vv. 355–356)}{ὅσαις δὲ προσήκει τὰ ἄριστα λέγειν, ταύτας νικᾶν λεγούσας.}{ὅσ\expo{αις} δ(ὲ) προσήκ(ει) τὰ ἄριστ(α) λέγ(ειν), ταύτ\expo{ας} νικᾶν λεγούσ\expo{ας} {\policeFr ÷}}{« Et que toutes celles à qui il appartient de donner les meilleurs avis l'emportent par leurs discours aux suffrages de l'assemblée. »}
\commScholie{Paraphrase de la formule finale du héraut.}
\entreeScholie{129}{358}{παραβαίνουσί τε τοὺς / ὅρκους τοὺς νενομισμένους}{Interlinéaire (supra lineam, f. 166v, v. 359)}{ὑφ’ ἡμῶν.}{ὑφ' ἡμ\expo{ῶν} {\policeFr ÷}}{« Par nous (dans le chant du Chœur). »}
\commScholie{Glose interlinéaire de R\expo{2} précisant l'agent implicite du participe parfait passif τοὺς νενομισμένους (v. 358) : les serments solennels institués et reconnus « par nous » (ὑφ᾽ ἡμῶν, la communauté des femmes).}
\entreeScholie{130}{363}{τἀπόρρητά τε τοῖσιν ἐχ- / θροῖς τοῖς ἡμετέροις λέγουσ᾽,}{Marge gauche (extérieure, f. 166v, en face des vv. 363–364)}{δηλοῖ ἐνταῦθα ὅτι λέγουσι καὶ ὁμοίως ἡμῖν τὰ ἀπόρρητα, οὐ μόνον ἀπηγορευμένων.}{δηλοῖ ἐνταῦθ(α) ὅτι λέγουσ(ι) \textbf{ς} ὁμοί\expo{ως} ἡμῖν τὰ ἀπόρρητ(α), οὐ μόν\expo{ον} ἀπηγορευμέν\expo{ων} {\policeFr ÷}}{« Il montre ici que les Attiques disent aussi comme nous τὰ ἀπόρρητα (“les secrets interdits”, ou les marchandises prohibées à l'exportation), et non pas seulement au génitif absolu ἀπηγορευμένων. »}
\commScholie{Remarque grammaticale d'atticiste sur l'emploi substantivé de τὰ ἀπόρρητα.}
\entreeScholie{131}{365}{ἢ Μήδους ἐπάγουσι τῆς / χώρας οὕνεκ᾽ ἐπὶ βλάβῃ,}{Marge gauche (extérieure, f. 166v, en face du v. 365)}{ἀντὶ τοῦ εἰπεῖν πολεμίους εἶπε Μήδους. κατ’ ἐξοχὴν δὲ εἴρηται, ὥσπερ καὶ παρ’ Ὁμήρῳ [Il. Ν, 1] «Ζεὺς ἐπεὶ οὖν Τρῶάς τε καὶ Ἕκτορα.»}{ἀν\auDessus{τ}{ι} τ\expo{οῦ} εἰπ(εῖν) πολεμί\expo{ους} εἶπ(ε) Μήδ\expo{ους} · κατ' ἐξοχ(ὴν) δ(ὲ) εἴρητ(αι), ὥσπ(ερ) \textbf{ς} παρ' Ὁμήρῳ «Ζεὺς ἐπεὶ οὖν Τρῶάς τε \textbf{ς} Ἕκτορα» {\policeFr ÷}}{« Au lieu de dire “les ennemis”, elle a dit “les Mèdes” (les Perses). C'est dit par excellence (κατ' ἐξοχήν), tout comme chez Homère (Iliade XIII, 1) : “Quand donc Zeus eut amené les Troyens et Hector près des nefs.” »}
\commScholie{Rapprochement rhétorique sur la désignation de l'ennemi héréditaire dans l'imprécation civique athénienne (commenté par L. Hauseux, vv 1-641.tex, l. 2982).}
\entreeScholie{132}{370}{ἡμῖν θεοὺς παραστατεῖν}{Interlinéaire (supra lineam, f. 166v, v. 370)}{συμμαχεῖν.}{συμμαχεῖν {\policeFr ÷}}{« Prêter main-forte (être alliés). »}
\commScholie{Glose interlinéaire de R\expo{2} expliquant l'infinitif poétique et religieux παραστατεῖν (« se tenir aux côtés de, assister ») par le verbe courant συμμαχεῖν (« être allié, combattre avec »).}
\entreeScholie{133}{373}{ἔδοξε τῇ βουλῇ τάδε / τῇ τῶν γυναικῶν· Τιμόκλει᾽ ἐπεστάτει,}{Marge gauche (extérieure, f. 166v, en face des vv. 373–374)}{ἐπιστάτις ἐστὶν ἡ Τιμόκλεια. ἀρχὴ δέ τίς ἐστιν ἐπιστάτις, ἥτις ἀνισταμένη συμβουλεύει περὶ πολέμου ἢ εἰρήνης ἢ περὶ ἄλλου τινός.}{ἐπιστάτ(ις) ἐστ(ὶν) ἡ Τιμόκλεια · ἀρχὴ δέ τίς ἐστ(ιν) ἐπιστάτ(ις), ἥτις ἀνισταμέν(η) συμβουλεύ(ει) π(ερὶ) πολέμ\expo{ου} ἢ εἰρήν\expo{ης} ἢ π(ερὶ) ἄλλ\expo{ου} τινός {\policeFr ÷}}{« Timocléia est la présidente (ἐπιστάτις) de l'assemblée des femmes. L'épistate est une magistrature (le président des prytanes tiré au sort pour un jour) qui se lève pour mettre en délibération les questions de guerre, de paix ou toute autre affaire. »}
\commScholie{Repris par la Souda ε 2623. Comme le note Louis Hauseux (Bibi, v. 374, vv 1-641.tex, l. 3023), le nom Τιμόκλει' est attesté à la fois par R\expo{1}, {\policeGr Σ}\expo{R} et la Souda (et maintenu par Bibi), tandis que le papyrus de Satyrus (Ψ\expo{5} = POxy 1176) offre la variante ancienne Ἀρ[χίκλ]ει'.}
\entreeScholie{134}{376}{χρηματίζειν πρῶτα περὶ Εὐριπίδου, / ὅ τι χρὴ παθεῖν ἐκεῖνον·}{Marge gauche (extérieure, f. 166v, en face du v. 376)}{ἐν γὰρ ταῖς ἄλλαις ἡμέραις περὶ τὰς θυσίας γίνονται.}{ἐν γ(ὰρ) τ\expo{αῖς} ἄλλ\expo{αις} ἡμέρ\expo{αις} π(ερὶ) τ\expo{ὰς} θυσί\expo{ας} γίνοντ(αι) {\policeFr ÷}}{« (“Le matin du jour du milieu, où nous avons du loisir”) : car durant les autres journées des Thesmophories, elles sont occupées par les sacrifices. »}
\commScholie{Cité par Alphonse Willems et Louis Hauseux (Annexes.tex, l. 1164 n. 1) : le jour de la Nêsteia (Jeûne) ne comporte pas de banquets sacrificiels, ce qui laisse aux femmes le loisir de tenir leur assemblée.}
\entreeScholie{135}{377}{χρηματίζειν}{Interlinéaire (supra lineam au-dessus de χρηματίζειν, f. 166v, v. 377)}{χρηματίζειν ἐστὶ τὸ λόγους διδόναι ἀλλήλοις.}{χρηματίζ(ειν) ἐστὶ τ(ὸ) λόγ\expo{ους} διδόν(αι) ἀλλήλ\expo{οις} {\policeFr ÷}}{« χρηματίζειν (“traiter les affaires inscrites à l'ordre du jour”) : délibérer entre soi, ouvrir la discussion. »}
\commScholie{Terme technique du droit parlementaire athénien.}
\clearpage
\renewcommand{\currHeadLE}{ΑΡΙΣΤΟΦΑΝΟΥΣ ΘΕΣΜΟΦΟΡΙΑΖΟΥΣΑΙ · CODEX RAVENNAS 429 (\R) · FOLIO 10/30}
\renewcommand{\currHeadRE}{f. 167r (Classense 169r) · vv. 380–422 · Collation diplomatique}
\renewcommand{\currHeadLO}{ΑΡΙΣΤΟΦΑΝΟΥΣ ΘΕΣΜΟΦΟΡΙΑΖΟΥΣΑΙ · CODEX RAVENNAS 429 (\R) · FOLIO 10/30}
\renewcommand{\currHeadRO}{f. 167r (Classense 169r) · vv. 380–422 · Collation diplomatique}
\noindent\begin{adjustbox}{max totalheight=30.4cm, max width=\textwidth, keepaspectratio}\begin{minipage}{\textwidth}
\bandeauRav{f. 167r}{Classense 169r}{380–422}{Folio 10/30 · Recto (page de droite) · 18 scholies \& gloses ({\policeGr Σ} 136–153)}
{\fontsize{8.8pt}{10.2pt}\selectfont
\ligneRav{380}{\auDessus{γυ}{ν}}{ἐγώ. : περίθου νυν τόνδε πρῶτον πρὶν λέγειν·}{\textbf{380} (Γυνή τις) R : \lecon{(Μίκα)} Rogers, cf. v. 760 \espApp \lecon{πρῶτον} R : πρότερον Suid. π 1166}
\ligneRav{}{}{σίγα σιώπα, πρόσεχε τὸν νοῦν· χρέμπτεται γὰρ ἤδη}{}
\ligneRav{}{}{ὅπερ ποιοῦσ᾽ οἱ ῥήτορες. μακρὰν ἔοικε λέξειν·}{}
\ligneRav{383}{καλλιλεξία}{φιλοτιμίᾳ μὲν οὐδεμιᾷ μὰ τὼ θεὼ}{\textbf{383} (\auDessus{γυ}{ν}) \Rac : (καλλιλεξία γυνή) \Rpc, \Mon : (\lecon{Μίκα}) Rogers}
\ligneRav{}{}{λέξουσ᾽ ἀνέστην ὦ γυναῖκες· ἀλλὰ γὰρ}{}
\ligneRav{385}{}{βαρέως φέρω τάλαινα πολὺν ἤδη χρόνον}{}
\ligneRav{386}{}{προπηλακιζομένας ὁρῶς \difference{ἡμᾶς} ὑπὸ}{\textbf{386} ὑμᾶς R : \lecon{ἡμᾶς} Cobet}
\ligneRav{}{}{Εὐριπίδου τοῦ τῆς λαχανοπωλητρίας}{}
\ligneRav{}{}{καὶ πολλὰ καὶ παντοῖ᾽ ἀκουούσας κακά·}{}
\ligneRav{389}{}{τί γὰρ οὗτος ἡμᾶς οὐκ ἐπισμῇ τῶν κακῶν;}{\textbf{389} \lecon{ἡμᾶς} R, Suid. ε 2592 : ἡμῖν Brunck}
\ligneRav{390}{}{ποῦ δ᾽ οὐχὶ διαβέβληχ᾽, ὅπουπερ \difference{ἔμβραχυ}}{\textbf{390} ἐν βραχὺ R : \lecon{ἔμβραχυ} Suid., Σ \italique{Gorg}. et \italique{Hipp.m}.}
\ligneRav{391}{}{εἰσὶν θεαταὶ καὶ τραγῳδοὶ καὶ χοροί,}{\textbf{391} \lecon{τραγῳδοὶ καὶ} R : τραγῳδικοὶ Schol. Plat. \italique{Theag} : τραγῳδία καὶ Suid. οι 128}
\ligneRav{392}{}{τὰς \difference{μοιχοτρόπους}, τὰς \difference{ἀνδρεραστρίας} καλῶν,}{\textbf{392} μυ- R : \lecon{μοι-} Suid. \espApp ἀνδρεα- R : \lecon{ἀνδρερα-} Suid.}
\ligneRav{393}{}{τὰς \difference{οἰνοπίπας}, τὰς προδότιδας, τὰς λάλους,}{\textbf{393} οἰνοπότιδας R : \lecon{-πίπας} Suid. : -πίπους Σ}
\ligneRav{394}{}{τὰς οὐδὲν ὑγιές, τὰς μέγ᾽ \difference{ἀνδράσιν} κακόν·}{\textbf{394} ἀνδράσι R : \lecon{-σιν} Daubuz}
\ligneRav{395}{}{ὥστ᾽ εὐθὺς εἰσιόντες ἀπὸ τῶν ἰκρίων}{}
\ligneRav{}{}{ὑποβλέπουσ᾽ ἡμᾶς σκοποῦνταί τ᾽ εὐθέως}{}
\ligneRav{}{}{μὴ μοιχὸς ἔνδον ᾖ τις ἀποκεκρυμμένος·}{}
\ligneRav{398}{}{δρᾶσαι δ᾽ ἔθ᾽ ἡμῖν οὐδὲν \difference{ὧνπερ} καὶ πρὸ τοῦ}{\textbf{398} ὥσπερ R : \lecon{ὧνπερ} Küster}
\ligneRav{}{}{ἔξεστι· τοιαῦθ᾽ οὗτος ἐδίδαξεν κακὰ}{}
\ligneRav{400}{}{τοὺς ἄνδρας ἡμῶν· ὥστ᾽ ἐάν \difference{τις ξυμπλέκῃ}}{\textbf{400} τις πλέκῃ R : \lecon{τις ξυμπλέκῃ} Willems : δή τις πλέκῃ Austin : alii alia}
\ligneRav{}{}{γυνὴ στέφανον, ἐρᾶν δοκεῖ· κἂν ἐκβάλῃ}{}
\ligneRav{}{}{σκεῦός τι κατὰ τὴν οἰκίαν πλανωμένη,}{}
\ligneRav{403}{}{\difference{ἁνὴρ} ἐρωτᾷ· `τῷ κατέαγεν ἡ χύτρα;}{\textbf{403} ἀνὴρ R : \lecon{ἁνὴρ} Brunck}
\ligneRav{}{}{οὐκ ἔσθ᾽ ὅπως οὐ τῷ Κορινθίῳ ξένῳ.'}{}
\ligneRav{405}{}{κάμνει κόρη τις, εὐθὺς \difference{ἁδελφὸς} λέγει·}{\textbf{405} ἀδελφὸς R : \lecon{ἁδελφὸς} Biset}
\ligneRav{}{}{`τὸ χρῶμα τοῦτό μ᾽ οὐκ ἀρέσκει τῆς κόρης.'}{}
\ligneRav{}{}{εἶἑν, γυνή τις ὑποβαλέσθαι βούλεται}{}
\ligneRav{}{}{ἀποροῦσα παίδων, οὐδὲ τοῦτ᾽ ἔστιν λαθεῖν·}{}
\ligneRav{409}{}{\difference{ἅνδρες} γὰρ ἤδη παρακάθηνται πλησίον·}{\textbf{409} ἄνδρες R : \lecon{ἅνδρες} Brunck.}
\ligneRav{410}{}{πρὸς τοὺς γέροντάς θ᾽ οἳ πρὸ τοῦ τὰς μείρακας}{}
\ligneRav{}{}{ἤγοντο, διαβέβληκεν, ὥστ᾽ οὐδεὶς γέρων}{}
\ligneRav{412}{}{γαμεῖν \difference{ἐθέλει} γυναῖκα διὰ τοὔπος τοδὶ}{\textbf{412} θέλει R : \lecon{ἐθέλει} Kappeyne van de Coppello \italique{Obs. crit.} I.2}
\ligneRav{}{}{`δέσποινα γὰρ γέροντι νυμφίῳ γυνή'·}{}
\ligneRav{}{}{εἶτα διὰ τοῦτον ταῖς γυναικωνίτισιν}{}
\ligneRav{415}{}{σφραγῖδας ἐπιβάλλουσιν ἤδη καὶ μοχλοὺς}{}
\ligneRav{}{}{τηροῦντες ἡμᾶς, καὶ προσέτι Μολοττικοὺς}{}
\ligneRav{417}{}{τρέφουσι μορμολυκεῖα τοῖς \difference{μοιχοῖς} κύνας·}{\textbf{417} μυχοῖς R : \lecon{μοιχοῖς} \Mon}
\ligneRav{}{}{καὶ ταῦτα μὲν ξυγγνώσθ᾽. ἃ δ᾽ ἦν ἡμῖν πρὸ τοῦ}{}
\ligneRav{419}{}{αὐταῖς \difference{ταμιεῦσαι καὶ} προαιρούσαις λαβεῖν}{\textbf{419} ταμιεύεσθαι R : \lecon{ταμιεῦσαι καὶ} Reiske}
\ligneRav{420}{}{ἄλφιτον ἔλαιον οἶνον, οὐδὲ ταῦτ᾽ ἔτι}{}
\ligneRav{}{}{ἔξεστιν. οἱ γὰρ ἄνδρες ἤδη κλῄδια}{}
\ligneRav{}{}{αὐτοὶ φοροῦσι κρυπτὰ κακοηθέστατα}{}
}
\vspace{3pt}{\color{ruleGold}\hrule height 0.6pt}\vspace{3pt}
{\fontsize{8.0pt}{9.4pt}\selectfont
\noindent{\bfseries\scshape\color{rubricRed}Loci paralleli et testimonia (f. 167r, vv. 380–422) :} \textbf{380} Suid. π 1166 περίθου νῦν τόνδε πρότερον \espTest \textbf{381} (χρέμπτεται) \Phryn 126.1 \espTest \textbf{383-4} (φιλοτιμίᾳ -- γυναῖκες) Suidas φ 431 \espTest \textbf{389} Suidas ε 2592 \espTest (ἐπισμῇ) Hesychius ε 5195 ἐπισμῇ \espTest \textbf{390-4} Suidas οι 128 (οἷς θέα pro εἰσὶν θεαταὶ et τραγῳδία pro τραγῳδοὶ) \espTest \textbf{390-1} scholium Plato \italique{Theages} 127c (πῶς pro ποῦ et τραγῳδικοὶ pro τραγῳδοὶ καὶ) = scholium Plato \italique{Gorgias} 457b = scholium Plato \italique{Hippias min.} 365d \espTest (ποῦ -- θεαταὶ) Suidas ε 955 (sed εἰσὶ pro εἰσὶν) \espTest \textbf{390-1} scholium Plato \italique{Theages} 127c = scholium Plato \italique{Gorgias} 457b = scholium Plato \italique{Hippias min.} 365d \espTest \textbf{392} (ἀνδρεραστρίας) Pollux III 71 ; \Phryn 34.18 ; \Eustathe 380.15 \espTest \textbf{393} (οἰνοπίπας) scholium \italique{Aves} 307 \espTest \textbf{395-6} (ὥστ᾽ εὐθὺς -- ἡμᾶς) Suidas ι 275 \espTest \textbf{395} (ἰκρίων) Hesychius ι 501 ἴκρια \espTest \textbf{404} (τῷ Κορινθίῳ ξένῳ) Hesychius κ 3629 Κορίνθιος ξένος \espTest \textbf{413} Menandros \italique{Monostich.} 191 (= Euripides fr. 804.3 ed. Kannicht) ; \Apostol V 93b \espTest \textbf{417} (μορμολυκεῖα) \EtymMagnum 590.53 \espTest \textbf{421-8} (οἱ γὰρ ἄνδρες -- ἐξαψαμένους) Suidas λ 54 (αὐτοὶ, τρεῖς et ᾡκότριψ Εὐριπίδης omissis, Λακωνίκα pro Λακωνίκ᾽, θριπηδέστατ᾽ pro θριπήδεστ᾽ et σφραγῖδα pro σφραγίδια) \espTest \textbf{421-7} (οἱ γὰρ ἄνδρες -- θριπήδεστ᾽) Suidas κ 159 (αὐτοὶ et ᾡκότριψ Εὐριπίδης omissis, Λακωνίκα pro Λακωνίκ᾽ et θριπηδέστατ᾽ pro θριπήδεστ᾽) \espTest \textbf{422} (κακοηθέστατα) \Eustathe 974.37, 1108.25, 1603.52\par\vspace{2.5pt}
{\color{ruleGold}\hrule height 0.35pt}\vspace{2.5pt}
\noindent{\bfseries\scshape\color{rubricRed}Suite (en commentaire) de l'apparat critique et paléographique :} \textbf{Note de personnage descriptive au v. 383 :} Là où \Rac{} avait simplement abrégé (\auDessus{γυ}{ν}), \Rpc{} (suivi par \Mon{}) a développé en marge \lecon{(καλλιλεξία γυνή)} (« la femme au beau langage », glose descriptive passée dans la rubrique de locuteur ; Rogers l'identifie à \textbf{Mika} d'après le v. 760). --- \textbf{Iotacismes οι → υ de \Run{} :} μυχοτρόπους (v. 392) et μυχοῖς (v. 417) pour μοιχο-, redressés par la Souda et \Mon{}. --- \textbf{Restitution par \textbf{Bibi} du fragment de Sappho (fr. 70/125) dans {\policeGr Σ}\expo{R} ad v. 401 :} La scholie du \italique{Ravennas} au v. 401 conserve un précieux vers de Sappho sous la forme corrompue \lecon{αυταόρα αἱστεφανηπλόκουν} ; \textbf{Louis Hauseux (\textbf{Bibi}, \textit{vv 1-641.tex}, ll. 3272–3294) restitue l'éolien αὐτὰ κόραις στεφαναπλόκεο} (ou -πλόκην), montrant que αυταόρα αἱ procède de la chute du κ et d'une médivision de αὐτὰ κόραις. --- \textbf{v. 390 :} ἐν βραχὺ R : \lecon{ἔμβραχυ} Souda, {\policeGr Σ} ad Plat. \italique{Gorg.} et \italique{Hipp. min.} --- \textbf{v. 392 :} μυ- R : \lecon{μοι-} (μοιχοτρόπους) Souda \espApp ἀνδρεα- R (haplographie du second ρ) : \lecon{ἀνδρερα-} (ἀνδρεραστρίας) Souda. --- \textbf{v. 393 :} οἰνοπότιδας R (substitution d'un mot banal au composé comique rare) : \lecon{οἰνοπίπας} Souda (-πίπους {\policeGr Σ}). --- \textbf{v. 400 :} τις πλέκῃ R : \lecon{τις ξυμπλέκῃ} Alphonse Willems (adopté par L. Hauseux ; δή τις πλέκῃ Austin). --- \textbf{{\policeGr Σ}\expo{R} ad v. 401 (Sappho fr. 70/125) :} αυταόρα αἱστεφανηπλόκουν {\policeGr Σ}\expo{R} : \lecon{αὐτὰ κόραις στεφαναπλόκεο} \lecon{Bibi}.\par}
\end{minipage}\end{adjustbox}
\newpage
\renewcommand{\currHeadLE}{SCHOLIES MARGINALES ({\policeGr Σ}\expo{R}) \& GLOSES INTERLINÉAIRES · FOLIO 10/30}
\renewcommand{\currHeadRE}{f. 167r (Classense 169r) · vv. 380–422 · {\policeGr Σ} 136 à {\policeGr Σ} 153}
\renewcommand{\currHeadLO}{SCHOLIES MARGINALES ({\policeGr Σ}\expo{R}) \& GLOSES INTERLINÉAIRES · FOLIO 10/30}
\renewcommand{\currHeadRO}{f. 167r (Classense 169r) · vv. 380–422 · {\policeGr Σ} 136 à {\policeGr Σ} 153}
\bandeauScholies{Folio 10/30 — f. 167r}{Classense 169r}{vv. 380–422}{Recto (page de droite) · 18 entrées ({\policeGr Σ} 136 à {\policeGr Σ} 153)}
{\fontsize{8.8pt}{10.5pt}\selectfont\noindent\textbf{\color{rubricRed}Disposition topographique des marges sur le f. 167r :} Marge extérieure à droite ; marge intérieure à gauche. — \textit{Fac-similé de référence :} \texttt{folios\_individuels\_ff162v-177r/10\_f167r\_classense169r\_vv0380-0422.jpg}.\par\vspace{3pt}}
\entreeScholie{136}{380}{περίθου νυν τόνδε πρῶτον πρὶν λέγειν.}{Marge supérieure et marge droite (extérieure, f. 167r, en haut du recto)}{ἀντὶ τοῦ στέφανον· ἔθος γὰρ ἦν τοῖς λέγουσι στεφανοῦσθαι πρῶτον. τὸ δὲ περίθου, ὅπερ ἡμῖν σύνηθες ἐπίθου λέγειν.}{ἀν\auDessus{τ}{ι} τ\expo{οῦ} στέφαν\expo{ον} · ἔθ\expo{ος} γ(ὰρ) ἦν τ\expo{οῖς} λέγουσ(ι) στεφανοῦσθ(αι) πρῶτ\expo{ον} · τ(ὸ) δ(ὲ) περίθ\expo{ου}, ὅπ(ερ) ἡμῖν σύνηθ\expo{ες} ἐπίθ\expo{ου} λέγ(ειν) {\policeFr ÷}}{« τόνδε : sous-entendu “cette couronne” (de myrte) ; car c'était la coutume que les orateurs prenant la parole à la tribune ceignent d'abord une couronne. Quant à l'attique περίθου (“mets autour de ta tête”), il correspond à ce que nous disons couramment ἐπίθου (“pose sur ta tête”). »}
\commScholie{Note institutionnelle et grammaticale exacte sur la couronne de l'orateur à l'Ecclésia (cf. Assemblée des femmes vv. 131, 148, 163).}
\entreeScholie{137}{383}{φιλοτιμίᾳ μὲν οὐδεμιᾷ μὰ τὼ θεὼ / λέξουσ᾽ ἀνέστην, ὦ γυναῖκες·}{Marge intérieure gauche (f. 167r, en face du v. 383)}{καλλιλεξία. γυνή.}{καλλιλεξία · γυνή {\policeFr ÷}}{« Élégance d'exorde oratoire (καλλιλεξία). — Ici parle la Première Femme (Mika). »}
\commScholie{Double note rhétorique et dramaturgique : le scholiaste admire le pastiche de proème démosthénien dans la bouche de Mika.}
\entreeScholie{138}{386}{προπηλακιζομένας ὁρῶσ᾽ ἡμᾶς [ὑμᾶς Scaliger] ὑπὸ}{Interlinéaire (supra lineam au-dessus de προπηλακιζομένας, f. 167r)}{ὑβριζομένας.}{ὑβριζομέν\expo{ας} {\policeFr ÷}}{« Outragées, couvertes de boue (ὑβριζομένας). »}
\commScholie{Glose du verbe imagé προ-πηλακίζω (de πηλός, « boue »).}
\entreeScholie{139}{389}{τί γὰρ οὗτος ἡμᾶς οὐκ ἐπισμῇ τῶν κακῶν;}{Marge droite (extérieure, f. 167r, grande scholie lexicographique de 10 lignes en face du v. 389)}{ἐπιχρίει, ἐπιξύει. σμῆξαι δέ ἐστι τὸ τὸν ῥύπον ἐπιξύσαι. μεταφορικῶς ἀπὸ τοῦ σμήχειν, οἷον ἐπιτρίβει· σφοδρὰν γάρ τινα τρίψιν καὶ τὴν σμῆξιν εἶναι. τάχα δὲ καὶ τὴν σμώδιγγα εἰρῆσθαι ἀπὸ τούτου. Διόδωρος οὐ μεταφορικὴν, ἀλλ’ ἐξ ἔθους γλωσσηματικὴν, καὶ μὴ εἶναι ἐπιτρίβειν ἀλλὰ τύπτειν ἢ σκώπτειν τὸ σμώχειν. Ἄλλως. ἐπισμῇ: καὶ παρὰ Κρατίνῳ ἐστὶν ἐν Κλεοβουλίναις τὸ ἐπισμῇ. ἔστι δὲ λοιδορία, ἀπὸ τοῦ καθάπτεσθαι τῶν σμωμένων.}{ἐπιχρίει, ἐπιξύει · σμῆξ(αι) δέ ἐστ(ι) τ(ὸ) τ(ὸν) ῥύπον ἐπιξύσαι · μεταφορικ\expo{ῶς} ἀπ(ὸ) τ\expo{οῦ} σμήχε(ιν), οἷ(ον) ἐπιτρίβει· σφοδρ(ὰν) γάρ τινα τρίψιν \textbf{ς} τ(ὴν) σμῆξ(ιν) εἶν(αι) · τάχα δ(ὲ) \textbf{ς} τ(ὴν) σμώδιγγα εἰρῆσθ(αι) ἀπ(ὸ) τούτου · Διόδωρ\expo{ος} οὐ μεταφορικ(ὴν), ἀλλ’ ἐξ ἔθ\expo{ους} γλωσσηματικ(ὴν), \textbf{ς} μὴ εἶν(αι) ἐπιτρίβε(ιν) ἀλλὰ τύπτ(ειν) ἢ σκώπτ(ειν) τ(ὸ) σμώχε(ιν) · ἄλλ\expo{ως} · ἐπισμῇ: \textbf{ς} παρὰ Κρατίνῳ ἐστ(ὶν) ἐν Κλεοβουλίναις τ(ὸ) ἐπισμῇ · ἔστι δ(ὲ) λοιδορία, ἀπ(ὸ) τ\expo{οῦ} καθάπτεσθ(αι) τ\expo{ῶν} σμωμένων {\policeFr ÷}}{« ἐπισμῇ : il nous enduit, il nous racle. σμῆξαι signifie gratter la crasse. Métaphoriquement, de σμήχειν (“frotter rudement”), au sens de “il nous écorche, il nous accable” (ἐπιτρίβει) : car la σμῆξις est un frottement violent, d'où vient peut-être aussi le mot σμῶδιξ (“meurtrissure, zébrure de coup”). Mais Diodore (de Tarse, grammairien du Ier s. av. J.-C.) soutient que l'emploi n'est pas métaphorique, mais idiomatique et glosématique, et que σμώχειν ne signifie pas “frotter”, mais “frapper” ou “railler”. — Autre rédaction : ἐπισμῇ se trouve aussi chez Cratinos dans ses Cléoboulines (fr. 99 K.-A.). C'est une injure, tirée du fait de mordre la peau de ceux que l'on frotte au savon acide. »}
\commScholie{Grande scholie alexandrine (Didyme / Symmaque) citant le grammairien Diodore et les Cléoboulines de Cratinos, reprise intégralement par la Souda ε 2608 et Photius ε 1705.}
\entreeScholie{140}{393}{τὰς οἰνοπίπας, τὰς προδότιδας, τὰς λάλους,}{Marge droite (extérieure, f. 167r, grande scholie de 7 lignes sous Σ 389)}{ὁ Σύμμαχός φησι παρὰ τὸ πιπίζειν τὸν οἶνον. ἐγὼ δὲ οὐχ ὁρῶ τὸ πιπίζειν ἐπὶ ταύτης τῆς ἐννοίας παρὰ τοῖς Ἀττικοῖς κείμενον, ἀλλὰ κατὰ μίμησιν φωνῆς λεγόμενον. ἄμεινον οὖν, ὥσπερ παιδοπίπης λέγεται, οὕτω καὶ οἰνοπίπης. Ἄλλως. τὰ πλεῖστα οἰνοπίπους ἔχει ἀντὶ τοῦ τὰς μεθύσους ἀπὸ τοῦ πιπίζειν, ὅπερ ἐστὶν ἐκμύζειν.}{ὁ Σύμμαχός φησ(ι) παρὰ τ(ὸ) πιπίζε(ιν) τ(ὸν) οἶν(ον) · ἐγὼ δ(ὲ) οὐχ ὁρῶ τ(ὸ) πιπίζε(ιν) ἐπὶ ταύτης τ\expo{ῆς} ἐννοίας παρὰ τ\expo{οῖς} Ἀττικ\expo{οῖς} κείμεν(ον), ἀλλὰ κατὰ μίμησ(ιν) φων\expo{ῆς} λεγόμεν(ον) · ἄμειν(ον) οὖν, ὥσπ(ερ) παιδοπίπης λέγετ(αι), οὕτω \textbf{ς} οἰνοπίπης · ἄλλ\expo{ως} · τὰ πλεῖστα οἰνοπίπους ἔχ(ει) ἀν\auDessus{τ}{ι} τ\expo{οῦ} τ\expo{ὰς} μεθύσους ἀπ(ὸ) τ\expo{οῦ} πιπίζε(ιν), ὅπερ ἐστ(ὶν) ἐκμύζε(ιν) {\policeFr ÷}}{« Symmaque dit que le mot est formé sur πιπίζειν τὸν οἶνον (“siroter le vin en pépiant”). Pour ma part, je ne vois pas que πιπίζειν soit employé en ce sens chez les auteurs attiques, où il désigne par onomatopée le pépiement des oiseaux. Il vaut donc mieux expliquer que, de même qu'on dit παρθενοπίπης / παιδοπίπης (“qui lorgne les garçons”, de ὀπιπεύω), on a dit οἰνοπίπης (“qui lorgne le vin”). — Autre rédaction : La plupart des exemplaires portent οἰνοπίπους, au sens d'ivrognesses, tiré de πιπίζειν qui signifie sucer, siroter. »}
\commScholie{Cas magistral analysé par Louis Hauseux (Bibi, v. 393, vv 1-641.tex, ll. 3086-3105) : alors que le scribe R\expo{1} a écrit la banalisation οἰνοπότιδας, cette scholie (confirmée par la Souda οι 128 et la scholie aux Oiseaux v. 307) nous permet de restituer avec certitude la leçon rare authentique d'Aristophane : τὰς οἰνοπίπας (« celles qui lorgnent le vin », adopté par Bibi).}
\entreeScholie{141}{395}{ἀπὸ τῶν ἰκρίων}{Marge droite (extérieure, f. 167r, en face du v. 395)}{ὡς ἔτι ἰκρίων ὄντων ἐν τῷ θεάτρῳ καὶ ἐν ταῖς ἐκκλησίαις ἐπὶ ξύλων καθημένων. πρὶν γὰρ γενέσθαι τὸ θέατρον, ξύλα ἐδέσμευον καὶ οὕτως ἐθεώρουν.}{ὡς ἔτι ἰκρί\expo{ων} ὄντ\expo{ων} ἐν τῷ θεάτρ(ῳ) \textbf{ς} ἐν τ\expo{αῖς} ἐκκλησί\expo{αις} ἐπὶ ξύλ\expo{ων} καθημέν\expo{ων} · πρὶν γ(ὰρ) γενέσθ(αι) τ(ὸ) θέατρ\expo{ον}, ξύλα ἐδέσμευ\expo{ον} \textbf{ς} οὕτ\expo{ως} ἐθεώρ\expo{ουν} {\policeFr ÷}}{« Comme s'il y avait encore des gradins de bois (ἴκρια) au théâtre et que dans les assemblées on s'asseyait sur des bancs de charpente ; car avant la construction du théâtre de pierre, on assemblait des poutres de bois et c'est ainsi qu'on assistait au spectacle. »}
\commScholie{Note archéologique célèbre sur les gradins de bois (ἴκρια) du théâtre de Dionysos avant la reconstruction en pierre de Lycurgue (reprise par la Souda ι 275).}
\entreeScholie{142}{399}{τοιαῦθ᾽ οὗτος ἐδίδαξεν κακὰ / τοὺς ἄνδρας ἡμῶν·}{Interlinéaire (supra lineam, f. 167r, v. 397)}{ἀντὶ τοῦ ἐν ὅλοις αὐτοῦ τοῖς δράμασι.}{ἀν\auDessus{τ}{ι} τ\expo{οῦ} ἐν ὅλ\expo{οις} αὐτ\expo{οῦ} τ\expo{οῖς} δράμασ(ι) {\policeFr ÷}}{« Partout (ὅπουπερ ἔμβραχυ, vv. 390 sqq.) : dans toutes ses pièces sans exception. »}
\commScholie{Glose explicative de R\expo{2} résumant la portée générale du réquisitoire de Mika (vv. 389–400) : ce n'est pas dans une seule pièce, mais « dans l'ensemble de ses drames » (ἐν ὅλοις αὐτοῦ τοῖς δράμασι) qu'Euripide enseigne aux maris à soupçonner les femmes.}
\entreeScholie{143}{401}{ὥστ᾽ ἐάν τις ξυμπλέκῃ [πλέκῃ R] / γυνὴ στέφανον, ἐρᾶν δοκεῖ·}{Marge droite (extérieure, f. 167r, en face du v. 401)}{νεωτέρων καὶ ἐρωτικῶν τὸ στεφανοπλοκεῖν. πρὸς τὸ ἔθος, ὅτι ἐστεφανηπλόκουν αἱ παλαιαί· Σαπφώ· «αὐτὰ δ’ ὥραις ἐστεφαναπλόκην [αυταόρα αιστεφανηπλόκουν R].»}{νεωτέρ\expo{ων} \textbf{ς} ἐρωτικ\expo{ῶν} τ(ὸ) στεφανοπλοκ(εῖν) · πρ(ὸς) τ(ὸ) ἔθ\expo{ος}, ὅτι ἐστεφανηπλόκ\expo{ουν} αἱ παλαιαί · Σαπφὼ «αυταόρα αιστεφανηπλόκ\expo{ουν}» {\policeFr ÷}}{« Tresser des couronnes est le propre des jeunes gens et des amoureux. — Autre note : En référence à l'usage ancien selon lequel les femmes d'autrefois tressaient des couronnes ; Sappho (fr. 125 Voigt) : “Moi-même, en ma jeunesse, je tressais des couronnes (αὐτὰ δ' ὥραις ἐστεφαναπλόκην / αὐτὰ κόραις στεφαναπλόκεον).” »}
\commScholie{Joyau philologique commenté par Louis Hauseux (vv 1-641.tex, ll. 3108-3115) : c'est cette unique scholie du Ravennas (reprise dans l'Etymologicum Genuinum) qui nous a conservé le fragment 125 de Sappho !}
\entreeScholie{144}{402}{κἂν ἐκβάλῃ / σκεῦός τι κατὰ τὴν οἰκίαν πλανωμένη,}{Interlinéaire (supra lineam au-dessus de πλανωμένη, f. 167r, v. 402)}{περιερχομένη.}{περιερχομέν(η) {\policeFr ÷}}{« En allant et venant (à travers la maison). »}
\commScholie{Glose lexicale de R\expo{2} expliquant l'emploi attique affaibli du participe πλανωμένη (« allant et venant au hasard dans la maison », κατὰ τὴν οἰκίαν) par περιερχομένη.}
\entreeScholie{145}{403}{ἁνὴρ ἐρωτᾷ· τῷ κατέαγεν ἡ χύτρα;}{Marge droite (extérieure, f. 167r, en face du v. 403)}{ἀντὶ τοῦ εἰς τίνα τὸν νοῦν ἐχούσης σου κατέαγεν ἡ χύτρα. ὡς διὰ τὸ ἐρᾶν μὴ προσεχουσῶν ἐπεφώνουν ταῦτα οἱ ἄνδρες, εἰ ἡ γυνή τι ἐκ τῶν χειρῶν ἀφίοι.}{ἀν\auDessus{τ}{ι} τ\expo{οῦ} εἰς τίνα τ(ὸν) νοῦν ἐχούσ\expo{ης} σ\expo{ου} κατέαγ(εν) ἡ χύτρα · ὡς διὰ τ(ὸ) ἐρᾶν μὴ προσεχουσ\expo{ῶν} ἐπεφών\expo{ουν} ταῦτ(α) οἱ ἄνδρ\expo{ες}, εἰ ἡ γυνή τι ἐκ τ\expo{ῶν} χειρ\expo{ῶν} ἀφίοι {\policeFr ÷}}{« “En l'honneur de qui la marmite s'est-elle cassée ?” (τῷ = τίνι) : c'est-à-dire “à quel galant pensais-tu donc quand la marmite s'est brisée ?” Comme les femmes amoureuses ont l'esprit ailleurs, les maris lançaient cette boutade dès que leur femme laissait tomber un objet des mains. »}
\commScholie{Excellente explication du datif d'intérêt τῷ (= τίνι) suivie au v. 404 de la réponse parodique : οὐκ ἔσθ' ὅπως οὐ τῷ Κορινθίῳ ξένῳ.}
\entreeScholie{146}{404}{οὐκ ἔσθ᾽ ὅπως οὐ τῷ Κορινθίῳ ξένῳ.}{Marge droite (extérieure, f. 167r, sous Σ 403)}{ἐκ Σθενεβοίας Εὐριπίδου. διαβάλλει δὲ ὡς μοιχοὺς τοὺς Κορινθίους.}{ἐκ Σθενεβοί\expo{ας} Εὐριπίδ\expo{ου} · διαβάλλ(ει) δ(ὲ) ὡς μοιχ\expo{οὺς} τ\expo{οὺς} Κορινθί\expo{ους} {\policeFr ÷}}{« Tiré de la Sthénébée d'Euripide (fr. 664 Kannicht, où l'hôte corinthien est Bellérophon). Et il raille en même temps les Corinthiens comme des séducteurs d'épouses. »}
\commScholie{Identifie la tragédie perdue d'Euripide (Sthénébée), où la reine d'Argos s'éprenait de son hôte venu de Corinthe, Bellérophon.}
\entreeScholie{147}{409}{εἶἑν, γυνή τις ὑποβαλέσθαι βούλεται / ἀποροῦσα παίδων, οὐδὲ τοῦτ᾽ ἔστιν λαθεῖν· / ἅνδρες γὰρ ἤδη παρακάθηνται πλησίον·}{Interlinéaire (supra lineam, f. 167r, v. 409)}{ὅτε τίκτουσιν αἱ γυναῖκες.}{ὅτε τίκτουσ(ιν) αἱ γυναῖκ\expo{ες} {\policeFr ÷}}{« Lorsque les femmes feignent d'être en couches. »}
\commScholie{Glose explicative sur la supposition d'enfant (allusion au Phénix ou à l'Égée d'Euripide).}
\entreeScholie{148}{410}{πρὸς τοὺς γέροντας, οἳ πρὸ τοῦ τὰς μείρακας / ἤγοντο,}{Marge droite (extérieure, f. 167r, en face des vv. 410–411)}{πρὸς τοὺς γέροντας διαβέβληκεν, οἳ πρὸ τοῦ τὰς μείρακας ἤγοντο.}{πρ(ὸς) τ\expo{οὺς} γέροντ\expo{ας} διαβέβληκ(εν), οἳ πρὸ τ\expo{οῦ} τ\expo{ὰς} μείρακ\expo{ας} ἤγοντ(ο) {\policeFr ÷}}{« Il nous a calomniées auprès des vieillards qui, autrefois, épousaient de toutes jeunes filles. »}
\commScholie{Paraphrase syntaxique de R\expo{2} rétablissant l'ordre des mots (ὑπερβατόν) des vv. 410–411 en rattachant le complément prépositionnel initial πρὸς τοὺς γέροντας au verbe principal διαβέβληκεν.}
\entreeScholie{149}{412}{γαμεῖν ἐθέλει γυναῖκα διὰ τοὔπος τοδὶ / «δέσποινα γὰρ γέροντι νυμφίῳ γυνή»·}{Marge droite (extérieure, f. 167r, en face des vv. 412–413)}{ἀντὶ τοῦ διὰ τὸν ἴαμβον τοῦτον. πᾶν γὰρ μέτρον οἱ παλαιοὶ ἔπος ἐκάλουν.}{ἀν\auDessus{τ}{ι} τ\expo{οῦ} διὰ τ(ὸν) ἴαμβ\expo{ον} τοῦτ\expo{ον} · πᾶν γ(ὰρ) μέτρ\expo{ον} οἱ παλαιοὶ ἔπ\expo{ος} ἐκάλ\expo{ουν} {\policeFr ÷}}{« διὰ τοῦτο τὸ ἔπος (“à cause de ce vers”) : au sens de “à cause de ce trimètre iambique” (du Phénix d'Euripide, fr. 804 K.) ; car les anciens appelaient ἔπος tout vers métrique (et pas seulement l'hexamètre épique). »}
\commScholie{Remarque terminologique d'Aristarque (fréquente dans les scholies d'Aristophane, cf. Grenouilles v. 862) sur le sens classique d'ἔπος = vers iambique ou lyrique.}
\entreeScholie{150}{416}{Μολοττικοὺς / τρέφουσι κύνας,}{Interlinéaire (supra lineam au-dessus de Μολοττικοὺς, f. 167r)}{ἠπειρωτικούς.}{ἠπειρωτικ\expo{ούς} {\policeFr ÷}}{« Des chiens d'Épire (les molosses épirotes, célèbres pour leur férocité). »}
\commScholie{Glose géographique exacte de R\expo{2} identifiant les fameux molosses de garde (Μολοττικοὺς κύνας) comme des chiens d'Épire (ἠπειρωτικούς), réputés dans toute la Grèce pour leur férocité.}
\entreeScholie{151}{417}{μορμολυκεῖα τοῖς μοιχοῖς,}{Interlinéaire (supra lineam au-dessus de μορμολυκεῖα, f. 167r)}{φόβον ἀπὸ τῶν προσωπείων.}{φόβ\expo{ον} ἀπ(ὸ) τ\expo{ῶν} προσωπεί\expo{ων} {\policeFr ÷}}{« Des épouvantails (inspirant la terreur), métaphore tirée des masques grimaçants de théâtre (de Mormô, croque-mitaine féminin). »}
\commScholie{Glose lexicographique de R\expo{2} reprise par la Souda μ 1252, expliquant l'apposition métaphorique μορμολυκεῖα (« épouvantails à amants », de Μορμώ) par analogie avec les masques grimaçants de théâtre (ἀπὸ τῶν προσωπαίων).}
\entreeScholie{152}{419}{αὐταῖς ταμιεῦσαι καὶ [ταμιεύεσθαι R] προαιρούσαις λαβεῖν}{Interlinéaire (supra lineam au-dessus de προαιρούσας λαβεῖν, f. 167r)}{ἐκ τοῦ ταμείου λαβεῖν.}{ἐκ τ\expo{οῦ} ταμεί\expo{ου} λαβ(εῖν) {\policeFr ÷}}{« Tirer du cellier (prendre dans la réserve domestique, promere). »}
\commScholie{Cité et approuvé par Alphonse Willems et Louis Hauseux (Annexes.tex, l. 821) pour défendre le texte de R προαιρούσας λαβεῖν.}
\entreeScholie{153}{420}{ἄλφιτον, ἔλαιον, οἶνον·}{Marge intérieure gauche (bas du f. 167r, en face du v. 420)}{ὅτι ἑνικῶς τὸ ἄλφιτον.}{ὅτι ἑνικ\expo{ῶς} τ(ὸ) ἄλφιτ\expo{ον} {\policeFr ÷}}{« Notez l'emploi remarquable de ἄλφιτον (“farine d'orge”) au singulier (au lieu du pluriel habituel τὰ ἄλφιτα). »}
\commScholie{Remarque grammaticale d'atticiste signalant l'emploi rare du singulier collectif ἄλφιτον.}
\clearpage
\renewcommand{\currHeadLE}{ΑΡΙΣΤΟΦΑΝΟΥΣ ΘΕΣΜΟΦΟΡΙΑΖΟΥΣΑΙ · CODEX RAVENNAS 429 (\R) · FOLIO 11/30}
\renewcommand{\currHeadRE}{f. 167v (Classense 169v) · vv. 423–465 · Collation diplomatique}
\renewcommand{\currHeadLO}{ΑΡΙΣΤΟΦΑΝΟΥΣ ΘΕΣΜΟΦΟΡΙΑΖΟΥΣΑΙ · CODEX RAVENNAS 429 (\R) · FOLIO 11/30}
\renewcommand{\currHeadRO}{f. 167v (Classense 169v) · vv. 423–465 · Collation diplomatique}
\noindent\begin{adjustbox}{max totalheight=30.4cm, max width=\textwidth, keepaspectratio}\begin{minipage}{\textwidth}
\bandeauRav{f. 167v}{Classense 169v}{423–465}{Folio 11/30 · Verso (page de gauche) · 12 scholies \& gloses ({\policeGr Σ} 154–165)}
{\fontsize{8.4pt}{9.7pt}\selectfont
\ligneRav{423}{\auDessus{γυ}{ν}}{Λακωνίκ᾽ ἄττα, τρεῖς ἔχοντα γομφίους·}{}
\ligneRav{}{}{πρὸ τοῦ μὲν οὖν ἦν ἀλλ᾽ ὑποῖξαι τὴν θύραν}{}
\ligneRav{425}{}{ποιησαμέναισι δακτύλιον τριωβόλου,}{}
\ligneRav{426}{}{νῦν δ᾽ οὗτος αὐτοὺς \difference{ᾡκότριψ} Εὐριπίδης}{\textbf{426} ὡκότριψ R : \lecon{ᾡκότριψ} Bentley : οἱκότριψ Kühner-Blass}
\ligneRav{427}{}{\difference{ἐδίδαξε} θριπήδεστ᾽ ἔχειν σφραγίδια}{\textbf{427} ἐδίδαξεν R : \lecon{ἐδίδαξε} Suid. λ 64}
\ligneRav{}{}{ἐξαψαμένους. νῦν οὖν ἐμοὶ τούτῳ δοκεῖ}{}
\ligneRav{429}{}{ὄλεθρόν τιν᾽ ἡμᾶς κυρκανᾶν \difference{ἁμωσγέπως},}{\textbf{429} ἀμωσγέπως R, Suid. : \lecon{ἁμωσγέπως} Blaydes}
\ligneRav{430}{}{ἢ φαρμάκοισιν ἢ μιᾷ γέ τῳ τέχνῃ,}{}
\ligneRav{}{}{ὅπως ἀπολεῖται. ταῦτ᾽ ἐγὼ φανερῶς λέγω,}{}
\ligneRav{}{}{τὰ δ᾽ ἄλλα μετὰ τῆς γραμματέως συγγράψομαι·}{}
\ligneRav{433}{χ/}{\hspace{1.1em}\difference{οὔπω} (οὐπώ\expo{πο}τε) ταύτης ἤκουσα}{\textbf{433} οὔπω τε \Rac : οὐπώ\expo{πο}τε \Rpc : \lecon{οὔπω} Dindorf}
\ligneRav{}{}{\hspace{1.1em}πολυπλοκωτέρας γυναικὸς}{}
\ligneRav{}{}{\hspace{1.1em}οὐδὲ δεινότερον λεγούσης·}{}
\ligneRav{}{}{\hspace{1.1em}πάντα γὰρ λέγει δίκαια,}{}
\ligneRav{437}{}{\hspace{1.1em}†\difference{πάσας} δ᾽ \difference{ἰδέας} ἐξήτασεν†,}{\textbf{437} \lecon{πάσας} εἰδέας \lecon{ἐξήτασεν} R (sed \lecon{ἰδέας} Suid. ε 22) : πᾶσαν ἰδέαν ἐξετάζει Hermann}
\ligneRav{438}{}{\hspace{1.1em}πάντα δ᾽ \difference{ἐβάστασε} φρενὶ πυκνῶς τε}{\textbf{438} ἐβάστασεν R : \lecon{-σε} Suid.}
\ligneRav{439}{}{\hspace{1.1em}ποικίλους λόγους \difference{ἀνηῦρεν}}{\textbf{439} ἀνεῦρεν R, Suid. : \lecon{ἀνηῦρεν} Meineke}
\ligneRav{}{}{\hspace{1.1em}εὖ διεζητημένους·}{}
\ligneRav{440}{}{\hspace{1.1em}ὥστ᾽ ἂν εἰ λέγοι παρ᾽ αὐτῆς}{\textbf{440} \lecon{αὐτῆς} R : αὐτὴν Zanetti}
\ligneRav{441}{}{\hspace{1.1em}\difference{Ξενοκλέης} ὁ Καρκίνου, δο-}{\textbf{441} Ξενοκλῆς R : \lecon{Ξενοκλέης} Bentley}
\ligneRav{}{}{\hspace{1.1em}κεῖν ἂν αὐτόν, ὡς ἐγᾦμαι,}{}
\ligneRav{}{}{\hspace{1.1em}πᾶσιν ὑμῖν}{}
\ligneRav{}{}{\hspace{1.1em}ἄντικρυς μηδὲν λέγειν·}{}
\ligneRav{443}{γυνὴ ἑτέρα}{ὀλίγων \difference{ἕνεκα καὐτὴ} παρῆλθον ῥημάτων·}{\textbf{443} γυνὴ ἑτέρα \Rpc \espApp ἕνεκ᾽ αὐτὴ R : \lecon{ἕνεκα καὐτὴ} Porson}
\ligneRav{}{}{τὰ μὲν γὰρ ἄλλ᾽ αὕτη κατηγόρηκεν εὖ·}{}
\ligneRav{445}{}{ἃ δ᾽ ἐγὼ πέπονθα, ταῦτα λέξαι βούλομαι·}{}
\ligneRav{446}{}{ἐμοὶ γὰρ \difference{ἁνὴρ} ἀπέθανεν μὲν \difference{ἔγκυος}}{\textbf{446} ἀνὴρ R : \lecon{ἁνὴρ} Brunck \espApp ἐν Κύπρῳ R : \lecon{ἔγκυος} Bibi}
\ligneRav{}{}{παιδάρια πέντε καταλιπών, ἁγὼ μόλις}{}
\ligneRav{}{}{στεφανηπλοκοῦσ᾽ ἔβοσκον ἐν ταῖς μυρρίναις·}{}
\ligneRav{}{}{τέως μὲν οὖν ἀλλ᾽ ἡμικάκως ἐβοσκόμην·}{}
\ligneRav{450}{}{νῦν δ᾽ οὗτος ἐν \difference{ταῖσιν} τραγῳδίαις ποιῶν}{\textbf{450} ταῖσι R : \lecon{ταῖσιν} Brunck}
\ligneRav{}{}{τοὺς ἄνδρας ἀναπέπεικεν οὐκ εἶναι θεούς·}{}
\ligneRav{}{}{ὥστ᾽ οὐκέτ᾽ ἐμπολῶμεν οὐδ᾽ εἰς ἥμισυ·}{}
\ligneRav{453}{}{νῦν οὖν \difference{ἁπάσαισιν} παραινῶ καὶ λέγω}{\textbf{453} ἁπάσαισι R : \lecon{ἁπάσαισιν} Gellius \italique{Ν.Α.} 15.20}
\ligneRav{}{}{τοῦτον κολάσαι τὸν ἄνδρα πολλῶν οὕνεκα·}{}
\ligneRav{455}{}{ἄγρια γὰρ ἡμᾶς ὦ γυναῖκες δρᾷ κακά,}{}
\ligneRav{456}{}{ἅτ᾽ ἐν ἀγρίοισι \difference{τοῖς} λαχάνοις αὐτὸς τραφείς·}{\textbf{456} \lecon{τοῖς} Plut., \italique{Mor.}, 853C, Gellius \italique{N.Α.}15.20 : om. R}
\ligneRav{}{}{ἀλλ᾽ εἰς ἀγορὰν ἄπειμι· δεῖ γὰρ ἀνδράσιν}{}
\ligneRav{}{}{πλέξαι στεφάνους συνθηματιαίους εἴκοσιν·}{}
\ligneRav{}{χορ/}{\hspace{1.1em}ἕτερον αὖ τι λῆμα τοῦτο}{}
\ligneRav{460}{}{\hspace{1.1em}κομψότερον ἔτ᾽ ἢ τὸ πρότερον}{}
\ligneRav{}{}{\hspace{1.1em}ἀναπέφηνεν·}{}
\ligneRav{461}{}{\hspace{1.1em}οἷα κατεστωμύλατο}{\textbf{461} \lecon{κατεστωμύλατο} R : κἀστωμύλατο Dobree}
\ligneRav{462}{}{\hspace{1.1em}οὐκ \difference{ἄκαιρα}, φρένας ἔχουσα}{\textbf{462} ἄκερα R : \lecon{ἄκαιρα} Zanetti}
\ligneRav{463}{}{\hspace{1.1em}καὶ πολύπλοκον νόημ᾽, οὐδ᾽}{\textbf{463} \lecon{πολύπλοκον} R : τι πολύπλοκον Dobree : πολύστροφον Hermann}
\ligneRav{}{}{\hspace{1.1em}ἀσύνετ᾽ ἀλλὰ πιθανὰ πάντα·}{}
\ligneRav{465}{}{\hspace{1.1em}δεῖ δὲ ταύτης τῆς ὕβρεως ἡ-}{}
\ligneRav{}{}{\hspace{1.1em}μῖν τὸν ἄνδρα}{}
\ligneRav{}{}{\hspace{1.1em}περιφανῶς δοῦναι δίκην·}{}
}
\vspace{3pt}{\color{ruleGold}\hrule height 0.6pt}\vspace{3pt}
{\fontsize{8.0pt}{9.4pt}\selectfont
\noindent{\bfseries\scshape\color{rubricRed}Loci paralleli et testimonia (f. 167v, vv. 423–465) :} \textbf{423} (γομφίους) \Phryn 61.1 \espTest \textbf{426} (οἰκότριψ = οἰκοτραφής) Moeris ο 25 ; (ᾡκότριψ) Hesychius ο 265-266 ; \AelDion ο 8-9 \espTest \textbf{428-9} (ἐμοὶ -- ἁμωσγέπως) Suidas κ 2771 (τινα pro τιν᾽ ἡμᾶς) \espTest \textbf{429} (κυρκανᾶν) Photius κ 1254 κυρκανᾶν = Hesychius κ 4691 κυρκανᾷ ; \Eustathe 1213.50 \espTest (ἁμωσγέπως) Hesychius α 4182 ἁμωσγέπως ; \Syn α 449 = Suidas α 1645 \espTest \textbf{432} (μετὰ τῆς γραμματέως συγγράψομαι) \Eustathe 1484.53 \espTest (τῆς γραμματέως) Pollux IV 19 \espTest \textbf{437-9} (πάσας -- ἀνηῦρεν) Suidas ε 22 πάσας δ' ἰδέας ἐξήτασε πάντα δ' ἐβάστασε φρενί, πυκνῶς τε ποικίλως λόγους ἀνεῦρεν \espTest \textbf{438} (πυκνῶς) Hesychius π 4335 πυκνῶς \espTest \textbf{442} (ἄντικρυς) Photius α 2210 ἄντικρυς ; \AelDion α 107 \espTest \textbf{446} Hesychius ε 239 ἔγκυος· ἐγγαστρωμένη γυνή \espTest \textbf{447-9} (μόλις -- τέως μὲν οὖν ἀλλ᾽) scholium\expo{T} Plato \italique{Hipparch.} 229d (μυρίναις pro μυρρίναις et μέν οὐκ ἄν pro μὲν οὖν ἀλλ᾽) \espTest \textbf{448} (μυρρίνη = μυρσίνη) Moeris μ 23 \espTest \textbf{449} Suidas τ 321 \espTest (ἡμικάκως) Pollux VI 161 \espTest \textbf{453-6} Aulus Gellius XV 20.7 \espTest \textbf{455-6} Plut. \italique{Moral.}, \italique{Compar. Arist. et Menandr.} (853c) \espTest \textbf{457-8} (δεῖ γὰρ -- εἴκοσιν) Suidas σ 1591 (sed ἀνδράσι et εἴκοσι pro ἀνδράσιν et εἴκοσιν) \espTest \textbf{458} Athenaeus XV 680c \espTest (στεφάνους συνθηματιαίους) Pollux VII 200 \espTest (συνθηματιαίους) Hesychius σ 2624 συνθηματιαῖοι \espTest \textbf{459-60} (ἕτερον -- ἀναπέφηνεν) Suidas λ 441 (ἀναπέφηνε pro ἀναπέφηνεν)\par\vspace{2.5pt}
{\color{ruleGold}\hrule height 0.35pt}\vspace{2.5pt}
\noindent{\bfseries\scshape\color{rubricRed}Suite (en commentaire) de l'apparat critique et paléographique :} \textbf{Correction au v. 433 :} \Rac{} avait écrit οὔπω τε, corrigé par \Rpc{} en οὐπώποτε (Dindorf restitue οὔπω). --- \textbf{Note d'interlocuteur au v. 443 :} \Rpc{} inscrit γυνὴ ἑτέρα. --- \textbf{Conjecture remarquable de \textbf{Bibi} au v. 446 :} ἀνὴρ μὲν ἐν Κύπρῳ R {\policeFr →} \lecon{ἁνὴρ μὲν ἔγκυος} (ἔγκυος se rapportant à la femme ou construit avec l'idée que son mari l'a laissée enceinte de son 5\ieme{} enfant avant de mourir, au lieu de la localisation géographique hors de propos ἐν Κύπρῳ). --- \textbf{v. 437 :} \lecon{πάσας} εἰδέας \lecon{ἐξήτασεν} R (mais \lecon{ἰδέας} Souda ε 22, sans iotacisme ; πᾶσαν ἰδέαν ἐξετάζει Hermann). --- \textbf{v. 443 :} \lecon{γυνὴ ἑτέρα} \Rpc{} \espApp ἕνεκ' αὐτὴ R (haplographie de κα-) : \lecon{ἕνεκα καὐτὴ} Porson. --- \textbf{v. 446 :} ἀνὴρ R : \lecon{ἁνὴρ} Brunck \espApp ἐν Κύπρῳ R : \lecon{ἔγκυος} \lecon{Bibi}. --- \textbf{v. 456 :} \lecon{τοῖς} Plutarque, \italique{Moralia} 853C, Aulu-Gelle 15.20 : om. R (devant λαχάνοις).\par}
\end{minipage}\end{adjustbox}
\newpage
\renewcommand{\currHeadLE}{SCHOLIES MARGINALES ({\policeGr Σ}\expo{R}) \& GLOSES INTERLINÉAIRES · FOLIO 11/30}
\renewcommand{\currHeadRE}{f. 167v (Classense 169v) · vv. 423–465 · {\policeGr Σ} 154 à {\policeGr Σ} 165}
\renewcommand{\currHeadLO}{SCHOLIES MARGINALES ({\policeGr Σ}\expo{R}) \& GLOSES INTERLINÉAIRES · FOLIO 11/30}
\renewcommand{\currHeadRO}{f. 167v (Classense 169v) · vv. 423–465 · {\policeGr Σ} 154 à {\policeGr Σ} 165}
\bandeauScholies{Folio 11/30 — f. 167v}{Classense 169v}{vv. 423–465}{Verso (page de gauche) · 12 entrées ({\policeGr Σ} 154 à {\policeGr Σ} 165)}
{\fontsize{8.8pt}{10.5pt}\selectfont\noindent\textbf{\color{rubricRed}Disposition topographique des marges sur le f. 167v :} Marge extérieure à gauche ; marge intérieure à droite. — \textit{Fac-similé de référence :} \texttt{folios\_individuels\_ff162v-177r/11\_f167v\_classense169v\_vv0423-0465.jpg}.\par\vspace{3pt}}
\entreeScholie{154}{423}{κλῄδια / αὐτοὶ φοροῦσι κρυπτὰ κακοηθέστατα / Λακωνίκ᾽ ἄττα, τρεῖς ἔχοντα γομφίους·}{Marge gauche (extérieure, f. 167v, grande scholie de 7 lignes en haut du verso)}{τὰ γὰρ ἀρχαῖα μονοβάλανά φασιν εἶναι. γομφίους δὲ, οὓς ἡμεῖς ὀδόντας. Μένανδρος Μισουμένῳ «Λακωνικὴ κλείς ἐστιν ὡς ἔοικέ μοι περιοιστέα.» καί φασιν ὅτι ἔξωθεν περικλείεται μοχλοῦ παρατιθεμένου ἤ τινος τοιούτου, ὥστε τοῖς ἔνδον μὴ εἶναι ἀνοῖξαι. περιβόητοί εἰσιν αἱ Λακωνικαὶ κλεῖδες.}{τὰ γ(ὰρ) ἀρχαῖα μονοβάλανά φασ(ιν) εἶν(αι) · γομφίους δ(ὲ), οὓς ἡμεῖς ὀδόντ\expo{ας} · Μένανδρ\expo{ος} Μισουμένῳ «Λακωνικὴ κλείς ἐστ(ιν) ὡς ἔοικέ μοι περιοιστέα · » καί φασ(ιν) ὅτι ἔξωθ(εν) περικλείετ(αι) μοχλ\expo{οῦ} παρατιθεμένου ἤ τιν\expo{ος} τοιούτου, ὥστε τ\expo{οῖς} ἔνδ(ον) μὴ εἶν(αι) ἀνοῖξ(αι) · περιβόητοί εἰσ(ιν) αἱ Λακωνικαὶ κλεῖδες {\policeFr ÷}}{« Car les serrures anciennes, dit-on, n'avaient qu'un seul pêne (μονοβάλανα). Il appelle γομφίους (“molaires”) ce que nous appelons les “dents” (panneton à trois dents) de la clé. Ménandre dit dans le Haï (Μισούμενος, fr. 8 Koerte) : “C'est une clé laconienne, à ce qu'il semble, qu'il me faut trimbaler.” Et l'on explique que la serrure laconienne se verrouille de l'extérieur au moyen d'un loquet ou mécanisme tel que ceux qui sont enfermés à l'intérieur ne peuvent plus l'ouvrir. Les clés laconiennes sont célèbres. »}
\commScholie{Scholie archéologique et littéraire majeure (reprise par la Souda λ 64) conservant une citation du Misoumenos de Ménandre et expliquant le mécanisme de la serrure laconienne à trois goupilles.}
\entreeScholie{155}{425}{πρὸ τοῦ μὲν οὖν ἦν ἀλλ᾽ ὑποῖξαι τὴν θύραν / ποιησαμέναισι δακτύλιον τριωβόλου,}{Interlinéaire (supra lineam, f. 167v, v. 425)}{σφραγίδιον παραπλήσιον ᾧ ἐσφράγιζεν ὁ ἀνήρ.}{σφραγίδι\expo{ον} παραπλήσι\expo{ον} ᾧ ἐσφράγιζ(εν) ὁ ἀνήρ {\policeFr ÷}}{« Une petite bague-cachet (d'une valeur de trois oboles) pareille à celle avec laquelle le mari scellait la porte du cellier. »}
\commScholie{Explique comment les femmes contrefaisaient aisément jadis l'empreinte du cachet marital.}
\entreeScholie{156}{427}{θριπήδεστ᾽ ἔχειν / σφραγίδια}{Marge gauche (extérieure, f. 167v, en face du v. 427)}{ξύλα ὑπὸ θριπῶν βεβρωμένα, οἷς ἐσφράγιζον. θρῖπες δὲ εἶδος σκωλήκων.}{ξύλα ὑπ(ὸ) θριπ\expo{ῶν} βεβρωμένα, οἷς ἐσφράγιζ\expo{ον} · θρῖπ\expo{ες} δ(ὲ) εἶδ\expo{ος} σκωλήκ\expo{ων} {\policeFr ÷}}{« Des morceaux de bois vermoulus (rongés par les vers à bois), dont les maris se servent désormais comme cachets impossibles à imiter. Les θρῖπες sont une espèce de vers qui rongent le bois. »}
\commScholie{Repris par la Souda θ 495, Photius θ 222 et Hésychius θ 754 : les sillons capricieux creusés par les vers à bois (θρῖπες) formaient un sceau naturel infalsifiable.}
\entreeScholie{157}{428}{ἐξαψαμένους}{Interlinéaire (supra lineam au-dessus d'ἐξαψαμένους, f. 167v, v. 428)}{ἐξάψαντας ἑαυτῶν.}{ἐξάψαντ\expo{ας} ἑαυτ\expo{ῶν} {\policeFr ÷}}{« Les ayant attachés (suspendus par un cordon) à leur propre personne. »}
\commScholie{Glose la valeur moyenne directe du participe ἐξαψαμένους.}
\entreeScholie{158}{429}{ὄλεθρόν τιν᾽ ἡμᾶς κυρκανᾶν ἁμωσγέπως,}{Interlinéaire (supra lineam au-dessus de κυκῶντας, f. 167v, v. 429)}{ταράττειν καὶ κινεῖν.}{ταράττ(ειν) \textbf{ς} κιν(εῖν) {\policeFr ÷}}{« Brasser et remuer (un breuvage mortel, un poison). »}
\commScholie{Glose lexicale de R\expo{2} expliquant l'infinitif expressif et familier κυρκανᾶν (« brasser, touiller, machiner » un trépas, v. 429) par les verbes usuels ταράττειν καὶ κινεῖν ; reprise par la Souda κ 2769 et Hésychius κ 4679.}
\entreeScholie{159}{441}{ὥστ᾽ ἂν εἰ λέγοι παρ᾽ αὐτῆς / Ξενοκλέης [Ξενοκλῆς R] ὁ Καρκίνου, δο- / κεῖν ἂν αὐτόν, ὡς ἐγᾦμαι,}{Marge gauche (extérieure, f. 167v, bloc de 5 lignes en face des vv. 440–441)}{διὰ τὸ εὐμήχανον ἐν τοῖς δράμασι λέγεσθαι τὸν Ξενοκλέα πρὸς τὰς τῆς γυναικὸς μηχανὰς παρελήφθη. χλευάζων δὲ λέγει οὐκ ἀπὸ σπουδῆς, ὡς ἐπὶ συμμέτρου. Ἄλλως. ὁ τραγῳδοποιός. ἤτοι δὲ διαβάλλει αὐτὸν ὡς ἀφυᾶ.}{διὰ τ(ὸ) εὐμήχαν(ον) ἐν τ\expo{οῖς} δράμασι λέγεσθ(αι) τ(ὸν) Ξενοκλέα πρ(ὸς) τ\expo{ὰς} τ\expo{ῆς} γυναικὸς μηχανὰς παρελήφθη · χλευάζων δ(ὲ) λέγει οὐκ ἀπ(ὸ) σπουδ\expo{ῆς}, ὡς ἐπὶ συμμέτρ\expo{ου} · ἄλλ\expo{ως} · ὁ τραγῳδοποιός · ἤτοι δ(ὲ) διαβάλλ(ει) αὐτ(ὸν) ὡς ἀφυᾶ {\policeFr ÷}}{« C'est parce que Xénoclès (fils de Carcinos) passait pour abuser des machines de théâtre (εὐμήχανος, cf. Paix v. 790 μηχανοδίφας) qu'il est mentionné ici en regard des machinations de l'oratrice. Mais le poète dit cela par dérision et non sérieusement. — Autre rédaction : Le poète tragique Xénoclès ; il le raille ici pour son manque total de talent (ὡς ἀφυᾶ). »}
\commScholie{Double scholie intéressante sur le tragique Xénoclès, fils de Carcinos (qui avait pourtant battu la tétralogie troyenne d'Euripide en 415 av. J.-C.).}
\entreeScholie{160}{442}{ἄντικρυς}{Interlinéaire (supra lineam au-dessus d'ἄντικρυς, f. 167v, v. 442)}{φανερῶς.}{φανερ\expo{ῶς} {\policeFr ÷}}{« Ouvertement, tout droit, franchement. »}
\commScholie{Glose adverbiale de R\expo{2} expliquant l'adverbe attique ἄντικρυς (« carrément, tout net, ouvertement ») par φανερῶς dans l'éloge choral de l'éloquence de Mika.}
\entreeScholie{161}{445}{παρῆλθον ῥημάτων,}{Interlinéaire (supra lineam, f. 167v, v. 445)}{ὑπὸ Εὐριπίδου.}{ὑπ(ὸ) Εὐριπίδ\expo{ου} {\policeFr ÷}}{« (Le mal que j'ai subi) du fait d'Euripide. »}
\commScholie{Précise l'agent du préjudice subi par la marchande de couronnes (Στεφανόπωλις dans l'édition de Bibi).}
\entreeScholie{162}{448}{στεφανηπλοκοῦσ᾽ ἔβοσκον ἐν ταῖς μυρρίναις.}{Interlinéaire (supra lineam, f. 167v, v. 448)}{διὰ μυρρινῶν στεφάνους ποιοῦσα.}{διὰ μυρριν\expo{ῶν} στεφάν\expo{ους} ποιοῦσα {\policeFr ÷}}{« En confectionnant des couronnes avec des branches de myrte (sur le marché aux myrtes, ἐν ταῖς μυρρίναις). »}
\commScholie{Explique la métonymie attique ἐν ταῖς μυρρίναις (le quartier du marché aux couronnes de myrte à l'Agora).}
\entreeScholie{163}{449}{τέως μὲν οὖν ἀλλ᾽ ἡμικάκως ἐβοσκόμην·}{Interlinéaire (supra lineam au-dessus de τέως, f. 167v, v. 449)}{ἀντὶ τοῦ πρότερον.}{ἀν\auDessus{τ}{ι} τ\expo{οῦ} πρότερ\expo{ον} {\policeFr ÷}}{« τέως : au sens de “jusqu'alors, auparavant”. »}
\commScholie{Cité aussi par le scholiaste de Platon (Hipparque 229d).}
\entreeScholie{164}{458}{ἀλλ᾽ εἰς ἀγορὰν ἄπειμι· δεῖ γὰρ ἀνδράσιν / πλέξαι στεφάνους συνθηματιαίους εἴκοσιν·}{Interlinéaire (supra lineam au-dessus de συνθηματιαίους, f. 167v, v. 458)}{οὓς ἡμεῖς ἐκδοτικοὺς λέγομεν.}{οὓς ἡμεῖς ἐκδοτικ\expo{οὺς} λέγομ(εν) {\policeFr ÷}}{« Ce que nous appelons aujourd'hui des couronnes “faites sur commande” (ἐκδοτικούς, réservées par contrat préalable). »}
\commScholie{Repris par la Souda σ 1456 et Photius σ 696 : glose précieuse donnant l'équivalent grec d'époque impériale (ἐκδοτικός) de l'attique classique συνθηματιαῖος.}
\entreeScholie{165}{459}{ἕτερον αὖ τι λῆμα τοῦτο}{Interlinéaire (supra lineam au-dessus de λῆμα, f. 167v, v. 459)}{φρόνημα.}{φρόνημα {\policeFr ÷}}{« La fierté d'âme, la résolution (φρόνημα). »}
\commScholie{Glose interlinéaire de R\expo{2} expliquant le substantif poétique λῆμα (« caractère résolu, audace, tempérament ») par φρόνημα au début de l'intermède choral.}
\clearpage
\renewcommand{\currHeadLE}{ΑΡΙΣΤΟΦΑΝΟΥΣ ΘΕΣΜΟΦΟΡΙΑΖΟΥΣΑΙ · CODEX RAVENNAS 429 (\R) · FOLIO 12/30}
\renewcommand{\currHeadRE}{f. 168r (Classense 170r) · vv. 466–508 · Collation diplomatique}
\renewcommand{\currHeadLO}{ΑΡΙΣΤΟΦΑΝΟΥΣ ΘΕΣΜΟΦΟΡΙΑΖΟΥΣΑΙ · CODEX RAVENNAS 429 (\R) · FOLIO 12/30}
\renewcommand{\currHeadRO}{f. 168r (Classense 170r) · vv. 466–508 · Collation diplomatique}
\noindent\begin{adjustbox}{max totalheight=30.4cm, max width=\textwidth, keepaspectratio}\begin{minipage}{\textwidth}
\bandeauRav{f. 168r}{Classense 170r}{466–508}{Folio 12/30 · Recto (page de droite) · 14 scholies \& gloses ({\policeGr Σ} 166–179)}
{\fontsize{8.8pt}{10.2pt}\selectfont
\ligneRav{466}{μ\auDessus{ν}{Η}}{τὸ μὲν ὦ γυναῖκες ὀξυθυμεῖσθαι σφόδρα}{}
\ligneRav{467}{}{Εὐριπίδῃ, τοιαῦτ᾽ \difference{ἀκουούσας} κακά,}{\textbf{467} ἀκουούσαις R : \lecon{ἀκουούσας} Zanetti}
\ligneRav{}{}{οὐ θαυμάσιόν ἐστ᾽, οὐδ᾽ ἐπιζεῖν τὴν χολήν·}{}
\ligneRav{469}{}{καὐτὴ γὰρ ἔγωγ᾽, οὕτως \difference{ὀναίμην} τῶν τέκνων,}{\textbf{469} ὀνοίμην R : \lecon{ὀναίμην} Suid. ο 330}
\ligneRav{470}{}{μισῶ τὸν ἄνδρ᾽ ἐκεῖνον, εἰ μὴ μαίνομαι·}{}
\ligneRav{471}{}{ὅμως δ᾽ ἐν \difference{ἀλλήλαισι} χρὴ δοῦναι λόγον·}{\textbf{471} ἀλλήλοισι R : \lecon{ἀλλήλαισι} Brunck}
\ligneRav{472}{}{αὐταὶ γάρ ἐσμεν, κοὐδεμί᾽ ἔκφορος λόγου·}{\textbf{472} \lecon{ἔκφορος} R : ἐκφορὰ Valckenaer ad Eur. \italique{Hipp.} 294}
\ligneRav{473}{}{τί ταῦτ᾽ ἔχουσαι \difference{᾽κεῖνον} αἰτιώμεθα}{\textbf{473} κεῖνον R, Suid τ 683 : \lecon{᾽κεῖνον} Meineke}
\ligneRav{474}{}{βαρέως τε φέρομεν, \difference{εἰ δύ᾽} ἡμῶν ἢ τρία}{\textbf{474} ἢ δύ᾽ R : \lecon{εἰ δύ᾽} Grynaeus (\italique{si} Divus)}
\ligneRav{475}{}{κακὰ ξυνειδὼς εἶπε δρώσας μυρία;}{}
\ligneRav{476}{}{ἐγὼ γὰρ αὐτὴ πρῶτον, ἵνα \difference{μἄλλην} λέγω,}{\textbf{476} μὴ ἄλλην R : \lecon{μἄλλην} Dindorf}
\ligneRav{477}{}{ξύνοιδ᾽ ἐμαυτῇ \difference{πολλά γ᾽· ἀλλ᾽} ἐκεῖνο δ᾽ οὖν}{\textbf{477} πολλ᾽ R : δεινὰ πολλ᾽ Bentley : \lecon{πολλά γ᾽· ἀλλ᾽} Bibi}
\ligneRav{478}{}{δεινότατον, ὅτε νύμφη μὲν \difference{ἦ} τρεῖς ἡμέρας,}{\textbf{478} ἦν R : \lecon{ἦ} Dindorf (1825)}
\ligneRav{479}{}{ὁ δ᾽ ἀνὴρ παρ᾽ ἐμοὶ \difference{καθηῦδεν}· ἦν δέ μοι φίλος,}{\textbf{479} καθεύδειν R : \lecon{καθηῦδεν} Scaliger}
\ligneRav{480}{}{ὅσπερ με διεκόρησεν \difference{οὖσαν ἑπτέτιν}·}{\textbf{480} διεκόρευσεν R : \lecon{-ρησεν} Pollux cf. \italique{Pax} 59 \espApp ἑπτ. οὖσ. R : \lecon{οὖσαν ἑπτ.} Pollux}
\ligneRav{481}{}{οὖτος πόθῳ \difference{μου ᾽κνυεν} ἐλθὼν τὴν θύραν·}{\textbf{481} μ᾽ οὔκνυεν R : \lecon{μου ᾽κνυεν} Suid. ε 520}
\ligneRav{482}{}{\difference{κᾆτ᾽} εὐθὺς ἔγνων· εἶτα καταβαίνω λάθρᾳ·}{\textbf{482} κᾶτ᾽ R : \lecon{κᾆτ᾽} Küster}
\ligneRav{}{}{ὁ δ᾽ ἀνὴρ ἐρωτᾷ `ποῖ σὺ καταβαίνεις;' `ὅποι;}{}
\ligneRav{}{}{στρόφος μ᾽ ἔχει τὴν γαστέρ᾽ ὦνερ κὠδύνη·}{}
\ligneRav{485}{}{ἐς τὸν κοπρῶν᾽ οὖν ἔρχομαι.' `βάδιζέ νυν.'}{}
\ligneRav{486}{}{κᾆθ᾽ ὁ μὲν ἔτριβεν κεδρίδας \difference{ἄννηθον} σφάκον·}{\textbf{486} ἄνηθον R, Suid. : \lecon{ἄννηθον} Dindorf (1830)}
\ligneRav{}{}{ἐγὼ δὲ καταχέασα τοῦ στροφέως ὕδωρ}{}
\ligneRav{488}{}{ἐξῆλθον ὡς τὸν μοιχόν· εἶτ᾽ \difference{ἠρειδόμην}}{\textbf{488} ἐρειδόμην R : \lecon{ἠρειδόμην} Küster}
\ligneRav{489}{}{παρὰ τὸν Ἀγυιᾶ κ\inRas{υν}{ύβ}δ᾽ ἐχομένη τῆς δάφνης·}{\textbf{489} κυνδ᾽ \Rac : \lecon{κύβδ᾽} \Rpc, \Mon}
\ligneRav{490}{}{ταῦτ᾽ οὐδεπώποτ᾽ εἶφ᾽, ὁρᾶτ᾽, Εὐριπίδης·}{}
\ligneRav{}{}{οὐδ᾽ ὡς ὑπὸ τῶν δούλων τε κὠρεωκόμων}{}
\ligneRav{}{}{σποδούμεθ᾽, ἢν μὴ 'χωμεν ἕτερον, οὐ λέγει·}{}
\ligneRav{493}{}{οὐδ᾽ ὡς ὅταν μάλισθ᾽ ὑπό του \difference{ληκώμεθα} (\inRas{ληκ}{κιν}ώμεθα)}{\textbf{493} \lecon{ληκώμεθα} \Rac (?), Suid. σ 1126 : \inRas{ληκ}{κιν}ώμεθα \Rpc}
\ligneRav{494}{}{τὴν νύχθ᾽, ἕωθεν \difference{σκόροδα διαμασώμεθα},}{\textbf{494} σκορόδια R : \lecon{σκόροδα δια} Bentley \espApp post νύχθ' add. ὅλην Suid. (reminiscentia Pherecratis ? cf. Photius λ 254)}
\ligneRav{495}{}{ἵν᾽ \difference{ὀσφρόμενος} \difference{ἁνὴρ} \difference{ἀπὸ} τείχους εἰσιὼν}{\textbf{495} ἀπὸ τοῦ R : \lecon{ἀπὸ} Biset \espApp ὀσφραινόμενος R : \lecon{ὀσφρόμενος} Bentley, Küster \espApp ἀνὴρ R : \lecon{ἁνὴρ} Dawes}
\ligneRav{}{}{μηδὲν κακὸν δρᾶν ὑποτοπῆται. ταῦθ᾽, ὁρᾷς,}{}
\ligneRav{}{}{οὐπώποτ᾽ εἶπεν. εἰ δὲ Φαίδραν λοιδορεῖ,}{}
\ligneRav{}{}{ἡμῖν τί τοῦτ᾽ ἔστ᾽; οὐδ᾽ ἐκεῖν᾽ εἴρηκέ πω,}{}
\ligneRav{499}{}{ὡς ἡ γυνὴ δεικνῦσα τἀνδρὶ τοὔγκυκλον}{}
\ligneRav{(500)}{}{\textlbrackdbl\sout{ὑπ᾽ αὐγὰς οἷόν ἐστιν ἐγκεκαλυμμένον}\textrbrackdbl}{\textbf{v. 500} ὑπ᾽ αὐγὰς οἷόν ἐστιν ἐγκεκαλυμμένον R (ὑπαυγάσ᾽ κτλ. Hermann) : \lecon{del.} van Herwerden, Bibi (ex gl. ᾦόν)}
\ligneRav{501}{}{τὸν \difference{μοιχὸν} ἐξέπεμψεν, οὐκ εἴρηκέ πω·}{\textbf{501} μυχὸν R : \lecon{μοιχὸν} Grynaeus}
\ligneRav{502}{}{ἑτέραν δ᾽ ἐγᾦδ᾽ \difference{ἣ 'φασκεν} ὠδίνειν γυνὴ}{\textbf{502} ήφασκεν R : \lecon{ἣ 'φασκεν} Zanetti}
\ligneRav{}{}{δέχ᾽ ἡμέρας, ἔως ἐπρίατο παιδίον·}{}
\ligneRav{504}{}{ὁ δ᾽ ἀνὴρ περιήρχετ᾽ \difference{ὠκυτόκι᾽} ὠνούμενος·}{\textbf{504} \lecon{περιήρχετ᾽} R : περιέτρεχ᾽ Bergk \espApp ὠκυτόκειa R : \lecon{ὠκυτόκι᾽} Pollux II.7}
\ligneRav{505}{}{τὸ δ᾽ εἰσέφερε γραῦς ἐν χύτρᾳ τὸ παιδίον,}{\textbf{505} \lecon{τὸ δ᾽} R : τότ᾽ Biset : τῇδ᾽ Daubuz}
\ligneRav{506}{}{ἵνα μὴ \expo{βοῴη} κηρίῳ βεβυσμένον·}{\textbf{506} \lecon{βοῴη} suppl. \Rpc}
\ligneRav{}{}{εἶθ᾽ ὡς ἔνευσεν ἡ φέρουσ᾽, εὐθὺς βοᾷ,}{}
\ligneRav{}{}{`ἄπελθ᾽ ἄπελθ᾽, ἤδη γὰρ ὦνέρ μοι δοκῶ}{}
}
\vspace{3pt}{\color{ruleGold}\hrule height 0.6pt}\vspace{3pt}
{\fontsize{8.0pt}{9.4pt}\selectfont
\noindent{\bfseries\scshape\color{rubricRed}Loci paralleli et testimonia (f. 168r, vv. 466–508) :} \textbf{466} (ὀξυθυμεῖσθαι) Photius ο 380 ὀξυθυμεῖσθαι \espTest \textbf{468} (θαυμάσιον = θαυμαστόν) Moeris θ 16 {\policeFr ≈} Hesychius θ 146 θαυμάσιος \espTest \textbf{469} Suidas ο 330 (sed ἐγώ pro ἔγωγ᾽) ; cf. (ὀναίμην) Hesychius ο 852 ὀναίμην ; \Syn ο 162 ὀναίμην = Photius ο 332 \espTest \textbf{470} Suidas μ 328 (τἄνδρ᾽ pro τὸν ἄνδρ᾽) \espTest \textbf{473-4} (τί ταῦτ᾽ -- φέρομεν) Suidas τ 683 \espTest \textbf{480} Pollux III 42 (οὗτος pro ὅσπερ et διεκόρησεν cod. A pro διεκόρευσεν) \espTest (διεκόρησεν) Hesychius δ 1558 διεκόρησε \espTest \textbf{480} Pollux III 42 (οὗτος pro ὅσπερ et διεκόρησεν pro διεκόρευσεν) \espTest \textbf{481} Suidas ε 520 ; Zonaras 669 ἔκνυεν; \Eustathe 1746.6 (μοὔκνυεν ἕλκων pro μου ᾽κνυεν ἐλθὼν) \espTest \textbf{484} \Eustathe 1746.45 \espTest (στρόφος) \Eustathe 503.20 \espTest \textbf{486} Suidas σ 1712 \espTest (κεδρίδας) Hesychius κ 1980 κεδρίδες ; Photius κ 1980 κ 517 \espTest \textbf{489} (Ἀγυιᾶ) Stephanus Byzantinus 22.15 ; cf. Moeris α 13 \espTest \textbf{493-4} Suidas λ 434 Ληκώμεθα: κινώμεθα. Ἀριστοφάνης: μάλισθ' ὅταν ὑπό του ληκώμεθα τὴν νύχθ' ὅλην. Post νύχθ' add ὅλην Suid. (reminiscentia Pherecratis ? cf. Photius λ 254) \espTest \textbf{504} (ὠκυτόκι᾽) Pollux II 7 ; \Phryn 129.3 \espTest \textbf{506} (κηρίῳ βεβυσμένον) Hesychius κ 2549 κηρίῳ βύσασα ; Photius κ 675 κηρίῳ βῦσαι\par\vspace{2.5pt}
{\color{ruleGold}\hrule height 0.35pt}\vspace{2.5pt}
\noindent{\bfseries\scshape\color{rubricRed}Suite (en commentaire) de l'apparat critique et paléographique :} \textbf{v. 477 :} Le vers de R (οὐδ' εἴρηκεν οὐδὲν πρὸς ἃ δρῶμεν· πολλ') est amputé d'une syllabe finale ; \textbf{\textbf{Bibi} restitue πολλά γ'· ἀλλ'} (chute de -ά γ'· ἀλλ' par saut du même au même entre les deux ΛΛ). --- \textbf{v. 489 :} \Rac{} avait écrit κυνδ' (confusion minuscule β/ν), corrigé en \lecon{κύβδ'} par \Rpc{} (suivi par \Mon{}). --- \textbf{v. 493 :} Comme le précise Louis Hauseux (\textit{vv 1-641.tex}, l. 3715), le copiste avait d'abord écrit la forme crue \lecon{ληκώμεθα} (attestée par la Souda σ 1126 et adoptée par L. Hauseux), puis a \textbf{gratté les trois premières lettres ληκ-} pour les remplacer par κιν- (\lecon{\inRas{ληκ}{κιν}ώμεθα} \Rpc{}) sous l'influence de la glose κινώμεθα. --- \textbf{v. 500 (genèse paléographique d'un vers interpolé) :} Le v. 500 (ὑπ' αὐγὰς οἷόν ἐστιν ἐγκεκαλυμμένον), athétisé par van Herwerden et L. Hauseux (\textit{vv 1-641.tex}, ll. 3781–3806), est né d'une glose ou citation marginale [τὸν μοιχόν·] ὑπαυγάσας. ᾦόν ἐστιν ἐγκεκαλυμμένον, dans laquelle ΩΙΟΝ (avec iota muet adscrit = ᾦόν, « un œuf enveloppé ») a été lu par iotacisme et bêtacisme ω/ο \textbf{ΟΙΟΝ (οἷόν)}, entraînant la réécriture métrique ὑπ' αὐγὰς οἷόν ἐστιν.... --- \textbf{v. 506 :} L'optatif \lecon{βοῴη}, omis par \Rac{}, est suppléé par \Rpc{}. --- \textbf{v. 477 :} πολλ' R (δεινὰ πολλ' Bentley) : \lecon{πολλά γ'· ἀλλ'} \lecon{Bibi}. --- \textbf{v. 481 :} μ' οὔκνυεν R (médivision et diphtongaison) : \lecon{μου 'κνυεν} Souda ε 520 (cf. Hésychius κ 3133). --- \textbf{v. 501 :} μυχὸν R (iotacisme) : \lecon{μοιχὸν} Grynaeus.\par}
\end{minipage}\end{adjustbox}
\newpage
\renewcommand{\currHeadLE}{SCHOLIES MARGINALES ({\policeGr Σ}\expo{R}) \& GLOSES INTERLINÉAIRES · FOLIO 12/30}
\renewcommand{\currHeadRE}{f. 168r (Classense 170r) · vv. 466–508 · {\policeGr Σ} 166 à {\policeGr Σ} 179}
\renewcommand{\currHeadLO}{SCHOLIES MARGINALES ({\policeGr Σ}\expo{R}) \& GLOSES INTERLINÉAIRES · FOLIO 12/30}
\renewcommand{\currHeadRO}{f. 168r (Classense 170r) · vv. 466–508 · {\policeGr Σ} 166 à {\policeGr Σ} 179}
\bandeauScholies{Folio 12/30 — f. 168r}{Classense 170r}{vv. 466–508}{Recto (page de droite) · 14 entrées ({\policeGr Σ} 166 à {\policeGr Σ} 179)}
{\fontsize{8.8pt}{10.5pt}\selectfont\noindent\textbf{\color{rubricRed}Disposition topographique des marges sur le f. 168r :} Marge extérieure à droite ; marge intérieure à gauche. — \textit{Fac-similé de référence :} \texttt{folios\_individuels\_ff162v-177r/12\_f168r\_classense170r\_vv0466-0508.jpg}.\par\vspace{3pt}}
\entreeScholie{166}{466}{τὸ μὲν, ὦ γυναῖκες, ὀξυθυμεῖσθαι σφόδρα / Εὐριπίδῃ, τοιαῦτ᾽ ἀκουούσας κακά,}{Marge supérieure et marge droite (extérieure, f. 168r, en tête du discours de Mnésiloque)}{ἐπίτηδες ὁ κηδεστὴς ὑπὲρ τὰς ἄλλας γυναῖκας γυναικίζεται, ἵνα διὰ τούτου λάθῃ.}{ἐπίτηδ\expo{ες} ὁ κηδεστ\expo{ὴς} ὑπ(ὲρ) τ\expo{ὰς} ἄλλ\expo{ας} γυναῖκ\expo{ας} γυναικίζετ(αι), ἵνα διὰ τούτ\expo{ου} λάθ(ῃ) {\policeFr ÷}}{« C'est à dessein que le parent par alliance (Mnésiloque) affecte de parler en femme encore plus que les autres femmes, afin de passer ainsi inaperçu (avant de prendre la défense d'Euripide). »}
\commScholie{Cité par Louis Hauseux (vv 1-641.tex, l. 3611) : note l'habileté rhétorique de l'insinuatio de Mnésiloque aux vv. 466-472.}
\entreeScholie{167}{472}{αὐταὶ γάρ ἐσμεν, κοὐδεμί᾽ ἔκφορος λόγου·}{Marge droite (extérieure) et interlinéaire (f. 168r, v. 472)}{ἀντὶ τοῦ μόναι. οὐδεμία ἐκφέρουσα τοὺς λόγους πρὸς τοὺς ἄνδρας.}{ἀν\auDessus{τ}{ι} τ\expo{οῦ} μόναι {\policeFr ÷} // οὐδεμία ἐκφέρουσα τ\expo{οὺς} λόγ\expo{ους} πρ(ὸς) τ\expo{οὺς} ἄνδρ\expo{ας} {\policeFr ÷}}{« αὐταί : au sens de “entre nous toutes seules” (μόναι, cf. Acharniens v. 504 αὐτοὶ γάρ ἐσμεν). — κοὐδεμί' ἐκφορὰ λόγου : et il n'y a aucun risque que nos paroles soient divulguées au dehors auprès des maris. »}
\commScholie{Repris par la Souda α 4470 et ε 720.}
\entreeScholie{168}{473}{τί ταῦτ᾽ ἔχουσαι ᾽κεῖνον αἰτιώμεθα}{Interlinéaire (supra lineam au-dessus de τί, f. 168r, v. 473)}{διὰ τί.}{διὰ τί {\policeFr ÷}}{« Pourquoi donc ? »}
\commScholie{Glose syntaxique de R\expo{2} précisant la valeur causale de l'interrogatif neutre adverbial τί (« pourquoi donc », διὰ τί) en tête de la phrase τί ταῦτ᾽ ἔχουσαι ᾽κεῖνον αἰτιώμεθα.}
\entreeScholie{169}{481}{οὗτος πόθῳ μου ᾽κνυεν ἐλθὼν τὴν θύραν·}{Marge droite (extérieure, f. 168r, en face du v. 481)}{ἔξυεν ἠρέμα, παρὰ τὸ κνᾶν.}{ἔξυεν ἠρέμα, παρὰ τ(ὸ) κνᾶν {\policeFr ÷}}{« Il grattait doucement la porte du doigt (signal discret de l'amant nocturne), dérivé de κνᾶν (“gratter”). »}
\commScholie{Repris par la Souda κ 1891 : distingue le discret κνύειν (gratter du bout de l'ongle pour avertir la femme adultère) du bruyant κόπτειν (heurter à la porte).}
\entreeScholie{170}{485}{ὁ δ᾽ ἀνὴρ ἐρωτᾷ «ποῖ σὺ καταβαίνεις;» — «ὅποι; / στρόφος μ᾽ ἔχει τὴν γαστέρ᾽ ὦνερ κὠδύνη· / ἐς τὸν κοπρῶν᾽ οὖν ἔρχομαι.»}{Interlinéaire (supra lineam, f. 168r, v. 485)}{λείπει τὸ ἔφη.}{λείπ(ει) τ(ὸ) ἔφ(ην) {\policeFr ÷}}{« Il manque (par ellipse) le verbe ἔφην (“et je dis à mon mari : Je me meurs de coliques !”). »}
\commScholie{Noté avec justesse : dans la vivacité du récit parlé, le verbe déclaratif est sous-entendu.}
\entreeScholie{171}{486}{κᾆθ᾽ ὁ μὲν ἔτριβεν κεδρίδας ἄννηθον σφάκον·}{Marge droite (extérieure, f. 168r, en face du v. 486)}{ἐλελίσφακον. ὡς ἐπιτήδεια παῦσαι στρόφον.}{ἐλελίσφακ\expo{ον} · ὡς ἐπιτήδεια παῦσ(αι) στρόφ\expo{ον} {\policeFr ÷}}{« σφάκον : la sauge officinale (ἐλελίσφακος, Salvia triloba) ; plantes propres à calmer les tranchées intestinales (στρόφος). »}
\commScholie{Cité par Louis Hauseux (vv 1-641.tex, l. 3667) avec le témoignage de Pline l'Ancien et Dioscoride sur les propriétés antispasmodiques de la sauge (ἐλελίσφακος) et des baies de genévrier (κεδρίς).}
\entreeScholie{172}{487}{ἐγὼ δὲ καταχέασα τοῦ στροφέως ὕδωρ}{Marge droite (extérieure, f. 168r, en face du v. 487)}{ἵνα ὀλισθηρὸς γενόμενος μὴ ψοφῇ.}{ἵνα ὀλισθηρ\expo{ὸς} γενόμεν\expo{ος} μὴ ψοφῇ {\policeFr ÷}}{« (“Ayant versé de l'eau sur le gond de la porte”) : afin que, rendu glissant (lubrifié), il ne grince pas en tournant. »}
\commScholie{Cité par Louis Hauseux (vv 1-641.tex, l. 3677) : le pivot de bois de la porte tournant dans sa crapaudine grinçait à sec ; les femmes adultères ou les esclaves l'arrosaient d'eau ou de vin pour sortir la nuit sans bruit (cf. Plaute, Curculio vv. 158-161).}
\entreeScholie{173}{489}{κύβδ᾽ ἐχομένη τῆς δάφνης / παρὰ τὸν Ἀγυιᾶ}{Marge droite (extérieure, f. 168r, en face du v. 489)}{ἀγυιεὺς οὕτω καλούμενος Ἀπόλλων τετράγωνος.}{ἀγυιεὺς οὕτω καλούμεν\expo{ος} Ἀπόλλ(ων) τετράγων\expo{ος} {\policeFr ÷}}{« On appelle Aguieus (Ἀγυιεύς, “Protecteur des rues”) le pilier conique ou quadrangulaire d'Apollon dressé devant la porte d'entrée des maisons athéniennes. »}
\commScholie{Repris par la Souda α 383 : détail pittoresque de la topographie domestique athénienne (le pilier d'Apollon Aguieus et le laurier planté devant la porte de rue).}
\entreeScholie{174}{491}{ὑπὸ τῶν δούλων τε κὠρεωκόμων}{Marge droite (extérieure, f. 168r, en face du v. 491)}{τῶν νῦν καλουμένων σταβλίτων ἢ βουρδωναρίων, τῶν τοὺς οὐρῆας κομούντων.}{τ\expo{ῶν} νῦν καλουμέν\expo{ων} σταβλίτ\expo{ων} ἢ βουρδωναρί\expo{ων}, τ\expo{ῶν} τ\expo{οὺς} οὐρῆ\expo{ας} κομούντ\expo{ων} {\policeFr ÷}}{« ὀρεωκόμων (“muletiers”) : ceux qu'on appelle aujourd'hui σταβλῖται (palefreniers d'écurie, lat. stabularii) ou βουρδωνάριοι (conducteurs de bardots/mulets, lat. burdonarii), c'est-à-dire ceux qui soignent les mulets (οὐρῆας). »}
\commScholie{Scholie célèbre pour l'histoire de la langue grecque tardive : la glose initiale ancienne (τῶν τοὺς οὐρῆας κομούντων) a été enrichie à l'époque protobyzantine (Ve-VIe s.) des deux emprunts latins σταβλίτης (stabularius) et βουρδωνάριος (burdonarius).}
\entreeScholie{175}{495}{ἵν᾽ ὀσφρόμενος ἁνὴρ ἀπὸ τείχους εἰσιὼν}{Marge droite (extérieure, f. 168r, en face du v. 495)}{διὰ γὰρ τὸν πρὸς Πελοποννησίους πόλεμον ἐφύλαττον Ἀθηναῖοι τὰ τείχη.}{διὰ γ(ὰρ) τ(ὸν) πρ(ὸς) Πελοποννησί\expo{ους} πόλεμ\expo{ον} ἐφύλαττ\expo{ον} Ἀθηναῖ(οι) τὰ τείχ(η) {\policeFr ÷}}{« Parce qu'en raison de la guerre du Péloponnèse (et de l'occupation permanente de Décélie par Agis depuis 413 av. J.-C.), les citoyens athéniens montaient la garde toute la nuit sur les remparts (cf. Thucydide VII, 28, 2). »}
\commScholie{Exactitude historique remarquable : depuis la fortification de Décélie en 413, les Athéniens couchaient en armes sur les remparts, ce qui favorisait les escapades nocturnes évoquées par Mnésiloque.}
\entreeScholie{176}{499}{ὡς ἡ γυνὴ δεικνῦσα τἀνδρὶ τοὔγκυκλον}{Interlinéaire (supra lineam au-dessus d'ἔγκυκλον, f. 168r, v. 499)}{στολῆς εἶδος.}{στολ\expo{ῆς} εἶδ\expo{ος} {\policeFr ÷}}{« Une espèce de manteau féminin. »}
\commScholie{Glose vestimentaire de R\expo{2} rappelant que l'ἔγκυκλον (τοὔγκυκλον, déjà glosé en Σ 261) désigne un ample manteau circulaire de femme, dont l'épouse infidèle se sert ici comme d'un écran pour faire évader son amant.}
\entreeScholie{177}{500}{ὑπ᾽ αὐγὰς οἷόν ἐστιν ἐγκεκαλυμμένον}{Marge droite (extérieure, f. 168r, en face du v. 500)}{ὅτι καὶ συστολαῖς βιαίοις χρῶνται καὶ ἐκτάσεσι. καὶ ἄνω [176] «ἀλλ’ ὧνπερ οὕνεκ’ ἦλθον ἀλλ’ εἰπεῖν βούλομαι.»}{ὅτι \textbf{ς} συστολ\expo{αῖς} βιαί\expo{οις} χρῶντ(αι) \textbf{ς} ἐκτάσεσ(ι) · \textbf{ς} ἄνω «ἀλλ' ὧνπ(ερ) οὕνεκ' ἦλθ\expo{ον} ἀλλ' εἰπ(εῖν) βούλομ(αι)» {\policeFr ÷}}{« Notez que les poètes comiques usent parfois d'abrègements métriques violents (συστολαῖς βιαίοις, ici sur la finale de ἐγκεκαλυμμένος/ον) et d'allongements, comme plus haut (v. 176) : “Mais ce pourquoi je suis venu, laisse-moi le dire.” »}
\commScholie{Scholie métrique capitale longuement analysée par Louis Hauseux (Bibi, v. 500, vv 1-641.tex, ll. 3748-3780) : le métricien antique constatait que le v. 500 pèche contre le mètre iambique ; Willems et Bibi montrent que tout le v. 500 (ὑπ' αὐγὰς οἷόν ἐστιν [= ῷόν ἐστιν] ἐγκεκαλυμμένον) est en réalité une glose marginale expliquant comment la femme dissimulait son amant sous son manteau déployé, glose qui s'est glissée dans la colonne avant l'époque d'Héliodore !}
\entreeScholie{178}{505}{ἑτέραν δ᾽ ἐγᾦδ᾽ ἣ ᾽φασκεν ὠδίνειν γυνὴ / δέχ᾽ ἡμέρας, ἕως ἐπρίατο παιδίον· / τὸ δ᾽ εἰσέφερε γραῦς ἐν χύτρᾳ τὸ παιδίον,}{Marge droite (extérieure, bas du f. 168r, en face du v. 505)}{ὅτι ἐν χύτρᾳ τὰ παιδία ἐξετίθεσαν.}{ὅτι ἐν χύτρᾳ τὰ παιδία ἐξετίθεσ(αν) {\policeFr ÷}}{« Notez que l'on exposait les nouveau-nés dans une marmite de terre cuite (ἐν χύτρᾳ, d'où le verbe ἐγχυτρίζω). »}
\commScholie{Repris par la Souda χ 605 : coutume athénienne bien attestée (cf. Grenouilles v. 1190 οἰδῶν τὼ πόδε ... ἐν ὀστράκῳ).}
\entreeScholie{179}{506}{τὸ δ᾽ εἰσέφερε γραῦς ἐν χύτρᾳ, τὸ παιδίον, / ἵνα μὴ βοῴη, κηρίω βεβυσμένον·}{Marge inférieure du f. 168r (en face des vv. 506–507)}{τὸ στόμα πεπληρωμένον. καὶ ἄλλως οὐ γάλα πρότερον τοῖς βρέφεσιν ἐδίδοσαν, ἀλλὰ μέλι ἀπολείχειν. Μένανδρος δὲ οὐκ ὀρθῶς ποιεῖ τὰ ἀρτίτοκα γάλακτος δεόμενα.}{τ(ὸ) στόμα πεπληρωμέν\expo{ον} · \textbf{ς} ἄλλ\expo{ως} οὐ γάλα πρότερ\expo{ον} τ\expo{οῖς} βρέφεσ(ιν) ἐδίδοσ(αν), ἀλλὰ μέλι ἀπολείχ(ειν) · Μένανδρ\expo{ος} δ(ὲ) οὐκ ὀρθ\expo{ῶς} ποιεῖ τὰ ἀρτίτοκα γάλακτ\expo{ος} δεόμενα {\policeFr ÷}}{« κηρίῳ βεβυσμένον : ayant la bouche bouchée (remplie) d'un rayon de miel. Et par ailleurs, on ne donnait pas d'abord du lait aux nouveau-nés, mais du miel à lécher ; aussi Ménandre (fr. 929 Koerte) commet-il un anachronisme médical lorsqu'il représente des nouveau-nés réclamant tout de suite du lait. »}
\commScholie{Scholie alexandrine savante (reprise par la Souda κ 1545 et discutée à propos des soins néonataux chez Soranos d'Éphèse) critiquant une inadvertance de Ménandre. Sur R, le verbe βοῴη omis par R\expo{1} a été rétabli par R\expo{2}.}
\clearpage
\renewcommand{\currHeadLE}{ΑΡΙΣΤΟΦΑΝΟΥΣ ΘΕΣΜΟΦΟΡΙΑΖΟΥΣΑΙ · CODEX RAVENNAS 429 (\R) · FOLIO 13/30}
\renewcommand{\currHeadRE}{f. 168v (Classense 170v) · vv. 509–550 · Collation diplomatique}
\renewcommand{\currHeadLO}{ΑΡΙΣΤΟΦΑΝΟΥΣ ΘΕΣΜΟΦΟΡΙΑΖΟΥΣΑΙ · CODEX RAVENNAS 429 (\R) · FOLIO 13/30}
\renewcommand{\currHeadRO}{f. 168v (Classense 170v) · vv. 509–550 · Collation diplomatique}
\noindent\begin{adjustbox}{max totalheight=30.4cm, max width=\textwidth, keepaspectratio}\begin{minipage}{\textwidth}
\bandeauRav{f. 168v}{Classense 170v}{509–550}{Folio 13/30 · Verso (page de gauche) · 9 scholies \& gloses ({\policeGr Σ} 180–188)}
{\fontsize{8.8pt}{10.2pt}\selectfont
\ligneRav{509}{μ\auDessus{ν}{Η}}{τέξειν.' τὸ γὰρ ἦτρον τῆς χύτρας ἐλάκτισεν·}{}
\ligneRav{510}{}{χὠ μὲν γεγηθὼς \difference{ἀπέτρεχ᾽}, ἡ δ᾽ ἐξέσπασεν}{\textbf{510} ἔτρεχεν R : \lecon{ἀπέτρεχ᾽} van Herwerden}
\ligneRav{511}{}{ἐκ τοῦ στόματος τοῦ παιδίου, τὸ δ᾽ ἀνέκραγεν·}{\textbf{511} \lecon{τοῦ παιδίου} R : τὸ κηρίον Hirschig}
\ligneRav{512}{}{εἶθ᾽ ἡ μιαρὰ γραῦς, \difference{ἣ ᾽φερεν} (\auDessus{εἰ}{η}\inRas{σ}{}έφερε) τὸ παιδίον,}{\textbf{512} \auDessus{εἰ}{η}\inRas{σ}{}έφερε \Rpc : \lecon{ἥ ᾽φερεν} Bentley}
\ligneRav{}{}{θεῖ μειδιῶσα πρὸς τὸν ἄνδρα καὶ λέγει,}{}
\ligneRav{514}{}{λέων λέων σοι γέγονεν, \difference{αὐτέκμαγμα} (αὖτ᾽ ἔ\inRas{γ}{κ}μαγμα) σόν,}{\textbf{514} αὖτ᾽ ἔκμαγμα (ex έγμαγμα) \Rpc, Pollux IX.131 : ἔκμαγμα Suid. λ 271, π 2114 : \lecon{αὐτέκμαγμα} Reiske}
\ligneRav{515}{}{τά τ᾽ ἄλλ᾽ ἁπαξάπαντα καὶ τὸ πόσθιον}{}
\ligneRav{516}{}{τῷ σῷ προσόμοιον, στρεβλὸν ὥσπερ κύτταρον.'}{\textbf{516} \lecon{κύτταρον} R, Suid. κ 2787 : κύτταρος Fritzsche}
\ligneRav{}{}{ταῦτ᾽ οὐ ποιοῦμεν τὰ κακά; νὴ τὴν Ἄρτεμιν}{}
\ligneRav{}{}{ἡμεῖς γε. κᾆτ᾽ Εὐριπίδῃ θυμούμεθα,}{}
\ligneRav{}{}{οὐδὲν παθοῦσαι μεῖζον ἢ δεδράκαμεν;}{}
\ligneRav{520}{—}{\hspace{1.1em}τουτὶ μέντοι θαυμαστόν,}{}
\ligneRav{521}{}{\hspace{1.1em}ὁπόθεν \difference{ηὑρέθη} τὸ χρῆμα,}{\textbf{521} εὑρέθη R : \lecon{ηὑρέθη} Meineke}
\ligneRav{}{}{\hspace{1.1em}χἤτις ἐξέθρεψε χώρα}{}
\ligneRav{}{}{\hspace{1.1em}τήνδε τὴν θρασεῖαν οὕτω·}{}
\ligneRav{}{}{\hspace{1.1em}τάδε γὰρ εἰπεῖν τὴν πανοῦργον}{}
\ligneRav{525}{}{\hspace{1.1em}κατὰ τὸ φανερὸν ὧδ᾽ ἀναιδῶς}{}
\ligneRav{526}{}{\hspace{1.1em}οὐκ ἂν ᾠόμην ἐν \difference{ἡμῖν}}{\textbf{526} ὑμῖν R : \lecon{ἡμῖν} Bonini}
\ligneRav{}{}{\hspace{1.1em}οὐδὲ τολμῆσαί ποτ᾽ ἄν·}{}
\ligneRav{527}{}{\hspace{1.1em}\difference{ἀλλὰ πᾶν} γένοιτ᾽ ἂν ἤδη·}{\textbf{527} ἀλλ᾽ ἅπαν R : \lecon{ἀλλὰ πᾶν} Cobet}
\ligneRav{}{}{\hspace{1.1em}τὴν παροιμίαν δ᾽ ἐπαινῶ}{}
\ligneRav{}{}{\hspace{1.1em}τὴν παλαιάν· ὑπὸ λίθῳ γὰρ}{}
\ligneRav{530}{}{\hspace{1.1em}παντί που χρὴ}{}
\ligneRav{}{}{\hspace{1.1em}μὴ δάκῃ ῥήτωρ ἀθρεῖν·}{}
\ligneRav{531}{χορ/}{ἀλλ᾽ οὐ γάρ ἐστι τῶν ἀναισχύτων φύσει γυναικῶν}{\textbf{531} (Κορ.) add. von Velsen}
\ligneRav{532}{}{οὐδὲν κάκιον εἰς ἅπαντα πλὴν ἄρ᾽ εἰ γυναῖκες·}{\textbf{532} \lecon{ἄρ᾽ εἰ} R : ἄλλαι Suid. α 1349}
\ligneRav{533}{γυνή τις}{οὔτοι μὰ τὴν \difference{Ἄγλαυρον} ὦ γυναῖκες εὖ φρονεῖτε,}{\textbf{533} (γυνὴ τις) \Rpc : \lecon{(Μίκα)} Rogers \espApp Ἄγραυλον R : \lecon{Ἄγλαυρον} Daubuz}
\ligneRav{}{}{ἀλλ᾽ ἢ πεφάρμαχθ᾽ ἢ κακόν τι μέγα πεπόνθατ᾽ ἄλλο,}{}
\ligneRav{535}{}{ταύτην ἐῶσαι τὴν φθόρον τοιαῦτα περιυβρίζειν}{}
\ligneRav{536}{}{ἡμᾶς ἁπάσας. εἰ μὲν οὖν τις ἔστιν· εἰ δὲ μή, ἡμεῖς}{\textbf{536} \lecon{μή ἡμεῖς} R : μἠμεῖς Bentley}
\ligneRav{537}{}{\difference{αὐταί} γε καὶ τὰ δουλάρια τέφραν ποθὲν λαβοῦσαι}{\textbf{537} αὐτοί R : \lecon{αὐταί} Bonini \espApp \lecon{γε} R : τε Σ T ad Hom.\italique{Il}.I.135}
\ligneRav{}{}{ταύτης ἀποψιλώσομεν τὸν χοῖρον, ἵνα διδαχθῇ}{}
\ligneRav{}{}{γυνὴ γυναῖκας οὖσα μὴ κακῶς λέγειν τὸ λοιπόν·}{}
\ligneRav{540}{μ\auDessus{ν}{Η}}{μὴ δῆτα τόν γε χοῖρον ὦ γυναῖκες. εἰ γὰρ οὔσης}{}
\ligneRav{}{}{παρρησίας κἀξὸν λέγειν ὅσαι πάρεσμεν ἀσταί,}{}
\ligneRav{542}{}{εἶτ᾽ εἶπον \difference{ἅ ᾽γίγνωσκον} ὑπὲρ Εὐριπίδου δίκαια,}{\textbf{542} ἅ γιγνωσκον R : \lecon{ἅ ᾽γίγνωσκον} Bentley}
\ligneRav{}{}{διὰ τοῦτο τιλλομένην με δεῖ δοῦναι δίκην ὑφ᾽ ὑμῶν;}{}
\ligneRav{}{\auDessus{γυ}{ν}}{οὐ γάρ σε δεῖ δοῦναι δίκην; ἥτις μόνη τέτληκας}{}
\ligneRav{545}{}{ὑπὲρ ἀνδρὸς ἀντειπεῖν, \difference{ὃς} ἡμᾶς πολλὰ κακὰ δέδρακεν}{\textbf{545} ὡς R : \lecon{ὃς} Grynaeus}
\ligneRav{546}{}{\difference{ἐπίτηδες} εὑρίσκων λόγους, ὅπου γυνὴ πονηρὰ}{\textbf{546} ἐξεπίτηδες R : \lecon{ἐπίτηδες} Scaliger}
\ligneRav{}{}{ἐγένετο, Μελανίππας ποιῶν Φαίδρας τε· Πηνελόπην δὲ}{}
\ligneRav{548}{}{οὐπώποτ᾽ \difference{ἐποίησ᾽}, ὅτι γυνὴ σώφρων ἔδοξεν εἶναι·}{\textbf{548} ἐποίησεν R : \lecon{ἐποίησ᾽} Scaliger}
\ligneRav{549}{μ\auDessus{ν}{Η}}{ἐγὼ γὰρ οἶδα ταἴτιον· μίαν γὰρ οὐκ \inRas{ἂρ}{ἂν} εἴποις}{\textbf{549} \lecon{ἂν} ex ἂρ corr. \Rpc}
\ligneRav{550}{}{τῶν νῦν γυναικῶν Πηνελόπην, Φαίδρας \difference{δ᾽ ἁπαξαπάσας}·}{\textbf{550} ἁπαξαπάσας R, M\expo{pc} : \lecon{δ᾽ ἁπαξαπάσας} M\expo{ac}}
}
\vspace{3pt}{\color{ruleGold}\hrule height 0.6pt}\vspace{3pt}
{\fontsize{8.0pt}{9.4pt}\selectfont
\noindent{\bfseries\scshape\color{rubricRed}Loci paralleli et testimonia (f. 168v, vv. 509–550) :} \textbf{509} (τὸ γὰρ ἦτρον -- ἐλάκτισεν) Suidas η 637 ; Zonaras 1012 ἦτρον \espTest (ἦτρον) Moeris η 14 \espTest \textbf{514-6} Suid. λ 271 {\policeFr ≈} Suid. π 2114 \espTest \textbf{514} Pollux IX 131 (λέων\expo{2} omisso et αὐτὸ ἔκμαγμα pro αὐτέκμαγμα) ; \Eustathe 1857.10, 462.40 ; cf. \AelDion ε 26 \espTest \textbf{515-6} (πόσθιον -- κύτταρον) Suidas κ 2787 \espTest \textbf{516} (στρεβλὸν) Hesychius σ 1977 στρεβλὸν \espTest {(κύτταρον)} Hesychius κ 4747 ; cf. Photius κ 1270 \espTest \textbf{527-30} Suidas α 1073 \espTest \textbf{528-30} \Stobee III 13.29 \espTest \textbf{531-2} Suidas α 1349 (sed ἄλλαι pro ἄρ᾽ εἰ) \espTest \textbf{536-7} (εἰ μὲν -- λαβοῦσαι) scholium\expo{bT} Homerus \italique{Il.} I 135-7 (ἔστι τις pro τις ἔστιν et αὐταί omisso)\par\vspace{2.5pt}
{\color{ruleGold}\hrule height 0.35pt}\vspace{2.5pt}
\noindent{\bfseries\scshape\color{rubricRed}Suite (en commentaire) de l'apparat critique et paléographique :} \textbf{Correction complexe \italique{in rasura} et supralinéaire au v. 512 :} \Rac{} avait répété εἰσέφερε d'après le v. 509 ; \Rpc{} a gratté le σ et tracé un η au-dessus de εἰ (\lecon{\auDessus{εἰ}{η}\inRas{σ}{}έφερε}), visant ἣ 'φερεν (Bentley). --- \textbf{Correction au v. 514 :} αὖτ' ἔκμαγμα corrigé par \Rpc{} à partir de έγμαγμα. --- \textbf{Correction au v. 549 :} ἂν corrigé par \Rpc{} à partir de ἂρ. --- \textbf{v. 533 :} (γυνὴ τις) \Rpc{} : \lecon{(Μίκα)} Rogers \espApp Ἄγραυλον R (métathèse de liquides ρ/λ) : \lecon{Ἄγλαυρον} Daubuz. --- \textbf{v. 537 :} αὐτοί R : \lecon{αὐταί} Bonini \espApp \lecon{γε} R (τε {\policeGr Σ} T ad Hom. \italique{Il.} I.135).\par}
\end{minipage}\end{adjustbox}
\newpage
\renewcommand{\currHeadLE}{SCHOLIES MARGINALES ({\policeGr Σ}\expo{R}) \& GLOSES INTERLINÉAIRES · FOLIO 13/30}
\renewcommand{\currHeadRE}{f. 168v (Classense 170v) · vv. 509–550 · {\policeGr Σ} 180 à {\policeGr Σ} 188}
\renewcommand{\currHeadLO}{SCHOLIES MARGINALES ({\policeGr Σ}\expo{R}) \& GLOSES INTERLINÉAIRES · FOLIO 13/30}
\renewcommand{\currHeadRO}{f. 168v (Classense 170v) · vv. 509–550 · {\policeGr Σ} 180 à {\policeGr Σ} 188}
\bandeauScholies{Folio 13/30 — f. 168v}{Classense 170v}{vv. 509–550}{Verso (page de gauche) · 9 entrées ({\policeGr Σ} 180 à {\policeGr Σ} 188)}
{\fontsize{8.8pt}{10.5pt}\selectfont\noindent\textbf{\color{rubricRed}Disposition topographique des marges sur le f. 168v :} Marge extérieure à gauche ; marge intérieure à droite. — \textit{Fac-similé de référence :} \texttt{folios\_individuels\_ff162v-177r/13\_f168v\_classense170v\_vv0509-0550.jpg}.\par\vspace{3pt}}
\entreeScholie{180}{509}{τὸ δ᾽ ἐλάκτισεν τὸ ἦτρον τῆς χύτρας.}{Marge supérieure et marge gauche (extérieure, f. 168v, en haut du verso)}{δέον εἰπεῖν «τὸ βρέφος τὸ ἦτρον τῆς μήτρας ἐλάκτισεν» εἶπε τῆς χύτρας, ἐπειδὴ ἐν χύτρᾳ εἰσῆλθε τὸ παιδίον. τὸ κάλυμμα τῆς μήτρας, ὃ καλοῦμεν ἔλυτρον.}{δέον εἰπ(εῖν) «τ(ὸ) βρέφ\expo{ος} τ(ὸ) ἦτρ\expo{ον} τ\expo{ῆς} μήτρ\expo{ας} ἐλάκτισ(εν)» εἶπ(ε) τ\expo{ῆς} χύτρ\expo{ας}, ἐπειδὴ ἐν χύτρᾳ εἰσῆλθ(ε) τ(ὸ) παιδί\expo{ον} · τ(ὸ) κάλυμμα τ\expo{ῆς} μήτρ\expo{ας}, ὃ καλοῦμ(εν) ἔλυτρ\expo{ον} {\policeFr ÷}}{« Alors qu'on attendait “le bébé a donné des coups de pied dans le bas-ventre de la matrice (τῆς μήτρας)”, il a dit par surprise comique “le ventre de la marmite (τῆς χύτρας)”, puisque l'enfant supposé vient d'être apporté dans une marmite ! — L'ἦτρον est l'enveloppe du bas-ventre (la paroi abdominale). »}
\commScholie{Repris par la Souda η 629. Explique à merveille le jeu de mots paronomastique et médical sur μήτρας / χύτρας.}
\entreeScholie{181}{516}{τά τ᾽ ἄλλ᾽ ἁπαξάπαντα καὶ τὸ πόσθιον / τῷ σῷ προσόμοιον, στρεβλὸν ὥσπερ κύτταρον·}{Marge gauche (extérieure, f. 168v, grande scholie alexandrine de 10 lignes en face du v. 516)}{Θεόφραστος τὴν προάνθησίν φησι τῆς πίτυος καὶ τῆς πεύκης. Ἐρατοσθένης ἀγγεῖον τοῦ σφηκὸς ἢ τῆς μελίσσης, ὅπου οἱ σκώληκες ἔπεισιν. ἀλλ’ οὐδὲ οὕτως πάνυ σώζεται ἡ ἀκολουθία. δεῖ οὖν χωρίζειν ἑκάτερον οὕτως, καὶ ὡς στρεβλὸν καὶ ὥσπερ κύτταρον, προσόμοιον κυττάρῳ πλήρει καὶ ἔχοντι τὸν σκώληκα. Ἄλλως. κύτταρος καλεῖται τὰ κοιλώματα τῶν μελικηρίδων· ὁμοίως δὲ καὶ τὸ πῶμα τῆς βαλάνου, ὅπου ἐγκάθηται ἡ βάλανος, ὅπερ νῦν ἀκουστέον· περὶ γὰρ αἰδοίου λέγει.}{Θεόφραστ\expo{ος} τ(ὴν) προάνθησίν φησ(ι) τ\expo{ῆς} πίτυ\expo{ος} \textbf{ς} τ\expo{ῆς} πεύκης · Ἐρατοσθένης ἀγγεῖ(ον) τ\expo{οῦ} σφηκὸς ἢ τ\expo{ῆς} μελίσσ\expo{ης}, ὅπ\expo{ου} οἱ σκώληκες ἔπεισ(ιν) · ἀλλ’ οὐδὲ οὕτ\expo{ως} πάνυ σώζετ(αι) ἡ ἀκολουθία · δεῖ οὖν χωρίζε(ιν) ἑκάτερ(ον) οὕτ\expo{ως}, \textbf{ς} ὡς στρεβλ(ὸν) \textbf{ς} ὥσπ(ερ) κύτταρ(ον), προσόμοι(ον) κυττάρῳ πλήρει \textbf{ς} ἔχοντι τ(ὸν) σκώληκα · ἄλλ\expo{ως} · κύτταρ\expo{ος} καλεῖτ(αι) τὰ κοιλώματα τ\expo{ῶν} μελικηρίδων· ὁμοίως δ(ὲ) \textbf{ς} τ(ὸ) πῶμα τ\expo{ῆς} βαλάνου, ὅπ\expo{ου} ἐγκάθητ(αι) ἡ βάλαν\expo{ος}, ὅπερ νῦν ἀκουστέον· π(ερὶ) γ(ὰρ) αἰδοίου λέγει {\policeFr ÷}}{« Théophraste (Histoire des plantes III, 3, 8) dit que le κύτταρος est le chaton (la préfloraison conique) du pin et du sapin. Ératosthène dit que c'est l'alvéole de la guêpe ou de l'abeille, dans lequel se logent les larves ; mais même ainsi la suite des idées n'est pas tout à fait claire. Il faut donc séparer chaque terme : à la fois “tordu” (στρεβλόν) et “pareil à un alvéole plein contenant sa larve”. — Autre rédaction : On appelle κύτταρος les cavités des rayons de miel, et pareillement la cupule du gland de chêne (τὸ πῶμα τῆς βαλάνου), dans laquelle est enchâssé le gland : c'est en ce dernier sens qu'il faut l'entendre ici, car il parle du gland de la verge enfantine (βαλάνου / αἰδοίου). »}
\commScholie{L'une des plus belles scholies lexicographiques du Ravennas (reprise par la Souda κ 2789, Photius κ 1244 et Hésychius κ 4738), citant Théophraste et le traité Περὶ τῆς ἀρχαίας κωμῳδίας d'Ératosthène de Cyrène, avant de donner la vraie clé du jeu de mots anatomique sur βάλανος (gland du chêne / gland du pénis) et sa cupule (κύτταρος / prépuce).}
\entreeScholie{182}{519}{κᾆτ᾽ Εὐριπίδῃ θυμούμεθα, / οὐδὲν παθοῦσαι μεῖζον ἢ δεδράκαμεν;}{Marge gauche (extérieure, f. 168v, en face des vv. 518–519)}{παρὰ τὰ ἐκ Τηλέφου Εὐριπίδου «εἰ δὴ θυμούμεθα παθόντες οὐδὲν μᾶλλον ἢ δεδρακότες.»}{παρὰ τὰ ἐκ Τηλέφ\expo{ου} Εὐριπίδ\expo{ου} «εἰ δὴ θυμούμεθ(α) παθόντ\expo{ες} οὐδ(ὲν) μᾶλλ\expo{ον} ἢ δεδρακότ\expo{ες}» {\policeFr ÷}}{« Parodie des vers du Télèphe d'Euripide (fr. 711 Kannicht) : “Et après cela nous nous indignons, alors que nous n'avons rien subi de plus que ce que nous avons fait nous-mêmes !” »}
\commScholie{Cité par Louis Hauseux (vv 1-641.tex, l. 3810) : Mnésiloque conclut habilement son plaidoyer en retournant contre les femmes le vers par lequel Télèphe apaisait les chefs achéens.}
\entreeScholie{183}{529}{ὑπὸ λίθῳ γὰρ παντί που / χρὴ μὴ δάκῃ ῥήτωρ ἀθρεῖν.}{Marge gauche (extérieure, f. 168v, en face des vv. 529–530)}{ἐκ τῶν εἰς Πράξιλλαν ἀναφερομένων «ὑπὸ παντὶ λίθῳ σκορπίον ὦ ’ταῖρε φυλάσσεο.» καὶ ἑτέρα «πάντα λίθον κίνει.»}{ἐκ τ\expo{ῶν} εἰς Πράξιλλ(αν) ἀναφερομέν\expo{ων} «ὑπ(ὸ) παντὶ λίθῳ σκορπί\expo{ον} ὦ 'ταῖρε φυλάσσεο» · \textbf{ς} ἑτέρα «πάντα λίθ\expo{ον} κίνει» {\policeFr ÷}}{« Parodie des chansons à boire (σκόλια) attribuées à la poétesse Praxilla de Sicyone (PMG 750 = Scolion 20 Page) : “Sous chaque pierre, ami, prends garde au scorpion !” Et cet autre proverbe : “Remue toute pierre.” »}
\commScholie{Cité par Louis Hauseux (vv 1-641.tex, l. 3904 n. 5) avec le parallèle d'Athénée XV, 695d : Aristophane substitue παρ' ὑπόνοιαν le mot ῥήτωρ (« un orateur venimeux ») à σκορπίος (« un scorpion »).}
\entreeScholie{184}{531}{ἀλλ᾽ οὐ γάρ ἐστι τῶν ἀναισχύτων [ἀναισχύντων R] φύσει γυναικῶν / οὐδὲν κάκιον εἰς ἅπαντα πλὴν ἄρ᾽ εἰ γυναῖκες·}{Marge intérieure droite (note métrique d'Héliodore, f. 168v, v. 531)}{τοῦτο ἰαμβικὸν τετράμετρον καταλῆγον.}{τοῦτ(ο) ἰαμβικ(ὸν) τετράμετρ\expo{ον} καταλῆγ\expo{ον} {\policeFr ÷}}{« Ce vers (inaugurant l'épirrhème de la Coryphée) est un tétramètre iambique catalectique. »}
\commScholie{Scholie métrique ancienne (issue de la colométrie d'Héliodore) signalant le passage aux tétramètres iambiques catalectiques (vv. 531-570). Au v. 531, le mot γυναικῶν qui fausse le mètre dans R\expo{1} est une glose intruse supprimée par Brunck et Bibi.}
\entreeScholie{185}{532}{πλὴν ἄρ᾽ εἰ γυναῖκες.}{Interlinéaire (supra lineam au-dessus de πλὴν ἄρ' εἰ γυναῖκες, f. 168v, v. 532)}{εἰ μὴ ἄλλη γυνή.}{εἰ μὴ ἄλλ(η) γυνή {\policeFr ÷}}{« Sinon les femmes elles-mêmes (il n'est rien de pire qu'une femme impudente, sinon une autre femme). »}
\commScholie{Explique la pointe comique du proverbe détourné (cf. L. Hauseux, vv 1-641.tex, l. 3922).}
\entreeScholie{186}{533}{οὔτοι μὰ τὴν Ἄγλαυρον [Ἄγραυλον R], ὦ γυναῖκες, εὖ φρονεῖτε,}{Marge gauche (extérieure, f. 168v, en face du v. 533)}{κατὰ τῆς Ἀγραύλου ὤμνυον, κατὰ δὲ τῆς Πανδρόσου σπανιώτερον· κατὰ δὲ τῆς Ἔρσης οὐχ εὑρήκαμεν.}{κατὰ τ\expo{ῆς} Ἀγραύλ\expo{ου} ὤμνυ\expo{ον}, κατὰ δ(ὲ) τ\expo{ῆς} Πανδρόσ\expo{ου} σπανιώτερ\expo{ον} · κατὰ δ(ὲ) τ\expo{ῆς} Ἔρσ\expo{ης} οὐχ εὑρήκαμ(εν) {\policeFr ÷}}{« Les femmes athéniennes juraient couramment par Agraulos, plus rarement par Pandrosos (cf. Lysistrata v. 439), mais nous n'avons jamais trouvé dans les textes de serment par Hersé (les trois filles de Cécrops). »}
\commScholie{Observation d'érudit alexandrin (Aristarque ou Didyme, citée par L. Hauseux, vv 1-641.tex, l. 3932) fondée sur le dépouillement exhaustif du corpus attique.}
\entreeScholie{187}{536}{εἰ δὲ μή, ἡμεῖς / αὐταί γε καὶ τὰ δουλάρια τέφραν ποθὲν λαβοῦσαι}{Marge gauche (extérieure, f. 168v, en face du v. 536)}{συντιμωρεῖσθαι αὐτὸν βουλόμεναι, ἐπειδὴ ἐπὶ σκηνῆς εἰσίν.}{συντιμωρεῖσθ(αι) αὐτ(ὸν) βουλόμεν(αι), ἐπειδὴ ἐπὶ σκην\expo{ῆς} εἰσίν {\policeFr ÷}}{« Voulant le châtier toutes ensemble, puisqu'elles sont présentes sur scène avec leurs petites servantes. »}
\commScholie{Note scénique sur la menace d'épilation punitive à la cendre chaude (τέφρα) proférée par la Coryphée.}
\entreeScholie{188}{547}{ἐπίτηδες εὑρίσκων λόγους, ὅπου γυνὴ πονηρὰ / ἐγένετο, Μελανίππας ποιῶν Φαίδρας τε·}{Marge gauche (extérieure, bas du f. 168v, scholie métrique d'Héliodore en face du v. 547)}{τοῦτο μόνον τὸ τετράμετρον ἰαμβεῖον ἀνάπαιστον ἔχει τὸν καταλήγοντα. ἄξιον οὖν αὐτὸν τηρῆσαι.}{τοῦτ(ο) μόν\expo{ον} τ(ὸ) τετράμετρ\expo{ον} ἰαμβεῖ\expo{ον} ἀνάπαιστ\expo{ον} ἔχ(ει) τ(ὸν) καταλήγοντ(α) · ἄξι\expo{ον} οὖν αὐτ(ὸν) τηρῆσ(αι) {\policeFr ÷}}{« Ce tétramètre iambique catalectique est le seul de la série à présenter un anapeste au pied pénultième avant la catalexe (Φαίδρας τε· Πη-νελόπην δέ). Il convient donc d'observer cette licence métrique (et de ne pas corriger le texte). »}
\commScholie{Précieuse scholie métrique d'Héliodore (comme en Grenouilles vv. 930, 932, 937) qui protège expressément la leçon transmise contre toute tentation de normalisation métrique.}
\clearpage
\renewcommand{\currHeadLE}{ΑΡΙΣΤΟΦΑΝΟΥΣ ΘΕΣΜΟΦΟΡΙΑΖΟΥΣΑΙ · CODEX RAVENNAS 429 (\R) · FOLIO 14/30}
\renewcommand{\currHeadRE}{f. 169r (Classense 171r) · vv. 551–589 · Collation diplomatique}
\renewcommand{\currHeadLO}{ΑΡΙΣΤΟΦΑΝΟΥΣ ΘΕΣΜΟΦΟΡΙΑΖΟΥΣΑΙ · CODEX RAVENNAS 429 (\R) · FOLIO 14/30}
\renewcommand{\currHeadRO}{f. 169r (Classense 171r) · vv. 551–589 · Collation diplomatique}
\noindent\begin{adjustbox}{max totalheight=30.4cm, max width=\textwidth, keepaspectratio}\begin{minipage}{\textwidth}
\bandeauRav{f. 169r}{Classense 171r}{551–589}{Folio 14/30 · Recto (page de droite) · 15 scholies \& gloses ({\policeGr Σ} 189–203)}
{\fontsize{9.3pt}{10.8pt}\selectfont
\ligneRav{551}{\auDessus{γυ}{ν}}{ἀκούετ᾽ ὦ γυναῖκες οἷ᾽ εἴρηκεν ἡ πανοῦργος}{}
\ligneRav{}{—}{ἡμᾶς ἁπάσας αὖθις αὖ. : καὶ νὴ Δί᾽ οὐδέπω γε}{}
\ligneRav{553}{}{εἴρηχ᾽ ὅσα ξύνοιδ᾽· ἐπεὶ βούλεσθε \difference{πλείον᾽} εἴπω;}{\textbf{553} πλεῖον R : \lecon{πλείον᾽} Bentley}
\ligneRav{554}{\auDessus{γυ}{ν}}{ἀλλ᾽ οὐκ ἂν ἔτ᾽ ἔχοις· ὅσα γὰρ \difference{ᾔδησθ᾽} ἐξέχεας ἅπαντα·}{\textbf{554} ᾔδεις R : \lecon{ᾔδησθ᾽} Hirschig (Ar. \italique{Vesp}.146)}
\ligneRav{555}{μ\auDessus{ν}{Η}}{μὰ Δί᾽ οὐδὲ τήν \difference{γε} μυριοστὴν μοῖραν ὧν ποιοῦμεν·}{\textbf{555} \lecon{add. γε} Faber}
\ligneRav{556}{}{\difference{ἐπεὶ τάδ᾽} \difference{οὐκ} εἴρηχ᾽, ὁρᾷς, ὡς \difference{στλεγγίδας} λαβοῦσαι}{\textbf{556} ἔπειτά γ᾽ R : \lecon{ἐπεὶ τάδ᾽} Dindorf \espApp οὐχ R : \lecon{οὐκ} \Mon \espApp τλεγγίδας R : \lecon{στλεγγίδας} \Mon}
\ligneRav{557}{—}{ἔπειτα σιφωνίζομεν \difference{τὸν κύσθον}. : ἐπιτριβείης·}{\textbf{557} τὸν σῖτον R, Σ, Suid. σ 512 : \lecon{τὸν κύσθον} Bibi cf. v.1114}
\ligneRav{558}{μ\auDessus{ν}{Η}}{\difference{ὥς τ᾽ αὖ τὰ} κρέ᾽ ἐξ Ἀπατουρίων ταῖς μαστροποῖς διδοῦσαι}{\textbf{558} ὥς τ᾽ αὐτὰ R : ὥς ταῦτα \Mon : \lecon{ὥς τ᾽ αὖ τὰ} Zanetti}
\ligneRav{559}{—}{ἔπειτα τὴν γαλῆν φαμεν– : τάλαιν᾽ ἐγώ· φλυαρεῖς·}{}
\ligneRav{560}{μ\auDessus{ν}{Η}}{οὐδ᾽ ὡς τὸν ἄνδρα τῷ πελέκει γυνὴ κατεσπόδησεν,}{\textbf{560} \lecon{τὸν ἄνδρα τῷ πελέκει γυνὴ} R : ἑτέρα τὸν ἄνδρα τῷ πελέκει Suid. κ 996}
\ligneRav{}{}{οὐκ εἶπον· οὐδ᾽ ὡς φαρμάκοις ἑτέρα τὸν ἄνδρ᾽ ἔμηνεν,}{}
\ligneRav{562}{—}{οὐδ᾽ ὡς ὑπὸ τῇ πυέλῳ κατώρυξέν ποτε— : ἐξόλοιο·}{\textbf{562} \lecon{ποτε} R : ποτ᾽ ed. pr.}
\ligneRav{563}{μ\auDessus{ν}{Η}}{Ἀχαρνικὴ τὸν πατέρα. : ταυτὶ δῆτ᾽ ἀνέκτ᾽ ἀκούειν;}{\textbf{563} \lecon{ἀχαρνικὴ} R : Ἁχαρνικὴ Dobree}
\ligneRav{}{μ\auDessus{ν}{Η}}{οὐδ᾽ ὡς σὺ τῆς δούλης τεκούσης ἄρρεν εἶτα σαυτῇ}{}
\ligneRav{565}{}{τοῦθ᾽ ὑπεβάλου, τὸ σὸν δὲ θυγάτριον παρῆκας αὐτῇ·}{}
\ligneRav{}{\auDessus{γυ}{ν}}{οὔ τοι μὰ τὼ θεὼ σὺ καταπροίξει λέγουσα ταυτί,}{}
\ligneRav{567}{—}{ἀλλ᾽ ἐκποκιῶ σου τὰς ποκάδας. : οὐ \difference{δὴ} μὰ Δία σύ γ᾽ ἅψει·}{\textbf{567} \lecon{ποκάδας} R, Suid. : πλοκ- Σ \espApp δὲ R : \lecon{δὴ} Bothe}
\ligneRav{}{\auDessus{γυ}{ν}}{καὶ μὴν ἰδού. : καὶ μὴν ἰδού. : λαβὲ θοἰμάτιον Φιλίστη·}{}
\ligneRav{569}{μ\auDessus{ν}{Η}}{\difference{πρόσθιγε} μόνον, κἀγώ σε νὴ τὴν Ἄρτεμιν— : τί δράσεις;}{\textbf{569} πρόσθες R : πρόσιθι Cobet : \lecon{πρόσθιγε} Willems}
\ligneRav{570}{μ\auDessus{ν}{Η}}{τὸν σησαμοῦνθ᾽ ὃν κατέφαγες, τοῦτον χεσεῖν ποιήσω·}{}
\ligneRav{571}{χορ/}{\difference{παύσασθε} λοιδορούμεναι· καὶ γὰρ γυνή τις ἡμῖν}{\textbf{571} (Χορ.) R : \lecon{(Κορ.)} von Velsen \espApp παύσασθαι R : \lecon{παύσασθε} Μυ2}
\ligneRav{}{}{ἐσπουδακυῖα προστρέχει. πρὶν οὖν ὁμοῦ γενέσθαι,}{}
\ligneRav{}{}{σιγᾶθ᾽, ἵν᾽ αὐτῆς κοσμίως πυθώμεθ᾽ ἅττα λέξει·}{}
\ligneRav{}{κλεισθ/}{φίλαι γυναῖκες ξυγγενεῖς τοὐμοῦ τρόπου·}{}
\ligneRav{575}{}{ὅτι μὲν φίλος εἴμ᾽ ὑμῖν, ἐπίδηλος ταῖς γνάθοις·}{}
\ligneRav{576}{}{γυναικομανῶ γὰρ προξενῶ θ᾽ ὑμῶν ἀεί·}{\textbf{576} \lecon{ὑμῶν} R : ὑμᾶς Suid. π 2541}
\ligneRav{}{}{καὶ νῦν ἀκούσας πρᾶγμα περὶ ὑμῶν μέγα}{}
\ligneRav{}{}{ὀλίγῳ τι πρότερον κατ᾽ ἀγορὰν λαλούμενον,}{}
\ligneRav{}{}{ἥκω φράσων τοῦτ᾽ ἀγγελῶν θ᾽ ὑμῖν, ἵνα}{}
\ligneRav{580}{}{\difference{σκοπῆτε καὶ τηρῆτε} \difference{μὴ καὶ} προσπέσῃ}{\textbf{580} σκοπεἴτε καὶ τηρεἴτε R : \lecon{σκοπῆτε καὶ τηρῆτε} Biset \espApp καὶ μὴ R : \lecon{μὴ καὶ} Porson (\italique{Aristophanica} p. 216)}
\ligneRav{581}{}{\difference{ὑμῖν} \difference{ἀφάρκτοις} πρᾶγμα δεινὸν καὶ μέγα·}{\textbf{581} ἡμῖν R : \lecon{ὑμῖν} \Mon{} \espApp ἀφράκτοις R : \lecon{ἀφάρκτοις} Dindorf (1830, métathèse attique ancienne -ρκ-)}
\ligneRav{}{χορ/}{τί δ᾽ ἔστιν ὦ παῖ; παῖδα γάρ σ᾽ εἰκὸς καλεῖν,}{}
\ligneRav{}{}{ἕως ἂν οὕτως τὰς γνάθους ψιλὰς ἔχῃς·}{}
\ligneRav{584}{κλεισθ/}{Εὐριπίδην \difference{φάσ᾽} ἄνδρα κηδεστήν τινα}{\textbf{584} ἔφας᾽ R : \lecon{φάσ᾽} Zanetti}
\ligneRav{585}{}{\difference{αὑτοῦ} γέροντα δεῦρ᾽ ἀναπέμψαι τήμερον·}{\textbf{585} αὐτου R : \lecon{αὑτοῦ} Biset}
\ligneRav{}{χορ/}{πρὸς ποῖον ἔργον ἢ τίνος γνώμης χάριν;}{}
\ligneRav{}{κλεισθ/}{ἵν᾽ ἅττα βουλεύοισθε καὶ μέλλοιτε δρᾶν,}{}
\ligneRav{}{}{ἐκεῖνος εἴη τῶν λόγων κατάσκοπος·}{}
\ligneRav{}{χορ/}{καὶ πῶς λέληθεν ἐν γυναιξὶν ὢν ἀνήρ;}{}
}
\vspace{3pt}{\color{ruleGold}\hrule height 0.6pt}\vspace{3pt}
{\fontsize{8.0pt}{9.4pt}\selectfont
\noindent{\bfseries\scshape\color{rubricRed}Loci paralleli et testimonia (f. 169r, vv. 551–589) :} \textbf{554} (ᾔδησθ᾽) \EtymGen (A) ᾔδεισθα \espTest \textbf{557} (σιφωνίζομεν τὸν σῖτον) Suidas σ 512 \espTest (σιφωνίζομεν) Pollux VI 19 \espTest \textbf{560} Suidas κ 996 (ἑτέρα pro γυνὴ et κατεσπόδησε vel κατεσπόδωσεν) \espTest \textbf{566} Suidas κ 730 (sed καταπροίξῃ pro καταπροίξει) \espTest \textbf{567} (ἐκποκιῶ σου τὰς ποκάδας) Suidas ε 588 (ἐκπιῶ cod. A) ; Zonaras 672 ἐκποκιῶ \espTest (ἐκποκιῶ) Heschius ε 1637 ἐκποκιῶ ; \EtymMagnum 323.25 \espTest (ποκάδας) Hesychius π 2750 ποκάδες = Photius π 1010 π\{λ\}οκάδες \espTest \textbf{570} (χεσεῖν) \EpimAlphab ε 201 \espTest \textbf{571-2} (πρὶν οὖν -- πυθώμεθ᾽) Suidas ο 292 (κοσμίως omisso) \espTest \textbf{574} (ξυγγενεῖς τοὐμοῦ τρόπου) Suidas ξ 83 \espTest \textbf{576} Suidas π 2541 προξενῶ θ' ὑμᾶς ἀεί\par\vspace{2.5pt}
{\color{ruleGold}\hrule height 0.35pt}\vspace{2.5pt}
\noindent{\bfseries\scshape\color{rubricRed}Suite (en commentaire) de l'apparat critique et paléographique :} \textbf{vv. 556–557 (locus classicus de l'édition Louis Hauseux) :} 1° Au \textbf{v.~556}, après ὡς, \textbf{\R{} écrit τλεγγίδας sans σ initial} (haplographie de σ après -ς, corrigée en στλεγγίδας par \Mon{}) ; 2° au \textbf{v.~557}, \R{}, {\policeGr Σ}\expo{R} et la Souda (σ 512) portent τὸν σῖτον (« siphonnant le blé avec des strigiles ») : Louis Hauseux (Bibi) restitue \lecon{τὸν κύσθον} (par métathèse du σ et iotacisme υ/ι, exactement comme au v.~1114 où \R{} écrit σκύτο pour κύστο = κύσθον et au v.~1188 où \Rac{} écrit τὸ σποστίον). --- \textbf{v. 557 :} τὸν σῖτον R, {\policeGr Σ}, Souda σ 512 : \lecon{τὸν κύσθον} Bibi (cf. v. 1114). --- \textbf{v. 567 :} \lecon{ποκάδας} R, Souda (πλοκάδας {\policeGr Σ}) \espApp δὲ R : \lecon{δὴ} Bothe. --- \textbf{v. 569 :} πρόσθες R (πρόσιθι Cobet) : \lecon{πρόσθιγε} Alphonse Willems (adopté par L. Hauseux). --- \textbf{v. 571 :} (Χορ.) R : \lecon{(Κορ.)} von Velsen \espApp παύσασθαι R (iotacisme -αι/-ε) : \lecon{παύσασθε} \Mon\expo{2}. --- \textbf{v. 581 :} ἡμῖν R : \lecon{ὑμῖν} \Mon{} \espApp ἀφράκτοις R : \lecon{ἀφάρκτοις} Dindorf (1830, métathèse attique ancienne -ρκ-).\par}
\end{minipage}\end{adjustbox}
\newpage
\renewcommand{\currHeadLE}{SCHOLIES MARGINALES ({\policeGr Σ}\expo{R}) \& GLOSES INTERLINÉAIRES · FOLIO 14/30}
\renewcommand{\currHeadRE}{f. 169r (Classense 171r) · vv. 551–589 · {\policeGr Σ} 189 à {\policeGr Σ} 203}
\renewcommand{\currHeadLO}{SCHOLIES MARGINALES ({\policeGr Σ}\expo{R}) \& GLOSES INTERLINÉAIRES · FOLIO 14/30}
\renewcommand{\currHeadRO}{f. 169r (Classense 171r) · vv. 551–589 · {\policeGr Σ} 189 à {\policeGr Σ} 203}
\bandeauScholies{Folio 14/30 — f. 169r}{Classense 171r}{vv. 551–589}{Recto (page de droite) · 15 entrées ({\policeGr Σ} 189 à {\policeGr Σ} 203)}
{\fontsize{8.8pt}{10.5pt}\selectfont\noindent\textbf{\color{rubricRed}Disposition topographique des marges sur le f. 169r :} Marge extérieure à droite ; marge intérieure à gauche. — \textit{Fac-similé de référence :} \texttt{folios\_individuels\_ff162v-177r/14\_f169r\_classense171r\_vv0551-0589.jpg}.\par\vspace{3pt}}
\entreeScholie{189}{556}{ἐπεὶ τάδ᾽ οὐκ εἴρηχ᾽, ὁρᾷς, ὡς στλεγγίδας λαβοῦσαι}{Interlinéaire (supra lineam au-dessus de τλεγγίδας de R\expo{1}, f. 169r, v. 556)}{ξύστρας.}{ξύστρ\expo{ας} {\policeFr ÷}}{« Des strigiles (grattoirs de bain ou pailles creuses à siphonner). »}
\commScholie{R\expo{1} ayant écrit τλεγγίδας (corrigé en στλεγγίδας dans M), la glose ξύστρας identifie l'instrument.}
\entreeScholie{190}{557}{ἔπειτα σιφωνίζομεν τὸν κύσθον [τὸν σῖτον R, Σ].}{Marge droite (extérieure, f. 169r, en face du v. 557)}{κατεχρήσατο εἰποῦσα σιφωνίζομεν σῖτον. κυρίως γὰρ σιφωνίζειν τὸ τὰ ὑγρὰ ἀποσπᾶν.}{κατεχρήσατ(ο) εἰποῦσα σιφωνίζομ(εν) σῖτ\expo{ον} · κυρί\expo{ως} γ(ὰρ) σιφωνίζ(ειν) τ(ὸ) τὰ ὑγρὰ ἀποσπᾶν {\policeFr ÷}}{« Il a commis une catachrèse comique en disant “nous siphonnons le blé” (σιφωνίζομεν τὸν σῖτον) : car au sens propre, σιφωνίζειν ne se dit que pour soutirer des liquides (du vin) avec un siphon. »}
\commScholie{Longuement discuté par Louis Hauseux (Bibi, v. 557, vv 1-641.tex, ll. 4013-4052) : l'incohérence relevée par le scholiaste entre les strigiles/siphons et le « blé » (τὸν σῖτον) conduit L. Hauseux à proposer la restitution grivoise τὸν κύσθον (cf. v. 1114) ou τὴν μᾶζαν glosée par τὸν σῖτον.}
\entreeScholie{191}{560}{οὐδ᾽ ὡς τὸν ἄνδρα τῷ πελέκει γυνὴ κατεσπόδησεν,}{Interlinéaire (sur κατεσπόδησεν) et marge droite (extérieure, f. 169r, v. 560)}{κατέκοψεν. τοῦτο διὰ τὴν Κλυταιμνήστραν, οὐκ ἀπὸ ἱστορίας.}{κατέκοψ(εν) {\policeFr ÷} // τοῦτ(ο) διὰ τ(ὴν) Κλυταιμνήστρ(αν), οὐκ ἀπ(ὸ) ἱστορί\expo{ας} {\policeFr ÷}}{« κατεσπόδησεν : elle l'a abattu, réduit en poussière (assommé à coups de hache). — Ceci est dit par allusion à Clytemnestre (assassinant Agamemnon), et non d'après un fait divers historique contemporain. »}
\commScholie{Repris par la Souda κ 1068.}
\entreeScholie{192}{561}{οὐκ εἶπον· οὐδ᾽ ὡς φαρμάκοις ἑτέρα τὸν ἄνδρ᾽ ἔμηνεν,}{Marge droite (extérieure, f. 169r, en face du v. 561)}{οὐκ ἀπὸ ἱστορίας παλαιᾶς εἴληφεν, ἀλλ’ ὡς ἐν τῇ Ἀττικῇ τούτου γενομένου.}{οὐκ ἀπ(ὸ) ἱστορί\expo{ας} παλαιᾶς εἴληφ(εν), ἀλλ' ὡς ἐν τῇ Ἀττικῇ τούτ\expo{ου} γενομέν\expo{ου} {\policeFr ÷}}{« Ceci n'est pas tiré d'une légende ancienne (mythologique, bien qu'on pense à Déjanire), mais rapporté comme un fait divers survenu en Attique (cf. l'affaire de la marâtre d'Antiphon, Contre la belle-mère). »}
\commScholie{Le scholiaste distingue l'exemple tragique du v. 560 (Clytemnestre) des faits divers criminels athéniens évoqués aux vv. 561-565.}
\entreeScholie{193}{562}{οὐδ᾽ ὡς ὑπὸ τῇ πυέλῳ κατώρυξέν ποτε—}{Marge droite (extérieure, f. 169r, en face du v. 562)}{ὡς καὶ τούτου πάλιν γενομένου.}{ὡς \textbf{ς} τούτ\expo{ου} πάλ(ιν) γενομέν\expo{ου} {\policeFr ÷}}{« Ceci encore est mentionné comme un fait divers réel (la femme d'Acharnes qui avait enterré son mari assassiné sous la baignoire). »}
\commScholie{Note explicative de R\expo{2} prolongeant la remarque de Σ 561 : le meurtre du mari enterré sous la baignoire (ὑπὸ τῇ πυέλῳ) fait allusion à un fait divers criminel récent à Athènes et non à un mythe tragique.}
\entreeScholie{194}{564}{Ἀχαρνικὴ τὸν πατέρα}{Marge droite (extérieure, f. 169r, en face du v. 564)}{ὡς καὶ τούτου πάλιν γενομένου.}{ὡς \textbf{ς} τούτ\expo{ου} πάλ(ιν) γενομέν\expo{ου} {\policeFr ÷}}{« Comme s'étant là encore réellement produit. »}
\commScholie{Suite de la note précédente (Σ 562) appliquée au parricide commis par une femme du dème d'Acharnes (Ἀχαρνικὴ τὸν πατέρα, v. 563), présenté lui aussi comme un fait divers contemporain.}
\entreeScholie{195}{565}{οὐδ᾽ ὡς σὺ τῆς δούλης τεκούσης ἄρρεν εἶτα σαυτῇ / τοῦθ᾽ ὑπεβάλου, τὸ σὸν δὲ θυγάτριον παρῆκας αὐτῇ·}{Marge droite (extérieure, f. 169r, en face du v. 565)}{ἐπιτετήδευκεν ὑφαρπάζουσαν τὴν ἑτέραν ποιῆσαι τὸν λόγον, ἐπεὶ ἔλειπε τὸ λαβεῖν.}{ἐπιτετήδευκ(εν) ὑφαρπάζουσ(αν) τ(ὴν) ἑτέρ(αν) ποιῆσ(αι) τ(ὸν) λόγ\expo{ον}, ἐπεὶ ἔλειπ(ε) τ(ὸ) λαβ(εῖν) {\policeFr ÷}}{« Le poète a fait exprès de représenter l'autre femme (Mika) coupant la parole à Mnésiloque avant qu'il n'achève sa phrase, où manquait le verbe “prendre”. »}
\commScholie{Remarque dramaturgique très juste sur l'interruption furieuse de Mika au v. 566.}
\entreeScholie{196}{566}{οὔ τοι μὰ τὼ θεὼ σὺ καταπροίξει λέγουσα ταυτί,}{Marge droite (extérieure, f. 169r, en face du v. 566)}{ὡς τὼ χεῖρε. οὐκέτι δὲ τοῖν θεοῖν, ἀλλὰ ταῖν θεαῖν.}{ὡς τὼ χεῖρε · οὐκέτι δ(ὲ) τοῖν θεοῖν, ἀλλὰ ταῖν θεαῖν {\policeFr ÷}}{« μὰ τὼ θεώ : au nominatif-accusatif duel féminin, on emploie la forme τώ comme dans τὼ χεῖρε ; mais au génitif-datif duel, les Attiques disent ταῖν θεαῖν (ou τοῖν θεοῖν). »}
\commScholie{Observation morphologique d'atticiste sur les formes du duel féminin de l'article en attique.}
\entreeScholie{197}{567}{ἀλλ᾽ ἐκποκιῶ σου τὰς ποκάδας.}{Marge droite (extérieure, f. 169r, bloc de 5 lignes en face du v. 567)}{ὅτι ἐπὶ γυναικείων τριχῶν ὡς τὸ πολὺ, καὶ ὅτι παρετυμολογεῖ τὸν πόκον οὐ κακῶς. Λυκόφρων δὲ ἁπλῶς τὰς τρίχας. ἐμπλοκὴ δέ τις φαίνεται, καὶ ὡς πλοκαμῖδες σπανίως ποτὲ ἐπὶ ἀνδρὸς λέγεται.}{ὅτι ἐπὶ γυναικείων τριχ\expo{ῶν} ὡς τ(ὸ) πολὺ, \textbf{ς} ὅτι παρετυμολογ(εῖ) τ(ὸν) πόκον οὐ κακ\expo{ῶς} · Λυκόφρ\expo{ων} δ(ὲ) ἁπλ\expo{ῶς} τ\expo{ὰς} τρίχας · ἐμπλοκὴ δέ τις φαίνετ(αι), \textbf{ς} ὡς πλοκαμῖδες σπ(αν)ίως ποτὲ ἐπὶ ἀνδρὸς λέγετ(αι) {\policeFr ÷}}{« Notez que ποκάδες s'emploie le plus souvent pour la chevelure des femmes, et qu'il joue non sans esprit sur l'étymologie de πόκος (“toison de laine” -> ἐκ-ποκιῶ, “je vais te plumer la toison”). Lycophron (dans son traité Sur la Comédie) dit que ποκάδες désigne simplement les cheveux ; mais il semble y avoir une idée d'entrelacement (πλοκή), et de même que πλοκαμῖδες, le mot se dit parfois aussi d'un homme. »}
\commScholie{Précieuse scholie alexandrine (citée par L. Hauseux, vv 1-641.tex, l. 4149, et reprise par la Souda π 1874) qui conserve un fragment du traité Περὶ κωμῳδίας de Lycophron de Chalcis (IIIe s. av. J.-C.) et confirme la leçon ποκάδας de R et de Bibi.}
\entreeScholie{198}{572}{παύσασθε λοιδορούμεναι· καὶ γὰρ γυνή τις ἡμῖν / ἐσπουδακυῖα προστρέχει. πρὶν οὖν ὁμοῦ γενέσθαι,}{Marge droite (extérieure, f. 169r, en face du v. 572)}{ἀντὶ τοῦ ἐγγὺς παρὰ τοῖς Ἀττικοῖς. καὶ Μένανδρος «ἔστιν δ’ ὁμοῦ τὸ χρῆμα.»}{ἀν\auDessus{τ}{ι} τ\expo{οῦ} ἐγγὺς παρὰ τ\expo{οῖς} Ἀττικ\expo{οῖς} · \textbf{ς} Μένανδρ\expo{ος} «ἔστ(ιν) δ' ὁμοῦ τ(ὸ) χρῆμα» {\policeFr ÷}}{« ὁμοῦ : au sens de “tout près” (ἐγγύς) chez les écrivains attiques. Ainsi Ménandre (fr. 850 Koerte) : “L'affaire (ou l'accouchement) est tout proche.” »}
\commScholie{Note d'atticiste classique (cf. Harpocraton et Photius s.v. ὁμοῦ) illustrée par un hémistiche de Ménandre.}
\entreeScholie{199}{574}{φίλαι γυναῖκες, ξυγγενεῖς τοὐμοῦ τρόπου,}{Interlinéaire (supra lineam, f. 169r, v. 574, entrée de Clisthène)}{τὰ αὐτά μοι πράττουσα.}{τὰ αὐτά μοι πράττουσ(α) {\policeFr ÷}}{« Qui avez les mêmes mœurs et les mêmes occupations que moi (dit par l'efféminé Clisthène). »}
\commScholie{Glose interlinéaire de R\expo{2} expliquant l'expression comique ξυγγενεῖς τοὐμοῦ τρόπου (« parentes de ma façon de vivre ») par laquelle l'efféminé Clisthène revendique sa solidarité de mœurs et de pratiques avec les femmes.}
\entreeScholie{200}{575}{ὅτι μὲν φίλος εἰμ᾽ ὑμῖν, ἐπίδηλος ταῖς γνάθοις·}{Marge droite (extérieure, f. 169r, en face du v. 575)}{ἐψίλωτο γὰρ ὁ Κλεισθένης τὰς γνάθους γυναικωδῶς.}{ἐψίλωτ(ο) γ(ὰρ) ὁ Κλεισθέν(ης) τ\expo{ὰς} γνάθ\expo{ους} γυναικωδ\expo{ῶς} {\policeFr ÷}}{« Car Clisthène avait les joues complètement épilées à la manière des femmes. »}
\commScholie{Explique le comique visuel du masque glabre de Clisthène (cible favorite d'Aristophane depuis les Acharniens v. 118 et les Nuées v. 355).}
\entreeScholie{201}{576}{γυναικομανῶ γὰρ προξενῶ θ᾽ ὑμῶν ἀεί.}{Interlinéaire (supra lineam au-dessus de προξενῶ, f. 169r, v. 576)}{προΐσταμαι.}{προΐσταμ(αι) {\policeFr ÷}}{« Je suis votre représentant attitré (votre patron et défenseur politique, πρόξενος). »}
\commScholie{Glose interlinéaire de R\expo{2} reprise par la Souda π 2539, expliquant l'emploi métaphorique de προξενῶ θ᾽ ὑμῶν (« je suis votre proxène, votre représentant attitré ») par προΐσταμαι (« je prends votre défense »).}
\entreeScholie{202}{578}{ὀλίγῳ τι πρότερον}{Interlinéaire (supra lineam, f. 169r, v. 578)}{πρὸ ὀλίγου.}{πρὸ ὀλίγ\expo{ου} {\policeFr ÷}}{« Il y a un petit instant (tout à l'heure, sur l'Agora). »}
\commScholie{Glose temporelle de R\expo{2} expliquant la locution adverbiale attique ὀλίγῳ τι πρότερον (« tout à l'heure, il y a bien peu de temps », avec datif de mesure) par le tour usuel πρὸ ὀλίγου.}
\entreeScholie{203}{585}{δεῦρ᾽ ἀναπέμψαι}{Marge droite (extérieure, bas du f. 169r, en face du v. 585)}{ὅτι ἀναπέμψαι κυρίως. διὸ καὶ ἄνοδος ἡ πρώτη λέγεται, παρ’ ἐνίοις καὶ κάθοδος, διὰ τὴν θέσιν τῶν θεσμοφορίων. ἐπεὶ καὶ ἄνοδον τὴν εἰς τὸ θεσμοφόριον ἄφιξιν λέγουσιν· ἐπὶ ὑψηλοῦ γὰρ κεῖται τὸ θεσμοφόριον.}{ὅτι ἀναπέμψ(αι) κυρί\expo{ως} · διὸ \textbf{ς} ἄνοδ\expo{ος} ἡ πρώτ(η) λέγετ(αι), παρ' ἐνί\expo{οις} \textbf{ς} κάθοδ\expo{ος}, διὰ τ(ὴν) θέσ(ιν) τ\expo{ῶν} θεσμοφορί\expo{ων} · ἐπεὶ \textbf{ς} ἄνοδ\expo{ον} τ(ὴν) εἰς τ(ὸ) θεσμοφόρι\expo{ον} ἄφιξ(ιν) λέγουσ(ιν) · ἐπὶ ὑψηλ\expo{οῦ} γ(ὰρ) κεῖτ(αι) τ(ὸ) θεσμοφόρι\expo{ον} {\policeFr ÷}}{« Notez l'emploi très exact du préverbe ἀνα- dans ἀναπέμψαι (“envoyer là-haut”). C'est pourquoi le premier jour de la fête s'appelle l'Anodos (“la Montée”, parfois aussi Kathodos pour la descente aux mégarons souterrains), en raison de la situation topographique du sanctuaire : on appelle en effet Anodos l'arrivée au Thesmophorion, parce que le Thesmophorion est situé sur une hauteur (sur la colline de la Pnyx / vers l'Éleusinion supérieur). »}
\commScholie{Scholie topographique de grande valeur (cf. aussi Σ 623) confirmant que le Thesmophorion athénien occupait un emplacement élevé surplombant l'Agora.}
\clearpage
\renewcommand{\currHeadLE}{ΑΡΙΣΤΟΦΑΝΟΥΣ ΘΕΣΜΟΦΟΡΙΑΖΟΥΣΑΙ · CODEX RAVENNAS 429 (\R) · FOLIO 15/30}
\renewcommand{\currHeadRE}{f. 169v (Classense 171v) · vv. 590–634 · Collation diplomatique}
\renewcommand{\currHeadLO}{ΑΡΙΣΤΟΦΑΝΟΥΣ ΘΕΣΜΟΦΟΡΙΑΖΟΥΣΑΙ · CODEX RAVENNAS 429 (\R) · FOLIO 15/30}
\renewcommand{\currHeadRO}{f. 169v (Classense 171v) · vv. 590–634 · Collation diplomatique}
\noindent\begin{adjustbox}{max totalheight=30.4cm, max width=\textwidth, keepaspectratio}\begin{minipage}{\textwidth}
\bandeauRav{f. 169v}{Classense 171v}{590–634}{Folio 15/30 · Verso (page de gauche) · 12 scholies \& gloses ({\policeGr Σ} 204–215)}
{\fontsize{8.8pt}{10.2pt}\selectfont
\ligneRav{590}{κλεισθ/}{ἄφευσεν αὐτὸν \difference{κἀπέτιλ᾽} Εὐριπίδης}{\textbf{590} \lecon{ἄφευσεν} R : ἀφηῦσεν Meineke \espApp καπέτιλλ᾽ R : \lecon{κἀπέτιλ᾽} Bekker}
\ligneRav{}{}{καὶ τἄλλ᾽ ἅπανθ᾽ ὥσπερ γυναῖκ᾽ ἐσκεύασεν·}{}
\ligneRav{}{μ\auDessus{ν}{Η}}{πείθεσθε τούτῳ ταῦτα; τίς δ᾽ οὕτως ἀνὴρ}{}
\ligneRav{}{}{ἠλίθιος ὅστις τιλλόμενος ἠνείχετο;}{}
\ligneRav{594}{}{οὐκ \difference{οἴομαι 'γωγ᾽} ὦ πολυτιμήτω θεώ·}{\textbf{594} οἴομ᾽ ἔγωγ᾽ R : \lecon{οἴομαι 'γωγ᾽} anonymus Parisinus}
\ligneRav{595}{κλεισθ/}{ληρεῖς· ἐγὼ γὰρ οὐκ ἂν ἦλθον ἀγγελῶν,}{}
\ligneRav{596}{}{εἰ μὴ \difference{᾽πεπύσμην} \difference{ταῦτα} τῶν σάφ᾽ εἰδότων·}{\textbf{596} πεπύσμην R, \PsiUn : \lecon{᾽πεπύσμην} Faber \espApp ταυτὶ R : \lecon{ταῦτα} Bentley}
\ligneRav{}{χορ/}{τὸ πρᾶγμα τουτὶ δεινὸν εἰσαγγέλλεται·}{}
\ligneRav{598}{}{ἀλλ᾽ ὦ γυναῖκες οὐκ \difference{ἐλινύειν} ἐχρῆν,}{\textbf{598} ἐλιννύειν R : \lecon{ἐλινύειν} Bourdin}
\ligneRav{}{}{ἀλλὰ σκοπεῖν τὸν ἄνδρα καὶ ζητεῖν ὅπου}{}
\ligneRav{600}{}{λέληθεν \difference{ἡμᾶς} κρυπτὸς ἐγκαθήμενος·}{\textbf{600} ὑμᾶς R : \lecon{ἡμᾶς} \Mon}
\ligneRav{601}{}{καὶ σὺ \difference{ξυνέξευρ᾽} αὐτόν, ὡς ἂν τὴν χάριν}{\textbf{601} συνέξευρ᾽ R, Suid. π 2540 codd. AG : \lecon{ξυνέξευρ᾽} Suid. codd. FVM}
\ligneRav{602}{}{ταύτην τε κἀκείνην ἔχῃς ὦ πρόξενε·}{\textbf{602} \lecon{ταύτην} R : αὐτὴν Suid.}
\ligneRav{}{κλεισθ/}{φέρ᾽ ἴδω· τίς ἡ πρώτη σύ; : ποῖ τις τρέψεται;}{}
\ligneRav{604}{κλεισθ/}{ζητητέαι γάρ ἐστε. : κακοδαίμων ἐγώ·}{\textbf{604} \lecon{ζητητέαι γάρ ἐστε} R : ζητητέον τἄρ᾽ ἐστί Maas \espApp 604 ante 603 transp. Maas, post 607 collocare malui}
\ligneRav{605}{\auDessus{γυ}{ν}}{ἔμ᾽ \difference{ἥτις} \difference{εἴμ᾽ ἤρου}; Κλεωνύμου γυνή·}{\textbf{605} εἰτις ἤρου; R : \lecon{ἥτις} (Zanetti) \lecon{εἴμ᾽ ἤρου;} Bentley}
\ligneRav{606}{κλεισθ/}{γιγνώσκεθ᾽ ὑμεῖς ἥτις ἔσθ᾽ \difference{ἡδὶ} γυνή;}{\textbf{606} ἥδε R : ἥδ᾽ ἡ \Mon : \lecon{ἡδὶ} Grynaeus}
\ligneRav{}{χορ/}{γιγνώσκομεν δῆτ᾽. ἀλλὰ τὰς ἄλλας ἄθρει·}{}
\ligneRav{}{κλεισθ/}{ἡδὶ δὲ δὴ τίς ἐστιν ἡ τὸ παιδίον}{}
\ligneRav{609}{—}{ἔχουσα; : \difference{τίτθη νὴ Δί᾽} ἐμή. : διοίχομαι·}{\textbf{609} τίτθην η δί᾽ R : \lecon{τίτθη νὴ Δί᾽} Ellebodius}
\ligneRav{610}{κλεισθ/}{αὕτη σὺ ποῖ στρέφει; μέν᾽ αὐτοῦ. τί τὸ κακόν;}{\textbf{610} τί τὸ κακόν; continuat Clistheni R : Parenti tribuit (cum versu sequenti) \Rpc}
\ligneRav{611}{μ\auDessus{ν}{Η}}{ἔασον οὐρῆσαί μ᾽· ἀναίσχυντός \difference{τις εἶ}·}{\textbf{611} εἶ R : \lecon{τις εἶ} Küster}
\ligneRav{612}{κλεισθ/}{σὺ δ᾽ οὖν ποίει τοῦτ᾽. \difference{ἀναμενῶ} γὰρ ἐνθάδε·}{\textbf{612} ἀναμένω R : \lecon{ἀναμενῶ} Grynaeus}
\ligneRav{}{χορ/}{ἀνάμενε δῆτα καὶ σκόπει γ᾽ αὐτὴν σφόδρα·}{}
\ligneRav{}{}{μόνην γὰρ αὐτὴν ὦνερ οὐ γιγνώσκομεν·}{}
\ligneRav{615}{κλεισθ/}{\difference{πολύν} γε χρόνον οὐρεῖς σύ. : νὴ Δί᾽ ὦ μέλε·}{\textbf{615} πολύ R : \lecon{πολύν} Gelenius}
\ligneRav{}{}{στραγγουριῶ γάρ· ἐχθὲς ἔφαγον κάρδαμα·}{}
\ligneRav{}{κλεισθ/}{τί καρδαμίζεις; οὐ βαδιεῖ δεῦρ᾽ ὡς ἐμέ;}{}
\ligneRav{}{μ\auDessus{ν}{Η}}{τί δῆτά μ᾽ ἕλκεις ἀσθενοῦσαν; : εἰπέ μοι,}{}
\ligneRav{619}{—}{τίς \difference{ἔστ᾽} ἀνήρ σοι; : τὸν ἐμὸν ἄνδρα πυνθάνει;}{\textbf{619} ἔσθ᾽ R : \lecon{ἔστ᾽} \Mon}
\ligneRav{620}{}{τὸν δεῖνα γιγνώσκεις, τὸν ἐκ Κοθωκιδῶν·}{\textbf{620} add. signum interrogationis Canini}
\ligneRav{621}{κλεισθ/}{τὸν δεῖνα; ποῖον; : ἔσθ᾽ ὁ δεῖν᾽, ὃς καί ποτε}{\textbf{621} (Κλεισθ.) continuat R : \lecon{(Μνη.)} \Mon}
\ligneRav{}{—}{τὸν δεῖνα τὸν τοῦ δεῖνα— : ληρεῖν μοι δοκεῖς·}{}
\ligneRav{}{—}{ἀνῆλθες ἤδη δεῦρο πρότερον; : νὴ Δία}{}
\ligneRav{624}{—}{\difference{ὁσέτη} γε. : καὶ τίς \difference{σοὐστὶ} συσκηνήτρια;}{\textbf{624} ὅς᾽ ἔτη R : \lecon{ὁσέτη} G. H. Schaefer ad Gregor. Cor. p.65 \espApp σοὐστὶν R : \lecon{σοὐστὶ} Dindorf (σοῦ ᾽στὶ iam Biset)}
\ligneRav{625}{μ\auDessus{ν}{Η}}{ἡ δεῖν᾽ \difference{ἔμοιγ᾽}. : οἴμοι τάλας. οὐδὲν λέγεις·}{\textbf{625} ἔμοιγε R : \lecon{ἔμοιγ᾽} edd. pr.}
\ligneRav{}{\auDessus{γυ}{ν}}{ἄπελθ᾽. ἐγὼ γὰρ βασανιῶ ταύτην καλῶς}{}
\ligneRav{}{}{ἐκ τῶν ἱερῶν τῶν πέρυσι· σὺ δ᾽ ἀπόστηθί μοι,}{}
\ligneRav{}{}{ἵνα μὴ 'πακούσῃς ὢν ἀνήρ. σὺ δ᾽ εἰπέ μοι}{}
\ligneRav{}{}{ὅ τι πρῶτον ἡμῖν τῶν ἱερῶν ἐδείκνυτο·}{}
\ligneRav{630}{μ\auDessus{ν}{Η}}{φέρ᾽ ἴδω, τί \expo{μέντοι} (\text{in mg. }\Rpc) πρῶτον ἦν; : ἐπίνομεν·}{\textbf{630} \lecon{μέντοι} \Rpc in mg., Suid. π 2563 : om. \Rac, \Mon}
\ligneRav{631}{\auDessus{γυ}{ν}}{τί δὲ \difference{μετὰ} τοῦτο δεύτερον; : προὐπίνομεν·}{\textbf{631} \lecon{δὲ} R, Suid. π 2563 : δαὶ Bentley \espApp με R : \lecon{μετὰ} Suid.}
\ligneRav{632}{\auDessus{γυ}{ν}}{ταυτὶ μὲν ἤκουσάς τινος· \difference{τί δ᾽ ἦν τρίτον};}{\textbf{632} τί δὲ τρίτον R : τί δ᾽ αὖ Bentley : τί δαὶ τρίτον Elmsley : \lecon{τί δ᾽ ἦν τρίτον} Teubner}
\ligneRav{}{μ\auDessus{ν}{Η}}{σκάφιον Ξένυλλ᾽ ᾔτησεν· οὐ γὰρ ἦν ἀμίς·}{}
\ligneRav{634}{\auDessus{γυ}{ν}}{οὐδὲν λέγεις· δεῦρ᾽ ἐλθὲ \difference{δεῦρ᾽ ὦ} (δευρ ὦ \Rdeux) \difference{Κλείσθενες}·}{\textbf{634} δεῦρο \Rac, \Mon : \lecon{δευρ ὦ} \Rdeux \espApp Κλεισόσθενες R : \lecon{Κλείσθενες} Bonini}
}
\vspace{3pt}{\color{ruleGold}\hrule height 0.6pt}\vspace{3pt}
{\fontsize{8.0pt}{9.4pt}\selectfont
\noindent{\bfseries\scshape\color{rubricRed}Loci paralleli et testimonia (f. 169v, vv. 590–634) :} \textbf{598} (ἐλινύειν) Hesychius ε 2095 ; cf. Moeris ε 61 (ἐλινύων = ἀναπαυόμενος) {\policeFr ∼} Hesychius ε 2097 \espTest \textbf{601-2} Suidas π 2540 \espTest \textbf{616} Suidas σ 1163 (χθὲς γὰρ pro γάρ ἐχθὲς) \espTest \textbf{620} (Κοθωκιδῶν) Photius κ 859 Κοθωκίδαι (ap Scholium) {\policeFr ∼} Hesychius κ 3218 \espTest \textbf{630-1} Suidas π 2563 (sed προπίνομεν pro προὐ- cod. V) \espTest \textbf{633} Pollux X 45 ἐν Πολυίδῳ (sic) Ἀριετοφάνης· σκάφιον Ξένυλλ᾽ ἤτησεν· οὐ γὰρ ἦν ἁμίς \espTest Photius σ 286 σκάφιον\par\vspace{2.5pt}
{\color{ruleGold}\hrule height 0.35pt}\vspace{2.5pt}
\noindent{\bfseries\scshape\color{rubricRed}Suite (en commentaire) de l'apparat critique et paléographique :} \textbf{Témoin papyrologique :} \textbf{\PsiUn{} (\italique{PSI} XI 1194)} couvre les \textbf{vv. 594–596}. --- \textbf{Supplément marginal au v. 630 :} La particule \lecon{μέντοι} (confirmée par la Souda π 2563) avait été omise dans la colonne par \Rac{} (et reste \textbf{omise dans \Mon{}}) ; elle a été rétablie dans la marge par \Rpc{}. --- \textbf{Dittographie syllabique au v. 634 :} En bas de page, \Run{} écrit Κλεισόσθενες (dittographie de -σο- dans Κλείσθενες, corrigée par Bonini), tandis que \Rdeux{} corrige δεῦρο (\Rac{}, \Mon{}) en \lecon{δευρ ὦ}. --- \textbf{v. 609 :} τίτθην η δί' R (médivision de ν final) : \lecon{τίτθη νὴ Δί'} Ellebodius. --- \textbf{v. 630 :} \lecon{μέντοι} \Rpc{} in mg., Souda π 2563 : om. \Rac{}, M.\par}
\end{minipage}\end{adjustbox}
\newpage
\renewcommand{\currHeadLE}{SCHOLIES MARGINALES ({\policeGr Σ}\expo{R}) \& GLOSES INTERLINÉAIRES · FOLIO 15/30}
\renewcommand{\currHeadRE}{f. 169v (Classense 171v) · vv. 590–634 · {\policeGr Σ} 204 à {\policeGr Σ} 215}
\renewcommand{\currHeadLO}{SCHOLIES MARGINALES ({\policeGr Σ}\expo{R}) \& GLOSES INTERLINÉAIRES · FOLIO 15/30}
\renewcommand{\currHeadRO}{f. 169v (Classense 171v) · vv. 590–634 · {\policeGr Σ} 204 à {\policeGr Σ} 215}
\bandeauScholies{Folio 15/30 — f. 169v}{Classense 171v}{vv. 590–634}{Verso (page de gauche) · 12 entrées ({\policeGr Σ} 204 à {\policeGr Σ} 215)}
{\fontsize{8.8pt}{10.5pt}\selectfont\noindent\textbf{\color{rubricRed}Disposition topographique des marges sur le f. 169v :} Marge extérieure à gauche ; marge intérieure à droite. — \textit{Fac-similé de référence :} \texttt{folios\_individuels\_ff162v-177r/15\_f169v\_classense171v\_vv0590-0634.jpg}.\par\vspace{3pt}}
\entreeScholie{204}{600}{λέληθεν ἡμᾶς κρυπτὸς ἐγκαθήμενος·}{Marge gauche (extérieure, f. 169v, bloc de 5 lignes en haut du verso)}{ἀντὶ τοῦ κεκρυμμένος. καλοῦνται μὲν γὰρ καὶ κρύπται παρὰ Πλάτωνι τῷ φιλοσόφῳ [Leg. 1, 633, A] καὶ παρ’ Εὐριπίδῃ καὶ ἐν ταῖς τῶν Λακεδαιμονίων πολιτείαις. καὶ ἐν Θάσῳ ἀρχή τις κρύπται. ἁπλῶς ἀντὶ τοῦ κεκρυμμένος.}{ἀν\auDessus{τ}{ι} τ\expo{οῦ} κεκρυμμέν\expo{ος} · καλοῦντ(αι) μ(ὲν) γ(ὰρ) \textbf{ς} κρύπτ(αι) παρὰ Πλάτωνι τῷ φιλοσόφῳ \textbf{ς} παρ' Εὐριπίδῃ \textbf{ς} ἐν τ\expo{αῖς} τ\expo{ῶν} Λακεδαιμονί\expo{ων} πολιτεί\expo{αις} · \textbf{ς} ἐν Θάσῳ ἀρχή τις κρύπτ(αι) · ἁπλ\expo{ῶς} ἀν\auDessus{τ}{ι} τ\expo{οῦ} κεκρυμμέν\expo{ος} {\policeFr ÷}}{« (Sur la variante κρυπτός, v. 600, pour ποιός τις οὗτος;) Au sens de “caché” (κεκρυμμένος). On appelle en effet κρύπται (“membres de la Cryptie / inspecteurs secrets”) certains jeunes gens chez Platon le philosophe (Lois I, 633b), chez Euripide et dans les Constitutions des Lacédémoniens (d'Aristote, fr. 538 Rose) ; et à Thasos aussi existe un collège de magistrats appelés κρύπται. Mais ici, simplement au sens d'homme dissimulé en cachette. »}
\commScholie{Scholie alexandrine d'une grande richesse institutionnelle (citant Platon, Aristote et une magistrature thasienne autrement connue par les seules inscriptions de Thasos !), qui atteste l'existence d'une variante antique κρυπτός au v. 600.}
\entreeScholie{205}{602}{ταύτην τε κἀκείνην ἔχῃς, ὦ πρόξενε·}{Interlinéaire (supra lineam au-dessus de Λύκι', f. 169v, v. 602)}{ὦ προστάτα.}{ὦ προστάτα {\policeFr ÷}}{« Ô Protecteur (Apollon Lycien, invoqué comme dieu tutélaire qui écarte le malheur). »}
\commScholie{Glose interlinéaire de R\expo{2} expliquant le titre honorifique ὦ πρόξενε (« ô notre patron, notre hôte tutélaire », en écho au verbe προξενῶ du v. 576) adressé à Clisthène par la Coryphée : « ô notre protecteur » (ὦ προστάτα).}
\entreeScholie{206}{603}{φέρ᾽ ἴδω· τίς ἡ πρώτη σύ; — ποῖ τις τρέψεται; / ζητητέαι γάρ ἐστε. — κακοδαίμων ἐγώ·}{Marge intérieure droite (f. 169v, en face des vv. 603–604)}{Μνησίλοχος τοῦτο διὰ τὸ μέλλειν ἔρευναν γίνεσθαι.}{Μνησίλοχ\expo{ος} τοῦτ(ο) διὰ τ(ὸ) μέλλ(ειν) ἔρευν(αν) γίνεσθ(αι) {\policeFr ÷}}{« Mnésiloque dit cela (en aparté) parce qu'une fouille générale va avoir lieu parmi les femmes. »}
\commScholie{Note dramaturgique : confirme au passage que le scholiaste R\expo{2} appelle constamment le personnage Μνησίλοχος.}
\entreeScholie{207}{612}{ἔασον οὐρῆσαί μ᾽· ἀναίσχυντός τις εἶ· / σὺ δ᾽ οὖν ποίει τοῦτ᾽. ἀναμενῶ γὰρ ἐνθάδε·}{Interlinéaire (supra lineam, f. 169v, v. 612)}{ἀντὶ τοῦ οὔρει.}{ἀν\auDessus{τ}{ι} τ\expo{οῦ} οὔρει {\policeFr ÷}}{« (Sur la variante ou l'explication de l'excuse de Mnésiloque : “J'ai envie d'uriner”) : va uriner. »}
\commScholie{Glose interlinéaire de R\expo{2} explicitant l'impératif euphémique de Clisthène σὺ δ᾽ οὖν ποίει τοῦτ᾽ (« fais donc cela », v. 612) en réponse à la demande de Mnésiloque ἔασον οὐρῆσαί μ᾽ (v. 611) : « c'est-à-dire : urine » (ἀντὶ τοῦ οὔρει).}
\entreeScholie{208}{616}{κάρδαμα ἐχθὲς ἔφαγον.}{Marge gauche (extérieure, f. 169v, bloc de 5 lignes en face du v. 616)}{ἐπισχετικὰ γάρ ἐστιν οὔρου τὰ κάρδαμα καὶ πτύσματος. καὶ διὰ τοῦτό φασι τοὺς Πέρσας αὐτοῖς χρῆσθαι. ὅτι μὲν γὰρ χρῶνται καὶ Ξενοφῶν φησι [Cyrop. 1, 2, 16.] «φυλάττονται πολλὰ οὐρεῖν, ὥσπερ καὶ πτύειν καὶ ἀπομύττεσθαι.»}{ἐπισχετικὰ γάρ ἐστ(ιν) οὔρ\expo{ου} τὰ κάρδαμ(α) \textbf{ς} πτύσματ\expo{ος} · \textbf{ς} διὰ τοῦτό φασ(ι) τ\expo{οὺς} Πέρσ\expo{ας} αὐτ\expo{οῖς} χρῆσθ(αι) · ὅτι μ(ὲν) γ(ὰρ) χρῶντ(αι) \textbf{ς} Ξενοφ(ῶν) φησ(ι) «φυλάττοντ(αι) πολλὰ οὐρ(εῖν), ὥσπ(ερ) \textbf{ς} πτύ(ειν) \textbf{ς} ἀπομύττεσθ(αι)» {\policeFr ÷}}{« Car le cresson alénois (τὰ κάρδαμα, Lepidium sativum) a la propriété de retenir l'urine et la salive (ou au contraire provoque la strangurie, cf. στραγγουριῶ v. 615) ; et c'est pourquoi, dit-on, les Perses en consomment. Que les Perses en mangent, Xénophon l'atteste dans la Cyropédie (I, 2, 8 et 16) : “Ils se gardent d'uriner souvent en public, tout comme de cracher et de se moucher.” »}
\commScholie{Repris par la Souda κ 369 et discuté par Louis Hauseux (vv 1-641.tex, l. 4461) : le cresson passait en médecine humorale pour dessécher les humeurs, mais pris en excès il irritait la vessie (στραγγουρία).}
\entreeScholie{209}{620}{τὸν δεῖνα γιγνώσκεις, τὸν ἐκ Κοθωκιδῶν;}{Marge gauche (extérieure, f. 169v, en face du v. 620)}{Κοθωκίδαι δῆμος τῆς Οἰνηίδος φυλῆς.}{Κοθωκίδ(αι) δῆμ\expo{ος} τ\expo{ῆς} Οἰνηίδ\expo{ος} φυλ\expo{ῆς} {\policeFr ÷}}{« Cothôcides (Κοθωκίδαι) est un dème attique appartenant à la tribu Oénéide (Οἰνηὶς φυλή). »}
\commScholie{Repris par Photius κ 859, Hésychius κ 3218 et la Souda κ 1906 (cité par L. Hauseux, vv 1-641.tex, l. 4492).}
\entreeScholie{210}{622}{τὸν δεῖνα τὸν τοῦ δεῖνα— — ληρεῖν μοι δοκεῖς·}{Marge gauche (extérieure, f. 169v, en face du v. 622)}{οἱ μὴ δυνάμενοί τι λέγειν τοιαῦτα λέγουσιν.}{οἱ μὴ δυνάμενοί τι λέγ(ειν) τοιαῦτ(α) λέγουσ(ιν) {\policeFr ÷}}{« Ceux qui sont pris de court et ne savent que répondre bredouillent des formules de ce genre (“Un tel, le fils d'un tel...”). »}
\commScholie{Note de psychologie comique très fine sur l'embarras de Mnésiloque interrogé sur le nom de son prétendu mari.}
\entreeScholie{211}{623}{ἀνῆλθες ἤδη δεῦρο πρότερον; — νὴ Δία / ὁσέτη γε.}{Marge gauche (extérieure, f. 169v, en face du v. 623)}{καὶ τοῦτο πρὸς τὴν ἄνοδον, ἐπειδὴ ἄνοδος ἦν πρὸς τὸ ἱερόν.}{\textbf{ς} τοῦτ(ο) πρ(ὸς) τ(ὴν) ἄνοδ\expo{ον}, ἐπειδὴ ἄνοδ\expo{ος} ἦν πρ(ὸς) τ(ὸ) ἱερόν {\policeFr ÷}}{« Ceci encore (ἀν-ῆλθες, “es-tu montée”) fait référence à l'Anodos, parce qu'il fallait monter pour accéder au sanctuaire du Thesmophorion. »}
\commScholie{Confirme l'observation topographique de Σ 585.}
\entreeScholie{212}{624}{καὶ τίς σοὐστὶ συσκηνήτρια;}{Marge gauche (extérieure, f. 169v, en face du v. 624)}{φίλη συνδίαιτος. σκηνὰς γὰρ ἑαυταῖς ἐποίουν τὸ ἱερόν.}{φίλ(η) συνδίαιτ\expo{ος} · σκην\expo{ὰς} γ(ὰρ) ἑαυτ\expo{αῖς} ἐποί\expo{ουν} τ(ὸ) ἱερόν {\policeFr ÷}}{« Une amie partageant la même tente (συσκηνήτρια) : car les femmes se dressaient des tentes (des barraques de feuillage, σκηναί) dans l'enceinte du sanctuaire pendant la fête. »}
\commScholie{Renseignement rituel capital sur le campement des femmes mariées sous des tentes (σκηναί, cf. v. 658) à l'intérieur du Thesmophorion.}
\entreeScholie{213}{630}{φέρ᾽ ἴδω, τί μέντοι πρῶτον ἦν; — ἐπίνομεν·}{Marge gauche (extérieure, f. 169v, en face du v. 630)}{διαβάλλει πάλιν τὰς γυναῖκας ὡς μεθύσους.}{διαβάλλ(ει) πάλ(ιν) τ\expo{ὰς} γυναῖκ\expo{ας} ὡς μεθύσ\expo{ους} {\policeFr ÷}}{« Il raille de nouveau les femmes comme portées sur la boisson (en imaginant que le premier rite secret des Thesmophories consiste à boire). »}
\commScholie{Sur R, le mot μέντοι omis au v. 630 par R\expo{1} a été suppléé par R\expo{2}.}
\entreeScholie{214}{631}{τί δὲ μετὰ τοῦτο δεύτερον; — προὐπίνομεν.}{Interlinéaire (supra lineam au-dessus de προὐπίνομεν, f. 169v, v. 631)}{ἀλλήλας ἐδεξιούμεθα.}{ἀλλήλ\expo{ας} ἐδεξιούμεθ(α) {\policeFr ÷}}{« Nous portions des santés les unes aux autres coupe en main (ἐδεξιούμεθα). »}
\commScholie{Glose exacte du verbe sympotique προ-πίνω.}
\entreeScholie{215}{633}{σκάφιον Ξένυλλ᾽ ᾔτησεν· οὐ γὰρ ἦν ἀμίς.}{Marge inférieure et marge gauche (bas du f. 169v, en face du v. 633)}{μὴ εὑρίσκων τί εἴπῃ ὁ κηδεστής φησιν ὅτι ἀντὶ ἀμίδος σκάφιον ᾔτησεν ἡ Ξένυλλα. ἀμὶς δέ ἐστιν οὐροδόχον ἀγγεῖον.}{μὴ εὑρίσκ\expo{ων} τί εἴπ(ῃ) ὁ κηδεστής φησ(ιν) ὅτι ἀν\auDessus{τ}{ι} ἀμίδ\expo{ος} σκάφι\expo{ον} ᾔτησ(εν) ἡ Ξένυλλα · ἀμὶς δέ ἐστ(ιν) οὐροδόχ\expo{ον} ἀγγεῖ\expo{ον} {\policeFr ÷}}{« Ne trouvant plus rien à inventer, le parent par alliance (Mnésiloque) raconte qu'à défaut de pot de chambre (ἀμίς), Xénylla a demandé un bassin à boire (σκάφιον). L'ἀμίς est un vase de nuit destiné à recevoir l'urine. »}
\commScholie{Repris par la Souda α 1590 et σ 546 (commenté par L. Hauseux, vv 1-641.tex, l. 4585).}
\clearpage
\renewcommand{\currHeadLE}{ΑΡΙΣΤΟΦΑΝΟΥΣ ΘΕΣΜΟΦΟΡΙΑΖΟΥΣΑΙ · CODEX RAVENNAS 429 (\R) · FOLIO 16/30}
\renewcommand{\currHeadRE}{f. 170r (Classense 172r) · vv. 635–678 · Collation diplomatique}
\renewcommand{\currHeadLO}{ΑΡΙΣΤΟΦΑΝΟΥΣ ΘΕΣΜΟΦΟΡΙΑΖΟΥΣΑΙ · CODEX RAVENNAS 429 (\R) · FOLIO 16/30}
\renewcommand{\currHeadRO}{f. 170r (Classense 172r) · vv. 635–678 · Collation diplomatique}
\noindent\begin{adjustbox}{max totalheight=30.4cm, max width=\textwidth, keepaspectratio}\begin{minipage}{\textwidth}
\bandeauRav{f. 170r}{Classense 172r}{635–678}{Folio 16/30 · Recto (page de droite) · 8 scholies \& gloses ({\policeGr Σ} 216–223)}
{\fontsize{8.8pt}{10.2pt}\selectfont
\ligneRav{635}{\auDessus{γυ}{ν}}{ὅδ᾽ ἐστὶν \difference{ἁνὴρ} ὃν λέγεις. : τί οὖν ποιῶ;}{\textbf{635} ἀνὴρ R : \lecon{ἁνὴρ} Bentley}
\ligneRav{}{\auDessus{γυ}{ν}}{ἀπόδυσον αὐτόν· οὐδὲν ὑγιὲς γὰρ λέγει·}{}
\ligneRav{637}{μ\auDessus{ν}{Η}}{κἄπειτ᾽ ἀποδύσετ᾽ ἐννέα παίδων μητέρα;}{\textbf{637-40} (Μνη.) R : \lecon{(Μνη.) (637-8) et (Μίκα) (639-40)} Bibi}
\ligneRav{638}{}{\difference{χαλᾷ} ταχέως τὸ στρόφιον ὦναίσχυντε σύ·}{\textbf{638} χάλᾳ R : χάλα Grynaeus : \lecon{χαλᾷ} Bibi}
\ligneRav{639}{\auDessus{γυ}{ν}}{ὡς καὶ στιβαρά τις φαίνεται καὶ καρτερά·}{\textbf{639-40} (Μνη.) 639 et (Γυνή) 640 R : \lecon{(Μίκα)} Bibi}
\ligneRav{640}{}{καὶ νὴ Δία \difference{τιτθούς} γ᾽ ὥσπερ ἡμεῖς οὐκ ἔχει·}{\textbf{640} τίτθος R : \lecon{τιτθούς} ed. Iuntina}
\ligneRav{641}{μ\auDessus{ν}{Η}}{\difference{στερίφη} γάρ εἰμι κοὐκ ἐκύησα πώποτε·}{\textbf{641} στεριφῇ R : \lecon{στερίφη} (sed στεριφὴ) \Mon}
\ligneRav{642}{κλεισθ/}{νῦν· τότε \difference{δὲ} μήτηρ ἦσθα παίδων ἐννέα·}{\textbf{642-3} (Μνησ.) (641-3) R : \lecon{(Κλεισθ.) (642-3)} Bibi \espApp δὴ R : \lecon{δὲ} \Mon}
\ligneRav{}{}{ἀνίστασ᾽ ὀρθός. ποῖ τὸ πέος ὠθεῖς κάτω;}{}
\ligneRav{644}{\auDessus{γυ}{ν}}{\difference{τοδὶ διέκυψε} καὶ μάλ᾽ εὔχρων ὦ τάλαν·}{\textbf{644} τὸ δὶ δὴ ἔκυψε R : \lecon{τοδὶ διέκυψε} Dobree, Dindorf}
\ligneRav{645}{κλεισθ/}{καὶ ποῦ 'στιν; : αὖθις εἰς τὸ πρόσθεν οἴχεται·}{}
\ligneRav{646}{κλεισθ/}{οὐκ ἐνγεταυθί. : ἀλλὰ δεῦρ᾽ ἥκει πάλιν·}{\textbf{646} \lecon{ἐνγεταυθί} R : ἐγγεταυθί Thiersch \espApp \lecon{ἀλλὰ} R : μἀλλὰ Bentley}
\ligneRav{647}{κλεισθ/}{ἰσθμόν τιν᾽ ἔχεις \difference{ὤνθρωπ᾽}· ἄνω τε καὶ κάτω}{\textbf{647} ῶνθρωπε R : ἄνθρωπε \Mon : ἄνθρωπ᾽ Suid. ι 640 : \lecon{ὤνθρωπ᾽} Ellebodius}
\ligneRav{}{}{τὸ πέος διέλκεις πυκνότερον Κορινθίων·}{}
\ligneRav{}{\auDessus{γυ}{ν}}{ὤ μιαρὸς οὗτος· ταῦτ᾽ ἄρ᾽ ὑπὲρ Εὐριπίδου}{}
\ligneRav{650}{—}{ἡμῖν ἐλοιδορεῖτο. : κακοδαίμων ἐγώ,}{}
\ligneRav{651}{}{εἰς οἷ᾽ ἐμαυτὸν \difference{εἰσεκύλισα} πράγματα·}{\textbf{651} εἰσεκύλησα R : \lecon{εἰσεκύλισα} Scaliger}
\ligneRav{}{\auDessus{γυ}{ν}}{ἄγε δὴ τί δρῶμεν; : τουτονὶ φυλάττετε}{}
\ligneRav{653}{}{καλῶς, ὅπως μὴ διαφυγὼν \difference{οἰχήσεται}·}{\textbf{653} οἴχεται R : \lecon{οἰχήσεται} Faber p.788}
\ligneRav{654}{}{ἐγὼ δὲ ταῦτα τοῖς \difference{πρυτάνεσιν} ἀγγελῶ·}{\textbf{654} πρυτάνεσσιν R : \lecon{πρυτάνεσιν} \Mon}
\ligneRav{655}{χορ/}{ἡμᾶς τοίνυν μετὰ τοῦτ᾽ ἤδη τὰς λαμπάδας ἁψαμένας χρὴ}{}
\ligneRav{}{}{ξυζωσαμένας εὖ κἀνδρείως τῶν θ᾽ ἱματίων ἀποδύσας}{}
\ligneRav{657}{}{ζητεῖν, εἴ που κἄλλος τις ἀνὴρ \difference{ἐπελήλυθε}, καὶ περιθρέξαι}{\textbf{657} εἰσελήλυθεν R : \lecon{ἐπελήλυθε} Handley}
\ligneRav{658}{}{τὴν Πύκνα πᾶσαν καὶ τὰς σκηνὰς καὶ τὰς διόδους \difference{διαθρῆσαι}·}{\textbf{658} ἀθρῆσαι R : \lecon{διαθρῆσαι} Daubuz}
\ligneRav{659}{ἡμιχ/}{εἶα δὴ πρώτιστα μὲν χρὴ κοῦφον ἐξορμᾶν πόδα}{\textbf{659} (Ἡμιχ. γυν.) \Rpc}
\ligneRav{660}{}{καὶ διασκοπεῖν σιωπῇ \difference{πανταχῇ}· μόνον δὲ χρὴ}{\textbf{660} πανταχῇ post χρὴ R : \lecon{ante μόνον} Küster}
\ligneRav{}{}{μὴ βραδύνειν, ὡς ὁ καιρός ἐστι μὴ μέλλειν ἔτι,}{}
\ligneRav{662}{}{ἀλλὰ τὴν πρώτην τρέχειν \difference{χρῆν} ὡς τάχιστ᾽ ἤδη κύκλῳ·}{\textbf{662} χρὴ R : \lecon{χρῆν} Bentley : χρῆ μ᾽ Austin}
\ligneRav{}{χορ/}{\hspace{1.1em}εἶά νυν ἴχνευε καὶ μά-}{}
\ligneRav{}{}{\hspace{1.1em} τευε ταχὺ πάντ᾽, εἴ τις ἐν τόποις ἑδραῖος}{}
\ligneRav{}{}{\hspace{1.1em}ἄλλος αὖ λέληθεν ὤν·}{}
\ligneRav{665}{}{\hspace{1.1em}πανταχῇ \difference{δὲ ῥῖψον} ὄμμα,}{\textbf{665} διάρριψον R : \lecon{δὲ ῥῖψον} Hermann}
\ligneRav{666}{}{\hspace{1.1em}καὶ τὰ τῇδε καὶ τὰ δεῦρο}{\textbf{666} \lecon{καὶ τὰ τῇδε} R (post καὶ τὰ τῇδε add. καὶ τὰ κεῖσε Kaibel)}
\ligneRav{}{}{\hspace{1.1em}πάντ᾽ ἀνασκόπει καλῶς·}{}
\ligneRav{667}{}{\hspace{1.1em}ἢν γὰρ \difference{ληφθῇ} δράσας ἀνόσια,}{\textbf{667} μὴ λάθῃ R : \lecon{ληφθῇ} Reisig : με λάθῃ Bergk}
\ligneRav{}{}{\hspace{1.1em}δώσει τε δίκην καὶ πρὸς τούτῳ}{}
\ligneRav{669}{}{\hspace{1.1em}τοῖς ἄλλοις \difference{ἀνδράσιν} (ἅπασιν \inRas{ἅπασιν}{}) ἔσται}{\textbf{669} ἅπασιν \inRas{ἅπασιν}{} ἔσται \Rpc : ἅπασιν ἔσται \Mon : \lecon{ἀνδράσιν ἔσται} Beer ap. Hermann (1845) : ἔσται ἅπασιν Bergk : alii alia}
\ligneRav{670}{}{\hspace{1.1em}παράδειγμ᾽ ὕβρεως}{}
\ligneRav{}{}{\hspace{1.1em}ἀδίκων τ᾽ ἔργων ἀθέων τε τρόπων·}{}
\ligneRav{}{}{\hspace{1.1em}φήσει δ᾽ εἶναί τε θεοὺς φανερῶς,}{}
\ligneRav{}{}{\hspace{1.1em}δείξει τ᾽ ἤδη}{}
\ligneRav{}{}{\hspace{1.1em}πᾶσιν ἀνθρώποις σεβίζειν δαίμονας}{}
\ligneRav{675}{}{\hspace{1.1em}†δικαίως τ᾽ ἐφέποντας † ὅσια καὶ νόμιμα}{\textbf{675} \lecon{ἐφέποντας} R : ἐφέπειν Hermann}
\ligneRav{}{}{\hspace{1.1em}μηδομένους ποιεῖν ὅ τι καλῶς ἔχει·}{}
\ligneRav{}{}{\hspace{1.1em}κἂν μὴ ποιῶσι ταῦτα τοιάδ᾽ ἔσται·}{}
}
\vspace{3pt}{\color{ruleGold}\hrule height 0.6pt}\vspace{3pt}
{\fontsize{8.0pt}{9.4pt}\selectfont
\noindent{\bfseries\scshape\color{rubricRed}Loci paralleli et testimonia (f. 170r, vv. 635–678) :} \textbf{641} (στερίφη) Hesychius σ 1769 στερίφη· στεῖρα, μὴ τεκοῦσα, μηδὲ τίκτουσα = Photius σ 532 ; cf. Moeris σ 46 (στερίφη = στερρά) \espTest \textbf{647-8} Suidas ι 640 \espTest \textbf{659} (εἶα δὴ) Hesychius ε 8 ἔα δή· ἄγε δή = Photius ε 1 = \Syn ε 1 = Suidas ε 3\par\vspace{2.5pt}
{\color{ruleGold}\hrule height 0.35pt}\vspace{2.5pt}
\noindent{\bfseries\scshape\color{rubricRed}Suite (en commentaire) de l'apparat critique et paléographique :} \textbf{Sigle tachygraphique marginal μν\auDessus{η}{σ} au v. 638 et redistribution des répliques par \textbf{Bibi} (vv. 637–643) :} Comme le précise Louis Hauseux (\textit{vv 1-641.tex}, l. 4618), R comporte en marge du v. 638 le sigle tachygraphique \lecon{μν\auDessus{η}{σ}} et aucune indication de changement d'interlocuteur jusqu'au v. 640 (Γυνή). Aux vv. 637–640, Louis Hauseux donne 637–638 à (Μνη.) (en lisant l'indicatif 3\ieme{} pers. \lecon{χαλᾷ} au v. 638) et 639–640 à \lecon{(Μίκα)} ; aux vv. 642–643, il restitue la réplique à \lecon{(Κλεισθ.)}. --- \textbf{Confusion onciale ΙΣ {\policeFr ↔} Π au v. 657 :} εἰσελήλυθεν R corrigé par \textbf{Handley} en \lecon{ἐπελήλυθε}. --- \textbf{Grattage d'un mot entier répété par dittographie au v. 669 :} \Rac{} avait écrit deux fois de suite ἅπασιν ; \Rpc{} a gratté le second (\lecon{ἅπασιν \inRas{ἅπασιν}{} ἔσται}), laissant sur le parchemin un espace blanc très net entre ἅπασιν et ἔσται (Beer ap. Hermann restitue ἀνδράσιν ἔσται). --- \textbf{vv. 637–640 :} μν\auDessus{η}{σ} (ad v. 638, nulla mut. pers. usque ad 640 Γυνή) R : \lecon{(Μνη.) (637–638) et (Μίκα) (639–640)} \lecon{Bibi}. --- \textbf{v. 638 :} χάλᾳ R (χάλα Grynaeus) : \lecon{χαλᾷ} \lecon{Bibi}. --- \textbf{vv. 642–643 :} (Μνησ.) (641–643) R : \lecon{(Κλεισθ.) (642–643)} \lecon{Bibi} \espApp \textbf{v. 642 :} δὴ R : \lecon{δὲ} \Mon{}. --- \textbf{v. 644 :} τὸ δὶ δὴ ἔκυψε R (médivision et iotacisme) : \lecon{τοδὶ διέκυψε} Dobree, Dindorf. --- \textbf{v. 651 :} εἰσεκύλησα R (iotacisme) : \lecon{εἰσεκύλισα} Scaliger. --- \textbf{v. 657 :} εἰσελήλυθεν R : \lecon{ἐπελήλυθε} Handley (confusion onciale ΙΣ/Π).\par}
\end{minipage}\end{adjustbox}
\newpage
\renewcommand{\currHeadLE}{SCHOLIES MARGINALES ({\policeGr Σ}\expo{R}) \& GLOSES INTERLINÉAIRES · FOLIO 16/30}
\renewcommand{\currHeadRE}{f. 170r (Classense 172r) · vv. 635–678 · {\policeGr Σ} 216 à {\policeGr Σ} 223}
\renewcommand{\currHeadLO}{SCHOLIES MARGINALES ({\policeGr Σ}\expo{R}) \& GLOSES INTERLINÉAIRES · FOLIO 16/30}
\renewcommand{\currHeadRO}{f. 170r (Classense 172r) · vv. 635–678 · {\policeGr Σ} 216 à {\policeGr Σ} 223}
\bandeauScholies{Folio 16/30 — f. 170r}{Classense 172r}{vv. 635–678}{Recto (page de droite) · 8 entrées ({\policeGr Σ} 216 à {\policeGr Σ} 223)}
{\fontsize{8.8pt}{10.5pt}\selectfont\noindent\textbf{\color{rubricRed}Disposition topographique des marges sur le f. 170r :} Marge extérieure à droite ; marge intérieure à gauche. — \textit{Fac-similé de référence :} \texttt{folios\_individuels\_ff162v-177r/16\_f170r\_classense172r\_vv0635-0678.jpg}.\par\vspace{3pt}}
\entreeScholie{216}{641}{στερίφη γάρ εἰμι κοὐκ ἐκύησα πώποτε.}{Interlinéaire (supra lineam au-dessus de στερίφη, f. 170r, v. 641)}{ἀντὶ τοῦ στείρα.}{ἀν\auDessus{τ}{ι} τ\expo{οῦ} στεῖρα {\policeFr ÷}}{« στερίφη (“ferme, stérile”) : au sens de στεῖρα (femme stérile qui n'a jamais enfanté). »}
\commScholie{Commenté par Louis Hauseux (vv 642-fin-bis.tex, l. 502) : double sens entre « qui a les chairs dures/fermes » (pour excuser l'absence de seins) et « stérile ».}
\entreeScholie{217}{646}{οὐκ ἐνγεταυθί. — ἀλλὰ δεῦρ᾽ ἥκει πάλιν·}{Marge intérieure gauche (f. 170r, en face du v. 646)}{οὐκ ἐνταῦθά γε. παρατετήρηται καὶ πρότερον.}{οὐκ ἐνταῦθά γε · παρατετήρητ(αι) \textbf{ς} πρότερ\expo{ον} {\policeFr ÷}}{« (Il fait passer son phallus par-dessous entre ses cuisses pour qu'on ne le voie pas par-devant) : “pas de ce côté-ci du moins !” Et cela a déjà été noté précédemment. »}
\commScholie{Note sur le lazzi scénique célèbre du phallus de cuir que Mnésiloque fait passer alternativement d'avant en arrière entre ses jambes pendant l'inspection de Clisthène et de Mika.}
\entreeScholie{218}{648}{ἰσθμόν τιν᾽ ἔχεις ὤνθρωπ᾽· ἄνω τε καὶ κάτω / τὸ πέος διέλκεις πυκνότερον Κορινθίων·}{Marge droite (extérieure, f. 170r, en face des vv. 647–648)}{ἐπεὶ τὰς ναῦς διὰ τοῦ Ἰσθμοῦ εἷλκον Κορίνθιοι, ὥστε μὴ περιέρχεσθαι. τοῦτο δὲ διισθμονίσαι ἐκάλουν.}{ἐπεὶ τ\expo{ὰς} ναῦς διὰ τ\expo{οῦ} Ἰσθμ\expo{οῦ} εἷλκ\expo{ον} Κορίνθι(οι), ὥστ(ε) μὴ περιέρχεσθ(αι) · τοῦτ(ο) δ(ὲ) διισθμονίσ(αι) ἐκάλ\expo{ουν} {\policeFr ÷}}{« Parce que les Corinthiens halaient les navires à travers l'Isthme de Corinthe (sur le diolkos) afin de leur éviter de contourner tout le Péloponnèse ; et ils appelaient cette manœuvre διισθμονίσαι (“faire franchir l'Isthme”). »}
\commScholie{Repris par la Souda ι 632 : explication limpide de la métaphore géographique et nautique appliquée par Clisthène au va-et-vient du phallus de Mnésiloque entre ses jambes.}
\entreeScholie{219}{649}{ὁ μιαρὸς οὗτος· ταῦτ᾽ ἄρ᾽ ὑπὲρ Εὐριπίδου / ἡμῖν ἐλοιδορεῖτο.}{Interlinéaire (supra lineam, f. 170r, v. 649)}{ἀντὶ τοῦ διὰ ταῦτα.}{ἀν\auDessus{τ}{ι} τ\expo{οῦ} διὰ ταῦτ(α) {\policeFr ÷}}{« C'est-à-dire : à cause de cela (sur ταῦτα δῆθ' ὑβρίζετο v. 650). »}
\commScholie{Glose explicative se rapportant à l'accusatif adverbial ταῦτα du v. 650.}
\entreeScholie{220}{656}{ξυζωσαμένας εὖ κἀνδρείως τῶν θ᾽ ἱματίων ἀποδύσας}{Interlinéaire (supra lineam au-dessus de τῶν θ' ἱματίων, f. 170r, v. 656)}{τῶν ἐγκύκλων ἀποδυσαμένας.}{τ\expo{ῶν} ἐγκύκλ\expo{ων} ἀποδυσαμέν\expo{ας} {\policeFr ÷}}{« Ayant quitté nos manteaux de dessus (les ἔγκυκλα, cf. Σ 261 et L. Hauseux, vv 642-fin-bis.tex, l. 558). »}
\commScholie{Confirme la distinction vestimentaire établie par Σ 261.}
\entreeScholie{221}{657}{ζητεῖν, εἴ που κἄλλος τις ἀνὴρ ἐπελήλυθε [εἰσελήλυθεν R], καὶ περιθρέξαι}{Marge droite (extérieure) et interlinéaire sur περιθρέξαι (f. 170r, v. 657)}{ἐπεὶ μόνων γυναικῶν ἐστιν ἡ ἐκκλησία. περιδραμεῖν.}{ἐπεὶ μόν\expo{ων} γυναικ\expo{ῶν} ἐστ(ιν) ἡ ἐκκλησία {\policeFr ÷} // περιδραμ(εῖν) {\policeFr ÷}}{« Puisque l'assemblée des Thesmophories est strictement réservée aux seules femmes. — περιθρέξαι : parcourir au pas de course en faisant le tour (περιδραμεῖν). »}
\commScholie{Double glose sur le motif religieux de l'inspection et sur l'aoriste dorien/attique poétique περιθρέξαι (de περιτρέχω).}
\entreeScholie{222}{658}{τὴν πύκνα πᾶσαν καὶ τὰς σκηνὰς καὶ τὰς διόδους διαθρῆσαι.}{Marge droite (extérieure, f. 170r, en face du v. 658)}{τὴν ἐκκλησίαν παρὰ τὸ πυκνοῦσθαι τοὺς ὄχλους. ἐπεποίηντο γὰρ ἐκκλησίαν. τὰς δὲ σκηνὰς διὰ μέσου εἴρηκε. καὶ γὰρ πρὸς τῇ πυκνὶ σκηναὶ ἦσαν, ὡς καὶ αὐτὸς Ἀριστοφάνης μέμνηται τούτου.}{τ(ὴν) ἐκκλησί(αν) παρὰ τ(ὸ) πυκνοῦσθ(αι) τ\expo{οὺς} ὄχλ\expo{ους} · ἐπεποίηντ(ο) γ(ὰρ) ἐκκλησί(αν) · τ\expo{ὰς} δ(ὲ) σκην\expo{ὰς} διὰ μέσ\expo{ου} εἴρηκ(ε) · \textbf{ς} γ(ὰρ) πρ(ὸς) τῇ πυκνὶ σκηναὶ ἦσ(αν), ὡς \textbf{ς} αὐτ\expo{ὸς} Ἀριστοφάν(ης) μέμνητ(αι) τούτ\expo{ου} {\policeFr ÷}}{« τὴν πύκνα : le lieu de l'assemblée (la Pnyx), ainsi nommé parce que la foule s'y presse de façon serrée (πυκνοῦσθαι). Elles s'étaient en effet constituées en assemblée. Quant aux “tentes” (τὰς σκηνάς), il le dit à la fois pour les tentes du Thesmophorion et parce qu'il y avait des baraques près de la Pnyx, comme Aristophane lui-même le mentionne ailleurs. »}
\commScholie{Repris par la Souda π 3129 : explique le double sens topographique entre la Pnyx civique et le campement de tentes des Thesmophories.}
\entreeScholie{223}{664}{μάτευε ταχὺ πάντ᾽, εἴ τις ἐν τόποις ἑδραῖος}{Interlinéaire (supra lineam, f. 170r, v. 664)}{παρακαθεζόμενος.}{παρακαθεζόμεν\expo{ος} {\policeFr ÷}}{« Assis furtivement auprès de nous (sur ὑποκαθιζόμενος / λαθὼν καθέζεται). »}
\commScholie{Glose interlinéaire de R\expo{2} expliquant l'adjectif prédicatif ἑδραῖος (« posté à demeure, assis en embuscade » dans l'enceinte sacrée, v. 664) par le participe παρακαθεζόμενος.}
\clearpage
\renewcommand{\currHeadLE}{ΑΡΙΣΤΟΦΑΝΟΥΣ ΘΕΣΜΟΦΟΡΙΑΖΟΥΣΑΙ · CODEX RAVENNAS 429 (\R) · FOLIO 17/30}
\renewcommand{\currHeadRE}{f. 170v (Classense 172v) · vv. 679–720 · Collation diplomatique}
\renewcommand{\currHeadLO}{ΑΡΙΣΤΟΦΑΝΟΥΣ ΘΕΣΜΟΦΟΡΙΑΖΟΥΣΑΙ · CODEX RAVENNAS 429 (\R) · FOLIO 17/30}
\renewcommand{\currHeadRO}{f. 170v (Classense 172v) · vv. 679–720 · Collation diplomatique}
\noindent\begin{adjustbox}{max totalheight=30.4cm, max width=\textwidth, keepaspectratio}\begin{minipage}{\textwidth}
\bandeauRav{f. 170v}{Classense 172v}{679–720}{Folio 17/30 · Verso (page de gauche) · 10 scholies \& gloses ({\policeGr Σ} 224–233)}
{\fontsize{9.3pt}{10.8pt}\selectfont
\ligneRav{679}{χορ/}{\hspace{1.1em}αὐτῶν ὅταν ληφθῇ τις †ὅσια δρῶν†,}{\textbf{679} \lecon{αὐτῶν} R : τούτων Richards \espApp \lecon{ὅσια} R : ἀνόσια Zanetti : ὅσια μὴ Burges : ἀνόσιόν τι Hermann}
\ligneRav{680}{}{\hspace{1.1em}μανίαις φλέγων λύσσῃ παράκο-}{}
\ligneRav{}{}{\hspace{1.1em} πος, εἴ τι δρῴη πᾶσιν ἐμφανὴς ὁρᾶν}{}
\ligneRav{683}{}{\hspace{1.1em}\difference{ἔσται} γυναιξὶ καὶ βροτοῖς,}{\textbf{683} ἐστὶν R : \lecon{ἔσται} Bothe \espApp γυναιξὶν : \lecon{-ξὶ} Brunck}
\ligneRav{}{}{\hspace{1.1em}ὅτι τά παράνομα τά τ᾽ ἀνόσια παρὼν}{}
\ligneRav{686}{}{\hspace{1.1em}θεὸς \difference{παραχρῆμ᾽ ἀποτίνεται}·}{\textbf{686} ἀποτίνεται· | παραχρῆμά τε τίνεται R : \lecon{παραχρῆμ᾽ ἀποτίνεται} Hermann}
\ligneRav{687}{χορ/}{ἀλλ᾽ ἔοιχ᾽ ἡμῖν ἅπαντά πως διεσκέφθαι καλῶς·}{\textbf{687} (Κορ.) add. von Velsen \espApp Κορ. add. Velsen \espApp \lecon{πως} R : που Blaydes}
\ligneRav{}{}{οὐχ ὁρῶμεν γοῦν ἔτ᾽ ἄλλον οὐδέν᾽ ἐγκαθήμενον}{}
\ligneRav{689a}{\auDessus{γυ}{ν}}{\difference{ἆ ἆ} (ἀά supra lin. R)}{\textbf{689a} ἀά supra lin. R : \lecon{ἆ ἆ} Thiersch}
\ligneRav{689b}{}{\difference{ποῖ ποῖ} σὺ φεύγεις; οὗτος οὗτος οὐ μενεῖς;}{\textbf{689b} ποῖ R : \lecon{ποῖ ποῖ} Biset}
\ligneRav{690}{}{τάλαιν᾽ ἐγὼ τάλαινα, καὶ τὸ παιδίον}{}
\ligneRav{}{}{ἐξαρπάσας μοι φροῦδος ἀπὸ τοῦ τιτθίου·}{}
\ligneRav{}{μ\auDessus{ν}{Η}}{κέκραχθι· τοῦτο δ᾽ οὐδέποτε σὺ ψωμιεῖς,}{}
\ligneRav{693}{}{ἢν μή μ᾽ \difference{ἀφῆτ᾽} (ἀφε\auDessus{ῖ}{ῆ}τ᾽ \Rdeux)· ἀλλ᾽ ἐνθάδ᾽ ἐπὶ τῶν μηρίων}{\textbf{693} \lecon{ἀφε\auDessus{ῖ}{ῆ}τ᾽} \Rdeux : ἀφεῖτ᾽ \Mon}
\ligneRav{}{}{πληγὲν μαχαίρᾳ τῇδε φοινίας φλέβας}{}
\ligneRav{695}{—}{καθαιματώσει βωμόν. : ὢ τάλαιν᾽ ἐγώ·}{}
\ligneRav{}{}{γυναῖκες, οὐκ ἀρήξετ᾽; οὐ πολλὴν βοὴν}{}
\ligneRav{697}{}{στήσεσθε \difference{καὶ} τροπαῖον, ἀλλὰ τοῦ μόνου}{\textbf{697} καὶ om. R : \lecon{καὶ} suppl. Σ ad \italique{Pl.} v.453 et Greg. Cor. \italique{de dialectis} p.22}
\ligneRav{}{—}{τέκνου με περιόψεσθ᾽ ἀποστερουμένην; : ἔα ἔα·}{}
\ligneRav{700}{}{\hspace{1.1em}ὦ πότνιαι Μοῖραι τί \difference{τόδε} δέρκομαι}{\textbf{700} δὲ R : \lecon{τόδε} Dobree : δὴ Reisig}
\ligneRav{}{}{\hspace{1.1em}νεοχμὸν αὖ τέρας;}{}
\ligneRav{702}{}{ὡς \difference{ἅπαντ᾽ ἄρ᾽ ἐστὶ} τόλμης †\difference{πλέα}† κἀναισχυντίας·}{\textbf{702} ἅπαν γὰρ ἐστὶν R : ἅπαντα Suid., cod. Laur. : \lecon{ἅπαντ᾽ ἄρ᾽ ἐστὶ} Bentley \espApp ἔργα R : \lecon{πλέα} Suid. : δοῦλα cod. Laur. : μεστὰ Porson}
\ligneRav{}{}{οἷον αὖ δέδρακεν ἔργον, οἷον αὖ φίλαι τόδε·}{}
\ligneRav{704}{μ\auDessus{ν}{Η}}{\difference{οἷον}; ὑμῶν \difference{ἐξαρήξω} τὴν ἄγαν αὐθαδίαν·}{\textbf{704} οἷον R : \lecon{οἷον;} Ellebodius \espApp ἐξάρξω R : ἐξαράξει Ellebodius : ἐξαράξω Bentley : \lecon{ἐξαρήξω} Bibi}
\ligneRav{705}{χορ/}{ταῦτα δῆτ᾽ οὐ δεινὰ πράγματ᾽ ἐστὶ καὶ περαιτέρω;}{}
\ligneRav{706}{\auDessus{γυ}{ν}}{δεινὰ δῆθ᾽, \difference{ὁτιή} γ᾽ ἔχει μου 'ξαρπάσας τὸ παιδίον·}{\textbf{706} ὅτι R : \lecon{ὁτιή} Faber : ὅστις Porson}
\ligneRav{}{χορ/}{\hspace{1.1em}τί ἂν οὖν εἴποι πρὸς ταῦτά τις, ὅτε}{}
\ligneRav{}{}{\hspace{1.1em}τοιαῦτα ποιῶν ὅδ᾽ ἀναισχυντεῖ;}{}
\ligneRav{}{μ\auDessus{ν}{Η}}{\hspace{1.1em}κοὔπω μέντοι γε πέπαυμαι·}{}
\ligneRav{710}{χορ/}{\hspace{1.1em}ἀλλ᾽ οὖν \difference{ἥκειν γ᾽} ὅθεν ἥκεις·}{\textbf{710} ἥκεις τ᾿ R : ἥκεις γ᾽ Suid. η 170 : ἥκων γ᾽ Willems : \lecon{ἥκειν γ᾽} Bibi ap. Phot. η 88}
\ligneRav{711}{}{\hspace{1.1em}φαύλως \difference{δ᾿} ἀποδρὰς \difference{οὐκ ἔξει}·}{\textbf{711} τ᾿ R, Suid : \lecon{δ᾿} Bentley \espApp οὐ λέξεις R, Suid. : \lecon{οὐκ ἔξει.} Bibi ap. Σ}
\ligneRav{}{}{\hspace{1.1em}οἷον δράσας διέδυς ἔργον·}{}
\ligneRav{}{}{\hspace{1.1em}λήψει δὲ κακόν·}{}
\ligneRav{}{μ\auDessus{ν}{Η}}{τοῦτο μέντοι μὴ γένοιτο μηδαμῶς, ἀπεύχομαι·}{}
\ligneRav{715}{χορ/}{\hspace{1.1em}τίς οὖν σοι, τίς ἂν σύμμαχος ἐκ θεῶν}{\textbf{715} \lecon{οὖν σοι} R : ἂν σοι Hamaker}
\ligneRav{}{}{\hspace{1.1em}ἀθανάτων ἔλθοι ξὺν ἀδίκοις ἔργοις;}{}
\ligneRav{717}{μ\auDessus{ν}{Η}}{μάτην λαλεῖτε· \difference{τήνδ᾽} ἐγὼ οὐκ ἀφήσω·}{\textbf{717} τὴν δ᾽ R : \lecon{τήνδ᾽} Biset}
\ligneRav{}{χορ/}{\hspace{1.1em}ἀλλ᾽ οὐ μὰ τὼ θεὼ τάχ᾽ οὐ}{}
\ligneRav{719}{}{\hspace{1.1em}χαίρων ἴσως \difference{ἐνυβριεῖς}}{\textbf{719} ἐνυβρίσεις R : \lecon{ἐνυβριεῖς} anonymus Romanus}
\ligneRav{720}{}{\hspace{1.1em}λόγους τε λέξεις ἀνοσίους}{}
}
\vspace{3pt}{\color{ruleGold}\hrule height 0.6pt}\vspace{3pt}
{\fontsize{8.0pt}{9.4pt}\selectfont
\noindent{\bfseries\scshape\color{rubricRed}Loci paralleli et testimonia (f. 170v, vv. 679–720) :} \textbf{696-7} (γυναῖκες -- τροπαῖον) \GregCor p.20-22 cf. schol.\expo{Ald} \italique{Plutus} 453c \espTest \textbf{702} (ἅπαντ' -- κἀναισχυντίας) Suidas α 2902 (sed ἅπαντα pro ἅπαν γὰρ ἐστὶν et πλέα pro ἔργα) ; \italique{Corpus paroemiographorum Graecorum}, Appendices, I, 35 (sed ἅπαντα pro ἅπαν γὰρ ἐστὶν et ἔργα omisso) ; \italique{Laurentiamus} LVIII 24 p. 40 (ed. Cohn) (sed ἅπαντα pro ἅπαν γὰρ ἐστὶν et δοῦλα pro ἔργα) \espTest \textbf{710-3} Suidas η 170 (sed λήψῃ pro λήψει) \espTest \textbf{710} Photius η 88 ἥκειν· ἀντὶ τοῦ ἥκων\par\vspace{2.5pt}
{\color{ruleGold}\hrule height 0.35pt}\vspace{2.5pt}
\noindent{\bfseries\scshape\color{rubricRed}Suite (en commentaire) de l'apparat critique et paléographique :} \textbf{Page particulièrement riche en corruptions d'onciale et de minuscule élucidées par Louis Hauseux (Bibi) :} 1° Au \textbf{v.~686}, dittographie avec glose dans \R{} (ἀποτίνεται· | παραχρῆμά τε τίνεται pour \lecon{παραχρῆμ' ἀποτίνεται} Hermann) ; 2° au \textbf{v.~689}, l'exclamation ἀά est écrite au-dessus de la ligne dans \R{} ; 3° au \textbf{v.~693}, correction supralinéaire de \Rdeux{} \lecon{ἀφε\auDessus{ῖ}{ῆ}τ'} (ἀφεῖτ' \Mon{}) ; 4° au \textbf{v.~702}, ΑΠΑΝΤΑΡ (ἅπαντ' ἄρ') est devenu ἅπαν γὰρ par confusion onciale Τ/Γ et la glose ἔργα a chassé \lecon{πλέα} (Souda) ; 5° trois restitutions consécutives de Bibi aux \textbf{vv.~704, 710 et 711} : \lecon{ἐξαρήξω} (ap. Hésychius ε 3599), \lecon{ἥκειν γ'} (ap. Photius η 88) et \lecon{οὐκ ἔξει.} (ap. {\policeGr Σ}\expo{R} ὅθεν οὐ φεύξῃ, confusion onciale Κ/Λ). --- \textbf{v. 679 :} \lecon{αὐτῶν} R (τούτων Richards) \espApp \lecon{ὅσια} R (maintenu par L. Hauseux avec interrogation ; ἀνόσια Zanetti ; ὅσια μὴ Burges ; ἀνόσιόν τι Hermann). --- \textbf{v. 697 :} καὶ om. R : \lecon{καὶ} suppl. {\policeGr Σ} ad \italique{Plut.} 453 et Grégoire de Corinthe, \italique{De dialectis}, p. 22. --- \textbf{v. 704 :} οἷον R : \lecon{οἷον;} Ellebodius \espApp ἐξάρξω R (ἐξαράξει Ellebodius ; ἐξαράξω Bentley) : \lecon{ἐξαρήξω} Bibi ap. Hésychius ε 3599. --- \textbf{v. 710 :} ἥκεις τ' R (ἥκεις γ' Souda η 170 ; ἥκων γ' Willems) : \lecon{ἥκειν γ'} Bibi ap. Photius η 88 (permutation minuscule ν → ς). --- \textbf{v. 711 :} τ' R, Souda : \lecon{δ'} Bentley \espApp οὐ λέξεις R, Souda : \lecon{οὐκ ἔξει.} Bibi ap. {\policeGr Σ}\expo{R} (ὅθεν οὐ φεύξῃ, confusion onciale Κ/Λ).\par}
\end{minipage}\end{adjustbox}
\newpage
\renewcommand{\currHeadLE}{SCHOLIES MARGINALES ({\policeGr Σ}\expo{R}) \& GLOSES INTERLINÉAIRES · FOLIO 17/30}
\renewcommand{\currHeadRE}{f. 170v (Classense 172v) · vv. 679–720 · {\policeGr Σ} 224 à {\policeGr Σ} 233}
\renewcommand{\currHeadLO}{SCHOLIES MARGINALES ({\policeGr Σ}\expo{R}) \& GLOSES INTERLINÉAIRES · FOLIO 17/30}
\renewcommand{\currHeadRO}{f. 170v (Classense 172v) · vv. 679–720 · {\policeGr Σ} 224 à {\policeGr Σ} 233}
\bandeauScholies{Folio 17/30 — f. 170v}{Classense 172v}{vv. 679–720}{Verso (page de gauche) · 10 entrées ({\policeGr Σ} 224 à {\policeGr Σ} 233)}
{\fontsize{8.8pt}{10.5pt}\selectfont\noindent\textbf{\color{rubricRed}Disposition topographique des marges sur le f. 170v :} Marge extérieure à gauche ; marge intérieure à droite. — \textit{Fac-similé de référence :} \texttt{folios\_individuels\_ff162v-177r/17\_f170v\_classense172v\_vv0679-0720.jpg}.\par\vspace{3pt}}
\entreeScholie{224}{680}{μανίαις φλέγων λύσσῃ παράκοπος,}{Interlinéaire (supra lineam au-dessus de φλέγων / μανίαις φλέγων, f. 170v, v. 680)}{ἀντὶ τοῦ διάπυρος ὢν ταῖς μανίαις.}{ἀν\auDessus{τ}{ι} τ\expo{οῦ} διάπυρ\expo{ος} ὢν τ\expo{αῖς} μανί\expo{αις} {\policeFr ÷}}{« μανίαις φλέγων : au sens d'étant tout enflammé (brûlant) de fureur insensée. »}
\commScholie{Glose l'emploi intransitif poétique de φλέγων dans le chant choral.}
\entreeScholie{225}{682}{πᾶσιν ἐμφανὴς ὁρᾶν / ἔσται [ἐστὶν R] γυναιξὶ καὶ βροτοῖς,}{Marge gauche (extérieure, f. 170v, en face des vv. 682–683)}{γελοίως εἶπεν, ὡς Ἀλέξανδρος «ἦσαν ἄνθρωποι πέντε καὶ γυναῖκες τρεῖς.» πέπαικται οὖν τοῦτο ὡς ἐν κωμῳδίᾳ.}{γελοί\expo{ως} εἶπ(εν), ὡς Ἀλέξανδρ\expo{ος} «ἦσ(αν) ἄνθρωπ(οι) πέντ(ε) \textbf{ς} γυναῖκ\expo{ες} τρεῖς» · πέπαικτ(αι) οὖν τοῦτ(ο) ὡς ἐν κωμῳδίᾳ {\policeFr ÷}}{« Il s'est exprimé de façon plaisante, comme le poète comique Alexandros (fr. 6 K.-A.) disant : “Il y avait là cinq êtres humains (ἄνθρωποι = hommes) et trois femmes.” C'est donc une plaisanterie propre à la comédie. »}
\commScholie{Scholie curieuse conservant une citation unique du poète comique Alexandros.}
\entreeScholie{226}{686}{ὅτι τὰ παράνομα τά τ᾽ ἀνόσια παρὼν / θεὸς παραχρῆμ᾽ ἀποτίνεται [παραχρῆμά τε τίνεται R]·}{Interlinéaire (supra lineam, f. 170v, v. 684)}{ἀποδίδωσιν.}{ἀποδίδωσ(ιν) {\policeFr ÷}}{« Il paie en retour (il expie). »}
\commScholie{Glose interlinéaire de R\expo{2} expliquant le verbe poétique ἀποτίνεται / τίνεται (« fait payer, punit en retour ») par ἀποδίδωσιν.}
\entreeScholie{227}{689}{[σκηνικὴ σημείωσις R\expo{2} :] ἆ ἆ [ἀά supra lin. R]· ποῖ ποῖ σὺ φεύγεις; οὗτος οὗτος οὐ μενεῖς;}{Intra-textuel (dans R\expo{1} au v. 688) et Marge gauche (extérieure, f. 170v, bloc de 7 lignes en face des vv. 688–692)}{[...] ἁρπάζει παιδίον μιᾶς γυναικὸς, καὶ καταφεύγει ἐν τῷ ἱερῷ, λέγων ὅτι ἐάν με τύψητε, κατασφάξω τοῦτο τὸ παιδίον. ἀντὶ δὲ παιδίου ἀσκὸν οἴνου λαμβάνει, καὶ αὐτὸν ἐνδύει κροκωτὸν, καὶ ποιεῖ αὐτῷ χεῖρας καὶ πόδας, καί φησιν αὐτὸ παιδίον. διαβάλλει δὲ τὰς γυναῖκας ὡς ἐν ταῖς μάχαις αὐτῶν τοὺς ἀσκοὺς τῶν οἴνων κατέχουσιν.}{ἁρπάζει παιδίον μι\expo{ᾶς} γυναικὸς, \textbf{ς} καταφεύγει ἐν τῷ ἱερῷ, λέγων ὅτι ἐάν με τύψητε, κατασφάξω τοῦτ(ο) τ(ὸ) παιδίον · ἀν\auDessus{τ}{ι} δ(ὲ) παιδίου ἀσκ(ὸν) οἴν\expo{ου} λαμβάνει, \textbf{ς} αὐτ(ὸν) ἐνδύει κροκωτ(ὸν), \textbf{ς} ποι(εῖ) αὐτῷ χεῖρ\expo{ας} \textbf{ς} πόδας, καί φησ(ιν) αὐτὸ παιδίον · διαβάλλ(ει) δ(ὲ) τ\expo{ὰς} γυναῖκ\expo{ας} ὡς ἐν τ\expo{αῖς} μάχαις αὐτ\expo{ῶν} τ\expo{οὺς} ἀσκοὺς τ\expo{ῶν} οἴν\expo{ων} κατέχουσ(ιν) {\policeFr ÷}}{« [Parodie du Télèphe d'Euripide :] Il (Mnésiloque) arrache le bébé d'une des femmes (Mika) et se réfugie sur l'autel du sanctuaire en criant : “Si vous me frappez, j'égorgerai cet enfant !” Mais au lieu d'un vrai bébé, c'est une outre pleine de vin (ἀσκὸν οἴνου) qu'il saisit, enveloppée d'une petite robe safran et munie de bras et de pieds postiches (chaussés de persiques), qu'elle appelle son bébé. Il raille ainsi les femmes comme tenant des outres de vin jusque dans leurs assemblées et leurs combats. »}
\commScholie{Grande scholie scénique fondamentale (reprise et discutée par L. Hauseux, vv 642-fin-bis.tex, ll. 1483-1510, en regard du célèbre cratère apulien de Würzburg H 5697 qui illustre exactement cette scène de l'outre-bébé coiffée et chaussée sur l'autel !).}
\entreeScholie{228}{693}{ἢν μή μ᾽ ἀφῆτ᾽· ἀλλ᾽ ἐνθάδ᾽ ἐπὶ τῶν μηρίων / πληγὲν μαχαίρᾳ τῇδε φοινίας φλέβας / καθαιματώσει βωμόν.}{Marge gauche (extérieure, f. 170v, en face du v. 693)}{τῶν βωμῶν, ἀπὸ τοῦ τὰ μηρία ἐπάνω ἀποκεῖσθαι. ἐκ δὲ τοῦ περιεχομένου τὸ περιέχον εἶπεν.}{τ\expo{ῶν} βωμ\expo{ῶν}, ἀπ(ὸ) τ\expo{οῦ} τὰ μηρία ἐπάνω ἀποκεῖσθ(αι) · ἐκ δ(ὲ) τ\expo{οῦ} περιεχομέν\expo{ου} τ(ὸ) περιέχ\expo{ον} εἶπ(εν) {\policeFr ÷}}{« ἐπὶ μηρίων : sur les autels, ainsi désignés parce qu'on y dépose par-dessus les cuisseaux des victimes (τὰ μηρία). Par métonymie, il a désigné le contenant (l'autel) par le contenu (les cuisseaux brûlés). »}
\commScholie{Analyse rhétorique exacte du trope paratragique ἐπὶ μηρίων.}
\entreeScholie{229}{697}{στήσεσθε καὶ τροπαῖον,}{Marge gauche (extérieure, f. 170v, en face du v. 697)}{τροπαῖον προπερισπωμένως ἀναγνωστέον παρὰ Ἀριστοφάνει καὶ παρὰ Θουκυδίδῃ, τρόπαιον δὲ προπαροξυτόνως παρὰ τοῖς νεωτέροις ποιηταῖς}{τροπαῖ\expo{ον} προπερισπωμέν\expo{ως} ἀναγνωστέ\expo{ον} παρὰ Ἀριστοφάν(ει) \textbf{ς} παρὰ Θουκυδίδ(ῃ), τρόπαι\expo{ον} δ(ὲ) προπαροξυτόν\expo{ως} παρὰ τ\expo{οῖς} νεωτέρ\expo{οις} ποιητ\expo{αῖς} {\policeFr ÷}}{« Il faut lire et accentuer τροπαῖον comme propérispomène (avec circonflexe sur la pénultième) chez Aristophane et chez Thucydide (selon l'ancien attique), mais τρόπαιον comme proparoxyton chez les poètes plus récents. »}
\commScholie{Célèbre règle d'accentuation d'Hérodien et d'Ælius Dionysius (τ 28 Erbse, cf. scholie au Ploutos v. 453 et Grégoire de Corinthe p. 22), scrupuleusement appliquée par Louis Hauseux (Bibi, v. 697, vv 642-fin-bis.tex, l. 853 : στήσεσθε καὶ τροπαῖον).}
\entreeScholie{230}{702}{ὡς ἅπαντ᾽ ἄρ᾽ ἐστὶ τόλμης †πλέα† [ἔργα Σ] κἀναισχυντίας·}{Marge intérieure droite (f. 170v, en face des vv. 700–701)}{παροιμία.}{παροιμία {\policeFr ÷}}{« C'est un proverbe. »}
\commScholie{Longuement analysé par Louis Hauseux (Bibi, v. 701, vv 642-fin-bis.tex, ll. 877-915) : le mot ἔργα qui fausse le mètre dans R\expo{1} est une glose explicative d'ἅπαντα qui a chassé διπλᾶ (cf. Oiseaux v. 717 et Souda α 1638).}
\entreeScholie{231}{710}{ἀλλ᾽ οὖν ἥκειν γ᾽ ὅθεν ἥκεις,}{Marge gauche (extérieure) et débris de glose intra-textuelle (f. 170v, v. 710)}{ἐκ τῶν ἀνδρῶν φησὶν, ὅθεν οὐκ ἔξεστιν. ὁ Σύμμαχος· γελοίως φησὶν ὅθεν ἥκεις φαύλως ἥκεις. ἥκεις, φησὶν, εἰς τὰς χεῖρας ἡμῶν, ὅθεν οὐ φεύξῃ.}{ἐκ τ\expo{ῶν} ἀνδρ\expo{ῶν} φησ(ὶν), ὅθ(εν) οὐκ ἔξεστ(ιν) · ὁ Σύμμαχ\expo{ος} · γελοί\expo{ως} φησ(ὶν) ὅθ(εν) ἥκ(εις) φαύλ\expo{ως} ἥκ(εις) · ἥκ(εις), φησ(ὶν), εἰς τ\expo{ὰς} χεῖρ\expo{ας} ἡμ\expo{ῶν}, ὅθ(εν) οὐ φεύξῃ {\policeFr ÷}}{« “D'où tu es venu” : de chez les hommes, dit-elle, d'où il n'est pas permis de venir ici. — Symmaque explique : il dit plaisamment “d'où tu es venu, tu es venu pour ton malheur”. — Autre glose (reprise par la Souda η 170) : “Tu es venu tomber entre nos mains, d'où tu ne t'échapperas pas !” »}
\commScholie{Discussion approfondie de Louis Hauseux (Bibi, vv. 710-711, vv 642-fin-bis.tex, ll. 1012-1125) qui croise cette double scholie de Symmaque avec la glose de Photius η 88 (ἥκειν· ἀντὶ τοῦ ἥκων) et la Souda η 170.}
\entreeScholie{232}{714}{λήψει δὲ κακόν. — τοῦτο μέντοι μὴ γένοιτο μηδαμῶς, ἀπεύχομαι·}{Interlinéaire (supra lineam, f. 170v, v. 714)}{τοῦτο τὸ παθεῖν.}{τοῦτ(ο) τ(ὸ) παθ(εῖν) {\policeFr ÷}}{« Subir ce sort (ou commettre cet attentat). »}
\commScholie{Glose interlinéaire de R\expo{2} précisant l'antécédent du démonstratif τοῦτο dans la réplique apotropaïque de Mnésiloque (v. 714) à la menace du Chœur (λήψει δὲ κακόν, v. 713) : « le fait de subir ce malheur » (τοῦτο τὸ παθεῖν).}
\entreeScholie{233}{717}{μάτην λαλεῖτε· τήνδ᾽ ἐγὼ οὐκ ἀφήσω·}{Marge inférieure du f. 170v (en face des vv. 716–717)}{κάθηται ἐπὶ βωμοῦ ὡς ἱκετεύων, ἡρπακώς τινος θυγάτριον· φαίνεται δὲ ἀσκὸν ἔχων.}{κάθητ(αι) ἐπὶ βωμ\expo{οῦ} ὡς ἱκετεύ\expo{ων}, ἡρπακώς τιν\expo{ος} θυγάτρι\expo{ον} · φαίνετ(αι) δ(ὲ) ἀσκ(ὸν) ἔχ\expo{ων} {\policeFr ÷}}{« Il est assis sur l'autel en posture de suppliant (comme Télèphe à Argos), après avoir arraché la petite fille de l'une des femmes ; mais aux yeux des spectateurs, on voit que c'est une outre qu'il tient dans les bras. »}
\commScholie{Note de mise en scène précisant le jeu de Mnésiloque sur l'autel.}
\clearpage
\renewcommand{\currHeadLE}{ΑΡΙΣΤΟΦΑΝΟΥΣ ΘΕΣΜΟΦΟΡΙΑΖΟΥΣΑΙ · CODEX RAVENNAS 429 (\R) · FOLIO 18/30}
\renewcommand{\currHeadRE}{f. 171r (Classense 173r) · vv. 721–764 · Collation diplomatique}
\renewcommand{\currHeadLO}{ΑΡΙΣΤΟΦΑΝΟΥΣ ΘΕΣΜΟΦΟΡΙΑΖΟΥΣΑΙ · CODEX RAVENNAS 429 (\R) · FOLIO 18/30}
\renewcommand{\currHeadRO}{f. 171r (Classense 173r) · vv. 721–764 · Collation diplomatique}
\noindent\begin{adjustbox}{max totalheight=30.4cm, max width=\textwidth, keepaspectratio}\begin{minipage}{\textwidth}
\bandeauRav{f. 171r}{Classense 173r}{721–764}{Folio 18/30 · Recto (page de droite) · 15 scholies \& gloses ({\policeGr Σ} 234–248)}
{\fontsize{8.8pt}{10.2pt}\selectfont
\ligneRav{721}{χορ/}{\hspace{1.1em}\difference{ξὺν} ἀθέοις ἔργοις}{\textbf{721} ἀθέοις ἔργοις γὰρ R : ἀθ. γὰρ ἔργ. Burges : ἐπ' ἀθ. ἔρ. γὰρ Enger : \lecon{ξὺν ἀθέοις ἔργοις} Bibi ap. Taillardat cf. v.716}
\ligneRav{722}{}{\difference{καὶ γὰρ} ἀνταμειψόμεσθά σ᾽ ὥσπερ εἰκὸς ἀντὶ τῶνδε·}{\textbf{722} γὰρ R : \lecon{καὶ γὰρ} Hermann}
\ligneRav{724}{}{\hspace{1.1em}\difference{τάχα δὲ} μεταβαλοῦσ᾽ ἐπὶ κακὸν ἑτερότρο-}{\textbf{724} τάχα δέ σε R : \lecon{τάχα δὲ} Fritzsche}
\ligneRav{725}{}{\hspace{1.1em} πον ἐπέχει \difference{τύχη}·}{\textbf{725} τις τύχη R : \lecon{τύχη} Bergk}
\ligneRav{726}{χορ/}{ἀλλὰ τάσδε μὲν λαβεῖν \expo{χ}ρῆν σ᾽ ἐκφέρειν τε τῶν ξύλων,}{\textbf{726} \lecon{\expo{χ}ρῆν σ᾽} \Rdeux : χρὴ μ᾽ \Mon}
\ligneRav{}{}{καὶ καταίθειν τὸν πανοῦργον πυρπολεῖν θ᾽ ὅσον τάχος·}{}
\ligneRav{}{\auDessus{γυ}{ν}}{ἴωμεν ἐπὶ τὰς κληματίδας ὦ Μανία·}{}
\ligneRav{729}{}{κἀγώ σ᾽ ἀποδείξω θυμάλωπα τήμερον·}{\textbf{729} \lecon{κἀγώ} R : ἦ ᾽γώ Sommerstein}
\ligneRav{730}{μ\auDessus{ν}{Η}}{ὕφαπτε καὶ κάταιθε· \difference{σὺ δὲ τὸ} Κρητικὸν}{\textbf{730} σὺ τόδε τὸ R : \lecon{σὺ δὲ τὸ} Suid. κ 2408}
\ligneRav{}{}{ἀπόδυθι ταχέως· τοῦ θανάτου δ᾽ ὦ παιδίον}{}
\ligneRav{}{}{μόνην γυναικῶν αἰτιῶ τὴν μητέρα·}{}
\ligneRav{733}{}{τουτὶ τί ἔστιν; ἀσκὸς \difference{ἐγένεθ᾽} ἡ κόρη}{\textbf{733} ἐγένηθ᾽ R : \lecon{ἐγένεθ᾽} Gryaneaus}
\ligneRav{}{}{οἴνου πλέως καὶ ταῦτα Περσικὰς ἔχων·}{}
\ligneRav{735}{}{ὦ θερμόταται γυναῖκες, ὦ ποτίσταται}{\textbf{735} \lecon{θερμόταται} R, Suid. θ 251}
\ligneRav{}{}{κἀκ παντὸς ὑμεῖς μηχανώμεναι πιεῖν,}{}
\ligneRav{}{}{ὦ μέγα καπήλοις ἀγαθὸν ἡμῖν δ᾽ αὖ κακόν,}{}
\ligneRav{}{}{κακὸν δὲ καὶ τοῖς σκευαρίοις καὶ τῇ κρόκῃ·}{}
\ligneRav{}{\auDessus{γυ}{ν}}{παράβαλλε πολλὰς κληματίδας ὦ Μανία·}{}
\ligneRav{740}{μ\auDessus{ν}{Η}}{παράβαλλε δῆτα· σὺ δ᾽ \difference{ἀπόκριναί} μοι τοδί,}{\textbf{740} ἀπόκρινέ R : \lecon{ἀπόκριναί} Μυ2}
\ligneRav{741}{—}{τουτὶ τεκεῖν φῄς; : \difference{καὶ δέκα μῆνας} αὔτ᾽ ἐγὼ}{\textbf{741} δέκα μῆνας R : \lecon{καὶ δέκα μῆνας} Biset : δέκα μῆνας γε Scaliger}
\ligneRav{}{—}{ἤνεγκον. : ἤνεγκας σύ; : νὴ τὴν Ἄρτεμιν·}{}
\ligneRav{743}{μ\auDessus{ν}{Η}}{τρικότυλον; ἢ πῶς. : εἰπέ μοι· τί μ᾽ \difference{ἠργάσω};}{\textbf{743} \_τρικο\douteux{τ}\douteux{υ}\douteux{λ}\douteux{ο}ν·η\douteux{π}\douteux{ω}\douteux{c} \  ειπ\douteux{ε}\douteux{μ}οι\douteux{τ}[ \PsiDeux : τρικότυλον ἢ πῶς; εἰπέ μοι --- τί μ᾽ εἰργάσω; R : \lecon{ἠργάσω} Hall-Geldart}
\ligneRav{}{}{ἀπέδυσας ὦναίσχυντέ μου τὸ παιδίον}{}
\ligneRav{745}{—}{\difference{τυννοῦτον} ὄν. : \difference{τυννοῦτο}; μικρὸν νὴ Δία·}{\textbf{745} \lecon{τυννοῦτον} et \lecon{τυννοῦτο} \PsiDeux, Brunck : τυνοῦτον et τυνοῦτο R \espApp Γϋ. ante μικρὸν R : \lecon{nulla nota mutat. pers. } \PsiDeux}
\ligneRav{746}{}{πόσ᾽ ἔτη δὲ \difference{γέγονε}; τρεῖς χοᾶς ἢ τέτταρας;}{\textbf{746} \lecon{γέγονε} \PsiDeux, Brunck : γέγονεν R}
\ligneRav{747}{\auDessus{γυ}{ν}}{σχεδὸν τοσοῦτον \difference{χὤσον} ἐκ Διονυσίων·}{\textbf{747} καὶ ὅσον \PsiDeux, R : ὅσον Suid. δ 1168 : \lecon{χὤσον} Bentley}
\ligneRav{}{—}{ἀλλ᾽ ἀπόδος αὐτό. : μὰ τὸν Ἀπόλλω τουτονί·}{}
\ligneRav{749}{\auDessus{γυ}{ν}}{ἐμπρήσομεν τοίνυν σε. : πάνυ γ᾽ \difference{ἐμπίμπρατε}·}{\textbf{749} ἐμπιμπρᾶται R : \lecon{ἐμπιμπράτε} (sic) \Mon}
\ligneRav{750}{}{αὕτη δ᾽ ἀποσφαγήσεται μάλ᾽ αὐτίκα·}{}
\ligneRav{}{\auDessus{γυ}{ν}}{μὴ δῆθ᾽, ἱκετεύω σ᾽· ἀλλ᾽ ἔμ᾽ ὅ τι χρῄζεις ποίει}{}
\ligneRav{}{—}{ὑπέρ γε τούτου. : φιλότεκνός τις εἶ φύσει·}{}
\ligneRav{}{}{ἀλλ᾽ οὐδὲν ἧττον ἥδ᾽ ἀποσφαγήσεται}{}
\ligneRav{754}{\auDessus{γυ}{ν}}{οἴμοι τέκνον· \difference{δός τὸ} \difference{σφαγεῖον} (σφαγ\inRas{ί}{εῖ}ον \Rpc) Μανία,}{\textbf{754} \lecon{δός τὸ} \PsiDeux, Lobeck (ad Soph. \italique{Ai.} 1066) : δός μοι τὸ R \espApp σφαγίον \Rac : \lecon{σφαγεῖον} \Rpc, \Mon}
\ligneRav{755}{}{ἵν᾽ οὖν τό γ᾽ αἷμα τοῦ τέκνου τοὐμοῦ λάβω·}{}
\ligneRav{}{μ\auDessus{ν}{Η}}{ὕπεχ᾽ αὐτό, χαριοῦμαι γὰρ ἕν γε τοῦτό σοι·}{}
\ligneRav{}{\auDessus{γυ}{ν}}{κακῶς ἀπόλοι᾽, ὡς φθονερὸς εἶ καὶ δυσμενής·}{}
\ligneRav{}{μ\auDessus{ν}{Η}}{τουτὶ τὸ δέρμα τῆς ἱερείας γίγνεται·}{}
\ligneRav{758}{\auDessus{γυ}{ν}}{τί τῆς \difference{ἱερείας} γίγνεται; : τουτί. λαβέ·}{\textbf{758-9} \lecon{ἱερείας} \PsiDeux, R : ἱερέας Austin \& Olson}
\ligneRav{760}{ἄλλη γυνή}{ταλαντάτη Μίκα τίς \difference{ἐξεκόρησέ} σε;}{\textbf{760} ἐξεκόρησεν R : \lecon{ἐξεκόρησε} \Mon}
\ligneRav{761}{}{τίς τὴν ἀγαπητὴν παῖδά \difference{σοὐξηράσατο};}{\textbf{761} [\douteux{ }\douteux{ }\douteux{ }\douteux{ }\douteux{ }\douteux{ }\douteux{ }\douteux{ }]απη\douteux{τ}η\douteux{ν}\douteux{π}\douteux{ }ι\douteux{ }α\douteux{ }\douteux{ }υξ[ \PsiDeux \espApp σοὐξῃρήσατο R : \lecon{σοὐξηράσατο} Fritzsche}
\ligneRav{}{\auDessus{γυ}{ν}}{ὁ πανοῦργος οὗτος. ἀλλ᾽ ἐπειδήπερ πάρει,}{}
\ligneRav{}{}{φύλαξον αὐτόν, ἵνα λαβοῦσα Κλεισθένη}{}
\ligneRav{}{}{τοῖσιν πρυτάνεσιν ἃ πεποίηχ᾽ οὗτος φράσω·}{}
}
\vspace{3pt}{\color{ruleGold}\hrule height 0.6pt}\vspace{3pt}
{\fontsize{8.0pt}{9.4pt}\selectfont
\noindent{\bfseries\scshape\color{rubricRed}Loci paralleli et testimonia (f. 171r, vv. 721–764) :} \textbf{729} (θυμάλωπα) Hesychius θ 862 θυμάλωψ ; cf. Photius θ 256 Θυμάλωψ, Hesychius θ 861 \espTest \textbf{730-1} (σὺ δὲ -- ταχέως) Suidas κ 2408 \espTest \textbf{731} (Κρητικὸν) Hesychius κ 4087 Κρητικόν· ἱματίδιον λεπτὸν καὶ βραχύ \espTest \textbf{734} (Περσικὰς) \Eustathe 1522.12 \espTest \textbf{735-7} Suidas θ 251 (ὦ et αὖ omissis, καλόν et δὲ pro ἀγαθόν et δ᾽) \espTest \textbf{735} \Eustathe 1441.23 ; \EtymSym α 241 (ed. Lasserre-Livadaras) \espTest (ποτίσταται) \EtymMagnum 31.13 \espTest \textbf{742} (ἤνεγκον) \Antiatt 98.11 ; \Phryn 73.1 \espTest \textbf{745} (τυννοῦτο) Photius τ 552 τυννοῦτο {\policeFr ∼} Hesychius τ 1645 \espTest \textbf{747} Suidas δ 1168 (καὶ omisso) \espTest \textbf{749} (ἐμπίμπρατε) Hesychius ε 2458 ἐμπίπραται {\policeFr ∼} Suidas ε 1018 {\policeFr ∼} Zonaras 709 \espTest \textbf{754} Photius σ 850 σφαγεῖον· [...] περισπᾶται ἡ δευτέρα συλλαβή\par\vspace{2.5pt}
{\color{ruleGold}\hrule height 0.35pt}\vspace{2.5pt}
\noindent{\bfseries\scshape\color{rubricRed}Suite (en commentaire) de l'apparat critique et paléographique :} \textbf{Contexte :} Suite et dénouement de la parodie du \italique{Télèphe} : les femmes apportent du bois pour brûler Mnésiloque ; en démaillotant l'otage, Mnésiloque découvre que le « bébé » est une outre pleine de vin chaussée de petites bottines perses (\italique{persikai}), qu'il « égorge » et boit goulûment au-dessus du bassin sacrificiel (\italique{sphageion}, vv. 721–764). --- \textbf{Témoin papyrologique :} \textbf{\PsiDeux{} (\italique{P.Oxy.} LVI 3839, fr. 1)} couvre les \textbf{vv. 742–766} et apporte quatre confirmations décisives face à R (vv. 743, 745, 746, 754). --- \textbf{Confrontation directe entre R et le papyrus \PsiDeux{} (vv. 743–754) :} 1° Au \textbf{v.~743}, \PsiDeux{} (\_τρικοτυλον· ηπωc ειπε μοι τ[) marque le changement d'interlocuteur avant εἰπέ μοι par un point haut et un espace blanc, tandis que \R{} plaçait le \italique{dikolon} après εἰπέ μοι ; 2° au \textbf{v.~745}, \PsiDeux{} donne \lecon{τυννοῦτον ... τυννοῦτο} (géminée expressive contre τυνοῦτον ... τυνοῦτο de \R{}) et n'a aucun signe de changement de personnage devant μικρὸν, prouvant que le sigle Γϋ. de \R{} est fautif ; 3° au \textbf{v.~754}, \PsiDeux{} confirme \lecon{δός τὸ σφαγεῖον} (contre δός μοι τὸ de \R{}), tandis que sur \R{} lui-même \Rpc{} corrige l'iotacisme σφαγίον (\Rac{}) en \lecon{σφαγεῖον}. --- \textbf{v. 721 :} ἀθέοις ἔργοις γὰρ R (ἀθ. γὰρ ἔργ. Burges ; ἐπ' ἀθ. ἔρ. γὰρ Enger) : \lecon{ξὺν ἀθέοις ἔργοις γὰρ} Bibi ap. Taillardat (cf. v. 716) \espApp γὰρ R : \lecon{καὶ γὰρ} Hermann.\par}
\end{minipage}\end{adjustbox}
\newpage
\renewcommand{\currHeadLE}{SCHOLIES MARGINALES ({\policeGr Σ}\expo{R}) \& GLOSES INTERLINÉAIRES · FOLIO 18/30}
\renewcommand{\currHeadRE}{f. 171r (Classense 173r) · vv. 721–764 · {\policeGr Σ} 234 à {\policeGr Σ} 248}
\renewcommand{\currHeadLO}{SCHOLIES MARGINALES ({\policeGr Σ}\expo{R}) \& GLOSES INTERLINÉAIRES · FOLIO 18/30}
\renewcommand{\currHeadRO}{f. 171r (Classense 173r) · vv. 721–764 · {\policeGr Σ} 234 à {\policeGr Σ} 248}
\bandeauScholies{Folio 18/30 — f. 171r}{Classense 173r}{vv. 721–764}{Recto (page de droite) · 15 entrées ({\policeGr Σ} 234 à {\policeGr Σ} 248)}
{\fontsize{8.8pt}{10.5pt}\selectfont\noindent\textbf{\color{rubricRed}Disposition topographique des marges sur le f. 171r :} Marge extérieure à droite ; marge intérieure à gauche. — \textit{Fac-similé de référence :} \texttt{folios\_individuels\_ff162v-177r/18\_f171r\_classense173r\_vv0721-0764.jpg}.\par\vspace{3pt}}
\entreeScholie{234}{726}{ἀλλὰ τάσδε μὲν λαβεῖν χρῆν σ᾽ ἐκφέρειν τε τῶν ξύλων,}{Interlinéaire (supra lineam, f. 171r, v. 726)}{πρὸς τὰς γυναῖκας.}{πρ(ὸς) τ\expo{ὰς} γυναῖκ\expo{ας} {\policeFr ÷}}{« S'adressant aux femmes du Chœur. »}
\commScholie{Précise le destinataire de l'ordre donné par Mika.}
\entreeScholie{235}{728}{ἴωμεν ἐπὶ τὰς κληματίδας ὦ Μανία·}{Marge droite (extérieure, f. 171r, en face du v. 728)}{ὄνομα κύριον δούλης τὸ Μανία. καὶ ἐκτεταμένον ἔχει τὸ α.}{ὄνομα κύρι\expo{ον} δούλ\expo{ης} τ(ὸ) Μανία · \textbf{ς} ἐκτεταμέν\expo{ον} ἔχ(ει) τ(ὸ) α {\policeFr ÷}}{« Μανία (Mania) est un nom propre d'esclave (d'origine phrygienne) ; et il a l'alpha final long (ἐκτεταμένον ἔχει τὸ α, au vocatif ὦ Μανία). »}
\commScholie{Remarque prosodique et onomastique très fine d'Hérodien/Symmaque : contrairement aux noms abstraits où -ία s'abrège parfois, le nom servile phrygien Μανία garde son -α final long.}
\entreeScholie{236}{729}{κἀγώ σ᾽ ἀποδείξω θυμάλωπα τήμερον·}{Interlinéaire (supra lineam au-dessus de θυμάλωπα, f. 171r, v. 729)}{σπινθῆρα.}{σπινθῆρα {\policeFr ÷}}{« Un tison ardent, un charbon embrasé (σπινθῆρα). »}
\commScholie{Glose lexicale de R\expo{2} expliquant le substantif expressif θυμάλωπα (acc. de θυμάλωψ, « tison fumant, charbon ardent », repris au v. 729 dans la menace κἀγώ σ᾽ ἀποδείξω θυμάλωπα) par σπινθῆρα (« tison étincelant »).}
\entreeScholie{237}{730}{σὺ δὲ τὸ Κρητικὸν / ἀπόδυθι ταχέως·}{Marge droite (extérieure, f. 171r, en face du v. 730)}{εἶδος ἱματίου. ἠμφίεσται δὲ ὡς παιδίον καὶ ὑποδήματα ἔχει ὁ ἀσκός.}{εἶδ\expo{ος} ἱματί\expo{ου} · ἠμφίεστ(αι) δ(ὲ) ὡς παιδί\expo{ον} \textbf{ς} ὑποδήματ(α) ἔχ(ει) ὁ ἀσκός {\policeFr ÷}}{« τὸ Κρητικόν : une espèce de petite tunique légère (mantelet crétois). L'outre à vin est habillée comme un poupon et porte même de petites chaussures. »}
\commScholie{Repris par la Souda κ 2398 (cf. Pollux VII, 77). Confirme le travestissement complet de l'outre.}
\entreeScholie{238}{734}{τουτὶ τί ἔστιν; ἀσκὸς ἐγένεθ᾽ ἡ κόρη / οἴνου πλέως καὶ ταῦτα Περσικὰς ἔχων·}{Marge droite (extérieure, f. 171r, en face du v. 734)}{ὁ ἀσκός. εἶδος δὲ ὑποδημάτων αἱ Περσικαί.}{ὁ ἀσκός · εἶδ\expo{ος} δ(ὲ) ὑποδημάτ\expo{ων} αἱ Περσικαί {\policeFr ÷}}{« C'est l'outre qui les porte ! Et les Περσικαί (“persiques”) sont une espèce de chaussons souples de femme. »}
\commScholie{Repris par la Souda π 1372 : chaussons blancs de femme que Mnésiloque retire à l'outre avant de l'égorger (détail visible sur le cratère de Würzburg !).}
\entreeScholie{239}{735}{ὦ θερμόταται γυναῖκες, ὦ ποτίσταται}{Interlinéaire (supra lineam au-dessus de ποτίσταται, f. 171r, v. 735)}{ὦ μέθυσοι.}{ὦ μέθυσοι {\policeFr ÷}}{« Ô buveuses invétérées, ô ivrognesses ! »}
\commScholie{Cas de méprise lexicographique brillamment démasqué par Louis Hauseux (vv 642-fin-bis.tex, ll. 1260-1277) : sur le Ravennas, la glose ὦ μέθυσοι explique évidemment le superlatif comique ὦ ποτίσταται (« ô très biberonnantes ») ; mais le compilateur de la Souda (θ 251) a rattaché par erreur cette glose au premier adjectif du vers, créant de toutes pièces l'entrée fantôme Θερμόταται· μέθυσοι !}
\entreeScholie{240}{738}{ὦ μέγα καπήλοις ἀγαθὸν, ἡμῖν δ᾽ αὖ κακόν, / κακὸν δὲ καὶ τοῖς σκευαρίοις καὶ τῇ κρόκῃ.}{Marge droite (extérieure, f. 171r, en face des vv. 737–738)}{ἅπαντα γὰρ ἕνεκα τοῦ πιεῖν ἐνέχυρα τίθεται καὶ πιπράσκεται.}{ἅπαντ(α) γ(ὰρ) ἕνεκα τ\expo{οῦ} πι(εῖν) ἐνέχυρα τίθετ(αι) \textbf{ς} πιπράσκετ(αι) {\policeFr ÷}}{« (“Malheur pour la vaisselle et pour la trame du métier à tisser”) : car les femmes mettent en gage et vendent tous les ustensiles du ménage et la laine pour s'acheter à boire chez le cabaretier. »}
\commScholie{Cité et complété par Louis Hauseux (vv 642-fin-bis.tex, l. 1352 n. 2).}
\entreeScholie{241}{742}{τουτὶ τεκεῖν φῄς; — καὶ δέκα μῆνας αὔτ᾽ ἐγὼ / ἤνεγκον.}{Marge droite (extérieure, f. 171r, en face des vv. 741–742)}{ἀλλ’ οὐ δέκα μῆνας κυοῦσιν αἱ γυναῖκες, ἀλλ’ ἐννέα. εἰώθασι δὲ οὕτω πλήρει ἀριθμῷ χρῆσθαι ἀντὶ τῶν ἐννέα μηνῶν.}{ἀλλ' οὐ δέκα μῆν\expo{ας} κυοῦσ(ιν) αἱ γυναῖκ\expo{ες}, ἀλλ' ἐννέα · εἰώθασ(ι) δ(ὲ) οὕτω πλήρ(ει) ἀριθμῷ χρῆσθ(αι) ἀν\auDessus{τ}{ι} τ\expo{ῶν} ἐννέα μην\expo{ῶν} {\policeFr ÷}}{« Mais les femmes ne portent pas leur enfant dix mois, mais neuf ! Toutefois, les anciens ont coutume d'employer ce chiffre rond (dix mois lunaires de 29 jours et demi) au lieu des neuf mois solaires. »}
\commScholie{Commenté par Louis Hauseux (vv 642-fin-bis.tex, l. 1378) avec le parallèle d'Aulu-Gelle (Nuits attiques III, 16) et de Ménandre (Plokion).}
\entreeScholie{242}{743}{τρικότυλον;}{Interlinéaire (supra lineam au-dessus de τρικότυλον, f. 171r, v. 743)}{τρεῖς κοτύλας οἴνου χωροῦν.}{τρεῖς κοτύλ\expo{ας} οἴν\expo{ου} χωροῦν {\policeFr ÷}}{« Contenant trois cotyles de vin (environ 0,8 litre). »}
\commScholie{Glose métrologique sur le double sens de τρικότυλον.}
\entreeScholie{243}{746}{πόσ᾽ ἔτη δὲ γέγονε; τρεῖς Χοᾶς ἢ τέτταρας;}{Marge droite (extérieure, f. 171r, en face du v. 746)}{τοῦτο πάλιν ὡς ἐπὶ ἀσκοῦ· ἀντὶ γὰρ τοῦ εἰπεῖν ἔτη εἶπε χοᾶς.}{τοῦτ(ο) πάλ(ιν) ὡς ἐπὶ ἀσκ\expo{οῦ} · ἀν\auDessus{τ}{ι} γ(ὰρ) τ\expo{οῦ} εἰπ(εῖν) ἔτ(η) εἶπ(ε) χοᾶς {\policeFr ÷}}{« Ceci encore est dit à double entente comme s'appliquant à une outre : au lieu de dire “trois ou quatre ans”, il a dit “trois ou quatre fêtes des Choës / trois ou quatre conges de vin (χοᾶς)” ! »}
\commScholie{Explique le calembour entre la fête printanière des Choës (deuxième jour des Anthestéries, où l'on couronnait les enfants de trois ans) et la mesure de capacité pour le vin (χοῦς).}
\entreeScholie{244}{747}{σχεδὸν τοσοῦτον χὤσον [καὶ ὅσον R] ἐκ Διονυσίων·}{Marge droite (extérieure, f. 171r, bloc de 7 lignes en face des vv. 747–750)}{εἰώθασι λέγειν, ἐάν τις ἐν Διονυσίοις ἢ Διασίοις ἢ Σκίροις μικρόν τι προσεναγισμάτων φησὶν ευπ\expo{α}. ἅμα δὲ ὅτι καὶ περὶ οἴνου ἐστὶν ὁ λόγος. Ἄλλως. οὕτω[...] λέγοντος ἃ δὲ ἔχει ἐκ Διονυσίων ἢ Ἀπατουρίων, τὸν ὑπερπίπτοντα τῶν μηνῶν ἀριθμὸν οὕτω δηλοῖ τις. τοῦτο ἔθος ἦν Ἀττικοῖς, λέγειν τὰ ἔτη κατὰ τὸν ὑπερπίπτοντα ἀριθμὸν ἀπὸ τῶν Διονυσίων.}{εἰώθασι λέγε(ιν), ἐάν τις ἐν Διονυσίοις ἢ Διασίοις ἢ Σκίροις μικρόν τι προσεναγισμάτων φησ(ὶν) ευπ\expo{α} · ἅμα δ(ὲ) ὅτι \textbf{ς} π(ερὶ) οἴν\expo{ου} ἐστ(ὶν) ὁ λόγος · ἄλλ\expo{ως} · οὕτω λέγοντ\expo{ος} ἃ δ(ὲ) ἔχ(ει) ἐκ Διονυσίων ἢ Ἀπατουρίων, τ(ὸν) ὑπερπίπτοντα τ\expo{ῶν} μην\expo{ῶν} ἀριθμ(ὸν) οὕτω δηλοῖ τις · τοῦτ(ο) ἔθ\expo{ος} ἦν Ἀττικ\expo{οῖς}, λέγε(ιν) τὰ ἔτη κατὰ τ(ὸν) ὑπερπίπτοντα ἀριθμ(ὸν) ἀπ(ὸ) τ\expo{ῶν} Διονυσίων {\policeFr ÷}}{« On a coutume de s'exprimer ainsi en prenant pour repère les Dionysies, les Diasies ou les Skira ; et en même temps, c'est parce qu'il est question de vin (placé sous le patronage de Dionysos). — Autre rédaction : En ajoutant “plus le temps écoulé depuis les dernières Dionysies (ou les Apatouries)”, on indique le nombre de mois excédentaires au-delà des années pleines. C'était l'usage des Attiques de compter les années en exprimant la fraction supplémentaire à partir des Grandes Dionysies. »}
\commScholie{Commenté par Louis Hauseux (vv 642-fin-bis.tex, l. 1429) : tout en feignant de donner l'âge exact du bébé en années et mois depuis les dernières Dionysies, Mika avoue que son outre est remplie du vin nouveau tiré aux Dionysies.}
\entreeScholie{245}{754}{οἴμοι τέκνον· δὸς τὸ [δός μοι τὸ R] σφαγεῖον, Μανία,}{Marge droite (extérieure, f. 171r, en face du v. 754)}{τοῦτο αἱματοδεκτικὸν ἀγγεῖον, ὃ εἶπεν ὁ ποιητὴς ἄμνιον.}{τοῦτ(ο) αἱματοδεκτικ(ὸν) ἀγγεῖ\expo{ον}, ὃ εἶπ(εν) ὁ ποιητ\expo{ὴς} ἄμνι\expo{ον} {\policeFr ÷}}{« Le σφαγεῖον est le bassin sacrificiel destiné à recueillir le sang de la victime égorgée, vase que le Poète (Homère, Odyssée III, 444) appelle ἄμνιον. »}
\commScholie{Repris par la Souda σ 1705. Comme le relève Louis Hauseux (vv 642-fin-bis.tex, ll. 1483-1494), ce bassin (σφαγεῖον, corrigé de σφαγίον sur R) est exactement celui que tend Mika sous le goulot de l'outre sur le cratère de Würzburg !}
\entreeScholie{246}{756}{ἵν᾽ οὖν τό γ᾽ αἷμα τοῦ τέκνου τοὐμοῦ λάβω.}{Marge droite (extérieure, f. 171r, en face du v. 756)}{τοῦτο ἅμα εἰπὼν ἐκχέει τὸν οἶνον ὁ κηδεστὴς ἐν τῷ βωμῷ.}{τοῦτ(ο) ἅμα εἰπ(ὼν) ἐκχέ(ει) τ(ὸν) οἶν\expo{ον} ὁ κηδεστ\expo{ὴς} ἐν τῷ βωμῷ {\policeFr ÷}}{« Tout en disant ces mots (au v. 757), le parent par alliance (Mnésiloque) verse (ou boit à la régalade) le vin de l'outre sur l'autel. »}
\commScholie{Discuté par Louis Hauseux (vv 642-fin-bis.tex, l. 1506 n. 1) : Mnésiloque perce ou débouche l'outre et en boit le vin en laissant à peine quelques gouttes pour le bassin de Mika.}
\entreeScholie{247}{758}{τουτὶ τὸ δέρμα τῆς ἱερείας γίγνεται.}{Marge droite (extérieure, f. 171r, en face du v. 758)}{ἀπὸ ἔθους. ἐκδερματώσαντες γὰρ τὰ ἱερεῖα διδόασι τὰ δέρματα ταῖς ἱερείαις. ἐκχέας οὖν τὸν οἶνον χαρίζεται αὐτῇ τὸν ἀσκόν.}{ἀπ(ὸ) ἔθ\expo{ους} · ἐκδερματώσαντ\expo{ες} γ(ὰρ) τὰ ἱερεῖα διδόασ(ι) τὰ δέρματ(α) τ\expo{αῖς} ἱερεί\expo{αις} · ἐκχέ\expo{ας} οὖν τ(ὸν) οἶν\expo{ον} χαρίζετ(αι) αὐτῇ τ(ὸν) ἀσκόν {\policeFr ÷}}{« D'après l'usage rituel (le δερματικόν) : car après avoir écorché les victimes, on remet les peaux en émolument aux prêtresses. Ayant donc vidé le vin, il fait cadeau à la prêtresse de la peau de l'outre vide ! »}
\commScholie{Repris par la Souda δ 262 : parodie irrésistible du droit sacerdotal sur la peau des victimes (τὸ δερματικόν).}
\entreeScholie{248}{760}{ταλαντάτη Μίκα, τίς ἐξεκόρησέ σε;}{Marge inférieure du f. 171r (en face du v. 760)}{ἐνταῦθα ἀπέδωκεν ὁ Ἀριστοφάνης τὸ ὄνομα τῆς γυναικὸς, ἧς ἥρπασε τὸ παιδίον ἤτοι τὸν ἀσκὸν ὁ κηδεστής.}{ἐνταῦθ(α) ἀπέδωκ(εν) ὁ Ἀριστοφάν(ης) τ(ὸ) ὄνομα τ\expo{ῆς} γυναικ\expo{ὸς}, ἧς ἥρπασ(ε) τ(ὸ) παιδί\expo{ον} ἤτοι τ(ὸν) ἀσκ(ὸν) ὁ κηδεστής {\policeFr ÷}}{« C'est ici seulement qu'Aristophane nous révèle le nom propre (Mika) de la femme à qui le parent par alliance avait arraché son enfant — c'est-à-dire son outre. »}
\commScholie{Observation dramaturgique très juste d'Aristarque/Didyme : selon un procédé familier à Aristophane (ainsi pour Dikaïopolis au v. 406 des Acharniens ou Strepsiade au v. 134 des Nuées), le nom de la Première Femme (Μίκα) n'est dévoilé qu'au v. 760.}
\clearpage
\renewcommand{\currHeadLE}{ΑΡΙΣΤΟΦΑΝΟΥΣ ΘΕΣΜΟΦΟΡΙΑΖΟΥΣΑΙ · CODEX RAVENNAS 429 (\R) · FOLIO 19/30}
\renewcommand{\currHeadRE}{f. 171v (Classense 173v) · vv. 765–799 · Collation diplomatique}
\renewcommand{\currHeadLO}{ΑΡΙΣΤΟΦΑΝΟΥΣ ΘΕΣΜΟΦΟΡΙΑΖΟΥΣΑΙ · CODEX RAVENNAS 429 (\R) · FOLIO 19/30}
\renewcommand{\currHeadRO}{f. 171v (Classense 173v) · vv. 765–799 · Collation diplomatique}
\noindent\begin{adjustbox}{max totalheight=30.4cm, max width=\textwidth, keepaspectratio}\begin{minipage}{\textwidth}
\bandeauRav{f. 171v}{Classense 173v}{765–799}{Folio 19/30 · Verso (page de gauche) · 8 scholies \& gloses ({\policeGr Σ} 249–256)}
{\fontsize{9.3pt}{10.8pt}\selectfont
\ligneRav{765}{μ\auDessus{ν}{Η}}{ἄγε δὴ τίς ἔσται μηχανὴ σωτηρίας;}{}
\ligneRav{}{}{τίς πεῖρα, τίς ἐπίνοι᾽; ὁ μὲν γὰρ αἴτιος}{}
\ligneRav{767}{}{κἄμ᾽ \difference{εἰσκυλίσας εἰς} τοιαυτὶ πράγματα}{\textbf{767} ἐσκυλίσας ἐς R : \lecon{εἰσκυλίσας εἰς} edd. pler.}
\ligneRav{768}{}{οὐ φαίνετ᾽, οὔπω. φέρε \difference{τίν᾽ οὖν ἂν} ἄγγελον}{\textbf{768} τίν᾽ οὖν R : τίν᾽ οὖν νῦν Biset : \lecon{τίν᾽ οὖν ἂν} Porson (\italique{Aristophanica} p.218)}
\ligneRav{}{}{πέμψαιμ᾽ ἐπ᾽ αὐτόν; οἶδ᾽ ἐγὼ καὶ δὴ πόρον}{}
\ligneRav{770}{}{ἐκ τοῦ Παλαμήδους· ὡς ἐκεῖνος, τὰς πλάτας}{}
\ligneRav{771}{}{ῥίψω γράφων. ἀλλ᾽ οὐ πάρεισιν αἱ πλάται·}{\textbf{771} \lecon{οὐ πάρεισιν} R : οὐ γάρ εἰσιν Fritzsche}
\ligneRav{772}{}{πόθεν οὖν γένοιτ᾽ ἄν μοι πλάται. πόθεν; πόθεν;}{\textbf{772} \lecon{γένοιτ᾽} R : γένοιντ᾽ Grynaeus \espApp πόθεν; R : \lecon {πόθεν; πόθεν;} Scaliger}
\ligneRav{773}{}{τί δ᾽ ἂν \difference{εἰ ταδὶ τἀγάλματ᾽} ἀντὶ τῶν πλατῶν}{\textbf{773} εἶτα διὰ τἀγάλματ᾽ R : ταδὶ τἀγάλματ᾽ Suid. : \lecon{εἰ ταδὶ τἀγάλματ᾽} Ellebodius}
\ligneRav{}{}{γράφων διαρρίπτοιμι; βέλτιον πολύ·}{}
\ligneRav{775}{}{ξύλον γέ τοι καὶ ταῦτα κἀκεῖν᾽ ἦν ξύλον·}{}
\ligneRav{}{}{\hspace{1.1em}ὦ χεῖρες ἐμαὶ}{}
\ligneRav{777}{}{\hspace{1.1em}ἐγχειρεῖν \difference{χρῆν} ἔργῳ πορίμῳ·}{\textbf{777} χρή R, Suid. π 2064 : \lecon{χρῆν} Bentley}
\ligneRav{}{}{\hspace{1.1em}ἄγε δὴ πινάκων ξεστῶν δέλτοι,}{}
\ligneRav{}{}{\hspace{1.1em}δέξασθε σμίλης ὁλκοὺς}{}
\ligneRav{780}{}{\hspace{1.1em}κήρυκας ἐμῶν μόχθων· οἴμοι}{}
\ligneRav{781}{}{\hspace{1.1em}τουτὶ \difference{τορῶ} μοχθηρόν·}{\textbf{781} τὸ ῥῶ R~: \lecon{τορῶ} Bibi ap. Hesych. τ 1162}
\ligneRav{782}{}{\hspace{1.1em}`\difference{χώρει}', `χώρει'. ποίαν αὔλακα;}{\textbf{782} \lecon{χώρει} bis R : χωρεῖ bis Brunck}
\ligneRav{783}{}{\hspace{1.1em}βάσκετ᾽ ἐπείγετε πάσας \difference{καθ᾽ ὁδοὺς}}{\textbf{783} καθόδους R : \lecon{καθ᾽ ὁδοὺς} Biset}
\ligneRav{784}{}{\hspace{1.1em}κείνᾳ \difference{ταύτᾳ}· ταχέως χρή·}{\textbf{784} ταῦτα R : \lecon{ταύτᾳ} Grynaeus (\italique{hac} Divus)}
\ligneRav{785}{χορ/}{ἡμεῖς τοίνυν ἡμᾶς αὐτὰς εὖ λέξωμεν παραβᾶσαι,}{}
\ligneRav{}{}{καίτοι πᾶς τις τὸ γυναικεῖον φῦλον κακὰ πόλλ᾽ ἀγορεύει,}{}
\ligneRav{}{}{ὡς πᾶν ἐσμὲν κακὸν ἀνθρώποις κἀξ ἡμῶν ἐστιν ἅπαντα,}{}
\ligneRav{788}{}{ἔριδες νείκη \difference{στάσις ἀργαλέα} λύπη πόλεμος. φέρε δή νυν,}{\textbf{788} στάσεισ· ἀργαλέα· R : \lecon{στάσις ἀργαλέα} Scaliger}
\ligneRav{789}{}{\difference{εἰ κακόν} ἐσμεν, τί γαμεῖθ᾽ ἡμᾶς, εἴπερ ἀληθῶς κακόν ἐσμεν,}{\textbf{789} εἰ καὶ κακόν R : \lecon{εἰ κακόν} Zanetti}
\ligneRav{790}{}{κἀπαγορεύετε μήτ᾽ ἐξελθεῖν μήτ᾽ \difference{ἐκκύψασαν} ἁλῶναι,}{\textbf{790} ἐγκύψασαν R : \lecon{ἐκκύψασαν} Reiske}
\ligneRav{}{}{ἀλλ᾽ οὑτωσὶ πολλῇ σπουδῇ τὸ κακὸν βούλεσθε φυλάττειν;}{}
\ligneRav{792}{}{κἂν ἐξέλθῃ τὸ γύναιόν ποι, \difference{κᾆθ᾽ εὕρητ᾽} αὐτὸ θύρασιν,}{\textbf{792} κᾆθ᾽ εὕροιτ' R : \lecon{κᾆθ᾽ εὕρητ᾽} Porson}
\ligneRav{793}{}{μανίας \difference{μαίνεσθ᾽}, οὓς \difference{χρῆν} σπένδειν καὶ χαίρειν, εἴπερ ἀληθῶς}{\textbf{793} μαίνεθ᾽ R : \lecon{μαίνεσθ᾽} \Mon \espApp χρή R : \lecon{χρῆν} Brunck}
\ligneRav{794}{}{ἔνδοθεν \difference{ηὕρετε} φροῦδον τὸ κακὸν καὶ μὴ \difference{κατελαμβάνετ᾽} ἔνδον·}{\textbf{794} εὕρετε R : \lecon{ηὕρετε} Meineke \espApp καταλαμβάνετ᾽ R : \lecon{κατελαμβάνετ᾽} Brunck}
\ligneRav{795}{}{κἂν \difference{καταδάρθωμεν} ἐν ἀλλοτρίων παίζουσαι καὶ κοπιῶσαι,}{\textbf{795} καταδαρθῶμεν R : \lecon{καταδάρθωμεν} Dindorf}
\ligneRav{}{}{πᾶς τις τὸ κακὸν τοῦτο ζητεῖ περὶ τὰς κλίνας περινοστῶν·}{}
\ligneRav{797}{}{κἂν ἐκ θυρίδος παρακύπτωμεν, \difference{τὸ κακὸν ζητεῖτε θεᾶσθαι}·}{\textbf{797} ζητεῖ τὸ κακὸν τεθεᾶσθαι R : \lecon{τὸ κακὸν ζητεῖτε θεᾶσθαι} J. Kaye}
\ligneRav{798}{}{κἂν \difference{αἰσχυνθεῖσ᾽} (αἰσχυνθ\inRas{η}{εὶ}ς \Rpc) ἀναχωρήσῃ, πολὺ μᾶλλον πᾶς ἐπιθυμεῖ}{\textbf{798} αἰσχυνθεὶς (ex αἰσχυνθῆς) \Rpc : \lecon{αἰσχυνθεῖσ᾽} Zanetti}
\ligneRav{799}{}{αὖθις \difference{τὸ κακὸν παρακύψαν ἰδεῖν}. οὕτως ἡμεῖς ἐπιδήλως}{\textbf{799} παρακύψασαν ἰδεῖν τὸ κακὸν R : \lecon{τὸ κακὸν παρακύψαν ἰδεῖν} Dindorf}
}
\vspace{3pt}{\color{ruleGold}\hrule height 0.6pt}\vspace{3pt}
{\fontsize{8.0pt}{9.4pt}\selectfont
\noindent{\bfseries\scshape\color{rubricRed}Loci paralleli et testimonia (f. 171v, vv. 765–799) :} \textbf{769-75} Suidas π 45 (οἶδα δ' pro οἶδ᾽, πότερον pro πόρον et γ' ἔκειντο pro γένοιντ᾽ ἄν) \espTest \textbf{770-1} Suidas π 1702 \espTest \textbf{777} Suidas π 2064 \espTest \textbf{778-80} (πινάκων -- μόχθων) Pollux X 58 (ἄγε δὴ et δέλτοι omissis) \espTest \textbf{778-9} Suidas σ 742\par\vspace{2.5pt}
{\color{ruleGold}\hrule height 0.35pt}\vspace{2.5pt}
\noindent{\bfseries\scshape\color{rubricRed}Suite (en commentaire) de l'apparat critique et paléographique :} \textbf{Témoin papyrologique :} \textbf{\PsiDeux{} (\italique{P.Oxy.} LVI 3839, fr. 1)} s'achève au \textbf{v. 766}. --- \textbf{Conjecture éclatante de \textbf{Bibi} au v. 781 (τὸ ῥῶ R {\policeFr →} τορῶ \textbf{Bibi}) :} Là où tous les éditeurs suivaient la coupure de mots de R (τὸ ῥῶ, « ce rhô »), Louis Hauseux restitue sans modifier une seule lettre de l'onciale ΤΟΡΩ le verbe \lecon{τορῶ} (« je grave, je creuse au burin », d'après la glose d'Hésychius τ 1162 τορεῖ· τορνεύει, γλύφει, τιτρώσκει), répondant directement à la question ποίαν αὔλακα; (« Quel sillon ? — Je grave : \italique{Va, va !} »). --- \textbf{Correction \italique{in rasura} au v. 798 :} \Rac{} avait écrit αἰσχυνθῆς, corrigé sur grattage par \Rpc{} en \lecon{αἰσχυνθ\inRas{η}{εὶ}ς} (pour le participe féminin élidé αἰσχυνθεῖσ' restitué par Zanetti). --- \textbf{v. 772 :} \lecon{γένοιτ'} R (γένοιντ' Grynaeus) \espApp πόθεν; R (haplographie finale) : \lecon{πόθεν; <πόθεν;>} Scaliger. --- \textbf{v. 781 :} τὸ ῥῶ R : \lecon{τορῶ} \lecon{Bibi} ap. Hésychius τ 1162. --- \textbf{v. 797 :} ζητεῖ τὸ κακὸν τεθεᾶσθαι R (inversion et médivision) : \lecon{τὸ κακὸν ζητεῖτε θεᾶσθαι} J. Kaye.\par}
\end{minipage}\end{adjustbox}
\newpage
\renewcommand{\currHeadLE}{SCHOLIES MARGINALES ({\policeGr Σ}\expo{R}) \& GLOSES INTERLINÉAIRES · FOLIO 19/30}
\renewcommand{\currHeadRE}{f. 171v (Classense 173v) · vv. 765–799 · {\policeGr Σ} 249 à {\policeGr Σ} 256}
\renewcommand{\currHeadLO}{SCHOLIES MARGINALES ({\policeGr Σ}\expo{R}) \& GLOSES INTERLINÉAIRES · FOLIO 19/30}
\renewcommand{\currHeadRO}{f. 171v (Classense 173v) · vv. 765–799 · {\policeGr Σ} 249 à {\policeGr Σ} 256}
\bandeauScholies{Folio 19/30 — f. 171v}{Classense 173v}{vv. 765–799}{Verso (page de gauche) · 8 entrées ({\policeGr Σ} 249 à {\policeGr Σ} 256)}
{\fontsize{8.8pt}{10.5pt}\selectfont\noindent\textbf{\color{rubricRed}Disposition topographique des marges sur le f. 171v :} Marge extérieure à gauche ; marge intérieure à droite. — \textit{Fac-similé de référence :} \texttt{folios\_individuels\_ff162v-177r/19\_f171v\_classense173v\_vv0765-0799.jpg}.\par\vspace{3pt}}
\entreeScholie{249}{766}{τίς πεῖρα, τίς ἐπίνοι᾽; ὁ μὲν γὰρ αἴτιος / κἄμ᾽ εἰσκυλίσας εἰς τοιαυτὶ πράγματα}{Interlinéaire (supra lineam, f. 171v, v. 766)}{ὁ Εὐριπίδης.}{ὁ Εὐριπίδ(ης) {\policeFr ÷}}{« C'est Euripide (qui m'a embarqué dans cette affaire). »}
\commScholie{Glose interlinéaire de R\expo{2} identifiant la périphrase ὁ μὲν γὰρ αἴτιος (« car le véritable responsable, celui qui m'a roulé dans un tel pétrin », vv. 766–767) comme désignant Euripide (ὁ Εὐριπίδης).}
\entreeScholie{250}{771}{ἐκ τοῦ Παλαμήδους· ὡς ἐκεῖνος, τὰς πλάτας / ῥίψω γράφων.}{Marge gauche (extérieure, f. 171v, grande scholie de 8 lignes en haut du verso)}{ὁ γὰρ Εὐριπίδης ἐν τῷ Παλαμήδει ἐποίησε τὸν Οἴακα τὸν ἀδελφὸν Παλαμήδους ἐπιγράψαι εἰς τὰς ναῦς τὸν θάνατον αὐτοῦ, ἵνα φερόμεναι αὗται ἔλθωσιν εἰς τὸν Ναύπλιον τὸν πατέρα αὐτοῦ καὶ ἀπαγγείλωσι τὸν θάνατον αὐτοῦ. ὥσπερ Οἴαξ τῷ Ναυπλίῳ γράφει ἐν τῷ Παλαμήδει Εὐριπίδου. ὁ γὰρ Οἴαξ ἐγχαράττει πολλαῖς πλάταις τὰ περὶ τὸν Παλαμήδην καὶ ἀφίησιν εἰς θάλασσαν, ὥστε μιᾷ γέ τινι τὸν Ναύπλιον προσπεσεῖν.}{ὁ γ(ὰρ) Εὐριπίδης ἐν τῷ Παλαμήδει ἐποίησε τ(ὸν) Οἴακα τ(ὸν) ἀδελφ(ὸν) Παλαμήδους ἐπιγράψαι εἰς τ\expo{ὰς} ναῦς τ(ὸν) θάνατ(ον) αὐτ\expo{οῦ}, ἵνα φερόμεν(αι) αὗτ(αι) ἔλθωσ(ιν) εἰς τ(ὸν) Ναύπλι(ον) τ(ὸν) πατέρα αὐτ\expo{οῦ} \textbf{ς} ἀπαγγείλωσι τ(ὸν) θάνατ(ον) αὐτ\expo{οῦ} · ὥσπ(ερ) Οἴαξ τῷ Ναυπλίῳ γράφει ἐν τῷ Παλαμήδει Εὐριπίδου · ὁ γ(ὰρ) Οἴαξ ἐγχαράττ(ει) πολλ\expo{αῖς} πλάταις τὰ π(ερὶ) τ(ὸν) Παλαμήδην \textbf{ς} ἀφίησ(ιν) εἰς θάλασσ(αν), ὥστε μιᾷ γέ τινι τ(ὸν) Ναύπλι(ον) προσπεσ(εῖν) {\policeFr ÷}}{« En effet, Euripide, dans son Palamède (joué en 415 av. J.-C., fr. 588a Kannicht), avait représenté Oïax, le frère de Palamède, inscrivant sur des rames de navires la mort de son frère, afin que ces rames, emportées par les flots, parviennent jusqu'à leur père Nauplios (en Eubée) et lui annoncent son meurtre. — Autre rédaction : Tout comme Oïax écrit à Nauplios dans le Palamède d'Euripide : car Oïax grave sur de nombreuses pales d'avirons (πλάταις) le sort de Palamède et les lance à la mer, afin que Nauplios tombe au moins sur l'une d'entre elles. »}
\commScholie{Cité et commenté par Louis Hauseux (vv 642-fin-bis.tex, ll. 1607-1616) : témoignage capital sur l'intrigue du Palamède d'Euripide.}
\entreeScholie{251}{773}{τί δ᾽ ἂν, εἰ ταδὶ τἀγάλματ᾽ ἀντὶ τῶν πλατῶν / γράφων διαρρίπτοιμι;}{Marge gauche (extérieure, f. 171v, sous Σ 771)}{πᾶν ἐφ’ ᾧ τις ἀγάλλεται. πινάκια δὲ ἦσαν ἀνακείμενα. πίνακες γὰρ ἦσαν ἐν τῷ ἱερῷ, καὶ ἀπὸ τούτων λαβὼν γράφει, καὶ ῥίπτει αὐτὰς λέγων «ἀπέλθετε, σημάνατε Εὐριπίδῃ.»}{πᾶν ἐφ' ᾧ τις ἀγάλλετ(αι) · πινάκια δ(ὲ) ἦσ(αν) ἀνακείμενα · πίνακ\expo{ες} γ(ὰρ) ἦσ(αν) ἐν τῷ ἱερῷ, \textbf{ς} ἀπ(ὸ) τούτ\expo{ων} λαβ(ὼν) γράφ(ει), \textbf{ς} ῥίπτ(ει) αὐτ\expo{ὰς} λέγ\expo{ων} «ἀπέλθετ(ε), σημάνατ(ε) Εὐριπίδῃ» {\policeFr ÷}}{« ἄγαλμα désigne au sens large toute offrande dont on se fait gloire (ἐφ' ᾧ τις ἀγάλλεται). Ici, ce sont des tablettes votives de bois (πινάκια) consacrées dans le sanctuaire : il y avait en effet des tablettes votives suspendues dans le Thesmophorion ; Mnésiloque les décroche, y grave son message et les jette au loin en disant : “Partez, allez prévenir Euripide !” »}
\commScholie{Repris par la Souda α 133 et discuté par Alphonse Willems (22e note) et Louis Hauseux (vv 642-fin-bis.tex, ll. 1643-1695) : τἀγάλματα ne désigne pas des statues en ronde bosse, mais des tablettes votives de bois peint (πινάκια) suspendues aux parois du sanctuaire.}
\entreeScholie{252}{777}{ὦ χεῖρες ἐμαί, / ἐγχειρεῖν χρῆν [χρὴ R] ἔργῳ πορίμῳ·}{Interlinéaire (supra lineam au-dessus de πορίμῳ, f. 171v, v. 777)}{ἄνυσιν ἔχοντι καὶ ἐπίνοιαν.}{ἄνυσ(ιν) ἔχοντ(ι) \textbf{ς} ἐπίνοι(αν) {\policeFr ÷}}{« πορίμῳ (“ingénieux, efficace”) : qui procure une issue salutaire et une invention féconde. »}
\commScholie{Glose interlinéaire de R\expo{2} reprise par la Souda π 2070, expliquant l'adjectif verbal πορίμῳ (« fertile en ressources, ingénieux, efficace ») appliqué par Mnésiloque à son stratagème des tablettes votives.}
\entreeScholie{253}{779}{ἄγε δὴ πινάκων ξεστῶν δέλτοι, / δέξασθε σμίλης ὁλκοὺς}{Marge gauche (extérieure, f. 171v, en face des vv. 778–779)}{σμίλης λέγει ἀντὶ τοῦ γραφείου. ὁλκοὺς λέγει τὰ σημεῖα, ἐπειδὴ ἕλκομεν γράφοντες.}{σμίλ\expo{ης} λέγ(ει) ἀν\auDessus{τ}{ι} τ\expo{οῦ} γραφεί\expo{ου} · ὁλκ\expo{οὺς} λέγ(ει) τὰ σημεῖα, ἐπειδὴ ἕλκομ(εν) γράφοντ\expo{ες} {\policeFr ÷}}{« Il dit σμίλης (“du ciseau à bois / burin”) au lieu du stylet à écrire (γραφεῖον) ; et il appelle ὁλκούς (“sillons tracés”) les caractères gravés, parce que nous tirons des traits (ἕλκομεν) en écrivant. »}
\commScholie{Repris par la Souda ο 166 et σ 742.}
\entreeScholie{254}{781}{οἴμοι, / τουτὶ τορῶ [τὸ ῥῶ R, Σ] μοχθηρόν· / χώρει χώρει· ποίαν αὔλακα;}{Marge gauche (extérieure, f. 171v, en face du v. 781)}{ὡς Εὐριπίδην γράφων, ἐν ᾧ τὸ ρ. ἐν τῷ γράφειν τὸ τοῦ Εὐριπίδου ὄνομα τοῦτό φησιν.}{ὡς Εὐριπίδ(ην) γράφ\expo{ων}, ἐν ᾧ τ(ὸ) ρ · ἐν τῷ γράφ(ειν) τ(ὸ) τ\expo{οῦ} Εὐριπίδ\expo{ου} ὄνομα τοῦτό φησ(ιν) {\policeFr ÷}}{« Comme s'il était en train de graver le nom Εὐριπίδης, dans lequel se trouve la lettre rhô (ρ). C'est en traçant le nom d'Euripide qu'il dit : “Aïe ! ce rhô est mal venu !” »}
\commScholie{Longuement discuté et rectifié par Louis Hauseux (vv 642-fin-bis.tex, ll. 1669-1684) : le scholiaste croit que Mnésiloque grave Εὐριπίδης, mais le vers suivant (χώρει, χώρει, « Accours, accours ! ») montre que Mnésiloque est en train de se dicter à lui-même le mot ΧΩΡΕΙ au moment même où son burin dérape sur la boucle du Ρ dans le bois !}
\entreeScholie{255}{785}{ἡμεῖς τοίνυν ἡμᾶς αὐτὰς εὖ λέξωμεν παραβᾶσαι,}{Interlinéaire (supra lineam au-dessus de παραβᾶσαι, f. 171v, v. 785)}{τῇ παραβάσει χρησάμεναι.}{τῇ παραβάσ(ει) χρησάμεν(αι) {\policeFr ÷}}{« En nous avançant vers le public pour prononcer la Parabase (παράβασις). »}
\commScholie{Note technique marquant le début de la Parabase du Chœur (vv. 785-845).}
\entreeScholie{256}{790}{κἀπαγορεύετε μήτ᾽ ἐξελθεῖν μήτ᾽ ἐκκύψασαν ἁλῶναι,}{Interlinéaire (supra lineam au-dessus d'ἐκκύψασαν ἁλῶναι, f. 171v, v. 790)}{μήτε καταληφθῆναι παρακύπτουσαν.}{μήτ(ε) καταληφθῆν(αι) παρακύπτουσ(αν) {\policeFr ÷}}{« Et de ne pas se laisser surprendre en train de pencher la tête à la fenêtre (παρακύπτουσαν). »}
\commScholie{Glose explicative de R\expo{2} paraphrasant la défense faite aux femmes par leurs maris jaloux μήτ᾽ ἐκκύψασαν ἁλῶναι (« ni de se laisser surprendre en train de se pencher à la fenêtre », παρακύπτουσαν).}
\clearpage
\renewcommand{\currHeadLE}{ΑΡΙΣΤΟΦΑΝΟΥΣ ΘΕΣΜΟΦΟΡΙΑΖΟΥΣΑΙ · CODEX RAVENNAS 429 (\R) · FOLIO 20/30}
\renewcommand{\currHeadRE}{f. 172r (Classense 174r) · vv. 800–837 · Collation diplomatique}
\renewcommand{\currHeadLO}{ΑΡΙΣΤΟΦΑΝΟΥΣ ΘΕΣΜΟΦΟΡΙΑΖΟΥΣΑΙ · CODEX RAVENNAS 429 (\R) · FOLIO 20/30}
\renewcommand{\currHeadRO}{f. 172r (Classense 174r) · vv. 800–837 · Collation diplomatique}
\noindent\begin{adjustbox}{max totalheight=30.4cm, max width=\textwidth, keepaspectratio}\begin{minipage}{\textwidth}
\bandeauRav{f. 172r}{Classense 174r}{800–837}{Folio 20/30 · Recto (page de droite) · 11 scholies \& gloses ({\policeGr Σ} 257–267)}
{\fontsize{9.3pt}{10.8pt}\selectfont
\ligneRav{800}{χορ/}{ὑμῶν ἐσμεν πολὺ βελτίους, βάσανός τε πάρεστιν ἰδέσθαι·}{}
\ligneRav{}{}{βάσανον δῶμεν πότεροι χείρους. ἡμεῖς μὲν γάρ φαμεν ὑμᾶς,}{}
\ligneRav{802}{}{\difference{ὑμεῖς δ᾽ ἡμᾶς}. \difference{σκεψώμεθα} δὴ κἀντιτιθῶμεν πρὸς ἕκαστον,}{\textbf{802} ὑμεῖς δ᾽ ὑμᾶς R : \lecon{ὑμεῖς δ᾽ ἡμᾶς} \Mon \espApp σκεψώμεσθα R : \lecon{σκεψώμεθα} \Mon}
\ligneRav{803}{}{παραβάλλουσαι τῆς τε γυναικὸς καὶ τἀνδρὸς τοὔνομ᾽ \difference{ἑκάστου}·}{\textbf{803} ἕκαστον R : \lecon{ἑκάστου} Grynaeus}
\ligneRav{804}{}{Ναυσιμάχης \difference{μέν γ᾽} ἥττων \difference{ἐστὶν} Χαρμῖνος· \difference{δῆλα δὲ} τἄργα·}{\textbf{804} μέν R : \lecon{μέν γ᾽} Dobree \espApp ἐστὶ R : \lecon{ἐστὶν} Küster \espApp δηλαδῆ R : \lecon{δῆλα δὲ} Grynaeus}
\ligneRav{805}{}{καὶ μὲν δὴ καὶ Κλεοφῶν \difference{χείρων} πάντως δήπου Σαλαβακχοῦς·}{\textbf{805} χεῖρον R : \lecon{χείρων} Zanetti}
\ligneRav{}{}{πρὸς Ἀριστομάχην δὲ χρόνου πολλοῦ, πρὸς ἐκείνην τὴν Μαραθῶνι,}{}
\ligneRav{807}{}{καὶ \difference{Στρατονίκην} ὑμῶν οὐδεὶς οὐδ᾽ ἐγχειρεῖ πολεμίζειν·}{\textbf{807} \lecon{Στρατονίκην} \PsiUn : στρατωνίκην R}
\ligneRav{808}{}{ἀλλ᾽ Εὐβούλης τῶν πέρυσίν \difference{τις βουλευτής} ἐστιν ἀμείνων}{\textbf{808} τίς βουλευτής R : \lecon{τις βουλευτής} Küster}
\ligneRav{809}{}{παραδοὺς ἑτέρῳ τὴν βουλείαν; οὐδ᾽ αὐτὸς τοῦτό γε \difference{φήσει}·}{\textbf{809} \lecon{αὐτὸς} R : Ἄνυτος Maas \espApp φήσεις R : \lecon{φήσει} Küster}
\ligneRav{810}{}{οὕτως ἡμεῖς πολὺ βελτίους τῶν ἀνδρῶν \difference{εὐχόμεθ᾽} εἶναι·}{\textbf{810} εὐχόμεσθ᾽ R : \lecon{εὐχόμεθ᾽} \Mon}
\ligneRav{811}{}{οὐδ᾽ ἂν κλέψασα γυνὴ \difference{ζεύγει} κατὰ πεντήκοντα τάλαντα}{\textbf{811} \lecon{ζεύγει post γυνὴ} R : post οὐδ᾽ ἂν van Herwerden}
\ligneRav{}{}{εἰς πόλιν ἔλθοι τῶν δημοσίων· ἀλλ᾽ ἢν τὰ μέγισθ᾽ ὑφέληται}{}
\ligneRav{813}{}{φορμὸν πυρῶν τἀνδρὸς κλέψασ᾽, αὐθημερὸν αὔτ᾽ ἀπέδωκεν}{\textbf{813} \lecon{αὔτ᾽ ἀπέδωκεν} R, M\expo{pc} : ἀνταπέδωκεν M\expo{ac}, Bentley}
\ligneRav{}{—}{\hspace{1.1em}ἀλλ᾽ ἡμεῖς ἂν πολλοὺς τούτων}{}
\ligneRav{815}{}{\hspace{1.1em}\difference{ἀποδείξαιμεν} ταῦτα \difference{ποιοῦντας}·}{\textbf{815} ἀποδείξομεν R : \lecon{ἀποδείξαιμεν} Bentley \espApp ποθοῦντας R : \lecon{ποιοῦντας} Zanetti (\italique{facientes} Divus)}
\ligneRav{}{}{\hspace{1.1em}καὶ πρὸς τούτοις γάστριδας ἡμῶν}{}
\ligneRav{}{}{\hspace{1.1em}ὄντας μᾶλλον καὶ λωποδύτας}{}
\ligneRav{}{}{\hspace{1.1em}καὶ βωμολόχους κἀνδραποδιστάς·}{}
\ligneRav{819}{}{\hspace{1.1em}καὶ μὲν δήπου καὶ τὰ πατρῷά γε}{\textbf{819} δήπου R : δή που Denniston}
\ligneRav{820}{}{\hspace{1.1em}χείρους ἡμῶν εἰσιν σῴζειν·}{}
\ligneRav{}{}{\hspace{1.1em}ἡμῖν μὲν γὰρ σῶν ἔτι καὶ νῦν}{}
\ligneRav{822}{}{\hspace{1.1em}\difference{ταὐτὸ κανοῦν καὶ τὸ σκιάδειον}·}{\textbf{822-3} τἀντίον ὁ κανὼν οἱ καλαθίσκοι / τὸ σκιάδειον· R : \lecon{ταὐτὸ κανοῦν καὶ τὸ σκιάδειον·} Bibi}
\ligneRav{824}{}{\hspace{1.1em}τοῖς δ᾽ ἡμετέροις \difference{ἀνδράσι} τούτοις}{\textbf{824} ἀνδράσιν R : \lecon{ἀνδράσι} \Mon}
\ligneRav{825}{}{\hspace{1.1em}\difference{ἀπόλωλεν} μὲν πολλοῖς ὁ κανὼν}{\textbf{825} ἀπόλωλε R : \lecon{ἀπόλωλεν} Daubuz, Bentley}
\ligneRav{}{}{\hspace{1.1em}ἐκ τῶν οἴκων αὐτῇ λόγχῃ,}{}
\ligneRav{}{}{\hspace{1.1em}πολλοῖς δ᾽ ἑτέροις ἀπὸ τῶν ὤμων}{}
\ligneRav{}{}{\hspace{1.1em}\{ ἐν ταῖς στρατιαῖς \}}{}
\ligneRav{}{}{\hspace{1.1em}ἔρριπται τὸ σκιάδειον·}{}
\ligneRav{830}{χορ/}{πόλλ᾽ ἂν αἱ γυναῖκες ἡμεῖς ἐν δίκῃ μεμψαίμεθ᾽ ἂν}{}
\ligneRav{}{}{τοῖσιν ἀνδράσιν δικαίως, ἓν δ᾽ ὑπερφυέστατον·}{}
\ligneRav{832}{}{\difference{χρῆν} γάρ, ἡμῶν εἰ τέκοι τις ἄνδρα χρηστὸν τῇ πόλει,}{\textbf{832} χρὴ R : \lecon{χρῆν} Brunck}
\ligneRav{}{}{ταξίαρχον ἢ στρατηγόν, λαμβάνειν τιμήν τινα,}{}
\ligneRav{834}{}{προεδρίαν τ᾽ αὐτῇ δίδοσθαι \difference{Στηνίοισι} καὶ Σκίροις}{\textbf{834} Τηνίοισι R : \lecon{Στηνίοισι} Σ, Suid. σ 1089}
\ligneRav{835}{}{ἔν τε ταῖς ἄλλαις ἑορταῖς αἷσιν ἡμεῖς ἤγομεν·}{}
\ligneRav{836}{}{\difference{εἰ δὲ} δειλὸν καὶ πονηρὸν ἄνδρα τις τέκοι γυνή,}{\textbf{836} ἤν δὲ R : \lecon{εἰ δὲ} Brunck}
\ligneRav{}{}{ἢ τριήραρχον πονηρὸν ἢ κυβερνήτην κακόν,}{}
}
\vspace{3pt}{\color{ruleGold}\hrule height 0.6pt}\vspace{3pt}
{\fontsize{8.0pt}{9.4pt}\selectfont
\noindent{\bfseries\scshape\color{rubricRed}Loci paralleli et testimonia (f. 172r, vv. 800–837) :} \textbf{800, 810} (βελτίους = βελτίονες) Moeris β 8 \espTest \textbf{805} (Σαλαβακχοῦς) Hésychius (σ 88) Σαλαβακχώ· πόρνης ὄνομα. ἀπὸ δὲ ταύτης καὶ τὰς καταφερεῖς εἰς τὰ ἀφροδίσια οὕτως ἔλεγον Ἀττικοί \espTest \textbf{808-9} (τις βουλευτής -- βουλείαν) Suidas β 431 (ἐστι pro ἐστιν et ἀμείνων omisso) \espTest \textbf{809} (παραδοὺς ἑτέρῳ τὴν βουλείαν) \Antiatt 84.30 \espTest (βουλείαν) \Harp B 19 ; β Photius β 233 βουλία ; Hesychius β 921 βουλεῖα \espTest \textbf{816} (γάστριδας) Photius γ 40 γάστρις = Suidas γ 72 ; cf. \AelDion γ 3 \espTest \textbf{822} %(τἀντίον ὁ κανὼν) Pollux VII 35 κανὼν ἱστοῦ τὸ καλούμενον ἀντίον \espTest \textbf{834} (Στηνίοισι) Suidas σ 1089 Στήνια καὶ Σκίρα (= Scholium) ; Hesychius σ 1825 στήνια\par\vspace{2.5pt}
{\color{ruleGold}\hrule height 0.35pt}\vspace{2.5pt}
\noindent{\bfseries\scshape\color{rubricRed}Suite (en commentaire) de l'apparat critique et paléographique :} \textbf{Témoin papyrologique :} \textbf{\PsiUn{} (\italique{PSI} XI 1194)} couvre les \textbf{vv. 804–809} et confirme Στρατονίκην au v. 807. --- \textbf{v. 815 (ποθοῦντας R pour πο(ι)οῦντας) :} L. Hauseux consacre un important excursus paléographique à cette faute d'onciale (ΠΟΟΥΝΤΑΣ {\policeFr →} ΠΟΘΟΥΝΤΑΣ, confusion des lettres rondes Ο/Θ sur la forme attique sans iota ποοῦντας). --- \textbf{vv. 822–823 (restitution de \textbf{Bibi}) :} τἀντίον ὁ κανὼν οἱ καλαθίσκοι / τὸ σκιάδειον· R {\policeFr →} \lecon{ταὐτὸ κανοῦν καὶ τὸ σκιάδειον·} \lecon{Bibi} (οἱ καλαθίσκοι est une glose marginale intruse expliquant κανοῦν, corrompu en κανὼν). --- \textbf{v. 834 :} Τηνίοισι R (chute du {\policeGr Σ} initial par haplographie après δίδοσθαι) {\policeFr →} \lecon{Στηνίοισι} {\policeGr Σ}\expo{R}, Souda σ 1089 (fête des \italique{Stènia}). --- \textbf{vv. 822–823 :} τἀντίον ὁ κανὼν οἱ καλαθίσκοι / τὸ σκιάδειον· R : \lecon{ταὐτὸ κανοῦν καὶ τὸ σκιάδειον·} \lecon{Bibi}. --- \textbf{v. 834 :} Τηνίοισι R : \lecon{Στηνίοισι} {\policeGr Σ}, Souda σ 1089.\par}
\end{minipage}\end{adjustbox}
\newpage
\renewcommand{\currHeadLE}{SCHOLIES MARGINALES ({\policeGr Σ}\expo{R}) \& GLOSES INTERLINÉAIRES · FOLIO 20/30}
\renewcommand{\currHeadRE}{f. 172r (Classense 174r) · vv. 800–837 · {\policeGr Σ} 257 à {\policeGr Σ} 267}
\renewcommand{\currHeadLO}{SCHOLIES MARGINALES ({\policeGr Σ}\expo{R}) \& GLOSES INTERLINÉAIRES · FOLIO 20/30}
\renewcommand{\currHeadRO}{f. 172r (Classense 174r) · vv. 800–837 · {\policeGr Σ} 257 à {\policeGr Σ} 267}
\bandeauScholies{Folio 20/30 — f. 172r}{Classense 174r}{vv. 800–837}{Recto (page de droite) · 11 entrées ({\policeGr Σ} 257 à {\policeGr Σ} 267)}
{\fontsize{8.8pt}{10.5pt}\selectfont\noindent\textbf{\color{rubricRed}Disposition topographique des marges sur le f. 172r :} Marge extérieure à droite ; marge intérieure à gauche. — \textit{Fac-similé de référence :} \texttt{folios\_individuels\_ff162v-177r/20\_f172r\_classense174r\_vv0800-0837.jpg}.\par\vspace{3pt}}
\entreeScholie{257}{804}{Ναυσιμάχης μέν γ᾽ ἥττων ἐστὶν Χαρμῖνος· δῆλα δὲ τἀργα.}{Marge supérieure et marge droite (extérieure, f. 172r, grande scholie historique de 7 lignes en haut du recto)}{πέπλασται τὰ ὀνόματα, ἐπεὶ εὐδοκίμουν οἱ Ἀθηναῖοι τότε τῷ ναυτικῷ. Χαρμῖνος δὲ περὶ Σάμον συνεστρατήγησε κατὰ τὸν καιρὸν τοῦτον τοῖς περὶ Φρύνιχον. Ναυσιμάχη δὲ πόρνη καὶ Σαλαβακχώ. Ἄλλως. συνεμάχησε γὰρ Φρυνίχῳ ἐν Θρᾴκῃ ὁ Χαρμῖνος, καὶ ἴσως οὐδὲν ἀξιόλογον ἔπραξεν. διὸ λέγει «δῆλα τὰ ἔργα.»}{πέπλαστ(αι) τὰ ὀνόματα, ἐπεὶ εὐδοκίμουν οἱ Ἀθηναῖοι τότε τῷ ναυτικῷ · Χαρμῖν\expo{ος} δ(ὲ) π(ερὶ) Σάμον συνεστρατήγησε κατὰ τ(ὸν) καιρ(ὸν) τοῦτ(ον) τ\expo{οῖς} π(ερὶ) Φρύνιχ(ον) · Ναυσιμάχη δ(ὲ) πόρνη \textbf{ς} Σαλαβακχώ · ἄλλ\expo{ως} · συνεμάχησε γ(ὰρ) Φρυνίχῳ ἐν Θρᾴκῃ ὁ Χαρμῖν\expo{ος}, \textbf{ς} ἴσ\expo{ως} οὐδὲν ἀξιόλογ(ον) ἔπραξ(εν) · διὸ λέγει «δῆλα τὰ ἔργα · » {\policeFr ÷}}{« Ces noms féminins (Nausimaché, Aristomaché, Stratoniké, Euboulé) ont été forgés pour les besoins du jeu de mots, parce que les Athéniens se distinguaient alors sur mer. Quant au stratège Charminos, il commandait vers cette époque autour de Samos aux côtés de Phrynichos (et subit une défaite navale devant Symè à l'hiver 412/411, cf. Thucydide VIII, 41-42). Et Nausimaché était aussi le nom d'une courtisane connue, tout comme Salabakchô. — Autre rédaction : Charminos combattit aux côtés de Phrynichos et ne fit sans doute rien de brillant, d'où l'ironie : “les faits le montrent assez !” »}
\commScholie{Scholie historique de premier ordre (commentée par L. Hauseux, Introduction.tex, l. 175, et vv 642-fin-bis.tex, ll. 1755-1764) : la défaite navale du stratège Charminos à Symè en janvier 411 av. J.-C. (Thucydide VIII, 42) fournit le terminus post quem irréfutable des Thesmophories ; et l'information Ναυσιμάχη δὲ πόρνη révèle le double sens grivois de Nausimaché.}
\entreeScholie{258}{805}{καὶ μὲν δὴ καὶ Κλεοφῶν χείρων πάντως δήπου Σαλαβακχοῦς.}{Marge droite (extérieure, f. 172r, en face du v. 805)}{Κλεοφῶντα λέγει τὸν λυροποιόν. διαβάλλει δὲ αὐτὸν ὡς κίναιδον.}{Κλεοφῶντ(α) λέγ(ει) τ(ὸν) λυροποιόν · διαβάλλ(ει) δ(ὲ) αὐτ(ὸν) ὡς κίναιδ\expo{ον} {\policeFr ÷}}{« Il parle du démagogue Cléophon le fabricant de lyres (ὁ λυροποιός) ; et il le raille comme un débauché efféminé (κίναιδος, inférieur même à la courtisane Salabakchô). »}
\commScholie{Commenté par Louis Hauseux (vv 642-fin-bis.tex, ll. 1765-1788).}
\entreeScholie{259}{806}{πρὸς Ἀριστομάχην δὲ χρόνου πολλοῦ,}{Interlinéaire (supra lineam, f. 172r, v. 806)}{πόρνης ὄνομα.}{πόρν\expo{ης} ὄνομα {\policeFr ÷}}{« Nom d'une courtisane (pour Σαλαβακχοῦς v. 805, cf. Hésychius σ 88 et Souda σ 50). »}
\commScholie{Glose onomastique de R\expo{2} se rapportant à Σαλαβακχοῦς (fin du v. 805, déjà mentionnée dans les Cavaliers, v. 765), célèbre courtisane athénienne dont le nom est opposé ici aux noms guerriers forgés par le Chœur.}
\entreeScholie{260}{808}{ἀλλ᾽ Εὐβούλης τῶν πέρυσίν τις βουλευτής ἐστιν ἀμείνων / παραδοὺς ἑτέρῳ τὴν βουλείαν;}{Marge droite (extérieure, f. 172r, grande scholie alexandrine de 11 lignes en face des vv. 808–809)}{καὶ ἡ Εὐβούλη πέπλασται. τὸ δὲ ὅλον τί βούλεται, οὐκ ἔστι σαφές. ἀλλ’ οὐδὲ παρὰ Πλάτωνι ἐν τῷ Ὑπερβόλῳ· λέγει γάρ τι περὶ τοῦ ἐπιλαγχάνειν τοῖς βουλευταῖς «εὐτυχεῖς, ὦ δέσποτα. — τί δ’ ἔστι; — βουλεύειν ὀλίγου ’λαχες πάνυ· ἀτὰρ οὐ λαχὼν ὅμως ἔλαχες, ἢν νοῦν ἔχῃς. πῶς ἢν ἔχω νοῦν; ὅτι πονηρῷ καὶ ξένῳ ἐπέλαχες ἀνδρί· [...] οὐδέπω γὰρ ἐλευθέρῳ.» καὶ ἐν ταὐτῷ «ἄπερρ’· ἐγὼ δ’ ὑμῖν τὸ πρᾶγμα δὴ φράσω. Ὑπερβόλῳ (βουλῆς γὰρ) ἄνδρες ἐπέλαχον.» τὸ βουλεῦσαι βουλείαν ἔλεγον.}{\textbf{ς} ἡ Εὐβούλη πέπλαστ(αι) · τ(ὸ) δ(ὲ) ὅλ(ον) τί βούλετ(αι), οὐκ ἔστι σαφές · ἀλλ’ οὐδὲ παρὰ Πλάτωνι ἐν τῷ Ὑπερβόλῳ· λέγει γάρ τι π(ερὶ) τ\expo{οῦ} ἐπιλαγχάνε(ιν) τ\expo{οῖς} βουλευτ\expo{αῖς} «εὐτυχεῖς ὦ δέσποτα · τί δ’ ἔστι; βουλεύειν ὀλίγου ’λαχες πάνυ · ἀτὰρ οὐ λαχὼν ὅμ\expo{ως} ἔλαχες, ἢν νοῦν ἔχῃς · πῶς ἢν ἔχω νοῦν; ὅτι πονηρῷ \textbf{ς} ξένῳ ἐπέλαχες ἀνδρί· οὐδέπω γ(ὰρ) ἐλευθέρῳ · » \textbf{ς} ἐν ταὐτῷ «ἄπερρ’· ἐγὼ δ’ ὑμῖν τ(ὸ) πρᾶγμα δὴ φράσω · Ὑπερβόλῳ (βουλ\expo{ῆς} γ(ὰρ)) ἄνδρες ἐπέλαχ(ον) · » τ(ὸ) βουλεῦσ(αι) βουλείαν ἔλεγ(ον) {\policeFr ÷}}{« Le nom d'Euboulé (“Bon-Conseil”) est lui aussi choisi pour le jeu de mots (face aux bouleutes de l'année précédente qui avaient cédé leur place aux Quatre-Cents). Mais ce qu'il veut dire au juste n'est pas clair, pas plus que chez Platon le Comique dans son Hyperbolos (fr. 182-183 K.-A.), où il parle du tirage au sort des bouleutes suppléants (ἐπιλαγχάνειν) : “Tu as de la chance, maître ! — Qu'y a-t-il ? — Il s'en est fallu de bien peu que tu ne sois tiré au sort pour siéger au Conseil ! Car sans avoir été tiré titulaire, tu l'as été tout de même, si tu as du bon sens. — Comment cela ? — Parce que tu as été tiré comme suppléant (ἐπέλαχες) d'un coquin et d'un étranger [...] qui n'est pas encore citoyen libre !” Et dans la même pièce : “Va-t'en au diable ! Quant à moi, je vais vous expliquer l'affaire : à Hyperbolos (car il est du Conseil), des citoyens ont été tirés comme suppléants.” — Et l'on appelait βουλεία la charge de bouleute. »}
\commScholie{Joyau inestimable de la critique alexandrine (Didyme / Symmaque, discuté par L. Hauseux, vv 642-fin-bis.tex, ll. 1798-1820) : c'est à cette seule scholie du Ravennas que nous devons la conservation des deux célèbres fragments de l'Hyperbolos de Platon le Comique sur l'institution des bouleutes suppléants (ἐπιλαχόντες) !}
\entreeScholie{261}{811}{οὐδ᾽ ἂν κλέψασα γυνὴ ζεύγει κατὰ πεντήκοντα τάλαντα / εἰς πόλιν ἔλθοι τῶν δημοσίων·}{Marge droite (extérieure, f. 172r, en face des vv. 811–812)}{λέγει ἀπὸ συμμάχων ἢ στρατείας. τὸ δὲ ζεύγει ἀντὶ τοῦ ἁμάξαις ἐκφέρειν δεῖ. [...] διὸ καὶ οὐκ εἶπεν οὕτως πεντήκοντα, ἀλλὰ κατὰ πεντήκοντα, τὸ συνεχὲς ἐκφαίνων.}{λέγ(ει) ἀπ(ὸ) συμμάχ\expo{ων} ἢ στρατεί\expo{ας} · τ(ὸ) δ(ὲ) ζεύγ(ει) ἀν\auDessus{τ}{ι} τ\expo{οῦ} ἁμάξ\expo{αις} ἐκφέρ(ειν) δεῖ [...] διὸ \textbf{ς} οὐκ εἶπ(εν) οὕτ\expo{ως} πεντήκοντ(α), ἀλλὰ κατὰ πεντήκοντ(α), τ(ὸ) συνεχ(ὲς) ἐκφαίν\expo{ων} {\policeFr ÷}}{« “Revenir à Athènes (ἐς πόλιν ἔλθοι, ou monter à l'Acropole en char)” : il veut dire au retour d'une ambassade chez les alliés ou d'une expédition militaire. Quant à ζεύγει (“en attelage à deux chevaux”), il l'emploie pour le train somptueux de ceux qui ont détourné les fonds publics. C'est pourquoi il n'a pas dit simplement “cinquante talents”, mais “par tranches de cinquante talents” (κατὰ πεντήκοντα), soulignant la répétition continuelle de leurs rapines. »}
\commScholie{Analyse très fine de la valeur distributive de κατὰ πεντήκοντα τάλαντα.}
\entreeScholie{262}{822}{ταὐτὸ κανοῦν καὶ τὸ σκιάδειον [τἀντίον, ὁ κανών, οἱ καλαθίσκοι, τὸ σκιάδειον R]·}{Marge droite (extérieure, f. 172r, en face du v. 823)}{σκέπασμά τι ἐπὶ τῆς κεφαλῆς βαλλόμενον, ἢ μέρος τοῦ ἱστίου.}{σκέπασμά τι ἐπὶ τ\expo{ῆς} κεφαλ\expo{ῆς} βαλλόμεν\expo{ον}, ἢ μέρ\expo{ος} τ\expo{οῦ} ἱστί\expo{ου} {\policeFr ÷}}{« Une sorte d'abri que l'on déploie au-dessus de la tête (l'ombrelle, τὸ σκιάδειον), ou bien une pièce du métier à tisser (pour ἀντίον, l'ensouple). »}
\commScholie{Au v. 822, le mot οἱ καλαθίσκοι qui fausse le rythme anapestique dans R\expo{1} est analysé par Louis Hauseux (vv 642-fin-bis.tex, l. 1891) comme une glose de κανοῦν insérée dans le texte.}
\entreeScholie{263}{825}{τοῖς δ᾽ ἡμετέροις ἀνδράσι τούτοις / ἀπόλωλεν μὲν πολλοῖς ὁ κανὼν / ἐκ τῶν οἴκων αὐτῇ λόγχῃ,}{Marge droite (extérieure, f. 172r, en face du v. 825)}{τὴν κάμακα τῆς λόγχης λέγει. οἱ μέν φασι κανόνα αὐτὸν λέγειν τὸ ξύλον τῆς λόγχης, οἱ δὲ τὸν κανόνα τῆς ἀσπίδος.}{τ(ὴν) κάμακα τ\expo{ῆς} λόγχ\expo{ης} λέγ(ει) · οἱ μέν φασ(ι) κανόνα αὐτ(ὸν) λέγ(ειν) τ(ὸ) ξύλ\expo{ον} τ\expo{ῆς} λόγχ\expo{ης}, οἱ δ(ὲ) τ(ὸν) κανόνα τ\expo{ῆς} ἀσπίδ\expo{ος} {\policeFr ÷}}{« Par ὁ κανών (“la règle de tisserand / la hampe droite”), elle désigne chez les maris le fût de la lance (τὴν κάμακα τῆς λόγχης). Les uns disent qu'Aristophane appelle κανών le bois de la lance (puisqu'il ajoute “avec la pointe de fer elle-même”, αὐτῇ λόγχῃ), les autres la barre transversale du bouclier. »}
\commScholie{Repris par la Souda κ 312. Le double sens entre la baguette de lice du métier à tisser des femmes (κανών) et la hampe de lance perdue à la guerre par les hommes est parfaitement saisi.}
\entreeScholie{264}{826}{αὐτῇ λόγχῃ}{Interlinéaire (supra lineam au-dessus de λόγχῃ, f. 172r, v. 826)}{τῷ σιδήρῳ.}{τῷ σιδήρῳ {\policeFr ÷}}{« Avec la pointe de fer de la lance (au sens propre de λόγχη par opposition à la hampe de bois, κανών). »}
\commScholie{Précision lexicale exacte de R\expo{2} distinguant, dans le tour comique ὁ κανὼν ... αὐτῇ λόγχῃ (« le fût avec la pointe elle-même »), la hampe de bois (κανών) et le fer de lance proprement dit (τῷ σιδήρῳ).}
\entreeScholie{265}{828}{ἐν ταῖς στρατιαῖς}{Marge droite (extérieure, f. 172r, en face du v. 828)}{συνέχεεν καὶ οὗτος ὡς Εὔπολις πολλάκις· στρατιὰ μὲν γὰρ τὸ πλῆθος, στρατεία δὲ ἡ στράτευσις.}{συνέχε(εν) \textbf{ς} οὗτ\expo{ος} ὡς Εὔπολ(ις) πολλάκ(ις) · στρατιὰ μ(ὲν) γ(ὰρ) τ(ὸ) πλῆθ\expo{ος}, στρατεία δ(ὲ) ἡ στράτευσ(ις) {\policeFr ÷}}{« Aristophane lui aussi, tout comme Eupolis le fait souvent, a confondu (pour les besoins du mètre anapestique) στρατιά et στρατεία : car au sens strict, στρατιά (avec ι bref) désigne le corps de troupes, tandis que στρατεία (avec ει long) désigne l'expédition militaire (la campagne). »}
\commScholie{Observation métrique et grammaticale d'Aristarque / Héliodore : le mètre anapestique exigeant deux brèves (-τĭαῖς), l'attique emploie couramment στρατιά au sens de « campagne militaire ».}
\entreeScholie{266}{829}{ἔρριπται τὸ σκιάδειον.}{Marge droite (extérieure, f. 172r, en face du v. 829)}{τὴν ἀσπίδα λέγει· σκιάζει γὰρ τὸν ἄνδρα. αἰνίττεται δὲ καὶ εἰς Κλεώνυμον.}{τ(ὴν) ἀσπίδα λέγ(ει) · σκιάζ(ει) γ(ὰρ) τ(ὸν) ἄνδρα · αἰνίττετ(αι) δ(ὲ) \textbf{ς} εἰς Κλεώνυμ\expo{ον} {\policeFr ÷}}{« Par “l'ombrelle” (τὸ σκιάδειον), elle désigne par métaphore le bouclier (τὴν ἀσπίδα), parce qu'il couvre et protège le guerrier de son ombre ; et elle vise particulièrement Cléonyme (le célèbre ῥίψασπις qui avait jeté son bouclier à Délion). »}
\commScholie{Cité par Louis Hauseux (vv 642-fin-bis.tex, ll. 1966-1972) et repris par la Souda σ 597.}
\entreeScholie{267}{834}{προεδρίαν τ᾽ αὐτῇ δίδοσθαι Στηνίοισι καὶ Σκίροις,}{Marge inférieure et marge droite (bas du f. 172r, en face du v. 834)}{ἀμφότεραι ἑορταὶ γυναικῶν, τὰ μὲν Στήνια πρὸ δυεῖν τῶν Θεσμοφορίων Πυανεψιῶνος θ', τὰ δὲ Σκίρα λέγεσθαί φασί τινες τὰ γινόμενα ἱερὰ ἐν τῇ ἑορτῇ ταύτῃ Δήμητρι καὶ Κόρῃ. οἱ δὲ, ὅτι ἐπὶ Σκίρῳ θύεται τῇ Ἀθηνᾷ.}{ἀμφότερ(αι) ἑορτ(αὶ) γυναικ\expo{ῶν}, τὰ μ(ὲν) Στήνια πρὸ δυεῖν τ\expo{ῶν} Θεσμοφορί\expo{ων} Πυανεψιῶν\expo{ος} θ', τὰ δ(ὲ) Σκίρα λέγεσθαί φασί τιν\expo{ες} τὰ γινόμεν(α) ἱερὰ ἐν τῇ ἑορτῇ ταύτ(ῃ) Δήμητρι \textbf{ς} Κόρῃ · οἱ δ(ὲ), ὅτι ἐπὶ Σκίρῳ θύετ(αι) τῇ Ἀθηνᾷ {\policeFr ÷}}{« Ce sont toutes deux des fêtes réservées aux femmes : les Sténies (Στήνια, où les femmes s'échangeaient des brocards nocturnes, cf. στηβάζειν / στηνιῶσαι) ont lieu deux jours avant les Thesmophories, le 9 du mois de Pyanepsion ; quant aux Skira (Σκίρα, le 12 de Skirophorion), certains disent qu'on appelle ainsi les rites célébrés en l'honneur de Déméter et de Coré, d'autres que la fête tire son nom du sacrifice offert à Athéna Skiras au lieu-dit Skiron (sur la voie d'Éleusis). »}
\commScholie{Scholie héortologique capitale (commentée par L. Hauseux, Bibi, v. 834, vv 642-fin-bis.tex, ll. 1982-1987, et l. 2812 n. 11) : alors que le scribe R\expo{1} avait commis la faute d'aphérèse Τηνίοισι, la scholie {\policeGr Σ}\expo{R} (ainsi que la Souda σ 1089 et Hésychius σ 1825) conserve la vraie leçon Στηνίοισι !}
\clearpage
\renewcommand{\currHeadLE}{ΑΡΙΣΤΟΦΑΝΟΥΣ ΘΕΣΜΟΦΟΡΙΑΖΟΥΣΑΙ · CODEX RAVENNAS 429 (\R) · FOLIO 21/30}
\renewcommand{\currHeadRE}{f. 172v (Classense 174v) · vv. 838–875 · Collation diplomatique}
\renewcommand{\currHeadLO}{ΑΡΙΣΤΟΦΑΝΟΥΣ ΘΕΣΜΟΦΟΡΙΑΖΟΥΣΑΙ · CODEX RAVENNAS 429 (\R) · FOLIO 21/30}
\renewcommand{\currHeadRO}{f. 172v (Classense 174v) · vv. 838–875 · Collation diplomatique}
\noindent\begin{adjustbox}{max totalheight=30.4cm, max width=\textwidth, keepaspectratio}\begin{minipage}{\textwidth}
\bandeauRav{f. 172v}{Classense 174v}{838–875}{Folio 21/30 · Verso (page de gauche) · 15 scholies \& gloses ({\policeGr Σ} 268–282)}
{\fontsize{9.3pt}{10.8pt}\selectfont
\ligneRav{838}{χορ/}{ὑστέρ\auDessus{ο}{α}ν αὐτὴν καθῆσθαι σκάφιον ἀποκεκαρμένην}{\textbf{838} \lecon{ὑστέρ\auDessus{ο}{α}ν} \Rpc : ὑστάτην Σ}
\ligneRav{839}{}{τῆς τὸν ἀνδρεῖον τεκούσης· τῷ γὰρ εἰκὸς \difference{ωὖπολι}}{\textbf{839} ὦ πόλεις R : ὦ πόλις Gelenius : \lecon{ωὖπολι} Bibi}
\ligneRav{840}{}{τὴν Ὑπερβόλου καθῆσθαι μητέρ᾽ ἠμφιεσμένην}{}
\ligneRav{}{}{λευκὰ καὶ κόμας καθεῖσαν πλησίον τῆς Λαμάχου·}{}
\ligneRav{842}{}{καὶ δανείζειν \difference{χρήμαθ᾽}, ἥν χρῆν, \difference{εἰ} \difference{δανείσειέν} τινι}{\textbf{842} χρήματα τ' R : \lecon{χρήμαθ᾽} Scaliger \espApp \lecon{ἥν χρῆν} R : ᾗ χρῆν Scaliger \espApp χρῆν ἦν R : \lecon{χρῆν εἰ} Brunck \espApp δανείσειε R : \lecon{δανείσειέν} Biset}
\ligneRav{}{}{καὶ τόκον πράττοιτο, διδόναι μηδέν᾽ ἀνθρώπων τόκον,}{}
\ligneRav{844}{}{ἀλλ᾽ \difference{ἀφαιρεῖσθαι} βίᾳ τὰ \difference{χρήματ᾽} εἰπόντας \difference{τοδί} (τ\auDessus{ὰ}{ο}δὶ \Rpc),}{\textbf{844} ἀφερεῖσθαι R : \lecon{ἀφαιρεῖσθαι} \Mon \espApp χρήματατ᾽ R : \lecon{χρήματ᾽} Grynaeus \espApp τ\auDessus{ὰ}{ο}δὶ \Rpc : \lecon{τοδί} \Mon}
\ligneRav{845}{}{`ἀξία γοῦν εἶ τόκου τεκοῦσα τοιοῦτον τόκον.'}{}
\ligneRav{}{μ\auDessus{ν}{Η}}{ἰλλὸς γεγένημαι προσδοκῶν· ὁ δ᾽ οὐδέπω·}{}
\ligneRav{}{}{τί δῆτ᾽ ἂν εἴη τοὐμποδών; οὐκ ἔσθ᾽ ὅπως}{}
\ligneRav{848}{}{οὐ τὸν \difference{Παλαμήδη} ψυχρὸν ὄντ᾽ αἰσχύνεται·}{\textbf{848} Παλαμήδην R : \lecon{Παλαμήδη} Elmsley (\italique{Quart. Rev.} 7 (1812) p.454)}
\ligneRav{}{}{τῷ δῆτ᾽ ἂν αὐτὸν προσαγαγοίμην δράματι;}{}
\ligneRav{850}{}{ἐγᾦδα· τὴν καινὴν Ἑλένην μιμήσομαι·}{}
\ligneRav{851}{}{\difference{πάντως δ᾽} ὑπάρχει μοι γυναικεία στολή·}{\textbf{851} πάντως R : \lecon{πάντως δ᾽} Bentley, cf. v.1012}
\ligneRav{852}{\auDessus{γυ}{ν}}{τί αὖ σὺ κυρκανᾷς; \difference{τί κοικύλλεις} ἔχων;}{\textbf{852} ἢ τί κοικύλλεις R, Suid. : \lecon{τί κοικύλλεις} Faber}
\ligneRav{853}{}{πικρὰν \difference{ἑλένην} ὄψει τάχ᾽, εἰ μὴ κοσμίως}{\textbf{853} Ἑλένη R : \lecon{ἑλένην} Bibi ap. Hesych. ε 1995}
\ligneRav{}{}{ἕξεις, ἕως ἂν τῶν πρυτάνεών τις φανῇ·}{}
\ligneRav{855}{μ\auDessus{ν}{Η}}{Νείλου μὲν αἵδε καλλιπάρθενοι ῥοαί,}{}
\ligneRav{}{}{ὃς ἀντὶ δίας ψακάδος Αἰγύπτου πέδον}{}
\ligneRav{}{}{λευκῆς νοτίζει μελανοσυρμαῖον λεών·}{}
\ligneRav{}{\auDessus{γυ}{ν}}{πανοῦργος εἶ νὴ τὴν Ἑκάτην τὴν φωσφόρον·}{}
\ligneRav{}{μ\auDessus{ν}{Η}}{ἐμοὶ δὲ γῆ μὲν πατρὶς οὐκ ἀνώνυμος}{}
\ligneRav{860}{—}{Σπάρτη, πατὴρ δὲ Τυνδάρεως· : \difference{σοί} (σ\auDessus{ί}{ύ} \Rdeux) \difference{γ᾽} ὦλεθρε}{\textbf{860} σύ (ex σί) \Rdeux : \lecon{σοί} Zanetti (\italique{tibi} Divus) \espApp τ' R : \lecon{γ᾽} Gelenius}
\ligneRav{861}{}{πατὴρ ἐκεῖνός \difference{ἐστι}; Φρυνώνδας μὲν οὖν·}{\textbf{861} ἐστιν R : \lecon{ἐστι} \Mon}
\ligneRav{862}{μ\auDessus{ν}{Η}}{Ἑλένη δ᾽ ἐκλήθην. : αὖθις αὖ \difference{γίγνει} γυνή,}{\textbf{862} γίγνη R : \lecon{γίγνει} Brunck}
\ligneRav{863}{}{πρὶν τῆς ἑτέρας δοῦναι γυναικίσεως δίκην;}{\textbf{863} \lecon{δοῦναι post ἑτέρας} R : post δίκην Suid. γ 496}
\ligneRav{864}{μ\auDessus{ν}{Η}}{ψυχαὶ δὲ πολλαὶ δι᾽ ἔμ᾽ ἐπὶ \difference{Σκαμανδρίοις}}{\textbf{864} Σκαμανδρίαις R : \lecon{Σκαμανδρίοις} Fritzsche ex Eur.\italique{Hel} 52 et 609}
\ligneRav{865}{—}{ῥοαῖσιν ἔθανον. : ὤφελες δὲ καὶ σύ γε·}{}
\ligneRav{}{μ\auDessus{ν}{Η}}{κἀγὼ μὲν ἐνθάδ᾽ εἴμ᾽· ὁ δ᾽ ἄθλιος πόσις}{}
\ligneRav{867}{}{οὑμὸς \difference{Μενέλεως} οὐδέπω προσέρχεται·}{\textbf{867} Μενέλαος R : \lecon{Μενέλεως} Brunck, cf. v.901}
\ligneRav{868}{—}{τί οὖν ἔτι ζῶ; : τῶν κοράκων πονηρίᾳ·}{\textbf{868} τῶν κορ. πον. continuat (Μνη.) R : \lecon{(Γυνή)} Σ, Ellebodius}
\ligneRav{}{μ\auDessus{ν}{Η}}{ἀλλ᾽ ὥσπερ αἰκάλλει τι καρδίαν ἐμήν·}{}
\ligneRav{870}{}{μὴ ψεῦσον ὦ Ζεῦ τῆς ἐπιούσης ἐλπίδος·}{\textbf{870} \lecon{ἐλπίδος} R : ἡμέρας ap. Greg. Cor. p.14-15, Σ ad Eur. \italique{Hec}.225 \espApp post v. 870 < οἶσθ᾽ οὖν ὃ δρᾶσον; — ˘ — | × — ˘ — > \lecon{om.} R : restituit ap. Σ Bibi (cf. Eur. \italique{Hel.}325)}
\ligneRav{871}{\auDessus{εὐ}{ρ}}{τίς τῶνδ᾽ ἐρυμνῶν δωμάτων ἔχει κράτος,}{\textbf{871} ὡς Περσεύς in mg. R : \lecon{ὡς Μενέλαος} Zanetti}
\ligneRav{872}{}{ὅστις \difference{ξένους} δέξαιτο ποντίῳ σάλῳ}{\textbf{872} ξένος R : \lecon{ξένους} Zanetti}
\ligneRav{873}{}{κάμνοντας ἐν χειμῶνι καὶ ναυαγίαις;}{\textbf{873} \lecon{κάμνοντας} R : καμόντας Lenting}
\ligneRav{}{μ\auDessus{ν}{Η}}{Πρωτέως τάδ᾽ ἐστὶ μέλαθρα. : ποίου Πρωτέως;}{}
\ligneRav{875}{}{\difference{ὤ τρισκακοδαίμων}, ψεύδεται νὴ τὼ θεώ,}{\textbf{875} ὦ τρισκακόδαιμον R : \lecon{ὤ τρισκακοδαίμων} Wilamowitz ad \italique{Lys.}449}
}
\vspace{3pt}{\color{ruleGold}\hrule height 0.6pt}\vspace{3pt}
{\fontsize{8.0pt}{9.4pt}\selectfont
\noindent{\bfseries\scshape\color{rubricRed}Loci paralleli et testimonia (f. 172v, vv. 838–875) :} \textbf{838} (σκάφιον ἀποκεκαρμένην) Photius σ 291 σκάφιον ἀποκεκαρμένην ; cf. Hesychius σ 879 σκάφιον ; Photius σ 292 ({\policeFr ≈} Scholium) ; \Eustathe 907.40, 1292.61 ; \AelDion σ 23 \espTest \textbf{846} Suidas ι 323 {\policeFr ∼} Zonaras 1103 ἰλλός \espTest (ἰλλός) Pollux II 52 ; \Eustathe 768.9 ; \AelDion ι 6 ; cf. Moeris ι 21 (ἰλλόν = στραβόν) ; Hesychius ι 578 ἰλλός ; Hesychius η 257 (ἠίλλοντο = ἐστράβιζον) \espTest \textbf{852} Suidas κ 2534 \espTest (κοικύλλεις) Hesychius κ 3232 κοικύλλειν {\policeFr ∼} Photius κ 860 \espTest \textbf{861} (Φρυνώνδας) Suidas φ 700 ; \Syn φ 211 Φρυνώνδας {\policeFr ∼} Hesychius φ 938 \espTest \textbf{863} Suidas γ 496 γυναίκισις et Zonaras 459 πρὶν τῆς ἑτέρας γυναικίσεως δίκην δοῦναι \espTest \textbf{868} Suidas τ 648 \espTest \textbf{869} (αἰκάλλει τι καρδίαν ἐμήν) Suidas αι 161 ; Zonaras 91 (τι omisso) \espTest \textbf{870} \GregCor p.14-15 (= scholium\expo{Mosch.} Euripides, \italique{Hecub.} 225) (sed ἡμέρας pro ἐλπίδος) ; Suidas ψ 62 \espTest (μὴ ψεῦσον) \Antiatt 107.31 \espTest \textbf{870} \GregCor pp. 14-15 (= scholium\expo{Mosch.} Euripides, \italique{Hecub.} 225) (sed ἡμέρας pro ἐλπίδος) ; Suidas ψ 62\par\vspace{2.5pt}
{\color{ruleGold}\hrule height 0.35pt}\vspace{2.5pt}
\noindent{\bfseries\scshape\color{rubricRed}Suite (en commentaire) de l'apparat critique et paléographique :} \textbf{Corrections supralinéaires de \Rpc{} / \Rdeux{} :} Au \textbf{v. 838}, \lecon{ὑστέρ\auDessus{ο}{α}ν} (\Rpc{} trace α au-dessus de ο ; {\policeGr Σ}\expo{R} lisait ὑστάτην) ; au \textbf{v. 844}, \lecon{τ\auDessus{ὰ}{ο}δὶ} (\Rpc{}, τοδί \Mon{}) ; au \textbf{v. 860}, \lecon{σύ} corrigé par \Rdeux{} à partir de σί (pour σοί Zanetti). --- \textbf{Deux conjectures capitales de Bibi (vv. 839 et 853) :} 1° Au \textbf{v.~839}, là où \R{} porte ὦ πόλεις (ὦ πόλις Gelenius), Louis Hauseux (Bibi) restitue le vocatif par crase \lecon{ωὖπολι} (ὦ Εὔπολι, interpellation adressée au poète Eupolis, mort au combat en 411, argument majeur pour dater les \italique{Thesmophories} de \textbf{410 av.~J.-C.}) ; 2° au \textbf{v.~853}, là où \R{} écrit πικρὰν Ἑλένην ὄψει τάχ' (v.~853) après τίς δ' αὖτ' ἐστὶν ἥδε; (v.~852), Bibi restitue l'accusatif \lecon{ἑλένην} d'après Hésychius ε 1995 (ἑλένη· λαμπάς, δετή, torche de roseaux tressés que tient Mnésiloque, en calembour scénique sur le nom d'\italique{Hélène}). --- \textbf{v. 838 :} \lecon{ὑστέρ\auDessus{ο}{α}ν} \Rpc{} : ὑστάτην {\policeGr Σ}. --- \textbf{v. 839 :} ὦ πόλεις R (ὦ πόλις Gelenius) : \lecon{ωὖπολι} Bibi. --- \textbf{v. 853 :} Ἑλένη R : \lecon{ἑλένην} Bibi ap. Hésychius ε 1995. --- \textbf{v. 868 :} τῶν κορ. πον. continuat (Μνη.) R : \lecon{(Γυνή)} {\policeGr Σ}, Ellebodius.\par}
\end{minipage}\end{adjustbox}
\newpage
\renewcommand{\currHeadLE}{SCHOLIES MARGINALES ({\policeGr Σ}\expo{R}) \& GLOSES INTERLINÉAIRES · FOLIO 21/30}
\renewcommand{\currHeadRE}{f. 172v (Classense 174v) · vv. 838–875 · {\policeGr Σ} 268 à {\policeGr Σ} 282}
\renewcommand{\currHeadLO}{SCHOLIES MARGINALES ({\policeGr Σ}\expo{R}) \& GLOSES INTERLINÉAIRES · FOLIO 21/30}
\renewcommand{\currHeadRO}{f. 172v (Classense 174v) · vv. 838–875 · {\policeGr Σ} 268 à {\policeGr Σ} 282}
\bandeauScholies{Folio 21/30 — f. 172v}{Classense 174v}{vv. 838–875}{Verso (page de gauche) · 15 entrées ({\policeGr Σ} 268 à {\policeGr Σ} 282)}
{\fontsize{8.8pt}{10.5pt}\selectfont\noindent\textbf{\color{rubricRed}Disposition topographique des marges sur le f. 172v :} Marge extérieure à gauche ; marge intérieure à droite. — \textit{Fac-similé de référence :} \texttt{folios\_individuels\_ff162v-177r/21\_f172v\_classense174v\_vv0838-0875.jpg}.\par\vspace{3pt}}
\entreeScholie{268}{838}{ὑστέραν αὐτὴν καθῆσθαι σκάφιον ἀποκεκαρμένην / τῆς τεκούσης τὸν ἀνδρεῖον·}{Marge supérieure et marge gauche (extérieure, f. 172v, en haut du verso, Planche 23)}{εἶδος κουρᾶς δουλικῆς. τὸ ἑξῆς. ὑστέραν τῆς τὸν ἀνδρεῖον τεκούσης. ἀντὶ τοῦ μιᾶς γυναικὸς ὑστέραν. ἐὰν δὲ ὑστάτην, ἁπάντων ὑστάτην.}{εἶδ\expo{ος} κουρ\expo{ᾶς} δουλικ\expo{ῆς} · τ(ὸ) ἑξ\expo{ῆς} · ὑστέρ(αν) τ\expo{ῆς} τ(ὸν) ἀνδρεῖ\expo{ον} τεκούσ\expo{ης} · ἀν\auDessus{τ}{ι} τ\expo{οῦ} μιᾶς γυναικ\expo{ὸς} ὑστέρ(αν) · ἐὰν δ(ὲ) ὑστάτ(ην), ἁπάντ\expo{ων} ὑστάτ(ην) {\policeFr ÷}}{« σκάφιον ἀποκεκαρμένην (“tondue en sébile”) : espèce de coupe de cheveux servile (rasée tout autour de la tête à l'aide d'une écuelle). — L'ordre syntaxique est : qu'elle soit assise au-dessous (ὑστέραν) de celle qui a enfanté un citoyen courageux. Le comparatif ὑστέραν s'emploie quand on compare à une seule femme, le superlatif ὑστάτην par rapport à toutes. »}
\commScholie{Repris par Photius σ 291-292, Hésychius σ 879 et la Souda σ 546 (commenté par L. Hauseux, vv 642-fin-bis.tex, ll. 1989-1995, et Planche 23 du Manuel).}
\entreeScholie{269}{840}{τὴν Ὑπερβόλου καθῆσθαι μητέρ᾽ ἠμφιεσμένην / λευκὰ καὶ κόμας καθεῖσαν πλησίον τῆς Λαμάχου,}{Marge gauche (extérieure, f. 172v, grande scholie alexandrine de 6 lignes en face du v. 840)}{πολλάκις εἶπον ὅτι προβάλλει ὁ Ὁμηρικὸς ἐν τῷ πρὸς Πολύβιον προτατικῷ τίς ἡ Ὑπερβόλου μήτηρ. οὐκ ἔχομεν δὲ εἰπεῖν. εἰ δὲ χρὴ τὰ τοιαῦτα ζητεῖν, ζητείτωσαν καὶ τὴν ἑξῆς, τὴν Λαμάχου· τὸν γὰρ πατέρα μόνον οἶδα Λαμάχου, Ξενόφαντον.}{πολλάκ(ις) εἶπ\expo{ον} ὅτι προβάλλ(ει) ὁ Ὁμηρικ\expo{ὸς} ἐν τῷ πρ(ὸς) Πολύβι\expo{ον} προτατικῷ τίς ἡ Ὑπερβόλ\expo{ου} μήτ(ηρ) · οὐκ ἔχομ(εν) δ(ὲ) εἰπ(εῖν) · εἰ δ(ὲ) χρὴ τὰ τοιαῦτ(α) ζητ(εῖν), ζητείτωσ(αν) \textbf{ς} τ(ὴν) ἑξ\expo{ῆς}, τ(ὴν) Λαμάχ\expo{ου} · τ(ὸν) γ(ὰρ) πατέρα μόν\expo{ον} οἶδα Λαμάχ\expo{ου}, Ξενόφαντ\expo{ον} {\policeFr ÷}}{« J'ai déjà dit bien souvent que “l'Homérique” (Séleucos d'Alexandrie ou Ptolémée Épithète), dans son Recueil de questions adressé à Polybe (ἐν τῷ πρὸς Πολύβιον προτατικῷ), pose la question : “Quel était le nom de la mère d'Hyperbolos ?” Or nous n'avons aucun moyen de le dire ! Et s'il faut absolument chercher ce genre de futilités, qu'ils cherchent donc aussi le nom de la suivante, la mère de Lamachos : car pour ma part, je ne connais que le nom du père de Lamachos, qui s'appelait Xénophantos (cf. Thucydide VI, 8, 2). »}
\commScholie{Savoureuse scholie personnelle de Didyme ou de Symmaque (cf. Σ 1175 sur le Προτατικός de Séleucos), qui se moque avec humour des « questions insolubles » (προβλήματα / προτατικά) posées dans les banquets de grammairiens alexandrins.}
\entreeScholie{270}{841}{πλησίον τῆς Λαμάχου,}{Marge gauche (extérieure, f. 172v, sous Σ 840)}{ἐπαινεῖ τὸν Λάμαχον νῦν· ἤδη γὰρ ἐτεθνήκει ἐν Σικελίᾳ τετάρτῳ ἔτει πρότερον.}{ἐπαιν(εῖ) τ(ὸν) Λάμαχ\expo{ον} νῦν · ἤδη γ(ὰρ) ἐτεθνήκ(ει) ἐν Σικελίᾳ τετάρτῳ ἔτ(ει) πρότερ\expo{ον} {\policeFr ÷}}{« Il fait maintenant l'éloge du stratège Lamachos (qu'il avait pourtant raillé de son vivant dans les Acharniens et la Paix) : car Lamachos était déjà mort héroïquement au combat en Sicile (devant Syracuse à l'été 414 av. J.-C.), quatre ans auparavant (en compte inclusif : 414-411). »}
\commScholie{Troisième confirmation chronologique majeure de la date de 411 av. J.-C. (après Σ 32 et Σ 190, cf. L. Hauseux, Introduction.tex, l. 175), et remarque littéraire très juste (cf. Grenouilles v. 1039 ὁ Λάμαχος ἥρως).}
\entreeScholie{271}{846}{ἰλλὸς γεγένημαι προσδοκῶν·}{Marge gauche (extérieure, f. 172v, en face du v. 846)}{τυφλὸς, διεστραμμένος τὴν ὄψιν. Σώφρων «ἰλλοτέρα τᾶν κορωνᾶν.»}{τυφλ\expo{ὸς}, διεστραμμέν\expo{ος} τ(ὴν) ὄψ(ιν) · Σώφρ(ων) «ἰλλοτέρα τᾶν κορωνᾶν» {\policeFr ÷}}{« ἰλλός : bigle, qui a les yeux de travers à force de guetter (διεστραμμένος τὴν ὄψιν). Ainsi Sophron de Syracuse (dans ses Mimes, fr. 64 K.-A.) : “Plus bigle que les corneilles (ἰλλοτέρα τᾶν κορωνᾶν).” »}
\commScholie{Repris par la Souda ι 319 et l'Etymologicum Magnum : conserve une précieuse citation dorienne des Mimes de Sophron de Syracuse.}
\entreeScholie{272}{848}{οὐκ ἔσθ᾽ ὅπως / οὐ τὸν Παλαμήδη ψυχρὸν ὄντ᾽ αἰσχύνεται·}{Interlinéaire (supra lineam, f. 172v, v. 848)}{ἄνω γὰρ ἐκέχρητο τῇ ἐκ Παλαμήδους ἐπινοίᾳ.}{ἄνω γ(ὰρ) ἐκέχρητ(ο) τῇ ἐκ Παλαμήδ\expo{ους} ἐπινοίᾳ {\policeFr ÷}}{« Parce que plus haut (vv. 769-784) il avait eu recours au stratagème des tablettes tiré du Palamède (et qu'Euripide a honte de la froideur de sa pièce). »}
\commScholie{Explique le jeu de mots sur ψυχρός (« froid, sans succès théâtral »).}
\entreeScholie{273}{850}{τὴν καινὴν Ἑλένην μιμήσομαι.}{Interlinéaire (supra lineam au-dessus de καινὴν, f. 172v, v. 850)}{ὑπόγυον δεδιδαγμένην.}{ὑπόγυ\expo{ον} δεδιδαγμέν(ην) {\policeFr ÷}}{« “La nouvelle Hélène” : celle qui vient tout récemment d'être représentée (l'année précédente, aux Grandes Dionysies de 412 av. J.-C., avec l'Andromède, cf. Σ 1012 et Σ 1060). »}
\commScholie{Quatrième repère chronologique exact de la pièce.}
\entreeScholie{274}{852}{τί αὖ σὺ κυρκανᾷς; τί κοικύλλεις ἔχων;}{Marge gauche (extérieure, f. 172v, en face du v. 852)}{τὸ κοικύλλειν, φασὶν, ἐπὶ τοῦ περιβλέπεσθαι δύναται καὶ ἐνθάδε. βέλτιον δὲ ἐπὶ τοῦ κακοτεχνεῖν. τὸ δὲ ἔχων παρέλκει Ἀττικῷ ἔθει.}{τ(ὸ) κοικύλλ(ειν), φασ(ὶν), ἐπὶ τ\expo{οῦ} περιβλέπεσθ(αι) δύνατ(αι) \textbf{ς} ἐνθάδ(ε) · βέλτι\expo{ον} δ(ὲ) ἐπὶ τ\expo{οῦ} κακοτεχν(εῖν) · τ(ὸ) δ(ὲ) ἔχ\expo{ων} παρέλκ(ει) Ἀττικῷ ἔθ(ει) {\policeFr ÷}}{« Le verbe κοικύλλειν, dit-on, a ici encore le sens de “regarder autour de soi en badaud en roulant les yeux” (περιβλέπεσθαι, cf. Hésychius κ 3232). Mais il vaut mieux l'entendre ici au sens de “manigancer quelque mauvais tour”. Quant au participe ἔχων, il est pléonastique selon l'idiome attique (marquant la durée insistante : “qu'as-tu à toujours lorgner ?”). »}
\commScholie{Cité par Louis Hauseux (vv 642-fin-bis.tex, l. 2131 n. 2) avec Hésychius κ 3232 et la Souda κ 2553.}
\entreeScholie{275}{855}{Νείλου μὲν αἵδε καλλιπάρθενοι ῥοαί,}{Marge gauche (extérieure, f. 172v, en face du v. 855)}{ἐξ Ἑλένης Εὐριπίδου πολλὰ τούτων.}{ἐξ Ἑλέν\expo{ης} Εὐριπίδ\expo{ου} πολλὰ τούτ\expo{ων} {\policeFr ÷}}{« Un grand nombre de ces vers sont repris textuellement de l'Hélène d'Euripide (vv. 1-3 sqq.). »}
\commScholie{Signale le début de la grande parodie de l'Hélène d'Euripide (dont le papyrus Ψ\expo{3} = POxy 3840 conserve les vv. 856-857).}
\entreeScholie{276}{857}{λευκῆς νοτίζει μελανοσυρμαῖον λεών.}{Marge gauche (extérieure, f. 172v, en face du v. 857)}{μέλανα καὶ συρμαϊζόμενον. ἡ δὲ συρμαία κρίθινόν ἐστι πότημα.}{μέλανα \textbf{ς} συρμαϊζόμεν\expo{ον} · ἡ δ(ὲ) συρμαία κρίθινόν ἐστ(ι) πότημα {\policeFr ÷}}{« Le peuple égyptien “à la peau noire et qui se purge à la syrmaïa” (μελανοσυρμαῖον, substitué par surprise comique à τὰς ἀρούρας d'Euripide). La συρμαία est une boisson purgative (à base d'orge, de radis et d'eau salée, cf. Hérodote II, 77 et 125). »}
\commScholie{Commenté par Louis Hauseux (vv 642-fin-bis.tex, ll. 2251-2265) avec la Souda σ 1667.}
\entreeScholie{277}{861}{ἐμοὶ δὲ γῆ μὲν πατρὶς οὐκ ἀνώνυμος, / Σπάρτη, πατὴρ δὲ Τυνδάρεως. — σοί γ᾽, ὦλεθρε, / πατὴρ ἐκεῖνός ἐστι; Φρυνώνδας μὲν οὖν.}{Marge gauche (extérieure, f. 172v, en face du v. 861)}{ἐπὶ πονηρίᾳ διαβάλλεται. τινὲς δὲ καὶ πατέρα αὐτῷ φασὶν ἀληθῶς εἶναι τοῦτον.}{ἐπὶ πονηρίᾳ διαβάλλετ(αι) · τιν\expo{ὲς} δ(ὲ) \textbf{ς} πατέρα αὐτῷ φασ(ὶν) ἀληθ\expo{ῶς} εἶν(αι) τοῦτ\expo{ον} {\policeFr ÷}}{« Phrynôndas était un fourbe proverbial décrié pour sa scélératesse ; certains prétendent même qu'il était réellement le père de Mnésiloque (pure Autoschediasma de biographe). »}
\commScholie{Repris par la Souda φ 765 : Phrynôndas est le type proverbial du coquin fieffé à Athènes (cf. Platon, Protagoras 327d, Isocrate, Contre Callimachos 57).}
\entreeScholie{278}{863}{Ἑλένη δ᾽ ἐκλήθην. — αὖθις αὖ γίγνει γυνή, / πρὶν τῆς ἑτέρας δοῦναι γυναικίσεως δίκην;}{Interlinéaire (supra lineam, f. 172v, v. 863)}{τῆς γυναικείας μιμήσεως.}{τ\expo{ῆς} γυναικεί\expo{ας} μιμήσ(εως) {\policeFr ÷}}{« De ton imitation de femme (de ton travestissement en Hélène). »}
\commScholie{Glose interlinéaire de R\expo{2} expliquant le nom d'action comique γυναικίσεως (« travestissement en femme, imitation féminine », dérivé de γυναικίζω, v. 268) par τῆς γυναικείας μιμήσεως ; reprise par la Souda γ 496.}
\entreeScholie{279}{868}{τί οὖν ἔτι ζῶ; — τῶν κοράκων πονηρίᾳ.}{Marge gauche (extérieure, f. 172v, en face du v. 868)}{ὅτι πονηροὶ οἱ κόρακες, καὶ ὅτι μέχρι νῦν οὐκ ἐσπάραξάν σε.}{ὅτι πονηρ(οὶ) οἱ κόρακ\expo{ες}, \textbf{ς} ὅτι μέχρι νῦν οὐκ ἐσπάραξάν σε {\policeFr ÷}}{« “Par la fainéantise (la lâcheté) des corbeaux !” : parce que les corbeaux sont de mauvais ouvriers et ne t'ont pas encore déchiqueté la carcasse ! »}
\commScholie{Crucial pour l'attribution des répliques (Louis Hauseux, Bibi, v. 868, vv 642-fin-bis.tex, l. 2350) : alors que R\expo{1} continue tout le v. 868 à Mnésiloque, la 2e personne οὐκ ἐσπάραξάν σε de {\policeGr Σ}\expo{R} (et de la Souda τ 648) prouve que la réplique mordante τῶν κοράκων πονηρίᾳ est lancée en interruption par Kritylla (Γυνή, adopté par Ellebodius et Bibi) !}
\entreeScholie{280}{869}{ἀλλ᾽ ὥσπερ αἰκάλλει τι καρδίαν ἐμήν·}{Interlinéaire (supra lineam au-dessus d'αἰκάλλει, f. 172v, v. 869)}{σαίνει τὸ ᾆσμα.}{σαίν(ει) τ(ὸ) ᾆσμα {\policeFr ÷}}{« αἰκάλλει : flatte doucement, caresse (σαίνει) mon cœur d'un espoir. »}
\commScholie{Glose interlinéaire de R\expo{2} expliquant la métaphore tragique αἰκάλλει (« caresse, flatte comme un chien qui remue la queue », glosé par σαίνει) reprise de l'Andromaque d'Euripide (v. 630).}
\entreeScholie{281}{870}{μὴ ψεῦσον, ὦ Ζεῦ, τῆς ἐπιούσης ἐλπίδος.}{Marge inférieure du f. 172v (grande scholie grammaticale de 4 lignes en face du v. 870)}{σύνηθες τὸ σχῆμα. Μένανδρος «οἶσθ’ ὃ ποίησον.» Σοφοκλῆς Πηλεῖ «μὴ ψεῦσον ὦ Ζεῦ· μή μ’ ἕλῃς ἄνευ δορός.» «οἶσθ’ οὖν ὃ δρᾶσον» Εὐριπίδης Πολυΐδῳ.}{σύνηθ\expo{ες} τ(ὸ) σχῆμα · Μένανδρ\expo{ος} «οἶσθ' ὃ ποίησ\expo{ον}» · Σοφοκλ\expo{ῆς} Πηλεῖ «μὴ ψεῦσ\expo{ον} ὦ Ζεῦ · μή μ' ἕλῃς ἄνευ δορός» · «οἶσθ' οὖν ὃ δρᾶσ\expo{ον}» Εὐριπίδ(ης) Πολυΐδῳ {\policeFr ÷}}{« Ce tour syntaxique (μὴ suivi de l'impératif aoriste ψεῦσον au lieu du subjonctif aoriste μὴ ψεύσῃς) est consacré par l'usage attique. Ainsi Ménandre : “Sais-tu ce que tu dois faire (οἶσθ' ὃ ποίησον) ?” ; Sophocle dans son Pélée (fr. 493 Radt) : “Ne me trompe pas (μὴ ψεῦσον), ô Zeus ! Ne me prends pas sans lance !” ; et Euripide dans son Polyïdos (fr. 645 Kannicht) : “Sais-tu donc ce que tu dois faire (οἶσθ' οὖν ὃ δρᾶσον) ?” »}
\commScholie{Admirable scholie grammaticale alexandrine (commentée par Louis Hauseux, vv 642-fin-bis.tex, ll. 2356-2398, en regard de Grégoire de Corinthe p. 14-15 et de la Souda ψ 62) qui conserve le vers de Sophocle (Pélée, fr. 493 R.) parodié ici par Aristophane !}
\entreeScholie{282}{871}{τίς τῶνδ᾽ ἐρυμνῶν δωμάτων ἔχει κράτος,}{Marge intérieure droite (dernière ligne du f. 172v, v. 871)}{ὁ Εὐριπίδης ἀναλαμβάνει τὸ πρόσωπον τοῦ Μενελάου καὶ ὑποκρίνεται.}{ὁ Εὐριπίδ(ης) ἀναλαμβάν(ει) τ(ὸ) πρόσωπ\expo{ον} τ\expo{οῦ} Μενελά\expo{ου} \textbf{ς} ὑποκρίνετ(αι) {\policeFr ÷}}{« Ici Euripide entre en scène en endossant le rôle de Ménélas (naufragé en haillons, comme dans l'Hélène v. 386) et joue la scène tragique. »}
\commScholie{Rubrique dramaturgique essentielle marquant l'entrée d'Euripide en Ménélas.}
\clearpage
\renewcommand{\currHeadLE}{ΑΡΙΣΤΟΦΑΝΟΥΣ ΘΕΣΜΟΦΟΡΙΑΖΟΥΣΑΙ · CODEX RAVENNAS 429 (\R) · FOLIO 22/30}
\renewcommand{\currHeadRE}{f. 173r (Classense 175r) · vv. 876–919 · Collation diplomatique}
\renewcommand{\currHeadLO}{ΑΡΙΣΤΟΦΑΝΟΥΣ ΘΕΣΜΟΦΟΡΙΑΖΟΥΣΑΙ · CODEX RAVENNAS 429 (\R) · FOLIO 22/30}
\renewcommand{\currHeadRO}{f. 173r (Classense 175r) · vv. 876–919 · Collation diplomatique}
\noindent\begin{adjustbox}{max totalheight=30.4cm, max width=\textwidth, keepaspectratio}\begin{minipage}{\textwidth}
\bandeauRav{f. 173r}{Classense 175r}{876–919}{Folio 22/30 · Recto (page de droite) · 11 scholies \& gloses ({\policeGr Σ} 283–293)}
{\fontsize{8.8pt}{10.2pt}\selectfont
\ligneRav{876}{\auDessus{γυ}{ν}}{ἐπεὶ τέθνηκε Πρωτέας ἔτη δέκα·}{}
\ligneRav{}{\auDessus{εὐ}{ρ}}{ποίαν δὲ χώραν εἰσεκέλσαμεν σκάφει;}{}
\ligneRav{}{μ\auDessus{ν}{Η}}{Αἴγυπτον. : ὦ δύστηνος οἷ πεπλώκαμεν·}{}
\ligneRav{879}{μ\auDessus{ν}{Η}}{πείθει τι \difference{τῷ κακῷ} κακῶς ἀπολουμένῳ}{\textbf{879} τῷ R : τούτῳ τῷ Grynaeus : \lecon{τῷ κακῷ} Scaliger}
\ligneRav{880}{}{ληροῦντι λῆρον; Θεσμοφόριον τουτογί·}{\textbf{880} \lecon{Θεσμοφόριον} R, cf.278}
\ligneRav{881}{\auDessus{εὐ}{ρ}}{αὐτὸς δὲ Πρωτεὺς ἔνδον ἔστ᾽ \difference{ἢ 'ξώπιος};}{\textbf{881} ἠξώπιος R, id. v.884 : \lecon{ἢ 'ξώπιος;} Biset}
\ligneRav{}{\auDessus{γυ}{ν}}{οὐκ ἔσθ᾽ ὅπως οὐ ναυτιᾷς ἔτ᾽ ὦ ξένε,}{}
\ligneRav{883}{}{\difference{ὅστις γ᾽} ἀκούσας ὅτι τέθνηκε Πρωτέας}{\textbf{883} ὅστις R : \lecon{ὅστις γ᾽} Scaliger}
\ligneRav{}{}{ἔπειτ᾽ ἐρωτᾷς `ἔνδον ἔστ᾽ ἢ 'ξώπιος;'}{}
\ligneRav{885}{\auDessus{εὐ}{ρ}}{αἰαῖ \difference{τέθνηκε}. ποῦ δ᾽ ἐτυμβεύθη τάφῳ;}{\textbf{885} τέθνηκεν R : \lecon{τέθνηκε} Bentley}
\ligneRav{}{μ\auDessus{ν}{Η}}{τόδ᾽ ἐστὶν αὐτοῦ σῆμ᾽, ἐφ᾽ ᾧ καθήμεθα·}{}
\ligneRav{887}{\auDessus{γυ}{ν}}{κακῶς τ᾽ ἄρ᾽ ἐξόλοιο κἀξολεῖ γέ τοι,}{\textbf{887} \lecon{κακῶς τ᾽ ἄρ᾽} R : κακῶς γ᾽ ἄρ᾽ Denniston (\italique{CR} XLIV(1930), p.215)}
\ligneRav{}{}{ὅστις γε τολμᾷς σῆμα τὸν βωμὸν καλεῖν·}{}
\ligneRav{889}{\auDessus{εὐ}{ρ}}{\difference{τί δὴ σὺ} θάσσεις τάσδε τυμβήρεις ἕδρας}{\textbf{889} τί δὲ σὺ R : τί δαὶ σὺ Faber : \lecon{τί δὴ σὺ} Scaliger}
\ligneRav{890}{—}{φάρει καλυπτὸς ὦ ξένη; : βιάζομαι}{}
\ligneRav{891}{}{γάμοισι Πρωτέως παιδὶ \difference{συμμεῖξαι} λέχος·}{\textbf{891} συμμῖξαι R : \lecon{συμμεῖξαι} van Leeuwen}
\ligneRav{892}{\auDessus{γυ}{ν}}{τί ὦ κακόδαιμον ἐξαπατᾷς \difference{αὖ τὸν} ξένον;}{\textbf{892} αὐτὸν (ex αὖτὸν) \Rpc : \lecon{αὖ τὸν} Casaubon}
\ligneRav{}{}{οὗτος πανουργῶν δεῦρ᾽ ἀνῆλθεν ὦ ξένε}{}
\ligneRav{}{}{ὡς τὰς γυναῖκας ἐπὶ κλοπῇ τοῦ χρυσίου·}{}
\ligneRav{895}{μ\auDessus{ν}{Η}}{\difference{κόνυζα} τοὐμὸν σῶμα βάλλουσα ψόγῳ·}{\textbf{895} βαϋζα R : βάϋζε Zanetti : \lecon{κόνυζα} Bibi}
\ligneRav{}{\auDessus{εὐ}{ρ}}{ξένη τίς ἡ γραῦς ἡ κακορροθοῦσά σε;}{}
\ligneRav{}{μ\auDessus{ν}{Η}}{αὕτη Θεονόη Πρωτέως. : μὰ τὼ θεὼ}{}
\ligneRav{}{}{εἰ μὴ Κρίτυλλά γ᾽ Ἀντιθέου Γαργηττόθεν·}{}
\ligneRav{}{—}{σὺ δ᾽ εἶ πανοῦργος. : ὁπόσα τοι βούλει λέγε·}{}
\ligneRav{900}{}{οὐ γὰρ γαμοῦμαι σῷ κασιγνήτῳ ποτὲ}{}
\ligneRav{901}{}{προδοῦσα \difference{Μενέλεων} τὸν ἐμὸν ἐν Τροίᾳ πόσιν·}{\textbf{901} Μενέλαον R : \lecon{Μενέλεων} Scaliger}
\ligneRav{}{\auDessus{εὐ}{ρ}}{γύναι τί εἶπας; στρέψον ἀνταυγεῖς κόρας·}{}
\ligneRav{}{μ\auDessus{ν}{Η}}{αἰσχύνομαί σε τὰς γνάθους ὑβρισμένη·}{}
\ligneRav{}{\auDessus{εὐ}{ρ}}{τουτὶ τί ἔστιν; ἀφασία τίς τοί μ᾽ ἔχει·}{}
\ligneRav{905}{}{ὦ θεοὶ τίν᾽ ὄψιν εἰσορῶ; τίς εἶ γύναι;}{}
\ligneRav{906}{μ\auDessus{ν}{Η}}{σὺ δ᾽ εἶ τίς; \difference{αὑτὸς} γὰρ σὲ κἄμ᾽ ἔχει λόγος·}{\textbf{906} αὐτὸς R : \lecon{αὑτὸς} Boissonade}
\ligneRav{}{\auDessus{εὐ}{ρ}}{Ἑλληνὶς εἶ τις ἢ 'πιχωρία γυνή;}{}
\ligneRav{}{μ\auDessus{ν}{Η}}{Ἑλληνίς. ἀλλὰ καὶ τὸ σὸν θέλω μαθεῖν·}{}
\ligneRav{909}{\auDessus{εὐ}{ρ}}{\difference{Ἑλένῃ σ᾽} ὁμοίαν δὴ \difference{μάλιστ᾽ εἶδον} γύναι·}{\textbf{909} Ἐλένης R : \lecon{Ἑλένῃ σ᾽} Eur. \italique{Hel.}563 \espApp μάλισθ᾽ ἴδον R : \lecon{μάλιστ᾽ εἶδον} \italique{Hel}.563}
\ligneRav{910}{μ\auDessus{ν}{Η}}{ἐγὼ δὲ \difference{Μενελέῳ σ᾽} ὅσα γ᾽ ἐκ τῶν \difference{λαιφίων}·}{\textbf{910} Μενέλαον R : Μενέλεων Suid. : \lecon{Μενελέῳ σ᾽} Thiersch \espApp ἀφύων R : ἰφύων Σ, Suid. : ἀμφίων Grégoire : \lecon{λαιφίων} Bibi}
\ligneRav{}{\auDessus{εὐ}{ρ}}{ἔγνως ἄρ᾽ ὀρθῶς ἄνδρα δυστυχέστατον·}{}
\ligneRav{}{μ\auDessus{ν}{Η}}{ὦ χρόνιος ἐλθὼν σῆς δάμαρτος ἐσχάρας}{}
\ligneRav{914}{}{\hspace{1.1em}λαβέ με λαβέ με πόσι, \difference{περίβαλε} δὲ χέρας·}{\textbf{914} περίβαλλε R : \lecon{περίβαλε} Biset}
\ligneRav{915}{}{\hspace{1.1em}φέρε σὲ κύσω. ἄπαγέ μ᾽ ἄπαγ᾽ ἄπαγέ με}{}
\ligneRav{}{—}{λαβὼν ταχὺ πάνυ. : κλαύσετ᾽ ἄρα νὴ τὼ θεὼ}{}
\ligneRav{}{}{ὅστις σ᾽ ἀπάξει τυπτόμενος τῇ λαμπάδι·}{}
\ligneRav{}{\auDessus{εὐ}{ρ}}{σὺ τὴν ἐμὴν γυναῖκα κωλύεις ἐμέ,}{}
\ligneRav{}{}{τὴν Τυνδάρειον παῖδ᾽, ἐπὶ Σπάρτην ἄγειν;}{}
}
\vspace{3pt}{\color{ruleGold}\hrule height 0.6pt}\vspace{3pt}
{\fontsize{8.0pt}{9.4pt}\selectfont
\noindent{\bfseries\scshape\color{rubricRed}Loci paralleli et testimonia (f. 173r, vv. 876–919) :} \textbf{878} Suidas οι 1 ὦ δύστηνος, οἷ πεπλώκαμεν; \espTest \textbf{882} (ναυττιᾶν = ναυσιᾶν) Moeris ν 9 \espTest \textbf{902} (στρέψον ἀνταυγεῖς κόρας) Suidas α 2617 \espTest (ἀνταυγεῖς κόρας) Photius α 2059 ἀνταυγεῖς κόρας \espTest \textbf{904} \Antiatt 83.9 ; cf. Hesychius α 8572 ἀφασία \espTest \textbf{910} Suidas ι 778 (sed Μενέλεων pro Μενέλαον, ὄσον pro ὅσα γ᾽ et ἰφύων pro ἀφύων) \espTest \textbf{912} (ἐσχάρας) Photius ε 2044 ἐσχάρας {\policeFr ⊂} Hesychius ε 6447 ἐσχάραι\par\vspace{2.5pt}
{\color{ruleGold}\hrule height 0.35pt}\vspace{2.5pt}
\noindent{\bfseries\scshape\color{rubricRed}Suite (en commentaire) de l'apparat critique et paléographique :} \textbf{Contexte :} Suite et sommet de la \textbf{parodie d'\italique{Hélène}} : stichomythie de reconnaissance entre Euripide-Ménélas et Mnésiloque-Hélène sous le regard incrédule et furieux de la gardienne (\textbf{Kritylla}, vv. 876–919). --- \textbf{Deux émendations paléographiques majeures de Bibi (vv. 895 et 910) :} 1° Au \textbf{v.~895}, là où \R{} porte βαϋζα (βάϋζε Zanetti), Louis Hauseux (Bibi) restitue \lecon{κόνυζα} (« l'inule visqueuse », plante de chasteté rituelle des Thesmophories : dans la minuscule de \R{}, le κ initial a été lu β bilobé et -ον- lié a été lu -α-) ; 2° au \textbf{v.~910}, là où \R{} porte ἐκ τῶν ἀφύων (ἰφύων {\policeGr Σ}\expo{R}, Souda), Bibi restitue \lecon{ἐκ τῶν λαιφίων} (« à en juger par tes haillons », ap. Hésychius λ 171 λαίφια· ῥάκη : en onciale ΕΚΤΩΝΛΑΙΦΙΩΝ, le Λ triangulaire devant Α est tombé par haplographie, puis l'iotacisme a produit ἀφύων / ἰφύων). --- \textbf{v. 895 :} βαϋζα R (βάϋζε Zanetti) : \lecon{κόνυζα} Bibi (confusion minuscule κ → β). --- \textbf{v. 910 :} Μενέλαον R (Μενέλεων Souda) : \lecon{Μενελέῳ σ'} Thiersch \espApp \textbf{ἀφύων R} (ἰφύων {\policeGr Σ}, Souda ; ἀμφίων Grégoire) : \lecon{λαιφίων} Bibi ap. Hésychius λ 171 (chute de {\policeGr Λ} devant Α en onciale).\par}
\end{minipage}\end{adjustbox}
\newpage
\renewcommand{\currHeadLE}{SCHOLIES MARGINALES ({\policeGr Σ}\expo{R}) \& GLOSES INTERLINÉAIRES · FOLIO 22/30}
\renewcommand{\currHeadRE}{f. 173r (Classense 175r) · vv. 876–919 · {\policeGr Σ} 283 à {\policeGr Σ} 293}
\renewcommand{\currHeadLO}{SCHOLIES MARGINALES ({\policeGr Σ}\expo{R}) \& GLOSES INTERLINÉAIRES · FOLIO 22/30}
\renewcommand{\currHeadRO}{f. 173r (Classense 175r) · vv. 876–919 · {\policeGr Σ} 283 à {\policeGr Σ} 293}
\bandeauScholies{Folio 22/30 — f. 173r}{Classense 175r}{vv. 876–919}{Recto (page de droite) · 11 entrées ({\policeGr Σ} 283 à {\policeGr Σ} 293)}
{\fontsize{8.8pt}{10.5pt}\selectfont\noindent\textbf{\color{rubricRed}Disposition topographique des marges sur le f. 173r :} Marge extérieure à droite ; marge intérieure à gauche. — \textit{Fac-similé de référence :} \texttt{folios\_individuels\_ff162v-177r/22\_f173r\_classense175r\_vv0876-0919.jpg}.\par\vspace{3pt}}
\entreeScholie{283}{876}{Πρωτέως τάδ᾽ ἐστὶ μέλαθρα. — ποίου Πρωτέως;}{Marge supérieure et marge droite (extérieure, f. 173r, en haut du recto)}{Πρωτέας οὕτω καλούμενος Ἀθηναῖος, ὃς ἐτεθνήκει πρόπαλαι.}{Πρωτέ\expo{ας} οὕτω καλούμεν\expo{ος} Ἀθηναῖ\expo{ος}, ὃς ἐτεθνήκ(ει) πρόπαλαι {\policeFr ÷}}{« Il y avait un citoyen athénien nommé Protéas (du dème d'Aixônè, stratège célèbre sous Périclès en 433/431 av. J.-C., cf. Thucydide I, 45 et II, 23), qui était mort depuis une dizaine d'années (d'où la réplique de Kritylla au v. 883 : τέθνηκε Πρωτέας ἔτη δέκα !). »}
\commScholie{Scholie historique capitale (reprise par la Souda π 2948) qui donne la clé du quiproquo entre le roi mythologique d'Égypte Protée (Πρωτεύς) invoqué par Euripide-Ménélas et le stratège athénien défunt Protéas d'Aixônè (Πρωτέας) auquel pense Kritylla !}
\entreeScholie{284}{878}{πεπλώκαμεν.}{Marge droite (extérieure, f. 173r, en face du v. 878)}{διαβάλλει τὴν Αἴγυπτον ὡς ἀλίμενον καὶ ὡς κίνδυνον ἐμποιοῦσαν ταῖς ναυσίν.}{διαβάλλ(ει) τ(ὴν) Αἴγυπτ\expo{ον} ὡς ἀλίμεν\expo{ον} \textbf{ς} ὡς κίνδυν\expo{ον} ἐμποιοῦσ(αν) τ\expo{αῖς} ναυσίν {\policeFr ÷}}{« Il décrie la côte d'Égypte comme dépourvue de ports abrités (ἀλίμενος) et dangereuse pour les navires. »}
\commScholie{Lieu commun géographique depuis l'Odyssée (IV, 355 sqq.).}
\entreeScholie{285}{882}{οὐκ ἔσθ᾽ ὅπως οὐ ναυτιᾷς ἔτ᾽ ὦ ξένε,}{Interlinéaire (supra lineam, f. 173r, v. 882)}{ἐκ τῆς θαλάσσης.}{ἐκ τ\expo{ῆς} θαλάσσ\expo{ης} {\policeFr ÷}}{« Au sortir de la mer (débarquement après le naufrage). »}
\commScholie{Glose interlinéaire de R\expo{2} expliquant l'origine du verbe ναυτιᾷς (« tu as encore le mal de mer », v. 882) : à la suite de la tempête et de la sortie « hors de la mer » (ἐκ τῆς θαλάσσης).}
\entreeScholie{286}{897}{αὐτὴ Θεονόη Πρωτέως.}{Marge droite (extérieure, f. 173r, en face du v. 897)}{ἣν εἶπεν Ὅμηρος [Od. Θ, 366] Εἰδοθέαν.}{ἣν εἶπ(εν) Ὅμηρ\expo{ος} Εἰδοθέ(αν) {\policeFr ÷}}{« Celle (la fille prophétesse de Protée) qu'Homère, dans l'Odyssée (IV, 366), a appelée Idothée (Εἰδοθέα), et qu'Euripide dans son Hélène a rebaptisée Théonoé. »}
\commScholie{Remarque mythographique d'Aristarque sur l'innovation onomastique d'Euripide dans l'Hélène (cf. Hélène vv. 11-14).}
\entreeScholie{287}{898}{μὰ τὼ θεὼ, / εἰ μὴ Κριτύλλα γ᾽ Ἀντιθέου Γαργηττόθεν·}{Marge droite (extérieure, f. 173r, en face du v. 898)}{ἀπὸ δήμου τῆς Αἰγηίδος φυλῆς, ἔνθα τὸ Εὐρυσθέως ἀπόκειται σῶμα ἐν τῷ ἀκρωτηρίῳ. ἐπίρρημα δὲ τὸ Γαργηττόθεν.}{ἀπ(ὸ) δήμ\expo{ου} τ\expo{ῆς} Αἰγηίδ\expo{ος} φυλ\expo{ῆς}, ἔνθ(α) τ(ὸ) Εὐρυσθέ\expo{ως} ἀπόκειτ(αι) σῶμα ἐν τῷ ἀκρωτηρίῳ · ἐπίρρημα δ(ὲ) τ(ὸ) Γαργηττόθ(εν) {\policeFr ÷}}{« Γαργηττόθεν (“du dème de Gargettos”) : dème appartenant à la tribu Égéide (Αἰγηὶς φυλή), où repose le corps d'Eurysthée (cf. Euripide, Héraclides v. 1031). Et Γαργηττόθεν est un adverbe de lieu d'origine en -θεν. »}
\commScholie{Cité par Louis Hauseux (vv 642-fin-bis.tex, l. 2706) et repris par la Souda γ 60.}
\entreeScholie{288}{902}{στρέψον ἀνταυγεῖς κόρας.}{Interlinéaire (supra lineam au-dessus du v. 902, f. 173r)}{τὰς κόρας σου ἀντιπροσώπους μοι στρέψον.}{τ\expo{ὰς} κόρ\expo{ας} σ\expo{ου} ἀντιπροσώπ\expo{ους} μοι στρέψ\expo{ον} {\policeFr ÷}}{« Tourne les prunelles de tes yeux en face des miens (regarde-moi en face). »}
\commScholie{Paraphrase de l'expression tragique ἀνταυγεῖς κόρας.}
\entreeScholie{289}{903}{αἰσχύνομαί σε τὰς γνάθους ὑβρισμένη.}{Marge droite (extérieure, f. 173r, en face du v. 903)}{ἐπειδὴ ξυρηθεὶς ἦν.}{ἐπειδὴ ξυρηθ(εὶς) ἦν {\policeFr ÷}}{« (“J'ai honte devant toi d'avoir les joues outragées”) : parce qu'il avait été rasé de près au début de la pièce (alors qu'Hélène chez Euripide parlait des griffures de deuil sur ses joues). »}
\commScholie{Cité par Louis Hauseux (vv 642-fin-bis.tex, l. 2730) : le vers d'Euripide est détourné comiquement vers le rasage humiliant subi par Mnésiloque aux vv. 215-235.}
\entreeScholie{290}{910}{ἐγὼ δὲ Μενελέῳ σ᾽, ὅσα γ᾽ ἐκ τῶν λαιφίων [ἀφύων R, ἰφύων Σ, Suid.]·}{Marge droite (extérieure, f. 173r, en face du v. 910)}{δέον εἰπεῖν ἐκ τῶν ὄψεων, εἶπεν ἐκ τῶν ἰφύων. ἴφυον δέ ἐστιν εἶδος ἀγρίου λαχάνου, ὅτι Εὐριπίδης λαχανοπώλιδος Κλειτοῦς υἱὸς ἦν δηλονότι.}{δέον εἰπ(εῖν) ἐκ τ\expo{ῶν} ὄψε\expo{ων}, εἶπ(εν) ἐκ τ\expo{ῶν} ἰφύ\expo{ων} · ἴφυ\expo{ον} δέ ἐστ(ιν) εἶδ\expo{ος} ἀγρί\expo{ου} λαχάν\expo{ου}, ὅτι Εὐριπίδ(ης) λαχανοπώλιδ\expo{ος} Κλειτοῦς υἱ(ὸς) ἦν δηλονότι {\policeFr ÷}}{« Alors qu'on attendait “autant qu'on en peut juger à tes traits (ἐκ τῶν ὄψεων / ἐκ τῶν ἀμφίων, à tes haillons)”, il a dit par surprise comique “à tes feuilles de lavande / d'asphodèle (ἐκ τῶν ἰφύων)”. L'ἴφυον est une espèce d'herbe potagère sauvage, parce qu'Euripide était, comme chacun sait, le fils de la marchande de légumes Clitô ! »}
\commScholie{Cas passionnant analysé par Louis Hauseux (Bibi, v. 910, vv 642-fin-bis.tex, ll. 2779-2815) : R\expo{1} porte la faute amétrique ἐκ τῶν ἀφύων (« aux anchois »), tandis que {\policeGr Σ}\expo{R} et la Souda ι 778 lisent ἐκ τῶν ἰφύων (« aux herbes potagères ») ; L. Hauseux examine la conjecture d'Henri Grégoire ἐκ τῶν λαιφίων (« à tes haillons », cf. Hélène v. 421).}
\entreeScholie{291}{912}{ὦ χρόνιος ἐλθὼν σῆς δάμαρτος ἐσχάρας [ἐς χέρας Eur. Hel. 566]·}{Marge droite (extérieure, f. 173r, en face des vv. 912–913)}{καὶ ταῦτα ἐκ τῆς Ἑλένης, καὶ καθόλου τὰ πλεῖστα.}{\textbf{ς} ταῦτ(α) ἐκ τ\expo{ῆς} Ἑλέν\expo{ης}, \textbf{ς} καθόλ\expo{ου} τὰ πλεῖστ(α) {\policeFr ÷}}{« Ces vers aussi sont tirés textuellement de l'Hélène d'Euripide (vv. 566 sqq.), et d'une manière générale la majeure partie de cette scène. »}
\commScholie{Signale le centon continu de la scène de reconnaissance de l'Hélène (vv. 558-569).}
\entreeScholie{292}{915}{λαβέ με λαβέ με πόσι, περίβαλε δὲ χέρας.}{Marge droite (extérieure, f. 173r, en face du v. 915)}{ταῦτα λέγει ἀκολουθῶν αὐτῷ καὶ βουλόμενος ἐκφυγεῖν.}{ταῦτ(α) λέγ(ει) ἀκολουθ\expo{ῶν} αὐτῷ \textbf{ς} βουλόμεν\expo{ος} ἐκφυγ(εῖν) {\policeFr ÷}}{« Il dit ces paroles tout en s'élançant derrière Euripide et en cherchant à s'enfuir du sanctuaire avec lui. »}
\commScholie{Note de jeu scénique (παρεπιγραφή commentée) de R\expo{2} précisant que Mnésiloque prononce le vers lyrique d'Hélène (λαβέ με λαβέ με πόσι) tout en s'accrochant à Euripide-Ménélas pour tenter de s'enfuir avec lui loin de Critylla.}
\entreeScholie{293}{917}{κλαύσετ᾽ ἄρα νὴ τὼ θεὼ / ὅστις σ᾽ ἀπάξει τυπτόμενος τῇ λαμπάδι.}{Marge inférieure et marge droite (bas du f. 173r, grande scholie alexandrine de Callistrate en face des vv. 916–917)}{τὸ χ παράκειται τῷ στίχῳ. καί φησι Καλλίστρατος [...] οὐ τῇ λαμπάδι εἶπεν, ἀλλὰ νὴ τῇ δᾳδὶ, ὡς δῆλον εἶναι ὅτι ἑκατέρως ἔλεγον. ἀνάξιος δὲ ἂν εἴη ὁ Καλλίστρατος πιστεύεσθαι ὡς ἂν περὶ διδασκάλου τοῦ Ἀριστοφάνους λέγων.}{τ(ὸ) χ παράκειτ(αι) τῷ στίχῳ · καί φησ(ι) Καλλίστρατ\expo{ος} [...] οὐ τῇ λαμπάδι εἶπ(εν), ἀλλὰ νὴ τῇ δᾳδὶ, ὡς δῆλ\expo{ον} εἶν(αι) ὅτι ἑκατέρ\expo{ως} ἔλεγ\expo{ον} · ἀνάξι\expo{ος} δ(ὲ) ἂν εἴη ὁ Καλλίστρατ\expo{ος} πιστεύεσθ(αι) ὡς ἂν π(ερὶ) διδασκάλ\expo{ου} τ\expo{οῦ} Ἀριστοφάν\expo{ους} λέγ\expo{ων} {\policeFr ÷}}{« Le signe critique chi (τὸ χ) est apposé en marge de ce vers. Callistrate (élève d' Aristophane de Byzance, IIe s. av. J.-C.) dit que [Callistratos le metteur en scène ?] n'a pas dit “avec la torche (τῇ λαμπάδι)”, mais “avec le flambeau de résine (τῇ δᾳδί)”, de sorte qu'il est clair qu'on employait indifféremment les deux mots. Mais Callistrate ne mérite sans doute aucune créance lorsqu'il parle ainsi du metteur en scène (διδάσκαλος) d'Aristophane. »}
\commScholie{Scholie alexandrine fameuse (commentée par L. Hauseux, vv 642-fin-bis.tex, l. 2868) : elle atteste la présence du signe critique χ (chi) dans l'édition alexandrine et cite le commentaire de Callistrate (confusion antique entre le grammairien Callistrate et Callistratos, le διδάσκαλος d'Aristophane).}
\clearpage
\renewcommand{\currHeadLE}{ΑΡΙΣΤΟΦΑΝΟΥΣ ΘΕΣΜΟΦΟΡΙΑΖΟΥΣΑΙ · CODEX RAVENNAS 429 (\R) · FOLIO 23/30}
\renewcommand{\currHeadRE}{f. 173v (Classense 175v) · vv. 920–958 · Collation diplomatique}
\renewcommand{\currHeadLO}{ΑΡΙΣΤΟΦΑΝΟΥΣ ΘΕΣΜΟΦΟΡΙΑΖΟΥΣΑΙ · CODEX RAVENNAS 429 (\R) · FOLIO 23/30}
\renewcommand{\currHeadRO}{f. 173v (Classense 175v) · vv. 920–958 · Collation diplomatique}
\noindent\begin{adjustbox}{max totalheight=30.4cm, max width=\textwidth, keepaspectratio}\begin{minipage}{\textwidth}
\bandeauRav{f. 173v}{Classense 175v}{920–958}{Folio 23/30 · Verso (page de gauche) · 11 scholies \& gloses ({\policeGr Σ} 294–304)}
{\fontsize{9.3pt}{10.8pt}\selectfont
\ligneRav{920}{\auDessus{γυ}{ν}}{οἴμ᾽ ὡς πανοῦργος καὐτὸς εἶναί μοι δοκεῖς}{}
\ligneRav{}{}{καὶ τοῦδέ τις ξύμβουλος. οὐκ ἐτὸς πάλαι}{}
\ligneRav{}{}{ᾐγυπτιάζετ᾽. ἀλλ᾽ ὅδε μὲν δώσει δίκην·}{}
\ligneRav{}{}{προσέρχεται γὰρ ὁ πρύτανις χὠ τοξότης·}{}
\ligneRav{}{\auDessus{εὐ}{ρ}}{τουτὶ πονηρόν· ἀλλ᾽ ὑπαποκινητέον·}{}
\ligneRav{925}{μ\auDessus{ν}{Η}}{ἐγὼ δ᾽ ὁ κακοδαίμων τί δρῶ; : μέν᾽ ἥσυχος·}{}
\ligneRav{}{}{οὐ γὰρ προδώσω σ᾽ οὐδέποτ᾽, ἤνπερ ἐμπνέω,}{}
\ligneRav{927}{}{ἢν μὴ \difference{προλίπωσ᾽} αἱ μυρίαι με μηχαναί·}{\textbf{927} προλείπωσ᾽ R : \lecon{προλίπωσ᾽} Scaliger}
\ligneRav{928}{μ\auDessus{ν}{Η}}{αὕτη μὲν ἡ μήρινθος οὐδὲν ἔσπασεν·}{\textbf{928} (Γυ.) R : \lecon{(Μνη.)} Scaliger}
\ligneRav{929}{πρυτ/}{ὅδ᾽ \difference{ἔσθ᾽} ὁ πανοῦργος ὃν \difference{ἔλεγ᾽} ἡμῖν Κλεισθένης;}{\textbf{929} ἔστ᾽ R : \lecon{ἔσθ᾽} \Mon \espApp ἔλεγεν R : \lecon{ἔλεγ᾽} \Mon}
\ligneRav{930}{}{οὗτος τί κύπτεις; δῆσον αὐτὸν εἰσάγων}{}
\ligneRav{}{}{ὦ τοξότ᾽ ἐν τῇ σανίδι, κἄπειτ᾽ ἐνθαδὶ}{}
\ligneRav{}{}{στήσας φύλαττε καὶ προσιέναι μηδένα}{}
\ligneRav{}{}{ἔα πρὸς αὐτόν, ἀλλὰ τὴν μάστιγ᾽ ἔχων}{}
\ligneRav{934}{—}{παῖ᾽ ἢν προσίῃ τις. : νὴ Δί᾽ ὡς \difference{νῦν δή γ᾽} ἀνὴρ}{\textbf{934} νῦν δῆτ' R : \lecon{νῦν δή γ᾽} Dobree : νυνδή γ' Cobet}
\ligneRav{935}{}{ὀλίγου μ᾽ ἀφείλετ᾽ αὐτὸν ἱστιορράφος·}{}
\ligneRav{}{μ\auDessus{ν}{Η}}{ὦ πρύτανι πρὸς τῆς δεξιᾶς, ἥνπερ φιλεῖς}{}
\ligneRav{}{}{κοίλην προτείνειν ἀργύριον ἤν τις διδῷ,}{}
\ligneRav{}{}{χάρισαι βραχύ τί μοι καίπερ ἀποθανουμένῳ·}{}
\ligneRav{939}{πρυτ/}{τί σοι \difference{χαρίσωμαι}; : γυμνὸν ἀποδύσαντά με,}{\textbf{939} χαρίσομαι R : \lecon{χαρίσωμαι} Casaubon}
\ligneRav{940}{}{κέλευε πρὸς τῇ σανίδι δεῖν τὸν τοξότην,}{}
\ligneRav{941}{}{ἵνα \difference{μὴ 'ν} κροκωτοῖς καὶ μίτραις γέρων ἀνὴρ}{\textbf{941} μὴν R : \lecon{μὴ 'ν} Küster}
\ligneRav{942}{}{γέλωτα παρέχω \inRas{γε}{} τοῖς κόραξιν ἑστιῶν·}{\textbf{942} \inRas{γε}{} τοῖς \Rdeux}
\ligneRav{943}{πρυτ/}{ἔχοντα ταῦτ᾽ \difference{ἔδοξε} τῇ βουλῇ σε δεῖν,}{\textbf{943} ἔδοξεν R : \lecon{ἔδοξε} Biset}
\ligneRav{944}{}{ἵνα τοῖς παροῦσι δῆλος ᾖς πανοῦργος ὤν·}{\textbf{944} \lecon{παροῦσι} R : παριοῦσι Brunck}
\ligneRav{945}{μ\auDessus{ν}{Η}}{ἰαππαπαιάξ· ὦ κροκώθ᾽ οἷ᾽ εἴργασαι·}{\textbf{945} \lecon{ἰαππαπαιάξ} R, Suid. ε 2807 : ἰατταταιάξ Biset}
\ligneRav{946}{}{κοὐκ \difference{ἔστ᾽} ἔτ᾽ ἐλπὶς οὐδεμία σωτηρίας·}{\textbf{946} ἔστιν R : \lecon{ἔστ᾽} Bentley}
\ligneRav{947}{χορ/}{ἄγε νυν ἡμεῖς \difference{παίσωμεν} ἅπερ νόμος ἐνθάδε \difference{ταῖσι} γυναιξίν,}{\textbf{947} πέσωμεν R : \lecon{παίσωμεν} Zanetti, cf. Σ \espApp ταῖσ R : \lecon{ταῖσι} Daubuz}
\ligneRav{948}{}{παραχώρει ὅταν ὄργια σεμνὰ θεαῖν ἱεραῖς ὥραις ἀνέχωμεν, ἅπερ καὶ}{\textbf{948} \lecon{θεαῖν} R : θεοῖν Cobet}
\ligneRav{}{}{\hspace{1.1em}Παύσων σέβεται καὶ νηστεύει,}{}
\ligneRav{950}{}{\hspace{1.1em}πολλάκις αὐταῖν ἐκ τῶν ὡρῶν}{\textbf{950} \lecon{αὐταῖν} R : αὐτοῖν Cobet}
\ligneRav{}{}{\hspace{1.1em}ἐς τὰς ὥρας ξυνεπευχόμενος}{}
\ligneRav{952}{}{\hspace{1.1em}τοιαῦτα \difference{μέλειν} θάμ᾽ ἑαυτῷ·}{\textbf{952} μέλλειν R : \lecon{μέλειν} Zanetti (\italique{curae esse} Divus) : μολεῖν Casaubon}
\ligneRav{}{χορ/}{\hspace{1.1em}ὅρμα χώρει·}{}
\ligneRav{}{}{\hspace{1.1em}κοῦφα ποσὶν ἄγ᾽ ἐς κύκλον,}{}
\ligneRav{955}{}{\hspace{1.1em}χειρὶ σύναπτε χεῖρα,}{\textbf{955} \lecon{χειρὶ} R : χερὶ Dobree}
\ligneRav{956}{}{\hspace{1.1em}†ῥυθμὸν χορείας ὕπαγε πᾶσα†·}{\textbf{956} \lecon{ὕπαγε πᾶσα} R : πᾶσ' ὕπαγε Austin}
\ligneRav{}{}{\hspace{1.1em}βαῖνε καρπαλίμοιν ποδοῖν·}{}
\ligneRav{}{}{\hspace{1.1em}ἐπισκοπεῖν δὲ πανταχῇ}{}
\ligneRav{}{}{\hspace{1.1em}κυκλοῦσαν ὄμμα χρὴ χοροῦ κατάστασιν·}{}
}
\vspace{3pt}{\color{ruleGold}\hrule height 0.6pt}\vspace{3pt}
{\fontsize{8.0pt}{9.4pt}\selectfont
\noindent{\bfseries\scshape\color{rubricRed}Loci paralleli et testimonia (f. 173v, vv. 920–958) :} \textbf{922} (ᾐγυπτιάζετ᾽) \SynB α 494 αἰγυπτιάζειν = Photius α 518 {\policeFr ⊂} Suidas αι 75 \espTest \textbf{926} (ἤνπερ ἐμπνέω) Suidas η 379 (ἤν pro ἤνπερ) \espTest \textbf{928} Suidas α 4475, μ 978 ; \Eustathe 760.44 ; scholium\expo{RV} \italique{Vespae} 175b \espTest \textbf{931} (σανίδι) Photius σ 68 σανίς, σ 65 σανίδα ; Hesychius σ 169 σανίς ; \AelDion σ 7 \espTest \textbf{935} Suidas ι 691 ὀλίγου ταύτην ἱστιορράφος \espTest (ἱστιορράφος) Pollux VII 161 \espTest \textbf{940} (πρὸς τῇ σανίδι δεῖν) Hesychius δ 926 δῆσαι πρὸς σανίδι· Ἀττικῶς· βασανίσαι ἐν τῷ δεσμωτηρίῳ \espTest \textbf{942} Suidas ε 3215 \espTest (ἑστιῶν) Hesychius ε 6411 ἑστιῶν· ἀριστοποιῶν = Photius ε 2031 = \Syn ε 875 = Suidas ε 3227 ; cf. \Phryn 66.7 \espTest \textbf{945} (ἰαππαπαιάξ) Suidas ε 2807 \espTest \textbf{947} (ὅταν -- ἀνέχωμεν) Suidas α 2364 (θεαῖν codd. GFS\expo{a.c.} vel θεῶν codd. ATS\expo{p.c} vel θεεῖν cod. I vel θεαῖνων cod. M pro θεοῖν)\par\vspace{2.5pt}
{\color{ruleGold}\hrule height 0.35pt}\vspace{2.5pt}
\noindent{\bfseries\scshape\color{rubricRed}Suite (en commentaire) de l'apparat critique et paléographique :} \textbf{Témoin papyrologique :} \textbf{\PsiDeux{} (\italique{P.Oxy.} LVI 3839, fr. 2)} couvre les \textbf{vv. 941–953}. --- \textbf{Grattage \italique{in rasura} au v. 942 :} \Rac{} avait inséré la particule γε devant τοῖς ; \textbf{\Rdeux{} l'a grattée (\inRas{γε}{} τοῖς)} pour rétablir le trimètre iambique. --- \textbf{Bêtacisme vocalique au v. 947 :} πέσωμεν R pour \lecon{παίσωμεν} (Zanetti, confirmé par {\policeGr Σ}\expo{R}). --- \textbf{v. 927 :} προλείπωσ' R (iotacisme) : \lecon{προλίπωσ'} Scaliger. --- \textbf{v. 947 :} πέσωμεν R : \lecon{παίσωμεν} Zanetti, cf. {\policeGr Σ} \espApp ταῖσ R : \lecon{ταῖσι} Daubuz.\par}
\end{minipage}\end{adjustbox}
\newpage
\renewcommand{\currHeadLE}{SCHOLIES MARGINALES ({\policeGr Σ}\expo{R}) \& GLOSES INTERLINÉAIRES · FOLIO 23/30}
\renewcommand{\currHeadRE}{f. 173v (Classense 175v) · vv. 920–958 · {\policeGr Σ} 294 à {\policeGr Σ} 304}
\renewcommand{\currHeadLO}{SCHOLIES MARGINALES ({\policeGr Σ}\expo{R}) \& GLOSES INTERLINÉAIRES · FOLIO 23/30}
\renewcommand{\currHeadRO}{f. 173v (Classense 175v) · vv. 920–958 · {\policeGr Σ} 294 à {\policeGr Σ} 304}
\bandeauScholies{Folio 23/30 — f. 173v}{Classense 175v}{vv. 920–958}{Verso (page de gauche) · 11 entrées ({\policeGr Σ} 294 à {\policeGr Σ} 304)}
{\fontsize{8.8pt}{10.5pt}\selectfont\noindent\textbf{\color{rubricRed}Disposition topographique des marges sur le f. 173v :} Marge extérieure à gauche ; marge intérieure à droite. — \textit{Fac-similé de référence :} \texttt{folios\_individuels\_ff162v-177r/23\_f173v\_classense175v\_vv0920-0958.jpg}.\par\vspace{3pt}}
\entreeScholie{294}{921}{οἴμ᾽ ὡς πανοῦργος καὐτὸς εἶναί μοι δοκεῖς / καὶ τοῦδέ τις ξύμβουλος.}{Interlinéaire (supra lineam au-dessus de προδότην / ξύμβουλος, f. 173v, v. 921)}{οὐ κυρίως, ἀλλ’ οἷον κοινόβουλος.}{οὐ κυρί\expo{ως}, ἀλλ' οἷον κοινόβουλ\expo{ος} {\policeFr ÷}}{« Non pas au sens propre, mais au sens de complice partageant le même complot (κοινόβουλος). »}
\commScholie{Glose lexicographique de R\expo{2} sur ξύμβουλος (v. 921), précisant que le mot n'est pas employé ici au sens institutionnel athénien de conseiller politique, mais au sens courant de « complice partageant le même dessein » (κοινόβουλος).}
\entreeScholie{295}{922}{οὐκ ἐτὸς πάλαι / ᾐγυπτιάζετ᾽.}{Marge gauche (extérieure, f. 173v, en face du v. 922)}{ἐπανουργεῖτε, ὡς δὴ τῶν Αἰγυπτίων πανούργων ὄντων.}{ἐπανουργεῖτ(ε), ὡς δὴ τ\expo{ῶν} Αἰγυπτί\expo{ων} πανούργ\expo{ων} ὄντ\expo{ων} {\policeFr ÷}}{« ᾐγυπτιάζετε (“vous faisiez les Égyptiens”) : vous jouiez aux fourbes, puisque les Égyptiens avaient la réputation proverbiale d'être des maîtres en fourberie (cf. Eschyle, fr. 373 Radt δεινοὶ πλέκειν τοι μηχανὰς Αἰγύπτιοι). »}
\commScholie{Cité et commenté par Louis Hauseux (vv 642-fin-bis.tex, ll. 2884-2890) et repris par la Souda αι 74.}
\entreeScholie{296}{926}{ἐμπνέω}{Interlinéaire (supra lineam au-dessus d'ἐμπνέω, f. 173v, v. 926)}{ἕως ζῶ.}{ἕως ζῶ {\policeFr ÷}}{« Tant que je respire (tant que je suis en vie). »}
\commScholie{Glose interlinéaire de R\expo{2} expliquant l'emploi absolu tragique du verbe ἐμπνέω (« tant que j'ai le souffle de vie », dans la parodie d'Hélène) par la locution courante ἕως ζῶ (« tant que je vis »).}
\entreeScholie{297}{928}{αὕτη μὲν ἡ μήρινθος οὐδὲν ἔσπασεν.}{Marge gauche (extérieure, f. 173v, en face du v. 928)}{ἡ σπάρτος. ἀντὶ τοῦ ἡ μηχανή. παροιμία ἐπὶ τῶν ἐπιχειρούντων τι ποιεῖν καὶ ἀποτυγχανόντων.}{ἡ σπάρτ\expo{ος} · ἀν\auDessus{τ}{ι} τ\expo{οῦ} ἡ μηχανή · παροιμία ἐπὶ τ\expo{ῶν} ἐπιχειρούντ\expo{ων} τι ποι(εῖν) \textbf{ς} ἀποτυγχανόντ\expo{ων} {\policeFr ÷}}{« ἡ μήρινθος : la ligne de pêche (la cordelette, ἡ σπάρτος), employée métaphoriquement au sens de “cette ruse, ce stratagème”. Proverbe (“cette ligne n'a rien ferré”) appliqué à ceux qui tentent un coup et échouent bredouilles. »}
\commScholie{Cité par Louis Hauseux (vv 642-fin-bis.tex, ll. 2913-2920) avec le parallèle des Guêpes v. 175, de la Souda α 4475 et μ 978, et d'Eustathe 760, 44. Sur R\expo{1}, la réplique était attribuée à Γυ. (Kritylla) ; Scaliger et Bibi la rendent à Mnésiloque (ou Euripide).}
\entreeScholie{298}{935}{ὦ πρύτανι, πρὸς τῆς δεξιᾶς ... / ... ἱστιορράφος.}{Marge gauche (extérieure, f. 173v, en face du v. 935)}{ἐπεὶ ἄνω αἰγυπτιάζειν αὐτοὺς ἔφη, οἱ δὲ Αἰγύπτιοι λινοποιοί εἰσιν. ὁ πανοῦργος. ἀπὸ μεταφορᾶς τῶν τὰ ἄρμενα ῥαπτόντων.}{ἐπεὶ ἄνω αἰγυπτιάζ(ειν) αὐτ\expo{οὺς} ἔφ(η), οἱ δ(ὲ) Αἰγύπτι(οι) λινοποιοί εἰσ(ιν) · ὁ πανοῦργ\expo{ος} · ἀπ(ὸ) μεταφορ\expo{ᾶς} τ\expo{ῶν} τὰ ἄρμεν(α) ῥαπτόντ\expo{ων} {\policeFr ÷}}{« ἱστιορράφος (“couseur de voiles, raccommodeur de ruses”) : parce que plus haut (v. 922) Kritylla les avait accusés de “faire les Égyptiens”, et que les Égyptiens sont des tisserands de lin ; ou simplement au sens de “fourbe qui coud des intrigues” (comme μηχανορράφος), par métaphore de ceux qui cousent les voiles de navire. »}
\commScholie{Repris par la Souda ι 645 : explique ingénieusement le lien entre l'Égypte (terre du lin, Hérodote II, 105) et le composé comique ἱστιο-ρράφος.}
\entreeScholie{299}{940}{γυμνὸν ἀποδύσαντά με / κέλευε πρὸς τῇ σανίδι δεῖν τὸν τοξότην,}{Marge gauche (extérieure, f. 173v, en face du v. 940)}{πρὸς σανίσιν ἐδεσμεύοντο πολλάκις, ὡς καὶ Κρατῖνος δηλοῖ. «ἐπίτρεψον οὖν τῷ τοξότῃ» φησὶ «γυμνόν με πρὸς τῇ σανίδι δῆσαι.»}{πρ(ὸς) σανίσ(ιν) ἐδεσμεύοντ(ο) πολλάκ(ις), ὡς \textbf{ς} Κρατῖν\expo{ος} δηλοῖ · «ἐπίτρεψ\expo{ον} οὖν τῷ τοξότ(ῃ)» φησὶ «γυμνόν με πρ(ὸς) τῇ σανίδι δῆσ(αι)» {\policeFr ÷}}{« On attachait souvent les condamnés au poteau de supplice (la planche, ἡ σανίς), comme le montre aussi Cratinos (fr. 364 K.-A.) : “Permets donc à l'archer scythe”, dit-il, “de m'attacher nu sur la planche.” »}
\commScholie{Scholie alexandrine précieuse (discutée par L. Hauseux, vv 642-fin-bis.tex, ll. 3015-3048, en regard d'Hésychius σ 175 et δ 926) conservant un fragment de Cratinos que parodie directement Mnésiloque.}
\entreeScholie{300}{942}{ἵνα μὴ ᾽ν κροκωτοῖς καὶ μίτραις γέρων ἀνὴρ / γέλωτα παρέχω τοῖς κόραξιν ἑστιῶν·}{Marge gauche (extérieure, f. 173v, bloc de 5 lignes en face des vv. 941–942)}{μετὰ κροκωτοῦ τοῖς κόραξιν ἑστίασιν παρέχων. ἀντὶ τοῦ ἑστίασιν παρέχων. Ἀττικὴ δὲ ἡ σύνταξις, ὅτι τὸ ἑστιᾶν καὶ δοτικῇ συντάσσουσιν, ὡς καὶ παρ’ Εὐπόλιδί ἐστιν ἐν Αὐτολύκῳ θατέρῳ «ὥστ’ οὔτ’ ἐκεῖ οὔτ’ ἐνθάδε δεῖ ξενίζεσθαι.»}{μετὰ κροκωτ\expo{οῦ} τ\expo{οῖς} κόραξ(ιν) ἑστίασ(ιν) παρέχ\expo{ων} · ἀν\auDessus{τ}{ι} τ\expo{οῦ} ἑστίασ(ιν) παρέχ\expo{ων} · Ἀττικὴ δ(ὲ) ἡ σύνταξ\expo{ις}, ὅτι τ(ὸ) ἑστιᾶν \textbf{ς} δοτικῇ συντάσσουσ(ιν), ὡς \textbf{ς} παρ' Εὐπόλιδί ἐστ(ιν) ἐν Αὐτολύκῳ θατέρῳ «ὥστ' οὔτ' ἐκεῖ οὔτ' ἐνθάδ(ε) δεῖ ξενίζεσθ(αι)» {\policeFr ÷}}{« Offrir dans ma robe safran un festin aux corbeaux. La construction est attique, car on construit le verbe ἑστιᾶν / ἑστιᾶσθαι avec le datif, comme on le trouve aussi chez Eupolis dans son second Autolycos (fr. 66 K.-A.) : “De sorte qu'il ne faut donner l'hospitalité ni là-bas ni ici.” »}
\commScholie{Conserve une citation de la seconde rédaction de l'Autolycos d'Eupolis.}
\entreeScholie{301}{945}{ἵνα τοῖς παροῦσι δῆλος ᾖς πανοῦργος ὤν· / — ἰαππαπαιάξ· ὦ κροκώθ᾽ οἷ᾽ εἴργασαι·}{Interlinéaire (supra lineam, f. 173v, v. 944)}{ὅτι ἐνδυσάμενος καὶ ἐλθὼν εἰς τὰς γυναῖκας κατάδηλος ἐγενόμην.}{ὅτι ἐνδυσάμεν\expo{ος} \textbf{ς} ἐλθ(ὼν) εἰς τ\expo{ὰς} γυναῖκ\expo{ας} κατάδηλ\expo{ος} ἐγενόμ(ην) {\policeFr ÷}}{« Parce qu'après avoir revêtu cette robe safran et m'être introduit parmi les femmes, j'ai été démasqué. »}
\commScholie{Glose explicative de R\expo{2} explicitant l'exclamation désespérée de Mnésiloque ὦ κροκώθ᾽ οἷ᾽ εἴργασαι (v. 945) : c'est pour avoir revêtu la robe safran d'Agathon et s'être introduit parmi les femmes qu'il a été démasqué et condamné au carcan.}
\entreeScholie{302}{947}{ἄγε νυν ἡμεῖς παίσωμεν ἅπερ νόμος ἐνθάδε ταῖσι γυναιξίν,}{Marge gauche (extérieure, f. 173v, en face du v. 947)}{τὸ παίξωμεν οἱ Ἀττικοὶ παίσωμεν ἔλεγον, καὶ τὸ παίσωμεν πάλιν παίξωμεν. καὶ [Hom. Od. Ψ, 134] «πολυπαίγμονος ὀρχηθμοῖο.»}{τ(ὸ) παίξωμ(εν) οἱ Ἀττικ(οὶ) παίσωμ(εν) ἔλεγ\expo{ον}, \textbf{ς} τ(ὸ) παίσωμ(εν) πάλ(ιν) παίξωμ(εν) · \textbf{ς} «πολυπαίγμον\expo{ος} ὀρχηθμοῖο» {\policeFr ÷}}{« Ce que la koinè dit παίξωμεν (“dansons, jouons”, de παίζω), les Attiques le disaient παίσωμεν (sans γ/ξ), et inversement pour παίω (“frapper”). Cf. Homère (Odyssée XXIII, 134) : “de la danse enjouée (πολυπαίγμονος ὀρχηθμοῖο).” »}
\commScholie{Témoignage décisif d'atticiste (Ælius Dionysius / Phrynichos) exploité par Zanetti et Louis Hauseux (Bibi, v. 947, vv 642-fin-bis.tex, l. 3117) : alors que R\expo{1} avait écrit par iotacisme πέσωμεν, la scholie {\policeGr Σ}\expo{R} atteste et garantit le subjonctif aoriste attique παίσωμεν !}
\entreeScholie{303}{948}{ὅταν ὄργια σεμνὰ θεαῖν [θεοῖν Σ] ἱεραῖς ὥραις ἀνέχωμεν,}{Marge gauche (extérieure, f. 173v, en face du v. 948)}{ὅταν τὰ μυστήρια αὔξωμεν, φησὶν, τῶν θεῶν. λέγει δὲ Δημήτρας καὶ Περσεφόνης. Ἄλλως. ὅταν τὰς τελετὰς αὔξωμεν· τοῖς γὰρ ὕμνοις ἐπαίρομεν αὐτάς. καὶ τοῦτο ἐστὶ τὸ ἀνέχωμεν.}{ὅτ(αν) τὰ μυστήρια αὔξωμ(εν), φησ(ὶν), τ\expo{ῶν} θε\expo{ῶν} · λέγ(ει) δ(ὲ) Δημήτρ\expo{ας} \textbf{ς} Περσεφόν\expo{ης} · ἄλλ\expo{ως} · ὅτ(αν) τ\expo{ὰς} τελετ\expo{ὰς} αὔξωμ(εν) · τ\expo{οῖς} γ(ὰρ) ὕμν\expo{οις} ἐπαίρομ(εν) αὐτάς · \textbf{ς} τοῦτό ἐστ(ι) τ(ὸ) ἀνέχωμ(εν) {\policeFr ÷}}{« “Lorsque nous célébrons et exaltons (ἀνέχωμεν) les saints mystères des Deux Déesses”, à savoir de Déméter et de Perséphone. — Autre rédaction : Lorsque nous exaltons les cérémonies mystiques, car nous les élevons par nos hymnes (et par les torches brandies) : tel est le sens d'ἀνέχωμεν. »}
\commScholie{Repris par la Souda α 2331.}
\entreeScholie{304}{949}{ἅπερ καὶ Παύσων σέβεται καὶ νηστεύει,}{Marge inférieure du f. 173v (en face du v. 949)}{Παύσων ὁ ζωγράφος εἰς πενίαν διαβάλλεται, καὶ διὰ τοῦτο εἶπε νηστεύει.}{Παύσ(ων) ὁ ζωγράφ\expo{ος} εἰς πενί(αν) διαβάλλετ(αι), \textbf{ς} διὰ τοῦτ(ο) εἶπ(ε) νηστεύ(ει) {\policeFr ÷}}{« Le peintre Pauson est constamment raillé pour son extrême pauvreté (cf. Acharniens v. 854, Ploutos v. 602) : c'est pourquoi le Chœur dit qu'il “observe fidèlement le jeûne” des Thesmophories (et souhaite jeûner toute l'année faute d'avoir de quoi manger !). »}
\commScholie{Cité et commenté par Louis Hauseux (vv 642-fin-bis.tex, ll. 3139-3162) avec le Recueil Milliet d'Adolphe Reinach.}
\clearpage
\renewcommand{\currHeadLE}{ΑΡΙΣΤΟΦΑΝΟΥΣ ΘΕΣΜΟΦΟΡΙΑΖΟΥΣΑΙ · CODEX RAVENNAS 429 (\R) · FOLIO 24/30}
\renewcommand{\currHeadRE}{f. 174r (Classense 176r) · vv. 959–1000 · Collation diplomatique}
\renewcommand{\currHeadLO}{ΑΡΙΣΤΟΦΑΝΟΥΣ ΘΕΣΜΟΦΟΡΙΑΖΟΥΣΑΙ · CODEX RAVENNAS 429 (\R) · FOLIO 24/30}
\renewcommand{\currHeadRO}{f. 174r (Classense 176r) · vv. 959–1000 · Collation diplomatique}
\noindent\begin{adjustbox}{max totalheight=30.4cm, max width=\textwidth, keepaspectratio}\begin{minipage}{\textwidth}
\bandeauRav{f. 174r}{Classense 176r}{959–1000}{Folio 24/30 · Recto (page de droite) · 6 scholies \& gloses ({\policeGr Σ} 305–310)}
{\fontsize{8.8pt}{10.2pt}\selectfont
\ligneRav{959}{χορ/}{\hspace{1.1em}ἅμα δὲ καὶ γένος Ὀλυμπίων θεῶν}{}
\ligneRav{960}{}{\hspace{1.1em}μέλπε καὶ γέραιρε φωνῇ πᾶσα χορομανεῖ τρόπῳ·}{}
\ligneRav{}{}{\hspace{1.1em}εἰ δέ τις προσδοκᾷ κακῶς ἐρεῖν}{}
\ligneRav{965}{}{\hspace{1.1em}ἐν ἱερῷ γυναῖκά μ᾽ οὖσαν ἄνδρας, οὐκ ὀρθῶς φρονεῖ·}{}
\ligneRav{966}{}{\hspace{1.1em}ἀλλὰ χρή}{\textbf{966} \lecon{χρὴ} R : χρῆν Bothe : χρή μ᾽ Meineke}
\ligneRav{967}{}{\hspace{1.1em}ὥσπερ ἔργον αὖ τι καινὸν}{\textbf{967} \lecon{ὥσπερ} R : ὡς πρὸς Austin}
\ligneRav{968}{}{\hspace{1.1em}πρῶτον εὐκύκλου χορείας \difference{εὐφυᾶ} στῆσαι βάσιν·}{\textbf{968} εὐφυῆ R : \lecon{εὐφυᾶ} Brunck}
\ligneRav{969}{}{\hspace{1.1em}πρόβαινε \difference{ποσὶ} τὸν \difference{εὐλύραν}}{\textbf{969} ποσὶν R : \lecon{ποσὶ} Bentley \espApp ἐλύραν R : \lecon{εὐλύραν} Zanetti}
\ligneRav{970}{}{\hspace{1.1em}μέλπουσα καὶ τὴν τοξοφόρον}{}
\ligneRav{}{}{\hspace{1.1em}Ἄρτεμιν ἄνασσαν ἁγνήν·}{}
\ligneRav{}{}{\hspace{1.1em}χαῖρ᾽ ὦ ἑκάεργε,}{}
\ligneRav{}{}{\hspace{1.1em}ὄπαζε δὲ νίκην·}{}
\ligneRav{}{}{\hspace{1.1em}Ἥραν τε τὴν τελείαν}{}
\ligneRav{975}{}{\hspace{1.1em}μέλψωμεν ὥσπερ εἰκός,}{}
\ligneRav{}{}{\hspace{1.1em}ἣ πᾶσι τοῖς χοροῖσιν ἐμπαίζει τε καὶ}{}
\ligneRav{}{}{\hspace{1.1em}κλῇδας γάμου φυλάττει·}{}
\ligneRav{}{}{\hspace{1.1em}Ἑρμῆν τε νόμιον ἄντομαι}{}
\ligneRav{}{}{\hspace{1.1em}καὶ Πᾶνα καὶ Νύμφας φίλας}{}
\ligneRav{}{}{\hspace{1.1em}ἐπιγελάσαι προθύμως}{}
\ligneRav{980}{}{\hspace{1.1em}ταῖς \difference{ἡμετέραισι}}{\textbf{980} ἡμετέραις R : \lecon{ἡμετέραισι} Hermann}
\ligneRav{}{}{\hspace{1.1em}χαρέντα χορείαις·}{}
\ligneRav{}{}{\hspace{1.1em}ἔξαιρε δὴ προθύμως}{}
\ligneRav{982}{}{\hspace{1.1em}διπλῆν \difference{χάριν} χορείας·}{\textbf{982} χαίρειν R : \lecon{χάριν} Ellebodius}
\ligneRav{}{}{\hspace{1.1em}παίσωμεν ὦ γυναῖκες οἷάπερ νόμος,}{}
\ligneRav{984}{}{\hspace{1.1em}\difference{νηστεύομεν} δὲ πάντως·}{\textbf{984} νηστεύωμεν R : \lecon{νηστεύομεν} Bentley}
\ligneRav{985}{}{\hspace{1.1em}ἀλλ᾽ \difference{εἶα πάλλ᾽} ἀνάστρεφ᾽ εὐρύθμῳ ποδί,}{\textbf{985} ειαπαλλ᾽ R : \lecon{εἶα πάλλ᾽} Blaydes}
\ligneRav{}{}{\hspace{1.1em}τόρευε πᾶσαν ᾠδήν·}{}
\ligneRav{987}{}{\hspace{1.1em}ἡγοῦ δέ γ᾽ ὧδ᾽ αὐτὸς σὺ}{\textbf{987} γ᾽ ὧδ᾽ suspectum}
\ligneRav{}{}{\hspace{1.1em}κισσοφόρε Βακχεῖε}{}
\ligneRav{}{}{\hspace{1.1em}δέσποτ᾽· ἐγὼ δὲ κώμοις}{}
\ligneRav{989}{}{\hspace{1.1em}σὲ \difference{φιλοχόροισι} μέλψω·}{\textbf{989} φιλοχόροισιν R : \lecon{φιλοχόροισι} Scaliger}
\ligneRav{990}{}{\hspace{1.1em}†εὔιον ὦ Διόνυσε}{\textbf{990} \lecon{Διόνυσε} R : Διὸς σὺ Enger : Διός τε Fritzsche}
\ligneRav{}{}{\hspace{1.1em}Βρόμιε καὶ Σεμέλας παῖ,}{}
\ligneRav{}{}{\hspace{1.1em}χοροῖς τερπόμενος}{}
\ligneRav{993}{}{\hspace{1.1em}κατ᾽ ὄρεα νυμφᾶν ἐρατοῖς ἐν ὕμνοις,}{\textbf{993} \lecon{ἐρατοῖς ἐν ὕμνοις} R : ἐρατοῖσιν ὕμνοις Wilamowitz}
\ligneRav{994a}{}{\hspace{1.1em}†εὔιον εὔιον εὐοῖ}{\textbf{994a} \lecon{εὔιον εὔιον} R : ὦ Εὔι Εὔι' Hermann}
\ligneRav{994b}{}{\hspace{1.1em}< × \widearc{˘˘} ˘ > ἀναχορεύων·}{\textbf{994b} lacunam statuit Hermann : ἥκεις γὰρ von Velsen : ἡδόμενος Austin cf. \italique{Plut}. 288}
\ligneRav{995}{}{\hspace{1.1em}ἀμφὶ δὲ \difference{σοὶ κτυπεῖται}}{\textbf{995} συῒ κτυπεῖται R : \lecon{σοὶ κτυπεῖται} Zanetti}
\ligneRav{996}{}{\hspace{1.1em}\difference{Κιθαιρώνιος} ἠχώ,}{\textbf{996} Κιθαρώνιος R : \lecon{Κιθαιρώνιος} Zanetti}
\ligneRav{}{}{\hspace{1.1em}μελάμφυλλά τ᾽ ὄρη}{}
\ligneRav{998}{}{\hspace{1.1em}δάσκια †καὶ νάπαι πετρώδεις† βρέμονται·}{\textbf{998} \lecon{καὶ νάπαι πετρώδεις} R : πετρώδεις τε νάπαι Enger}
\ligneRav{}{}{\hspace{1.1em}κύκλῳ δὲ περὶ σὲ κισσὸς}{}
\ligneRav{1000}{}{\hspace{1.1em}εὐπέταλος ἕλικι θάλλει·}{}
}
\vspace{3pt}{\color{ruleGold}\hrule height 0.6pt}\vspace{3pt}
{\fontsize{8.0pt}{9.4pt}\selectfont
\noindent{\bfseries\scshape\color{rubricRed}Loci paralleli et testimonia (f. 174r, vv. 959–1000) :} \textbf{968} (εὐκύκλου) Hesychius ε 6929 εὔκυκλος· εὔχωρος = Photius ε 2236 \espTest \textbf{973} (τελείαν) Hesychius τ 400 Τελεία \espTest \textbf{986} Suidas τ 789\par\vspace{2.5pt}
{\color{ruleGold}\hrule height 0.35pt}\vspace{2.5pt}
\noindent{\bfseries\scshape\color{rubricRed}Suite (en commentaire) de l'apparat critique et paléographique :} \textbf{Contexte :} Grande ode chorale de danse en ronde (\italique{choros kyklikos}) célébrant Apollon, Artémis, Héra, Hermès, Pan, les Nymphes et Dionysos couronné de lierre (vv. 959–1000). --- Page entièrement lyrique où la colométrie de \Run{} remonte directement à l'édition alexandrine d'Aristophane de Byzance (avec \italique{koronis} et \italique{paragraphoi} strophiques). --- \textbf{Fautes phonétiques pittoresques de \Run{} :} Au \textbf{v. 969}, haplographie vocalique ἐλύραν pour \lecon{εὐλύραν} (Zanetti) ; au \textbf{v. 982}, χαίρειν pour \lecon{χάριν} (Ellebodius) ; au \textbf{v. 995},iotacisme/bêtacisme συῒ κτυπεῖται (« retentit pour un porc » !) au lieu de \lecon{σοὶ κτυπεῖται} (« retentit pour toi [Dionysos] », Zanetti).\par}
\end{minipage}\end{adjustbox}
\newpage
\renewcommand{\currHeadLE}{SCHOLIES MARGINALES ({\policeGr Σ}\expo{R}) \& GLOSES INTERLINÉAIRES · FOLIO 24/30}
\renewcommand{\currHeadRE}{f. 174r (Classense 176r) · vv. 959–1000 · {\policeGr Σ} 305 à {\policeGr Σ} 310}
\renewcommand{\currHeadLO}{SCHOLIES MARGINALES ({\policeGr Σ}\expo{R}) \& GLOSES INTERLINÉAIRES · FOLIO 24/30}
\renewcommand{\currHeadRO}{f. 174r (Classense 176r) · vv. 959–1000 · {\policeGr Σ} 305 à {\policeGr Σ} 310}
\bandeauScholies{Folio 24/30 — f. 174r}{Classense 176r}{vv. 959–1000}{Recto (page de droite) · 6 entrées ({\policeGr Σ} 305 à {\policeGr Σ} 310)}
{\fontsize{8.8pt}{10.5pt}\selectfont\noindent\textbf{\color{rubricRed}Disposition topographique des marges sur le f. 174r :} Marge extérieure à droite ; marge intérieure à gauche. — \textit{Fac-similé de référence :} \texttt{folios\_individuels\_ff162v-177r/24\_f174r\_classense176r\_vv0959-1000.jpg}.\par\vspace{3pt}}
\entreeScholie{305}{968}{ἀλλὰ χρῆν / ὥσπερ ἔργον αὖ τι καινὸν / πρῶτον εὐκύκλου χορείας εὐφυᾶ [εὐφυῆ R] στῆσαι βάσιν·}{Marge droite (extérieure, f. 174r, en face des vv. 967–968)}{ὅτι ἐν θεάτρῳ ἦσαν. ἐπειδὴ μέλλουσιν ἐλθεῖν εἰς τὴν ᾠδήν.}{ὅτι ἐν θεάτρῳ ἦσ(αν) · ἐπειδὴ μέλλουσ(ιν) ἐλθ(εῖν) εἰς τ(ὴν) ᾠδήν {\policeFr ÷}}{« Parce qu'elles évoluaient dans l'orchestra du théâtre ; et parce qu'elles vont maintenant passer au chant hymnique (après les mesures préparatoires de la ronde). »}
\commScholie{Note chorégraphique sur la transition entre la mise en place de la ronde chorale (vv. 947-968) et l'hymne aux dieux (vv. 969-1000).}
\entreeScholie{306}{969}{πρόβαινε ποσὶ τὸν Εὐλύραν / μέλπουσα}{Interlinéaire (supra lineam au-dessus de τὸν Εὐλύραν, f. 174r, v. 969)}{τὸν Ἀπόλλωνα.}{τ(ὸν) Ἀπόλλωνα {\policeFr ÷}}{« Apollon (“le Dieu à la belle lyre”). »}
\commScholie{Glose interlinéaire de R\expo{2} identifiant sous l'épithète cultuelle lyrique τὸν Εὐλύραν (« le Dieu à la belle lyre », v. 969) le dieu Apollon (τὸν Ἀπόλλωνα), invoqué en tête de la ronde chorale.}
\entreeScholie{307}{971}{καὶ τὴν τοξοφόρον / Ἄρτεμιν ἄνασσαν ἁγνήν [ἁγνάν Dindorf]·}{Interlinéaire (supra lineam au-dessus d'ἁγνάν, f. 174r, v. 971)}{τὴν παρθένον.}{τ(ὴν) παρθέν\expo{ον} {\policeFr ÷}}{« La Vierge chaste (Artémis). »}
\commScholie{Glose interlinéaire de R\expo{2} expliquant l'épithète rituelle ἁγνήν / ἁγνάν (« pure, chaste ») accolée à Artémis archère (τὴν τοξοφόρον Ἄρτεμιν) par τὴν παρθένον (« la Vierge »).}
\entreeScholie{308}{974}{Ἥραν τε τὴν Τελείαν}{Marge droite (extérieure, f. 174r, en face des vv. 973–974)}{Ἥρα τελεία καὶ Ζεὺς τέλειος ἐτιμῶντο ἐν τοῖς γάμοις ὡς πρυτάνεις ὄντες τῶν γάμων. τέλος δὲ ὁ γάμος. διὸ καὶ προτέλεια ἐκαλεῖτο ἡ θυσία ἡ πρὸ τῶν γάμων γινομένη.}{Ἥρα τελεία \textbf{ς} Ζεὺς τέλει\expo{ος} ἐτιμῶντ(ο) ἐν τ\expo{οῖς} γάμ\expo{οις} ὡς πρυτάν(εις) ὄντ\expo{ες} τ\expo{ῶν} γάμ\expo{ων} · τέλ\expo{ος} δ(ὲ) ὁ γάμ\expo{ος} · διὸ \textbf{ς} προτέλεια ἐκαλεῖτ(ο) ἡ θυσία ἡ πρὸ τ\expo{ῶν} γάμ\expo{ων} γινομέν(η) {\policeFr ÷}}{« Héra Téléia et Zeus Téléios étaient honorés dans les mariages en tant que présidents (πρυτάνεις) de l'union conjugale : car le mariage est appelé un accomplissement (τέλος). C'est pourquoi l'on appelait aussi προτέλεια le sacrifice célébré avant les noces. »}
\commScholie{Repris mot pour mot par la Souda τ 270 et Photius τ 140 : excellente notice d'histoire religieuse sur le culte d'Héra Téléia, qui garde les clés du mariage (ἣ κλῇδας γάμου φυλάττει, v. 976).}
\entreeScholie{309}{977}{Ἑρμῆν τε νόμιον ἄντομαι}{Marge droite (extérieure, f. 174r, en face du v. 977)}{ἔφορος γὰρ τῶν θρεμμάτων ὁ θεός.}{ἔφορ\expo{ος} γ(ὰρ) τ\expo{ῶν} θρεμμάτ\expo{ων} ὁ θεός {\policeFr ÷}}{« Hermès Nomios (“Pastoral”) : car ce dieu veille sur les troupeaux au pâturage (τῶν θρεμμάτων). »}
\commScholie{Explique l'association d'Hermès Nomios avec Pan et les Nymphes (vv. 977-978).}
\entreeScholie{310}{986}{τόρευε [τόρνευε Bentley ex Σ] πᾶσαν ᾠδήν·}{Interlinéaire (supra lineam, f. 174r, v. 986)}{τορῶς καὶ τρανῶς λέγε τὴν ᾠδήν.}{τορ\expo{ῶς} \textbf{ς} τραν\expo{ῶς} λέγε τ(ὴν) ᾠδήν {\policeFr ÷}}{« Entonne le chant d'une voix perçante et distincte (τορῶς καὶ τρανῶς). »}
\commScholie{Glose étymologique rapprochant τόρνευε de τορῶς.}
\clearpage
\renewcommand{\currHeadLE}{ΑΡΙΣΤΟΦΑΝΟΥΣ ΘΕΣΜΟΦΟΡΙΑΖΟΥΣΑΙ · CODEX RAVENNAS 429 (\R) · FOLIO 25/30}
\renewcommand{\currHeadRE}{f. 174v (Classense 176v) · vv. 1001–1043 · Collation diplomatique}
\renewcommand{\currHeadLO}{ΑΡΙΣΤΟΦΑΝΟΥΣ ΘΕΣΜΟΦΟΡΙΑΖΟΥΣΑΙ · CODEX RAVENNAS 429 (\R) · FOLIO 25/30}
\renewcommand{\currHeadRO}{f. 174v (Classense 176v) · vv. 1001–1043 · Collation diplomatique}
\noindent\begin{adjustbox}{max totalheight=30.4cm, max width=\textwidth, keepaspectratio}\begin{minipage}{\textwidth}
\bandeauRav{f. 174v}{Classense 176v}{1001–1043}{Folio 25/30 · Verso (page de gauche) · 18 scholies \& gloses ({\policeGr Σ} 311–328)}
{\fontsize{9.3pt}{10.8pt}\selectfont
\ligneRav{1001}{τοξ/}{ἐνταῦτα νῦν \difference{οἰμῶξι} πρὸς τὴν αἰτρίαν·}{\textbf{1001} οἰμώξει R : \lecon{οἰμῶξι} Brunck}
\ligneRav{1002}{μ\auDessus{ν}{Η}}{ὦ τοξόθ᾽ ἱκετεύω σε. : μή μ᾽ \difference{ἰκετεῦσι} σύ·}{\textbf{1002} ἱκετεῦση R : \lecon{ἱκετεῦσι} Brunck}
\ligneRav{1003}{μ\auDessus{ν}{Η}}{χάλασον τὸν ἧλον. : ἀλλὰ ταῦτα \difference{δρᾶσ᾽} ἐγώ·}{\textbf{1003} δρᾷς R : \lecon{δρᾶσ᾽} Voss}
\ligneRav{1004}{μ\auDessus{ν}{Η}}{οἴμοι κακοδαίμων, μᾶλλον \difference{ἐπικρούεις} σύ γε·}{\textbf{1004} ἐπικρούσεις R : \lecon{ἐπικρούεις} \Mon}
\ligneRav{1005}{τοξ/}{ἔτι \difference{μᾶλλο} βοῦλις; : ἀτταταῖ \difference{ἰατταταῖ}·}{\textbf{1005} μᾶλλον R : \lecon{μᾶλλο} Bentley \espApp τατταταί R : \lecon{ἰατταταῖ} Faber}
\ligneRav{1006}{—}{κακῶς ἀπόλοιο. : σῖγα κακοδαίμω\expo{ν} γέρ\inRas{ω}{ο}ν·}{\textbf{1006} κακοδαίμω\expo{ν} γέρ\inRas{ω}{ο}ν \Rpc, M\expo{pc}}
\ligneRav{1007}{}{πέρ᾽ ἐγὼ \difference{'ξινίγκι} πορμός, \difference{ἴνα} πυλάξι σοι·}{\textbf{1007} ξεινίγκι R : \lecon{'ξινίγκι} Enger \espApp ἵνα R : \lecon{ἴνα} Brunck \espApp \lecon{πυλάξι} R : πυλάξῃ Σ}
\ligneRav{}{μ\auDessus{ν}{Η}}{ταυτὶ τὰ βέλτιστ᾽ ἀπολέλαυκ᾽ Εὐριπίδου·}{}
\ligneRav{1009}{}{ἔα· θεοί, Ζεῦ σῶτερ, εἰσὶν ἐλπίδες·}{\textbf{1009} εἰσὶν R : εἴς ἔτ᾽ Naber (Faber ?) \cf 946}
\ligneRav{1010}{}{\difference{ἁνὴρ} ἔοικεν οὐ προδώσειν, ἀλλά μοι}{\textbf{1010} ἀνὴρ R : \lecon{ἁνὴρ} Bothe}
\ligneRav{1011}{}{σημεῖον \difference{ὑπεδήλωσε} Περσεὺς ἐκδραμών,}{\textbf{1011} ὑπεδήλωσεν R : \lecon{ὑπεδήλωσε} Biset}
\ligneRav{}{}{ὅτι δεῖ με γίγνεσθ᾽ Ἀνδρομέδαν· πάντως δέ μοι}{}
\ligneRav{1013}{}{τὰ δέσμ᾽ ὑπάρχει. δῆλον \difference{οὖν τοῦτ᾽} ἔσθ᾽ ὅτι}{\textbf{1013} οὖν R : οὖν νῦν Bizet : \lecon{οὖν τοῦτ᾽} Dobree}
\ligneRav{1014}{}{ἥξει με \difference{σώσων}· οὐ γὰρ ἂν παρέπτετο·}{\textbf{1014} σῶσον R : \lecon{σώσων} \Mon}
\ligneRav{1015}{\auDessus{εὐ}{ρ}}{\hspace{1.1em}φίλαι παρθένοι φίλαι,}{\textbf{1015-21} (Εὐρ.) R : \lecon{(Μνη.)} Tyrwhitt}
\ligneRav{1016}{}{\hspace{1.1em}πῶς ἂν ἀπέλθοιμι καὶ}{\textbf{1016} \lecon{ἀπέλθοιμι} R : ἐπέλθοιμι Ellebodius}
\ligneRav{1017}{}{\hspace{1.1em}τὸν Σκύθην \difference{λάθοιμι};}{\textbf{1017} λάβοιμι R : \lecon{λάθοιμι} Ellebodius}
\ligneRav{1019}{}{\hspace{1.1em}κλύεις; ὦ \difference{προσᾴδουσ'} \difference{ἀϋτὰς} ἐν ἄντροις,}{\textbf{1019} προσᾳδοῦσσαι τὰς R, Σ : \lecon{προσᾴδουσ'} Elmsley, \lecon{ἀϋτὰς} Burges (ἀϋταῖς Sommerstein) : alii alia}
\ligneRav{1020}{}{\hspace{1.1em}κατάνευσον, ἔασον ὡς}{}
\ligneRav{}{}{\hspace{1.1em}τὴν γυναῖκά μ᾽ ἐλθεῖν·}{}
\ligneRav{1022}{μ\auDessus{ν}{Η}}{\hspace{1.1em}ἄνοικτος ὅς μ᾽ ἔδησε τὸν}{\textbf{1022} (Μνη.) praefixit \Rpc, \Mon}
\ligneRav{1023}{}{\hspace{1.1em}πολυπονώτατον βροτῶν·}{\textbf{1023} \lecon{πολυπονώτατον} R, Σ : πολυστονώτατον Burges}
\ligneRav{}{}{\hspace{1.1em}μόλις δὲ γραῖαν ἀποφυγὼν}{}
\ligneRav{1025}{}{\hspace{1.1em}σαπρὰν ἀπωλόμην ὅμως·}{}
\ligneRav{1026}{}{\hspace{1.1em}ὅδε γὰρ ὁ Σκύθης φύλαξ πάλαι ἐφέστηκ᾽· ὀλοὸν}{\textbf{1026} ὁ Σκύθης R : del. Jackson}
\ligneRav{}{}{\hspace{1.1em}ἄφιλον ἐκρέμασε κόραξι δεῖπνον·}{}
\ligneRav{1030}{}{\hspace{1.1em}ὁρᾷς; οὐ χοροῖσιν οὐδ᾽ ὑφ᾽ ἡλίκων νεανίδων}{}
\ligneRav{1031}{}{\hspace{1.1em}ψῆφον \difference{ἕστηκ᾽ ἔχουσα κημόν},}{\textbf{1031} κημὸν ἕστηκ᾽ ἔχουσ᾽ R : \lecon{ἕστηκ᾽ ἔχουσα κημόν} Bibi, cf. Σ ad \italique{Vesp.} 674c}
\ligneRav{}{}{\hspace{1.1em}ἀλλ᾽ ἐν πυκνοῖς δεσμοῖσιν ἐμπεπλεγμένη}{}
\ligneRav{}{}{\hspace{1.1em}κήτει βορὰ Γλαυκέτῃ πρόκειμαι,}{}
\ligneRav{1034}{}{\hspace{1.1em}γαμηλίῳ μὲν οὐ \difference{ξὺν}}{\textbf{1034-5} ξυμπαιῶνι R : \lecon{ξὺν παιῶνι} Zanetti (\italique{cum carmine} Divus)}
\ligneRav{1035}{}{\hspace{1.1em}\difference{παιῶνι} δεσμίῳ δὲ}{\textbf{1035} cf. 1034}
\ligneRav{}{}{\hspace{1.1em}γοᾶσθέ μ᾽ ὦ γυναῖκες ὡς}{}
\ligneRav{}{}{\hspace{1.1em}μέλεα μὲν πέπονθα μέλεος,}{}
\ligneRav{}{}{\hspace{1.1em}ὢ τάλας ἐγὼ τάλας,}{}
\ligneRav{1039}{}{\hspace{1.1em}ἀπὸ δὲ συγγόνων \difference{ἀλλ᾽} ἄνομα πάθεα,}{\textbf{1039} ἀλλὰν R : \lecon{ἀλλ᾽} Scaliger}
\ligneRav{1040}{}{\hspace{1.1em}†φῶτά \difference{λιτομένα}†,}{\textbf{1040} λιτομέναν R : ἀντομένα (vel -αν) Σ : \lecon{λιτομένα} Enger}
\ligneRav{1041}{}{\hspace{1.1em}πολυδάκρυτον Ἀίδα γόον \difference{ψεφαυγοῦς}·}{\textbf{1041} φεύγουσαν R : χέουσα Rau (\italique{Paratragödia}, p.75) : φλέγουσαν Musgrave (ad Eur. \italique{Or.}1402=1396) : \lecon{ψεφαυγοῦς} Bibi, cf. Σ ad v.1054, Hesych. ψ 135}
\ligneRav{}{}{\hspace{1.1em}αἰαῖ αἰαῖ ἒ ἒ,}{}
\ligneRav{}{}{\hspace{1.1em}ὃς ἔμ᾽ ἀπεξύρησε πρῶτον,}{}
}
\vspace{3pt}{\color{ruleGold}\hrule height 0.6pt}\vspace{3pt}
{\fontsize{8.0pt}{9.4pt}\selectfont
\noindent{\bfseries\scshape\color{rubricRed}Loci paralleli et testimonia (f. 174v, vv. 1001–1043) :} Tradition directe : Codex Ravennas 429 (R) et scholies anciennes ({\policeGr Σ}\expo{R}).\par\vspace{2.5pt}
{\color{ruleGold}\hrule height 0.35pt}\vspace{2.5pt}
\noindent{\bfseries\scshape\color{rubricRed}Suite (en commentaire) de l'apparat critique et paléographique :} \textbf{Témoin papyrologique :} \textbf{\PsiQuatre{} (\italique{P.Oxy.} LXXIII 4935, fr. 1)} débute au \textbf{v. 1043}. --- \textbf{Traitement du grec barbare de l'Archer scythe (vv. 1001–1007) :} \Run{} tend parfois à normaliser les finales scythes en grec standard (οἰμώξει, ἱκετεῦση, μᾶλλον, ἵνα) ; au \textbf{v. 1006}, \Rpc{} (suivi de \Mon\expo{p.c.}) ajoute un ν en exposant et corrige \italique{in rasura} l'oméga de γέρων en omicron : \lecon{κακοδαίμω\expo{ν} γέρ\inRas{ω}{ο}ν}. --- \textbf{Attribution des vv. 1015–1021 :} R attribue le début de la monodie à (Εὐρ.), corrigé en \lecon{(Μνη.)} par Tyrwhitt, tandis qu'au \textbf{v. 1022} \Rpc{} (et \Mon{}) préfixe (Μνη.). --- \textbf{Deux interventions décisives de Bibi (vv. 1031 et 1041) :} 1° Au \textbf{v.~1031}, \R{} a inversé l'ordre des mots (κημὸν ἕστηκ' ἔχουσ') et Bibi rétablit \lecon{ἕστηκ' ἔχουσα κημόν} (cf. {\policeGr Σ} ad \italique{Vesp.} 674c et van Leeuwen 1904, pp.~XVIII–XIX) ; 2° au \textbf{v.~1041}, là où \R{} porte le participe banal φεύγουσαν, Louis Hauseux (Bibi) restitue le génitif du composé tragique rare \lecon{ψεφαυγοῦς} (« de l'obscure [nuit/grotte] », attesté par {\policeGr Σ}\expo{R} ad v.~1054 et Hésychius ψ 135 ψεφαυγής· σκοτεινή). --- \textbf{v. 1007 :} ξεινίγκι R : \lecon{'ξινίγκι} Enger \espApp ἵνα R : \lecon{ἴνα} Brunck \espApp \lecon{πυλάξι} R (πυλάξῃ {\policeGr Σ}). --- \textbf{v. 1017 :} λάβοιμι R (confusion minuscule θ/β) : \lecon{λάθοιμι} Ellebodius. --- \textbf{v. 1019 :} προσᾳδοῦσσαι τὰς R, {\policeGr Σ} : \lecon{προσᾴδουσ'} Elmsley, \lecon{ἀϋτὰς} Burges (ἀϋταῖς Sommerstein). --- \textbf{v. 1023 :} \lecon{πολυπονώτατον} R, {\policeGr Σ} (πολυστονώτατον Burges). --- \textbf{v. 1031 :} κημὸν ἕστηκ' ἔχουσ' R : \lecon{ἕστηκ' ἔχουσα κημόν} Bibi, cf. {\policeGr Σ} ad \italique{Vesp.} 674c. --- \textbf{v. 1040 :} λιτομέναν R (ἀντομένα vel -αν {\policeGr Σ}, confusion onciale ΛΙ/ΑΝ) : \lecon{λιτομένα} Enger. --- \textbf{v. 1041 :} φεύγουσαν R (χέουσα Rau ; φλέγουσαν Musgrave) : \lecon{ψεφαυγοῦς} Bibi, cf. {\policeGr Σ} ad v. 1054, Hésychius ψ 135.\par}
\end{minipage}\end{adjustbox}
\newpage
\renewcommand{\currHeadLE}{SCHOLIES MARGINALES ({\policeGr Σ}\expo{R}) \& GLOSES INTERLINÉAIRES · FOLIO 25/30}
\renewcommand{\currHeadRE}{f. 174v (Classense 176v) · vv. 1001–1043 · {\policeGr Σ} 311 à {\policeGr Σ} 328}
\renewcommand{\currHeadLO}{SCHOLIES MARGINALES ({\policeGr Σ}\expo{R}) \& GLOSES INTERLINÉAIRES · FOLIO 25/30}
\renewcommand{\currHeadRO}{f. 174v (Classense 176v) · vv. 1001–1043 · {\policeGr Σ} 311 à {\policeGr Σ} 328}
\bandeauScholies{Folio 25/30 — f. 174v}{Classense 176v}{vv. 1001–1043}{Verso (page de gauche) · 18 entrées ({\policeGr Σ} 311 à {\policeGr Σ} 328)}
{\fontsize{8.8pt}{10.5pt}\selectfont\noindent\textbf{\color{rubricRed}Disposition topographique des marges sur le f. 174v :} Marge extérieure à gauche ; marge intérieure à droite. — \textit{Fac-similé de référence :} \texttt{folios\_individuels\_ff162v-177r/25\_f174v\_classense176v\_vv1001-1043.jpg}.\par\vspace{3pt}}
\entreeScholie{311}{1001}{ἐνταῦτα νῦν οἰμῶξι πρὸς τὴν αἰτρίαν.}{Marge supérieure et marge gauche (extérieure, f. 174v, en haut du verso, entrée de l'Archer scythe)}{ἀντὶ τοῦ πρὸς τὴν αἰθρίαν. βαρβαρίζει δὲ ὁ τοξότης.}{ἀν\auDessus{τ}{ι} τ\expo{οῦ} πρ(ὸς) τ(ὴν) αἰθρί(αν) · βαρβαρίζ(ει) δ(ὲ) ὁ τοξότ(ης) {\policeFr ÷}}{« πρὸς τὴν αἰτρίαν : au lieu de πρὸς τὴν αἰθρίαν (“à la belle étoile, en plein air”). L'archer public scythe écorche le grec en barbare (βαρβαρίζει, en prononçant les aspirées sourdes θ, φ, χ comme des occlusives simples τ, π, κ). »}
\commScholie{Première note linguistique de R\expo{2} sur le sabir gréco-scythe du Toxotês (vv. 1001-1225), qui remplace systématiquement les aspirées par des ténues (αἰτρίαν = αἰθρίαν, πέρ' = φέρ', πορμός = φορμός, πυλάξι = φυλάξω) et confond les désinences verbales.}
\entreeScholie{312}{1003}{χάλασον τὸν ἧλον. — ἀλλὰ ταῦτα δρᾶσ᾽ [δρᾷς R] ἐγώ·}{Marge gauche (extérieure, f. 174v, en face des vv. 1003–1005)}{ἀντὶ τοῦ δρῶ. τοῦτο δὲ εἰπὼν πλέον αὐτὸν ἐπισφίγγει.}{ἀν\auDessus{τ}{ι} τ\expo{οῦ} δρῶ · τοῦτ(ο) δ(ὲ) εἰπ(ὼν) πλέον αὐτ(ὸν) ἐπισφίγγ(ει) {\policeFr ÷}}{« (Sur δρᾶσον / ναὶ μᾶλλον :) au sens de “je le fais”. Mais tout en disant cela (à Mnésiloque qui le suppliait de desserrer la cheville du carcan, χάλασον τὸν ἧλον), le Scythe la lui resserre encore plus fort ! »}
\commScholie{Note de jeu scénique indispensable à l'intelligence des vv. 1003-1005 (d'où le cri de Mnésiloque au v. 1005 : οἴμοι κακοδαίμων, μᾶλλον ἐπικρούεις σὺ γε !).}
\entreeScholie{313}{1007}{πέρ᾽ ἐγὼ ᾽ξινίγκι πορμός, ἴνα πυλάξι σοι.}{Interlinéaire (supra lineam au-dessus de πορμός et d'ἴνα πυλάξι, f. 174v, v. 1007)}{πορμός ἀντὶ τοῦ φορμὸς, ψίαθος. ἀντὶ τοῦ ἵνα φυλάξω σε.}{πορμός ἀν\auDessus{τ}{ι} τ\expo{οῦ} φορμ\expo{ὸς}, ψίαθ\expo{ος} {\policeFr ÷} // ἀν\auDessus{τ}{ι} τ\expo{οῦ} ἵνα φυλάξω σε {\policeFr ÷}}{« πορμός : au lieu de φορμός, une natte de jonc (ψίαθος, sur laquelle le garde va s'allonger). — ἴνα πυλάξι σοι : au lieu de ἵνα φυλάξω σε (“afin que je te garde”). »}
\commScholie{Traduit en grec standard le sabir de l'Archer scythe (commenté par L. Hauseux, vv 642-fin-bis.tex, ll. 3452-3454).}
\entreeScholie{314}{1009}{ἔα· θεοί, Ζεῦ σῶτερ, εἰσὶν ἐλπίδες· / ἁνὴρ ἔοικεν οὐ προδώσειν, ἀλλά μοι / σημεῖον ὑπεδήλωσε Περσεὺς ἐκδραμών,}{Marge intérieure droite (f. 174v, en face des vv. 1009–1010)}{ὁρᾷ Εὐριπίδην προσιόντα.}{ὁρᾷ Εὐριπίδ(ην) προσιόντ(α) {\policeFr ÷}}{« Mnésiloque aperçoit furtivement Euripide qui s'approche (costumé en Persée et lui faisant signe). »}
\commScholie{Note scénique cruciale expliquant pourquoi Mnésiloque entonne soudain la monodie d'Andromède aux vv. 1015 sqq.}
\entreeScholie{315}{1011}{σημεῖον ὑπεδήλωσε Περσεὺς ἐκδραμών,}{Interlinéaire (supra lineam, f. 174v, v. 1011)}{ἀντὶ τοῦ ὡς Περσεύς.}{ἀν\auDessus{τ}{ι} τ\expo{οῦ} ὡς Περσεύς {\policeFr ÷}}{« En tant que Persée (déguisé avec les attributs de Persée : talonnières ailées et gorgonéion). »}
\commScholie{Glose syntaxique de R\expo{2} précisant la valeur appositive et prédicative du nom propre Περσεὺς dans Περσεὺς ἐκδραμών (v. 1011) : Euripide a traversé la scène en courant « sous les traits de Persée » (ὡς Περσεύς).}
\entreeScholie{316}{1012}{ὅτι δεῖ με γίγνεσθ᾽ Ἀνδρομέδαν·}{Marge gauche (extérieure, f. 174v, en face du v. 1012)}{Ἀνδρομέδαν: πιθανῶς· συνδεδίδακται γὰρ τῇ Ἑλένῃ.}{Ἀνδρομέδ(αν) : πιθαν\expo{ῶς} · συνδεδίδακτ(αι) γ(ὰρ) τῇ Ἑλένῃ {\policeFr ÷}}{« Ἀνδρομέδαν : fort à propos, car l'Andromède d'Euripide a été représentée dans la même trilogie (la même année, en 412 av. J.-C.) que l'Hélène ! »}
\commScholie{Renseignement didascalique fondamental d'Aristophane de Byzance / Didyme (cité par Louis Hauseux, Introduction.tex, l. 176, et vv 642-fin-bis.tex, l. 3484 n. 1) : l'Hélène et l'Andromède d'Euripide faisaient partie de la même trilogie représentée en 412 av. J.-C. !}
\entreeScholie{317}{1015}{φίλαι παρθένοι φίλαι,}{Marge gauche (extérieure, f. 174v, en face des vv. 1015–1017)}{παρὰ τὰ ἐξ Ἀνδρομέδας Εὐριπίδου «φίλαις παρθένοι, φίλαι μοι.» τὰ δὲ ἐπιφερόμενα πρὸς τὸ αὐτὸ χρήσιμον.}{παρὰ τὰ ἐξ Ἀνδρομέδ\expo{ας} Εὐριπίδ\expo{ου} «φίλ(αι) παρθέν(οι), φίλ(αι) μοι» · τὰ δ(ὲ) ἐπιφερόμεν(α) πρ(ὸς) τ(ὸ) αὐτ(ὸ) χρήσιμ\expo{ον} {\policeFr ÷}}{« Parodie des vers de l'Andromède d'Euripide (fr. 117 Kannicht) : “Chères jeunes filles (les jeunes filles éthiopiennes du chœur), qui m'êtes chères !” Quant à la suite (“comment pourrais-je échapper au Scythe ?”), il l'adapte à sa propre situation. »}
\commScholie{Distingue avec précision l'incipit tragique authentique d'Euripide et sa continuation burlesque par Mnésiloque.}
\entreeScholie{318}{1018}{κλύεις; ὦ προσᾴδουσ᾽ ἀϋτὰς ἐν ἄντροις,}{Marge gauche (extérieure, f. 174v, bloc de 6 lignes en face des vv. 1018–1021)}{πάλιν ἐξ Ἀνδρομέδας. πρὸς τὴν ἠχὼ Ἀνδρομέδα λέγει «προσαιδοῦσσαι τὰς ἐν ἄντροις ἀπόπαυσον ἔασον ἀχοῖ με σὺν φίλαις γόου πόθον λαβεῖν.» διὰ τὸ λαμβάνειν ἀποσπάσματα ἀσύνδετον τὸ ὅλον γίνεται.}{πάλ(ιν) ἐξ Ἀνδρομέδ\expo{ας} · πρ(ὸς) τ(ὴν) ἠχὼ Ἀνδρομέδα λέγ(ει) «προσαιδοῦσσαι τὰς ἐν ἄντρ\expo{οις} ἀπόπαυσ\expo{ον} ἔασ\expo{ον} ἀχοῖ με σὺν φίλ\expo{αις} γό\expo{ου} πόθ\expo{ον} λαβ(εῖν)» · διὰ τ(ὸ) λαμβάν(ειν) ἀποσπάσματ(α) ἀσύνδετ\expo{ον} τ(ὸ) ὅλ\expo{ον} γίνετ(αι) {\policeFr ÷}}{« De nouveau tiré de l'Andromède (fr. 118 Kannicht). Andromède s'adresse à Écho et dit : “Ô toi qui accompagnes mes cris au fond des grottes, cesse, laisse-moi, ô Écho, rassasier avec mes amies mon désir de lamentations !” Et parce qu'Aristophane prélève des lambeaux détachés (ἀποσπάσματα) de la tragédie, l'ensemble de la monodie paraît décousu (ἀσύνδετον). »}
\commScholie{Témoignage majeur (commenté par L. Hauseux, vv 642-fin-bis.tex, ll. 3502-3520) : non seulement cette scholie nous restitue le texte complet du fragment 118 d'Euripide, mais elle formule un jugement critique d'une grande finesse sur la technique de montage en centon (διὰ τὸ λαμβάνειν ἀποσπάσματα ἀσύνδετον τὸ ὅλον γίνεται).}
\entreeScholie{319}{1021}{κατάνευσον, ἔασον ὡς / τὴν γυναῖκά μ᾽ ἐλθεῖν·}{Interlinéaire (supra lineam, f. 174v, v. 1021)}{ὡς ἐν κωμῳδίᾳ.}{ὡς ἐν κωμῳδίᾳ {\policeFr ÷}}{« Ajouté sur le ton de la comédie (pour revenir à sa femme au logis). »}
\commScholie{Fine remarque dramaturgique de R\expo{2} soulignant la rupture comique (ὡς ἐν κωμῳδίᾳ) au milieu de la monodie paratragique d'Andromède : au lieu de demander à rejoindre son époux ou son père, le vieux Mnésiloque demande à rentrer « chez sa femme » (ὡς τὴν γυναῖκά μ᾽ ἐλθεῖν).}
\entreeScholie{320}{1022}{ἄνοικτος ὅς μ᾽ ἔδησε, / τὸν πολυπονώτατον βροτῶν·}{Marge gauche (extérieure, f. 174v, en face des vv. 1022–1023)}{παρὰ τὸ τοῦ χοροῦ ἐν Ἀνδρομέδᾳ «ἄνοικτος ὃς τεκών σε τὴν πολυπονωτάτην βροτῶν μεθῆκεν ᾍδᾳ πατρὸς ὑπερθανεῖν.»}{παρὰ τ(ὸ) τ\expo{οῦ} χορ\expo{οῦ} ἐν Ἀνδρομέδᾳ «ἄνοικτ\expo{ος} ὃς τεκών σε τ(ὴν) πολυπονωτάτ(ην) βροτ\expo{ῶν} μεθῆκ(εν) ᾍδᾳ πατρ(ὸς) ὑπερθαν(εῖν)» {\policeFr ÷}}{« Parodie de la réplique du Chœur dans l'Andromède (fr. 119 Kannicht) : “Impitoyable est celui qui, t'ayant engendrée, t'a livrée à Hadès, toi la plus infortunée des mortelles, pour mourir au nom de ta patrie (ou de ton père Céphée) !” »}
\commScholie{Autre fragment capital de l'Andromède d'Euripide conservé par la seule scholie du Ravennas (commenté par L. Hauseux, vv 642-fin-bis.tex, l. 3524).}
\entreeScholie{321}{1024}{μόλις δὲ γραῖαν ἀποφυγὼν / σαπρὰν}{Interlinéaire (supra lineam au-dessus de γραῖαν, f. 174v, v. 1024)}{τὴν τηρήσασαν αὐτὸν ἐν τοῖς θεσμοφορίοις.}{τ(ὴν) τηρήσασ(αν) αὐτ(ὸν) ἐν τ\expo{οῖς} θεσμοφορί\expo{οις} {\policeFr ÷}}{« La vieille femme (Kritylla) qui le gardait prisonnier dans le Thesmophorion. »}
\commScholie{Glose explicative de R\expo{2} identifiant la « vieille décrépite » (γραῖαν σαπράν, vv. 1024–1025) à qui Mnésiloque vient tout juste d'échapper comme étant Critylla, qui le gardait dans le Thesmophorion avant l'arrivée du Prytane et de l'Archer.}
\entreeScholie{322}{1025}{ἀπωλόμην ὅμως.}{Marge gauche (extérieure, f. 174v, en face du v. 1025)}{δέον εἰπεῖν ἐσώθην, ἀπωλόμην εἶπε χάριν γέλωτος.}{δέον εἰπ(εῖν) ἐσώθ(ην), ἀπωλόμ(ην) εἶπ(ε) χάρ(ιν) γέλωτ\expo{ος} {\policeFr ÷}}{« Alors qu'après “ayant à grand-peine échappé à la vieille” on attendait “j'ai été sauvé (ἐσώθην)”, il a dit par surprise comique “je n'en suis pas moins perdu (ἀπωλόμην)”, pour faire rire ! »}
\commScholie{Souligne l'effet παρὰ προσδοκίαν de ἀπωλόμην ὅμως.}
\entreeScholie{323}{1030}{ὁρᾷς; οὐ χοροῖσιν οὐδ᾽ ὑφ᾽ ἡλίκων νεανίδων}{Marge gauche (extérieure, f. 174v, en face du v. 1030)}{πάλιν ἐξ Ἀνδρομέδας. τὸ δὲ ἐπιφερόμενον σκώπτοντός ἐστι τοὺς Ἀθηναίους ὡς φιλοδίκους.}{πάλ(ιν) ἐξ Ἀνδρομέδ\expo{ας} · τ(ὸ) δ(ὲ) ἐπιφερόμεν\expo{ον} σκώπτοντός ἐστ(ι) τ\expo{οὺς} Ἀθηναί\expo{ους} ὡς φιλοδίκ\expo{ους} {\policeFr ÷}}{« De nouveau tiré de l'Andromède (fr. 120 Kannicht). Mais ce qui suit immédiatement (au v. 1031 : “tenant la corbeille à bulletins de vote”) raille les vieillards athéniens pour leur manie des procès (ὡς φιλοδίκους). »}
\commScholie{Distingue le vers tragique (v. 1030) de la pointe comique du v. 1031.}
\entreeScholie{324}{1031}{ψῆφον ἕστηκ᾽ ἔχουσα κημόν,}{Marge gauche (extérieure, f. 174v, en face du v. 1031)}{παρὰ τὴν ψηφοθήκην. οὐ δικάζω, φησὶν ὁ γέρων.}{παρὰ τ(ὴν) ψηφοθήκ(ην) · οὐ δικάζω, φησ(ὶν) ὁ γέρ(ων) {\policeFr ÷}}{« (Le κημός est l'entonnoir d'osier coiffant) l'urne à bulletins de vote (ψηφοθήκη). “Je ne siège plus au tribunal comme héliaste”, dit le vieillard. »}
\commScholie{Longuement discuté par Louis Hauseux (Bibi, v. 1031, vv 642-fin-bis.tex, ll. 3553-3995) : contre la transposition de R\expo{1} κημὸν ἕστηκ' ἔχουσ', Bibi restitue l'ordre métrique ψῆφον ἕστηκ' ἔχουσα κημόν et examine la variante παρὰ τὴν ξιφοθήκην / ψηφοθήκην et les parallèles des Guêpes (vv. 99, 674, 754).}
\entreeScholie{325}{1033}{κήτει βορὰ Γλαυκέτῃ πρόκειμαι.}{Marge gauche (extérieure, f. 174v, en face du v. 1033)}{ἐπεὶ ὀψοφάγος καὶ γαστρίμαργος ὁ Γλαυκέτης, ὡς ἐν Εἰρήνῃ [1008] δηλοῦται.}{ἐπεὶ ὀψοφάγ\expo{ος} \textbf{ς} γαστρίμαργ\expo{ος} ὁ Γλαυκέτ(ης), ὡς ἐν Εἰρήνῃ δηλοῦτ(αι) {\policeFr ÷}}{« (Il substitue παρ' ὑπόνοιαν le nom propre Glaukétès à “monstre marin”, κήτει βορά) : parce que Glaukétès était un glouton fameux et un amateur effréné de poissons (ὀψοφάγος), comme le montre Aristophane dans la Paix (v. 1008). »}
\commScholie{Cité par Louis Hauseux (vv 642-fin-bis.tex, ll. 4004-4035) : dans l'Andromède d'Euripide (fr. 122 K., cité par Σ Oiseaux v. 348), l'héroïne disait ἐκθεῖναι κήτει φορβάν ; Aristophane ajoute malicieusement le nom du vorace Glaukétès en apposition à κήτει (« exposé en pâture au cétacé... Glaukétès ! »).}
\entreeScholie{326}{1034}{γαμηλίῳ μὲν οὐ ξὺν / παιῶνι, δεσμίῳ δὲ,}{Marge gauche (extérieure, f. 174v, en face des vv. 1034–1035)}{καὶ τοῦτο ἐξ Ἀνδρομέδας.}{\textbf{ς} τοῦτ(ο) ἐξ Ἀνδρομέδ\expo{ας} {\policeFr ÷}}{« Ces vers aussi sont repris textuellement de l'Andromède d'Euripide (fr. 122 Kannicht). »}
\commScholie{Indication de source littéraire de R\expo{2} signalant que l'opposition lyrique entre le péan d'hyménée et le chant de captivité (γαμηλίῳ μὲν οὐ ξὺν παιῶνι, δεσμίῳ δέ, vv. 1034–1035 = fr. 122 Kannicht) est elle aussi empruntée à l'Andromède d'Euripide.}
\entreeScholie{327}{1040}{†φῶτα λιτομένα†,}{Marge inférieure et marge gauche (bas du f. 174v, bloc de 5 lignes en face des vv. 1037–1040)}{ἀπὸ κοινοῦ λάβωμεν πέπονθα ἄνομα πάθεα. καὶ ταῦτα ἐξ Ἀνδρομέδας. δεομένη τοῦ βαρβάρου. προεῖπε δὲ «γοᾶσθέ μ’ ὦ γυναῖκες.» γρ. καὶ φῶτα ἀντομένα· καὶ τοῦτο ἔχει νοῦν πρὸς τὰ ἑξῆς, φῶτ’ ἀντομέναν, ὃς ἔμ’ ἀπεξύρησε πρῶτον.}{ἀπ(ὸ) κοιν\expo{οῦ} λάβωμ(εν) πέπονθ(α) ἄνομα πάθεα · \textbf{ς} ταῦτ(α) ἐξ Ἀνδρομέδ\expo{ας} · δεομέν(η) τ\expo{οῦ} βαρβάρ\expo{ου} · προεῖπ(ε) δ(ὲ) «γοᾶσθέ μ' ὦ γυναῖκ\expo{ες}» · γρ(άφεται) \textbf{ς} φῶτα ἀντομένα · \textbf{ς} τοῦτ(ο) ἔχ(ει) νοῦν πρ(ὸς) τὰ ἑξ\expo{ῆς}, φῶτ' ἀντομέναν, ὃς ἔμ' ἀπεξύρησ(ε) πρῶτ\expo{ον} {\policeFr ÷}}{« Sous-entendons en facteur commun : “j'ai souffert des maux impies” (πέπονθα ἄνομα πάθεα). Ces vers aussi viennent de l'Andromède. — λιτομένα : suppliant le barbare ; et il avait dit plus haut : “Pleurez-moi, ô femmes !” — On lit aussi la variante (γράφεται) φῶτα ἀντομένα (“suppliant l'homme”) ; et cette leçon offre un bon sens par rapport à la proposition relative qui suit (au v. 1043) : “suppliant l'homme (Euripide) qui m'a d'abord rasé de près”. »}
\commScholie{Scholie critique fondamentale longuement analysée par Louis Hauseux (Bibi, v. 1040, vv 642-fin-bis.tex, ll. 4062-4105) : la variante antique φῶτα ἀντομένα(ν) transmise par {\policeGr Σ}\expo{R} (face à φῶτα λιτομέναν de R\expo{1}) montre que dans l'archétype en onciale ΦΩΤΑΛΙΤΟΜΕΝΑ et ΦΩΤΑΑΝΤΟΜΕΝΑ procédaient d'une même séquence graphique.}
\entreeScholie{328}{1041}{πολυδάκρυτον Ἀίδα γόον ψεφαυγοῦς,}{Interlinéaire (supra lineam, dernière ligne du f. 174v, v. 1041)}{ἀντὶ τοῦ πολυδακρύτου.}{ἀν\auDessus{τ}{ι} τ\expo{οῦ} πολυδακρύτ\expo{ου} {\policeFr ÷}}{« πολυδάκρυτον (qualifiant γόον par hypallage) : au sens de πολυδακρύτου (se rapportant à Hadès aux mille larmes). »}
\commScholie{Note l'hypallage tragique de l'adjectif composé.}
\clearpage
\renewcommand{\currHeadLE}{ΑΡΙΣΤΟΦΑΝΟΥΣ ΘΕΣΜΟΦΟΡΙΑΖΟΥΣΑΙ · CODEX RAVENNAS 429 (\R) · FOLIO 26/30}
\renewcommand{\currHeadRE}{f. 175r (Classense 177r) · vv. 1044–1081 · Collation diplomatique}
\renewcommand{\currHeadLO}{ΑΡΙΣΤΟΦΑΝΟΥΣ ΘΕΣΜΟΦΟΡΙΑΖΟΥΣΑΙ · CODEX RAVENNAS 429 (\R) · FOLIO 26/30}
\renewcommand{\currHeadRO}{f. 175r (Classense 177r) · vv. 1044–1081 · Collation diplomatique}
\noindent\begin{adjustbox}{max totalheight=30.4cm, max width=\textwidth, keepaspectratio}\begin{minipage}{\textwidth}
\bandeauRav{f. 175r}{Classense 177r}{1044–1081}{Folio 26/30 · Recto (page de droite) · 15 scholies \& gloses ({\policeGr Σ} 329–343)}
{\fontsize{9.3pt}{10.8pt}\selectfont
\ligneRav{1044}{μ\auDessus{ν}{Η}}{\hspace{1.1em}ὃς ἐμὲ κροκόεντ᾽ ἐνέδυσεν·}{\textbf{1044} \lecon{κροκόεντ᾽} R : κροκόεν τόδ' Bergk \espApp ἐνέδυσεν R : ἀμφέδυσεν Hermann ap Σ}
\ligneRav{1045}{}{\hspace{1.1em}ἐπὶ δὲ τοῖσδε τόδ᾽ ἀνέπεμψεν}{}
\ligneRav{}{}{\hspace{1.1em}ἱερόν, ἔνθα γυναῖκες·}{}
\ligneRav{1047}{}{\hspace{1.1em}ἰώ μοι μοίρας}{\textbf{1047} [ι]\douteux{ω}μοι\douteux{ρ}[ας \PsiQuatre : \lecon{ἰώ μοι μοίρας} R : ἰώ μοίρας Hermann}
\ligneRav{}{}{\hspace{1.1em}\difference{ἄτεγκτε} δαίμων·}{\textbf{1047b} ἀνέτικτε R : \lecon{ἄτεγκτε} Ellebodius : ἅν ἔτικτε Casaubon \espApp \lecon{δαίμων} R : δαῖμον Burges}
\ligneRav{}{}{\hspace{1.1em}ὦ κατάρατος ἐγώ·}{}
\ligneRav{}{}{\hspace{1.1em}τίς ἐμὸν οὐκ ἐπόψεται}{}
\ligneRav{}{}{\hspace{1.1em}πάθος ἀμέγαρτον ἐπὶ κακῶν παρουσίᾳ;}{}
\ligneRav{1050}{}{\hspace{1.1em}εἴθε με πυρφόρος αἰθέρος ἀστὴρ—}{}
\ligneRav{1051}{}{\hspace{1.1em}τὸν \difference{βάρβαρον} ἐξολέσειεν·}{\textbf{1051} \lecon{βάρβαρον} \PsiQuatre{}, R, Willems}
\ligneRav{1052}{}{\hspace{1.1em}οὐ γὰρ ἔτ᾽ ἀθανάταν φλόγα \difference{λεύσσειν}}{\textbf{1052} λεύσειν R : \lecon{λεύσσειν} Biset}
\ligneRav{}{}{\hspace{1.1em}ἐστὶν ἐμοὶ φίλον, ὡς ἐκρεμάσθην,}{}
\ligneRav{1054}{}{\hspace{1.1em}λαιμότμητ᾽ ἄχη \difference{δία μένων} αἰόλαν}{\textbf{1054} δαιμόνων R, Σ : δαιμονῶν Brunck : δαιμόνι᾽ amicus Fritzschii : \lecon{διὰ μένων} Willems}
\ligneRav{1055}{}{\hspace{1.1em}νέκυσιν ἐπὶ πορείαν·}{}
\ligneRav{1056}{ευρ\auDessus{ι}{π} ηχω}{χαῖρ᾽ ὦ φίλη παῖ· τὸν δὲ πατέρα Κηφέα}{\textbf{1056} ευρ\auDessus{ι}{π} ηχω in mg. \Rpc, \Mon}
\ligneRav{}{}{ὅς σ᾽ ἐξέθηκεν ἀπολέσειαν οἱ θεοί·}{}
\ligneRav{1058}{μ\auDessus{ν}{Η}}{σὺ δ᾽ εἶ τίς ἥτις τοὐμὸν \difference{ᾤκτιρας} \difference{πάθος};}{\textbf{1058} ὤκτειρας R : \lecon{ᾤκτιρας} van Herwerden \espApp πάθος R : \lecon{πάθος;} Biset}
\ligneRav{}{\auDessus{εὐ}{ρ}}{Ἠχὼ, λόγων ἀντῳδὸς ἐπικοκκάστρια,}{}
\ligneRav{1060}{}{ἥπερ πέρυσιν ἐν τῷδε \difference{ταὐτῷ} χωρίῳ}{\textbf{1060} ταυτῷ R : \lecon{ταὐτῷ} Dindorf (1830)}
\ligneRav{1061}{}{Εὐριπίδῃ \difference{καὐτὴ} ξυνηγωνιζόμην·}{\textbf{1061} καὶ αὐτὴ R : \lecon{καὐτὴ} Biset}
\ligneRav{1062}{}{ἀλλ᾽ ὦ τέκνον σὲ μὲν \difference{τὸ σαυτῆς} χρὴ ποιεῖν,}{\textbf{1062} τοσαύτης R : \lecon{τὸ σαυτῆς} Dindorf (1839)}
\ligneRav{1063}{—}{\difference{κλαίειν} ἐλεινῶς. : σὲ δ᾽ ἐπικλαίειν ὕστερον·}{\textbf{1063} \lecon{κλαίειν} et \lecon{ἐπικλαίειν} R : κλάειν et ἐπικλάειν Brunck}
\ligneRav{}{\auDessus{εὐ}{ρ}}{ἐμοὶ μελήσει ταῦτά γ᾽. ἀλλ᾽ ἄρχου λόγων·}{}
\ligneRav{1065}{ὡς ἀνδρομ.}{\hspace{1.1em}ὦ νὺξ ἱερὰ}{\textbf{1065} ὁ μνησίλοχος ὡς ἀνδρομέδα in mg. \Rpc, \Mon}
\ligneRav{}{}{\hspace{1.1em}ὡς μακρὸν ἵππευμα διώκεις}{}
\ligneRav{}{}{\hspace{1.1em}ἀστεροειδέα νῶτα διφρεύουσ᾽}{}
\ligneRav{}{}{\hspace{1.1em}αἰθέρος ἱερᾶς}{}
\ligneRav{1069}{—}{\hspace{1.1em}τοῦ σεμνοτάτου δι᾽ Ὀλύμπου; : †δι᾽ Ὀλύμπου;†}{\textbf{1069} \lecon{δι᾽ Ὀλύμπου} R : †δι᾽† ὀλυμπίας Bibi ap. Hesych. (ο 649)}
\ligneRav{1070}{μ\auDessus{ν}{Η}}{\hspace{1.1em}τί ποτ᾽ Ἀνδρομέδα \difference{περίαλλα} κακῶν}{\textbf{1070} περὶ ἄλλα R : \lecon{περίαλλα} Ellebodius}
\ligneRav{1071}{—}{\hspace{1.1em}μέρος ἐξέλαχον; : \difference{μόρος ἔξελ᾽ ἄχος}·}{\textbf{1071} μέρος ἐξέλαχον; R : \lecon{μόρος ἔξελ᾽ ἄχος.} Bibi}
\ligneRav{1072}{μ\auDessus{ν}{Η}}{\hspace{1.1em}θανάτου τλήμων : \difference{†}θανάτου τλήμων\difference{†}·}{\textbf{1072} \lecon{†θανάτου τλήμων†} (θανάτου τλήμων R : θαλάττης τέμνον Bibi)}
\ligneRav{1073}{μ\auDessus{ν}{Η}}{\hspace{1.1em}ἀπολεῖς μ᾽ ὦ γραῦ \difference{στωμυλλομένη}·}{\textbf{1073} στωμυλλωμένη R : \lecon{στωμυλλομένη} Grynaeus}
\ligneRav{1074}{\auDessus{εὐ}{ρ}}{\hspace{1.1em}\difference{στωμυλιομενής}·}{\textbf{1074} στωμυλωμένη R : \lecon{στωμυλιομενής} Bibi}
\ligneRav{1075}{μ\auDessus{ν}{Η}}{\hspace{1.1em}νὴ Δί᾽, ὀχληρά γ᾽ εἰσήρρηκας}{}
\ligneRav{}{—}{\hspace{1.1em}λίαν. : †λίαν†·}{}
\ligneRav{}{μ\auDessus{ν}{Η}}{\hspace{1.1em}ὦγάθ᾽ ἔασόν με μονῳδῆσαι,}{}
\ligneRav{1078}{—}{\hspace{1.1em}καὶ χαριεῖ μοι παῦσαι. : \difference{᾽παΰσαι}·}{\textbf{1078} (Μνη.) χαριεῖ μοι παῦσαι --- (Εὐρ.) παῦσαι R : \lecon{(Μνη.) χαριεῖ μοι παῦσαι --- (Εὐρ.) ᾽παΰσαι} Bibi}
\ligneRav{1079}{μ\auDessus{ν}{Η}}{\hspace{1.1em}βάλλ᾽ ἐς κόρακας. : \difference{βάλλε σύ κόρακας}·}{\textbf{1079} βάλλ᾽ ἐς κόρακας R : \lecon{βάλλε σύ κόρακας} Bibi}
\ligneRav{1080}{μ\auDessus{ν}{Η}}{\hspace{1.1em}†\difference{τί τὸ κακόν}†; : †τί τὸ κακόν;† : ληρεῖς. : †ληρεῖς†·}{\textbf{1080} \lecon{τί τὸ κακόν} bis R : τί κακόν bis Bentley}
\ligneRav{}{μ\auDessus{ν}{Η}}{\hspace{1.1em}οἴμωζ᾽. : †οἴμωζ᾽†·}{}
}
\vspace{3pt}{\color{ruleGold}\hrule height 0.6pt}\vspace{3pt}
{\fontsize{8.0pt}{9.4pt}\selectfont
\noindent{\bfseries\scshape\color{rubricRed}Loci paralleli et testimonia (f. 175r, vv. 1044–1081) :} \textbf{1049} (παρουσίᾳ) Moeris π 14 \espTest \textbf{1059} Suidas η 686 (ἐπικοκκάστρια codd. GM\expo{a.c.} vel ἐπικοκάστρια cod. I vel ἐπικαγχάστρια cod. A vel ἐπικαχάστρια codd. FV\expo{mg.} vel ἐπικοστράτ textu V) \espTest (ἀντῳδὸς ἐπικοκκάστρια) \Eustathe 1761.26 ἐκτίθεται καινοφώνους λέξεις οἷον [ ... ] καὶ τὸ ἐπικοκκάζειν ὕβρει ἤ κατά τινας ἐπιμηκάζειν, ὃθεν παρὰ Ἀριστοφάνει τὸ `ἀοιδὸς ἐπικοκκάστρια' \espTest \textbf{1066-8} (ὡς μακρὸν -- αἰθέρος) Suidas ι 531 {\policeFr ∼} Zonaras 1117 \espTest \textbf{1067} (νῶτα) Moeris ν 6 \espTest \textbf{1067-8} Suidas αι 122 {\policeFr ∼} Zonaras 67, 78 (sed ἀστρο- pro ἀστερο-)\par\vspace{2.5pt}
{\color{ruleGold}\hrule height 0.35pt}\vspace{2.5pt}
\noindent{\bfseries\scshape\color{rubricRed}Suite (en commentaire) de l'apparat critique et paléographique :} \textbf{Témoin papyrologique :} \textbf{\PsiQuatre{} (\italique{P.Oxy.} LXXIII 4935, fr. 1)} couvre les \textbf{vv. 1043–1051}. --- \textbf{Deux indications scéniques marginales de \Rpc{} (copiées par \Mon{}) :} au \textbf{v.~1056}, \lecon{ευρ\auDessus{ι}{π} ηχω} (Εὐριπίδης Ἠχώ) ; au \textbf{v.~1065}, \lecon{ὁ μνησίλοχος ὡς ἀνδρομέδα}. --- \textbf{Apport du papyrus \PsiQuatre{} (vv. 1047 et 1051) :} Au v. 1047, \PsiQuatre{} lit [ι]ω μοιρ[ας sans la dittographie μοι de R (ἰώ μοι μοίρας) ; au v. 1051, \PsiQuatre{} confirme βάρβαρον de R (défendu par Alphonse Willems et L. Hauseux). --- \textbf{Cinq restitutions et conjectures de Bibi dans le dialogue avec Écho (vv. 1069, 1071, 1072, 1074, 1078, 1079) :} 1° \textbf{v.~1071} μέρος ἐξέλαχον; \R{} → \lecon{μόρος ἔξελ' ἄχος.} Bibi ; 2° \textbf{v.~1072} θανάτου τλήμων \R{} (conj. \lecon{θαλάττης τέμνον} Bibi) ; 3° \textbf{vv.~1073–1074} \R{} écrit στωμυλλωμένη avec deux λ au v.~1073 mais στωμυλωμένη avec un seul λ au v.~1074, dissymétrie que Bibi explique en restituant \lecon{στωμυλιομενής} ; 4° \textbf{v.~1078} (Εὐρ.) παῦσαι \R{} → \lecon{(Εὐρ.) 'παΰσαι} Bibi ; 5° \textbf{v.~1079} βάλλ' ἐς κόρακας \R{} → \lecon{βάλλε σύ κόρακας} Bibi. --- \textbf{v. 1044 :} \lecon{κροκόεντ'} R (κροκόεν τόδ' Bergk) \espApp ἐνέδυσεν R : \lecon{ἀμφέδυσεν} Hermann ap. {\policeGr Σ}. --- \textbf{v. 1047 :} [ι]ω μοιρ[ας \PsiQuatre{} : ἰώ μοι μοίρας R : \lecon{ἰὼ μοίρας} Hermann \espApp δαίμων R : \lecon{δαῖμον} Burges \espApp ἀνέτικτε R (médivision et épenthèse de ν) : \lecon{ἄτεγκτε} Ellebodius (ἅν ἔτικτε Casaubon). --- \textbf{v. 1051 :} \lecon{βάρβαρον} \PsiQuatre{}, R, Willems. --- \textbf{v. 1054 :} δαιμόνων R, {\policeGr Σ} (δαιμονῶν Brunck ; δαιμόνι' amicus Fritzschii) : \lecon{διὰ μένων} Alphonse Willems (adopté par L. Hauseux). --- \textbf{v. 1056 :} \lecon{ευρ\auDessus{ι}{π} ηχω} in mg. \Rpc{}, M. --- \textbf{v. 1065 :} \lecon{ὁ μνησίλοχος ὡς ἀνδρομέδα} in mg. \Rpc{}, M. --- \textbf{v. 1069 :} \lecon{δι' Ὀλύμπου} R (†δι'† ὀλυμπίας Bibi ap. Hésychius ο 649). --- \textbf{v. 1071 :} μέρος ἐξέλαχον; R : \lecon{μόρος ἔξελ' ἄχος.} Bibi. --- \textbf{v. 1072 :} \lecon{θανάτου τλήμων} R (θαλάττης τέμνον Bibi conj. in app.). --- \textbf{v. 1074 :} στωμυλωμένη R : \lecon{στωμυλιομενής} Bibi. --- \textbf{v. 1078 :} (Εὐρ.) παῦσαι R : \lecon{(Εὐρ.) 'παΰσαι} Bibi. --- \textbf{v. 1079 :} βάλλ' ἐς κόρακας R : \lecon{βάλλε σύ κόρακας} Bibi.\par}
\end{minipage}\end{adjustbox}
\newpage
\renewcommand{\currHeadLE}{SCHOLIES MARGINALES ({\policeGr Σ}\expo{R}) \& GLOSES INTERLINÉAIRES · FOLIO 26/30}
\renewcommand{\currHeadRE}{f. 175r (Classense 177r) · vv. 1044–1081 · {\policeGr Σ} 329 à {\policeGr Σ} 343}
\renewcommand{\currHeadLO}{SCHOLIES MARGINALES ({\policeGr Σ}\expo{R}) \& GLOSES INTERLINÉAIRES · FOLIO 26/30}
\renewcommand{\currHeadRO}{f. 175r (Classense 177r) · vv. 1044–1081 · {\policeGr Σ} 329 à {\policeGr Σ} 343}
\bandeauScholies{Folio 26/30 — f. 175r}{Classense 177r}{vv. 1044–1081}{Recto (page de droite) · 15 entrées ({\policeGr Σ} 329 à {\policeGr Σ} 343)}
{\fontsize{8.8pt}{10.5pt}\selectfont\noindent\textbf{\color{rubricRed}Disposition topographique des marges sur le f. 175r :} Marge extérieure à droite ; marge intérieure à gauche. — \textit{Fac-similé de référence :} \texttt{folios\_individuels\_ff162v-177r/26\_f175r\_classense177r\_vv1044-1081.jpg}.\par\vspace{3pt}}
\entreeScholie{329}{1044}{ὃς ἐμὲ κροκόεντ᾽ ἐνέδυσεν·}{Interlinéaire (supra lineam en haut du f. 175r, v. 1044)}{ἀντὶ τοῦ κροκωτὸν ἀμφέδυσεν.}{ἀν\auDessus{τ}{ι} τ\expo{οῦ} κροκωτ(ὸν) ἀμφέδυσ(εν) {\policeFr ÷}}{« κροκόεντ' ἐνέδυσεν : au sens de “m'a revêtu d'une robe safran” (κροκωτὸν ἀμφέδυσεν). »}
\commScholie{Cité par Louis Hauseux (Bibi, v. 1044, vv 642-fin-bis.tex, l. 4161) : le papyrus Ψ\expo{4} (POxy 4935) recoupe exactement ce passage.}
\entreeScholie{330}{1050}{εἴθε με πυρφόρος αἰθέρος ἀστὴρ / τὸν βάρβαρον ἐξολέσειεν·}{Marge droite (extérieure, f. 175r, en face du v. 1050)}{ἴσως τὸν ἥλιον ἀστέρα ἐκάλεσεν, ἢ μᾶλλον τὸν κεραυνόν.}{ἴσ\expo{ως} τ(ὸν) ἥλι\expo{ον} ἀστέρα ἐκάλεσ(εν), ἢ μᾶλλ\expo{ον} τ(ὸν) κεραυνόν {\policeFr ÷}}{« Peut-être a-t-il appelé “astre porteur de feu de l'éther” (πυρφόρος αἰθέρος ἀστήρ) le soleil, ou plutôt la foudre (τὸν κεραυνόν). »}
\commScholie{Cité et discuté avec verve par Alphonse Willems et Louis Hauseux (Annexes.tex, ll. 1279-1310) : dans l'Andromède d'Euripide (fr. 122 K.), l'héroïne souhaitait être foudroyée (τὸν δύσμορον ἐξολέσειεν) ; Mnésiloque détourne le souhait de foudre sur la tête du garde scythe (τὸν βάρβαρον ἐξολέσειεν, leçon de R confirmée par le papyrus Ψ\expo{4} = POxy 4935 et maintenue par Willems et Bibi !).}
\entreeScholie{331}{1051}{τὸν βάρβαρον ἐξολέσειεν.}{Marge intérieure gauche (variante διχῶς, f. 175r, v. 1051)}{διχῶς, τὸν ἄθλιον.}{διχ\expo{ῶς}, τ(ὸν) ἄθλι\expo{ον} {\policeFr ÷}}{« Il y a deux leçons dans les manuscrits (διχῶς) : τὸν βάρβαρον (“le barbare”) et τὸν δύσμορον (glosé ici τὸν ἄθλιον, “le malheureux”, qui était le mot d'Euripide). »}
\commScholie{Remarquable note critique alexandrine utilisant le terme technique d'Aristarque διχῶς (« double leçon ») : l'original d'Euripide portait τὸν δύσμορον (glosé τὸν ἄθλιον), qu'Aristophane a transformé παρ' ὑπόνοιαν en τὸν βάρβαρον (confirmé par R et par le papyrus du IIe s. Ψ\expo{4} = POxy 4935 !).}
\entreeScholie{332}{1054}{λαιμότμητ᾽ ἄχη διὰ μένων [δαιμόνων R, Σ] αἰόλαν / νέκυσιν ἐπὶ πορείαν·}{Marge droite (extérieure) et glose déplacée σκοτεινήν (f. 175r, v. 1054)}{λείπει ἔσχον, ἐκ τῶν δαιμόνων τὰ ἄχη ἔσχον. σκοτεινήν.}{λείπ(ει) ἔσχ\expo{ον}, ἐκ τ\expo{ῶν} δαιμόν\expo{ων} τὰ ἄχ(η) ἔσχ\expo{ον} {\policeFr ÷} // σκοτεινήν {\policeFr ÷}}{« Il manque le verbe ἔσχον : “j'ai reçu des dieux ces tourments qui me tranchent la gorge”. — σκοτεινήν (“ténébreuse”). »}
\commScholie{Découverte philologique majeure de Louis Hauseux (Bibi, vv. 1041 et 1054, vv 642-fin-bis.tex, ll. 4119-4135) : la glose σκοτεινήν (« obscure, ténébreuse »), égarée au v. 1054 dans R\expo{2} et rattachée à tort par la Souda (αι 245) à l'adjectif αἰόλαν, expliquait en réalité au v. 1041 l'épithète rarissime d'Hadès restituée au texte : ψεφαυγοῦς (de ψέφας, « ténèbres ») !}
\entreeScholie{333}{1056}{χαῖρ᾽, ὦ φίλη παῖ· τὸν δὲ πατέρα Κηφέα / ... Ἠχὼ λόγων ἀντῳδὸς ἐπικοκκάστρια.}{Marge droite (extérieure, f. 175r, en face des vv. 1056–1059)}{ὑποκρίνεται Εὐριπίδης τὸ πρόσωπον τῆς Ἠχοῦς. πατὴρ γὰρ τῆς Ἀνδρομέδας ὁ Κηφεύς.}{ὑποκρίνετ(αι) Εὐριπίδ(ης) τ(ὸ) πρόσωπ\expo{ον} τ\expo{ῆς} Ἠχοῦς · πατ(ὴρ) γ(ὰρ) τ\expo{ῆς} Ἀνδρομέδ\expo{ας} ὁ Κηφεύς {\policeFr ÷}}{« Euripide joue ici le rôle de la nymphe Écho ; et Céphée (roi d'Éthiopie) est le père d'Andromède. »}
\commScholie{Rubrique dramaturgique expliquant la seconde entrée d'Euripide (en nymphe Écho).}
\entreeScholie{334}{1059}{Ἠχὼ λόγων ἀντῳδὸς ἐπικοκκάστρια,}{Marge droite (extérieure, f. 175r, grande scholie alexandrine de 6 lignes en face du v. 1059)}{εἰωθυῖα γελᾶν, γελάστρια. ἐπεὶ εἰσήγαγε κακοστένακτον τὴν Ἠχὼ ὁ Εὐριπίδης ἐν τῇ Ἀνδρομέδᾳ, εἰς τοῦτο παίζει. ἐζήλωσε δὲ αὐτὸν Πτολεμαῖος ὁ Φιλοπάτωρ ἐν ᾗ πεποίηκε τραγῳδίᾳ Ἀδώνιδι, περὶ ἧς ὁ ἐρώμενος αὐτῷ Ἀγαθοκλῆς γέγραφεν, ὁ ἀδελφὸς τῆς ἐρωμένης αὐτοῦ πάλιν Ἀγαθοκλείας.}{εἰωθυῖα γελᾶν, γελάστρια · ἐπεὶ εἰσήγαγ(ε) κακοστένακτ\expo{ον} τ(ὴν) Ἠχὼ ὁ Εὐριπίδ(ης) ἐν τῇ Ἀνδρομέδᾳ, εἰς τοῦτ(ο) παίζ(ει) · ἐζήλωσ(ε) δ(ὲ) αὐτ(ὸν) Πτολεμαῖ\expo{ος} ὁ Φιλοπάτ(ωρ) ἐν ᾗ πεποίηκ(ε) τραγῳδίᾳ Ἀδώνιδι, π(ερὶ) ἧς ὁ ἐρώμεν\expo{ος} αὐτῷ Ἀγαθοκλ\expo{ῆς} γέγραφ(εν), ὁ ἀδελφ(ὸς) τ\expo{ῆς} ἐρωμέν\expo{ης} αὐτ\expo{οῦ} πάλ(ιν) Ἀγαθοκλεί\expo{ας} {\policeFr ÷}}{« ἐπικοκκάστρια : qui a l'habitude de rire et de railler, moqueuse (γελάστρια, qui caquette en écho). Comme Euripide, dans son Andromède, avait mis en scène la nymphe Écho gémissant sur les malheurs d'autrui (κακοστένακτον), Aristophane s'en moque ici. Et le roi Ptolémée IV Philopator (221-204 av. J.-C.) a imité cet effet d'Euripide dans la tragédie d'Adonis (TrGF 119 F 1) qu'il a composée, pièce sur laquelle son favori Agathocle a écrit un commentaire, lui qui était le frère d'Agathocléia, la maîtresse du roi. »}
\commScholie{Scholie alexandrine extraordinaire (longuement analysée par Louis Hauseux, vv 642-fin-bis.tex, ll. 4277-4292, et reprise par la Souda ε 2406) : non seulement elle nous révèle un détail unique de l'histoire littéraire lagide (la tragédie Adonis de Ptolémée IV Philopator commentée par son mignon Agathocle), mais l'épithète κακοστένακτον éclaire tout le jeu parodique d'Écho !}
\entreeScholie{335}{1060}{ἥπερ πέρυσιν ἐν τῷδε ταὐτῷ χωρίῳ}{Marge droite (extérieure, f. 175r, en face du v. 1060)}{ἐπεὶ πέρυσιν ἐδιδάχθη ἡ Ἀνδρομέδα.}{ἐπεὶ πέρυσ(ιν) ἐδιδάχθ(η) ἡ Ἀνδρομέδα {\policeFr ÷}}{« (“Qui l'an dernier, en ce même lieu du théâtre de Dionysos, ai partagé les plaintes d'Euripide”) : parce que l'Andromède d'Euripide a été représentée l'année précédente (en 412 av. J.-C.). »}
\commScholie{Cité par Louis Hauseux (Introduction.tex, l. 176, et vv 642-fin-bis.tex, l. 4297) : cinquième et décisive confirmation de la date de 411 av. J.-C. pour les Thesmophories (8e année avant la représentation des Grenouilles en 405, cf. Σ Grenouilles v. 53).}
\entreeScholie{336}{1064}{ἐμοὶ μελήσει ταῦτά γ᾽. ἀλλ᾽ ἄρχου λόγων·}{Marge droite (extérieure, f. 175r, en face du v. 1064)}{ἠρέμα ὁ κηδεστὴς τοῦτο πρὸς τὸν Εὐριπίδην.}{ἠρέμα ὁ κηδεστ\expo{ὴς} τοῦτ(ο) πρ(ὸς) τ(ὸν) Εὐριπίδ(ην) {\policeFr ÷}}{« Le parent par alliance (Mnésiloque) dit cela à voix basse (en aparté, ἠρέμα) à Euripide. »}
\commScholie{Discuté par Louis Hauseux (vv 642-fin-bis.tex, l. 4309) : dans le texte, c'est plutôt Euripide-Écho qui prie Mnésiloque de chanter sa monodie (σε μονῳδῆσαι) pour qu'elle puisse lui renvoyer l'écho.}
\entreeScholie{337}{1065}{{\policeFr [rubrique :]} ὁ Μνησίλοχος ὡς Ἀνδρομέδα — ὦ νὺξ ἱερά, / ὡς μακρὸν ἵππευμα διώκεις}{Intra-textuel (rubrique en semi-onciale de R\expo{2} sur deux lignes) et Marge droite (Σ 1065, f. 175r)}{ὁ Μνησίλοχος ὡς Ἀνδρομέδα. τοῦ προλόγου Ἀνδρομέδας εἰσβολή.}{ὁ Μνησίλοχ\expo{ος} ὡς Ἀνδρομέδα {\policeFr ÷} // τ\expo{οῦ} προλόγ\expo{ου} Ἀνδρομέδ\expo{ας} εἰσβολή {\policeFr ÷}}{« Mnésiloque dans le rôle d'Andromède. — C'est le début même (l'incipit, εἰσβολή) du prologue de l'Andromède d'Euripide (fr. 114 Kannicht). »}
\commScholie{Cité par Louis Hauseux (vv 642-fin-bis.tex, l. 4326) : grâce à cette scholie (confirmée par Σ Grenouilles v. 53), nous savons que l'Andromède d'Euripide s'ouvrait non par un monologue iambique mais par cette célèbre monodie anapestique d'Andromède attachée à son rocher dans la nuit étoilée !}
\entreeScholie{338}{1066}{ὡς μακρὸν ἵππευμα διώκεις}{Interlinéaire (supra lineam au-dessus d'ἵππευμα, f. 175r, v. 1066)}{πορείαν.}{πορεί(αν) {\policeFr ÷}}{« Course en char (le trajet nocturne du char de la Nuit à travers l'Olympe étoilé). »}
\commScholie{Glose interlinéaire de R\expo{2} expliquant la métaphore tragique ἵππευμα (« course de char équestre » de la Nuit sacrée à travers la voûte étoilée, dans l'ouverture de l'Andromède) par le terme courant πορείαν (« trajet, course »).}
\entreeScholie{339}{1069}{δι᾽ Ὀλύμπου.}{Marge intérieure gauche (f. 175r, en face du v. 1069)}{τοῦτο ἡ ἠχὼ λέγει.}{τοῦτ(ο) ἡ ἠχὼ λέγ(ει) {\policeFr ÷}}{« Ceci (la répétition de δι' Ὀλύμπου) est prononcé par Écho. »}
\commScholie{Précise le début du jeu d'écho entre Mnésiloque et Euripide.}
\entreeScholie{340}{1070}{τί ποτ᾽ Ἀνδρομέδα περίαλλα κακῶν / μέρος ἐξέλαχον,}{Marge droite (extérieure, f. 175r, en face des vv. 1070–1071)}{καὶ τοῦτο ἐκ τοῦ προλόγου.}{\textbf{ς} τοῦτ(ο) ἐκ τ\expo{οῦ} προλόγ\expo{ου} {\policeFr ÷}}{« Ces vers aussi sont tirés du prologue de l'Andromède (fr. 115 Kannicht). »}
\commScholie{Indication de source littéraire de R\expo{2} (faisant suite à Σ 1065) précisant que les anapestes plaintifs τί ποτ᾽ Ἀνδρομέδα περίαλλα κακῶν μέρος ἐξέλαχον (vv. 1070–1071 = fr. 115 Kannicht) appartiennent eux aussi au prologue de l'Andromède d'Euripide.}
\entreeScholie{341}{1072}{θανάτου τλήμων.}{Marge droite (extérieure, f. 175r, en face du v. 1072)}{λείπει μέλλουσα τυχεῖν.}{λείπ(ει) μέλλουσα τυχ(εῖν) {\policeFr ÷}}{« Dans la citation tronquée (θανάτου τλήμων, “malheureuse, promise à la mort”), il manque la fin du vers d'Euripide : μέλλουσα τυχεῖν (“sur le point de rencontrer la mort”). »}
\commScholie{Cité par Louis Hauseux (vv 642-fin-bis.tex, l. 4423) : cette précieuse scholie nous restitue la fin authentique du vers d'Euripide (θανάτου τλήμων μέλλουσα τυχεῖν, fr. 115, 2 Kannicht), interrompu par la réplique d'Écho !}
\entreeScholie{342}{1073}{ἀπολεῖς μ᾽, ὦ γραῦ, στωμυλλομένη.}{Interlinéaire (supra lineam au-dessus de στωμυλλομένη, f. 175r, v. 1073)}{φλυαροῦσα.}{φλυαροῦσα {\policeFr ÷}}{« Par ton bavardage incessant (en jacassant). »}
\commScholie{Glose lexicale de R\expo{2} expliquant le participe familier στωμυλλομένη (« jacassant, bavardant sans fin ») adressé par Mnésiloque à Écho par φλυαροῦσα ; reprise par la Souda σ 1154.}
\entreeScholie{343}{1077}{ὦγάθ᾽, ἔασόν με μονῳδῆσαι,}{Interlinéaire (supra lineam au-dessus d'ὦγάθ', f. 175r, v. 1077)}{τὸ πλῆρες ὠγαθή.}{τ(ὸ) πλῆρ\expo{ες} ὠγαθή {\policeFr ÷}}{« La forme pleine (élidée ici devant ἔασον) est le féminin ὦ ἀγαθή (“ô ma bonne”, puisque Mnésiloque s'adresse à la vieille Écho). »}
\commScholie{Remarque grammaticale très fine : l'élision ὦγάθ' devant voyelle représente ici le vocatif féminin ὦ ἀγαθή et non le masculin ὦ ἀγαθέ !}
\clearpage
\renewcommand{\currHeadLE}{ΑΡΙΣΤΟΦΑΝΟΥΣ ΘΕΣΜΟΦΟΡΙΑΖΟΥΣΑΙ · CODEX RAVENNAS 429 (\R) · FOLIO 27/30}
\renewcommand{\currHeadRE}{f. 175v (Classense 177v) · vv. 1082–1127 · Collation diplomatique}
\renewcommand{\currHeadLO}{ΑΡΙΣΤΟΦΑΝΟΥΣ ΘΕΣΜΟΦΟΡΙΑΖΟΥΣΑΙ · CODEX RAVENNAS 429 (\R) · FOLIO 27/30}
\renewcommand{\currHeadRO}{f. 175v (Classense 177v) · vv. 1082–1127 · Collation diplomatique}
\noindent\begin{adjustbox}{max totalheight=30.4cm, max width=\textwidth, keepaspectratio}\begin{minipage}{\textwidth}
\bandeauRav{f. 175v}{Classense 177v}{1082–1127}{Folio 27/30 · Verso (page de gauche) · 18 scholies \& gloses ({\policeGr Σ} 344–361)}
{\fontsize{8.4pt}{9.7pt}\selectfont
\ligneRav{1082}{μ\auDessus{ν}{Η}}{\hspace{1.1em}ὀτότυζ᾽. : †ὀτότυζ᾽†·}{\textbf{1082} \lecon{ὀτότυζ᾽} R : ὀλόφυγ᾽ Bibi}
\ligneRav{1083}{τοξ/}{\hspace{1.1em}\difference{οὖτος} τί λαλ\inRas{εῖ}{ῖ}ς; \hfill \textbf{[Ἠχώ :]} \difference{οὖτος} τί λαλ\inRas{εῖ}{ῖ}ς;}{\textbf{1083} οῦτος οὗτος R : \lecon{οὖτος} bis Brunck \espApp λαλεῖς bis \Rac : \lecon{λαλῖς} bis \Rdeux, Σ}
\ligneRav{1084}{τοξ/}{\hspace{1.1em}πρυτάνεις καλέσω· \hfill \textbf{[Ἠχώ :]} πρυτάνεις καλέσω·}{}
\ligneRav{1085}{τοξ/}{\hspace{1.1em}σί κακόν; \hfill \textbf{[Ἠχώ :]} σί κακόν;}{\textbf{1085} σί bis R, sch. : τί Coulon}
\ligneRav{1086}{τοξ/}{\hspace{1.1em}\difference{†}π\inRas{ρ}{}ῶτε\difference{†} τὸ πωνή; \hfill \textbf{[Ἠχώ :]} \difference{†}π\inRas{ρ}{}ῶτε\difference{†} τὸ πωνή;}{\textbf{1086} π\inRas{ρ}{}ωτε bis \Rpc : πωτε bis \Mon}
\ligneRav{1087}{τοξ/}{\hspace{1.1em}σὺ λαλ\auDessus{ῖ}{εῖ}ς; \hfill \textbf{[Ἠχώ :]} σὺ λαλ\auDessus{ῖ}{εῖ}ς;}{\textbf{1087} \lecon{λαλῖς} bis R : λαλεῖς bis \Rpc}
\ligneRav{1088}{τοξ/}{\hspace{1.1em}\difference{κλαύσαι}· \hfill \textbf{[Ἠχώ :]} \difference{κλαύσαι}·}{\textbf{1088} κλαύσαιμι bis R : κλαύσει bis Brunck : \lecon{κλαύσαι} bis Rogers : κλαῦσ' ἔτι bis Austin}
\ligneRav{1089}{τοξ/}{\hspace{1.1em}\difference{κἀκκάσκι} μοι; \hfill \textbf{[Ἠχώ :]} \difference{†κἀκκάσκι μοι;†}}{\textbf{1089} κακκασκι bis R : κἀγκάσκις bis van Leeuwen (<ς> add. Fritzsche) : κἀγκάσκι bis Coulon : alii alia}
\ligneRav{1090}{χορ/}{\hspace{1.1em}μὰ Δί᾽ ἀλλὰ γυνὴ πλησίον αὕτη·}{\textbf{1090} (Μνη.) edd. plerique : \lecon{(Χορ.)} Bibi, cf. Σ}
\ligneRav{1091}{\auDessus{εὐ}{ρ}}{\hspace{1.1em}\difference{πλήξασ᾿ αὕτη}·}{\textbf{1091} πλησίον αὕτη R : \lecon{πλήξασ᾿ αὕτη} Bibi}
\ligneRav{1092}{τοξ/}{\hspace{1.1em}ποῦ \difference{᾿στ᾽ ἠ} μιαρά; : καὶ δὴ φεύγει·}{\textbf{1092} 'σθ' ἡ R : \lecon{᾿στ᾽ ἠ} Brunck \espApp (Μνη.) edd. plerique : \lecon{(Χορ.)} Bibi}
\ligneRav{1093}{τοξ/}{\hspace{1.1em}ποῖ ποῖ \difference{πεύγεις}; : \difference{<ποῖ ποῖ πέτει>}·}{\textbf{1093} φεύγεις R : \lecon{πεύγεις} Enger \espApp (Εὐρ.) ποῖ ποῖ πέτει om. R : \lecon{suppl.} Bibi}
\ligneRav{1094}{τοξ/}{\hspace{1.1em}οὐ καιρήσεις. : \difference{οὐχ αἱρήσεις}·}{\textbf{1094} (Τοξ.) οὐ καιρήσεις om. R : \lecon{suppl.} Brunck ap. Σ \espApp οὐκ᾽ αἰρήσεις; R : \lecon{οὐχ αἱρήσεις.} Wilamowitz (\italique{Kleine Schriften} IV, p.486n.1)}
\ligneRav{1095}{τοξ/}{\hspace{1.1em}ἔτι γὰρ γρύζεις; : †ἔτι γὰρ γρύζεις;†}{}
\ligneRav{1096}{τοξ/}{\hspace{1.1em}λαβὲ \difference{τὴ} μιαρά. : †λαβὲ τὴ μιαρά.†}{\textbf{1096} τῇ et τῆ R : \lecon{τὴ} bis Bothe}
\ligneRav{}{τοξ/}{\hspace{1.1em}λάλο καὶ κατάρατο γύναικο·}{}
\ligneRav{1098}{ὡς περσευς}{ὦ θεοὶ τίν᾽ εἰς γῆν βαρβάρων ἀφίγμεθα}{\textbf{1098} ευρ\auDessus{ι}{π} ως περσευς in mg. add. \Rpc, \Mon}
\ligneRav{}{}{ταχεῖ πεδίλῳ; διὰ μέσου γὰρ αἰθέρος}{}
\ligneRav{1100}{}{τέμνων κέλευθον πόδα τίθημ᾽ ὑπόπτερον}{}
\ligneRav{1101}{}{Περσεὺς πρὸς Ἄργος \difference{ναυστολῶν} τὸ Γοργόνος}{\textbf{1101} ναυτολῶν R : \lecon{ναυστολῶν} Canini (\italique{navigans} iam Divus)}
\ligneRav{1102}{—}{\textlbrackdbl{\policeFr in marg. inf. \Rdeux{} : }κάρα κομίζων : τί λέγι; τὴ \difference{Γόργος} πέρι\textrbrackdbl}{\textbf{1102} Γοργόνος R : \lecon{Γόργος} Σ, Bothe}
\ligneRav{1103}{—}{\textlbrackdbl{\policeFr in marg. inf. \Rdeux{} : }τὸ γραμματέο σὺ τὴ \difference{κεπαλή}; : τὴν Γοργόνος\textrbrackdbl}{\textbf{1103} κεφαλή R : \lecon{κεπαλή} Bothe}
\ligneRav{}{—}{ἔγωγε φημί. : Γόργο τοι κἀγὼ λέγι·}{}
\ligneRav{1105}{\auDessus{εὐ}{ρ}}{ἔα· τίν᾽ ὄχθον τόνδ᾽ ὁρῶ καὶ παρθένον}{}
\ligneRav{}{}{θεαῖς ὁμοίαν ναῦν ὅπως ὡρμισμένην;}{}
\ligneRav{1107}{κ\auDessus{η}{δ} ἀνδρ.}{ὦ ξένε \difference{κατοίκτιρόν} με τὴν παναθλίαν,}{\textbf{1107} κ\auDessus{η}{δ} ως ανδρομ\auDessus{ε}{δ} in mg. \Rpc, \Mon \espApp κατοικτειρόν R : \lecon{κατοίκτιρόν} van Herwerden}
\ligneRav{1108}{—}{λῦσόν με δεσμῶν. : οὐκὶ μὶ λαλῆσι σύ;}{\textbf{1108} \lecon{μὶ λαλῆσι} R : μὴ λαλῆσι Brunck}
\ligneRav{1109}{}{κατάρατο τολμᾷς ἀποτανουμένη \difference{λαλῖς};}{\textbf{1109} λαλᾶς R : \lecon{λαλῖς;} Bergk (cf. vv. 1082 et 1087) : λαλᾷς; Hall-Geldart}
\ligneRav{1110}{\auDessus{εὐ}{ρ}}{ὦ παρθέν᾽ \difference{οἰκτίρω} σὲ κρεμαμένην ὁρῶν·}{\textbf{1110} οἰκτείρω R : \lecon{οἰκτίρω} van Herwerden}
\ligneRav{}{τοξ/}{οὐ παρτέν᾽ ἐστίν, ἀλλ᾽ ἀμαρτωλὴ γέρων}{}
\ligneRav{}{—}{καὶ κλέπτο καὶ πανοῦργο. : ληρεῖς ὦ Σκύθα·}{}
\ligneRav{}{}{αὕτη γάρ ἐστιν Ἀνδρομέδα παῖς Κηφέως·}{}
\ligneRav{1114}{τοξ/}{σκέψαι τὸ \difference{κύστο}· \difference{μή τι} \difference{μικκὸν} παίνεται;}{\textbf{1114} σκύτο R : \lecon{κύστο} Scaliger grat. metr. : σῦκο Sommerstein \espApp μῆτι R : \lecon{μή τι} \Mon (\italique{numquid} Divus) \espApp μικτὸν R : \lecon{μικκὸν} Bentley : μικρὸ Enger}
\ligneRav{1115}{\auDessus{εὐ}{ρ}}{φέρε \difference{δεῦρό} μοι τὴν χεῖρ᾽, ἵν᾽ ἅψωμαι κόρης·}{\textbf{1115} δεῦρο δεῦρο R : \lecon{δεῦρο} Biset}
\ligneRav{}{}{φέρε Σκύθ᾽· ἀνθρώποισι γὰρ νοσήματα}{}
\ligneRav{}{}{ἅπασίν ἐστιν· ἐμὲ δὲ καὐτὸν τῆς κόρης}{}
\ligneRav{}{—}{ταύτης ἔρως εἴληφεν. : οὐ ζηλῶσί σε·}{}
\ligneRav{1119}{}{ἀτὰρ εἰ \difference{τὸ πρωκτὸ} δεῦρο περιεστραμμένον,}{\textbf{1119} τῶ πρωκτῶ R : \lecon{τὸ πρωκτὸ} Biset}
\ligneRav{1120}{}{οὐκ \difference{ἐπτόνησά σ᾽} αὐτὸ πυγίζεις ἄγων·}{\textbf{1120} ἐπτόνησας R : \lecon{ἐπτόνησά σ᾽} Bentley}
\ligneRav{}{\auDessus{εὐ}{ρ}}{τί δ᾽ οὐκ ἐᾷς λύσαντά μ᾽ αὐτὴν ὦ Σκύθα}{}
\ligneRav{1122}{}{πεσεῖν \difference{ἐς} εὐνὴν καὶ γαμήλιον λέχος;}{\textbf{1122} ἐστ᾽ R : ἐς τ᾽ \Mon : \lecon{ἐς} Grynaeus}
\ligneRav{}{τοξ/}{εἰ σπόδρ᾽ ἐπιτυμεῖς τὴ γέροντο πύγισο,}{}
\ligneRav{1124}{}{τὴ σανίδο τρήσας \difference{ἐξόπιστο} πρώκτισον·}{\textbf{1124} ἐξόπισθο R : \lecon{ἐξόπιστο} Daubuz}
\ligneRav{1125}{\auDessus{εὐ}{ρ}}{μὰ Δί᾽ ἀλλὰ λύσω \difference{δεσμά}. : μαστιγῶ σ᾽ ἄρα·}{\textbf{1125} δέμας R : \lecon{δεσμά} Grynaeus}
\ligneRav{1126}{\auDessus{εὐ}{ρ}}{καὶ μὴν ποιήσω τοῦτο : \difference{τὸ κεπαλή σ᾽} (τὸ κ\inRas{αὶ }{ε}παλῆς) ἄρα}{\textbf{1126} τὸ κεπαλῆς (ex καὶ παλῆς) \Rpc, \Mon : \lecon{τὸ κεπαλή σ᾽} Brunck}
\ligneRav{1127}{}{τὸ ξιπομάκαιραν \difference{ἀποκεκόψι} (ἀποκεκόψ\inRas{᾽}{οι}) τουτοΐ·}{\textbf{1127} ἀποκέκοιψ᾽ \Rac : ἀποκεκόψοι \Rpc : ἀποκεκόψο \Mon : \lecon{ἀποκεκόψι} Brunck}
}
\vspace{3pt}{\color{ruleGold}\hrule height 0.6pt}\vspace{3pt}
{\fontsize{8.0pt}{9.4pt}\selectfont
\noindent{\bfseries\scshape\color{rubricRed}Loci paralleli et testimonia (f. 175v, vv. 1082–1127) :} \textbf{1125} (δεσμά) Moeris δ 28\par\vspace{2.5pt}
{\color{ruleGold}\hrule height 0.35pt}\vspace{2.5pt}
\noindent{\bfseries\scshape\color{rubricRed}Suite (en commentaire) de l'apparat critique et paléographique :} \textbf{Omission des vv. 1102–1103 par \Run{} et ajout par \Rdeux{} dans la marge inférieure :} Le copiste \Run{} avait sauté les \textbf{vv. 1102–1103} (Γόργος πέρι λέγις τὸ γραμματέο / σὺ τὴν κεπαλή; — omis également par \Mon{} qui n'a pas regardé le bas de page) ; le réviseur \textbf{\Rdeux{} les a rajoutés tout en bas du folio dans la marge inférieure}, avec un signe de renvoi. --- \textbf{Analyse critique des vv. 1085–1086 par Louis Hauseux (\textit{vv 642-fin-bis.tex}, ll. 4520, 4588–4612) :} 1° Au \textbf{v.~1085}, \R{} et {\policeGr Σ}\expo{R} (ἀντὶ τοῦ 'σοὶ κακόν') donnent \lecon{σί κακόν;} bis, que Bibi défend et conserve contre τί κακόν; (Coulon) ; 2° au \textbf{v.~1086}, \Rac{} avait écrit deux fois πρῶτε et une main postérieure a gratté les deux ρ (\lecon{π\inRas{ρ}{}ωτε} bis, copié πωτε bis par \Mon{}) d'après {\policeGr Σ}\expo{R} (πόθεν ἡ φωνή) : Louis Hauseux (Bibi) imprime \lecon{†πῶτε† τὸ πωνή;} bis en démontrant que l'Archer scythe ne confond jamais ο et ω. --- \textbf{Deux notes scéniques marginales de \Rpc{} :} Au \textbf{v. 1098}, \lecon{ευρ\auDessus{ι}{π} ως περσευς} ; au \textbf{v. 1107}, \lecon{κ\auDessus{η}{δ} ως ανδρομ\auDessus{ε}{δ}}. --- \textbf{Prothèse de σ au v. 1114 :} R écrit \lecon{σκύτο} pour \lecon{κύστο} (Scaliger, = κύσθον en prononciation scythe), parallèle exact invoqué par Louis Hauseux pour sa conjecture τὸν κύσθον au v. 557. --- \textbf{v. 1082 :} \lecon{ὀτότυζ'} R (ὀλόφυγ' Bibi conj. in app.). --- \textbf{v. 1083 :} οῦτος οὗτος R : \lecon{οὖτος} bis Brunck \espApp λαλεῖς bis \Rac{} : \lecon{λαλῖς} bis \Rdeux{}, {\policeGr Σ}. --- \textbf{v. 1085 :} \textbf{σί κακόν; bis} R, {\policeGr Σ}\expo{R} (ἀντὶ τοῦ 'σοὶ κακόν', défendu par Bibi contre τί κακόν; Coulon). --- \textbf{v. 1086 :} π\inRas{ρ}{}ωτε bis \Rpc{} (πωτε bis \Mon{}) : \lecon{†πῶτε†} bis Bibi (rejetant l'interprétation πόθεν de {\policeGr Σ}\expo{R}). --- \textbf{vv. 1090 et 1092 :} (Μνη.) edd. plerique : \lecon{(Χορ.)} Bibi, cf. {\policeGr Σ}\expo{R}. --- \textbf{v. 1091 :} πλησίον αὕτη R : \lecon{πλήξασ' αὕτη} Bibi. --- \textbf{v. 1093 :} φεύγεις R : \lecon{πεύγεις} Enger \espApp \lecon{<(Εὐρ.) ποῖ ποῖ πέτει>} om. R : \textbf{suppl. Bibi}. --- \textbf{v. 1094 :} <(Τοξ.) οὐχ αἱρήσεις> om. R : \textbf{suppl. Brunck} ap. {\policeGr Σ} \espApp οὐκ' αἰρήσεις; R : \lecon{οὐχ αἱρήσεις.} Wilamowitz (\italique{Kleine Schriften} IV, p. 486 n. 1). --- \textbf{v. 1098 :} \lecon{ευρ\auDessus{ι}{π} ως περσευς} in mg. add. \Rpc{}, M. --- \textbf{vv. 1102–1103 :} om. \Run{}, M : \textbf{in mg. inferiore add. \Rdeux{}} \espApp \textbf{v. 1102 :} Γοργόνος R : \lecon{Γόργος} {\policeGr Σ}, Bothe \espApp \textbf{v. 1103 :} κεφαλή R : \lecon{κεπαλή} Bothe. --- \textbf{v. 1107 :} \lecon{κ\auDessus{η}{δ} ως ανδρομ\auDessus{ε}{δ}} in mg. \Rpc{}, M \espApp κατοικτειρόν R : \lecon{κατοίκτιρόν} van Herwerden. --- \textbf{v. 1115 :} δεῦρο δεῦρο R (dittographie) : \lecon{δεῦρο} Biset. --- \textbf{v. 1125 :} δέμας R (métathèse de σ) : \lecon{δεσμά} Grynaeus. --- \textbf{v. 1127 :} ἀποκέκοιψ' \Rac{} (ἀποκεκόψοι \Rpc{}, ἀποκεκόψο \Mon{}) : \lecon{ἀποκεκόψι} Brunck.\par}
\end{minipage}\end{adjustbox}
\newpage
\renewcommand{\currHeadLE}{SCHOLIES MARGINALES ({\policeGr Σ}\expo{R}) \& GLOSES INTERLINÉAIRES · FOLIO 27/30}
\renewcommand{\currHeadRE}{f. 175v (Classense 177v) · vv. 1082–1127 · {\policeGr Σ} 344 à {\policeGr Σ} 361}
\renewcommand{\currHeadLO}{SCHOLIES MARGINALES ({\policeGr Σ}\expo{R}) \& GLOSES INTERLINÉAIRES · FOLIO 27/30}
\renewcommand{\currHeadRO}{f. 175v (Classense 177v) · vv. 1082–1127 · {\policeGr Σ} 344 à {\policeGr Σ} 361}
\bandeauScholies{Folio 27/30 — f. 175v}{Classense 177v}{vv. 1082–1127}{Verso (page de gauche) · 18 entrées ({\policeGr Σ} 344 à {\policeGr Σ} 361)}
{\fontsize{8.8pt}{10.5pt}\selectfont\noindent\textbf{\color{rubricRed}Disposition topographique des marges sur le f. 175v :} Marge extérieure à gauche ; marge intérieure à droite. — \textit{Fac-similé de référence :} \texttt{folios\_individuels\_ff162v-177r/27\_f175v\_classense177v\_vv1082-1127.jpg}.\par\vspace{3pt}}
\entreeScholie{344}{1083}{οὖτος, τί λαλῖς [λαλεῖς R\expo{ac}, λαλῖς R\expo{2}, Σ];}{Marge supérieure et marge gauche (extérieure, f. 175v, en haut du verso, Planche 23)}{χωρὶς τοῦ ε γράφεται· ὁ γὰρ Σκύθης βαρβαρίζει.}{χωρὶς τ\expo{οῦ} ε γράφετ(αι) · ὁ γ(ὰρ) Σκύθ(ης) βαρβαρίζ(ει) {\policeFr ÷}}{« Il faut écrire λαλῖς sans ε (par un simple ι, et non λαλεῖς) : car le Scythe écorche la langue grecque en barbare. »}
\commScholie{Cité par Louis Hauseux (Bibi, v. 1083, vv 642-fin-bis.tex, ll. 4566-4570) : sur le Ravennas, R\expo{1} avait écrit par automatisme orthographique λαλεῖς (bis) ; le réviseur R\expo{2} a corrigé en λαλῖς et inscrit cette scholie qui prouve le soin philologique des éditeurs alexandrins pour noter la prononciation barbare du Scythe !}
\entreeScholie{345}{1085}{σί κακόν;}{Interlinéaire (supra lineam au-dessus de σί κακόν, f. 175v, v. 1085)}{ἀντὶ τοῦ σοὶ κακόν.}{ἀν\auDessus{τ}{ι} τ\expo{οῦ} σοὶ κακόν {\policeFr ÷}}{« σί κακόν (prononcé par le Scythe et répété par Écho) : au lieu de σοὶ κακόν (“malheur à toi ! / la peste soit de toi !”). »}
\commScholie{Argument décisif de Louis Hauseux (Bibi, v. 1085, vv 642-fin-bis.tex, ll. 4520 et 4589) contre la « forfaiture » des éditeurs modernes (comme Coulon) qui corrigent σί κακόν de R\expo{1} en τί κακόν : la leçon σί κακόν est garantie à la fois par R\expo{1} et par la glose expresse de {\policeGr Σ}\expo{R} ἀντὶ τοῦ σοὶ κακόν !}
\entreeScholie{346}{1086}{†πῶτε† τὸ πωνή;}{Marge gauche (extérieure, f. 175v, en face du v. 1086)}{πόθεν ἡ φωνή; θαυμάζει τὴν ἠχώ.}{πόθ(εν) ἡ φωνή; θαυμάζ(ει) τ(ὴν) ἠχώ {\policeFr ÷}}{« “D'où vient cette voix (πόθεν ἡ φωνή) ?” Le Scythe s'ébahit devant le phénomène de l'écho. »}
\commScholie{Analysé en détail par Louis Hauseux (Bibi, v. 1086, vv 642-fin-bis.tex, ll. 4602-4615) : R\expo{1} avait écrit πρωτε (bis), dont les deux ρ ont été grattés postérieurement (π[ρ]ωτε) pour s'accorder avec l'interprétation du scholiaste (πόθεν) ; mais comme le Scythe ne confond jamais ο et ω, L. Hauseux appose les cruces †πῶτε†.}
\entreeScholie{347}{1089}{κἀκκάσκι [κακκάσκι R] μοι;}{Interlinéaire (supra lineam au-dessus de κακκάσκι μοι, f. 175v, v. 1089)}{καταγελᾷς μοι.}{καταγελᾷς μοι {\policeFr ÷}}{« κακκάσκι μοι (= κἀγχάσκεις μοι / καταχάσκεις μου) : “tu te moques de moi à gorge déployée, tu me ris au nez ?” »}
\commScholie{Cité par Louis Hauseux (vv 642-fin-bis.tex, l. 4645) avec le parallèle des Guêpes v. 721.}
\entreeScholie{348}{1090}{μὰ Δί᾽, ἀλλὰ γυνὴ πλησίον αὕτη.}{Marge gauche (extérieure, f. 175v, en face du v. 1090)}{τοῦτο ὁ χορὸς ἢ ὁ γέρων.}{τοῦτ(ο) ὁ χορ(ὸς) ἢ ὁ γέρ(ων) {\policeFr ÷}}{« Cette réplique (“Non par Zeus, mais c'est cette femme tout près d'ici !”) est dite soit par le Chœur (la Coryphée), soit par le vieillard (Mnésiloque). »}
\commScholie{Adopté par Louis Hauseux (Bibi, vv. 1090 et 1092, vv 642-fin-bis.tex, ll. 4655-4665) : contre la plupart des éditeurs modernes qui attribuent la réplique à Mnésiloque (lequel n'a aucun intérêt à dénoncer son sauveur Euripide !), Bibi suit la première indication de {\policeGr Σ}\expo{R} (τοῦτο ὁ χορός) et donne les vv. 1090 et 1092 à la Coryphée !}
\entreeScholie{349}{1094}{οὐ καιρήσεις. — οὐχ αἱρήσεις·}{Marge gauche (extérieure, f. 175v, en face du v. 1094)}{ἀντὶ τοῦ οὐ χαιρήσεις.}{ἀν\auDessus{τ}{ι} τ\expo{οῦ} οὐ χαιρήσ(εις) {\policeFr ÷}}{« οὐ καιρήσεις (prononcé par le Scythe sans aspiration) : au sens de οὐ χαιρήσεις (“tu ne t'en tireras pas à bon compte, tu vas le payer !”). »}
\commScholie{Miracle de conservation indirecte mis en lumière par Brunck et Louis Hauseux (Bibi, v. 1094, vv 642-fin-bis.tex, ll. 4281-4288 et 4701-4710) : alors que R\expo{1} avait sauté par haplographie la menace du Scythe <οὐ καιρήσεις> devant la réponse d'Écho οὐκ αἱρήσεις (« tu ne m'attraperas pas ! »), la glose de {\policeGr Σ}\expo{R} ἀντὶ τοῦ οὐ χαιρήσεις nous restitue à la lettre près le texte disparu de R\expo{1} !}
\entreeScholie{350}{1095}{ἔτι γὰρ γρύζεις;}{Interlinéaire (supra lineam au-dessus de γρύζεις, f. 175v, v. 1095)}{λαλεῖς.}{λαλεῖς {\policeFr ÷}}{« Tu oses encore souffler mot (grogner, parler) ? »}
\commScholie{Glose lexicale de R\expo{2} expliquant le verbe expressif γρύζεις (proprement « grogner, pousser le moindre son “γρῦ” ») dans la menace de l'Archer ἔτι γὰρ γρύζεις; par λαλεῖς (« tu parles encore ? »).}
\entreeScholie{351}{1096}{λαβὲ τὴ μιαρά. — †λαβὲ τὴ μιαρά.†}{Interlinéaire (supra lineam, f. 175v, v. 1096)}{κάτεχε.}{κάτεχ(ε) {\policeFr ÷}}{« Arrête-la, saisis-la ! »}
\commScholie{Glose interlinéaire de R\expo{2} expliquant l'impératif λαβέ (« saisis, arrête ! ») lancé par l'Archer scythe qui cherche la voix insaisissable d'Écho par κάτεχε (« retiens-la ! »).}
\entreeScholie{352}{1097}{λάλο καὶ κατάρατο γύναικο·}{Interlinéaire (supra lineam au-dessus de τὴν μιαράν, f. 175v, v. 1097)}{τὴν λάλον καὶ κατάρατον γυναῖκα.}{τ(ὴν) λάλον \textbf{ς} κατάρατ\expo{ον} γυναῖκα {\policeFr ÷}}{« Cette femme bavarde et maudite. »}
\commScholie{Glose grammaticale de R\expo{2} rétablissant en grec attique correct (τὴν λάλον καὶ κατάρατον γυναῖκα) le sabir barbare de l'Archer scythe qui confond les désinences au neutre/masculin (λάλο καὶ κατάρατο γύναικο, v. 1097).}
\entreeScholie{353}{1098}{{\policeFr [rubrique :]} Εὐριπίδης ὡς Περσεύς — ὦ θεοὶ τίν᾽ εἰς γῆν βαρβάρων ἀφίγμεθα / ταχεῖ πεδίλῳ;}{Intra-textuel (rubrique de R\expo{2}) et Marge gauche (extérieure, f. 175v, en face du v. 1098)}{εἰς Περσέα ἐξ Ἀνδρομέδας τρία τὰ πρῶτα. καὶ λοιπὸν ἐπέζευξε τὰ ἑξῆς.}{εὐρ\auDessus{ι}{π} ὡς περσεύς {\policeFr ÷} // εἰς Περσέα ἐξ Ἀνδρομέδ\expo{ας} τρία τὰ πρῶτ(α) · \textbf{ς} λοιπ(ὸν) ἐπέζευξ(ε) τὰ ἑξ\expo{ῆς} {\policeFr ÷}}{« Euripide sous le costume de Persée. — Ces trois premiers vers (vv. 1098-1100) appartiennent à l'entrée en scène de Persée dans l'Andromède d'Euripide (fr. 124 Kannicht) ; et le poète y a raccordé la suite. »}
\commScholie{Cité par Louis Hauseux (vv 642-fin-bis.tex, l. 4740) : identifie avec précision le fragment 124 de l'Andromède d'Euripide.}
\entreeScholie{354}{1099}{ταχεῖ πεδίλῳ;}{Interlinéaire (supra lineam au-dessus de ταχεῖ πεδίλῳ, f. 175v, v. 1099)}{ἐλέγετο γὰρ ὁ Περσεὺς εἰς τοὺς πόδας ἔχειν τὰ πτερά.}{ἐλέγετ(ο) γ(ὰρ) ὁ Περσεὺς εἰς τ\expo{οὺς} πόδ\expo{ας} ἔχ(ειν) τὰ πτερά {\policeFr ÷}}{« (“D'une sandale rapide”) : car selon le mythe, Persée portait aux pieds les sandales ailées (d'Hermès). »}
\commScholie{Glose mythologique et scénique de R\expo{2} expliquant l'expression tragique ταχεῖ πεδίλῳ (« d'une sandale rapide », v. 1099 = fr. 124 Kannicht) par les sandales ailées (τὰ πτερά) que porte Persée (ici incarné par Euripide).}
\entreeScholie{355}{1103}{τί λέγι; τὴ Γόργος πέρι / τὸ γραμματέο σὺ τὴ κεπαλή; — τὴν Γοργόνος}{Marge inférieure du f. 175v (sous les vv. 1102–1103 rétablis en marge inférieure par R\expo{2} avec signe de renvoi, Planche 20)}{τί λέγει περὶ τῆς κεφαλῆς τοῦ Γόργου τοῦ γραμματέως, ἐπεὶ προεῖπεν ἐκεῖνος «Γοργόνος κάρα κομίζων.» ὁ δὲ Γόργος γραμματεὺς, ἀλλὰ καὶ βάρβαρος.}{τί λέγ(ει) π(ερὶ) τ\expo{ῆς} κεφαλ\expo{ῆς} τ\expo{οῦ} Γόργ\expo{ου} τ\expo{οῦ} γραμματέ\expo{ως}, ἐπεὶ προεῖπ(εν) ἐκεῖν\expo{ος} «Γοργόν\expo{ος} κάρα κομίζ\expo{ων}» · ὁ δ(ὲ) Γόργ\expo{ος} γραμματεὺς, ἀλλὰ \textbf{ς} βάρβαρ\expo{ος} {\policeFr ÷}}{« “Que dit-il au sujet de la tête de Gorgos le greffier (τοῦ Γόργου τοῦ γραμματέως) ?”, parce qu'Euripide-Persée vient de dire (au v. 1101) : “apportant la tête de la Gorgone (Γοργόνος κάρα κομίζων)”. Et ce Gorgos était un greffier (secrétaire du Conseil), mais aussi d'origine barbare. »}
\commScholie{Double intérêt paléographique et philologique majeur (Louis Hauseux, Bibi, vv. 1102-1103, vv 642-fin-bis.tex, ll. 4753-4765, et Planche 20 du Manuel) : 1° les vv. 1102-1103, sautés par R\expo{1} dans la colonne par homéotéleute sur Γοργόνος (et absents de M !), ont été rétablis par R\expo{2} dans la marge inférieure du f. 175v ; 2° la scholie de R\expo{2} conserve le nom exact Γόργος (adopté par Bibi) là où le texte marginal de R\expo{2} avait écrit Γοργόνος.}
\entreeScholie{356}{1105}{ἔα· τίν᾽ ὄχθον τόνδ᾽ ὁρῶ καὶ παρθένον / θεαῖς ὁμοίαν ναῦν ὅπως ὡρμισμένην;}{Marge gauche (extérieure, f. 175v, en face des vv. 1105–1106)}{πάλιν ἐξ Ἀνδρομέδας «ἔα, τίν’ ὄχθον τόνδ’ ὁρῶ περίρρυτον ἀφρῷ θαλάσσης, παρθένου τ’ εἰκώ τινα.»}{πάλ(ιν) ἐξ Ἀνδρομέδ\expo{ας} «ἔα, τίν' ὄχθ\expo{ον} τόνδ' ὁρῶ περίρρυτ\expo{ον} ἀφρῷ θαλάσσ\expo{ης}, παρθέν\expo{ου} τ' εἰκώ τινα» {\policeFr ÷}}{« De nouveau tiré de l'Andromède d'Euripide (fr. 125 Kannicht) : “Ah ! quel est ce tertre rocheux que je vois tout entouré de l'écume de la mer, et cette image d'une vierge...” »}
\commScholie{Cité par Louis Hauseux (vv 642-fin-bis.tex, ll. 4777-4790) : conserve les vers originaux prononcés par Persée en apercevant Andromède enchaînée à son rocher.}
\entreeScholie{357}{1109}{οὐκὶ μὶ λαλῆσι σύ; / κατάρατο τολμᾷς ἀποτανουμένη λαλῖς [λαλᾶς R];}{Interlinéaire (supra lineam, f. 175v, v. 1109)}{ἀποθανουμένη τολμᾷς λαλῆσαι.}{ἀποθανουμέν(η) τολμᾷς λαλῆσ(αι) {\policeFr ÷}}{« Alors que tu vas mourir, tu oses encore bavarder ? (glose de la menace du Scythe aux vv. 1108-1109). »}
\commScholie{Glose grammaticale de R\expo{2} traduisant en grec correct (ἀποθανουμένη τολμᾷς λαλῆσαι) la réplique de l'Archer scythe pleine de psilose et de barbarismes (κατάρατο τολμᾷς ἀποτανουμένη λαλῖς, v. 1109).}
\entreeScholie{358}{1114}{σκέψαι τὸ κύστο· μή τι μικκὸν [μικρὸν R] παίνεται;}{Marge gauche (extérieure, f. 175v, en face du v. 1114)}{δείκνυσιν αὐτῷ τὸ αἰδοῖον.}{δείκνυσ(ιν) αὐτῷ τ(ὸ) αἰδοῖ\expo{ον} {\policeFr ÷}}{« Le Scythe soulève la robe de Mnésiloque et montre à Euripide son membre viril (en disant : “Regarde donc cette vulve [τὸ κύστο = τὸν κύσθον] : te paraît-elle petite ?”). »}
\commScholie{Note de jeu scénique (commentée par L. Hauseux à propos du v. 557 et du v. 1114).}
\entreeScholie{359}{1118}{οὐ ζηλῶσί σε·}{Interlinéaire (supra lineam au-dessus d'οὐ ζηλῶσί σε, f. 175v, v. 1118)}{οὐ ζηλῶ σε, φησὶν, εἰ ἐρᾷς.}{οὐ ζηλῶ σε, φησ(ὶν), εἰ ἐρᾷς {\policeFr ÷}}{« “Je ne t'envie pas”, dit le Scythe (οὐ ζηλῶσί σε pour οὐ ζηλῶ σε), “si tu es amoureux de cette créature !” »}
\commScholie{Explique le barbarisme morphologique du Scythe (3e pers. plur. ζηλῶσι employée pour la 1re pers. sing. ζηλῶ).}
\entreeScholie{360}{1119}{ἀτὰρ εἰ τὸ πρωκτὸ δεῦρο περιεστραμμένον, / οὐκ ἐπτόνησά σ᾽ αὐτὸ πυγίζεις ἄγων·}{Marge gauche (extérieure, f. 175v, en face des vv. 1119–1120)}{εἰ μὴ τὸν νῶτον ἦν, φησὶ, πρὸς τῇ σανίδι, ἀλλὰ πρὸς ἡμᾶς ἐτέτραπτο, οὐκ ἄν σοι ἐφθόνησα ἀπαγαγόντι περανεῖν.}{εἰ μὴ τ(ὸν) νῶτ\expo{ον} ἦν, φησ(ὶ), πρ(ὸς) τῇ σανίδι, ἀλλὰ πρ(ὸς) ἡμᾶς ἐτέτραπτ(ο), οὐκ ἄν σοι ἐφθόνησ(α) ἀπαγαγόντ(ι) περανεῖν {\policeFr ÷}}{« “Si son dos n'était pas plaqué contre la planche du carcan”, dit le Scythe, “mais que son derrière fût tourné de notre côté, je ne t'aurais pas refusé de l'emmener pour le besogner (περανεῖν) !” »}
\commScholie{Traduit en grec attique la boutade grivoise du Scythe (commentée par L. Hauseux, vv 642-fin-bis.tex, ll. 4828-4850).}
\entreeScholie{361}{1127}{τὸ κεπαλή σ᾽ ἄρα / τὸ ξιπομάκαιραν [τῇ ξιπομάκαιρᾳ R] ἀποκεκόψι τουτοΐ·}{Marge inférieure du f. 175v (en face du v. 1127)}{ἔλεγε τὸ ξιφομάχαιρα ὡς δρεπανομάχαιρα.}{ἔλεγ(ε) τ(ὸ) ξιφομάχαιρα ὡς δρεπανομάχαιρα {\policeFr ÷}}{« On disait ξιφομάχαιρα (“sabre-coutelas”) sur le modèle de δρεπανομάχαιρα (“épée-faucille”, la harpè de Persée). »}
\commScholie{Repris par la Souda ξ 74 : le Scythe oppose son propre cimeterre (ξιφομάχαιρα) à la harpè courbe de Persée.}
\clearpage
\renewcommand{\currHeadLE}{ΑΡΙΣΤΟΦΑΝΟΥΣ ΘΕΣΜΟΦΟΡΙΑΖΟΥΣΑΙ · CODEX RAVENNAS 429 (\R) · FOLIO 28/30}
\renewcommand{\currHeadRE}{f. 176r (Classense 178r) · vv. 1128–1168 · Collation diplomatique}
\renewcommand{\currHeadLO}{ΑΡΙΣΤΟΦΑΝΟΥΣ ΘΕΣΜΟΦΟΡΙΑΖΟΥΣΑΙ · CODEX RAVENNAS 429 (\R) · FOLIO 28/30}
\renewcommand{\currHeadRO}{f. 176r (Classense 178r) · vv. 1128–1168 · Collation diplomatique}
\noindent\begin{adjustbox}{max totalheight=30.4cm, max width=\textwidth, keepaspectratio}\begin{minipage}{\textwidth}
\bandeauRav{f. 176r}{Classense 178r}{1128–1168}{Folio 28/30 · Recto (page de droite) · 3 scholies \& gloses ({\policeGr Σ} 362–364)}
{\fontsize{9.3pt}{10.8pt}\selectfont
\ligneRav{1128}{\auDessus{εὐ}{ρ}}{αἰαῖ· τί δράσω; πρὸς τίνας στρεφθῶ λόγους;}{}
\ligneRav{1129}{}{ἀλλ᾽ οὐκ \difference{ἂν ἐνδέξαιτο} βάρβαρος φύσις·}{\textbf{1129} ἂν δέξαιτο R : ἂν ἀνδέξαιτο Ellebodius : \lecon{ἂν ἐνδέξαιτο} Küster}
\ligneRav{1130}{}{σκαιοῖσι γάρ τοι καινὰ προσφέρων σοφὰ}{}
\ligneRav{}{}{μάτην ἀναλίσκοις ἄν, ἀλλ᾽ ἄλλην τινὰ}{}
\ligneRav{1132}{}{\difference{τούτῳ} πρέπουσαν μηχανὴν προσοιστέον·}{\textbf{1132} τοῦτο R : \lecon{τούτῳ} Zanetti}
\ligneRav{1133}{τοξ/}{μιαρὸς ἀλώπηξ, \difference{οἶον} \difference{ἐπιτήκιζί} μοι·}{\textbf{1133} οἷον R : \lecon{οἶον} Brunck \espApp ἐπιτηκίζει R : \lecon{ἐπιτήκιζί} anonymus Parisinus}
\ligneRav{}{μ\auDessus{ν}{Η}}{μέμνησο Περσεῦ μ᾽ ὡς καταλείπεις ἀθλίαν·}{}
\ligneRav{1135}{τοξ/}{ἔτι γὰρ σὺ \difference{τὴ μάστιγαν} (τῆ μαστ\inRas{η}{ι}γᾶν) ἐπιτυμεῖς \difference{λαβεῖν;}}{\textbf{1135} \lecon{τῆ μαστιγᾶν} (ex μαστηγαν) \Rpc \espApp λαβεῖν· R : \lecon{λαβεῖν;} Scaliger}
\ligneRav{}{χορ/}{\hspace{1.1em}Παλλάδα τὴν φιλόχορον ἐμοὶ}{}
\ligneRav{1138}{}{\hspace{1.1em}δεῦρο καλεῖν νόμος \difference{εἰς} χορόν,}{\textbf{1138} ἐς R : \lecon{εἰς} von Velsen}
\ligneRav{1139}{}{\hspace{1.1em}παρθένον ἄζυγα \difference{κούρην},}{\textbf{1139} κόρην R : \lecon{κούρην} Bothe}
\ligneRav{1140}{}{\hspace{1.1em}ἣ πόλιν ἡμετέραν ἔχει}{}
\ligneRav{}{}{\hspace{1.1em}καὶ κράτος φανερὸν μόνη}{}
\ligneRav{}{}{\hspace{1.1em}κλῃδοῦχός τε καλεῖται·}{}
\ligneRav{}{}{\hspace{1.1em}φάνηθ᾽ ὦ τυράννους στυγοῦσ᾽ ὥσπερ εἰκός·}{}
\ligneRav{1145}{}{\hspace{1.1em}δῆμός τοί σε καλεῖ γυναι-}{}
\ligneRav{}{}{\hspace{1.1em} κῶν· ἔχουσα δέ μοι μόλοις}{}
\ligneRav{}{}{\hspace{1.1em}εἰρήνην φιλέορτον·}{}
\ligneRav{1148}{}{\hspace{1.1em}\difference{ἥκετέ τ᾽} εὔφρονες ἵλαοι,}{\textbf{1148} ἥκετ᾽ R : \lecon{ἥκετέ τ᾽} Enger : ἥκετε δ᾽ Hermann}
\ligneRav{}{}{\hspace{1.1em}πότνιαι, ἄλσος ἐς ὑμέτερον,}{}
\ligneRav{1150}{}{\hspace{1.1em}οὗ δὴ ἀνδράσιν οὐ \difference{θεμιτεύετον} εἰσορᾶν}{\textbf{1150-1} \lecon{οὗ δὴ ἀνδράσιν} R : οὗ δῆτ' ἀνδράσιν Burges : ἄνδρας ἵν(α) Hermann \espApp θεμιτὸν R : θέμις Burges : θεμιτ' Bothe : \lecon{θεμιτεύετον} Bibi}
\ligneRav{}{}{\hspace{1.1em}ὄργια σεμνὰ θεαῖν ἵνα λαμπάσι·}{}
\ligneRav{}{}{\hspace{1.1em}φαίνετον ἄμβροτον ὄψιν·}{}
\ligneRav{1155}{}{\hspace{1.1em}μόλετον, ἔλθετον, ἀντόμεθ᾽ ὦ}{}
\ligneRav{1156}{}{\hspace{1.1em}Θεσμοφόρω \difference{πολυποτνία},}{\textbf{1156} πολὺ πότνια R : \lecon{πολυποτνία} Hermann}
\ligneRav{}{}{\hspace{1.1em}εἰ καὶ πρότερόν ποτ᾽ ἐπηκόω ἤλθετον}{}
\ligneRav{1159}{}{\hspace{1.1em}νῦν ἀφίκεσθον †ἱκετεύομεν ἐνθάδ᾽ ἡμῖν†·}{\textbf{1159} \lecon{ἐνθάδ᾽ ἡμῖν} R : ἐνθάδε χἠμῖν Hermann}
\ligneRav{1160}{\auDessus{εὐ}{ρ}}{γυναῖκες εἰ βούλεσθε τὸν λοιπόν χρόνον}{}
\ligneRav{}{}{σπονδὰς ποιήσασθαι πρὸς ἐμέ, νυνὶ πάρα,}{}
\ligneRav{}{}{ἐφ᾽ ᾧτ᾽ ἀκοῦσαι μηδὲν ὑπ᾽ ἐμοῦ μηδαμὰ}{}
\ligneRav{}{}{κακὸν τὸ λοιπόν. ταῦτ᾽ ἐπικηρυκεύομαι·}{}
\ligneRav{}{χορ/}{χρείᾳ δὲ ποίᾳ τόνδ᾽ ἐπεσφέρεις λόγον;}{}
\ligneRav{1165}{\auDessus{εὐ}{ρ}}{ὅδ᾽ ἐστὶν \difference{οὑν} τῇ σανίδι κηδεστὴς ἐμός·}{\textbf{1165} οὖν R : \lecon{οὑν} Biset}
\ligneRav{1166}{}{ἢν οὖν κομίσωμαι τοῦτον, οὐδὲν μή ποτε}{\textbf{1166} \lecon{μή} R : μου Bentley}
\ligneRav{1167}{}{κακῶς \difference{ἀκούσητ᾽}· ἢν δὲ μὴ \difference{πίθησθέ} μοι,}{\textbf{1167} ἀκούσαιτ' R : \lecon{ἀκούσητ᾽} Elmsley (ad \italique{Ach.} v.295) \espApp πείθησθέ R : \lecon{πίθησθέ} Hirschig (ad \italique{Vesp} v.155)}
\ligneRav{}{}{ἃ νῦν ὑποικουρεῖτε τοῖσιν ἀνδράσιν}{}
}
\vspace{3pt}{\color{ruleGold}\hrule height 0.6pt}\vspace{3pt}
{\fontsize{8.0pt}{9.4pt}\selectfont
\noindent{\bfseries\scshape\color{rubricRed}Loci paralleli et testimonia (f. 176r, vv. 1128–1168) :} \textbf{1139} (ἄζυγα) \Phryn 100.3\par\vspace{2.5pt}
{\color{ruleGold}\hrule height 0.35pt}\vspace{2.5pt}
\noindent{\bfseries\scshape\color{rubricRed}Suite (en commentaire) de l'apparat critique et paléographique :} \textbf{Dittographie anticipatrice exponctuée après le v. 1143 :} Entre les vv. 1143 et 1144, \Rac{} avait écrit par anticipation déformée du v. 1144 (στυγοῦσ' ὥσπερ τὸν τύραννον) les mots \lecon{στυγνὰς ὡσσε} ; \Rpc{} les a supprimés (\italique{del.}), mais le copiste de \Mon{} les a néanmoins transcrits sous la forme \lecon{στυγνὰς ὡς ἐ}. --- \textbf{Glose interne et conjecture métrique de \textbf{Bibi} au v. 1151 :} R comporte la glose interne ὅπου (expliquant un ἵνα ou οὗ). Là où R porte \lecon{οὗ δὴ ἀνδράσιν οὐ θεμιτὸν} (déficient de deux syllabes dans la responsion lyrique), \textbf{Louis Hauseux (\textbf{Bibi}) conserve intact οὗ δὴ ἀνδράσιν de R et restitue le duel verbal οὐ θεμιτεύετον} (« où toutes deux [les Déesses] interdisent par leur loi sacrée aux hommes de voir leurs orgies... », haplographie de la finale -εύετον en -όν). --- \textbf{vv. 1150–1151 :} \lecon{οὗ δὴ ἀνδράσιν} R (avec glose interne ὅπου ; οὗ δῆτ' ἀνδράσιν Burges ; ἄνδρας ἵν(α) Hermann) \espApp θεμιτὸν R (θέμις Burges ; θεμιτ' Bothe) : \lecon{θεμιτεύετον} \lecon{Bibi}. --- \textbf{v. 1153 :} λαμπάσι R : \lecon{λαμπάσι.} \lecon{Bibi}.\par}
\end{minipage}\end{adjustbox}
\newpage
\renewcommand{\currHeadLE}{SCHOLIES MARGINALES ({\policeGr Σ}\expo{R}) \& GLOSES INTERLINÉAIRES · FOLIO 28/30}
\renewcommand{\currHeadRE}{f. 176r (Classense 178r) · vv. 1128–1168 · {\policeGr Σ} 362 à {\policeGr Σ} 364}
\renewcommand{\currHeadLO}{SCHOLIES MARGINALES ({\policeGr Σ}\expo{R}) \& GLOSES INTERLINÉAIRES · FOLIO 28/30}
\renewcommand{\currHeadRO}{f. 176r (Classense 178r) · vv. 1128–1168 · {\policeGr Σ} 362 à {\policeGr Σ} 364}
\bandeauScholies{Folio 28/30 — f. 176r}{Classense 178r}{vv. 1128–1168}{Recto (page de droite) · 3 entrées ({\policeGr Σ} 362 à {\policeGr Σ} 364)}
{\fontsize{8.8pt}{10.5pt}\selectfont\noindent\textbf{\color{rubricRed}Disposition topographique des marges sur le f. 176r :} Marge extérieure à droite ; marge intérieure à gauche. — \textit{Fac-similé de référence :} \texttt{folios\_individuels\_ff162v-177r/28\_f176r\_classense178r\_vv1128-1168.jpg}.\par\vspace{3pt}}
\entreeScholie{362}{1130}{αἰαῖ· τί δράσω; πρὸς τίνας στρεφθῶ λόγους; / ἀλλ᾽ οὐκ ἂν ἐνδέξαιτο βάρβαρος φύσις· / σκαιοῖσι γάρ τοι καινὰ προσφέρων σοφὰ / μάτην ἀναλίσκοις ἄν.}{Marge droite (extérieure, f. 176r, en haut du recto, en face des vv. 1130–1131)}{παρὰ τὰ ἐκ Μηδείας [298].}{παρὰ τὰ ἐκ Μηδεί\expo{ας} {\policeFr ÷}}{« Parodie des vers de la Médée d'Euripide (vv. 298-299 : σκαιοῖσι μὲν γὰρ καινὰ προσφέρων σοφὰ / δόξεις ἀχρεῖος κοὐ σοφὸς πεφυκέναι). »}
\commScholie{Identification exacte de la citation de la Médée d'Euripide (vv. 298-299).}
\entreeScholie{363}{1150}{οὗ δὴ ἀνδράσιν οὐ θεμιτεύετον [θεμιτὸν R] εἰσορᾶν / ὄργια σεμνὰ θεαῖν, ἵνα λαμπάσι}{Intra-textuel / interlinéaire (au-dessus de ἵνα, f. 176r, v. 1151)}{ὅπου.}{ὅπ\expo{ου} {\policeFr ÷}}{« ἵνα : au sens local de “là où” (ὅπου, dans l'enceinte sacrée où il n'est pas permis aux hommes de voir les mystères). »}
\commScholie{Cité par Louis Hauseux (vv 642-fin-bis.tex, ll. 5000-5003) avec les parallèles de la Συναγωγή ι 82, Hésychius ι 648, Photius ι 133 et la Souda ι 360.}
\entreeScholie{364}{1168}{ἃ νῦν ὑποικουρεῖτε τοῖσιν ἀνδράσιν}{Interlinéaire (supra lineam, bas du f. 176r, v. 1168)}{λάθρα ποιεῖτε.}{λάθρα ποιεῖτ(ε) {\policeFr ÷}}{« Ce que vous faites en cachette au logis (sans que nous le dénoncions désormais aux maris, promet Euripide). »}
\commScholie{Explique les conditions du traité de paix conclu entre Euripide et le Chœur des femmes.}
\clearpage
\renewcommand{\currHeadLE}{ΑΡΙΣΤΟΦΑΝΟΥΣ ΘΕΣΜΟΦΟΡΙΑΖΟΥΣΑΙ · CODEX RAVENNAS 429 (\R) · FOLIO 29/30}
\renewcommand{\currHeadRE}{f. 176v (Classense 178v) · vv. 1169–1210 · Collation diplomatique}
\renewcommand{\currHeadLO}{ΑΡΙΣΤΟΦΑΝΟΥΣ ΘΕΣΜΟΦΟΡΙΑΖΟΥΣΑΙ · CODEX RAVENNAS 429 (\R) · FOLIO 29/30}
\renewcommand{\currHeadRO}{f. 176v (Classense 178v) · vv. 1169–1210 · Collation diplomatique}
\noindent\begin{adjustbox}{max totalheight=30.4cm, max width=\textwidth, keepaspectratio}\begin{minipage}{\textwidth}
\bandeauRav{f. 176v}{Classense 178v}{1169–1210}{Folio 29/30 · Verso (page de gauche) · 15 scholies \& gloses ({\policeGr Σ} 365–379)}
{\fontsize{8.8pt}{10.2pt}\selectfont
\ligneRav{1169}{\auDessus{εὐ}{ρ}}{ἀπὸ τῆς στρατιᾶς παροῦσιν ὑμῶν διαβαλῶ·}{}
\ligneRav{1170}{χορ/}{τὰ μὲν παρ᾽ ἡμῶν ἴσθι σοι πεπεισμένα·}{}
\ligneRav{1171}{}{τὸν βάρβαρον δὲ τοῦτον αὐτὸς \difference{πεῖθε} σύ·}{\textbf{1171} πεῖσαι R : \lecon{πεῖθε} Biset}
\ligneRav{1172}{\auDessus{εὐ}{ρ}}{ἐμὸν \inRas{γ᾽}{} ἔργον ἐστίν· καὶ σὸν ὦλάφιον ἅ σοι}{\textbf{1172} Εὐριπίδης ἐν σχήματι προαγωγοῦ γραός Σ \espApp \lecon{ἐμὸν} \inRas{γ᾽}{} \Rdeux}
\ligneRav{1173}{}{καθ᾽ ὁδὸν \expo{ἔφραζον} (ἔφ\expo{ραζ}ον \Rpc) ταῦτα μεμνῆσθαι ποιεῖν·}{\textbf{1173} \lecon{ἔφραζον} (ex ἔφρον) \Rpc}
\ligneRav{1174}{}{πρῶτον μὲν οὖν δίελθε \difference{κἀνακάλπασον}·}{\textbf{1174} κἀνακόλπασον R : \lecon{κἀνακάλπασον} Hermann, cf. Σ, Suid. κ 623, Zonar. p.1171}
\ligneRav{1175}{}{σὺ δ᾽ ὦ \difference{Τερηδὸν} ἐπαναφύσα Περσικόν·}{\textbf{1175} Τερηδὼν R, Suid. τ 340 AGFM\expo{ac} : \lecon{Τερηδὸν} Suid. VM\expo{pc}}
\ligneRav{1176}{τοξ/}{τί τὸ βόμβο τοῦτο; \difference{κῶμό τις} ἀνεγεῖρί μοι;}{\textbf{1176} κωμοτίς ἀν- R : κωμοδρίαν Σ : \lecon{κῶμό τις ἀν-} van Velsen}
\ligneRav{}{\auDessus{εὐ}{ρ}}{ἡ παῖς ἔμελλε προμελετᾶν ὦ τοξότα·}{}
\ligneRav{1178}{}{ὀρχησομένη γὰρ \difference{ἔρχεθ᾽} ὡς ἄνδρας τινάς·}{\textbf{1178} ἔρχετ᾽ R : \lecon{ἔρχεθ᾽} Brunck}
\ligneRav{1179}{τοξ/}{ὀρκῆσι καὶ μελετῆσι, οὐ \difference{κωλῦσ᾽} ἐγώ·}{\textbf{1179} κωλύσ᾽ R : \lecon{κωλῦσ᾽} Blaydes (κωλῦσι = κωλύσω), cf. ad 1003, 1196, 1201, 1225}
\ligneRav{1180}{}{\difference{ὠς} ἐλαπρός, \difference{ὤσπερ} ψύλλο κατὰ τὸ κῴδιο·}{\textbf{1180} ὡς et ὥσπερ R : \lecon{ὠς} et \lecon{ὤσπερ} Brunck}
\ligneRav{}{\auDessus{εὐ}{ρ}}{φέρε θοἰμάτιον ἄνωθεν ὦ τέκνον τοδί·}{}
\ligneRav{1182}{}{καθιζομένη δ᾽ ἐπὶ \difference{τοῖσι} γόνασι τοῦ Σκύθου}{\textbf{1182} τοῖς R : \lecon{τοῖσι} Bentley}
\ligneRav{}{—}{τὼ πόδε πρότεινον, ἵν᾽ ὑπολύσω. : ναῖκι ναὶ}{}
\ligneRav{1184}{}{κάτησο κάτησο, ναίκι ναί, τυγάτριον·}{\textbf{1184} \lecon{ναίκι ναὶ} R : ναίκι ναίκι Bothe}
\ligneRav{1185}{}{οἴμ᾽ \difference{ὠς στέριπο} (ὡ σ\inRas{π}{τ}εριπὸ) τὸ τιττί᾽, \difference{ὤσπερ} \difference{γογγυλί}·}{\textbf{1185} ὡ στεριπὸ (ex σπεριπο) \Rpc : \lecon{ὡς στεριπὸ} (sic) \Mon \espApp ὥσπερ R : \lecon{ὤσπερ} Brunck \espApp \lecon{γογγυλί} \PsiTrois, Enger : γογγύλῃ R}
\ligneRav{}{\auDessus{εὐ}{ρ}}{αὔλει σὺ θᾶττον· ἔτι δέδοικας τὸν Σκύθην;}{}
\ligneRav{1187}{τοξ/}{καλό γε τὸ πυγή· \difference{κλαῦσί γ᾽} (κλαυσ\inRas{᾽ ε}{}ϊ γ᾽) \difference{ἢν} μὴ ᾽νδον μένῃς·}{\textbf{1187} κλαυσ\inRas{᾽ ε}{}ϊ γ᾽ \Rdeux : κλαυσεῖ γ᾽ \Mon \espApp ἂν R : \lecon{ἢν} Dindorf (1839)}
\ligneRav{}{}{\hspace{2.2em}\textbf{[παρεπιγραφή R\expo{1}, \PsiTrois{} :]} \lecon{ἀνακύπτη· καὶ παρακύπτι· ἀπεψωλημένος·}}{\textbf{post 1187} \lecon{ἀνακύπτη· καὶ παρακύπτι· ἀπεψωλημένος·} R, μ]ενο[ \PsiTrois}
\ligneRav{1188}{}{εἶεν· καλὴ τὸ σκῆμα περὶ τὸ \inRas{σ}{}πόστιον·}{\textbf{1188} τὸ \inRas{σ}{}ποστίον \Rpc : τὸ ποστίον \Mon}
\ligneRav{1189}{\auDessus{εὐ}{ρ}}{καλῶς ἔχει. λαβὲ θοἰμάτιον· ὥρα \difference{'στὶ} νῷν}{\textbf{1189} 'στιν R : \lecon{'στὶ} \Mon}
\ligneRav{1190}{—}{ἤδη βαδίζειν. : \difference{οὐκὶ} \difference{πιλῆσι} πρῶτά με;}{\textbf{1190} τί οὐκ ἐπιλήσει R : \lecon{οὐκὶ} (Biset) \lecon{πιλῆσι} Porson (ad \italique{Aristophanica} p.222), cf. Σ `οὐχὶ φιλήσειν'}
\ligneRav{1191}{\auDessus{εὐ}{ρ}}{πάνυ γε· φίλησον αὐτόν. : \difference{ὂ ὂ ὂ} παπαπαπαῖ,}{\textbf{1191} ὂ R : \lecon{ὂ ὂ ὂ} \Mon}
\ligneRav{1192}{}{\difference{ὠς} γλυκερὸ τὸ γλῶσσ᾽, \difference{ὤσπερ} Ἀττικὸς \difference{μέλι}·}{\textbf{1192} ὡς et ὥσπερ R : \lecon{ὠς} et \lecon{ὤσπερ} Brunck \espApp \lecon{μέλι} \PsiTrois, \Mon : μέλις R}
\ligneRav{}{—}{τί οὐ κατεύδει παρ᾽ ἐμέ; : χαῖρε τοξότα,}{}
\ligneRav{1194}{—}{οὐ γὰρ γένοιτ᾽ ἂν τοῦτο. : \difference{ναὶ ναὶ} \difference{γρᾴδιο}·}{\textbf{1194} ναὶ R : ναίκι Bentley : \lecon{ναὶ ναὶ} Brunck \espApp γράδιον R : \lecon{γρᾴδιο} Bothe}
\ligneRav{1195}{—}{ἐμοὶ \inRas{σὺ}{} \difference{κάρισο σὺ} τοῦτο : δώσεις οὖν δραχμήν;}{\textbf{1195} \inRas{σὺ}{} κάρισοσ οὐ \Rpc : κάρισο σοῦ \Mon : \lecon{κάρισο σὺ} Brunck}
\ligneRav{}{τοξ/}{ναὶ ναῖκι δῶσι. : τἀργύριον τοίνυν φέρε·}{}
\ligneRav{1197}{τοξ/}{ἀλλ᾽ οὐκ \difference{ἔκὠδέν}· ἀλλὰ τὸ συμβήνην λαβέ·}{\textbf{1197} έκ' ὠδέν R : ἐχω οὐδέν Suid. σ 1273, Photius σ 660 : \lecon{ἔκὠδέν} Bergk \espApp \lecon{συμβήνην} R : συβήνην Σ, Suid., Photius}
\ligneRav{1198}{γϋ. (Bon.)}{ἔπειτα \difference{κομίζι σ᾽ αὐτίκ᾽} (κ\inRas{ω}{ο}μίζεις αὐτοῖς)· \textbf{[σκ᾽.]} \difference{ἀκολούτι} τέκνον·}{\textbf{1198} κ\inRas{ω}{ο}μίζεις αὐτοῖς \Rpc : \lecon{(Τοξότ.) κομίζι σ᾽ αὐτίκ᾽} Bibi \espApp γϋ. ante v. 1198 (Boninus) et σκ᾽. ante ἀκολ. in R \espApp ἀκουλούτι R : \lecon{ἀκολούτι} \Mon}
\ligneRav{}{}{σὺ δὲ τοῦτο τήρει τὴ γέροντο, γρᾴδιο·}{}
\ligneRav{1200}{—}{ὄνομα δέ σοι τί ἔστιν; : Ἀρτεμισία·}{}
\ligneRav{1201}{τοξ/}{μεμνῆσι τοίνυν τοὔνομ᾽· \difference{Ἀρταμουξία}·}{\textbf{1201} ἀρτομουξία R : \lecon{Ἀρταμουξία} ed. Lipsiensis (1814)}
\ligneRav{}{\auDessus{εὐ}{ρ}}{Ἑρμῆ δόλιε ταυτὶ μὲν ἔτι καλῶς ποιεῖς·}{}
\ligneRav{1203}{}{σὺ μὲν οὖν ἀπότρεχε παιδάριον τουτὶ λαβών·}{\textbf{1203} \lecon{τουτὶ} R : ταυτὶ Dobree, cf. Σ}
\ligneRav{}{}{ἐγὼ δὲ λύσω τόνδε. σὺ δ᾽ ὅπως ἀνδρικῶς}{}
\ligneRav{1205}{}{ὅταν λυθῇς τάχιστα φεύξει καὶ τενεῖς}{}
\ligneRav{}{}{ὡς τὴν γυναῖκα καὶ τὰ παιδί᾽ οἴκαδε·}{}
\ligneRav{}{μ\auDessus{ν}{Η}}{ἐμοὶ μελήσει ταῦτά γ᾽ ἢν ἅπαξ λυθῶ·}{}
\ligneRav{1208}{\auDessus{εὐ}{ρ}}{λέλυσο\inRas{ν}{}· σὸν ἔργον, φεῦγε πρὶν τὸν τοξότην}{\textbf{1208} ca. 6 litt. ]· σ\douteux{ο}[νεργον \PsiQuatre : \lecon{λέλυσο}\inRas{ν}{} \Rpc : λέλυσαι Bentley}
\ligneRav{}{—}{ἥκοντα καταλαβεῖν. : ἐγὼ δὴ τοῦτο δρῶ·}{}
\ligneRav{1210}{τοξ/}{ὦ γρᾴδι᾽ ὠς καρίεντό σοι τὸ \difference{τυγά\inRas{θ}{τ}ριο},}{\textbf{1210} τυγά\inRas{θ}{τ}ριον \Rpc (ex τυγάθριον \Rac) : τυγάτριον \Mon : \lecon{τυγάτριο} Bibi}
}
\vspace{3pt}{\color{ruleGold}\hrule height 0.6pt}\vspace{3pt}
{\fontsize{8.0pt}{9.4pt}\selectfont
\noindent{\bfseries\scshape\color{rubricRed}Loci paralleli et testimonia (f. 176v, vv. 1169–1210) :} \textbf{1169} (στρατιᾶς) Photius σ 608 στρατιάν \espTest \textbf{1174} (κἀνακάλπασον) Photius α 1514 ἀνακαλπάζει ; cf. Zonaras 1171 καλπάζειν· τὸ ἄκρως βαδίζειν {\policeFr ≈} Scholium = Suidas κ 623 καλπάζειν \espTest \textbf{1175} Suidas τ \espTest (Περσικόν) Pollux IV 100 ; cf. Hesychius π 2005 (?) \espTest \textbf{1183} (τὼ πόδε) Hesychius τ 1758 τὼ πόδε· τοὺς πόδας δυϊκῶς \espTest \textbf{1187} (πυγή) Hesychius π 4279 πυγή· τὸ κάθισμα \espTest \textbf{1197} Photius σ 660 συβήνη = Suidas σ 1273 (sed ἐχω οὐδέν pro έκ' ὠδέν, τὸ omisso et συβήνην pro συμβήνην) \espTest (συμβήνην) Hesychius σ 2137 συμβίνης· καπροβόλον. ἐμβόλιον\par\vspace{2.5pt}
{\color{ruleGold}\hrule height 0.35pt}\vspace{2.5pt}
\noindent{\bfseries\scshape\color{rubricRed}Suite (en commentaire) de l'apparat critique et paléographique :} \textbf{v. 1172 :} {\policeGr Σ}\expo{R} Εὐριπίδης ἐν σχήματι προαγωγοῦ γραός \espApp \textbf{ἐμὸν \inRas{γ'}{}} \Rdeux{}. --- \textbf{v. 1174 :} κἀνακόλπασον R (confusion minuscule α → ο) : \lecon{κἀνακάλπασον} Hermann, cf. {\policeGr Σ}, Souda κ 623, Zonaras p. 1171 (dans {\policeGr Σ}\expo{R}, ἁβρῶς R : \lecon{ἄκρως} \lecon{Bibi} ex Zonaras). --- \textbf{v. 1176 :} κωμοτίς ἀν- R (κωμοδρίαν {\policeGr Σ}) : \lecon{κῶμό τις ἀν-} von Velsen. --- \textbf{v. 1190 :} τί οὐκ ἐπιλήσει R : \lecon{οὐκὶ} (Biset) \lecon{πιλῆσι} Porson (\italique{Aristophanica}, p. 222, cf. {\policeGr Σ} οὐχὶ φιλήσειν). --- \textbf{v. 1197 :} έκ' ὠδέν R (ἐχω οὐδέν Souda σ 1273, Photius σ 660) : \lecon{ἔκὠδέν} Bergk \espApp \lecon{συμβήνην} R, défendu par \lecon{Bibi} (ap. Hésychius σ 2137) : συβήνην {\policeGr Σ}, Souda, Photius. --- \textbf{v. 1198 :} γϋ. ante v. 1198 (Boninus) et σκ'. ante ἀκολ. in R \espApp κ\inRas{ω}{ο}μίζεις αὐτοῖς \Rpc{} (κομιεῖς αὐτό Biset ; αὖθις Brunck ; (Τοξότ.) κομίσι σ' αὐτός Enger) : \lecon{(Τοξότ.) κομίζι σ' αὐτίκ'} \lecon{Bibi} \espApp ἀκουλούτι R : \lecon{ἀκολούτι} \Mon{} (ἀκουλούτει Brunck). --- \textbf{v. 1203 :} \lecon{τουτὶ} R (ταυτὶ Dobree, cf. {\policeGr Σ}).\par}
\end{minipage}\end{adjustbox}
\newpage
\renewcommand{\currHeadLE}{SCHOLIES MARGINALES ({\policeGr Σ}\expo{R}) \& GLOSES INTERLINÉAIRES · FOLIO 29/30}
\renewcommand{\currHeadRE}{f. 176v (Classense 178v) · vv. 1169–1210 · {\policeGr Σ} 365 à {\policeGr Σ} 379}
\renewcommand{\currHeadLO}{SCHOLIES MARGINALES ({\policeGr Σ}\expo{R}) \& GLOSES INTERLINÉAIRES · FOLIO 29/30}
\renewcommand{\currHeadRO}{f. 176v (Classense 178v) · vv. 1169–1210 · {\policeGr Σ} 365 à {\policeGr Σ} 379}
\bandeauScholies{Folio 29/30 — f. 176v}{Classense 178v}{vv. 1169–1210}{Verso (page de gauche) · 15 entrées ({\policeGr Σ} 365 à {\policeGr Σ} 379)}
{\fontsize{8.8pt}{10.5pt}\selectfont\noindent\textbf{\color{rubricRed}Disposition topographique des marges sur le f. 176v :} Marge extérieure à gauche ; marge intérieure à droite. — \textit{Fac-similé de référence :} \texttt{folios\_individuels\_ff162v-177r/29\_f176v\_classense178v\_vv1169-1210.jpg}.\par\vspace{3pt}}
\entreeScholie{365}{1170}{τὰ μὲν παρ᾽ ἡμῶν ἴσθι σοι πεπεισμένα·}{Intra-textuel / interlinéaire (f. 176v, en haut du verso, v. 1170)}{Ἀττικῶς ἀντὶ τοῦ πεπείσμεθα.}{Ἀττικ\expo{ῶς} ἀν\auDessus{τ}{ι} τ\expo{οῦ} πεπείσμεθ(α) {\policeFr ÷}}{« Tour attique personnel/neutre (τὰ παρ' ἡμῶν πεπεισμένα) au sens de “sache que de notre côté nous sommes convaincues”. »}
\commScholie{Cité par Louis Hauseux (vv 642-fin-bis.tex, l. 5051).}
\entreeScholie{366}{1172}{ἐμὸν ἔργον ἐστίν· καὶ σὸν, ὦλάφιον, ἅ σοι / καθ᾽ ὁδὸν ἔφραζον,}{Marge gauche (extérieure, f. 176v, deux scholies en face du v. 1172)}{Εὐριπίδης ἐν σχήματι προαγωγοῦ γραός. ἑταίρας ὄνομα Ἐλάφιον ὡς Χρυσίον, [Μύρτι]ον καὶ τὰ ὅμοια.}{Εὐριπίδ(ης) ἐν σχήματ(ι) προαγωγ\expo{οῦ} γραός {\policeFr ÷} // ἑταίρ\expo{ας} ὄνομα Ἐλάφι\expo{ον} ὡς Χρυσί\expo{ον}, [Μύρτι]\expo{ον} \textbf{ς} τὰ ὅμοια {\policeFr ÷}}{« Euripide entre ici sous le déguisement d'une vieille entremetteuse (προαγωγὸς γραῦς, qui se fera appeler Artémisia au v. 1200). — Ἐλάφιον (“Petit Faon, Biche”) est un nom de courtisane (danseuse), formé au neutre diminutif comme Χρυσίον, Μύρτιον et autres noms semblables (cf. Σ 289). »}
\commScholie{Cité par Louis Hauseux (vv 642-fin-bis.tex, l. 5059) : c'est d'après cette scholie (ἐν σχήματι προαγωγοῦ γραός) qu'Euphrosynus Boninus en 1515 a ajouté de sa propre main sur le f. 176v (vv. 1189 et 1198) le sigle γϋ. (= γυνή), comme l'ont démontré Velsen et Martin.}
\entreeScholie{367}{1174–1175}{πρῶτον μὲν οὖν δίελθε κἀνακάλπασον [κἀνακόλπασον R]. / σὺ δ᾽, ὦ Τερηδὸν [Τερηδών R], ἐπαναφύσα Περσικόν·}{Marge gauche (extérieure) et interlinéaire sur Τερηδών (f. 176v, vv. 1174–1175)}{σημειωτέον τοῦτο. σημαίνει δὲ τὸ ἁβρῶς βαδίζειν. ἡ αὐλητρίς.}{σημειωτέ\expo{ον} τοῦτ(ο) · σημαίν(ει) δ(ὲ) τ(ὸ) ἁβρ\expo{ῶς} βαδίζ(ειν) {\policeFr ÷} // ἡ αὐλητρίς {\policeFr ÷}}{« Il faut noter ce mot rare (σημειωτέον τοῦτο, κἀνακάλπασον / κἀνακόλπασον) : il signifie marcher d'un pas souple et gracieux (ἁβρῶς βαδίζειν, ou plutôt ἄκρως βαδίζειν, marcher sur la pointe des pieds au trot). — Τερηδών : le joueur (ou la joueuse) d'aulos. »}
\commScholie{Longuement analysé par Louis Hauseux (Bibi, vv. 1174-1175, vv 642-fin-bis.tex, ll. 5062-5140) : alors que R\expo{1} porte la faute κἀνακόλπασον, Photius α 1514 (ἀνακαλπάζει), la Souda κ 623 et Zonaras p. 1171 (καλπάζειν· τὸ ἄκρως βαδίζειν, où ἄκρως a été altéré en ἁβρῶς dans {\policeGr Σ}\expo{R}) permettent de restituer avec certitude κἀνακάλπασον (« reviens en trottinant sur tes pointes ») !}
\entreeScholie{368}{1175}{σὺ δ᾽, ὦ Τερηδὸν [Τερηδών R], ἐπαναφύσα Περσικόν·}{Marge gauche (extérieure, f. 176v, grande scholie alexandrine de Juba et Séleucos en face du v. 1175) et parepigraphê intra-textuelle αὐλεῖ}{βαρβαρικὸν καὶ Περσικὸν ὄρχημα ὄκλασμα καλεῖται, περὶ οὗ Ἰόβας μακρὸν πεποίηται λόγον ἐν τοῖς περὶ τῆς θεατρικῆς ἱστορίας, ὥστε λελύσθαι τὴν Σελεύκου πρότασιν. προτείνει γὰρ ἐν τῷ πρὸς Ζήνωνα προτατικῷ τὸ Περσικὸν ὄρχημα.}{βαρβαρικ(ὸν) \textbf{ς} Περσικ(ὸν) ὄρχημ(α) ὄκλασμα καλεῖτ(αι), π(ερὶ) οὗ Ἰόβ(ας) μακρ(ὸν) πεποίητ(αι) λόγ\expo{ον} ἐν τ\expo{οῖς} π(ερὶ) τ\expo{ῆς} θεατρικ\expo{ῆς} ἱστορί\expo{ας}, ὥστ(ε) λελύσθ(αι) τ(ὴν) Σελεύκ\expo{ου} πρότασ(ιν) · προτείν(ει) γ(ὰρ) ἐν τῷ πρ(ὸς) Ζήνωνα προτατικῷ τ(ὸ) Περσικ(ὸν) ὄρχημα {\policeFr ÷}}{« On appelle ὄκλασμα (“danse accroupie”) une danse barbare et perse (τὸ Περσικὸν ὄρχημα), au sujet de laquelle le roi Juba II de Maurétanie (FGrHist 275 F 85) a fait un long développement dans son Histoire théâtrale (Περὶ τῆς θεατρικῆς ἱστορίας), de sorte que se trouve résolue la question posée par Séleucos (d'Alexandrie, dit l'Homérique) : celui-ci propose en effet comme problème “Qu'est-ce que la danse perse ?” dans son Recueil de questions adressé à Zénon (ἐν τῷ πρὸς Ζήνωνα προτατικῷ). » — [Et dans la colonne : αὐλεῖ (« on joue de l'aulos »)].}
\commScholie{L'une des scholies les plus célèbres et les plus précieuses pour la datation et les sources du commentaire antique d'Aristophane (citée par L. Hauseux, vv 642-fin-bis.tex, ll. 5141-5155) : en citant l'Histoire théâtrale du roi érudit Juba II de Maurétanie (env. 50 av. J.-C. – 23 ap. J.-C.) et le Προτατικός de Séleucos d'Alexandrie (sous Tibère), le compilateur (Symmaque, fin du Ier s. ap. J.-C.) signe la strate la plus récente du commentaire alexandrin avant sa transmission au Ravennas !}
\entreeScholie{369}{1176}{τί τὸ βόμβο τοῦτο; κῶμό τις ἀνεγεῖρί μοι;}{Marge intérieure droite (variante διχῶς, f. 176v, v. 1176)}{διχῶς, κωμοδρίαν.}{διχ\expo{ῶς}, κωμοδρίαν {\policeFr ÷}}{« Il y a double leçon dans les manuscrits (διχῶς) : κῶμό τις (κωμοτίς R\expo{1}) et κωμοδρίαν. »}
\commScholie{Cité par Louis Hauseux (vv 642-fin-bis.tex, ll. 5157-5160) : second emploi du terme technique d'Aristarque διχῶς sur ce folio.}
\entreeScholie{370}{1180}{ὠς ἐλαπρός, ὤσπερ ψύλλο κατὰ τὸ κῴδιο·}{Interlinéaire (supra lineam, f. 176v, v. 1180)}{κατὰ δέρμα.}{κατὰ δέρμα {\policeFr ÷}}{« À même la peau (ou sur sa tunique qu'elle retire pour danser). »}
\commScholie{Glose interlinéaire de R\expo{2} expliquant le diminutif dorisant/scythique κατὰ τὸ κῴδιο (pour κατὰ τὸ κῴδιον, « sur la petite peau de mouton / la toison », v. 1180) par l'équivalent courant κατὰ δέρμα.}
\entreeScholie{371}{1185}{οἴμ᾽ ὠς στέριπο [ὡς στεριπὸ R, ὡς στέριφα Σ] τὸ τιττί᾽, ὤσπερ γογγυλί [γογγύλη R]·}{Interlinéaire (supra lineam, f. 176v, v. 1185)}{ὡς στέριφα τὰ τιτθία.}{ὡς στέριφα τὰ τιτθία {\policeFr ÷}}{« Quels petits seins fermes et durs (comme des raves) ! »}
\commScholie{Lemme/glose confirmant le texte de R (et du papyrus Ψ\expo{3} = POxy 3840).}
\entreeScholie{372}{1186}{αὔλει σὺ θᾶττον· ἔτι δέδοικας τὸν Σκύθην;}{Interlinéaire (supra lineam, f. 176v, v. 1186)}{πρὸς τὴν αὐλητρίδα.}{πρ(ὸς) τ(ὴν) αὐλητρίδα {\policeFr ÷}}{« Dit en s'adressant au joueur (ou à la joueuse) d'aulos (Térédon, pour la première partie du vers) et à la danseuse Élaphion (pour la seconde). »}
\commScholie{Note de jeu scénique de R\expo{2} précisant le changement d'interlocuteur au v. 1186 : après avoir fait rasseoir l'Archer, Euripide (déguisé en vieille entremetteuse) se tourne vers la joueuse de flûte Terédon (πρὸς τὴν αὐλητρίδα) pour presser la cadence.}
\entreeScholie{373}{1187}{καλό γε τὸ πυγή· κλαῦσί γ᾽ ἢν [ἂν R] μὴ ᾽νδον μένῃς· — [παρεπιγραφή R\expo{1}, \PsiTrois{} :] ἀνακύπτη· καὶ παρακύπτι· ἀπεψωλημένος·}{Marge gauche (extérieure, f. 176v, Σ 1187) et parepigraphê scénique après le v. 1187 (Planche 19)}{πρὸς τὸ πέος λέγει ἐπιπλήττων ἐντεινομένῳ.}{πρ(ὸς) τ(ὸ) πέ\expo{ος} λέγ(ει) ἐπιπλήττ\expo{ων} ἐντεινομένῳ {\policeFr ÷} // [ἀνακύπτ(ει) \textbf{ς} παρακύπτ(ει) ἀπεψωλημέν\expo{ος}]}{« κλαύσι γ', ἂν μὴ 'νδον μένῃς (“Tu vas pleurer si tu ne restes pas dedans !”) : le Scythe dit cela en s'adressant à son propre membre viril (πρὸς τὸ πέος λέγει), qu'il réprimande parce qu'il se dresse en érection sous sa tunique ! — [Didascalie scénique : Il se redresse et se penche de côté, le gland décalotté en érection]. »}
\commScholie{Cité par Louis Hauseux (vv 642-fin-bis.tex, ll. 5203-5218) et étudié sur la Planche 19 du Manuel : le papyrus du IIe s. ap. J.-C. Ψ\expo{3} (POxy 3840) conserve lui aussi une trace de parepigraphê après le v. 1187 (]νο[), prouvant l'origine alexandrine de cette quatrième didascalie scénique !}
\entreeScholie{374}{1190}{οὐκὶ πιλῆσι πρῶτά με;}{Interlinéaire (supra lineam au-dessus de τί οὐκ ἐπιλήσει de R\expo{1}, f. 176v, v. 1190)}{οὐχὶ φιλήσειν.}{οὐχὶ φιλήσ(ει) {\policeFr ÷}}{« οὐκὶ πιλῆσι (prononcé par le Scythe sans aspiration) : au sens de οὐχὶ φιλήσει (“elle ne va pas m'embrasser d'abord ?”). »}
\commScholie{Encore un cas éclatant (souligné par Biset, Porson et Louis Hauseux, Bibi, v. 1190, vv 642-fin-bis.tex, l. 5230) où la glose de R\expo{2} (οὐχὶ φιλήσει) a conservé la vraie leçon d'Aristophane (οὐκὶ πιλῆσι πρῶτά με;) alors que R\expo{1} l'avait corrompue par mécoupure et iotacisme en τί οὐκ ἐπιλήσει !}
\entreeScholie{375}{1197}{ἀλλ᾽ οὐκ ἔκὠδέν· ἀλλὰ τὸ συμβήνην [συβήνην Σ, Suid.] λαβέ·}{Marge gauche (extérieure, f. 176v, deux scholies en face du v. 1197)}{τὴν τοξοθήκην. συβήνη αὐλοθήκη. λέγουσι δὲ καὶ τὸν φαρετρεῶνα συβήνην.}{τ(ὴν) τοξοθήκ(ην) · συβήν(η) αὐλοθήκ(η) · λέγουσ(ι) δ(ὲ) \textbf{ς} τ(ὸν) φαρετρεῶνα συβήν(ην) {\policeFr ÷}}{« τὴν τοξοθήκην : l'étui à arc (le carquois scythe, gorytos). Au sens propre, la συβήνη (en peau de truie, σῦς) est l'étui de cuir pour les flûtes (αὐλοθήκη) ; mais on appelle aussi συβήνη le carquois à flèches (φαρετρεῶνα). »}
\commScholie{Longuement discuté par Louis Hauseux (vv 642-fin-bis.tex, ll. 5289-5325) en regard d'Hésychius σ 2136, Photius σ 660 et la Souda σ 1273 : l'archer scythe qui n'a pas une obole en poche remet en gage à la vieille entremetteuse son carquois de cuir (συβήνη).}
\entreeScholie{376}{1199}{σὺ δὲ τοῦτο τήρει τὴ γέροντο, γρᾴδιο·}{Marge gauche (extérieure, f. 176v, en face des vv. 1199–1201)}{παρακελεύεται αὐτῷ τηρεῖν τὸν γέροντα, αὐτὸς δὲ εἰσέρχεται μετὰ τῆς ὀρχηστρίδος.}{παρακελεύετ(αι) αὐτῷ τηρ(εῖν) τ(ὸν) γέροντ(α), αὐτ\expo{ὸς} δ(ὲ) εἰσέρχετ(αι) μετὰ τ\expo{ῆς} ὀρχηστρίδ\expo{ος} {\policeFr ÷}}{« Le Scythe recommande à la vieille (c'est-à-dire à Euripide lui-même !) de surveiller le vieux prisonnier (Mnésiloque), tandis que lui-même sort de scène avec la petite danseuse (pour aller coucher avec elle sur sa natte). »}
\commScholie{Souligne le ressort comique suprême du dénouement : le gardien confie la garde du prisonnier à son complice venu le délivrer !}
\entreeScholie{377}{1202}{ὦ Ἑρμῆ δόλιε, ταυτὶ μὲν ἔτι καλῶς ποιεῖς.}{Marge gauche (extérieure, f. 176v, en face du v. 1202)}{ἐπειδὴ τῶν δόλων ἔφορος ὁ Ἑρμῆς.}{ἐπειδὴ τ\expo{ῶν} δόλ\expo{ων} ἔφορ\expo{ος} ὁ Ἑρμῆς {\policeFr ÷}}{« Il invoque Hermès Dolios (“le Rusé”), parce qu'Hermès est le dieu tutélaire des ruses et des tromperies heureuses. »}
\commScholie{Repris par la Souda ε 3036 (le papyrus Ψ\expo{4} = POxy 4935 + 5132 recoupe les vv. 1201-1210).}
\entreeScholie{378}{1203}{σὺ μὲν οὖν ἀπότρεχε, παιδάριον, τουτὶ λαβών·}{Marge gauche (extérieure, f. 176v, en face du v. 1203)}{τὴν τοξοθήκην καὶ τὰ ὄργανα τῆς ὀρχηστρίδος.}{τ(ὴν) τοξοθήκ(ην) \textbf{ς} τὰ ὄργανα τ\expo{ῆς} ὀρχηστρίδ\expo{ος} {\policeFr ÷}}{« τουτί (ou ταυτί) : le carquois du Scythe (que celui-ci vient de laisser en gage) ainsi que les instruments de la danseuse. »}
\commScholie{Cité par Louis Hauseux (vv 642-fin-bis.tex, l. 5381) : Euripide renvoie d'abord le jeune joueur d'aulos (παιδάριον, Térédon) en lui confiant le carquois confisqué au Scythe.}
\entreeScholie{379}{1205}{σὺ δ᾽ ὅπως ἀνδρικῶς / ὅταν λυθῇς τάχιστα φεύξει καὶ τενεῖς}{Interlinéaire (supra lineam au-dessus de θεύσει, f. 176v, v. 1205)}{δραμῇ.}{δραμῇ {\policeFr ÷}}{« θεύσει (futur de θέω) : tu courras à toutes jambes (δραμῇ). »}
\commScholie{Glose interlinéaire de R\expo{2} expliquant l'emploi intransitif familier du futur τενεῖς (de τείνω, « tendre le pas, filer à toute allure », associé à φεύξει au v. 1205) par le futur moyen δραμῇ (« tu courras »).}
\clearpage
\renewcommand{\currHeadLE}{ΑΡΙΣΤΟΦΑΝΟΥΣ ΘΕΣΜΟΦΟΡΙΑΖΟΥΣΑΙ · CODEX RAVENNAS 429 (\R) · FOLIO 30/30}
\renewcommand{\currHeadRE}{f. 177r (Classense 179r) · vv. 1211–1231 · Collation diplomatique}
\renewcommand{\currHeadLO}{ΑΡΙΣΤΟΦΑΝΟΥΣ ΘΕΣΜΟΦΟΡΙΑΖΟΥΣΑΙ · CODEX RAVENNAS 429 (\R) · FOLIO 30/30}
\renewcommand{\currHeadRO}{f. 177r (Classense 179r) · vv. 1211–1231 · Collation diplomatique}
\noindent\begin{adjustbox}{max totalheight=30.4cm, max width=\textwidth, keepaspectratio}\begin{minipage}{\textwidth}
\bandeauRav{f. 177r}{Classense 179r}{1211–1231}{Folio 30/30 · Recto (page de droite) · 7 scholies \& gloses ({\policeGr Σ} 380–386)}
{\fontsize{9.3pt}{10.8pt}\selectfont
\ligneRav{1211}{τοξ/}{κοὐ \difference{δύσκολ᾽} ἀλλὰ πρᾶο. ποῦ τὸ γρᾴδιο;}{\textbf{1211} δύσκολλ᾽ R : \lecon{δύσκολ᾽} Grynaeus}
\ligneRav{}{}{οἴμ᾽ ὠς ἀπόλωλο· ποῦ τὸ γέροντ᾽ ἐντευτενί;}{}
\ligneRav{1213}{}{ὦ γρᾴδι᾽, ὦ \difference{γρᾶ᾽}. οὐκ ἐπαινῶ γρᾴδιο·}{\textbf{1213} γρᾶίο R : \lecon{γρᾶ᾽} Bentley}
\ligneRav{}{}{Ἀρταμουξία·}{}
\ligneRav{1214}{}{διέβαλλέ μ᾽ \difference{ἡ γραῦς}. ἀπότρεκ᾽ ὠς τάκιστα σύ·}{\textbf{1214} \lecon{διέβαλλε} R : διέβαλέ Σ, Suid. δ 892, Eust. p. 1406.19 \espApp ὁ γραῦς R : \lecon{ἡ γραῦς} Suid., Eust. \espApp ὡς τάχιστα R : ὠς τάκιστα Brunck}
\ligneRav{1215}{}{ὀρτῶς δὲ \difference{συμβήνη} 'στί· καταβηνῆσι γάρ·}{\textbf{1215} συβήνη R : \lecon{συμβήνη} Bibi, cf. v.1197, Hesych. σ 2137}
\ligneRav{1216}{}{οἴμοι, τί \difference{δρᾶσι}; ποῖ τὸ \difference{γρᾴδιο};}{\textbf{1216} δράσει R : \lecon{δρᾶσι} Blaydes (1842) \espApp γράδιο R : \lecon{γρᾴδιο} \Mon}
\ligneRav{}{}{Ἀρταμουξία·}{}
\ligneRav{1217}{χορ/}{τὴν \difference{γραῦν} ἐρωτᾷς, ἣ 'φερεν τὰς πηκτίδας;}{\textbf{1217} γραῦ R : \lecon{γραῦν} Suid. π 1502}
\ligneRav{}{τοξ/}{ναὶ ναῖκι. εἶδες αὐτό; : ταύτῃ γ᾽ οἴχεται}{}
\ligneRav{1219}{}{αὐτή τ᾽ ἐκείνη καὶ \difference{γέρων} τις εἵπετο·}{\textbf{1219} γέρον R : \lecon{γέρων} \Mon}
\ligneRav{1220}{τοξ/}{κροκῶτ᾽ ἔκοντο τὴ γέροντο; : φήμ᾽ ἐγώ·}{}
\ligneRav{}{}{ἔτ᾽ ἂν καταλάβοις, εἰ διώκοις ταυτῃί·}{}
\ligneRav{1222}{τοξ/}{ὦ μιαρὸ \difference{γρᾷο}· πότερα τρέξι τὴν \difference{ὀδό};}{\textbf{1222} γραῦ et ὁδό R : \lecon{γρᾷο} et \lecon{ὀδό} Brunck}
\ligneRav{}{}{Ἀρταμουξία·}{}
\ligneRav{}{χορ/}{ὀρθὴν ἄνω δίωκε. ποῖ θεῖς; οὐ πάλιν}{}
\ligneRav{1224}{}{τῃδὶ \difference{διώξει}; \difference{᾽ς} τοὔμπαλιν τρέχεις σύ γε·}{\textbf{1224} δ᾽ ὤξεις R : \lecon{διώξει} Elmsley (ad \italique{Ach.} 278) \lecon{᾽ς} Cobet (\italique{Mnemos.} 2 (1853), p.213)}
\ligneRav{1225}{τοξ/}{κακόδαιμον, ἀλλὰ τρέξι. Ἀρταμουξία·}{}
\ligneRav{1226}{χορ/}{\hspace{1.1em}τρέχε νυν \difference{κατ᾽ αὐτοὺς} \difference{ἐς} κόρακας ἐπουρίσας·}{\textbf{1226} κατὰ τοὺς R : \lecon{κατ᾽ αὐτοὺς} Jackson (p.82) \lecon{ἐς} O. Bachmann}
\ligneRav{1227}{}{\hspace{1.1em}ἀλλὰ \difference{πέπαισται} μετρίως ἡμῖν·}{\textbf{1227} πέπυσται R : \lecon{πέπαισται} Σ}
\ligneRav{}{}{\hspace{1.1em}ὥσθ᾽ ὥρα δή 'στι βαδίζειν}{}
\ligneRav{}{}{\hspace{1.1em}οἴκαδ᾽ ἑκάστῃ. τὼ Θεσμοφόρω δ᾽}{}
\ligneRav{1230}{}{\hspace{1.1em}ἡμῖν ἀγαθὴν}{}
\ligneRav{1231}{}{\hspace{1.1em}τούτων χάριν \difference{ἀνταποδοίτην} (ἀνταδοῖτον R)·}{\textbf{1231} ἀνταδοῖτον R : ἀνταποδοῖτον Grynaeus : \lecon{ἀνταποδοίτην} Bibi}
}
\vspace{3pt}{\color{ruleGold}\hrule height 0.6pt}\vspace{3pt}
{\fontsize{8.0pt}{9.4pt}\selectfont
\noindent{\bfseries\scshape\color{rubricRed}Loci paralleli et testimonia (f. 177r, vv. 1211–1231) :} \textbf{1214} Suidas δ 892 διέβαλέ μ' ἡ γραῦς ; \Eustathe 1406.19 \espTest \textbf{1217} Suidas π 1502 \espTest \textbf{1231} Subscriptio \PsiUn \ fr. 6 ΑΡΙΣΤΟΦΑΝ]ΟΥΣ ΘΕΣΜΟΦΟΡΙ]ΑΖΥΣΑΙ ]ΧΗΗΙ (= 1201, versuum numerus)\par\vspace{2.5pt}
{\color{ruleGold}\hrule height 0.35pt}\vspace{2.5pt}
\noindent{\bfseries\scshape\color{rubricRed}Suite (en commentaire) de l'apparat critique et paléographique :} \textbf{Témoin papyrologique :} \textbf{\PsiUn{} (\italique{PSI} XI 1194, fr. 6)} conserve après le \textbf{v. 1231} la \italique{subscriptio} stichométrique antique en chiffres acrophoniques attiques : ΑΡΙΣΤΟΦΑΝ]ΟΥΣ ΘΕΣΜΟΦΟΡΙ]ΑΖΥΣΑΙ ]ΧΗΗΙ (= \textbf{1201 vers}). --- \textbf{Harmonisation de συμβήνη au v. 1215 par Bibi :} Là où R avait conservé συμβήνην avec -μ- au v. 1197 mais écrit συβήνη sans -μ- au v. 1215, \textbf{Bibi rétablit συμβήνη} au v. 1215 pour maintenir le jeu de mots avec συμβινεῖν / καταβηνῆσι (cf. Hésychius σ 2137). --- \textbf{Notes marginales de locuteur au v. 1219 :} R porte Σκ. en marge devant le v. 1219 et Χορ. devant καὶ γέρον à l'intérieur du vers. --- \textbf{Faute de minuscule au v. 1227 :} \Run{} écrit πέπυσται au lieu de \lecon{πέπαισται} (« nous avons assez joué », conservé par {\policeGr Σ}\expo{R}). --- \textbf{Haplographie syllabique au dernier vers (v. 1231) :} \Run{} écrit ἀνταδοῖτον (saut de la syllabe -πο-) au lieu de ἀνταποδοῖτον (Grynaeus) / \lecon{ἀνταποδοίτην} (optatif duel adopté dans le texte par L. Hauseux). --- \textbf{v. 1214 :} \lecon{διέβαλλε} R (διέβαλέ {\policeGr Σ}, Souda δ 892, Eustathe p. 1406.19) \espApp ὁ γραῦς R : \lecon{ἡ γραῦς} Souda, Eustathe \espApp ὡς τάχιστα R : \lecon{ὠς τάκιστα} Brunck (adopté par L. Hauseux). --- \textbf{v. 1215 :} συβήνη R : \lecon{συμβήνη} Bibi (cf. v. 1197, Hésychius σ 2137). --- \textbf{v. 1219 :} Σκ. ante versum in mg. et Χορ. ante καὶ γέρον in R \espApp γέρον R : \lecon{γέρων} \Mon{}. --- \textbf{v. 1227 :} πέπυσται R : \lecon{πέπαισται} {\policeGr Σ}.\par}
\end{minipage}\end{adjustbox}
\newpage
\renewcommand{\currHeadLE}{SCHOLIES MARGINALES ({\policeGr Σ}\expo{R}) \& GLOSES INTERLINÉAIRES · FOLIO 30/30}
\renewcommand{\currHeadRE}{f. 177r (Classense 179r) · vv. 1211–1231 · {\policeGr Σ} 380 à {\policeGr Σ} 386}
\renewcommand{\currHeadLO}{SCHOLIES MARGINALES ({\policeGr Σ}\expo{R}) \& GLOSES INTERLINÉAIRES · FOLIO 30/30}
\renewcommand{\currHeadRO}{f. 177r (Classense 179r) · vv. 1211–1231 · {\policeGr Σ} 380 à {\policeGr Σ} 386}
\bandeauScholies{Folio 30/30 — f. 177r}{Classense 179r}{vv. 1211–1231}{Recto (page de droite) · 7 entrées ({\policeGr Σ} 380 à {\policeGr Σ} 386)}
{\fontsize{8.8pt}{10.5pt}\selectfont\noindent\textbf{\color{rubricRed}Disposition topographique des marges sur le f. 177r :} Marge extérieure à droite ; marge intérieure à gauche. — \textit{Fac-similé de référence :} \texttt{folios\_individuels\_ff162v-177r/30\_f177r\_classense179r\_vv1211-1231\_fin.jpg}.\par\vspace{3pt}}
\entreeScholie{380}{1214}{διέβαλλέ μ᾽ ἡ γραῦς. ἀπότρεκ᾽ ὡς τάκιστα σύ·}{Marge supérieure et marge droite (extérieure, f. 177r, en haut du dernier folio)}{ἐξηπάτησε. παρὰ τοῖς Ἴωσιν. ὁ αὐτὸς ἐν Ὄρνισιν [1648] «διαβάλλεταί ς’ ὁ θεῖος.» ἑαυτῷ παρακελεύεται.}{ἐξηπάτησ(ε) · παρὰ τ\expo{οῖς} Ἴωσ(ιν) · ὁ αὐτ\expo{ὸς} ἐν Ὄρνισ(ιν) «διαβάλλεταί σ' ὁ θεῖ\expo{ος}» {\policeFr ÷} // ἑαυτῷ παρακελεύετ(αι) {\policeFr ÷}}{« διέβαλλέ με (ou διέβαλέ με) : “elle m'a trompé, elle m'a joué” (ἐξηπάτησε) ; emploi ionien et attique familier, comme chez Aristophane lui-même dans les Oiseaux (v. 1648) : “Ton oncle te roule (διαβάλλεταί σ' ὁ θεῖος) !” — ἀπότρεκ' ὡς τάκιστα σύ : le Scythe s'exhorte lui-même à courir au plus vite. »}
\commScholie{Cité par Louis Hauseux (vv 642-fin-bis.tex, l. 5414) et repris par la Souda δ 892 et Eustathe 1406, 19 : remarque lexicologique d'Aristarque sur le sens de διαβάλλω = ἐξαπατάω (« duper ») avec renvoi exact aux Oiseaux v. 1648.}
\entreeScholie{381}{1216–1217}{οἴμοι, τί δρᾶσι [δράσει R]; ποῖ τὸ γρᾴδιο [γράδιο R]; / Ἀρταμουξία· — τὴν γραῦν ἐρωτᾷς, ἣ ᾽φερεν τὰς πηκτίδας;}{Marge droite (extérieure, f. 177r, deux scholies en face des vv. 1217–1218)}{ὁ χορὸς ἐμπαίζει τῷ Σκύθῃ. ἵνα μὴ διώκων καταλάβῃ αὐτούς· συνέθετο γὰρ Εὐριπίδης ἐλθὼν φιλίαν πρὸς αὐτάς. πηκτὶς ὄργανον κιθαρῳδικόν. ἔφερεν δὲ τοῦτο· ἦν γὰρ μετ’ αὐτῆς ἡ αὐλητρίς.}{ὁ χορ(ὸς) ἐμπαίζ(ει) τῷ Σκύθ(ῃ) · ἵνα μὴ διώκ\expo{ων} καταλάβ(ῃ) αὐτ\expo{ούς} · συνέθετ(ο) γ(ὰρ) Εὐριπίδ(ης) ἐλθ(ὼν) φιλί(αν) πρ(ὸς) αὐτάς {\policeFr ÷} // πηκτὶς ὄργαν\expo{ον} κιθαρῳδικόν · ἔφερ(εν) δ(ὲ) τοῦτ(ο) · ἦν γ(ὰρ) μετ' αὐτ\expo{ῆς} ἡ αὐλητρίς {\policeFr ÷}}{« Le Chœur se joue du Scythe (en l'envoyant dans la mauvaise direction), afin qu'en poursuivant les fugitifs il ne les rattrape pas : car Euripide était venu conclure un pacte d'amitié avec les femmes (aux vv. 1160-1171). — σύβινα / πηκτίς (v. 1217) : la πηκτίς est un instrument de musique à cordes (harpe lydienne) que portait la vieille, puisqu'elle était accompagnée de la musicienne. »}
\commScholie{Explique le jeu de scène final de la Coryphée qui égare l'Archer scythe.}
\entreeScholie{382}{1222}{πότερα τρέξι τὴν ὁδό;}{Interlinéaire (supra lineam, f. 177r, v. 1222)}{ποίαν ὁδὸν, φησὶ, περιπατήσω.}{ποί(αν) ὁδ(ὸν), φησ(ὶ), περιπατήσω {\policeFr ÷}}{« “Par quel chemin”, dit le Scythe, “dois-je courir ?” »}
\commScholie{Glose grammaticale de R\expo{2} rétablissant en grec attique (ποίαν ὁδόν, φησί, περιπατήσω, « par quel chemin dois-je aller ? ») la question prononcée en sabir barbare par l'Archer scythe (πότερα τρέξι τὴν ὁδό, v. 1222).}
\entreeScholie{383}{1224}{ὀρθὴν ἄνω δίωκε. — ποῖ θεῖς; οὐ πάλιν / τῃδὶ διώξει; ᾽ς τοὔμπαλιν τρέχεις σύ γε·}{Marge droite (extérieure, f. 177r, en face des vv. 1223–1224)}{ὡς αὐτοῦ τὴν ὁδὸν ἐκείνην θέλοντος ἀπελθεῖν. οἵαν ἀπῆλθον οἱ περὶ τὸν Εὐριπίδην.}{ὡς αὐτ\expo{οῦ} τ(ὴν) ὁδ(ὸν) ἐκείν(ην) θέλοντ\expo{ος} ἀπελθ(εῖν) · οἵ(αν) ἀπῆλθ\expo{ον} οἱ π(ερὶ) τ(ὸν) Εὐριπίδ(ην) {\policeFr ÷}}{« Comme le Scythe fait mine de partir du bon côté (celui par où Euripide et Mnésiloque se sont réellement enfuis), la Coryphée le rappelle et le lance dans la direction opposée ! »}
\commScholie{Indication de jeu scénique très vivante sur le lazzi final de poursuite.}
\entreeScholie{384}{1226}{τρέχε νυν κατ᾽ αὐτοὺς ἐς κόρακας ἐπουρίσας.}{Interlinéaire (supra lineam au-dessus d'ἐπουρίσας, f. 177r, v. 1226)}{ἐπιδραμών.}{ἐπιδραμ(ών) {\policeFr ÷}}{« ἐπουρίσας : en filant bon train vent arrière (ἐπιδραμών). »}
\commScholie{Glose de la métaphore nautique ἐπουρίζω (de οὖρος, « vent favorable »).}
\entreeScholie{385}{1227}{ἀλλὰ πέπαισται μετρίως ἡμῖν·}{Interlinéaire (supra lineam au-dessus de πέπυσται de R\expo{1}, f. 177r, v. 1227)}{πέπαικται.}{πέπαικτ(αι) {\policeFr ÷}}{« πέπαισται (en attique) : “nous avons suffisamment joué et plaisanté” (πέπαικται dans la koinè, cf. Σ 947). »}
\commScholie{Dernière preuve éclatante de la valeur critique de {\policeGr Σ}\expo{R} (Louis Hauseux, Bibi, v. 1227, vv 642-fin-bis.tex, ll. 5468-5469) : alors que R\expo{1} avait défiguré par iotacisme le verbe en πέπυσται (« on s'est informé »), la glose de R\expo{2} πέπαικται restitue à coup sûr le parfait passif attique πέπαισται (« nous avons assez joué ») !}
\entreeScholie{386}{1230}{τὼ Θεσμοφόρω δ᾽ / ἡμῖν ἀγαθὴν / τούτων χάριν ἀνταποδοίτην [ἀνταδοῖτον R]·}{Interlinéaire (supra lineam au-dessus de τὼ Θεσμοφόρω, f. 177r, v. 1230, avant la korônis finale)}{ἡ Δημήτηρ καὶ ἡ Περσεφόνη.}{ἡ Δημήτ(ηρ) \textbf{ς} ἡ Περσεφόν(η) {\policeFr ÷}}{« τὼ Θεσμοφόρω (au duel) : Déméter et Perséphone. »}
\commScholie{Dernière glose de R\expo{2} sur Les Thesmophories, immédiatement suivie dans la colonne de la korônis finale (\textbf{>—}) et, sur la même page du f. 177r, de l'Hypothesis de L'Assemblée des femmes.}
\end{document}
